\documentclass[lang=cn,11pt]{elegantbook}

\title{概率论与数理统计}
\subtitle{学习笔记}

\author{夏同}
\institute{华南理工大学}
\date{\today}
\version{1.0}
\bioinfo{自定义}{信息}

\extrainfo{居敬持志，守正出奇。—— 张红兵}

\setcounter{tocdepth}{3}

\logo{logo-blue.png}
\cover{751_01.png}

% 本文档命令
\usepackage{array}

\makeatletter
\@ifundefined{setCJKfamilyfont}{}{%
  \IfFontExistsTF{KaiTi}{%
    \setCJKfamilyfont{codexkaiti}[AutoFakeBold=2.5]{KaiTi}%
    \renewcommand{\kaishu}{\CJKfamily{codexkaiti}}%
  }{}%
  \IfFontExistsTF{FangSong}{%
    \setCJKfamilyfont{codexfangsong}[AutoFakeBold=2.5]{FangSong}%
    \renewcommand{\fangsong}{\CJKfamily{codexfangsong}}%
  }{}%
}
\makeatother

\newcommand{\ccr}[1]{\makecell{{\color{#1}\rule{1cm}{1cm}}}}

% 修改标题页的橙色带
% \definecolor{customcolor}{RGB}{32,178,170}
% \colorlet{coverlinecolor}{customcolor}

\begin{document}

\maketitle
\frontmatter

\tableofcontents
\makeenvindex

\mainmatter

\chapter{随机事件与概率}
{\scriptsize
\[
\begin{cases}
\text{基本概念} & 
\begin{cases}
\text{随机试验} \ E: \text{可重复性、结果不唯一、所有结果已知、不确定性} \\
\text{样本空间} \ \Omega: \text{所有可能结果构成的集合} \\
\text{随机事件} \ A: \text{样本空间的子集} 
\begin{cases}
\text{关系：包含、相等、互斥、对立} \\
\text{表示：和事件、积事件、差事件、对立事件} \\
\text{运算：交换律、结合律、分配律、德摩根律}
\end{cases}
\end{cases} \\
\text{概率概念} & 
\begin{cases}
\text{统计概念：频率} \\
\text{公理化概念：非负性、规范性和可列可加性} \\
\text{古典模型(等可能概型)：有限元素、可能性相同} \\
\text{计算公式}
\begin{cases}
\text{加法公式}
\begin{cases}
P(A \cup B) = P(A) + P(B) - P(AB) \\
P(A \cup B) = P(A) + P(B) \ (\text{若事件} \ A \ \text{与} \ B \ \text{互斥})
\end{cases} \\
\text{减法公式}
\begin{cases}
P(A - B) = P(A) - P(AB) = P(A\overline{B}) \\
P(A - B) = P(A) - P(B) \ (\text{若事件} \ A \supset B)
\end{cases} \\
\text{求逆公式：} \ P(\overline{A}) = 1 - P(A)
\end{cases}
\end{cases} \\
\text{条件概率概念} & 
\begin{cases}
 \text{定义：当} \ P(A) > 0 \ \text{时}, \ P(B|A) = \dfrac{P(AB)}{P(A)} \\
 \text{且满足概率性质与上述公式} \\
 \quad (\text{非负性、规范性和可列可加性}) \\
 \text{计算：定义法、缩小样本空间法} \\
 \text{应用}
 \begin{cases}
 \text{乘法公式：当} \ P(A) > 0 \ \text{时}, \ P(AB) = P(B|A)P(A) \\
 \text{全概率公式(由因求果)：} \\
 \qquad P(A) = \sum_{i=1}^{n} P(A|B_i)P(B_i) \\
 \text{贝叶斯公式(执果寻因)：} \\
 \qquad P(B_i|A) = \dfrac{P(A|B_i)P(B_i)}{\sum_{j=1}^{n} P(A|B_j)P(B_j)} \\
 \quad (\text{其中} \ B_1,B_2,\dots,B_n \ \text{是一组完备事件组}, \ P(B_i) > 0)
 \end{cases}
\end{cases} \\
\text{独立性} & 
\begin{cases}
 \text{定义}
 \begin{cases}
 \text{两个事件:} \ P(AB) = P(A)P(B) \\
 \text{等价地，当} \ 0 < P(A) < 1 \ \text{时}, \ P(B|A) = P(B|\overline{A}) \\
 \text{多个事件：相互独立} \Rightarrow \text{两两独立}, \\
 \text{但两两独立} \nRightarrow \text{相互独立}
 \end{cases} \\
\text{性质}
\begin{cases}
\text{相互独立的事件中部分事件间相互独立} \\
\text{相互独立的事件中部分事件的逆事件与其余事件仍相互独立}
\end{cases} \\
 \text{应用}
 \begin{cases}
 \text{计算多个事件的和事件概率转化为逆事件概率乘积} \\
 \text{伯努利试验：只有两个结果} \ A \ \text{与} \ \overline{A}, \\
 n \ \text{次试验相互独立}
 \end{cases}
\end{cases}
\end{cases}
\]
}
\newpage

\begin{introduction}
  \item 随机试验与样本空间~\ref{def:random-experiment}
  \item 随机事件及其运算~\ref{def:random-event}
  \item 完备事件组~\ref{def:complete-event-group}
  \item 概率的公理化定义~\ref{def:probability-axiom}
  \item 概率的基本性质~\ref{prop:probability-basic}
  \item 古典概型~\ref{def:classical-prob}
  \item 几何概型~\ref{def:geometric-prob}
  \item 条件概率~\ref{def:conditional-prob}
  \item 乘法公式~\ref{thm:multiplication}
  \item 全概率公式与贝叶斯公式~\ref{thm:total-probability}
  \item 事件的独立性~\ref{def:independence}
  \item 伯努利概型~\ref{def:bernoulli}
\end{introduction}

\section{基本概念}

\begin{definition}[随机试验] \label{def:random-experiment}
如果一个试验满足以下三个条件：
\begin{enumerate}
    \item 试验可以在相同的条件下重复进行；
    \item 试验所有可能结果明确可知，且不止一个；
    \item 每一次试验会出现哪一个结果，事先并不能确定，
\end{enumerate}
则称此试验为\textbf{随机试验}。随机试验简称试验，用字母 \(E\) 或 \(E_1, E_2, \dots\) 表示。
\end{definition}

我们是通过研究随机试验来研究随机现象的。

\begin{note}
在不少情况下，不能确切知道某一随机试验的全部可能结果，但可以知道它不超出某个范围。这时，也可以用这个范围来作为该试验的全部可能结果的集合。例如，需要记录某个城市一天的交通事故数量，则试验结果将是非负整数 \(x\)。无法确定 \(x\) 的可能取值的确切范围，但可以把这个范围取为 \([0, +\infty)\)，它总能包含一切可能的试验结果，尽管明知某些结果，如 \(x > 10000\) 是不会出现的。甚至把这个范围取为 \((-\infty, +\infty)\) 也无妨。这里体现了一定的数学抽象，它可以带来很大的方便。
\end{note}

\begin{definition}[样本空间] \label{def:sample-space}
随机试验的每一个可能结果称为\textbf{样本点}，通常记为 \(\omega\)。全体样本点构成的集合称为\textbf{样本空间}（或基本事件空间），记为 \(\Omega\)（或 \(S\)），可写作
\[
\Omega=\{\omega:\omega \text{ 是该随机试验的一个可能结果}\}.
\]
\end{definition}

由一个样本点构成的事件称为\textbf{基本事件}。

\begin{definition}[随机事件] \label{def:random-event}
在一次试验中可能出现，也可能不出现的结果称为\textbf{随机事件}，简称事件，用大写字母 \(A, B, C\) 等表示。随机事件是样本空间 \(\Omega\) 的子集。在每次试验中，当且仅当这一子集中的一个样本点出现时，称这一事件发生。
\end{definition}

\begin{enumerate}
    \item 对于仅由一个基本结果构成的单点集，称为\textbf{基本事件}；
    \item 将每次试验中一定发生的事件称为\textbf{必然事件}，记为 \(\Omega\)；
    \item 每次试验中一定不发生的事件称为\textbf{不可能事件}，记为 \(\varnothing\)。
\end{enumerate}

\begin{note}
概念之间的联系：随机试验产生基本结果，基本结果 \(\in\) 随机事件 \(\subset\) 样本空间。随机事件在一次试验中是否发生虽然不能事先确定，但是在大量重复试验的情况下，它的发生呈现出一定的规律性，这门课程正是要研究这种规律性。
\end{note}

\section{事件的关系与运算}

事件是一个集合，因此事件之间的关系和运算与集合论中集合之间的关系和运算完全类似，可以利用集合论的知识来理解事件的关系与运算。

\subsection{事件的关系}

\begin{enumerate}
    \item \textbf{包含}——如果事件 \(A\) 发生必导致事件 \(B\) 发生，则称事件 \(B\) 包含事件 \(A\)（或 \(A\) 被 \(B\) 包含），记为 \(A \subset B\)。也称 \(A\) 是 \(B\) 的子事件。
    \item \textbf{相等}——如果 \(A \subset B\) 且 \(B \subset A\)，则称事件 \(A\) 与 \(B\) 相等，记为 \(A = B\)。\(A\) 与 \(B\) 相等，事实上也就是说，\(A\) 与 \(B\) 由一些完全相同的试验结果构成，它不过是同一事件表面上看起来不同的两个说法而已。
    \item \textbf{相容}——若 \(AB \neq \varnothing\)，则称事件 \(A\) 和 \(B\) 相容。
    \item \textbf{互斥（互不相容）}——若 \(AB = \varnothing\)，则称事件 \(A\) 与 \(B\) 互不相容，也叫互斥。如果一些事件中任意两个事件都互斥，则称这些事件是两两互斥的。任意一对基本事件都是互斥的。
    \item \textbf{对立}——事件 \(A, B\) 不能同时发生，且至少一个发生，称事件 \(A, B\) 互为对立事件，或互为逆事件，即 \(A \cap B = \varnothing\)，\(A \cup B = \Omega\)。
\end{enumerate}

\begin{definition}[完备事件组] \label{def:complete-event-group}
如果 \(\bigcup_{i=1}^{n} A_i = \Omega\)（或 \(\bigcup_{i=1}^{\infty} A_i = \Omega\)），且 \(A_i A_j = \varnothing\)（对一切 \(i \neq j\)；\(i, j = 1, 2, \dots, n(\dots)\)），称有限个（或可列个）事件 \(A_1, A_2, \dots, A_n(\dots)\) 构成一个\textbf{完备事件组}。
\end{definition}

完备事件组在概率论中具有重要作用，全概率公式和贝叶斯公式都需要完备事件组作为前提条件。

\begin{note}
互斥与对立的区别：互斥只要求 \(AB = \varnothing\)，而对立还要求 \(A \cup B = \Omega\)。因此对立一定是互斥的，但互斥不一定是对立的。
\end{note}

\subsection{事件的运算}

\begin{enumerate}
    \item \textbf{积事件（交）}——称"事件 \(A\) 与 \(B\) 同时发生"的事件为事件 \(A\) 与 \(B\) 的\textbf{积事件}（或交事件），记为 \(A \cap B\) 或 \(AB\)。称"有限个（或可列个）事件 \(A_1, A_2, \dots, A_n(\dots)\) 同时发生"的事件为这些事件的积事件，记为 \(\bigcap_{i=1}^{n} A_i\)（或 \(\bigcap_{i=1}^{\infty} A_i\)）。
    \item \textbf{和事件（并）}——称"事件 \(A\) 与 \(B\) 至少有一个发生"的事件为事件 \(A\) 与 \(B\) 的\textbf{和事件}（或并事件），记为 \(A \cup B\)。称"有限个（或可列个）事件 \(A_1, A_2, \dots, A_n(\dots)\) 至少有一个发生"的事件为这些事件的和事件，记为 \(\bigcup_{i=1}^{n} A_i\)（或 \(\bigcup_{i=1}^{\infty} A_i\)）。
    \item \textbf{差事件}——称"事件 \(A\) 发生而事件 \(B\) 不发生"的事件为事件 \(A\) 与 \(B\) 的\textbf{差事件}，记为 \(A - B\)。
    \item \textbf{逆事件（对立事件）}——称"事件 \(A\) 不发生"的事件为事件 \(A\) 的\textbf{逆事件}或\textbf{对立事件}，记为 \(\overline{A}\)。
\end{enumerate}

由定义易知
\[
A - B = A - AB = A\overline{B}, \quad B = \overline{A} \Leftrightarrow AB = \varnothing \text{ 且 } A \cup B = \Omega.
\]

\subsection{事件的运算律}

\begin{enumerate}
    \item \textbf{吸收律}——若 \(A \subset B\)，则 \(A \cup B = B\)，\(A \cap B = A\)。
    \item \textbf{交换律}——\(A \cup B = B \cup A\)，\(A \cap B = B \cap A\)。
    \item \textbf{结合律}——\((A \cup B) \cup C = A \cup (B \cup C)\)，\((A \cap B) \cap C = A \cap (B \cap C)\)。
    \item \textbf{分配律}——\(A \cap (B \cup C) = (A \cap B) \cup (A \cap C)\)，\(A \cup (B \cap C) = (A \cup B) \cap (A \cup C)\)。特别地，
    \[
    A \cap (B - C) = (A \cap B) - (A \cap C).
    \]
    \item \textbf{对偶律（德·摩根律）}——\(\overline{A \cup B} = \overline{A} \cap \overline{B}\)，\(\overline{A \cap B} = \overline{A} \cup \overline{B}\)。一般地，对于 \(n \ge 1\)，
    \[
    \overline{\bigcup_{i=1}^{n} A_i} = \bigcap_{i=1}^{n} \overline{A_i}; \quad \overline{\bigcap_{i=1}^{n} A_i} = \bigcup_{i=1}^{n} \overline{A_i}.
    \]
\end{enumerate}

\begin{note}
德摩根律记忆技巧：长线变短线，开口换方向。事件运算顺序约定为先进行逆运算，然后进行交运算，最后进行并或差运算。事件的关系、运算与集合的关系、运算相当，且具有相同的运算法则，所以我们可以对比着理解记忆，并要学会用集合关系去考虑事件关系。
\end{note}

\begin{exercise}\label{exer:1-1}
已知 \(A\)、\(B\)、\(C\) 为三个随机事件，则 \(A\)、\(B\)、\(C\) 不都发生的事件为（\quad ）。\\
A. $\overline{A}\overline{B}\overline{C}$\quad B. $\overline{ABC}$\quad C. $A \cup B \cup C$\quad D. $ABC$
\end{exercise}

\begin{solution}
应选 B。

这类题最稳妥的做法是先找它的\textbf{对立事件}。

“\(A\)、\(B\)、\(C\) 不都发生”并不是说“三个都不发生”，而是说“\textbf{不是三个全都发生}”。因此它的对立事件恰好是
\[
\text{A、B、C 都发生} \Longleftrightarrow ABC.
\]

所以所求事件应为
\[
\overline{ABC}.
\]

再看选项：

A 表示“三个都不发生”，范围太小，不对。

B 正是 \(\overline{ABC}\)，正确。

C 表示“至少有一个发生”，与题意不同。

D 是对立事件本身，也不对。

故答案选 B。
\end{solution}

\begin{exercise}\label{exer:1-2}
设有随机事件 \(A_i\ (i=1,2,\dots,n)\)，\(A = \bigcap_{i=1}^{n} A_i\)，则 \(A\) 的对立事件 \(\overline{A} =\ \underline{\hspace{3cm}}\)。
\end{exercise}

\begin{solution}
先把 \(A\) 的含义说清楚：
\[
A=\bigcap_{i=1}^{n}A_i
\]
表示 \(A_1,A_2,\dots,A_n\) \textbf{同时发生}。

那么它的对立事件 \(\overline A\) 就表示“这些事件\textbf{不可能全部同时发生}”，也就是“至少有一个 \(A_i\) 不发生”。

用德摩根律写成符号就是
\[
\overline{A} = \overline{\bigcap_{i=1}^{n} A_i} = \overline{A_1 \cap A_2 \cap \dots \cap A_n} = \overline{A_1} \cup \overline{A_2} \cup \dots \cup \overline{A_n} = \bigcup_{i=1}^{n} \overline{A_i}.
\]

因此应填
\[
\bigcup_{i=1}^{n}\overline{A_i}.
\]
\end{solution}

\begin{exercise}\label{exer:1-28}
设 \(A,B,C\) 是任意三个事件，则下列选项中正确的是（\quad）。\par
A. 若 \(A\cup C=B\cup C\)，则 \(A=B\)\quad
B. 若 \(A-C=B-C\)，则 \(A=B\)\quad
C. 若 \(AC=BC\)，则 \(A=B\)\quad
D. 若 \(AB=\varnothing\) 且 \(\overline{A}\,\overline{B}=\varnothing\)，则 \(\overline{A}=B\)
\end{exercise}

\begin{solution}
应选 D。

\textbf{方法一}（直接法）：

由德摩根律，
\[
\overline{A}\,\overline{B}=\overline{A\cup B}.
\]
题设给出 \(\overline{A}\,\overline{B}=\varnothing\)，故
\[
\overline{A\cup B}=\varnothing \quad\Longrightarrow\quad A\cup B=\Omega.
\]

另一方面，题设又有 \(AB=\varnothing\)。这说明 \(A\) 与 \(B\) 互不相容；而 \(A\cup B=\Omega\) 又说明二者合起来恰好覆盖整个样本空间。因此 \(A,B\) 互为对立事件，于是
\(\overline{A}=B\)。

\textbf{方法二}（排除法）：

对 A，取 \(C=\Omega\)，则 \(A\cup C=B\cup C=\Omega\) 恒成立，但 \(A\neq B\) 仍然可能发生，所以 A 错。

对 B，取 \(C=\Omega\)，则
\[
A-C=A-\Omega=\varnothing,\qquad B-C=B-\Omega=\varnothing,
\]
故 \(A-C=B-C\) 恒成立，但仍不能推出 \(A=B\)，所以 B 错。

对 C，取 \(C=\varnothing\)，则
\[
AC=BC=\varnothing
\]
恒成立，同样不能推出 \(A=B\)，所以 C 错。

因此只有 D 正确。
\end{solution}

\begin{note}
本题反映了事件运算与数的运算的不同之处。事件的并、交、差运算不满足消去律。
\end{note}

\begin{exercise}\label{exer:1-29}
用事件 \(A,B,C\) 的运算关系表示下列事件：\par
(1) \(A,B,C\) 都不发生；\par
(2) \(A,B,C\) 不都发生；\par
(3) \(A,B,C\) 不多于一个发生。
\end{exercise}

\begin{solution}
先把文字条件翻成集合语言，再直接套德摩根律。

(1) “都不发生”就是三者的对立事件同时发生：
\[
\overline{A}\,\overline{B}\,\overline{C}=\overline{A\cup B\cup C}.
\]

(2) “不都发生”就是“不是三者同时发生”：
\[
\overline{ABC}=\overline{A}\cup\overline{B}\cup\overline{C}.
\]

(3) “不多于一个发生”等价于“任意两个都不同时发生”，所以可写成
\[
\overline{AB}\cap\overline{AC}\cap\overline{BC},
\]
再展开就是四种互斥情形：
\[
\overline{A}\,\overline{B}\,\overline{C}\cup A\overline{B}\,\overline{C}\cup\overline{A}\,B\,\overline{C}\cup\overline{A}\,\overline{B}\,C.
\]
\end{solution}

\begin{exercise}\label{exer:1-30}
设随机变量 \(X,Y\)，事件 \(A=\{X\ge c\}\)，事件 \(B=\{Y\ge c\}\)。用 \(A,B\) 表示下列事件：\par
\(D=\{\min(X,Y)<c\}\)；\quad \(E=\{\max(X,Y)\ge c\}\)；\quad \(F=\{\min(X,Y)\ge c\}\)；\quad \(G=\{\max(X,Y)<c\}\)。
\end{exercise}

\begin{solution}
把 \(A=\{X\ge c\}\), \(B=\{Y\ge c\}\) 代回去即可。
\begin{itemize}
    \item \(D=\{\min(X,Y)<c\}\)：至少有一个变量小于 \(c\)，所以 \(D=\overline{A}\cup\overline{B}\)；
    \item \(E=\{\max(X,Y)\ge c\}\)：至少有一个变量不小于 \(c\)，所以 \(E=A\cup B\)；
    \item \(F=\{\min(X,Y)\ge c\}\)：两个都要不小于 \(c\)，所以 \(F=A\cap B\)；
    \item \(G=\{\max(X,Y)<c\}\)：两个都要小于 \(c\)，所以 \(G=\overline{A}\cap\overline{B}=\overline{A\cup B}\)。
\end{itemize}
\end{solution}

\begin{exercise}\label{exer:1-31}
设 \(A\) 和 \(B\) 是任意两个事件，则下列两个命题（\quad）。\par
(1) 若 \(P(A)=P(B)\)，则 \(A=B\)；\par
(2) 若 \(P(AB)=0\)，则 \(AB=\varnothing\)。
\par
A. (1) 正确，(2) 不正确\quad
B. (1) 不正确，(2) 正确\quad
C. (1)、(2) 均正确\quad
D. (1)、(2) 均不正确
\end{exercise}

\begin{solution}
应选 D。

判断这种“命题真假”题，最直接的方法就是给出反例。

取试验为“在区间 \([0,2]\) 上随机投一点”，令
\[
A=\{\text{点落入 }[0,1]\},\qquad B=\{\text{点落入 }[1,2]\}.
\]

先看命题 (1)。有
\(P(A)=P(B)=\dfrac12\)，但显然 \(A\ne B\)。
所以“若 \(P(A)=P(B)\)，则 \(A=B\)”是错误的。

再看命题 (2)。这里
\[
AB=A\cap B=\{1\},\qquad P(AB)=P(\{1\})=0,
\]
但 \(AB=\{1\}\ne\varnothing\)。
所以“若 \(P(AB)=0\)，则 \(AB=\varnothing\)”也是错误的。

两个命题都不正确，故应选 D。
\end{solution}

\begin{note}
\(P(A)=0\) 不能推出 \(A=\varnothing\)，\(P(A)=1\) 不能推出 \(A=\Omega\)。概率只反映可能性大小，不反映事件的集合构成。
\end{note}

\begin{exercise}\label{exer:1-32}
已知事件 \(A,B\) 满足 \(P(AB)=P(\overline{A}\,\overline{B})\)，记 \(P(A)=p\)，则 \(P(B)=\underline{\hspace{2cm}}\)。
\end{exercise}

\begin{solution}
应填 \(1-p\)。

题目给出的等式
\[
P(AB)=P(\overline{A}\,\overline{B})
\]
表示“\(A,B\) 同时发生”的概率，等于“\(A,B\) 同时都不发生”的概率。要由此求 \(P(B)\)，关键是先把右边化成 \(P(A)\)、\(P(B)\) 和 \(P(AB)\) 的形式。

由德摩根律，
\[
\overline{A}\,\overline{B}=\overline{A\cup B},
\]
所以
\[
P(\overline{A}\,\overline{B})=P(\overline{A\cup B})=1-P(A\cup B).
\]

再利用加法公式
\[
P(A\cup B)=P(A)+P(B)-P(AB),
\]
便得到
\[
P(\overline{A}\,\overline{B})=1-P(A)-P(B)+P(AB),
\]
从而由题设 \(P(AB)=P(\overline{A}\,\overline{B})\) 可知
\[
P(AB)=1-P(A)-P(B)+P(AB).
\]
两边同时消去 \(P(AB)\)，得 \(P(A)+P(B)=1\)。又因为题目记 \(P(A)=p\)，所以 \(P(B)=1-P(A)=1-p\)。

因此应填 \(1-p\)。
\end{solution}

\begin{exercise}\label{exer:1-33}
已知 \(P(A)=0.7\)，\(P(A-B)=0.3\)，则 \(P(\overline{AB})=\underline{\hspace{2cm}}\)。
\end{exercise}

\begin{solution}
应填 \(0.6\)。

题目给出了 \(P(A)\) 和 \(P(A-B)\)，而 \(A-B\) 表示“\(A\) 发生但 \(B\) 不发生”，因此可先用减法公式
\[
P(A-B)=P(A)-P(AB).
\]
代入数据得 \(0.3=0.7-P(AB)\)，所以 \(P(AB)=0.7-0.3=0.4\)，进而
\[
P(\overline{AB})=1-P(AB)=1-0.4=0.6.
\]

因此应填 \(0.6\)。
\end{solution}

\begin{exercise}\label{exer:1-34}
设 \(A,B\) 为两个随机事件，且 \(P(A)=0.4\)，\(P(A\cup B)=0.7\)。若 \(A,B\) 互不相容，则 \(P(B)=\underline{\hspace{2cm}}\)；若 \(A,B\) 相互独立，则 \(P(B)=\underline{\hspace{2cm}}\)；若 \(A\) 发生则 \(B\) 必发生，则 \(P(B)=\underline{\hspace{2cm}}\)。
\end{exercise}

\begin{solution}
应填 \(0.3\)；\(0.5\)；\(0.7\)。

三种情形都先用同一个加法公式
\[
P(A\cup B)=P(A)+P(B)-P(AB).
\]
\begin{itemize}
    \item 互不相容时 \(P(AB)=0\)，所以 \(P(B)=0.7-0.4=0.3\)；
    \item 相互独立时 \(P(AB)=P(A)P(B)\)，代入得 \(0.7=0.4+P(B)-0.4P(B)\)，所以 \(P(B)=0.5\)；
    \item “\(A\) 发生则 \(B\) 必发生”就是 \(A\subset B\)，于是 \(A\cup B=B\)，故 \(P(B)=0.7\)。
\end{itemize}
\end{solution}

\section{概率的定义}

\subsection{描述性定义}

通常将随机事件 \(A\) 发生的可能性大小的度量（非负值）称为事件 \(A\) 发生的\textbf{概率}，记为 \(P(A)\)。

\subsection{统计性定义}

在相同条件下做重复试验，事件 \(A\) 出现的次数 \(k\) 和总的试验次数 \(n\) 之比 \(\frac{k}{n}\) 称为事件 \(A\) 在这 \(n\) 次试验中出现的\textbf{频率}。当试验次数 \(n\) 充分大时，频率将"稳定"于某常数 \(p\)。\(n\) 越大，频率偏离这个常数 \(p\) 的可能性越小。这个常数 \(p\) 就称为事件 \(A\) 的\textbf{概率}。

\begin{note}
概率的统计性定义实质上是说，用频率 \(\frac{k}{n}\) 作为事件 \(A\) 的概率 \(P(A)\) 的估计。其直观理解为某事件出现的可能性大小，可由其在多次重复试验中出现的频率去刻画。从上述可以看出，频率只是概率的估计，而非概率本身。概率的统计性定义的重要性主要基于以下两点：
\begin{enumerate}
    \item 它提供了估计概率的方法。比如在一批产品中抽取样品，来估计该批产品的合格率（合格率是客观的数据，抽取样品计算出来的合格率只是一种估计）。
    \item 它提供了一种检验某结论是否正确的准则。比如，你说某批产品的合格率是 95\%，我们做试验，抽取样品进行计算，得出的结果是合格率为 20\%，远远低于你所说的 95\%，于是毫不犹豫地拒绝你的结论。
\end{enumerate}
\end{note}

\subsection*{公理化定义}

\begin{definition}[概率的公理化定义] \label{def:probability-axiom}
设随机试验的样本空间为 \(\Omega\)，如果对每一个事件 \(A\) 都有一个确定的实数 \(P(A)\)，且事件函数 \(P(\cdot)\) 满足：
\begin{enumerate}
    \item \textbf{非负性}：\(P(A) \ge 0\)；
    \item \textbf{规范性}：\(P(\Omega) = 1\)；
    \item \textbf{可列可加性}：对任意可列个两两互不相容事件 \(A_1, A_2, \dots, A_n, \dots\)（即 \(A_i A_j = \varnothing\)，\(i \neq j\)；\(i, j = 1, 2, \dots\)），有
    \[
    P\!\left(\bigcup_{i=1}^{\infty} A_i\right) = \sum_{i=1}^{\infty} P(A_i),
    \]
\end{enumerate}
则称 \(P(\cdot)\) 为\textbf{概率}，\(P(A)\) 为事件 \(A\) 的\textbf{概率}。
\end{definition}

\begin{remark}
(1) 数学上所说的"公理"，就是一些不加证明而承认的前提，上述公理化定义只是界定了概率这个概念所必须满足的一般性质，它不解决具体场合下的概率计算。\newline
(2) 概率 \(P(\cdot)\) 是事件的函数。\newline
(3) 虽然公理化定义不解决具体场合下的概率计算，但是我们却经常用它来判断事件函数 \(P(\cdot)\) 是否是概率，这种题型在考研试题中也会经常遇到。
\end{remark}

由概率的公理化定义，可以推导出以下基本性质：

\begin{property}\label{prop:probability-basic}
概率的基本性质：
\begin{enumerate}
    \item \(P(\varnothing) = 0\)；
    \item \textbf{有限可加性}：若 \(A_1, A_2, \dots, A_n\) 两两互不相容，则
    \[
    P\!\left(\bigcup_{i=1}^{n} A_i\right) = \sum_{i=1}^{n} P(A_i);
    \]
    \item \textbf{减法公式}：若 \(B \subset A\)，则 \(P(A - B) = P(A) - P(B)\)。更一般地，
    \[
    P(A - B) = P(A) - P(AB);
    \]
    \item \textbf{单调性}：若 \(B \subset A\)，则 \(P(B) \le P(A)\)；
    \item \textbf{有界性}：对任意事件 \(A\)，\(0 \le P(A) \le 1\)；
    \item \textbf{求逆公式}：\(P(\overline{A}) = 1 - P(A)\)；
    \item \textbf{加法公式}：\(P(A \cup B) = P(A) + P(B) - P(AB)\)。当 \(A\)、\(B\) 互斥时，\(P(A \cup B) = P(A) + P(B)\)。
    \item 设 \(A_1, A_2, A_3\) 为任意三个事件，则有
    \[
    P(A_1 \cup A_2 \cup A_3) = P(A_1) + P(A_2) + P(A_3) - P(A_1A_2) - P(A_1A_3) - P(A_2A_3) + P(A_1A_2A_3).
    \]
    \item 若 \(A_1, A_2, \dots, A_n\) 是两两互不相容的事件，则有
    \[
    P(A_1 \cup A_2 \cup \dots \cup A_n) = P(A_1) + P(A_2) + \dots + P(A_n).
    \]
\end{enumerate}
\end{property}
\begin{proof}[性质 1 的证明：\(P(\varnothing)=0\)]
由公理 3（可列可加性），取 \(A_i = \varnothing\ (i=1,2,\dots)\)，则
\[
P\!\left(\bigcup_{i=1}^{\infty} \varnothing\right) = \sum_{i=1}^{\infty} P(\varnothing).
\]
左边 \(P(\varnothing) = \sum_{i=1}^{\infty} P(\varnothing)\)，即 \(P(\varnothing)\) 等于无穷多个自身的和，故只能有 \(P(\varnothing) = 0\)。
\end{proof}

\begin{note}
概率为 0 并不意味着为空集。即 \(P(A) = 0 \nRightarrow A = \varnothing\)，而 \(A = \varnothing \Rightarrow P(A) = 0\)。同样的，$P(A) = 0$不能断言$A = \varnothing$，而$P(A) = 1$不能断言$A = \Omega$.
\end{note}

\begin{exercise}\label{exer:1-3}
设 \(A, B\) 为两个随机事件，\(P(A) = 0.4\)，\(P(A \cup B) = 0.7\)。\par
(1) 求 \(P(B - A)\)。\par
(2) 若事件 \(A\) 和事件 \(B\) 互不相容，求 \(P(B)\)。
\end{exercise}

\begin{solution}
(1) 由加法公式
\[
P(A \cup B) = P(A) + P(B) - P(AB) = P(A) + P(B - A),
\]
故 \(P(B - A) = P(A \cup B) - P(A) = 0.7 - 0.4 = 0.3\)。

(2) 若 \(A\) 与 \(B\) 互不相容，则 \(P(AB) = 0\)，此时加法公式简化为
\[
P(A \cup B) = P(A) + P(B),
\]
故 \(P(B) = P(A \cup B) - P(A) = 0.7 - 0.4 = 0.3\)。
\end{solution}


\begin{exercise}\label{exer:1-27}
设 \(A\) 和 \(B\) 是任意两个概率不为 0 的互不相容事件，则下列结论中肯定正确的是（\quad）。\par
A. \(\overline{A}\) 与 \(\overline{B}\) 不相容\quad
B. \(\overline{A}\) 与 \(\overline{B}\) 相容\quad
C. \(P(AB) = P(A)P(B)\)\quad
D. \(P(A-B) = P(A)\)
\end{exercise}

\begin{solution}
应选 D。

先看 D。由于题设说明 \(A\) 与 \(B\) 互不相容，所以 \(AB=\varnothing,\ P(AB)=0\)。而 \(A-B=A\overline{B}\)，它表示“\(A\) 发生且 \(B\) 不发生”。因为 \(A\) 与 \(B\) 不可能同时发生，故 \(A\) 发生时必然已经落在 \(A-B\) 中，于是
\[
P(A-B)=P(A)-P(AB)=P(A).
\]
因此 D 正确。

再排除其余选项。

对 A、B，仅由“\(A,B\) 互不相容”不能唯一决定 \(\overline{A}\) 与 \(\overline{B}\) 是否相容。比如若 \(A,B\) 不是对立事件，则 \(\overline{A}\) 与 \(\overline{B}\) 往往还有公共部分；只有在 \(A,B\) 恰好互为对立事件时，\(\overline{A}\) 与 \(\overline{B}\) 才互不相容。所以 A、B 都不能必然成立。

对 C，因为
\(P(AB)=0\)，但题中又给出 \(P(A)>0,\ P(B)>0\)，于是 \(P(A)P(B)>0\)。所以
\[
P(AB)\neq P(A)P(B),
\]
故 C 错。

综上，应选 D。
\end{solution}
\begin{exercise}\label{exer:1-35}
考虑一元二次方程 \(x^2+Bx+C=0\)，其中 \(B,C\) 分别是将一枚骰子接连抛两次先后出现的点数。则该方程有实根的概率为（\quad）。\par
A. \(\dfrac{19}{36}\)\quad
B. \(\dfrac{17}{36}\)\quad
C. \(\dfrac{15}{36}\)\quad
D. \(\dfrac{13}{36}\)
\end{exercise}

\begin{solution}
应选 A。

方程有实根的条件为 \(B^2-4C\ge 0\)。总样本数 \(6^2=36\)。列表统计满足条件的样本点数：
\[
\begin{array}{c|cccccc}
B\setminus C & 1 & 2 & 3 & 4 & 5 & 6 \\
\midrule
1 & - & - & - & - & - & - \\
2 & 0 & - & - & - & - & - \\
3 & + & + & - & - & - & - \\
4 & + & + & + & + & - & - \\
5 & + & + & + & + & + & + \\
6 & + & + & + & + & + & + \\
\end{array}
\]
``+''和``0''对应的样本点共 \(1+2+4+6+6=19\) 个，故 \(P=\dfrac{19}{36}\)。
\end{solution}

\begin{exercise}\label{exer:1-36}
甲口袋有 5 个白球、3 个黑球，乙口袋有 4 个白球、6 个黑球。从两个口袋中各任取一个球，则取到的两个球颜色相同的概率为（\quad）。\par
A. \(\dfrac{19}{40}\)\quad
B. \(\dfrac{21}{40}\)\quad
C. \(\dfrac{1}{4}\)\quad
D. \(\dfrac{9}{40}\)
\end{exercise}

\begin{solution}
应选 A。

把“颜色相同”拆成两种互不相交的情况：

1. 两个球都取到白球；
2. 两个球都取到黑球。

题目是“不放回依次取球”，所以我们按有序取法来数。

全部可能取法共有 \(8\times 10=80\) 种。

有利取法中，“两白”有 \(5\times 4=20\) 种，“两黑”有 \(3\times 6=18\) 种，因此
\[
\text{有利取法总数}=20+18=38,\qquad P=\frac{38}{80}=\frac{19}{40}.
\]

因此应选 A。
\end{solution}

\begin{exercise}\label{exer:1-45}
已知 \(P(B)=0.5\)，\(P(A\cup B)=0.7\)，则 \(P(A\overline{B})=\underline{\hspace{2cm}}\)，\(P(\overline{A}\overline{B})=\underline{\hspace{2cm}}\)。
\end{exercise}

\begin{solution}
应填 \(0.2\)；\(0.3\)。

由加法公式：
\[
P(A\overline{B}) = P(A) - P(AB) = P(A\cup B) - P(B) = 0.7 - 0.5 = 0.2,
\]
\[
P(\overline{A}\overline{B}) = P(\overline{A\cup B}) = 1 - P(A\cup B) = 1 - 0.7 = 0.3.
\]
\end{solution}

\begin{exercise}\label{exer:1-46}
对任意事件 \(A,B\)，下面结论正确的是（\quad）。\par
A. 若 \(P(AB)=0\)，则 \(\overline{A}\,\overline{B}=\Omega\)\quad
B. \(P(A-B)=P(A)-P(B)\)\quad
C. \(P(A\overline{B})=P(A)-P(B)+P(\overline{A}B)\)\quad
D. \(P(A\cup B)=P(A)+P(B)-P(A)P(B)\)
\end{exercise}

\begin{solution}
应选 C。

先看 A：\(P(AB)=0\) 只说明交集概率为 0，不代表交集一定为空，因此不能推出 \(\overline{A}\,\overline{B}=\Omega\)。

再看 B：\(A-B=A\overline{B}\)，所以
\[
P(A-B)=P(A)-P(AB),
\]
并不等于 \(P(A)-P(B)\)。

再看 D：加法公式里应当是 \(P(AB)\)，只有在独立时才可写成 \(P(A)P(B)\)，题目并未给独立性。

因此只有 C 正确。由
\[
P(A\overline{B})=P(A)-P(AB),\qquad P(\overline{A}B)=P(B)-P(AB),
\]
消去 \(P(AB)\) 得
\[
P(A\overline{B}) = P(A) - P(B) + P(\overline{A}B).
\]
\end{solution}

\begin{exercise}\label{exer:1-47}
设 \(A\) 和 \(B\) 是任意两个随机事件，则（\quad）。\par
A. \(P(A+B)\le P(A\cup B)\le P(A)+P(B)\)\quad
B. \(P(A)+P(B)-1\le P(AB)\le P(A\cup B)\le P(A)+P(B)\)\quad
C. \(P(A)+P(B)-1\le P(A(A\cup B))\le P(AB)\le P(A\cup(A\cup B))\)\quad
D. \(1-P(A)-P(B)\le P(AB)\le P(A)+P(B)\le P(A\cup B)\)
\end{exercise}

\begin{solution}
应选 B。

A 错误：\(P(A+B)=P(A\cup B)\)（概率中 "+" 即并）。C 错误：\(P(A(A\cup B))=P(A)\ge P(AB)\)。D 错误：\(P(A)+P(B)\ge P(A\cup B)\)，故 \(P(A)+P(B)\le P(A\cup B)\) 不成立。

B 正确：由 \(P(AB)=P(A)+P(B)-P(A\cup B)\)：
\begin{itemize}
    \item 左边：\(P(A\cup B)\le 1\)，故 \(P(AB)\ge P(A)+P(B)-1\)；
    \item 中间：\(AB\subset A\cup B\)，故 \(P(AB)\le P(A\cup B)\)；
    \item 右边：\(P(A\cup B)\le P(A)+P(B)\)（加法公式的直接推论）。
\end{itemize}
\end{solution}

\begin{exercise}\label{exer:1-48}
设 \(A\) 与 \(B\) 是两个随机事件，且 \(P(AB)=0\)，则（\quad）。\par
A. \(P(A-B)=P(A)\)\quad
B. \(A\) 与 \(B\) 相互独立\quad
C. \(P(A)=0\) 或 \(P(B)=0\)\quad
D. \(A\) 与 \(B\) 互不相容
\end{exercise}

\begin{solution}
应选 A。

先利用减法公式 \(P(A-B)=P(A)-P(AB)\)。题目已知 \(P(AB)=0\)，所以
\[
P(A-B)=P(A)-0=P(A),
\]
故 A 正确。

下面排除其余选项。

B 说 \(A\) 与 \(B\) 相互独立。独立要求 \(P(AB)=P(A)P(B)\)，
而题目只给了 \(P(AB)=0\)。除非 \(P(A)=0\) 或 \(P(B)=0\)，否则一般并不满足独立，所以 B 不一定成立。

C 说 \(P(A)=0\) 或 \(P(B)=0\)。这也不对。完全可能出现 \(P(A)>0,P(B)>0\)，但交集只是一个概率为 0 的细小集合。

D 说 \(A\) 与 \(B\) 互不相容。互不相容要求 \(AB=\varnothing\)，而 \(P(AB)=0\) 只说明交集的概率为 0，并不能推出交集一定是空集，因此 D 也不对。

故答案选 A。
\end{solution}


\begin{exercise}[提高]\label{exer:1-49}
设事件 \(A,B\) 相互独立，事件 \(B,C\) 互不相容，事件 \(A\) 与 \(C\) 不能同时发生，且 \(P(A)=P(B)=0.5\)，\(P(C)=0.2\)。则事件 \(A,B\) 和 \(C\) 中仅 \(C\) 发生或仅 \(C\) 不发生的概率为 \underline{\hspace{2cm}}。
\end{exercise}

\begin{solution}
应填 \(0.45\)。“仅 \(C\) 发生”即 \(C\overline{A}\,\overline{B}\)。由于 \(AC=\varnothing\) 且 \(BC=\varnothing\)，故 \(C\subset\overline{A}\) 且 \(C\subset\overline{B}\)，即 \(C\subset\overline{A}\,\overline{B}\)，从而 \(C\overline{A}\,\overline{B}=C\)，\(P(C)=0.2\)。

“仅 \(C\) 不发生”即 \(\overline{C}AB\)。由独立性 \(P(AB)=P(A)P(B)=0.25\)，且 \(AB\subset\overline{C}\)（因为 \(AC=\varnothing\)），故 \(\overline{C}AB=AB\)，\(P(\overline{C}AB)=0.25\)。

又 \(C\) 与 \(AB\) 互斥（\(C\subset\overline{A}\)，\(AB\subset A\)），故
\[
P(C\overline{A}\,\overline{B}\cup\overline{C}AB) = P(C)+P(AB) = 0.2+0.25 = 0.45.
\]
\end{solution}

\begin{exercise}[提高]\label{exer:1-50}
设 \(A,B\) 为任意两个随机事件，则（\quad）。\par
\begin{tabular}{ll}
A. \(P(AB)\le P(A)P(B)\) & B.\(P(AB)\ge P(A)P(B)\) \\[4pt]
C. \(P(AB)\le \dfrac{P(A)+P(B)}{2}\) & D. \(P(AB)\ge \dfrac{P(A)+P(B)}{2}\)
\end{tabular}
\end{exercise}

\begin{solution}
应选 C。由 \(AB\subset A\) 和 \(AB\subset B\)，得 \(P(AB)\le P(A)\) 且 \(P(AB)\le P(B)\)，因此
\[
2P(AB)\le P(A)+P(B) \implies P(AB)\le\frac{P(A)+P(B)}{2}.
\]

A、B 的反例：当 \(A\subset B\) 时 \(P(AB)=P(A)\)，与 \(P(A)P(B)\) 没有固定大小关系；当 \(AB=\varnothing\) 时 \(P(AB)=0\)。D 更不可能恒成立。
\end{solution}


\section{古典概型}

\begin{definition}[古典概型] \label{def:classical-prob}
如果随机试验的样本空间满足：
\begin{enumerate}
    \item 只有有限个样本点（基本事件）；
    \item 每个样本点（基本事件）发生的可能性都一样，
\end{enumerate}
则称随机试验的概率模型为\textbf{古典概型}（等可能概型）。
\end{definition}

如果古典概型的基本事件总数为 \(n\)，事件 \(A\) 包含 \(k\) 个基本事件，则 \(A\) 的概率为
\[
P(A) = \frac{k}{n} = \frac{\text{事件} \ A \ \text{所含基本事件的个数}}{\text{基本事件总数}}.
\]

由上式计算得出的概率称为 \(A\) 的\textbf{古典概率}。

\begin{note}
计算古典概率的关键是基本事件和样本空间的选定以及基本事件数的计算：
\begin{enumerate}
    \item \textbf{列举法}（直接数数法）：基本事件数不多时常用。
    \item \textbf{集合对应法}：
    \begin{enumerate}
        \item \textbf{加法原理}——完成一件事有 \(n\) 类办法，第一类办法中有 \(m_1\) 种方法，第二类办法中有 \(m_2\) 种方法，\(\dots\)，第 \(n\) 类办法中有 \(m_n\) 种方法，则完成此事共有 \(\sum_{i=1}^{n} m_i\) 种方法。
        \item \textbf{乘法原理}——完成一件事有 \(n\) 个步骤，第一步有 \(m_1\) 种方法，第二步有 \(m_2\) 种方法，\(\dots\)，第 \(n\) 步有 \(m_n\) 种方法，则完成此事共有 \(\prod_{i=1}^{n} m_i\) 种方法。
        \item \textbf{排列}——从 \(n\) 个不同的元素中取出 \(m(\le n)\) 个元素，并按照一定顺序排成一列，叫作排列。所有排列的个数叫作排列数，记作
        \[
        \mathrm{P}_n^m = n(n-1)(n-2)\cdots(n-m+1) = \frac{n!}{(n-m)!}.
        \]
        当 \(m=n\) 时，\(\mathrm{P}_n^n = \frac{n!}{0!} = n!\)，叫作全排列。
        \item \textbf{组合}——从 \(n\) 个不同的元素中取出 \(m(\le n)\) 个元素，并成一组，叫作组合。所有组合的个数叫作组合数，记作
        \[
        \mathrm{C}_n^m = \frac{\mathrm{P}_n^m}{m!} = \frac{n!}{m!(n-m)!}, \quad \mathrm{C}_n^m = \mathrm{C}_n^{n-m}, \quad 0!=1, \quad \mathrm{C}_n^0=1.
        \]
    \end{enumerate}
    \item \textbf{逆数法}：先求 \(\overline{A}\) 中的基本事件数 \(n_{\overline{A}}\)，用基本事件总数 \(n\) 减去 \(n_{\overline{A}}\) 便得 \(A\) 中的基本事件数，这种方法常用于计算含有"至少"字样的事件的概率。
\end{enumerate}
\end{note}

\subsection*{随机分配问题}

随机分配也叫随机占位，突出一个"放"字，即将 \(n\) 个可辨质点随机地分配到 \(N\) 个盒子中，区分每盒可以容纳任意多个质点（ 不同分法的总数\(N^n\)）和最多可以容纳一个质点（不同分法的总数\(\mathrm{P}_N^n = N(N-1)\cdots(N-n+1)\)）。

\begin{note}
每盒容纳任意多个质点：每个质点均可放到 \(N\) 个盒子中的任何一个，即有 \(N\) 种放法，于是 \(n\) 个可辨质点放到 \(N\) 个盒子中共有 \(N^n\) 种不同放法。\newline
每盒容纳至多一个质点：质点可辨，且一个盒子至多容纳一个质点，故 \(n\) 个质点放到 \(N(N\ge n)\) 个盒子中的所有不同放法即从 \(N\) 个元素中选取 \(n\) 个元素的排列数 \(\mathrm{P}_N^n\)。
\end{note}

\begin{exercise}\label{exer:1-4}
将 \(n\) 个球随机放入 \(N(n\le N)\) 个盒子中，每个盒子可以放任意多个球。求下列事件的概率：\par
(1) \(A=\{\text{某指定} \ n \ \text{个盒子各有一球}\}\)；\par
(2) \(B=\{\text{恰有} \ n \ \text{个盒子各有一球}\}\)；\par
(3) \(C=\{\text{指定} \ k(k\le n) \ \text{个盒子各有一球}\}\)。
\end{exercise}

\begin{solution}
设 \(n\) 个球、\(N\) 个盒子是可分辨的（例如编号），由于每个盒子可以放任意多个球，因此每个球都有 \(N\) 种不同的放置方法。将 \(n\) 个球随机放入 \(N\) 个盒子的一种放法作为基本事件，则基本事件总数为 \(N^n\)。

(1) 事件 \(A\) 所含的基本事件是 \(n\) 个不同球的一种排列，故
\[
n_A = n!, \quad P(A) = \frac{n!}{N^n}.
\]

(2) 事件 \(B\) 中的基本事件可以设想为先从 \(N\) 个盒子中选出 \(n\) 个（共有 \(\mathrm{C}_N^n\) 种不同选法），而后把 \(n\) 个球随机放入这 \(n\) 个盒子中，每盒一球（共有 \(n!\) 种不同放法），因此
\[
n_B = \mathrm{C}_N^n \cdot n!, \quad P(B) = \frac{\mathrm{C}_N^n \cdot n!}{N^n}.
\]

(3) 事件 \(C\) 的基本事件可以设想为先从 \(n\) 个球中选出 \(k\) 个球（有 \(\mathrm{C}_n^k\) 种不同选法），再将这 \(k\) 个球随机放入指定的 \(k\) 个盒子中，每盒一球（有 \(k!\) 种不同放法），最后将余下的 \(n-k\) 个球随机放入其余的 \(N-k\) 个盒子中，每个球都有 \(N-k\) 种放置方法，因此共有 \((N-k)^{n-k}\) 种不同放法，故
\[
n_C = \mathrm{C}_n^k \cdot k! \cdot (N-k)^{n-k}, \quad P(C) = \frac{\mathrm{C}_n^k k! (N-k)^{n-k}}{N^n} = \frac{n! (N-k)^{n-k}}{(n-k)! N^n}.
\]
\end{solution}

\begin{note}
许多问题的结构形式与分球入盒问题相同，都属于随机占位问题。例如生日问题（\(n\) 个人生日，相当于 \(n\) 个球随机放入 365 个盒子中，每盒可以放多个球）；住房分配问题（\(n\) 个人被分配到 \(N\) 个房间中去，每个房间可住多个人）；乘客下车问题（\(n\) 个乘客在 \(N\) 个车站下车的各种可能情况）；等等。对这些问题的求解都可以用"将 \(n\) 个球等可能地投放到 \(N\) 个盒子中"的思路来考虑。
\end{note}

\subsection*{简单随机抽样问题}

设 \(\Omega = \{\omega_1, \omega_2, \dots, \omega_N\}\) 含 \(N\) 个元素，称 \(\Omega\) 为总体。如果各元素被抽到的可能性相同，则总体 \(\Omega\) 的抽样称作\textbf{简单随机抽样}，突出一个"取"字。

简单随机抽样分为先后有放回、先后无放回及任取这三种不同的方式。

\begin{enumerate}
  \item 先后有放回取 \(n\) 次 ，抽取法总数 \(N^n\) 
  \item 先后无放回取 \(n\) 次 ，抽取法总数 \(\mathrm{P}_N^n = N(N-1)\cdots(N-n+1)\)
  \item 任取 \(n\) 个，抽取法总数 \(\mathrm{C}_N^n\)
\end{enumerate}

\begin{note}
先后有放回取 \(n\) 次：既考虑抽到何元素，又考虑各元素出现的顺序，每次从 \(\Omega\) 中随意抽取一个元素，并在抽下一元素前将其放回 \(\Omega\)，于是每次都有 \(N\) 个元素可被抽取，即有 \(N\) 种抽取方法，抽取 \(n\) 次，即 \(N^n\)。\newline
先后无放回取 \(n\) 次：既考虑抽到何元素，又考虑各元素出现的顺序，只是抽出的元素均不再放回 \(\Omega\)，于是每次抽取时都比上一次少了一个元素，抽 \(n(n\le N)\) 次，即 \(\mathrm{P}_N^n = N(N-1)\cdots(N-n+1)\)。\newline
任取 \(n\) 个：一次性取 \(n\) 个元素，相当于将 \(n\) 个元素无序且无放回地取走，共抽取法总数为 \(\mathrm{C}_N^n = \dfrac{\mathrm{P}_N^n}{n!}\)。
\end{note}

\begin{exercise}\label{exer:1-5}
袋中有 5 个球，3 个白球，2 个黑球。\par
(1) 先后有放回取 2 个球；\par
(2) 先后无放回取 2 个球；\par
(3) 任取 2 个球。\par
求取的 2 个球中至少 1 个是白球的概率。
\end{exercise}

\begin{solution}
用对立事件思想，计算"两球全黑"的种数，再用总数减去它。
\begin{enumerate}
    \item 先后有放回取 2 个球：
    总数 \(5^2 = 25\)，两球全黑 \(2^2 = 4\)，故
    \[
    P = \frac{5^2 - 2^2}{5^2} = \frac{21}{25}.
    \]
    \item 先后无放回取 2 个球：
    总数 \(5 \cdot 4 = 20\)，两球全黑 \(2 \cdot 1 = 2\)，故
    \[
    P = \frac{5 \cdot 4 - 2 \cdot 1}{5 \cdot 4} = \frac{18}{20} = \frac{9}{10}.
    \]
    \item 任取 2 个球：
    总数 \(\mathrm{C}_5^2 = 10\)，两球全黑 \(\mathrm{C}_2^2 = 1\)，故
    \[
    P = \frac{\mathrm{C}_5^2 - \mathrm{C}_2^2}{\mathrm{C}_5^2} = \frac{9}{10}.
    \]
\end{enumerate}
\end{solution}

\begin{note}
上述 (2) 与 (3) 概率相同，这是因为 \(5 \cdot 4 = \mathrm{C}_5^2 \cdot 2!\)，左边上下有序，右边上下无序，相当于把顺序"消掉"了，故"先后无放回取 \(k\) 个球"与"任取 \(k\) 个球"的概率相同。由于 (3) 较方便，因此计算 (2) 时可按 (3) 来计算。
\end{note}

\begin{exercise}\label{exer:1-6}
袋中有 100 个球，40 个白球，60 个黑球。\par
(1) 先后有放回取 20 个球，求取出 15 个白球、5 个黑球的概率；\par
(2) 先后无放回取 20 个球，求取出 15 个白球、5 个黑球的概率；\par
(3) 先后有放回取 20 个球，求第 20 次取到白球的概率；\par
(4) 先后无放回取 20 个球，求第 20 次取到白球的概率。
\end{exercise}
\begin{solution}
\begin{enumerate}
    \item \[
    p_1 = \frac{\mathrm{C}_{20}^{15} \cdot 40^{15} \cdot 60^{5}}{100^{20}}.
    \]
    \item 按照"任取 20 个球"来计算，即
    \[
    p_2 = \frac{\mathrm{C}_{40}^{15} \mathrm{C}_{60}^{5}}{\mathrm{C}_{100}^{20}}.
    \]
    \item 有放回取球，每次抽取的样本空间没有变化，故每次取到白球的概率始终为 \(\dfrac{40}{100} = \dfrac{2}{5}\)，即
    \[
    p_3 = \frac{40}{100} = \frac{2}{5}.
    \]
    \item 看作随机占位问题：设有 100 个盒子，每个盒子中放 1 个球，则要求第 20 个盒子中放入白球即可。故
    \[
    p_4 = \frac{\mathrm{C}_{40}^1 \cdot 99!}{100!} = \frac{40}{100} = \frac{2}{5}.
    \]
\end{enumerate}
\end{solution}

\begin{note}
\(p_4 = p_3\)。无放回抽取时，每次取到白球的概率也不会变，这可理解成依概率摸球（抓阄模型）。类似地，设有 100 个灰球，白的成分与黑的成分之比为 \(40:60\)，每次取到球中白的成分始终是 \(40\%\)，不受抽取顺序的影响。
\end{note}

\begin{exercise}\label{exer:1-26}
从一副扑克牌（52 张）中任取 3 张（不重复），则取出 3 张牌中至少有 2 张花色相同的概率为 \underline{\hspace{2cm}}。
\end{exercise}

\begin{solution}
应填 \(\dfrac{256}{425}\)。

题目求“至少有 2 张花色相同”，直接分类会比较繁琐，更适合用对立事件。

对立事件是：
\(\text{3 张牌花色各不相同}\)。

总取法数为 \(\mathrm C_{52}^{3}\)。

下面计算“花色各不相同”的取法数。先从 4 种花色中选出 3 种：
\(\mathrm C_4^3\)。
然后在每一种被选中的花色中各取 1 张牌，每种花色都有 13 张可选，所以共有
\(13^3\)
种取法。

因此
\[
P(\text{花色各不相同})
=\frac{\mathrm C_4^3\cdot 13^3}{\mathrm C_{52}^{3}}
=\frac{4\cdot 2197}{22100}
=\frac{8788}{22100}.
\]

于是所求概率为
\[
P(\text{至少 2 张同花色})
=1-\frac{8788}{22100}
=\frac{13312}{22100}
=\frac{256}{425}.
\]
\end{solution}


\begin{exercise}[提高]\label{exer:1-51}
从 1 到 9 的 9 个整数中有放回地随机取 3 次，每次取一个数，则取出的 3 个数之积能被 10 整除的概率为 \underline{\hspace{2cm}}。
\end{exercise}

\begin{solution}
应填 \(\dfrac{52}{243}\)。积能被 10 整除当且仅当积同时含有因子 2 和因子 5。从 \(\{1,\dots,9\}\) 中，5 的倍数仅有 \(\{5\}\)，2 的倍数为 \(\{2,4,6,8\}\)。

设 \(A\) 为“至少取出一个 5 的倍数”，\(B\) 为“至少取出一个 2 的倍数”。则所求就是 \(P(AB)\)。先算三个补事件：
\[
P(\overline A)=\left(\frac89\right)^3,\qquad P(\overline B)=\left(\frac59\right)^3,\qquad P(\overline A\,\overline B)=\left(\frac49\right)^3.
\]
由容斥原理：
\begin{align*}
P(AB) &= 1 - P(\overline{A}) - P(\overline{B}) + P(\overline{A}\,\overline{B}) = 1 - \frac{8^3}{9^3} - \frac{5^3}{9^3} + \frac{4^3}{9^3} \\
&= 1 - \frac{512}{729} - \frac{125}{729} + \frac{64}{729} = \frac{156}{729} = \frac{52}{243}.
\end{align*}
\end{solution}


\subsection{几何概型}

\begin{definition}[几何概型] \label{def:geometric-prob}
如果随机试验的样本空间满足：
\begin{enumerate}
    \item 样本空间（基本事件空间）\(\Omega\) 是一个可度量的有界区域；
    \item 每个样本点（基本事件）发生的可能性都一样，即样本点落入 \(\Omega\) 的某一可度量的子区域 \(S\) 的可能性大小与 \(S\) 的几何度量成正比，而与 \(S\) 的位置及形状无关，
\end{enumerate}
则称随机试验的概率模型为\textbf{几何概型}。
\end{definition}

如果 \(S_A\) 是样本空间 \(\Omega\) 的一个可度量的子区域，则事件 \(A = \{\text{样本点落入区域 } S_A\}\) 的概率为
\[
P(A) = \frac{S_A \text{ 的几何度量}}{\Omega \text{ 的几何度量}}.
\]

\begin{note}
古典概型与几何概型的区别：基本事件有限、等可能发生的随机试验为古典概型；基本事件无限且具有几何度量、等可能发生的随机试验为几何概型。
\end{note}

\begin{exercise}\label{exer:1-7}
在区间 \((0,1)\) 中随机地取两个数，则这两个数之差的绝对值小于 \(0.5\) 的概率为 \underline{\hspace{2cm}}。
\end{exercise}

\begin{solution}
应填 \(0.75=\dfrac34\)。

把在 \((0,1)\) 中随机取两个数记为 \(x,y\)，则样本点 \((x,y)\) 在单位正方形
\[
D=\{(x,y)\mid 0<x<1,\ 0<y<1\}
\]
内均匀分布，所以问题转化为一个几何概型。

条件 \(|x-y|<0.5\) 表示点 \((x,y)\) 必须落在两条直线
\[
y=x+0.5,\qquad y=x-0.5
\]
之间的带状区域里。

更方便的做法是求它的对立区域 \(|x-y|\ge 0.5\)。在单位正方形内，这部分正好由左上角和右下角两个全等直角三角形组成，每个三角形的直角边长都是 \(0.5\)，所以每个面积为
\[
\frac12\cdot 0.5\cdot 0.5=\frac18.
\]
两块合起来面积为 \(\frac18+\frac18=\frac14\)，故所求概率为 \(1-\frac14=\frac34\)。
\end{solution}

\begin{exercise}\label{exer:1-24}
长方形 \(ABCD\) 的长 \(AB=10\,\mathrm{cm}\)，宽 \(BC=5\,\mathrm{cm}\)，在长方形 \(ABCD\) 内任取一点 \(E\)，则 \(P\!\left(\angle AEB \ge \dfrac{\pi}{2}\right) = \underline{\hspace{2cm}}\)。
\end{exercise}

\begin{solution}
应填 \(\dfrac{\pi}{4}\)。

本题的关键是利用圆周角定理。对于固定的线段 \(AB\)，满足 \(\angle AEB=90^\circ\) 的点 \(E\) 全都落在以 \(AB\) 为直径的圆上；进一步，\(\angle AEB\ge \frac{\pi}{2}\) 当且仅当点 \(E\) 落在这个圆对应的半圆内部或圆周上。

因此，满足题意的区域就是“以 \(AB\) 为直径的半圆”。

因为 \(AB=10\text{ cm}\)，所以半径 \(r=\frac{AB}{2}=5\text{ cm}\)，半圆面积 \(S_{\text{半圆}}=\frac12\pi r^2=\frac{25\pi}{2}\)，矩形面积 \(S_{\text{矩形}}=AB\cdot BC=50\)。
又由于 \(BC=5\) 恰好等于半径，所以这个半圆刚好完整地落在矩形内部，不会超出去。

于是按几何概型，
\[
P\!\left(\angle AEB\ge \frac{\pi}{2}\right)
=\frac{S_{\text{半圆}}}{S_{\text{矩形}}}
=\frac{25\pi/2}{50}
=\frac{\pi}{4}.
\]
\end{solution}


\begin{exercise}\label{exer:1-37}
从 \([0,1]\) 中随机地取两个数，其积大于 \(\dfrac{1}{4}\)，其和小于 \(\dfrac{5}{4}\) 的概率为 \underline{\hspace{2cm}}。
\end{exercise}

\begin{solution}
应填 \(\dfrac{15}{32}-\dfrac{1}{2}\ln 2\)。设两个数分别为 \(x,y\)，则 \((x,y)\) 在正方形 \(\Omega=[0,1]^2\) 内均匀分布。所求事件为
\[A = \left\{(x,y)\,\middle|\, x+y<\frac{5}{4},\ xy>\frac{1}{4}\right\}.\]
曲线 \(xy=\dfrac{1}{4}\) 与 \(x+y=\dfrac{5}{4}\) 交于 \(x=\dfrac{1}{4}\) 和 \(x=1\)。对 \(x\in\left[\dfrac{1}{4},1\right]\)，\(y\) 的范围从 \(\dfrac{1}{4x}\) 到 \(\dfrac{5}{4}-x\)：
\begin{align*}
P(A) &= \int_{1/4}^{1}\left(\frac{5}{4}-x-\frac{1}{4x}\right)\mathrm{d}x = \left[\frac{5}{4}x-\frac{1}{2}x^2-\frac{1}{4}\ln x\right]_{1/4}^{1} \\
&= \frac{3}{4}-\left(\frac{9}{32}+\frac{1}{2}\ln 4\right) = \frac{15}{32}-\frac{1}{2}\ln 2.
\end{align*}
\end{solution}

\begin{exercise}\label{exer:1-38}
随机地向半圆 \(0<y<\sqrt{2ax-x^2}\)（\(a>0\)）内投掷一点，求该点和原点连线与 \(x\) 轴的夹角 \(\theta\le\dfrac{\pi}{4}\) 的概率。
\end{exercise}
\begin{solution}
应填 \(\dfrac{1}{2}+\dfrac{1}{\pi}\)。半圆方程 \(y=\sqrt{2ax-x^2}\) 即 \((x-a)^2+y^2=a^2\)（\(y>0\)），是以 \((a,0)\) 为圆心、\(a\) 为半径的上半圆。直线 \(y=x\) 与半圆交于 \((0,0)\) 和 \((a,a)\)。\(\theta\le\pi/4\) 对应 \(y\le x\)。在半圆内，该区域由两部分组成：
\begin{itemize}
    \item 三角形区域 \((0,0),(a,0),(a,a)\)：面积为 \(\dfrac{a^2}{2}\)；
    \item 半圆从 \(x=a\) 到 \(x=2a\) 的部分：面积为 \(\dfrac{1}{4}\pi a^2\)。
\end{itemize}
半圆总面积 \(S_\Omega=\dfrac{1}{2}\pi a^2\)，故
\[P = \frac{a^2/2+\pi a^2/4}{\pi a^2/2} = \frac{1}{2}+\frac{1}{\pi}.\]
\end{solution}


\section{条件概率}

\subsection{条件概率的定义}

\begin{definition}[条件概率] \label{def:conditional-prob}
设 \(A, B\) 是两个事件，且 \(P(A) > 0\)，称
\[
P(B|A) = \frac{P(AB)}{P(A)}
\]
为在事件 \(A\) 发生的条件下事件 \(B\) 发生的\textbf{条件概率}。
\end{definition}

\begin{note}
\textbf{自学追问：} 为什么分母是 \(P(A)\) 而不是 \(P(B)\)？因为题目说的是“在 \(A\) 已经发生的条件下”，此时新的样本空间已经缩小为 \(A\)，所以总量应当由 \(P(A)\) 来衡量；而 \(AB\) 只是这个新样本空间中又满足 \(B\) 的那一部分。

\textbf{易错点：} \(P(B|A)\) 不是 \(P(AB)\)，也一般不等于 \(P(A|B)\)。判断条件概率时，先看竖线右边谁是“已知条件”，分母就由谁决定。
\end{note}

\begin{property}\label{prop:conditional-prob}
条件概率 \(P(\cdot|A)\) 满足概率的三条公理：
\begin{enumerate}
    \item \textbf{非负性}：对于每一事件 \(B\)，有 \(P(B|A) \ge 0\)；
    \item \textbf{规范性}：\(P(\Omega|A) = 1\)；
    \item \textbf{可列可加性}：设 \(B_1, B_2, \dots, B_n, \dots\) 是两两互斥的事件，则
    \[
    P\!\left(\bigcup_{i=1}^{\infty} B_i \Big| A\right) = \sum_{i=1}^{\infty} P(B_i|A).
    \]
\end{enumerate}
因此，概率的一切性质和公式对条件概率都适用。
\end{property}
\begin{proof}
设 \(P(A) > 0\)。(1) 非负性：\(P(B|A) = P(AB)/P(A) \geqslant 0\)。\par
(2) 规范性：\(P(\Omega|A) = P(A)/P(A) = 1\)。\par
(3) 可列可加性：设 \(\{B_i\}\) 两两互斥，则
\[
P\!\left(\bigcup_{i=1}^{\infty} B_i \Big| A\right) = \frac{P\!\bigl[(\bigcup B_i)A\bigr]}{P(A)} = \frac{\sum P(B_i A)}{P(A)} = \sum P(B_i|A).
\]
\end{proof}

\begin{note}
对于条件概率 \(P(B|A)\)，可以理解为样本空间由原来的 \(\Omega\) 缩小为 \(A\)，即
\[
P(B) = P(B|\Omega) \xrightarrow{\Omega \to A} P(B|A).
\]
但要注意：\textbf{“样本空间缩小了”并不意味着任意事件的条件概率一定变大或变小}；条件概率的变化方向取决于所附加条件与目标事件之间的关系。例如，当 \(A\subset B\) 且 \(P(B)>0\) 时，
\[
P(A|B)=\frac{P(A)}{P(B)}\ge P(A),
\]
这是因为此时 \(A\cap B=A\)。条件概率的本质，是在“已知某事件已经发生”的前提下重新衡量另一个事件发生的可能性，例如 \(P(\overline{B}|A)=1-P(B|A)\)，以及
\[
P[(B-C)|A] = P(B|A) - P(BC|A).
\]
\end{note}

\subsection{乘法公式}

\begin{theorem}[乘法公式] \label{thm:multiplication}
设 \(P(A) > 0\)，则
\[
P(AB) = P(B|A)\,P(A).
\]
一般地，设 \(A_1, A_2, \dots, A_n\) 为 \(n\) 个事件（\(n \ge 2\)），且 \(P(A_1 A_2 \cdots A_{n-1}) > 0\)，则
\[
P(A_1 A_2 \cdots A_n) = P(A_1)\,P(A_2|A_1)\,P(A_3|A_1 A_2) \cdots P(A_n|A_1 A_2 \cdots A_{n-1}).
\]
\end{theorem}
\begin{proof}
对 \(n=2\)，由条件概率的定义直接可得 \(P(AB)=P(B|A)\,P(A)\)。
一般情形用数学归纳法。设 \(n-1\) 时成立，则
\begin{align*}
P(A_1 A_2 \cdots A_n) &= P\!\bigl[(A_1 \cdots A_{n-1}) \cdot A_n\bigr] \\
&= P(A_n | A_1 \cdots A_{n-1}) \cdot P(A_1 \cdots A_{n-1}),
\end{align*}
对 \(P(A_1 \cdots A_{n-1})\) 应用归纳假设即得。
\end{proof}

\begin{note}
\textbf{使用提醒：} 乘法公式适合处理“若干事件同时发生”的问题。做题时若题目信息天然带有先后顺序，例如“先抽到什么，再抽到什么”“先命中，再继续命中”，就应优先把交事件拆成“前面的概率 \(\times\) 后面的条件概率”。

\textbf{易错点：} 把 \(P(A_1A_2\cdots A_n)\) 直接写成 \(P(A_1)\cdots P(A_n)\) 时，其实额外假设了这些事件相互独立；若题目没有给独立性，只能使用条件概率连乘的形式。
\end{note}

\subsection{全概率公式与贝叶斯公式}

前面已经介绍了完备事件组的概念（见定义~\ref{def:complete-event-group}），它是全概率公式和贝叶斯公式的前提条件。

\begin{theorem}[全概率公式] \label{thm:total-probability}
设 \(B_1, B_2, \dots, B_n\) 是一组完备事件组，\(A\) 为任一事件，且 \(P(B_i) > 0\ (i=1,2,\dots,n)\)，则
\[
P(A) = \sum_{i=1}^{n} P(B_i)\,P(A|B_i).
\]
\end{theorem}

全概率公式的实质是\textbf{全集分解法}：如果事件 \(A\) 的发生总是与多个"原因" \(B_i\) 相联系，则
\begin{proof}
因 \(B_1, \dots, B_n\) 是完备事件组，故两两互斥且 \(B_1 \cup B_2 \cup \cdots \cup B_n = \Omega\)。于是
\[
A = A\Omega = A\!\left(\bigcup_{i=1}^{n} B_i\right) = \bigcup_{i=1}^{n} (AB_i),
\]
且 \(AB_1, AB_2, \dots, AB_n\) 两两互斥。由有限可加性和条件概率的定义：
\[
P(A) = \sum_{i=1}^{n} P(AB_i) = \sum_{i=1}^{n} P(B_i)\,P(A|B_i).
\]
\end{proof}

\begin{note}
\textbf{使用提醒：} 全概率公式的关键不在“背公式”，而在于先把样本空间按一组互斥且完备的原因 \(B_1,\dots,B_n\) 分层。若题目存在“先来自哪一类”“先由哪一种原因导致”的结构，就适合先按原因分解，再分别求每一层对结果 \(A\) 的贡献。

\textbf{易错点：} 完备事件组必须两两互斥且并为全集。若几个“原因”之间还有重叠，或没有覆盖全部可能情形，就不能直接代入全概率公式。
\end{note}

\begin{theorem}[贝叶斯公式] \label{thm:bayes}
设 \(B_1, B_2, \dots, B_n\) 是一组完备事件组，\(A\) 为任一事件，且 \(P(A) > 0\)，\(P(B_k) > 0\ (k=1,2,\dots,n)\)，则
\[
P(B_k|A) = \frac{P(B_k)\,P(A|B_k)}{P(A)} = \frac{P(B_k)\,P(A|B_k)}{\sum_{i=1}^{n} P(B_i)\,P(A|B_i)}, \quad k=1,2,\dots,n.
\]
\end{theorem}
\begin{proof}
由条件概率的定义、乘法公式和全概率公式：
\[
P(B_k|A) = \frac{P(AB_k)}{P(A)} = \frac{P(B_k)\,P(A|B_k)}{\displaystyle\sum_{i=1}^{n} P(B_i)\,P(A|B_i)}.
\]
\end{proof}

\begin{note}
全概率公式与贝叶斯公式的记忆技巧：
\begin{itemize}
    \item \textbf{全概率公式}——由因求果。将 \(A\) 视为结果，\(B_i\) 视为导致 \(A\) 发生的原因，已知各原因的概率 \(P(B_i)\) 及各原因导致结果 \(A\) 的概率 \(P(A|B_i)\)，求结果 \(A\) 发生的概率 \(P(A)\)。
    \item \textbf{贝叶斯公式}——执果寻因。已知结果 \(A\) 已经发生，反推是由第 \(k\) 个原因 \(B_k\) 导致的概率 \(P(B_k|A)\)。
\end{itemize}

\textbf{自学追问：} 什么时候要“先全概率、后贝叶斯”？当贝叶斯公式分母里的 \(P(A)\) 没有直接给出时，通常先用全概率公式把 \(P(A)\) 求出来，再代回贝叶斯公式。

\textbf{易错点：} 贝叶斯公式求的是后验概率 \(P(B_k|A)\)，它与题目常常先给出的 \(P(A|B_k)\) 不是同一个量。前者是“结果已知，反推原因”；后者是“原因已知，预测结果”。
\end{note}

\begin{exercise}\label{exer:1-8}
某人外出可以乘坐飞机、火车、轮船、汽车四种交通工具，其概率依次为 \(0.05, 0.15, 0.30, 0.50\)，而乘坐这几种交通工具能如期到达的概率依次为 \(0.80, 0.70, 0.60, 0.90\)。求：
\par
(1) 该人如期到达的概率；\par
(2) 已知该人误期到达，求他是乘坐火车的概率。
\end{exercise}

\begin{solution}
设 \(A_i\ (i=1,2,3,4)\) 分别表示乘坐飞机、火车、轮船、汽车，\(B\) 表示如期到达。

(1) 利用全概率公式：
\[
P(B) = \sum_{i=1}^{4} P(A_i)\,P(B|A_i) = 0.05 \times 0.80 + 0.15 \times 0.70 + 0.30 \times 0.60 + 0.50 \times 0.90 = 0.775.
\]

(2) 利用贝叶斯公式：
\[
P(A_2|\overline{B}) = \frac{P(A_2)\,P(\overline{B}|A_2)}{P(\overline{B})} = \frac{0.15 \times (1-0.70)}{1-0.775} = \frac{0.045}{0.225} = 0.2.
\]
\end{solution}

\begin{exercise}\label{exer:1-9}
设 \(A, B\) 为两个随机事件，则（\quad）。\\
\begin{tabular}{ll}
A. \(P(AB)+P(\overline{A}\,\overline{B})\le 1\) & B. \(P(AB)+P(\overline{A}\,\overline{B})\ge 1\) \\[4pt]
C. \(P(A-B)\le P(A)-P(B)\) & D. \(P(A\cup B)+P(\overline{A}\cup\overline{B})\le 1\)
\end{tabular}
\end{exercise}

\begin{solution}
应选 A。由于 \(AB \subset A \cup B\)，即 \(P(AB) \le P(A \cup B)\)，从而
\[
P(AB) \le P(A \cup B) = 1 - P(\overline{A \cup B}) = 1 - P(\overline{A}\,\overline{B}),
\]
即 \(P(AB) + P(\overline{A}\,\overline{B}) \le 1\)，故选 A。

同时有
\[
P(A \cup B) + P(\overline{A} \cup \overline{B}) - 1 = P(A \cup B) + P(\overline{A\,B}) - 1 \ge 0,
\]
故选项 B、D 不正确。又 \(AB \subset B\)，有 \(P(AB) \le P(B)\)，从而
\[
P(A-B) = P(A) - P(AB) \ge P(A) - P(B),
\]
知选项 C 不正确。
\end{solution}

\begin{exercise}\label{exer:1-10}
设随机事件 \(A, B\) 满足 \(P(A)=P(B)=\dfrac{1}{2}\) 和 \(P(A \cup B)=1\)，则（）。\par
A. \(A \cup B = \Omega\)\quad
B. \(AB = \varnothing\)\quad
C. \(P(\overline{A} \cup \overline{B})=1\)\quad
D. \(P(A-B)=0\)
\end{exercise}

\begin{solution}
应选 C。

先根据加法公式
\[
P(A\cup B)=P(A)+P(B)-P(AB)
\]
代入题设 \(P(A)=P(B)=\dfrac12,\ P(A\cup B)=1\)，得
\[
1=\frac12+\frac12-P(AB)\implies P(AB)=0.
\]

现在逐项判断。

对 A，\(P(A\cup B)=1\) 只说明事件 \(A\cup B\) 的概率为 1，并不能推出集合恒等式 \(A\cup B=\Omega\)。概率为 1 的事件不一定就等于整个样本空间，所以 A 不对。

对 B，同理，\(P(AB)=0\) 只说明交事件是零概率事件，不能必然推出 \(AB=\varnothing\)。零概率事件在连续型情形下完全可能非空，所以 B 不对。

对 D，由
\[
P(A-B)=P(A)-P(AB)=\frac12-0=\frac12\neq 0,
\]
故 D 不对。

最后看 C。利用德摩根律，\(\overline{A}\cup\overline{B}=\overline{AB}\)，于是
\[
P(\overline{A}\cup\overline{B})=P(\overline{AB})=1-P(AB)=1-0=1.
\]

因此正确选项是 C。
\end{solution}

\begin{exercise}\label{exer:1-11}
设 \(A, B, C\) 是三个随机事件，\(A\) 与 \(C\) 互斥，\(P(AB)=\dfrac{1}{2}\)，\(P(C)=\dfrac{1}{3}\)，则 \(P(AB|\overline{C})=\underline{\hspace{2cm}}\)。
\end{exercise}

\begin{solution}
应填 \(\dfrac34\)。

要求 \(P(AB\mid \overline C)\)，先写条件概率公式：
\[
P(AB\mid \overline C)=\frac{P(AB\overline C)}{P(\overline C)}.
\]

先看分母。已知 \(P(C)=\dfrac13\)，所以 \(P(\overline C)=1-\frac13=\frac23\)。

再看分子。题目告诉我们 \(A\) 与 \(C\) 互斥，所以 \(AC=\varnothing\)。而 \(ABC\subset AC\)，因此 \(ABC=\varnothing,\ P(ABC)=0\)。

于是
\[
P(AB\overline C)=P(AB)-P(ABC)=\frac12-0=\frac12.
\]

代入条件概率公式，
\[
P(AB\mid \overline C)=\frac{1/2}{2/3}=\frac34.
\]
\end{solution}

\begin{note}
由于 \(AC = \varnothing\)，则 \(ABC = \varnothing \cdot B = \varnothing\)，得到 \(P(ABC)=0\)。
\end{note}

\begin{exercise}\label{exer:1-12}
已知 \(P(A)=\dfrac{1}{4}\)，\(P(B|A)=\dfrac{1}{3}\)，\(P(A|B)=\dfrac{1}{2}\)，则 \(P(A \cup B)=\underline{\hspace{2cm}}\)。
\end{exercise}
\begin{solution}
应填 \(\dfrac{1}{3}\)。先用乘法公式把交事件算出来，再反推 \(B\)：
\[
P(AB) = P(A)\,P(B|A) = \frac{1}{4} \times \frac{1}{3} = \frac{1}{12}, \quad P(B) = \frac{P(AB)}{P(A|B)} = \frac{1/12}{1/2} = \frac{1}{6},
\]
所以
\[
P(A \cup B) = P(A) + P(B) - P(AB) = \frac{1}{4} + \frac{1}{6} - \frac{1}{12} = \frac{1}{3}.
\]
\end{solution}

\begin{exercise}\label{exer:1-13}
某人给四位亲友各写一封信，然后随机地装入 4 个写好地址的信封中，且每个信封装一封信。求：
\par
(1) 4 封信都装对的概率；\par
(2) 4 封信都装错的概率。
\end{exercise}
\begin{solution}
设事件 \(A_i\ (i=1,2,3,4)\) 为"第 \(i\) 封信装对了"，\(A\) 为"4 封信都装对了"，\(B\) 为"4 封信都装错了"。\\
(1) 利用乘法公式：
\[
\begin{aligned}
P(A) &= P(A_1 A_2 A_3 A_4) = P(A_1)\,P(A_2|A_1)\,P(A_3|A_1 A_2)\,P(A_4|A_1 A_2 A_3) \\
&= \frac{1}{4} \times \frac{1}{3} \times \frac{1}{2} \times \frac{1}{1} = \frac{1}{24}.
\end{aligned}
\]
(2) 利用逆向思维和对偶律：
\[
\begin{aligned}
P(B) &= P(\overline{A_1}\,\overline{A_2}\,\overline{A_3}\,\overline{A_4}) = 1 - P(A_1 \cup A_2 \cup A_3 \cup A_4) \\
&= 1 - \sum_{i=1}^{4} P(A_i) + \sum_{1 \le i < j \le 4} P(A_i A_j) - \sum_{1 \le i < j < k \le 4} P(A_i A_j A_k) + P(A_1 A_2 A_3 A_4),
\end{aligned}
\]
其中 \(P(A_i)=\dfrac{1}{4}\)，\(P(A_i A_j)=\dfrac{1}{4} \times \dfrac{1}{3}=\dfrac{1}{12}\)，\(P(A_i A_j A_k)=\dfrac{1}{4} \times \dfrac{1}{3} \times \dfrac{1}{2}=\dfrac{1}{24}\)，于是
\[
P(B) = 1 - 4 \times \frac{1}{4} + 6 \times \frac{1}{12} - 4 \times \frac{1}{24} + \frac{1}{24} = \frac{3}{8}.
\]
\end{solution}

\begin{note}
装信过程可看作无放回抽取信件，依次装入信封。第 (2) 问直接计算较为困难，采用逆向思维方法更为简捷。
\end{note}

\begin{exercise}\label{exer:1-14}
从 \(1, 2, 3, 4\) 中任取一个数，记为 \(X\)，再从 \(1, \dots, X\) 中任取一个数，记为 \(Y\)，则 \(P\{Y=2\}=\underline{\hspace{2cm}}\)。
\end{exercise}

\begin{solution}
应填 \(\dfrac{13}{48}\)。

由全概率公式：
\[
\begin{aligned}
P\{Y=2\} &= \sum_{x=1}^{4} P\{X=x\}\,P\{Y=2|X=x\} \\
&= \frac{1}{4} \times 0 + \frac{1}{4} \times \frac{1}{2} + \frac{1}{4} \times \frac{1}{3} + \frac{1}{4} \times \frac{1}{4} = \frac{13}{48}.
\end{aligned}
\]
\end{solution}


\begin{exercise}\label{exer:1-15}
一批产品有 10 个正品和 2 个次品，任意抽取两次，每次抽一个，抽出后不再放回，则第二次抽出的是次品的概率为 \(\underline{\hspace{2cm}}\)。
\end{exercise}

\begin{solution}
应填 \(\dfrac16\)。

这道题可以用全概率公式来做，也可以把它看成抓阄模型。这里先按全概率公式完整写一遍。

记
\[
A_1=\{\text{第一次抽到次品}\},\qquad A_2=\{\text{第二次抽到次品}\}.
\]
按第一次的结果分类，有
\[
P(A_2)=P(A_1)P(A_2\mid A_1)+P(\overline{A_1})P(A_2\mid \overline{A_1}).
\]

各项分别为：
\[
P(A_1)=\frac{2}{12},\qquad P(\overline{A_1})=\frac{10}{12}.
\]
若第一次已经抽到次品，则剩下 \(11\) 件中只有 \(1\) 件次品，所以
\[
P(A_2\mid A_1)=\frac{1}{11}.
\]
若第一次抽到正品，则剩下 \(11\) 件中还有 \(2\) 件次品，所以
\[
P(A_2\mid \overline{A_1})=\frac{2}{11}.
\]

代入可得
\[
P(A_2)=\frac{2}{12}\cdot\frac{1}{11}+\frac{10}{12}\cdot\frac{2}{11}
=\frac{2+20}{132}
=\frac{22}{132}
=\frac16.
\]

所以第二次抽到次品的概率为 \(\dfrac16\)。
\end{solution}

\begin{note}
此例为\textbf{抓阄模型}：无放回抽取时，每次取到次品的概率不变，均为 \(\dfrac{2}{12} = \dfrac{1}{6}\)。从第 1 次到第 12 次取到次品的概率均为 \(\dfrac{1}{6}\)。
\end{note}

\begin{exercise}\label{exer:1-16}
设有两批数量相同的零件，已知有一批产品全部合格，另一批产品有 25\% 不合格。从两批产品中任取 1 只，经检验是正品，放回原处，并在原所在批次再取 1 只，求这只产品是次品的概率。
\end{exercise}

\begin{solution}
设 \(H_i\ (i=1,2)\) 为"第一次从第 \(i\) 批产品中抽取"，\(A\) 为"所取产品为正品"。

第一次抽取：
\[
P(H_1)=P(H_2)=\frac{1}{2},\quad P(A|H_1)=1,\quad P(A|H_2)=\frac{3}{4},
\]
\[
P(A) = P(H_1)\,P(A|H_1) + P(H_2)\,P(A|H_2) = \frac{1}{2} \times 1 + \frac{1}{2} \times \frac{3}{4} = \frac{7}{8}.
\]

由贝叶斯公式更新概率：
\[
P(H_1|A) = \frac{P(H_1)\,P(A|H_1)}{P(A)} = \frac{4}{7}, \quad P(H_2|A) = \frac{3}{7}.
\]

第二次抽取：设所取为次品的概率为
\[
P(\overline{A}) = P(H_1|A) \cdot P(\overline{A}|H_1) + P(H_2|A) \cdot P(\overline{A}|H_2) = \frac{4}{7} \times 0 + \frac{3}{7} \times \frac{1}{4} = \frac{3}{28}.
\]
\end{solution}

\begin{note}
本题是全概率公式与贝叶斯公式的复合应用：第一次用全概率公式计算 \(P(A)\)，再用贝叶斯公式更新各原因的概率，最后第二次再用全概率公式计算结果。两次抽取的关键区别在于：第一次在完全不知情的情况下等可能地抽取；第二次在已知第一次取正品的条件下，各批次被抽到的概率已经发生了变化。
\end{note}

\begin{exercise}\label{exer:1-21}
设 \(P(\overline{A}) = 0.3\)，\(P(B) = 0.4\)，\(P(A\overline{B}) = 0.5\)，则 \(P(B|A\cup\overline{B}) = \underline{\hspace{2cm}}\)。
\end{exercise}

\begin{solution}
应填 \(0.25\)。

首先，\(P(A) = 1 - P(\overline{A}) = 0.7\)。根据减法公式：
\[P(AB) = P(A) - P(A\overline{B}) = 0.7 - 0.5 = 0.2.\]
根据加法公式：
\[P(A\cup\overline{B}) = P(A) + P(\overline{B}) - P(A\overline{B}) = 0.7 + 0.6 - 0.5 = 0.8.\]
利用分配律化简分子：
\[B \cap (A\cup\overline{B}) = (B\cap A)\cup(B\cap\overline{B}) = AB,\]
故
\[P(B|A\cup\overline{B}) = \frac{P(AB)}{P(A\cup\overline{B})} = \frac{0.2}{0.8} = 0.25.\]
\end{solution}

\begin{exercise}\label{exer:1-22}
设 \(A, B\) 为任意两事件且 \(A\subset B\)，\(P(B) > 0\)，则下列结论一定正确的是（\quad）。\par
A. \(P(A) \ge P(A|B)\)\quad
B. \(P(A) > P(A|B)\)\quad
C. \(P(A) < P(A|B)\)\quad
D. \(P(A) \le P(A|B)\)
\end{exercise}

\begin{solution}
应选 D。

由 \(A\subset B\)，可知 \(AB=A\)，从而
\[
P(A\mid B)=\frac{P(AB)}{P(B)}=\frac{P(A)}{P(B)}\ge P(A).
\]

这就说明一定成立的是 \(P(A)\le P(A\mid B)\)，
即选项 D。

注意这里不能保证严格大于，因为当 \(B=\Omega\) 或 \(P(B)=1\) 时，也可能出现 \(P(A\mid B)=P(A)\)。
所以 C 不一定对，D 才是“一定正确”的表述。
\end{solution}

\begin{exercise}\label{exer:1-23}
假如在数字通信中传送信号 0 与 1 的概率分别为 0.7 和 0.3。由于随机干扰，当传送信号 0 时接收到信号 0 的概率是 0.8；当传送信号 1 时接收到信号 1 的概率是 0.9。求：\par
(1) 接收到信号 0 的概率；\par
(2) 当接收到信号 0 时，传送的信号是 0 的概率。
\end{exercise}

\begin{solution}
设
\[
A=\{\text{传送信号 }0\},\qquad B=\{\text{接收到信号 }0\}.
\]
则
\[
P(A)=0.7,\qquad P(\overline A)=0.3,\qquad P(B|A)=0.8.
\]
题目还告诉我们：当传送信号 1 时接收到信号 1 的概率为 \(0.9\)，因此
\[
P(\overline B|\overline A)=0.9,
\]
从而
\[
P(B|\overline A)=1-0.9=0.1.
\]

(1) 先求“接收到 0”的总概率。这是典型的“由因求果”问题，应按传送的是 0 还是 1 两种原因分类，利用全概率公式：
\[
\begin{aligned}
P(B)
&=P(A)P(B|A)+P(\overline A)P(B|\overline A)\\
&=0.7\times 0.8+0.3\times 0.1\\
&=0.56+0.03\\
&=0.59.
\end{aligned}
\]

(2) 再求“已知接收到 0，原来传送的是 0 的概率”，即 \(P(A|B)\)。这是“执果寻因”问题，应用贝叶斯公式：
\[
P(A|B)=\frac{P(A)P(B|A)}{P(B)}.
\]
把第 (1) 问得到的 \(P(B)=0.59\) 代入，得
\[
P(A|B)=\frac{0.7\times 0.8}{0.59}=\frac{0.56}{0.59}=\frac{56}{59}.
\]

因此，接收到信号 0 的概率为 \(0.59\)，在接收到信号 0 的条件下原来传送信号 0 的概率为 \(\dfrac{56}{59}\)。
\end{solution}

\begin{note}
本题是全概率公式与贝叶斯公式的典型应用：传送信号为"因"，接收信号为"果"。第 (1) 问"由因求果"用全概率公式，第 (2) 问"执果寻因"用贝叶斯公式。
\end{note}

\begin{exercise}\label{exer:1-39}
甲、乙两人独立地对同一目标射击一次，其命中率分别为 0.6 和 0.7。已知目标被击中，则它是甲射中的概率为 \underline{\hspace{2cm}}。
\end{exercise}

\begin{solution}
应填 \(\dfrac{15}{22}\)。

设
\[
A=\{\text{甲射中}\},\qquad B=\{\text{乙射中}\},\qquad C=\{\text{目标被击中}\}.
\]
则
\[
P(A)=0.6,\qquad P(B)=0.7,\qquad C=A\cup B.
\]

先求事件 \(C\) 的概率。因为甲、乙独立射击，
\[
P(AB)=P(A)P(B)=0.6\times 0.7=0.42.
\]
再由加法公式，
\[
\begin{aligned}
P(C)
&=P(A\cup B)\\
&=P(A)+P(B)-P(AB)\\
&=0.6+0.7-0.42\\
&=0.88.
\end{aligned}
\]

现在要求的是“已知目标被击中，甲射中的概率”，即 \(P(A|C)\)。注意这里的“甲射中”包括“甲单独命中”和“甲乙同时命中”两种情形，因此 \(A\subset C\)，从而 \(P(AC)=P(A)\)。所以
\[
P(A|C)=\frac{P(AC)}{P(C)}=\frac{P(A)}{P(C)}=\frac{0.6}{0.88}=\frac{15}{22}.
\]

所以所求概率为 \(\dfrac{15}{22}\)。
\end{solution}

\begin{exercise}\label{exer:1-40}
口袋中有 1 个球，不知其颜色是黑还是白。现再往口袋中放入 1 个白球，然后从中任意取出 1 个，发现取出的是白球。求口袋中原来那个球是白球的概率。
\end{exercise}

\begin{solution}
应填 \(\dfrac{2}{3}\)。

设
\[
A=\{\text{取出的球是白球}\},\qquad B=\{\text{原来那个球是白球}\}.
\]
由于题目没有给出“原球更可能是黑还是白”的额外信息，默认黑、白两种可能等可能，即
\[
P(B)=P(\overline B)=\frac12.
\]

下面分别计算在两种“原因”下观察到白球的概率：

1. 若原来那个球是白球，即事件 \(B\) 发生，那么放入一个白球后袋中共有 2 个白球，所以
\(P(A|B)=1\)。

2. 若原来那个球是黑球，即事件 \(\overline B\) 发生，那么放入一个白球后袋中是 1 白 1 黑，因此
\(P(A|\overline B)=\dfrac12\)。

题目告诉我们结果 \(A\) 已经发生，现在要反推原因 \(B\) 的概率，应用贝叶斯公式：
\[
P(B|A)=\frac{P(B)P(A|B)}{P(B)P(A|B)+P(\overline B)P(A|\overline B)}.
\]
代入数据：
\[
\begin{aligned}
P(B|A)
&=\frac{\frac12\times 1}{\frac12\times 1+\frac12\times\frac12}\\
&=\frac{\frac12}{\frac12+\frac14}
=\frac{\frac12}{\frac34}
=\frac23.
\end{aligned}
\]

所以，在已经取出白球的条件下，原来那一个球是白球的概率为 \(\dfrac{2}{3}\)。
\end{solution}

\begin{exercise}\label{exer:1-41}
已知装有同种零件的产品两箱，第一箱内装 50 件，其中一等品 10 件；第二箱内装 30 件，其中一等品 18 件。现从两箱中任意挑选一箱，从中先后取出两件产品（取后不放回）。求：先取出的产品是一等品的概率 \(p\)；已知先取出的产品是一等品，则第二次取出的产品仍然是一等品的概率 \(q\)。
\end{exercise}

\begin{solution}
设 \(A_i\) 为"第 \(i\) 次取出一等品"（\(i=1,2\)），\(C_j\) 为"挑选第 \(j\) 箱"（\(j=1,2\)）。\(C_1\cup C_2=\Omega\)，\(C_1C_2=\varnothing\)，\(P(C_1)=P(C_2)=\dfrac{1}{2}\)。

由全概率公式：
\begin{align*}
p = P(A_1) &= P(C_1)\,P(A_1|C_1)+P(C_2)\,P(A_1|C_2) = \frac{1}{2}\times\frac{10}{50}+\frac{1}{2}\times\frac{18}{30} = \frac{2}{5}, \\
P(A_1A_2) &= P(C_1)\,P(A_1A_2|C_1)+P(C_2)\,P(A_1A_2|C_2) \\
&= \frac{1}{2}\times\frac{10\times 9}{50\times 49}+\frac{1}{2}\times\frac{18\times 17}{30\times 29} = \frac{276}{1421}.
\end{align*}
故 \(q = P(A_2|A_1) = \dfrac{P(A_1A_2)}{P(A_1)} = \dfrac{276}{1421}\times\dfrac{5}{2} = \dfrac{690}{1421}\)。
\end{solution}

\begin{exercise}\label{exer:1-42}
每箱产品有 10 件，其次品数从 0 到 2 是等可能的。开箱检验时从中任取 1 件，检验出次品则拒收。由于检验有误，将正品误认为次品的概率为 2\%，次品被漏查的概率为 5\%。求该箱产品通过验收的概率。
\end{exercise}

\begin{solution}
应填 \(0.887\)。

设 \(A_i\)（\(i=0,1,2\)）为"箱中有 \(i\) 件次品"，\(B\) 为"通过验收"。\(P(A_i)=\dfrac{1}{3}\)。

取到正品时通过验收的概率为 0.98，取到次品时通过验收的概率为 0.05。由全概率公式（两层）：
\begin{align*}
P(B|A_0) &= 1\times 0.98 = 0.98, \\
P(B|A_1) &= \frac{9}{10}\times 0.98 + \frac{1}{10}\times 0.05 = 0.887, \\
P(B|A_2) &= \frac{8}{10}\times 0.98 + \frac{2}{10}\times 0.05 = 0.794, \\
P(B) &= \frac{1}{3}(0.98+0.887+0.794) = \frac{2.661}{3} = 0.887.
\end{align*}
\end{solution}

\begin{exercise}\label{exer:1-52}
设水杯成箱出售，每箱 12 个，每箱中含 0 个、1 个、2 个残次品的概率分别为 0.7、0.19、0.11。某人欲购一箱，售货员随机取一箱，从中任取 3 只查看，若无残次品则购买。\par
(1) 此人购买一箱的概率；\par
(2) 在此人购买一箱的条件下，该箱无残次品的概率。
\end{exercise}

\begin{solution}
设 \(A_i\)（\(i=0,1,2\)）为"箱中有 \(i\) 件残次品"，\(B\) 为"购买此箱"（即取出的 3 只均无残次品）。

(1) 由全概率公式：
\begin{align*}
P(B) &= \sum_{i=0}^{2}P(A_i)\,P(B|A_i), \\
P(B|A_0) &= 1, \quad P(B|A_1) = \frac{\mathrm{C}_{11}^3}{\mathrm{C}_{12}^3} = \frac{3}{4}, \quad P(B|A_2) = \frac{\mathrm{C}_{10}^3}{\mathrm{C}_{12}^3} = \frac{6}{11}, \\
P(B) &= 0.7\times 1 + 0.19\times\frac{3}{4} + 0.11\times\frac{6}{11} = 0.9025.
\end{align*}

(2) 由贝叶斯公式：
\[
P(A_0|B) = \frac{P(A_0)\,P(B|A_0)}{P(B)} = \frac{0.7\times 1}{0.9025} = \frac{280}{361} \approx 0.7756.
\]
\end{solution}

\begin{exercise}\label{exer:1-53}
统计资料显示，在出口罐头导致的索赔事件中，有 50\% 是质量问题，30\% 是数量问题，20\% 是包装问题。在这三类问题各引起的索赔事件中，不经过法律诉讼而通过协商解决的概率分别为 40\%、60\% 和 75\%。\par
(1) 索赔事件通过协商解决的概率；\par
(2) 已知一索赔事件通过协商解决，求该事件不是包装问题的概率。
\end{exercise}

\begin{solution}
设 \(A_1,A_2,A_3\) 分别为质量、数量、包装问题引起的索赔，\(B\) 为"通过协商解决"。

(1) 由全概率公式：
\[
P(B) = \sum_{i=1}^{3}P(A_i)\,P(B|A_i) = 0.5\times 0.40 + 0.3\times 0.60 + 0.2\times 0.75 = 0.53.
\]

(2) 由贝叶斯公式：
\[
P(A_3|B) = \frac{P(A_3)\,P(B|A_3)}{P(B)} = \frac{0.2\times 0.75}{0.53} = \frac{15}{53},
\]
故所求概率为 \(1-P(A_3|B) = \dfrac{38}{53}\)。
\end{solution}

\begin{exercise}\label{exer:1-54}
在某城市中，下雨和晴天的天数各占一半，而天气无论在雨天还是晴天都有 \(\dfrac{2}{3}\) 的概率预报准确。如果预报下雨，王明同学就一定带雨伞；如果预报无雨，则不带伞。设 \(A=\{\text{天下雨}\}\)，\(B=\{\text{预报有雨}\}\)，\(C=\{\text{王明带雨伞}\}\)。\par
(1) 事件 \(\overline{A}\cap\overline{B}\cap C\)、\(\overline{A}\cap B\cap C\) 的含义是什么？哪个为不可能事件？\par
(2) 求他带雨伞而没有下雨的概率。
\end{exercise}

\begin{solution}
由题意：\(P(A)=\dfrac{1}{2}\)，\(P(B|A)=\dfrac{2}{3}\)（预报准确），\(P(\overline{B}|\overline{A})=\dfrac{2}{3}\)（预报准确），故 \(P(B|\overline{A})=\dfrac{1}{3}\)。又 \(P(C|B)=1\)（预报有雨必带伞），\(P(C|\overline{B})=0\)（预报无雨不带伞）。

(1) \(\overline{A}\cap\overline{B}\cap C\)：天没下雨、预报无雨、却带了雨伞。由于 \(P(C|\overline{B})=0\)，这是一个\textbf{不可能事件}。\par
\(\overline{A}\cap B\cap C\)：天没下雨、预报有雨、带了雨伞。这是可能事件。

(2) "带雨伞而没有下雨"即 \(\overline{A}\cap C\)。由于 \(C\) 完全由 \(B\) 决定，\(\overline{A}\cap\overline{B}\cap C\) 为不可能事件，故
\[
P(\overline{A}\cap C) = P(\overline{A}\cap B\cap C) = P(\overline{A})\,P(B|\overline{A})\,P(C|\overline{A}B) = \frac{1}{2}\times\frac{1}{3}\times 1 = \frac{1}{6}.
\]
\end{solution}


\begin{exercise}[提高]\label{exer:1-55}
设某同学的手机在一天内收到短信数服从泊松分布 \(P(\lambda)\)，每个短信是否为垃圾短信与其到达时间独立，也与其他短信是否为垃圾短信相互独立。假设每个短信是垃圾短信的概率为 \(p\)。\par
(1) 已知一天内收到了 \(n\) 条短信，求其中恰有 \(k\) 条垃圾短信的概率（\(0\le k\le n\)）。\par
(2) 求一天内收到 \(k\) 条垃圾短信的概率（\(k=0,1,2,\dots\)）。
\end{exercise}

\begin{solution}
(1) 记 \(X\) 为一天内收到短信的总数，\(Y\) 为其中垃圾短信的条数。已知 \(X=n\) 时，每条短信独立地以概率 \(p\) 为垃圾短信，故
\[
P(Y=k\mid X=n) = \mathrm{C}_n^k\,p^k\,(1-p)^{n-k}.
\]

(2) 由全概率公式：
\begin{align*}
P(Y=k) &= \sum_{n=k}^{\infty} P(Y=k\mid X=n)\,P(X=n) = \sum_{n=k}^{\infty} \mathrm{C}_n^k\,p^k\,(1-p)^{n-k}\,\frac{\lambda^n}{n!}\,e^{-\lambda} \\
&= \frac{p^k\,e^{-\lambda}\,\lambda^k}{k!}\sum_{n=k}^{\infty}\frac{[\lambda(1-p)]^{n-k}}{(n-k)!} = \frac{(\lambda p)^k}{k!}\,e^{-\lambda}\,e^{\lambda(1-p)} = \frac{(\lambda p)^k}{k!}\,e^{-\lambda p}.
\end{align*}
即 \(Y\sim P(\lambda p)\)。
\end{solution}

\begin{note}
泊松分布的分解定理（稀疏性）：以概率 \(p\) 独立地对泊松过程 \(P(\lambda)\) 进行"稀疏"，得到的新过程仍为泊松分布，参数变为 \(\lambda p\)。本题是全概率公式与泊松分布结合的典型应用。
\end{note}


\begin{exercise}[提高]
某工厂进行产品质量盲测。现有外观相同的 \(A,B\) 两个初始样本盒，\(A\) 盒装有质量指标 \(1,5,10\) 的产品各一件，\(B\) 盒装有指标 \(2,5,8\) 的产品各一件。甲、乙两名质检员先后从待选盒中随机选择一盒，再从中随机抽取一件产品比对。\par
(1) 求甲所抽产品质量指标的数学期望；\par
(2) 现新增 \(n\) 个均含 3 件产品的新样本盒，设指标 5 的产品为标准件。甲先从这 \(n+2\) 个盒中随机选择一个，乙再从剩下的 \(n+1\) 个盒中随机选择一个抽取。\par
记事件 \(E_1\) 为"甲抽到标准件"，事件 \(E_2\) 为"乙抽到标准件"。\par
i) 记新增样本盒中的标准件总个数为 \(m\)，求事件 \(E_1\) 与事件 \(E_2\) 发生的概率；\par
ii) 设事件 \(E\) 为"甲选择 \(A\) 盒"。证明："事件 \(E_1,E_2,E\) 两两独立"当且仅当"每个新盒中恰有 1 件标准件"，并简述该独立性条件在盲检中的实际意义。
\end{exercise}

\begin{solution}
(1) 使用全概率公式。甲抽到 \(A\) 盒的概率为 \(\frac{1}{2}\)，抽到 \(B\) 盒的概率也为 \(\frac{1}{2}\)。若抽到 \(A\) 盒：
\[E(X|A) = \frac{1+5+10}{3} = \frac{16}{3}.\]
若抽到 \(B\) 盒：
\[E(X|B) = \frac{2+5+8}{3} = 5.\]
总期望
\[E(X) = P(A)\,E(X|A) + P(B)\,E(X|B) = \frac{1}{2}\cdot\frac{16}{3} + \frac{1}{2}\cdot 5 = \frac{31}{6}.\]

(2) 新增盒子后的抽样分析。

新增 \(n\) 个盒，加上原本的 \(A,B\) 盒，共 \(n+2\) 盒。原 \(A\) 盒有 1 件标准件（指标 5），原 \(B\) 盒有 1 件标准件，故新增前有 2 件标准件。已知新增盒子中共有 \(m\) 件标准件，因此总体中共有 \(m+2\) 件标准件。产品总数为 \(3(n+2)\) 件。

i) 这本质上是经典的"抽签模型"。甲在所有产品中随机抽一件，抽到标准件的概率等于标准件在总体中的比例：
\[P(E_1) = \frac{m+2}{3(n+2)}.\]
根据抽签原理（不放回抽样中，只要不知道前面抽取的结果，后面抽取的概率分布不变），乙抽到标准件的概率与甲完全相同：
\[P(E_2) = \frac{m+2}{3(n+2)}.\]

ii) 独立性证明与实际意义。

\textbf{第一步：}利用 \((E_1,E)\) 独立，求出 \(m\) 与 \(n\) 的关系。

若三者两两独立，则必有 \(P(E_1|E) = P(E_1)\)。已知 \(P(E) = \frac{1}{n+2}\)。在事件 \(E\) 发生（甲抽了 \(A\) 盒）的前提下，\(A\) 盒里只有 1 件标准件，所以
\[P(E_1|E) = \frac{1}{3}.\]
联立得
\[\frac{m+2}{3(n+2)} = \frac{1}{3} \implies m+2 = n+2 \implies m = n,\]
即 \(n\) 个新盒中总共含有 \(n\) 件标准件。

\textbf{第二步：}利用 \((E_1,E_2)\) 独立，求证每盒标准件的分布。

设第 \(i\) 个盒子（\(i=1,2,\dots,n+2\)）含有 \(s_i\) 件标准件，显然 \(s_i \in \{0,1,2,3\}\)，且 \(\sum_{i=1}^{n+2} s_i = m+2 = n+2\)。

计算 \(P(E_1 E_2)\)：甲抽到第 \(i\) 盒、乙抽到第 \(j\) 盒（\(j \neq i\)）的概率为 \(\frac{1}{n+2}\cdot\frac{1}{n+1}\)；在各自盒子抽中标准件的概率分别为 \(\frac{s_i}{3}\) 和 \(\frac{s_j}{3}\)。于是
\[
P(E_1 E_2) = \sum_{i=1}^{n+2}\sum_{j \neq i} \frac{1}{(n+2)(n+1)}\cdot\frac{s_i}{3}\cdot\frac{s_j}{3}
= \frac{1}{9(n+1)(n+2)}\sum_{i=1}^{n+2}\sum_{j \neq i}s_is_j.
\]
由代数恒等式 \(\displaystyle\sum_{i}\sum_{j \neq i}s_is_j = \bigl(\sum s_i\bigr)^2 - \sum s_i^2 = (n+2)^2 - \sum s_i^2\)，代入得
\[
P(E_1 E_2) = \frac{(n+2)^2 - \sum s_i^2}{9(n+1)(n+2)}.
\]
因 \(E_1, E_2\) 独立，\(P(E_1 E_2) = P(E_1)P(E_2) = \frac{1}{3}\times\frac{1}{3} = \frac{1}{9}\)，建立方程
\[
\frac{(n+2)^2 - \sum s_i^2}{(n+1)(n+2)} = 1,
\]
化简得 \(\displaystyle\sum s_i^2 = (n+2)^2 - (n+1)(n+2) = n+2\)。

于是得到方程组
\[
\begin{cases}
\displaystyle\sum s_i = n+2, \\[6pt]
\displaystyle\sum s_i^2 = n+2.
\end{cases}
\]
因为 \(s_i\) 都是非负整数，\(\sum(s_i^2-s_i)=0 \implies \sum s_i(s_i-1)=0\)。由于 \(s_i(s_i-1) \ge 0\) 恒成立，每一项都必须为 0，即 \(s_i = 0\) 或 \(1\)。但如果存在某个 \(s_k = 0\)，则必须有其他 \(s_l \ge 2\) 才能凑够总和 \(n+2\)，这与 \(s_l(s_l-1)=0\) 矛盾。因此唯一解是：对所有的 \(i\)，\(s_i = 1\)，即每个盒子中恰好有 1 件标准件。\hfill\(\square\)

\textbf{实际意义：}在工业质量盲检中，如果各批次间的良品率存在差异，那么"抽到了哪个批次"或"前一个人的抽样结果"就会成为一种先验信息，影响对剩余批次的概率判断（即抽样不再独立）。本结论表明：只有当每一批次的质量结构绝对均等时，盲选抽检才能实现统计学上的真正独立。这也解释了为什么在进行产品质量双盲测试时，必须保证各组对照样本在结构上高度一致。
\end{solution}


\section{事件的独立性}

\subsection{独立性的定义}

\begin{definition}[事件的独立性] \label{def:independence}
设 \(A, B\) 为两个事件，如果
\[
P(AB) = P(A)\,P(B),
\]
则称事件 \(A\) 与 \(B\) \textbf{相互独立}，简称 \(A\) 与 \(B\) 独立。
\end{definition}
顺带剧透一下，我们后面会讨论多个事件的独立性概念。
\begin{definition}[$n$ 个事件的相互独立] \label{def:n-event-independence}
设 \(A_1, A_2, \dots, A_n\) 为 \(n\ (n \ge 2)\) 个事件，如果对其中任意有限个事件 \(A_{i_1}, A_{i_2}, \dots, A_{i_k}\ (2 \le k \le n)\)，都有
\[
P(A_{i_1} A_{i_2} \cdots A_{i_k}) = P(A_{i_1})\,P(A_{i_2}) \cdots P(A_{i_k}),
\]
则称这 \(n\) 个事件\textbf{相互独立}。
\end{definition}

特别地，对于三个事件 \(A, B, C\)，相互独立要求同时满足以下四个等式：
\[
\begin{cases}
P(AB) = P(A)\,P(B), \\
P(AC) = P(A)\,P(C), \\
P(BC) = P(B)\,P(C), \\
P(ABC) = P(A)\,P(B)\,P(C).
\end{cases}
\]
若仅满足前三个等式（即任意两个事件独立），则称 \(A, B, C\) \textbf{两两独立}。

\begin{note}
从条件概率的角度理解独立性：若 \(A, B\) 相互独立且 \(P(B) > 0\)，则 \(P(A|B) = P(A)\)。这意味着样本空间缩小为 \(B\) 后，事件 \(A\) 发生的概率并没有改变——\(B\) 的发生对 \(A\) 没有任何影响，这正是"独立"一词的含义。
\end{note}

\subsection{独立性的判定}

\begin{property}\label{prop:independence-equiv}
\(A\) 与 \(B\) 独立（\(P(AB)=P(A)\,P(B)\)）的等价条件（设 \(0 < P(A) < 1\)）：
\begin{enumerate}
    \item \(A\) 与 \(B\) 独立 \(\Leftrightarrow\) \(A\) 与 \(\overline{B}\) 独立 \(\Leftrightarrow\) \(\overline{A}\) 与 \(B\) 独立 \(\Leftrightarrow\) \(\overline{A}\) 与 \(\overline{B}\) 独立；
    \item \(A\) 与 \(B\) 独立 \(\Leftrightarrow\) \(P(B|A) = P(B|\overline{A}) = P(B)\)；
    \item \(A\) 与 \(B\) 独立 \(\Leftrightarrow\) \(P(B|A) + P(\overline{B}|\overline{A}) = 1\)。
\end{enumerate}
\end{property}

\begin{remark}
(1) 关于等价条件 (1) 的证明：由 \(A, B\) 独立，有 \(P(AB)=P(A)\,P(B)\)，于是
\[
P(A\,\overline{B}) = P(A) - P(AB) = P(A) - P(A)\,P(B) = P(A)\,[1-P(B)] = P(A)\,P(\overline{B}),
\]
故 \(A, \overline{B}\) 独立，其余同理。\newline
(2) 推广：将相互独立的事件组中的任何几个事件换成各自的对立事件，所得的新事件组仍相互独立。
\end{remark}

\begin{property}\label{prop:independence-compute}
独立性的判定与运算：
\begin{enumerate}
    \item 若 \(P(A)=0\) 或 \(P(A)=1\)，则 \(A\) 与任意事件 \(B\) 相互独立。
    \item 若 \(P(A)>0, P(B)>0\)，且 \(A, B\) 互斥，则 \(A, B\) 一定\textbf{不}独立；反之，若 \(A, B\) 独立，则 \(A, B\) 一定\textbf{不}互斥（除非 \(P(A)=0\) 或 \(P(B)=0\)）。
    \item 若一组事件相互独立，则其中部分事件也相互独立；对不含相同事件的独立组分别作运算，所得新事件仍独立。例如 \(A, B, C, D\) 相互独立，则 \(AB\) 与 \(CD\) 独立，\(A\) 与 \(BC-D\) 独立。
\end{enumerate}
\end{property}

\begin{note}
常见误区：
\begin{enumerate}
    \item "独立性"是概率意义上的数学定义，与逻辑因果上的"独立"无关，只看 \(P(AB)=P(A)\,P(B)\) 是否成立。
    \item 两两独立 \(\nRightarrow\) 相互独立。三个事件两两独立只保证 \(P(AB)=P(A)\,P(B)\)、\(P(AC)=P(A)\,P(C)\)、\(P(BC)=P(B)\,P(C)\)，但不保证 \(P(ABC)=P(A)\,P(B)\,P(C)\)。
\end{enumerate}
\end{note}

\begin{exercise}\label{exer:1-17}
将一枚硬币独立地掷两次，设 \(A_1=\{\text{掷第一次出现正面}\}\)，\(A_2=\{\text{掷第二次出现正面}\}\)，\(A_3=\{\text{正反面各出现一次}\}\)，\(A_4=\{\text{正面出现两次}\}\)，则事件（\quad）。\par
\begin{tabular}{ll}
A. \(A_1,A_2,A_3\) 相互独立 & B. \(A_2,A_3,A_4\) 相互独立 \\[4pt]
C. \(A_1,A_2,A_3\) 两两独立 & D. \(A_2,A_3,A_4\) 两两独立
\end{tabular}
\end{exercise}

\begin{solution}
应选 C。

根据古典概型：
\[
P(A_1)=P(A_2)=P(A_3)=\frac{1}{2}, \quad P(A_1A_2)=P(A_1A_3)=P(A_2A_3)=\frac{1}{4},
\]
满足两两独立条件。但
\[
P(A_1A_2A_3)=0 \neq P(A_1)\,P(A_2)\,P(A_3)=\frac{1}{8},
\]
故 \(A_1, A_2, A_3\) 两两独立但不相互独立。

又 \(A_3 \cap A_4 = \varnothing\)，\(P(A_3A_4)=0 \neq P(A_3)\,P(A_4)=\frac{1}{8}\)，故 \(A_2, A_3, A_4\) 不两两独立。
\end{solution}

\begin{exercise}\label{exer:1-18}
设 \(A, B, C\) 为三个两两独立的随机事件，\(P(A)=P(B)=P(C)\)，\(P(A \cup B \cup C)=\dfrac{23}{25}\)，\(P(\overline{A} \cup \overline{B} \cup \overline{C})=\dfrac{14}{25}\)，求 \(P(A)\)。
\end{exercise}

\begin{solution}
由两两独立性，
\[
\begin{aligned}
P(\overline{A} \cup \overline{B} \cup \overline{C}) &= P(\overline{A})+P(\overline{B})+P(\overline{C}) - P(\overline{A}\,\overline{B}) - P(\overline{B}\,\overline{C}) - P(\overline{A}\,\overline{C}) + P(\overline{A}\,\overline{B}\,\overline{C}) \\
&= 3P(\overline{A}) - 3[P(\overline{A})]^2 + P(\overline{A}\,\overline{B}\,\overline{C}).
\end{aligned}
\]
代入数据：
\[
3P(\overline{A}) - 3[P(\overline{A})]^2 + \frac{2}{25} = \frac{14}{25} \implies 3P(\overline{A}) - 3[P(\overline{A})]^2 = \frac{12}{25}.
\]
令 \(x = P(\overline{A})\)，则 \(3x - 3x^2 = \dfrac{12}{25}\)，即 \(25x^2 - 25x + 4 = 0\)，解得 \(x = \dfrac{1}{5}\) 或 \(\dfrac{4}{5}\)。

又 \(P(\overline{A}) \le P(\overline{A} \cup \overline{B} \cup \overline{C}) = \dfrac{14}{25}\)，故 \(P(\overline{A}) = \dfrac{1}{5}\)，即 \(P(A) = \dfrac{4}{5}\)。
\end{solution}

\begin{exercise}\label{exer:1-43}
设 \(A,B,C\) 为三个随机事件，且 \(A\) 与 \(C\) 相互独立，\(B\) 与 \(C\) 相互独立，则 \(A\cup B\) 与 \(C\) 相互独立的充分必要条件是（\quad）\footnote{验证独立性不要看直觉，会被骗，还是要验证定义}。\par
A. \(A\) 与 \(B\) 相互独立\quad
B. \(A\) 与 \(B\) 互不相容\quad
C. \(AB\) 与 \(C\) 相互独立\quad
D. \(AB\) 与 \(C\) 互不相容
\end{exercise}

\begin{solution}
应选 C。

\(A\cup B\) 与 \(C\) 相互独立，即 \(P[(A\cup B)C] = P(A\cup B)\,P(C)\)。
\[
\begin{aligned}
P[(A\cup B)C] &= P(AC\cup BC) = P(AC)+P(BC)-P(ABC) \\
&= P(A)\,P(C)+P(B)\,P(C)-P(ABC), \\
P(A\cup B)\,P(C) &= [P(A)+P(B)-P(AB)]\,P(C) \\
&= P(A)\,P(C)+P(B)\,P(C)-P(AB)\,P(C),
\end{aligned}
\]
因此充要条件为 \(P(ABC)=P(AB)\,P(C)\)，即 \(AB\) 与 \(C\) 相互独立。
\end{solution}

\begin{exercise}\label{exer:1-56}
已知 \(A,B\) 相互独立，\(P(A)=0.5\)，\(P(B)=0.4\)，\(C=\overline{A}\cup B\)。求 \(P(\overline{C})\)，\(P(C\mid(A\cup B))\)。
\end{exercise}

\begin{solution}
\(\overline{C}=\overline{\overline{A}\cup B}=A\,\overline{B}\)，由独立性：
\[
P(\overline{C}) = P(A\,\overline{B}) = P(A)(1-P(B)) = 0.5\times 0.6 = 0.3.
\]

又 \(P(A\cup B)=P(A)+P(B)-P(A)P(B)=0.5+0.4-0.2=0.7\)。注意到
\[
C\cap(A\cup B) = (\overline{A}\cup B)\cap(A\cup B) = B,
\]
故
\[
P(C\mid(A\cup B)) = \frac{P(B)}{P(A\cup B)} = \frac{0.4}{0.7} = \frac{4}{7}.
\]
\end{solution}

\begin{exercise}\label{exer:1-57}
下列命题中，正确的是（\quad）。\par
A. 若 \(A,B\) 互不相容，则 \(\overline{A},\overline{B}\) 也互不相容\quad
B. 若 \(A,B\) 相容，则 \(\overline{A},\overline{B}\) 也相容\quad
C. 若 \(A,B\) 独立，则 \(\overline{A},\overline{B}\) 也独立\quad
D. 若 \(A,B\) 独立，则 \(\overline{A},\overline{B}\) 不相容
\end{exercise}

\begin{solution}
应选 C。

由独立性 \(P(AB)=P(A)P(B)\)：
\[
P(\overline{A}\,\overline{B}) = P(\overline{A\cup B}) = 1-P(A)-P(B)+P(AB) = (1-P(A))(1-P(B)) = P(\overline{A})P(\overline{B}),
\]
故 \(\overline{A},\overline{B}\) 也独立。

A 错误：例如 \(A=\varnothing\) 时 \(\overline{A}=\Omega\)，\(\overline{A}\,\overline{B}=\overline{B}\) 一般不为空。B 错误：类似地，\(\overline{A},\overline{B}\) 可相容也可不相容。D 错误：独立与相容没有必然关系。
\end{solution}


\begin{exercise}[提高]\label{exer:1-58}
设 \(A,B,C\) 相互独立，且 \(P(A)\neq 0\)，\(0<P(C)<1\)。则下面四个事件中不独立的是（\quad）。\par
A. \(\overline{A}\cup B\) 与 \(C\)\quad
B. \(\overline{AC}\) 与 \(\overline{C}\)\quad
C. \(\overline{A-B}\) 与 \(\overline{C}\)\quad
D. \(\overline{AB}\) 与 \(\overline{C}\)
\end{exercise}

\begin{solution}
应选 B。

A、C、D 中左边的事件均由 \(A,B\) 构成，而 \(A,B\) 各自与 \(C\) 独立，因此由 \(A,B\) 经过并、交、差运算构成的新事件也与 \(C\) 独立。

B 中，\(\overline{AC}\cap\overline{C}=\overline{C}\)（因为 \(\overline{AC}\supset\overline{C}\)）。若两者独立，需 \(P(\overline{AC})\,P(\overline{C})=P(\overline{C})\)，即 \(P(\overline{AC})=1\)。但
\[
P(\overline{AC}) = P(\overline{A})+P(\overline{C})-P(\overline{A})P(\overline{C}) < 1 \quad (\text{因为 }P(\overline{A})<1, P(\overline{C})<1),
\]
矛盾。故 \(\overline{AC}\) 与 \(\overline{C}\) 不独立。

等价地，若独立则需 \(P(\overline{C})=P(\overline{AC}\,\overline{C})=P(\overline{C})\) 恒成立，但由独立性 \(P(AC)=P(A)P(C)>0\)，可得 \(\overline{AC}\neq\Omega\)，从而 \(P(\overline{AC})<1\)。
\end{solution}

\section{独立重复试验与伯努利概型}

\begin{definition}[\(n\) 重伯努利概型] \label{def:bernoulli}
如果随机试验满足以下条件：
\begin{enumerate}
    \item 每次试验只有两个可能结果，记为 \(A\)（成功）和 \(\overline{A}\)（失败）；
    \item 每次试验中 \(P(A)=p\)（\(0 < p < 1\)）保持不变；
    \item 各次试验相互独立，
\end{enumerate}
则将这种试验独立重复进行 \(n\) 次的概率模型称为 \(\boldsymbol{n}\) \textbf{重伯努利概型}。
\end{definition}

在 \(n\) 重伯努利概型中，事件 \(A\) 恰好发生 \(k\) 次的概率为
\[
P\{A \text{ 恰好发生 } k \text{ 次}\} = \mathrm{C}_n^k\, p^k\, (1-p)^{n-k}, \quad k=0,1,\dots,n.
\]

如果用 \(X\) 表示 \(n\) 重伯努利概型中事件 \(A\) 发生的次数，则 \(X\) 服从参数为 \(n, p\) 的\textbf{二项分布}，记为 \(X \sim B(n,p)\)。

\begin{note}
要善于判定 \(n\) 重伯努利概型：题目中出现“将……独立重复进行 \(n\) 次”“对……重复观察 \(n\) 次”等字样，或可以转化为 \(n\) 次独立重复试验的问题，都应考虑使用二项分布计算概率。例如射击问题、抛硬币问题、产品抽样检验问题等。
\end{note}

\begin{exercise}\label{exer:1-19}
对同一目标接连进行 3 次独立重复射击，假设至少命中目标一次的概率为 \(\dfrac{7}{8}\)，则每次射击命中目标的概率 \(p=\underline{\hspace{2cm}}\)。
\end{exercise}

\begin{solution}
应填 \(\dfrac{1}{2}\)。

先把“至少命中一次”转成“全都不中”的补事件，这样最省事。

记 \(A_i=\{\text{第 }i\text{ 次命中}\}\ (i=1,2,3)\)，\(A_1, A_2, A_3\) 相互独立且 \(P(A_i)=p\)。由
\[
P(A_1 \cup A_2 \cup A_3) = 1 - P(\overline{A_1}\,\overline{A_2}\,\overline{A_3}) = 1 - (1-p)^3 = \frac{7}{8},
\]
得 \((1-p)^3 = \dfrac{1}{8}\)，故 \(1-p = \dfrac{1}{2}\)，\(p = \dfrac{1}{2}\)。
\end{solution}

\begin{exercise}\label{exer:1-20}
某人向同一目标独立重复射击，每次命中目标的概率为 \(p\ (0<p<1)\)，则此人第 4 次射击恰好第 2 次命中目标的概率为（\quad）。\par
A. \(3p(1-p)^2\)\quad
B. \(6p(1-p)^2\)\quad
C. \(3p^2(1-p)^2\)\quad
D. \(6p^2(1-p)^2\)
\end{exercise}

\begin{solution}
应选 C。

"第 4 次恰好第 2 次命中"意味着：前 3 次中有 1 次命中，第 4 次命中。设前 3 次命中次数为 \(X \sim B(3,p)\)，则
\[
P\{X=1\} = \mathrm{C}_3^1\, p\,(1-p)^2 = 3p(1-p)^2.
\]
第 4 次命中的概率为 \(p\)，故所求概率为
\[
3p(1-p)^2 \cdot p = 3p^2(1-p)^2.
\]
\end{solution}

\begin{note}
一般地，在独立重复试验中，第 \(n\) 次试验恰好第 \(k\) 次成功的概率为
\[
\mathrm{C}_{n-1}^{k-1}\, p^k\, (1-p)^{n-k},
\]
其中前 \(n-1\) 次中有 \(k-1\) 次成功，第 \(n\) 次成功。这实际上是负二项分布（帕斯卡分布）的一个特例。
\end{note}

\begin{exercise}\label{exer:1-25}
在四次独立试验中，事件 \(A\) 至少发生一次的概率为 \(\dfrac{80}{81}\)，则事件 \(A\) 在每次试验中发生的概率为（\quad）。\quad 
A. \(\dfrac{1}{3}\)\quad
B. \(\dfrac{2}{3}\)\quad
C. \(\dfrac{2}{3}\sqrt[4]{5}\)\quad
D. \(1 - \dfrac{2}{3}\sqrt[4]{5}\)
\end{exercise}

\begin{solution}
应选 B。

设 \(X\) 表示事件 \(A\) 在 4 次独立试验中发生的次数，则 \(X\sim B(4,p)\)。

题目给出“至少发生一次”的概率，所以先用对立事件来写：
\[
P\{X\ge 1\}=1-P\{X=0\}=1-(1-p)^4=\frac{80}{81}.
\]

移项得 \((1-p)^4=\frac{1}{81}=\left(\frac13\right)^4\)。由于 \(0\le 1-p\le 1\)，可得
\[
1-p=\frac13,\qquad p=1-\frac13=\frac23.
\]

因此应选 B。
\end{solution}

\begin{exercise}\label{exer:1-44}
三人独立地破译一个密码，他们能单独译出的概率分别为 \(\dfrac{1}{5},\dfrac{1}{3},\dfrac{1}{4}\)，则此密码被译出的概率为 \underline{\hspace{2cm}}。
\end{exercise}

\begin{solution}
应填 \(\dfrac{3}{5}\)。

记事件 \(A_i\) 为“第 \(i\) 个人译出密码”\((i=1,2,3)\)，则
\[
P(A_1)=\frac15,\qquad P(A_2)=\frac13,\qquad P(A_3)=\frac14.
\]

题目要求“密码被译出”的概率，也就是“至少有一个人译出”的概率。直接求并事件不太方便，最简洁的方法是求它的对立事件：
\[
B=A_1\cup A_2\cup A_3,
\]
\[
\overline{B}=\overline{A_1}\,\overline{A_2}\,\overline{A_3},
\]
其中 \(\overline{B}\) 表示“三个人都没有译出密码”。

由于三人独立，故其失败事件也相互独立，于是
\[
P(\overline{B})
=P(\overline{A_1})P(\overline{A_2})P(\overline{A_3})
=\frac45\times\frac23\times\frac34
=\frac25.
\]

所以
\[
P(B)=1-P(\overline{B})=1-\frac25=\frac35.
\]

因此应填 \(\dfrac35\)。
\end{solution}

\begin{note}
互不相容可简化并事件的概率计算，相互独立可简化交事件的概率计算。这里利用对偶律将事件的并转化为事件的交，从而利用独立性进行计算，这一方法经常用到。
\end{note}

\begin{exercise}\label{exer:1-59}
设甲、乙两人轮流射击同一目标，直到射中为止，每次命中率分别为 0.3 和 0.4，甲先射击。\par
(1) 目标是由甲射中的概率 \(p_1\)；\par
(2) 击中目标者射击了 \(k\) 次的概率 \(p_2\)。
\end{exercise}

\begin{solution}
(1) 甲在第一轮命中的概率为 0.3；若第一轮都未命中（概率 \(0.7\times 0.6=0.42\)），则回到初始状态。故
\[
p_1 = 0.3 + 0.42\times p_1 \implies p_1 = \frac{0.3}{1-0.42} = \frac{15}{29} \approx 0.5172.
\]

(2) "击中者射击了 \(k\) 次"意味着该射手在前 \(k-1\) 轮均未命中，第 \(k\) 次命中。若甲击中：前 \(k-1\) 轮甲乙均未命中（概率 \(0.42^{k-1}\)），甲第 \(k\) 次命中（概率 0.3）。若乙击中：前 \(k-1\) 轮均未命中（概率 \(0.42^{k-1}\)），甲第 \(k\) 次未命中（概率 0.7），乙第 \(k\) 次命中（概率 0.4）。故
\[
p_2 = 0.42^{k-1}\times 0.3 + 0.42^{k-1}\times 0.7\times 0.4 = 0.42^{k-1}\times 0.58 \quad (k=1,2,\dots).
\]
\end{solution}

\begin{note}
\textbf{【高频题型模板】全概率公式 + 贝叶斯公式（Ch1 最高频大题）：}

\textbf{识别信号：} 题目分多种"原因"或"来源"，先问总概率，再问"已知结果反推原因"。

\textbf{机械化步骤：}
\begin{enumerate}
    \item \textbf{找完备事件组 $B_1,B_2,\dots,B_n$}（互不相容且并为全集）——通常是"来自哪个工厂/哪种情况"。
    \item \textbf{列出已知：} $P(B_i)$（先验概率）和 $P(A|B_i)$（各原因下事件 $A$ 发生的概率）。
    \item \textbf{全概率公式（正向）：} $P(A)=\sum_{i=1}^n P(B_i)\,P(A|B_i)$。
    \item \textbf{贝叶斯公式（反向）：} $P(B_k|A)=\dfrac{P(B_k)\,P(A|B_k)}{P(A)}$。
\end{enumerate}
\textbf{易错点：} 分母 $P(A)$ 必须用全概率公式展开，不能只写某一项。
\end{note}



\chapter{一维随机变量及其分布}

{\small
\[
\begin{cases}
\text{随机变量分布函数 } F(x) = P(X \le x) \text{ 的性质}
\begin{cases}
(1)\ \text{单调不减} \\
(2)\ 0 \le F(x) \le 1,\ F(-\infty)=0,\ F(+\infty)=1 \\
(3)\ \text{右连续: } F(x+0) = F(x)
\end{cases} \\[1em]
\text{随机变量分类及概率分布}
\begin{cases}
\text{离散型: 分布律} \\
\text{连续型: 概率密度} \\
\text{混合型: 分布函数}
\end{cases}
\ \text{满足概率分布的性质} \\[1em]
\text{常见随机变量及其分布}
\begin{cases}
\text{离散型}
\begin{cases}
(0-1)\text{分布: } P(X=k) = p^k(1-p)^{1-k}\ (k=0,1) \\
\text{二项分布: } P(X=k) = C_n^k p^k(1-p)^{n-k}\ (k=0,1,2,\dots,n) \\
\text{泊松分布: } P(X=k) = \frac{\lambda^k}{k!}e^{-\lambda}\ (k=0,1,2\cdots)
\end{cases} \\[2em]
\text{连续型}
\begin{cases}
\text{均匀分布}
\begin{cases}
X \sim U(a,b) \\
\text{性质: 区间概率等于该区间占整个区间的比例}
\end{cases} \\
\text{指数分布}
\begin{cases}
X \sim E(\lambda) \\
\text{性质: 无记忆性}
\end{cases} \\
\text{正态分布}
\begin{cases}
X \sim N(\mu,\sigma^2) \\
\text{性质}
\begin{cases}
\text{参数含义: } \mu \text{ 为位置参数(注意对称性),} \\
\sigma^2 \text{ 为形状参数} \\
\text{利用标准化求解概率}
\end{cases}
\end{cases}
\end{cases}
\end{cases} \\[1em]
\text{随机变量函数的概率分布}
\begin{cases}
\text{离散型: 关键取值及对应概率} \\
\text{连续型}
\begin{cases}
\text{定义法: 求解函数的分布函数} \\
\text{公式法: 适用单调的随机变量}
\end{cases}
\end{cases}
\end{cases}
\]
}

\begin{introduction}
  \item 随机变量的定义与分类
  \item 分布函数的概念与性质
  \item 离散型随机变量及其分布律
  \item 连续型随机变量及其概率密度
  \item 常见离散分布：(0-1)分布、二项分布、泊松分布
  \item 常见连续分布：均匀分布、指数分布、正态分布
  \item 随机变量函数的分布求法
\end{introduction}

\section{随机变量与分布函数}

\subsection{随机变量}

\begin{definition}[随机变量] \label{def:random-variable}
设随机试验 \(E\) 的样本空间为 \(\Omega\)。如果对每一个 \(\omega \in \Omega\)，都有唯一的实数 \(X(\omega)\) 与之对应，并且对任意实数 \(x\)，集合
\[
\{\omega \mid X(\omega) \le x,\ \omega \in \Omega\}
\]
都是随机事件，则称定义在 \(\Omega\) 上的实值单值函数 \(X(\omega)\) 为\textbf{随机变量}，简记为 \(X\)。
\end{definition}

一般用大写字母 \(X, Y, Z, \dots\) 或希腊字母 \(\xi, \eta, \zeta, \dots\) 来表示随机变量。

\begin{note}
随机变量是"其值随试验结果而定"的变量，其定义有三个要点：
\begin{enumerate}
    \item 它是定义在样本空间 \(\Omega\) 上的实值单值函数——定义域是样本空间，而非实数集；
    \item 对任意实数 \(x\)，\(\{X \le x\}\) 必须是随机事件（即可定义概率）；
    \item 取值在试验之前不能预知，且取各值有一定的概率。
\end{enumerate}
随机事件主要关心“某件事是否发生”，而随机变量进一步把试验结果\textbf{数量化}，使我们不仅能讨论事件发生与否，还能研究数值的分布、均值、方差以及函数变换等整体规律。
\end{note}

\begin{note}
随机变量一定范围内的取值本质上与随机事件等价。例如，"掷骰子出现的点数"是一个随机变量 \(X\)，则 \(X=3\)、\(X \le 4\)、\(X\) 为偶数等都是随机事件。
\end{note}

\subsection{分布函数}

\begin{definition}[分布函数] \label{def:distribution-function}
设 \(X\) 是随机变量，\(x\) 是任意实数，称函数
\[
F(x) = P\{X \le x\} \quad (-\infty < x < +\infty)
\]
为随机变量 \(X\) 的\textbf{分布函数}，也称 \(X\) 服从 \(F(x)\) 分布，记为 \(X \sim F(x)\)。
\end{definition}

\begin{note}
\textbf{自学追问：} 为什么先定义分布函数？因为无论随机变量是离散型还是连续型，\(F(x)=P\{X\le x\}\) 都能统一描述“概率累计到了哪里”。后面的分布律、概率密度、分位数，本质上都是围绕这个累计概率函数展开的。

\textbf{易错点：} 分布函数是“到 \(x\) 为止的累计概率”，不是点 \(x\) 的概率。只有在离散型情形，\(F(x)\) 在某点的跳跃量才对应这个点的概率。
\end{note}

当 \(x\) 从 \(-\infty\) 到 \(+\infty\) 取遍所有实数时，\(\{X \le x\}\) 就相应从不可能事件 \(\varnothing\) 取到必然事件 \(\Omega\)，因此分布函数\textbf{完整地描述了随机变量的概率规律}。


\subsection{分布函数的性质}

\begin{property}\label{prop:distribution-function}
设 \(F(x)\) 为随机变量 \(X\) 的分布函数，则它满足以下三条性质：
\begin{enumerate}
    \item \textbf{单调不减}：对任意实数 \(x_1 < x_2\)，有 \(F(x_1) \le F(x_2)\)；
    \item \textbf{有界性}：\(0 \le F(x) \le 1\)，且
    \[
    F(-\infty) = \lim_{x \to -\infty} F(x) = 0, \quad F(+\infty) = \lim_{x \to +\infty} F(x) = 1;
    \]
    \item \textbf{右连续}：对任意 \(x_0 \in \mathbf{R}\)，有 \(\lim_{x \to x_0^+} F(x) = F(x_0)\)。
\end{enumerate}
\end{property}

\begin{remark}
上述三条性质也是\textbf{充要条件}：若函数 \(F(x)\) 满足性质 (1)--(3)，则它一定是某个随机变量的分布函数。因此，这三条性质常用于判断一个函数能否作为分布函数。
\end{remark}


\subsection{分布函数与事件概率}

分布函数最重要的应用是通过它来计算各种事件的概率。记 \(F(a^-) = \lim_{x \to a^-} F(x)\) 为 \(F\) 在点 \(a\) 处的左极限值，则：

\begin{property}\label{prop:df-probability}
设随机变量 \(X\) 的分布函数为 \(F(x)\)，则：
\begin{enumerate}
    \item \(P\{X \le a\} = F(a)\)；
    \item \(P\{X < a\} = F(a^-)\)；
    \item \(P\{X = a\} = F(a) - F(a^-)\)；
    \item \(P\{a < X < b\} = F(b^-) - F(a)\)；
    \item \(P\{a < X \le b\} = F(b) - F(a)\)；
    \item \(P\{a \le X < b\} = F(b^-) - F(a^-)\)；
    \item \(P\{a \le X \le b\} = F(b) - F(a^-)\)；
    \item \(P\{X > a\} = 1 - F(a)\)。
\end{enumerate}
\end{property}

\begin{note}
记忆技巧：事件中包含等号 \(\le\) 或 \(=\) 的一端取函数值 \(F(\cdot)\)，不包含等号的一端（开区间端）取左极限 \(F(\cdot^-)\)。例如：
\begin{itemize}
    \item \(P\{a < X \le b\}\) 中，\(b\) 端有等号取 \(F(b)\)，\(a\) 端无等号取 \(F(a^-)\)；
    \item \(P\{X = a\} = F(a) - F(a^-)\) 即函数值减去左极限。
\end{itemize}
特别地，当 \(F(x)\) 在点 \(a\) 处连续时，\(F(a^-) = F(a)\)，从而 \(P\{X = a\} = 0\)。

\textbf{易错点：} 在离散型情形，\(F(a)\) 与 \(F(a^-)\) 往往不同；只要端点上的等号判断错了，最后就会差一个跳跃量。
\end{note}

\begin{exercise}\label{exer:2-1}
设随机变量 \(X\) 的分布函数为
\[
F(x) = \begin{cases}
0, & x < 0, \\
\dfrac{1}{2}, & 0 \le x < 1, \\
1 - e^{-x}, & x \ge 1.
\end{cases}
\]
求 \(P\{X = 0\}\)、\(P\{0 < X \le 2\}\) 和 \(P\{X > 1\}\)。
\end{exercise}

\begin{solution}
\textbf{题型识别：} 混合型分布函数（既有跳跃又有连续段），求点概率和区间概率。

\textbf{依据：} 点概率 $P\{X=a\} = F(a) - F(a^-)$（左极限）；区间概率 $P\{a < X \leqslant b\} = F(b) - F(a)$。

\textbf{分步计算：}
\begin{enumerate}
    \item \(P\{X = 0\} = F(0) - F(0^-) = \dfrac{1}{2} - 0 = \dfrac{1}{2}\)。

    $F(0) = \frac{1}{2}$（取 $0 \le x < 1$ 段），$F(0^-)$（从左逼近 $0$）取 $x<0$ 段的值 $0$。

    \item \(P\{0 < X \le 2\} = F(2) - F(0) = (1 - e^{-2}) - \dfrac{1}{2} = \dfrac{1}{2} - e^{-2}\)。

    注意左端是严格不等号"$>$"，所以用 $F(2) - F(0)$；若左端是"$\geqslant$"，则应写 $F(2) - F(0^-)$。

    \item \(P\{X > 1\} = 1 - F(1) = 1 - (1 - e^{-1}) = e^{-1}\)。
\end{enumerate}
\end{solution}

\begin{exercise}\label{exer:2-2}
下列函数中，可以作为某个随机变量的分布函数的是（\quad）。\par
\begin{tabular}{ll}
A. \(F(x) = \dfrac{1}{\pi}\arctan x + \dfrac{1}{2},\ x \in \mathbf{R}\) & B. \(F(x) = \dfrac{1}{1+x^2},\ x \in \mathbf{R}\) \\[6pt]
C. \(F(x) = \begin{cases} 0, & x < 0 \\ 1-x, & 0 \le x < 1 \\ 1, & x \ge 1 \end{cases}\) & D. \(F(x) = \begin{cases} \sin x, & 0 \le x \le \pi \\ 0, & \text{其他} \end{cases}\)
\end{tabular}
\end{exercise}

\begin{solution}
应选 A。

逐条验证分布函数的三条充要条件：
\begin{itemize}
    \item A：\(\arctan x\) 单调递增，故 \(F(x)\) 单调不减；\(F(-\infty)=0\)，\(F(+\infty)=1\)；\(\arctan x\) 连续，故右连续。\textbf{三条均满足}。
    \item B：\(F(x)\) 在 \((-\infty,+\infty)\) 上\textbf{不是}单调不减的（例如 \(F(0)=1\)，\(F(1)=\frac{1}{2}\)），排除。
    \item C：在 \([0,1)\) 上 \(F(x)=1-x\) 单调\textbf{递减}，不满足单调不减，排除。
    \item D：\(\sin x\) 在 \([0,\pi]\) 上不是单调函数，不满足单调不减，排除。
\end{itemize}
\end{solution}

\begin{exercise}\label{exer:2-3}
设随机变量 \(X\) 的分布函数为
\[
F(x) = \begin{cases}
a + b e^{-x}, & x > 0, \\
0, & x \le 0.
\end{cases}
\]
求常数 \(a, b\) 及概率 \(P\{|X| < 2\}\)。
\end{exercise}

\begin{solution}
\(F(x)\) 中含有 2 个未知参数，依据分布函数的性质建立两个独立方程解之。

由分布函数的有界性，有
\[
F(+\infty) = \lim_{x \to +\infty} F(x) = \lim_{x \to +\infty} (a + b e^{-x}) = a = 1,
\]
得 \(a = 1\)。又由分布函数的右连续性，在 \(x=0\) 处有
\[
F(0) = \lim_{x \to 0^+} F(x) = a + b = 0,
\]
得 \(b = -1\)。所以
\[
F(x) = \begin{cases}
1 - e^{-x}, & x > 0, \\
0, & x \le 0.
\end{cases}
\]
从而
\[
P\{|X| < 2\} = P\{-2 < X < 2\} = F(2^-) - F(-2) = 1 - e^{-2} - 0 = 1 - e^{-2}.
\]
\end{solution}

\section{离散型随机变量}

\begin{definition}[离散型随机变量] \label{def:discrete-rv}
如果随机变量 $X$ 只可能取有限个或可列无限个值 $x_1, x_2, \dots$，则称 $X$ 为\textbf{离散型随机变量}。称
\[
P\{X = x_i\} = p_i, \quad i = 1, 2, \dots
\]
为 $X$ 的\textbf{分布律}（或分布列、概率分布），记为 $X \sim p_i$。
\end{definition}

分布律常用表格形式表示：
\[
\begin{array}{c|cccc}
X & x_1 & x_2 & \cdots \\
\hline
P & p_1 & p_2 & \cdots
\end{array}
\]

\begin{property}\label{prop:discrete-rv-condition}
数列 $\{p_i\}\ (i=1,2,\dots)$ 是某离散型随机变量的分布律的\textbf{充要条件}：
\begin{enumerate}
    \item $p_i \ge 0\ (i = 1, 2, \dots)$；
    \item $\displaystyle\sum_{i=1}^{\infty} p_i = 1$。
\end{enumerate}
\end{property}

\begin{property}\label{prop:discrete-rv-df}
设离散型随机变量 $X$ 的分布律为 $P\{X = x_i\} = p_i\ (i=1,2,\dots)$，则：
\begin{enumerate}
    \item $X$ 的分布函数为 $\displaystyle F(x) = \sum_{x_i \le x} p_i$，它是\textbf{阶梯函数}，在 $x = x_i$ 处有跳跃，跳跃高度为 $p_i$；
    \item $P\{X = x_i\} = F(x_i) - F(x_i^-)$；
    \item 对实数轴上的任一集合 $B$，有 $\displaystyle P\{X \in B\} = \sum_{x_i \in B} p_i$。
\end{enumerate}
\end{property}

\begin{note}
离散型随机变量最直观的特征是：\textbf{概率落在若干个可数的点上}。因此它的分布函数是阶梯函数，在每个取值点处都会跳一下，而且跳跃高度恰好等于该点的概率值。

\textbf{自学追问：} 如何快速判断一个随机变量是不是离散型？看它是否只能取有限个或可列个值，例如掷骰子的点数、样本中的废品数、首次成功前的试验次数等，都属于离散型。

\textbf{易错点：} “取值有无穷多个”并不等于“不是离散型”。只要这些值能够按 \(x_1,x_2,\dots\) 列出来，它仍然可能是离散型随机变量。
\end{note}

\section{连续型随机变量}

\begin{definition}[连续型随机变量] \label{def:continuous-rv}
如果随机变量 $X$ 的分布函数 $F(x)$ 可以表示为
\[
F(x) = \int_{-\infty}^{x} f(t)\,\mathrm{d}t \quad (x \in \mathbf{R}),
\]
其中 $f(x)$ 是非负可积函数，则称 $X$ 为\textbf{连续型随机变量}，称 $f(x)$ 为 $X$ 的\textbf{概率密度函数}（简称概率密度），记为 $X \sim f(x)$。
\end{definition}

\begin{note}
\textbf{自学追问：} 连续型随机变量的关键词不是“图像连续”，而是“区间概率可以由密度积分得到”。它真正关心的是某个区间上的概率，而不是落在某一个点上的概率。

\textbf{易错点：} \(f(a)\) 不是 \(P\{X=a\}\)。对连续型随机变量，总有 \(P\{X=a\}=0\)；概率要通过区间上的积分来计算。
\end{note}
\begin{property}\label{prop:pdf-property}
函数 $f(x)$ 为某一随机变量的概率密度的\textbf{充要条件}：$f(x) \ge 0$且 $\displaystyle\int_{-\infty}^{+\infty} f(x)\,\mathrm{d}x = 1$。
\end{property}

\begin{property}\label{prop:continuous-rv-calc}
设连续型随机变量 $X \sim f(x)$，则：
\begin{enumerate}
    \item 对任意实数 $a < b$，$\displaystyle P\{a < X \le b\} = \int_a^b f(x)\,\mathrm{d}x = F(b) - F(a)$；
    \item \textbf{连续型随机变量取一点的概率为零}：$P\{X = c\} = 0$，从而
    \[
    P\{a < X < b\} = P\{a \le X < b\} = P\{a < X \le b\} = P\{a \le X \le b\};
    \]
    \item 若 $f(x)$ 在点 $x$ 处连续，则 $F'(x) = f(x)$。更一般地，
    \[
    f(x) = \begin{cases}
    0, & x \text{ 为 } F(x) \text{ 的不可导点}, \\
    F'(x), & x \text{ 为 } F(x) \text{ 的可导点}.
    \end{cases}
    \]
\end{enumerate}
\end{property}

\begin{remark}
(1) 概率密度没有 $0 \le f(x) \le 1$ 的限制。例如 $X \sim U[0, 0.1]$，则 $f(x) = 10\ (x \in [0, 0.1])$，$f(x)$ 可以大于 1。\newline
(2) 概率密度反映的是概率在点 $x$ 处的"密集程度"而非概率值本身。$P\{x < X \le x+h\} \approx f(x) \cdot h$（当 $h$ 很小时），因此 $f(x)$ 类似于"单位长度上的概率"。从几何上看，$X$ 落入某区间的概率等于该区间上概率密度曲线下的曲边梯形面积。\newline
(3) 连续型随机变量的分布函数一定是\textbf{连续}函数（但反之不成立：$F(x)$ 连续的随机变量不一定是连续型的）。\newline
(4) 同一个分布函数的概率密度函数并不唯一——改变有限个点的值不影响积分结果。
\end{remark}

\begin{note}
除离散型和连续型外，还存在其他类型的随机变量。例如
\[
F(x) = \begin{cases}
0, & x < 0, \\
\dfrac{x^3}{2}, & 0 \le x < 1, \\
1, & x \ge 1,
\end{cases}
\]
该分布函数连续但不绝对连续（因 $\int_0^1 F'(x)\,\mathrm{d}x = \dfrac{1}{2} \neq 1$），故对应的随机变量既非离散型又非连续型。
\end{note}

\begin{exercise}[分布函数性质判断] \label{exer:2-32}
已知 $F(x)$ 是分布函数，下列函数：\\
（1）$aF(x)$ （$a>0, a\neq1$）；~~~~（2）$F(x)+F(-x)$；~~~~（3）$F(x)-F(-x)$；~~~~（4）$F(x)\cdot F(-x)$。\\
其中\textbf{不能作为}分布函数的个数是（\quad）。
\begin{tabular}{llll}
A. 1 & B. 2 & C. 3 & D. 4
\end{tabular}
\end{exercise}
\begin{solution}
应用分布函数的充要条件求解。由于
\begin{itemize}
    \item (1) $aF(+\infty)=a\neq1$；
    \item (2) $F(-\infty)+F(+\infty)=1\neq0$；
    \item (3) $F(-\infty)-F(+\infty)=-1\neq0$；
    \item (4) $F(+\infty)\cdot F(-\infty)=0\neq1$；
\end{itemize}
故(1)、(2)、(3)、(4)都不能作为分布函数，答案选 \textbf{(D)}。
\end{solution}


\begin{exercise}\label{exer:2-4}
设连续型随机变量 $X$ 的分布函数为
\[
F(x) = \begin{cases}
0, & x < 0, \\
Ax^2 + \dfrac{2}{3}x, & 0 \le x < 1, \\
1, & x \ge 1.
\end{cases}
\]
求常数 $A$、$X$ 的概率密度函数 $f(x)$ 以及 $P\{|X| < 0.5\}$。
\end{exercise}

\begin{solution}
由连续型随机变量分布函数的连续性，有
\[
\lim_{x \to 1^-} F(x) = A + \frac{2}{3} = 1 \implies A = \frac{1}{3}.
\]
在 $[0,1)$ 上，$F(x) = \dfrac{x^2}{3} + \dfrac{2x}{3}$，故
\[
f(x) = F'(x) = \frac{2x}{3} + \frac{2}{3}, \quad x \in [0,1),
\]
即
\[
f(x) = \begin{cases}
\dfrac{2x}{3} + \dfrac{2}{3}, & 0 \le x < 1, \\
0, & \text{其他}.
\end{cases}
\]
\[
P\{|X| < 0.5\} = P\{-0.5 < X < 0.5\} = F(0.5) - F(-0.5) = \left(\frac{1}{12} + \frac{1}{3}\right) - 0 = \frac{5}{12}.
\]
\end{solution}

\begin{exercise}\label{exer:2-5}
已知随机变量 $X$ 的概率分布为
\[
\begin{array}{c|ccc}
X & 1 & 2 & 3 \\
\hline
P\{X=k\} & \theta^2 & 2\theta(1-\theta) & (1-\theta)^2
\end{array}
\]
且 $P\{X \ge 2\} = \dfrac{3}{4}$。求未知参数 $\theta$ 及 $X$ 的分布函数 $F(x)$。
\end{exercise}

\begin{solution}
由 $P\{X \ge 2\} = P\{X=2\} + P\{X=3\} = 2\theta(1-\theta) + (1-\theta)^2 = 1 - \theta^2 = \dfrac{3}{4}$，解得 $\theta = \pm\dfrac{1}{2}$。

又 $P\{X=2\} = 2\theta(1-\theta) \ge 0$，当 $\theta = -\dfrac{1}{2}$ 时 $P\{X=2\} = -\dfrac{3}{2} < 0$，故 $\theta = \dfrac{1}{2}$。从而得 $X$ 的概率分布为
\[
\begin{array}{c|ccc}
X & 1 & 2 & 3 \\
\hline
P\{X=k\} & \dfrac{1}{4} & \dfrac{1}{2} & \dfrac{1}{4}
\end{array}
\]
$X$ 的分布函数为
\[
F(x) = P\{X \le x\} = \sum_{x_i \le x} P\{X = x_i\} = \begin{cases}
0, & x < 1, \\
\dfrac{1}{4}, & 1 \le x < 2, \\
\dfrac{3}{4}, & 2 \le x < 3, \\
1, & x \ge 3.
\end{cases}
\]
\end{solution}

\begin{exercise}\label{exer:2-25}
已知离散型随机变量 $X$ 的分布律为 $P\{X=k\} = c(0.75)^k,\ k=1,2,3,\dots$，则常数 $c$ 的值为\underline{\hspace{2cm}}。
\end{exercise}

\begin{solution}
应填 \(\dfrac{1}{3}\)。

离散型随机变量的分布律必须满足“所有概率之和等于 1”，即
\[
\sum_{k=1}^{\infty}P\{X=k\}=1.
\]

把题目给出的分布律代入，得 \(\sum_{k=1}^{\infty}c(0.75)^k=1\)。这里是首项为 \(0.75\)、公比也为 \(0.75\) 的等比级数，所以 \(\sum_{k=1}^{\infty}(0.75)^k=\dfrac{0.75}{1-0.75}=3\)，从而 \(3c=1\)，即 \(c=\dfrac13\)。
\end{solution}

\begin{exercise}\label{exer:2-26}
设随机变量 $X$ 的分布函数为
\[
F(x) = \begin{cases}
0, & x<-1, \\
0.4, & -1 \le x<1, \\
0.8, & 1 \le x<3, \\
1, & x \ge 3.
\end{cases}
\]
求：(1) $X$ 的分布律；(2) $P\{X<2 \mid X \neq 1\}$。
\end{exercise}

\begin{solution}
(1) 根据离散型随机变量的性质：$P\{X=k\} = F(k) - F(k^-)$。
\begin{align*}
P\{X=-1\} &= F(-1) - F(-1^-) = 0.4 - 0 = 0.4, \\
P\{X=1\} &= F(1) - F(1^-) = 0.8 - 0.4 = 0.4, \\
P\{X=3\} &= F(3) - F(3^-) = 1 - 0.8 = 0.2.
\end{align*}
故 $X$ 的分布律为：
\[
\begin{array}{c|ccc}
X & -1 & 1 & 3 \\
\hline
P & 0.4 & 0.4 & 0.2
\end{array}
\]
(2) 由条件概率公式：
\[
P\{X<2 \mid X \neq 1\} = \frac{P\{X<2,\, X \neq 1\}}{P\{X \neq 1\}} = \frac{P\{X=-1\}}{1-P\{X=1\}} = \frac{0.4}{0.6} = \frac{2}{3}.
\]
\end{solution}

\begin{exercise}[连续型分布函数性质] 
设 $F(x)$ 为某连续型随机变量的分布函数，$f(x)$ 为相应的密度函数，则（\quad）。
\begin{enumerate}
    \item[(A)] $F(x)$ 连续但不一定可导；
    \item[(B)] $F(x)$ 可导；
    \item[(C)] $f(x)$ 连续；
    \item[(D)] $0\leqslant f(x)\leqslant1$。
\end{enumerate}
\end{exercise}
\begin{solution}
答案选\textbf{(A)}。分析：
\begin{itemize}
    \item 连续型随机变量的分布函数 $F(x)$ 一定是连续的；
    \item 但 $F(x)$ 仅在密度 $f(x)$ 的连续点处可导；
    \item 密度函数 $f(x)$ 可以不连续（例如分段均匀分布）；
    \item 密度函数 $f(x)$ 可以大于1（例如 $U(0,0.5)$ 的密度为2）。
\end{itemize}
\end{solution}


\begin{exercise}\label{exer:2-6}
已知 $X$ 的分布函数为 $F(x)$，概率密度为 $f(x)$。当 $x \le 0$ 时，$f(x)$ 连续且 $f(x) = F(x)$。若 $F(0) = 1$，求 $f(x)$。
\end{exercise}

\begin{solution}
先利用分布函数的基本性质。因为 \(F(x)\) 单调不减，且题目给出 \(F(0)=1\)，所以对一切 \(x>0\) 都有 \(F(x)=1\)。
这意味着随机变量的全部概率质量都已经集中在 \((-\infty,0]\) 上，因此对 \(x>0\)，密度应为 0。

下面处理 \(x\le 0\) 的部分。连续型随机变量满足 \(F'(x)=f(x)\)，而题目又给出在 \(x\le 0\) 时 \(f(x)=F(x)\)，因此在 \(x\le 0\) 上，\(F(x)\) 满足微分方程
\[
F'(x)=F(x),\qquad F(0)=1.
\]
解这个一阶微分方程得
\[
F(x)=Ce^x,\qquad C=1,\qquad F(x)=e^x\quad(x\le 0).
\]

于是分布函数为
\[
F(x) = \begin{cases} e^x, & x < 0, \\ 1, & x \ge 0, \end{cases} \quad f(x) = \begin{cases} e^x, & x \le 0, \\ 0, & x > 0. \end{cases}
\]

因此所求密度函数是
\[
f(x)=
\begin{cases}
e^x,&x\le 0,\\
0,&x>0.
\end{cases}
\]
\end{solution}

\begin{note}
当分布函数是不连续的分段函数时，为保证右连续性，分段小区间除第一个外都应是左闭右开的（等号跟在"大于等于"端）。
\end{note}

\begin{exercise}\label{exer:2-7}
已知随机变量 $X$ 的概率密度为
\[
f(x) = \begin{cases}
2\left(1 - \dfrac{1}{x^2}\right), & 1 \le x \le 2, \\
0, & \text{其他}.
\end{cases}
\]
求 $X$ 的分布函数 $F(x)$。
\end{exercise}

\begin{solution}
分段积分：
\begin{itemize}
    \item 当 $x < 1$ 时，$F(x) = \displaystyle\int_{-\infty}^x 0\,\mathrm{d}t = 0$；
    \item 当 $1 \le x < 2$ 时，
    \[
    F(x) = \int_1^x 2\left(1 - \frac{1}{t^2}\right)\mathrm{d}t = 2\left[t + \frac{1}{t}\right]_1^x = 2\left(x + \frac{1}{x} - 2\right);
    \]
    \item 当 $x \ge 2$ 时，$F(x) = \displaystyle\int_1^2 2\left(1 - \frac{1}{t^2}\right)\mathrm{d}t = 2\left[\frac{3}{2} - 1\right] = 1$。
\end{itemize}
故
\[
F(x) = \begin{cases}
0, & x < 1, \\
2\left(x + \dfrac{1}{x} - 2\right), & 1 \le x < 2, \\
1, & x \ge 2.
\end{cases}
\]
\end{solution}

\begin{exercise}\label{exer:2-23}
下列函数中，可以作为某个随机变量的概率密度函数的是（\quad）。
\begin{align*}
&\text{A. } f(x) = \begin{cases} 2(1-|x|), & |x| \le 1 \\ 0, & |x|>1 \end{cases}
&&\text{B. } f(x) = \begin{cases} \frac{1}{\sqrt{2\pi}\sigma} e^{-\frac{(x-\mu)^2}{2\sigma^2}}, & x \ge 0 \\ 0, & x<0 \end{cases} \\
&\text{C. } f(x) = \begin{cases} x, & -1<x<0 \\ \frac{3}{4}x, & 0 \le x<2 \\ 0, & \text{其他} \end{cases}
&&\text{D. } f(x) = \begin{cases} \lambda e^{-\lambda x}, & x>0 \\ 0, & x \le 0 \end{cases} \ (\lambda>0)
\end{align*}
\end{exercise}

\begin{solution}
应选 D。

逐条验证概率密度函数的充要条件（非负性和归一性）：
\begin{itemize}
    \item A：$\displaystyle\int_{-\infty}^{\infty} f(x)\,\mathrm{d}x = 2\int_0^1 2(1-x)\,\mathrm{d}x = 4\left[x-\frac{x^2}{2}\right]_0^1 = 2 \neq 1$，错误；
    \item B：正态分布的概率密度在 $x<0$ 时截断为0，但 $\displaystyle\int_0^{+\infty} f(x)\,\mathrm{d}x = \frac{1}{2} \neq 1$，错误；
    \item C：当 $-1<x<0$ 时，$x<0$，故 $f(x)=x<0$，不满足非负性，错误；
    \item D：指数分布的标准形式，满足 $\displaystyle\int_0^{+\infty} \lambda e^{-\lambda x}\,\mathrm{d}x = 1$，正确。
\end{itemize}
\end{solution}

\begin{exercise}\label{exer:2-24}
设随机变量 $X$ 的概率密度为
\[
f(x) = \begin{cases} A\sin x, & 0 \le x \le \pi \\ 0, & \text{其他} \end{cases}
\]
求：(1) 系数 $A$；(2) $X$ 的分布函数 $F(x)$；(3) $P\{0<X<1\}$。
\end{exercise}
\begin{solution}
(1) 由归一性条件：
\[
\int_{0}^{\pi} A\sin x\,\mathrm{d}x = A[-\cos x]_{0}^{\pi} = A(1+1) = 2A = 1,~~A = \dfrac{1}{2}
\]
(2) 当 $x < 0$ 时，$F(x) = 0$；当 $0 \le x < \pi$ 时，
\[
F(x) = \int_{0}^{x} \frac{1}{2}\sin t\,\mathrm{d}t = \frac{1}{2}[-\cos t]_{0}^{x} = \frac{1}{2}(1-\cos x);
\]
当 $x \ge \pi$ 时，$F(x) = 1$。故
\[
F(x) = \begin{cases}
0, & x<0, \\[6pt]
\dfrac{1}{2}(1-\cos x), & 0 \le x < \pi, \\[6pt]
1, & x \ge \pi.
\end{cases}
\]
(3)
\[
P\{0<X<1\} = \int_{0}^{1} \frac{1}{2}\sin t\,\mathrm{d}t = \frac{1}{2}(1-\cos 1).
\]
\end{solution}



\begin{exercise}\label{exer:2-8}
设 $F_1(x), F_2(x)$ 是随机变量的分布函数，$f_1(x), f_2(x)$ 是相应的概率密度，则下列选项中正确的是（\quad）\footnote{验证概率密度和分布函数的定义即可}。\\
A. $F_1(x) + F_2(x)$ 是分布函数\quad
B. $F_1(x) \cdot F_2(x)$ 是分布函数\quad
C. $f_1(x) + f_2(x)$ 是概率密度\quad
D. $f_1(x) \cdot f_2(x)$ 是概率密度
\end{exercise}
\begin{solution}
应选 B。逐条验证：
\begin{itemize}
    \item A：$F_1(+\infty) + F_2(+\infty) = 2 \neq 1$，不满足边界条件，排除；
    \item B：$F_1(x) \cdot F_2(x)$ 单调不减（两个非负单调不减函数之积仍单调不减），右连续，且 $F_1(-\infty)F_2(-\infty) = 0$，$F_1(+\infty)F_2(+\infty) = 1$。\textbf{三条均满足}；
    \item C：$\displaystyle\int_{-\infty}^{+\infty}[f_1(x) + f_2(x)]\,\mathrm{d}x = 2 \neq 1$，排除；
    \item D：$\displaystyle\int_{-\infty}^{+\infty} f_1(x)f_2(x)\,\mathrm{d}x \neq 1$（一般不等于两个积分的乘积），排除。
\end{itemize}
\end{solution}


\section{常见离散型随机变量及分布}

\begin{note}
\textbf{使用提醒：} 识别分布时不要先背公式，先判断题目到底在“数什么”：
\begin{itemize}
    \item 固定 \(n\) 次独立重复试验中的成功总次数，优先考虑二项分布；
    \item 单位时间、单位长度或单位面积中的稀有事件出现次数，优先考虑泊松分布；
    \item 第一次成功出现在第几次，优先考虑几何分布；
    \item 区间内每个位置等可能，优先考虑均匀分布；
    \item 等待时间、寿命、服务时间一类题目，且常伴随无记忆性，优先考虑指数分布；
    \item 大量微小独立因素叠加后的对称钟形变量，优先考虑正态分布。
\end{itemize}

\textbf{易错点：} 同样是“成功次数”，若抽样是不放回进行的，通常就不能直接套二项分布，而应先检查是否属于超几何分布。
\end{note}

\subsection{(0-1)分布} \label{subsec:0-1-dist}
\begin{definition}[(0-1)分布]\label{def:0-1-dist}
随机变量 $X$ 只可能取 0 与 1 两个值，它的分布律为
\[
P\{X = k\} = p^k (1-p)^{1-k},\quad k=0,1\ (0 < p < 1).
\]
则称 $X$ 服从参数为 $p$ 的(0-1)分布，记为 $X \sim B(1,p)$。
\end{definition}

\begin{note}
(0-1)分布是描述伯努利试验结果的计数变量——事件发生计为1，不发生计为0。它是所有离散型分布的基础。
\end{note}


\begin{exercise}[（0-1）分布的分布函数] \label{exer:2-44}
设 $X$ 服从 0-1 分布，其分布律为 $P\{X=k\}=p^k(1-p)^{1-k}, k=0,1, 0<p<1$，求 $X$ 的分布函数。
\end{exercise}
\begin{solution}
0-1分布（伯努利分布）的分布函数为：
\[
F(x)=
\begin{cases}
0, & x<0, \\[6pt]
1-p, & 0\leqslant x<1, \\[6pt]
1, & x\geqslant1.
\end{cases}
\]
解释：
\begin{itemize}
    \item 当 $x<0$ 时，$F(x)=P\{X\leqslant x\}=0$；
    \item 当 $0\leqslant x<1$ 时，$F(x)=P\{X=0\}=1-p$；
    \item 当 $x\geqslant1$ 时，$F(x)=P\{X=0\}+P\{X=1\}=1-p+p=1$。
\end{itemize}
\end{solution}


\subsection{二项分布} \label{subsec:binomial-dist}
\begin{definition}[二项分布]\label{def:binomial-dist}
若随机变量 $X$ 的分布律为
\[
P\{X = k\} = \mathrm{C}_n^k p^k (1-p)^{n-k},\quad k=0,1,\dots,n,
\]
其中 $0 < p < 1$，则称 $X$ 服从参数为 $(n,p)$ 的二项分布，记为 $X \sim B(n,p)$。
\end{definition}

\begin{property}\label{prop:binomial-trans}
若 $X \sim B(n,p)$，则令 $Y = n - X$，有 $Y \sim B(n,1-p)$。
\end{property}

\begin{note}
二项分布的三个识别信号是：\textbf{固定试验次数 \(n\)、各次相互独立、每次只有“成功/失败”两种结果}，并且题目问的是“成功出现了多少次”。当 \(n=1\) 时，二项分布退化为(0-1)分布。

\textbf{易错点：} 若题目问的是“第一次成功出现在第几次”，那对应的是几何分布；若抽样不放回，也不能直接把每次试验当作独立伯努利试验。
\end{note}

\begin{exercise}[二项分布参数求解] \label{exer:2-33}
设 $X\sim B(2,p)$，$Y\sim B(4,p)$，且 $P(X\geqslant1)=\dfrac{5}{9}$，则 $P(Y\geqslant1)=$（\quad）。
\begin{tabular}{llll}
A. $\dfrac{65}{81}$ & B. $\dfrac{16}{81}$ & C. 1 & D. $\dfrac{4}{7}$
\end{tabular}
\end{exercise}
\begin{solution}
应先把“至少 1 次”翻译成补事件。因为
\[
P(X\geqslant1)=1-P(X=0),
\]
所以由题设
\[
\frac{5}{9}=1-P(X=0)=1-\mathrm{C}_2^0p^0(1-p)^2=1-(1-p)^2
\implies (1-p)^2=\frac{4}{9}\implies 1-p=\frac{2}{3},\quad p=\frac{1}{3}.
\]

接着求 \(Y\sim B(4,p)\) 时“至少 1 次成功”的概率，仍然用同一个补事件模板：
\[
\begin{aligned}
P(Y\geqslant1)
&=1-P(Y=0)=1-\mathrm{C}_4^0p^0(1-p)^4\\
&=1-\left(\frac{2}{3}\right)^4\\
&=1-\frac{16}{81}
=\frac{65}{81}.
\end{aligned}
\]

故答案选 \textbf{(A)}。
\end{solution}


\begin{exercise}[概率密度与二项分布结合] \label{exer:2-41}
设随机变量 $X$ 的概率密度为 $f(x)=\begin{cases}2x, & 0<x<1, \\ 0, & \text{其他},\end{cases}$，$Y$ 表示对 $X$ 的 3 次独立重复试验中事件 $\left\{X\leqslant\dfrac{1}{2}\right\}$ 出现的次数，则 $P\{Y=2\}=\underline{\hspace{2cm}}$。
\end{exercise}
\begin{solution}
首先计算单次试验中事件发生的概率：
\[
p = P\left\{X\leqslant\frac{1}{2}\right\} = \int_0^{\frac{1}{2}}2x\,\mathrm{d}x = \left[x^2\right]_0^{\frac{1}{2}} = \frac{1}{4}.
\]
$Y$ 服从二项分布 $B\left(3,\dfrac{1}{4}\right)$，故
\[
P\{Y=2\} = \mathrm{C}_3^2\left(\frac{1}{4}\right)^2\left(\frac{3}{4}\right) = 3\cdot\frac{1}{16}\cdot\frac{3}{4} = \frac{9}{64}.
\]
答案：$\boldsymbol{\dfrac{9}{64}}$。
\end{solution}



\subsection{泊松分布} \label{subsec:poisson-dist}
\begin{definition}[泊松分布] \label{def:poisson-dist}
若随机变量 $X$ 的分布律为
\[
P\{X = k\} = \frac{\lambda^k}{k!} e^{-\lambda},\quad \lambda > 0;\ k=0,1,2,\dots,
\]
则称 $X$ 服从参数为 $\lambda$ 的泊松分布，记为 $X \sim P(\lambda)$。
\end{definition}

\begin{theorem}[泊松定理] \label{thm:poisson-approx}
设 $\lambda > 0$ 是一个常数，$n$ 是任意正整数，设 $\lim_{n \to \infty} np_n = \lambda$，则对于任一固定的非负整数 $k$，有
\[
\lim_{n \to \infty} \mathrm{C}_n^k p_n^k (1-p_n)^{n-k} = \frac{\lambda^k e^{-\lambda}}{k!}.
\]
\end{theorem}
\begin{proof}[证明思路]
将二项分布概率改写为
\[
\mathrm{C}_n^k p_n^k(1-p_n)^{n-k} = \frac{n(n-1)\cdots(n-k+1)}{k!}\,p_n^k\,(1-p_n)^{n-k}.
\]
令 \(np_n \to \lambda\)，则：
\begin{itemize}
    \item \(\dfrac{n(n-1)\cdots(n-k+1)}{n^k} \to 1\)；
    \item \(\left(np_n\right)^k \to \lambda^k\)；
    \item \((1-p_n)^n = \left(1-\dfrac{np_n}{n}\right)^n \to e^{-\lambda}\)（重要极限）。
\end{itemize}
三项相乘，取极限即得 \(\lambda^k e^{-\lambda}/k!\)。
\end{proof}

\begin{note}
泊松定理表明：当$n$很大，$p$很小，且$\lambda = np$适中时，二项分布可用泊松分布近似，即
\[
\mathrm{C}_n^k p^k (1-p)^{n-k} \approx \frac{\lambda^k}{k!} e^{-\lambda}.
\]
一般地，当$n \ge 20,\ p \le 0.05$时，用泊松近似公式逼近二项分布效果比较好，特别当$n \ge 100,\ np \le 10$时，逼近效果更佳。这一结论在计算机出现前极大简化了二项分布的计算。

\textbf{使用提醒：} 泊松分布常用来描述单位时间、单位面积或单位长度中的稀有事件出现次数，参数 \(\lambda\) 同时也等于平均发生次数。

\textbf{易错点：} 泊松分布描述的是“次数”，不是“等待时间”；等待时间问题通常转向指数分布。
\end{note}

\begin{exercise}\label{exer:2-9}
设 $X$ 服从泊松分布，且 $P\{X=1\} = P\{X=2\}$，则 $P\{X=4\} = \underline{\hspace{2cm}}$。
\end{exercise}

\begin{solution}
这类参数题的关键是把已知的两项概率写开，再消去公共因子。

由泊松分布的分布律
\[
P\{X=k\}=\frac{\lambda^k}{k!}e^{-\lambda},
\]
题设条件 \(P\{X=1\}=P\{X=2\}\) 可写成
\[
\frac{\lambda^1}{1!}e^{-\lambda}=\frac{\lambda^2}{2!}e^{-\lambda}.
\]
两边同时除以公共因子 \(e^{-\lambda}\)，再利用 \(\lambda>0\)，得
\[
\lambda=\frac{\lambda^2}{2}\implies 1=\frac{\lambda}{2}\implies \lambda=2.
\]

把 \(\lambda=2\) 代回 \(P\{X=4\}\)：
\[
\begin{aligned}
P\{X=4\}
&=\frac{2^4}{4!}e^{-2}\\
&=\frac{16}{24}e^{-2}
=\frac{2}{3}e^{-2}.
\end{aligned}
\]

因此应填 \(\dfrac{2}{3}e^{-2}\)。
\end{solution}

\begin{exercise}[泊松分布]\label{exer:2-27}
假设某段时间内来百货公司的顾客数服从参数为 $\lambda$ 的泊松分布，而每位顾客购买电视机的概率均为 $p$，且顾客之间是否购买电视机相互独立。求该段时间内百货公司售出 $k$ 台电视机的概率（假设每位顾客至多买一台电视机）。
\end{exercise}

\begin{solution}
设 $A_i$ 表示这段时间内到达百货公司的顾客数为 $i\ (i=0,1,2,\dots)$，则 $P(A_i) = \dfrac{\lambda^i}{i!}e^{-\lambda}$。
当顾客数为 $i$ 时，售出的电视机数 $Y|A_i \sim B(i,p)$，故
\[
P\{Y=k \mid A_i\} = \mathrm{C}_i^k\, p^k(1-p)^{i-k}, \quad k \le i.
\]
由全概率公式：
\begin{align*}
P\{Y=k\} &= \sum_{i=0}^{\infty} P(A_i)\,P\{Y=k \mid A_i\} \\
&= \sum_{i=k}^{\infty} \frac{\lambda^i}{i!}e^{-\lambda} \cdot \frac{i!}{(i-k)!k!}\, p^k(1-p)^{i-k} \\
&= \frac{e^{-\lambda}(\lambda p)^k}{k!} \sum_{i=k}^{\infty} \frac{[\lambda(1-p)]^{i-k}}{(i-k)!}.
\end{align*}
令 $m = i-k$，则
\[
\sum_{i=k}^{\infty} \frac{[\lambda(1-p)]^{i-k}}{(i-k)!} = \sum_{m=0}^{\infty} \frac{[\lambda(1-p)]^m}{m!} = e^{\lambda(1-p)}.
\]
因此
\[
P\{Y=k\} = \frac{e^{-\lambda}(\lambda p)^k}{k!} \cdot e^{\lambda(1-p)} = \frac{(\lambda p)^k}{k!} e^{-\lambda p}, \quad k=0,1,2,\dots.
\]
即 $Y \sim P(\lambda p)$。这表明：若顾客数服从泊松分布，每位顾客独立地以概率 $p$ 购买，则购买总数仍服从泊松分布，参数变为 $\lambda p$。
\end{solution}


\begin{exercise}[超几何分布] \label{exer:2-45}
20 个产品中有 5 个不合格品，若从中随机取出 8 个，求其中不合格品数 $X$ 的概率分布。
\end{exercise}
\begin{solution}
$X$ 服从超几何分布 $H(N=20, M=5, n=8)$，其分布律为
\[
P\{X=k\} = \frac{\mathrm{C}_5^k \mathrm{C}_{15}^{8-k}}{\mathrm{C}_{20}^8}, \quad k=0,1,2,3,4,5.
\]
具体计算：
\[
\begin{array}{c|cccccc}
\hline
X & 0 & 1 & 2 & 3 & 4 & 5 \\
\hline
P & \dfrac{\mathrm{C}_5^0\mathrm{C}_{15}^8}{\mathrm{C}_{20}^8} & \dfrac{\mathrm{C}_5^1\mathrm{C}_{15}^7}{\mathrm{C}_{20}^8} & \dfrac{\mathrm{C}_5^2\mathrm{C}_{15}^6}{\mathrm{C}_{20}^8} & \dfrac{\mathrm{C}_5^3\mathrm{C}_{15}^5}{\mathrm{C}_{20}^8} & \dfrac{\mathrm{C}_5^4\mathrm{C}_{15}^4}{\mathrm{C}_{20}^8} & \dfrac{\mathrm{C}_5^5\mathrm{C}_{15}^3}{\mathrm{C}_{20}^8} \\
\hline
\text{数值} & 0.0511 & 0.2554 & 0.3973 & 0.2384 & 0.0542 & 0.0036 \\
\hline
\end{array}
\]
其中 $\mathrm{C}_{20}^8 = 125970$。答案：$X$ 的超几何分布如上表所示。
\end{solution}


\subsection{几何分布} \label{subsec:geometric-dist}
\begin{definition}[几何分布] \label{def:geometric-dist}
若随机变量 $X$ 的分布律为
\[
P\{X = k\} = (1-p)^{k-1} p,\quad k=1,2,\dots;\ 0 < p < 1,
\]
则称 $X$ 服从参数为 $p$ 的几何分布，记为 $X \sim G(p)$。
\end{definition}

\begin{note}
几何分布描述了在$n$重伯努利试验中，事件$A$首次出现所需的试验次数。它具有与指数分布类似的无记忆性。
\end{note}

\begin{exercise}[几何分布的无记忆性] \label{exer:2-58}
设 $X$ 为取正整数值的离散型随机变量，证明：$X$ 是服从参数为 $p$ 的几何分布的充要条件是 $X$ 具有无记忆性，即 $P(X=k+n\mid X>k)$ 与 $k$ 无关（$k,n$ 为任意正整数）。
\end{exercise}

\begin{solution}
\textbf{必要性证明}：若 $X$ 服从几何分布，则 $P(X=k)=p(1-p)^{k-1},\ k=1,2,\dots$。对任意正整数 $k,n$：
\[
\begin{aligned}
P(X=k+n\mid X>k) &= \frac{P(X=k+n, X>k)}{P(X>k)} = \frac{P(X=k+n)}{P(X>k)} \\
&= \frac{p(1-p)^{k+n-1}}{\sum_{j=k+1}^{\infty}p(1-p)^{j-1}} = \frac{p(1-p)^{k+n-1}}{(1-p)^k} = p(1-p)^{n-1}.
\end{aligned}
\]
结果与 $k$ 无关，故几何分布具有无记忆性。\\
\textbf{充分性证明}：设 $P(X=k+1\mid X>k)=p$ 与 $k$ 无关。记 $q=1-p$，则
\[
P(X>k) = q\cdot P(X>k-1) = q^k\cdot P(X>0).
\]
由 $P(X>0)=1$ 可得
\[
P(X=k) = P(X>k-1) - P(X>k) = q^{k-1} - q^k = p(1-p)^{k-1}.
\]
这正是参数为 $p$ 的几何分布。
\end{solution}



\subsection{超几何分布} \label{subsec:hypergeo-dist}
\begin{definition}[超几何分布] \label{def:hypergeo-dist}
若随机变量 $X$ 的分布律为
\[
P\{X=k\} = \frac{\mathrm{C}_M^k \mathrm{C}_{N-M}^{n-k}}{\mathrm{C}_N^n},\quad \max\{0,n-N+M\} \le k \le \min\{M,n\},
\]
其中 $M,N,n$ 为正整数且 $M \le N,\ n \le N$，则称 $X$ 服从参数为 $(n,N,M)$ 的超几何分布，记为 $X \sim H(n,N,M)$。
\end{definition}

\begin{note}
超几何分布描述的是不放回抽样场景：设总体中有$N$个产品，其中$M$个不合格品。若从中随机不放回地抽取$n$个，则其中不合格品的个数$X$服从超几何分布。当$n \ll N$时（抽取个数远小于总体），每次抽取后总体不合格品率$p=\frac{M}{N}$改变甚微，此时不放回抽样可近似看作放回抽样，超几何分布可用二项分布$H(n,N,M) \approx B(n,\frac{M}{N})$近似。
\end{note}

\subsection{极值分布} \label{subsec:extreme-dist}
\begin{property}\label{prop:extreme-value}
对于 $n$ 个相互独立的随机变量 $\xi_1,\xi_2,\dots,\xi_n$，其分布函数分别为 $F_1(x),F_2(x),\dots,F_n(x)$，则：
\begin{itemize}
    \item 最大值 $\eta_1 = \max\{\xi_1,\xi_2,\dots,\xi_n\}$ 的分布函数为
    \[
    F_{\eta_1}(x) = P\{\eta_1 \le x\} = P\left(\bigcap_{k=1}^n \xi_k \le x\right) = \prod_{k=1}^n F_k(x);
    \]
    \item 最小值 $\eta_2 = \min\{\xi_1,\xi_2,\dots,\xi_n\}$ 的分布函数为
    \[
    F_{\eta_2}(x) = P\{\eta_2 \le x\} = 1 - P\{\eta_2 > x\} = 1 - P\left(\bigcap_{k=1}^n \xi_k > x\right) = 1 - \prod_{k=1}^n (1-F_k(x)).
    \]
\end{itemize}
\end{property}


\begin{exercise}[最小值分布（系统寿命）] \label{exer:2-61}
一辆新车装有 4 个轮胎并配有 1 个备胎，当使用的轮胎损坏时就用备胎更换。设所有轮胎的寿命均服从参数为 $0.2$ 的指数分布且相互独立，记 $T$ 为开始使用备胎的时间，求 $T$ 的分布函数。
\end{exercise}

\begin{solution}
设四个在用轮胎的寿命分别为 $X_1,X_2,X_3,X_4$，则 $X_i\stackrel{\text{i.i.d.}}{\sim}E(0.2)$。

开始使用备胎的时间为
\[
T = \min\{X_1,X_2,X_3,X_4\}.
\]

单个轮胎的分布函数：
\[
F_X(x) = 1-e^{-0.2x},\quad x\geqslant0.
\]

$T$ 的分布函数：
\[
\begin{aligned}
F_T(t) &= P\{T\leqslant t\} = 1-P\{T>t\} = 1-P\{X_1>t, X_2>t, X_3>t, X_4>t\} \\
&= 1-\left[P\{X_1>t\}\right]^4 = 1-\left[1-F_X(t)\right]^4 = 1-\left[e^{-0.2t}\right]^4 = 1-e^{-0.8t},\quad t\geqslant0.
\end{aligned}
\]
当 $t<0$ 时，$F_T(t)=0$。故
\[
F_T(t)=
\begin{cases}
0, & t<0, \\[4pt]
1-e^{-0.8t}, & t\geqslant0.
\end{cases}
\]
即 $T\sim E(0.8)$。
\end{solution}



\section{常见连续型随机变量及分布}


\begin{exercise}[偶函数概率密度的性质证明] \label{exer:2-48}
设连续型随机变量 $X$ 的概率密度 $f(x)$ 是偶函数，证明下列结论：
\begin{enumerate}
    \item $F(-a)=1-F(a)$；
    \item $P(|X|<a)=2F(a)-1$；
    \item $P(|X|>a)=2[1-F(a)]$。
\end{enumerate}
\end{exercise}
\begin{solution}
思路是先把“$f(x)$ 是偶函数”翻译成“分布关于原点对称”，再把后两问都化成分布函数的差。

(1) 由分布函数定义，
\[
F(-a)=\int_{-\infty}^{-a}f(x)\,\mathrm{d}x.
\]
令 $t=-x$，则
\[
\begin{aligned}
F(-a)
&=\int_{+\infty}^{a}f(-t)(-\,\mathrm{d}t)
=\int_{a}^{+\infty}f(-t)\,\mathrm{d}t \\
&=\int_{a}^{+\infty}f(t)\,\mathrm{d}t
=1-\int_{-\infty}^{a}f(t)\,\mathrm{d}t
=1-F(a),
\end{aligned}
\]
其中第三步用到了偶函数条件 $f(-t)=f(t)$。

(2) 由题意
\[
P(|X|<a)=P(-a<X<a)=F(a)-F(-a).
\]
再代入第 (1) 问结果，得到
\[
P(|X|<a)=F(a)-[1-F(a)]=2F(a)-1.
\]

(3) 同理，
\[
P(|X|>a)=1-P(|X|\leqslant a)=1-[F(a)-F(-a)].
\]
再由 $F(-a)=1-F(a)$ 可得
\[
P(|X|>a)=1-\bigl[F(a)-(1-F(a))\bigr]=2[1-F(a)].
\]

因此三条结论都得证。这个题的核心是：\textbf{偶密度 $\Rightarrow$ 分布关于原点对称}。
\end{solution}

\subsection{均匀分布} \label{subsec:uniform-dist}
\begin{definition}[均匀分布]\label{def:uniform-dist}
若随机变量 $X$ 的概率密度为
\[
f(x) = \begin{cases}
\frac{1}{b-a}, & a < x < b, \\
0, & \text{其他},
\end{cases}
\]
则称 $X$ 在区间 $(a,b)$ 内服从均匀分布，记为 $X \sim U(a,b)$。其分布函数为
\[
F(x) = \begin{cases}
0, & x < a, \\
\frac{x-a}{b-a}, & a \le x < b, \\
1, & x \ge b.
\end{cases}
\]
\end{definition}

\begin{note}
均匀分布描述了随机变量$X$在区间$(a,b)$内任意子区间取值的概率只与子区间长度成正比，与其位置无关。几何概型是均匀分布的实际背景。

\textbf{使用提醒：} 做均匀分布题时，先找清楚支持区间，再把所求概率翻译成“长度比”。

\textbf{易错点：} 真正有概率的是区间长度，不是某个孤立点；因此端点单点概率仍为 0。
\end{note}

\begin{exercise}\label{exer:2-13}
设随机变量 $X$ 满足 $1 \le X \le 4$，且 $P\{X=1\} = \frac{1}{4}$，$P\{X=4\} = \frac{1}{4}$，当 $X$ 在区间 $(1,4)$ 内取值时，$X$ 服从均匀分布，求 $X$ 的分布函数 $F(x)$。
\end{exercise}
\begin{solution}
当 $x < 1$ 时，$F(x) = 0$；
当 $1 \le x < 4$ 时，
\[
\begin{aligned}
F(x) = P\{X \le 1\} + P\{1 < X \le x\} = \frac{1}{4} + \frac{x-1}{4-1} \cdot \frac{1}{2} = \frac{5 + 2x}{12};
\end{aligned}
\]
当 $x \ge 4$ 时，$F(x) = 1$。
故
\[
F(x) = \begin{cases}
0, & x < 1, \\
\frac{5 + 2x}{12}, & 1 \le x < 4, \\
1, & x \ge 4.
\end{cases}
\]
\end{solution}

\begin{exercise}\label{exer:2-16}
设 $X \sim U(a,b)$（$a>0$），且 $P\{X>4\} = \frac{1}{2}$，$P\{0<X<3\} = \frac{1}{4}$，求 $a,b$ 及 $P\{1<X<5\}$。
\end{exercise}
\begin{solution}
均匀分布题的核心是把概率翻译成区间长度比。

由于 \(X\sim U(a,b)\)，所以
\[
P\{X>4\}=\frac{b-4}{b-a}=\frac12.
\]
化简得
\[
2(b-4)=b-a \quad \Longrightarrow \quad a+b=8.
\]

又因为 \(a>0\)，所以事件 \(\{0<X<3\}\) 在支持区间上等价于 \(\{a<X<3\}\)。于是
\[
P\{0<X<3\}=P\{a<X<3\}=\frac{3-a}{b-a}=\frac14.
\]
化简得
\[
4(3-a)=b-a \quad \Longrightarrow \quad b=12-3a.
\]

联立
\[
\begin{cases}
a+b=8,\\
b=12-3a,
\end{cases}
\]
可得 \(a=2,\ b=6\)。因为此时 \(a=2>1\)，所以事件 \(\{1<X<5\}\) 等价于 \(\{2<X<5\}\)。因此
\[
P\{1<X<5\}=P\{2<X<5\}=\frac{5-2}{6-2}=\frac34.
\]

所以 \(a=2,\ b=6,\ P\{1<X<5\}=\frac34\)。
\end{solution}


\begin{exercise}\label{exer:2-29}
设随机变量 $X$ 在区间 $(0,8)$ 上服从均匀分布。
\begin{enumerate}
    \item[(1)] 求方程 $y^2 + 2y + X = 0$ 有实根的概率；
    \item[(2)] 若对随机变量 $X$ 进行 4 次独立观察，记 $Y$ 为上述方程有实根的次数，求 $Y$ 的数学期望和方差。
\end{enumerate}
\end{exercise}
\begin{solution}
(1) 先把“方程有实根”翻译成关于随机变量 \(X\) 的事件。二次方程 \(y^2+2y+X=0\) 有实根，当且仅当判别式 \(\Delta=2^2-4X\ge 0\)，即 \(X\le 1\)。又因为 \(X\sim U(0,8)\)，对均匀分布来说，区间概率等于“长度比”，所以
\[
P\{\text{方程有实根}\}=P(X\le 1)=\frac{1-0}{8}=\frac{1}{8}.
\]

(2) 现在对 \(X\) 作 4 次独立观察。每观察一次，就相当于做一次“方程是否有实根”的判断：
\begin{itemize}
    \item 成功：方程有实根；
    \item 成功概率：\(p=P(X\le 1)=\dfrac18\)。
\end{itemize}
由于 4 次观察相互独立，所以“4 次中成功的次数” \(Y\) 服从二项分布
\[
Y\sim B\!\left(4,\frac18\right).
\]
于是直接用二项分布的数字特征公式
\[
E(Y) = np = 4 \times \frac{1}{8} = \frac{1}{2}, \quad D(Y) = np(1-p) = 4 \times \frac{1}{8} \times \frac{7}{8} = \frac{7}{16}.
\]

所以所求数学期望为 \(\dfrac12\)，方差为 \(\dfrac{7}{16}\)。
\end{solution}


\begin{exercise}[均匀分布与二项分布结合] \label{exer:2-40}
设随机变量 $X$ 服从区间 $(2,5)$ 上的均匀分布，则对 $X$ 进行 3 次独立观测中，至少有 2 次的观测值大于 3 的概率为$\underline{\hspace{2cm}}$。
\end{exercise}
\begin{solution}
首先计算单次观测值大于3的概率：
\[
P\{X>3\} = \frac{5-3}{5-2} = \frac{2}{3}.
\]
设 $Y$ 表示3次独立观测中观测值大于3的次数，则 $Y\sim B\left(3,\dfrac{2}{3}\right)$。所求概率：
\[
P\{Y\geqslant2\} = P\{Y=2\} + P\{Y=3\} = \mathrm{C}_3^2\left(\frac{2}{3}\right)^2\left(\frac{1}{3}\right) + \mathrm{C}_3^3\left(\frac{2}{3}\right)^3 = 3\cdot\frac{4}{9}\cdot\frac{1}{3} + \frac{8}{27} = \frac{12}{27} + \frac{8}{27} = \frac{20}{27}.
\]
答案：$\boldsymbol{\dfrac{20}{27}}$。
\end{solution}

\subsection{指数分布} \label{subsec:exp-dist}
\begin{definition}[指数分布] \label{def:exp-dist}
若随机变量 $X$ 的概率密度为
\[
f(x) = \begin{cases}
\lambda e^{-\lambda x}, & x > 0, \\
0, & x \le 0,
\end{cases}
\]
其中 $\lambda > 0$，则称 $X$ 服从参数为 $\lambda$ 的指数分布，记为 $X \sim E(\lambda)$。其分布函数为
\[
F(x) = \begin{cases}
1 - e^{-\lambda x}, & x \ge 0, \\
0, & x < 0.
\end{cases}
\]
\end{definition}

\begin{property}[无记忆性] \label{prop:memoryless}
指数分布具有无记忆性：对于任意 $s,t > 0$，有
\[
P\{X > s + t \mid X > s\} = P\{X > t\}.
\]
这意味着，对于一个已经工作了$s$小时的元件，其再工作至少$t$小时的概率，与从最开始起就至少工作$t$小时的概率相同。
\end{property}
\begin{proof}
设 \(X \sim E(\lambda)\)，则 \(P\{X > x\} = e^{-\lambda x}\ (x \geq 0)\)。由条件概率的定义：
\begin{align*}
P\{X > s + t \mid X > s\} &= \frac{P\{X > s + t\}}{P\{X > s\}} \\
&= \frac{e^{-\lambda(s+t)}}{e^{-\lambda s}} = e^{-\lambda t} = P\{X > t\}.
\end{align*}
\end{proof}

\begin{note}
\textbf{使用提醒：} 指数分布描述的是“等待时间、寿命、服务时间”一类连续变量，而不是“发生了多少次”。若题目给出的是平均等待时间，应先换成参数关系 \(EX=\dfrac{1}{\lambda}\)。

\textbf{易错点：} \(\lambda\) 是速率，不是平均寿命；平均寿命是 \(\dfrac{1}{\lambda}\)。另外，无记忆性是指数分布的特殊性质，不要轻易对别的分布套用。
\end{note}


\begin{exercise}\label{exer:2-28}
设某医院门诊的内科看病窗口有 5 个，患者平均就诊时间为 10 分钟，假定每个患者的就诊时间服从指数分布。求：
\begin{enumerate}
    \item[(1)] 患者就诊时间超过 10 分钟的概率；
    \item[(2)] 若某时刻 5 个窗口都有病人就诊，则恰有 2 人的看病时间超过 10 分钟的概率。
\end{enumerate}
\end{exercise}

\begin{solution}
由题意，平均就诊时间为 10 分钟，故 $\dfrac{1}{\lambda}=10$，即 $\lambda=0.1$。记患者就诊时间为 $X$，则 $X \sim E(0.1)$。\\
(1) 
\[
P\{X \ge 10\} = \int_{10}^{\infty} 0.1\, e^{-0.1x}\,\mathrm{d}x = e^{-1}.
\]
(2) 设 5 个窗口都有病人就诊时，恰有 2 人看病时间超过 10 分钟为事件 $B$。每位患者就诊时间超过 10 分钟的概率为 $e^{-1}$，则
\[
P(B) = \mathrm{C}_5^2\, (e^{-1})^2\, (1-e^{-1})^3 = 10\, e^{-2}\,(1-e^{-1})^3.
\]
\end{solution}

\begin{exercise}\label{exer:2-10}
设随机变量 $X$ 服从参数为 $\lambda(\lambda>0)$ 的指数分布，且 $P\{X \le 1\} = \frac{1}{2}$。
\begin{enumerate}
    \item[(1)] 求参数 $\lambda$；
    \item[(2)] 求 $P\{X > 2 \mid X > 1\}$。
\end{enumerate}
\end{exercise}
\begin{solution}
(1) 根据指数分布的分布函数，$P\{X \le 1\} = 1 - e^{-\lambda} = \frac{1}{2}$，解得 $\lambda = \ln 2$。

(2) 由无记忆性，
\[
P\{X > 2 \mid X > 1\} = P\{X > 1\} = 1 - P\{X \le 1\} = \frac{1}{2}.
\]
或者直接计算：
\[
P\{X > 2 \mid X > 1\} = \frac{P\{X > 2\}}{P\{X > 1\}} = \frac{e^{-2\ln 2}}{e^{-\ln 2}} = e^{-\ln 2} = \frac{1}{2}.
\]
\end{solution}

\begin{exercise}[（指数分布无记忆性应用）]
设某电子元件的寿命服从参数 $\lambda=\dfrac{1}{4}$ 的指数分布（单位：万小时），某系统并联了两个这种电子元件，计算：
\begin{enumerate}
    \item[(1)] 已知一个元件已经工作了三万小时，求再正常工作四万小时的概率；
    \item[(2)] 求系统寿命大于五万小时的概率。
\end{enumerate}
\end{exercise}
\begin{solution}
设元件寿命为 $X$，则 $X\sim E(\lambda)$，其中 $\lambda=\dfrac{1}{4}$。\\
(1) 由指数分布的无记忆性：
\[
P\{X>3+4 \mid X>3\} = P\{X>4\} = e^{-\lambda\cdot4} = e^{-1}.
\]
(2) 并联系统寿命大于5万小时，即至少一个元件正常工作。设两个元件寿命为 $X_1, X_2$，则系统寿命为 $\max\{X_1,X_2\}$：
\[
\begin{aligned}
P\{\max\{X_1,X_2\}>5\} &= 1-P\{X_1\leqslant5, X_2\leqslant5\} \\
&= 1-\left(1-e^{-\lambda\cdot5}\right)^2 \\
&= 1-\left(1-e^{-\frac{5}{4}}\right)^2.
\end{aligned}
\]
\end{solution}


\begin{exercise}[指数分布与系统更换]
设一批元件的寿命独立同分布，其概率密度函数为
\[
f(x)=
\begin{cases}
e^{-x}, & x\geqslant0, \\[4pt]
0, & x<0,
\end{cases}
\]
系统开始工作时只有一个元件工作，在它损坏后立即更换一个新元件接替工作，求系统从开始工作到时刻 $T(T>0)$ 为止，恰更换了一个元件的概率。
\end{exercise}
\begin{solution}
设两个元件的寿命分别为 $X_1$ 和 $X_2$，则 $X_1, X_2 \stackrel{\text{i.i.d.}}{\sim} E(1)$，"恰更换了一个元件"等价于事件：
\[
\{X_1 < T,\ X_1+X_2 > T\}.
\]
所求概率：
\[
\begin{aligned}
P\{X_1 < T,\ X_1+X_2 > T\} &= \int_0^T f_{X_1}(x)\cdot P\{X_2 > T-x\}\,dx = \int_0^T e^{-x}\cdot e^{-(T-x)}\,dx \\
&= \int_0^T e^{-T}\,dx = Te^{-T}.
\end{aligned}
\]
\end{solution}



\begin{exercise}\label{exer:2-15}
某元件的工作寿命 $X$（小时）服从参数为 $\lambda(\lambda>0)$ 的指数分布。
\begin{enumerate}
    \item[(1)] 求该元件正常工作 $t$ 小时的概率；
    \item[(2)] 已知该元件已正常工作 10 小时，求在此基础上再工作 10 小时的概率（$\lambda=0.01$）。
\end{enumerate}
\end{exercise}
\begin{solution}
由题意 $X \sim F(x) = \begin{cases} 1 - e^{-\lambda x}, & x \ge 0, \\ 0, & \text{其他}, \end{cases}$

(1) $P\{X>t\} = 1 - F(t) = e^{-\lambda t}$。

(2) 由无记忆性，
\[
P\{X>20 \mid X>10\} = P\{X>10\} = e^{-0.01 \times 10} = e^{-0.1}.
\]
或者直接计算：
\[
P\{X>20 \mid X>10\} = \frac{P\{X>20\}}{P\{X>10\}} = \frac{e^{-0.2}}{e^{-0.1}} = e^{-0.1}.
\]
\end{solution}

\subsection{正态分布} \label{subsec:normal-dist}
\begin{definition}[正态分布]\label{def:normal-dist}
若随机变量 $X$ 的概率密度为
\[
f(x) = \frac{1}{\sqrt{2\pi}\sigma} \exp\left\{-\frac{(x-\mu)^2}{2\sigma^2}\right\},\quad -\infty < x < +\infty,
\]
其中 $\mu$ 为常数，$\sigma > 0$ 为常数，则称 $X$ 服从参数为 $(\mu,\sigma^2)$ 的正态分布，记为 $X \sim N(\mu,\sigma^2)$。
\end{definition}

\begin{definition}[标准正态分布] \label{def:std-normal}
当 $\mu = 0,\ \sigma = 1$ 时，称 $X$ 服从标准正态分布，记为 $X \sim N(0,1)$。其概率密度和分布函数分别记为
\[
\varphi(x) = \frac{1}{\sqrt{2\pi}} e^{-\frac{x^2}{2}},\quad \Phi(x) = \int_{-\infty}^x \varphi(t) dt.
\]
\end{definition}

\begin{property}\label{prop:normal-properties}
标准正态分布具有以下性质：
\begin{enumerate}
    \item $\Phi(0) = \frac{1}{2},\ \Phi(-a) = 1 - \Phi(a)$；
    \item 若 $X \sim N(\mu,\sigma^2)$，则 $F(x) = \Phi\left(\frac{x-\mu}{\sigma}\right)$； 若 $X \sim N(\mu,\sigma^2)$，则 $\frac{X-\mu}{\sigma} \sim N(0,1)$；
    \item 若 $X \sim N(\mu,\sigma^2)$，则 $P\{X \le \mu\} = P\{X > \mu\} = \frac{1}{2}$；
    \item $P\{a < X < b\} = \Phi\left(\frac{b-\mu}{\sigma}\right) - \Phi\left(\frac{a-\mu}{\sigma}\right)$。
\end{enumerate}
\end{property}

\begin{note}
\textbf{自学追问：} 正态分布为什么总要先标准化？因为概率表通常只给标准正态分布函数 \(\Phi\)，把 \(X\) 化成 \(\dfrac{X-\mu}{\sigma}\) 以后，所有正态题都能统一到同一张表上处理。

\textbf{使用提醒：} \(\mu\) 决定分布中心，\(\sigma\) 决定离散程度。题目一旦出现“距均值多远”“关于均值对称”“比较波动大小”等信息，就优先考虑标准化。

\textbf{易错点：} 方差是 \(\sigma^2\)，标准差才是 \(\sigma\)。标准化分母必须写 \(\sigma\)，不能写成 \(\sigma^2\)。
\end{note}

\begin{exercise}[正态分布对称点] 
设随机变量 $X\sim N(3,4)$，且 $P(X<C)=P(X>C)$，则常数 $C=\underline{\hspace{2cm}}$。
\end{exercise}
\begin{solution}
对于正态分布 $N(\mu,\sigma^2)$，其概率密度函数关于均值 $\mu$ 对称，因此
\[
P(X<\mu)=P(X>\mu)=\frac{1}{2}.
\]
由题设条件 $P(X<C)=P(X>C)$ 可知，$C$ 必为分布的对称中心，即
\[
C=\mu=3.
\]
\end{solution}

\begin{exercise}
设随机变量 $X \sim N(10,\sigma^2)$，且 $P\{10 < X < 20\} = 0.35$，则 $P\{X < 0\} = \underline{\hspace{2cm}}$。
\end{exercise}
\begin{solution}
由正态分布的对称性，
\[
P\{10 < X < 20\} = \Phi\left(\frac{20-10}{\sigma}\right) - \Phi(0) = \Phi\left(\frac{10}{\sigma}\right) - 0.5 = 0.35,
\]
得 $\Phi\left(\frac{10}{\sigma}\right) = 0.85$。因此
\[
P\{X < 0\} = \Phi\left(\frac{0-10}{\sigma}\right) = 1 - \Phi\left(\frac{10}{\sigma}\right) = 0.15.
\]
\end{solution}

\begin{exercise}\label{exer:2-17}
设随机变量 $X$ 服从正态分布 $N(\mu,\sigma^2)$，则随着 $\sigma$ 的增大，$P\{|X-\mu|<1\}$（\quad）。\\
(A) 单调增加 \quad (B) 单调减少 \quad (C) 保持不变 \quad (D) 先增后减
\end{exercise}
\begin{solution}
应选 (B)。

这类题不要直接盯着 \(\sigma\) 比大小，而要先把概率标准化。由
\[
Z=\frac{X-\mu}{\sigma}\sim N(0,1),
\]
可得
\[
\begin{aligned}
P\{|X-\mu|<1\}
&=P\left\{-1<X-\mu<1\right\}\\
&=P\left\{-\frac{1}{\sigma}<\frac{X-\mu}{\sigma}<\frac{1}{\sigma}\right\}\\
&=P\left\{-\frac{1}{\sigma}<Z<\frac{1}{\sigma}\right\}\\
&=2\Phi\left(\frac{1}{\sigma}\right)-1.
\end{aligned}
\]

现在只需考察 \(\dfrac{1}{\sigma}\) 的变化。随着 \(\sigma\) 增大，
\[
\frac{1}{\sigma}
\]
单调减小；而 \(\Phi(u)\) 是单调增加函数，所以
\[
\Phi\left(\frac{1}{\sigma}\right)
\]
随 \(\sigma\) 的增大而单调减小，从而
\[
2\Phi\left(\frac{1}{\sigma}\right)-1
\]
也单调减小。

因此 \(P\{|X-\mu|<1\}\) 随 \(\sigma\) 的增大而\textbf{单调减少}，应选 \textbf{(B)}。
\end{solution}

\begin{exercise}\label{exer:2-12}
设随机变量 $X$ 服从正态分布 $N(0,\sigma^2)$，对给定的 $\alpha(0<\alpha<1)$，数 $u_\alpha$ 满足 $P\{X>u_\alpha\} = \alpha$，若 $P\{|X|<x\} = \alpha$，则 $x$ 等于（\quad）。
(A) $u_{\frac{\alpha}{2}}$ \quad (B) $u_{\frac{1-\alpha}{2}}$ \quad (C) $u_{1-\frac{\alpha}{2}}$ \quad (D) $u_{1-\alpha}$
\end{exercise}

\begin{solution}
应选 (B)。

题目给的是中间概率
\[
P\{|X|<x\}=P\{-x<X<x\}=\alpha.
\]
由于正态分布关于 0 对称，所以区间外剩余的总概率为 \(1-\alpha\)，左右两边各占 \(\dfrac{1-\alpha}{2}\)。

也可以直接由分布函数写成
\[
\Phi\left(\frac{x}{\sigma}\right)-\Phi\left(-\frac{x}{\sigma}\right)=\alpha,
\]
再利用对称性 \(\Phi(-t)=1-\Phi(t)\)，得到
\[
2\Phi\left(\frac{x}{\sigma}\right)-1=\alpha,
\]
即
\[
\Phi\left(\frac{x}{\sigma}\right)=\frac{1+\alpha}{2}.
\]
于是
\[
P\{X>x\}=1-\Phi\left(\frac{x}{\sigma}\right)=\frac{1-\alpha}{2}.
\]

而 \(u_\beta\) 的定义是满足 \(P\{X>u_\beta\}=\beta\) 的点，与上式对照可知 \(x=u_{\frac{1-\alpha}{2}}\)。

故答案选 \textbf{(B)}。
\end{solution}


\begin{exercise}[正态分布标准差比较]
设 $X\sim N(\mu_1,\sigma_1^2), Y\sim N(\mu_2,\sigma_2^2)$，且 $P(|X-\mu_1|<1)>P(|Y-\mu_2|<1)$，则（\quad）。
(A)$\sigma_1<\sigma_2$\quad (B)$\sigma_1>\sigma_2$\quad (C) $\mu_1<\mu_2$\quad (D) $\mu_1>\mu_2$。
\end{exercise}
\begin{solution}
答案选\textbf{(A)}。标准化得：
\[
P\left\{\frac{|X-\mu_1|}{\sigma_1}<\frac{1}{\sigma_1}\right\}>P\left\{\frac{|Y-\mu_2|}{\sigma_2}<\frac{1}{\sigma_2}\right\}
\Rightarrow 2\Phi\left(\frac{1}{\sigma_1}\right)-1 > 2\Phi\left(\frac{1}{\sigma_2}\right)-1
\]
\[
\Rightarrow \Phi\left(\frac{1}{\sigma_1}\right)>\Phi\left(\frac{1}{\sigma_2}\right)\implies
\frac{1}{\sigma_1}>\frac{1}{\sigma_2} \Rightarrow \sigma_1<\sigma_2.
\]
\end{solution}

\begin{exercise}[正态分布标准差比较] \label{exer:2-34}
设随机变量 $X$ 服从正态分布 $N(\mu_1,\sigma_1^2)$，随机变量 $Y$ 服从正态分布 $N(\mu_2,\sigma_2^2)$，$\sigma_1,\sigma_2>0$，且 $P\{|X-\mu_1|<1\}>P\{|Y-\mu_2|<1\}$，则必有（\quad）。
\begin{tabular}{ll}
A. $\sigma_1<\sigma_2$ & B. $\sigma_1>\sigma_2$ \\
C. $\mu_1<\mu_2$ & D. $\mu_1>\mu_2$
\end{tabular}
\end{exercise}
\begin{solution}
正态分布标准化：
\[
P\left\{\frac{|X-\mu_1|}{\sigma_1}<\frac{1}{\sigma_1}\right\}>P\left\{\frac{|Y-\mu_2|}{\sigma_2}<\frac{1}{\sigma_2}\right\},
\]
\[
2\Phi\left(\frac{1}{\sigma_1}\right)-1 > 2\Phi\left(\frac{1}{\sigma_2}\right)-1 \implies \Phi\left(\frac{1}{\sigma_1}\right)>\Phi\left(\frac{1}{\sigma_2}\right),
\]
由 $\Phi(x)$ 单调递增，得 $\dfrac{1}{\sigma_1}>\dfrac{1}{\sigma_2}$，即 $\sigma_1<\sigma_2$，答案选 \textbf{(A)}。
\end{solution}


\begin{exercise}[正态分布参数求解] \label{exer:2-39}
设随机变量 $X$ 服从正态分布 $N(\mu,\sigma^2)(\sigma>0)$，且二次方程 $y^2+4y+X=0$ 无实根的概率为 $\dfrac{1}{2}$，则 $\mu=\underline{\hspace{2cm}}$。
\end{exercise}

\begin{solution}
先把“方程无实根”翻译成对随机变量 \(X\) 的条件。二次方程 \(y^2+4y+X=0\) 无实根的充要条件是判别式 \(\Delta=16-4X<0\)，即 \(X>4\)。因此题目给出的条件“无实根的概率为 \(\frac12\)”就等价于 \(P\{X>4\}=\frac12\)。

对正态分布 \(N(\mu,\sigma^2)\) 来说，分布关于均值 \(\mu\) 对称，所以 \(P\{X>\mu\}=\frac12\)。现在题目又告诉我们 \(P\{X>4\}=\frac12\)，这说明 4 正好就是这个正态分布的对称中心，也就是均值 \(\mu\)，故 \(\mu=4\)。

所以应填 \(\boxed{4}\)。
\end{solution}

\section{常见分布的期望与方差}

\begin{table}[h]
\centering
\caption{常见分布的数学期望与方差} \label{tab:dist-expect-var}
\begin{tabular}{c|c|c|c}
\toprule
名称与记号 & 分布列或概率密度 & 数学期望 & 方差 \\
\midrule
(0-1) 分布 $B(1,p)$ & $P\{X=1\}=p,\ P\{X=0\}=1-p$ & $p$ & $p(1-p)$ \\
二项分布 $B(n,p)$ & $P\{X=k\}=\mathrm{C}_n^k p^k (1-p)^{n-k}$ & $np$ & $np(1-p)$ \\
泊松分布 $P(\lambda)$ & $P\{X=k\}=e^{-\lambda}\frac{\lambda^k}{k!}$ & $\lambda$ & $\lambda$ \\
几何分布 $G(p)$ & $P\{X=k\}=(1-p)^{k-1}p$ & $\frac{1}{p}$ & $\frac{1-p}{p^2}$ \\
超几何分布 $H(n,N,M)$ & $P\{X=k\}=\frac{\mathrm{C}_M^k \mathrm{C}_{N-M}^{n-k}}{\mathrm{C}_N^n}$ & $\frac{nM}{N}$ & $\frac{nM}{N}\left(1-\frac{M}{N}\right)\frac{N-n}{N-1}$ \\
均匀分布 $U(a,b)$ & $f(x)=\frac{1}{b-a},\ a < x < b$ & $\frac{a+b}{2}$ & $\frac{(b-a)^2}{12}$ \\
指数分布 $E(\lambda)$ & $f(x)=\lambda e^{-\lambda x},\ x > 0$ & $\frac{1}{\lambda}$ & $\frac{1}{\lambda^2}$ \\
正态分布 $N(\mu,\sigma^2)$ & $f(x)=\frac{1}{\sqrt{2\pi}\sigma}e^{-\frac{(x-\mu)^2}{2\sigma^2}}$ & $\mu$ & $\sigma^2$ \\
\bottomrule
\end{tabular}
\end{table}

\begin{proof}[(0-1) 分布 $B(1,p)$ 的期望与方差推导]
由分布列 $P\{X=1\}=p,\ P\{X=0\}=1-p=q$，可得数学期望：
\[
E(X) = \sum_{k=0}^1 k P\{X=k\} = 0 \cdot (1-p) + 1 \cdot p = p.
\]
计算二阶矩：
\[
E(X^2) = \sum_{k=0}^1 k^2 P\{X=k\} = 0^2 \cdot (1-p) + 1^2 \cdot p = p.
\]
故方差为：
\[
D(X) = E(X^2) - [E(X)]^2 = p - p^2 = p(1-p).
\]
\end{proof}

\begin{proof}[二项分布 $B(n,p)$ 的期望与方差推导]
利用组合数性质 $k \mathrm{C}_n^k = n \mathrm{C}_{n-1}^{k-1}$，计算期望：
\[
\begin{aligned}
E(X) &= \sum_{k=0}^n k \mathrm{C}_n^k p^k (1-p)^{n-k} = \sum_{k=1}^n n \mathrm{C}_{n-1}^{k-1} p^k (1-p)^{n-k} \\
&= np \sum_{k=1}^n \mathrm{C}_{n-1}^{k-1} p^{k-1} (1-p)^{(n-1)-(k-1)}.
\end{aligned}
\]
令 $j = k-1$，由二项式定理，后方的求和等于 $(p + 1 - p)^{n-1} = 1$，故 $E(X) = np$。

为求方差，先计算 $E[X(X-1)]$，利用性质 $k(k-1)\mathrm{C}_n^k = n(n-1)\mathrm{C}_{n-2}^{k-2}$：
\[
\begin{aligned}
E[X(X-1)] &= \sum_{k=0}^n k(k-1) \mathrm{C}_n^k p^k (1-p)^{n-k} \\
&= \sum_{k=2}^n n(n-1) \mathrm{C}_{n-2}^{k-2} p^k (1-p)^{n-k} \\
&= n(n-1)p^2 \sum_{j=0}^{n-2} \mathrm{C}_{n-2}^j p^j (1-p)^{(n-2)-j} = n(n-1)p^2.
\end{aligned}
\]
由于 $D(X) = E(X^2) - [E(X)]^2 = E[X(X-1)] + E(X) - [E(X)]^2$，代入得：
\[
D(X) = n(n-1)p^2 + np - n^2p^2 = -np^2 + np = np(1-p).
\]
\end{proof}

\begin{proof}[泊松分布 $P(\lambda)$ 的期望与方差推导]
根据泰勒展开式 $\sum_{k=0}^\infty \frac{\lambda^k}{k!} = e^\lambda$。计算期望：
\[
E(X) = \sum_{k=0}^\infty k e^{-\lambda} \frac{\lambda^k}{k!} = \lambda e^{-\lambda} \sum_{k=1}^\infty \frac{\lambda^{k-1}}{(k-1)!} = \lambda e^{-\lambda} e^\lambda = \lambda.
\]
为求方差，先计算 $E[X(X-1)]$：
\[
E[X(X-1)] = \sum_{k=0}^\infty k(k-1) e^{-\lambda} \frac{\lambda^k}{k!} = \lambda^2 e^{-\lambda} \sum_{k=2}^\infty \frac{\lambda^{k-2}}{(k-2)!} = \lambda^2 e^{-\lambda} e^\lambda = \lambda^2.
\]
故方差为：
\[
D(X) = E[X(X-1)] + E(X) - [E(X)]^2 = \lambda^2 + \lambda - \lambda^2 = \lambda.
\]
\end{proof}

\begin{proof}[几何分布 $G(p)$ 的期望与方差推导]
记 $q = 1-p$。利用幂级数求和 $\sum_{k=1}^\infty q^k = \frac{q}{1-q}$，两边对 $q$ 求导得 $\sum_{k=1}^\infty k q^{k-1} = \frac{1}{(1-q)^2}$。
计算期望：
\[
E(X) = \sum_{k=1}^\infty k q^{k-1} p = p \cdot \frac{1}{(1-q)^2} = \frac{p}{p^2} = \frac{1}{p}.
\]
对幂级数再求一次导得 $\sum_{k=2}^\infty k(k-1) q^{k-2} = \frac{2}{(1-q)^3}$。计算 $E[X(X-1)]$：
\[
E[X(X-1)] = \sum_{k=1}^\infty k(k-1) q^{k-1} p = pq \sum_{k=2}^\infty k(k-1) q^{k-2} = pq \frac{2}{(1-q)^3} = \frac{2q}{p^2}.
\]
故方差为：
\[
D(X) = E[X(X-1)] + E(X) - [E(X)]^2 = \frac{2(1-p)}{p^2} + \frac{1}{p} - \frac{1}{p^2} = \frac{1-p}{p^2}.
\]
\end{proof}

\begin{proof}[超几何分布 $H(n,N,M)$ 的期望与方差推导]
采用引入示性变量（Indicator variables）的求法最为简便。设 $X_i$ 为第 $i$ 次抽出的物品是指定物品（次品/正品）的示性变量，若抽出则 $X_i=1$，否则 $X_i=0$。
易知 $P(X_i=1) = \frac{M}{N}$，故 $X_i \sim B(1, \frac{M}{N})$，则 $X = \sum_{i=1}^n X_i$。
期望为：
\[
E(X) = E\left(\sum_{i=1}^n X_i\right) = \sum_{i=1}^n E(X_i) = n \frac{M}{N}.
\]
对于方差，需考虑 $X_i$ 之间的协方差。当 $i \neq j$ 时：
\[
E(X_i X_j) = P(X_i=1, X_j=1) = \frac{M}{N} \cdot \frac{M-1}{N-1}.
\]
\[
\mathrm{Cov}(X_i, X_j) = E(X_i X_j) - E(X_i)E(X_j) = \frac{M(M-1)}{N(N-1)} - \frac{M^2}{N^2} = -\frac{M(N-M)}{N^2(N-1)}.
\]
由于 $D(X_i) = \frac{M}{N}\left(1-\frac{M}{N}\right)$，故总方差为：
\[
\begin{aligned}
D(X) &= \sum_{i=1}^n D(X_i) + \sum_{i \neq j} \mathrm{Cov}(X_i, X_j) \\
&= n \frac{M}{N}\left(1-\frac{M}{N}\right) + n(n-1)\left[-\frac{M(N-M)}{N^2(N-1)}\right] \\
&= n \frac{M}{N}\left(1-\frac{M}{N}\right) \left[1 - \frac{n-1}{N-1}\right] = \frac{nM}{N}\left(1-\frac{M}{N}\right)\frac{N-n}{N-1}.
\end{aligned}
\]
\end{proof}

\begin{proof}[均匀分布 $U(a,b)$ 的期望与方差推导]
根据定义，计算期望：
\[
E(X) = \int_{a}^{b} x \frac{1}{b-a} \,\mathrm{d}x = \frac{1}{b-a} \left[ \frac{1}{2}x^2 \right]_{a}^{b} = \frac{b^2-a^2}{2(b-a)} = \frac{a+b}{2}.
\]
计算二阶矩：
\[
E(X^2) = \int_{a}^{b} x^2 \frac{1}{b-a} \,\mathrm{d}x = \frac{1}{b-a} \left[ \frac{1}{3}x^3 \right]_{a}^{b} = \frac{b^3-a^3}{3(b-a)} = \frac{b^2+ab+a^2}{3}.
\]
故方差为：
\[
\begin{aligned}
D(X) = E(X^2) - [E(X)]^2 &= \frac{b^2+ab+a^2}{3} - \frac{a^2+2ab+b^2}{4} \\
&= \frac{4b^2+4ab+4a^2-3a^2-6ab-3b^2}{12} = \frac{(b-a)^2}{12}.
\end{aligned}
\]
\end{proof}

\begin{proof}[指数分布 $E(\lambda)$ 的期望与方差推导]
利用分部积分法，计算期望：
\[
E(X) = \int_{0}^{+\infty} x \lambda e^{-\lambda x} \,\mathrm{d}x = \left[ -x e^{-\lambda x} \right]_{0}^{+\infty} + \int_{0}^{+\infty} e^{-\lambda x} \,\mathrm{d}x = 0 + \left[ -\frac{1}{\lambda} e^{-\lambda x} \right]_{0}^{+\infty} = \frac{1}{\lambda}.
\]
计算二阶矩（再次利用分部积分）：
\[
\begin{aligned}
E(X^2) &= \int_{0}^{+\infty} x^2 \lambda e^{-\lambda x} \,\mathrm{d}x = \left[ -x^2 e^{-\lambda x} \right]_{0}^{+\infty} + \int_{0}^{+\infty} 2x e^{-\lambda x} \,\mathrm{d}x \\
&= 0 + \frac{2}{\lambda} \int_{0}^{+\infty} x \lambda e^{-\lambda x} \,\mathrm{d}x = \frac{2}{\lambda} E(X) = \frac{2}{\lambda^2}.
\end{aligned}
\]
故方差为：
\[
D(X) = E(X^2) - [E(X)]^2 = \frac{2}{\lambda^2} - \frac{1}{\lambda^2} = \frac{1}{\lambda^2}.
\]
\end{proof}

\begin{proof}[正态分布 $N(\mu,\sigma^2)$ 的期望与方差推导]
做变量代换 $t = \frac{x-\mu}{\sigma}$，则 $\mathrm{d}x = \sigma \,\mathrm{d}t$。计算期望：
\[
E(X) = \int_{-\infty}^{+\infty} x \frac{1}{\sqrt{2\pi}\sigma} e^{-\frac{(x-\mu)^2}{2\sigma^2}} \,\mathrm{d}x = \int_{-\infty}^{+\infty} (\sigma t + \mu) \frac{1}{\sqrt{2\pi}} e^{-\frac{t^2}{2}} \,\mathrm{d}t.
\]
由于 $t e^{-t^2/2}$ 是奇函数，在对称区间积分等于 $0$；而 $\frac{1}{\sqrt{2\pi}} \int e^{-t^2/2} \,\mathrm{d}t = 1$（标准正态密度全空间积分为1），故：
\[
E(X) = \frac{\sigma}{\sqrt{2\pi}} \int_{-\infty}^{+\infty} t e^{-\frac{t^2}{2}} \,\mathrm{d}t + \mu \int_{-\infty}^{+\infty} \frac{1}{\sqrt{2\pi}} e^{-\frac{t^2}{2}} \,\mathrm{d}t = 0 + \mu \cdot 1 = \mu.
\]
为求方差，直接计算中心矩 $E[(X-\mu)^2]$。同样代换 $t = \frac{x-\mu}{\sigma}$：
\[
D(X) = \int_{-\infty}^{+\infty} (x-\mu)^2 \frac{1}{\sqrt{2\pi}\sigma} e^{-\frac{(x-\mu)^2}{2\sigma^2}} \,\mathrm{d}x = \frac{\sigma^2}{\sqrt{2\pi}} \int_{-\infty}^{+\infty} t^2 e^{-\frac{t^2}{2}} \,\mathrm{d}t.
\]
使用分部积分法，令 $u=t, \mathrm{d}v = t e^{-t^2/2} \,\mathrm{d}t$，则 $v = -e^{-t^2/2}$：
\[
\int_{-\infty}^{+\infty} t \cdot \left( t e^{-\frac{t^2}{2}} \right) \,\mathrm{d}t = \left[ -t e^{-\frac{t^2}{2}} \right]_{-\infty}^{+\infty} + \int_{-\infty}^{+\infty} e^{-\frac{t^2}{2}} \,\mathrm{d}t = 0 + \sqrt{2\pi}.
\]
代回原式得：
\[
D(X) = \frac{\sigma^2}{\sqrt{2\pi}} \cdot \sqrt{2\pi} = \sigma^2.
\]
\end{proof}


\section{典型例题}

\begin{exercise}[拉普拉斯分布变换] 
设随机变量 $X$ 的概率密度为
\[
f(x)=\frac{1}{2\lambda}e^{-\frac{|x|}{\lambda}},\quad -\infty<x<+\infty,\ \lambda>0.
\]
\begin{enumerate}
    \item[(1)] 求方程 $y^2+2Xy+\lambda^2=0$ 有实根的概率；
    \item[(2)] 求 $Z=e^{-|X|}$ 的概率密度。
\end{enumerate}
\end{exercise}
\begin{solution}
(1) 方程有实根的条件是判别式
\[
\Delta = (2X)^2 - 4\lambda^2 = 4(X^2-\lambda^2) \geqslant 0 \Leftrightarrow |X|\geqslant\lambda.
\]
所求概率：
\[
\begin{aligned}
P\{|X|\geqslant\lambda\} &= \int_{-\infty}^{-\lambda} f(x)\,dx + \int_{\lambda}^{\infty} f(x)\,dx = 2\int_{\lambda}^{\infty} \frac{1}{2\lambda}e^{-\frac{x}{\lambda}}\,dx \\
&= \frac{1}{\lambda}\int_{\lambda}^{\infty} e^{-\frac{x}{\lambda}}\,dx = e^{-1}.
\end{aligned}
\]
(2) $Z=e^{-|X|}$ 的取值范围是 $(0,1]$。当 $0<z<1$ 时：
\[
\begin{aligned}
F_Z(z) &= P\{Z\leqslant z\} = P\{e^{-|X|}\leqslant z\} = P\{|X|\geqslant -\ln z\} \\
&= 2\int_{-\ln z}^{\infty} \frac{1}{2\lambda}e^{-\frac{x}{\lambda}}\,dx = z^{\frac{1}{\lambda}}.
\end{aligned}
\]

当 $z\leqslant0$ 时，$F_Z(z)=0$；当 $z\geqslant1$ 时，$F_Z(z)=1$。

求导得密度函数：
\[
f_Z(z) =
\begin{cases}
\dfrac{1}{\lambda}z^{\frac{1}{\lambda}-1}, & 0<z<1, \\[6pt]
0, & \text{其他}.
\end{cases}
\]
\end{solution}



\begin{exercise}\label{exer:2-14}
设随机变量 $X$ 的概率密度为 $f(x) = \begin{cases} 2^{-x}\ln 2, & x > 0, \\ 0, & \text{其他}, \end{cases}$ 对 $X$ 进行独立重复的观测，直到第 2 个大于 3 的观测值出现时停止，记 $Y$ 为观测次数，求 $Y$ 的概率分布。
\end{exercise}

\begin{solution}
记 $p$ 为"观测值大于 3"的概率，则
\[
p = P\{X > 3\} = \int_3^{+\infty} 2^{-x}\ln 2\,\mathrm{d}x = \frac{1}{8}.
\]
依题意，$Y$ 为离散型随机变量，取值为 $2,3,\dots$，则 $Y$ 的概率分布为
\[
P\{Y=k\} = \mathrm{C}_{k-1}^1 p(1-p)^{k-2} \cdot p = (k-1)p^2(1-p)^{k-2} = (k-1)\left(\frac{1}{8}\right)^2\left(\frac{7}{8}\right)^{k-2},\quad k=2,3,\dots.
\]
\end{solution}



\begin{exercise}[分布函数计算] \label{exer:2-35}
设随机变量 $X$ 的分布函数
\[
F(x)=
\begin{cases}
0, & x<0, \\
\dfrac{1}{2}, & 0\leqslant x<1, \\
1-e^{-x}, & x\geqslant1,
\end{cases}
\]
则 $P\{0\leqslant X\leqslant1\}=\underline{\hspace{2cm}}$，$P\{0<X<1\}=\underline{\hspace{2cm}}$。
\end{exercise}

\begin{solution}
\[
\begin{aligned}
P\{0\leqslant X\leqslant1\}&=F(1)-F(0^-)=1-e^{-1}; \\
P\{0<X<1\}&=F(1^-)-F(0)=\frac{1}{2}-\frac{1}{2}=0.
\end{aligned}
\]
\end{solution}

\begin{exercise}[离散型分布求解] \label{exer:2-36}
离散型随机变量 $X$ 的分布函数为
\[
F(x)=
\begin{cases}
0, & x<-1, \\
0.4, & -1\leqslant x<2, \\
0.8, & 2\leqslant x<3, \\
1, & x\geqslant3,
\end{cases}
\]
则 $X$ 的分布列为$\underline{\hspace{2cm}}$，$P\{X\leqslant3.2\}=\underline{\hspace{2cm}}$，方程 $x^2+Xx+1=0$ 有实根的概率为$\underline{\hspace{2cm}}$。
\end{exercise}

\begin{solution}
由分布函数可得分布列：
\[
P\{X=-1\}=F(-1)-F(-1^-)=0.4-0=0.4,
\]
\[
P\{X=2\}=F(2)-F(2^-)=0.8-0.4=0.4,
\]
\[
P\{X=3\}=F(3)-F(3^-)=1-0.8=0.2.
\]
故分布列：$X\sim\begin{pmatrix}-1 & 2 & 3 \\ 0.4 & 0.4 & 0.2\end{pmatrix}$。

$P\{X\leqslant3.2\}=1$（因为最大取值为3）。

方程有实根 $\Leftrightarrow \Delta=X^2-4\geqslant0 \Leftrightarrow X\geqslant2$ 或 $X\leqslant-2$（后者不可能），所以
\[
P(X\geqslant2)=P\{X=2\}+P\{X=3\}=0.4+0.2=0.6.
\]
答案：分布列如上；$\boldsymbol{1}$；$\boldsymbol{0.6}$。
\end{solution}

\begin{exercise}[连续型概率密度参数求解] \label{exer:2-37}
设连续型随机变量 $X$ 的概率密度为
\[
f(x)=
\begin{cases}
\dfrac{A}{x^2}, & x>100, \\[6pt]
0, & x\leqslant100,
\end{cases}
\]
则 $A=\underline{\hspace{2cm}}$，$P\{100<X\leqslant150\}=\underline{\hspace{2cm}}$，$P\{X\geqslant1000\}=\underline{\hspace{2cm}}$，$P\{X=1000\}=\underline{\hspace{2cm}}$，$X$ 的分布函数 $F(x)=\underline{\hspace{2cm}}$。
\end{exercise}

\begin{solution}
由归一性条件 $\displaystyle\int_{100}^{+\infty}\frac{A}{x^2}\,\mathrm{d}x=1$ 得
\[
A\left[-\frac{1}{x}\right]_{100}^{+\infty} = \frac{A}{100} = 1 \Rightarrow A=100.
\]
相关概率计算：
\[
P\{100<X\leqslant150\} = \int_{100}^{150}\frac{100}{x^2}\,\mathrm{d}x = 100\left(\frac{1}{100}-\frac{1}{150}\right) = \frac{1}{3},
\]
\[
P\{X\geqslant1000\} = \int_{1000}^{+\infty}\frac{100}{x^2}\,\mathrm{d}x = 100\cdot\frac{1}{1000} = \frac{1}{10},
\]
\[
P\{X=1000\}=0 \quad (\text{连续型随机变量单点概率为0}).
\]
分布函数：
\[
F(x) = \int_{100}^{x}\frac{100}{t^2}\,dt = 1-\frac{100}{x}, \quad x\geqslant100,
\]
故
\[
F(x)=
\begin{cases}
0, & x<100, \\[6pt]
1-\dfrac{100}{x}, & x\geqslant100.
\end{cases}
\]
答案：$\boldsymbol{100}$；$\boldsymbol{\dfrac{1}{3}}$；$\boldsymbol{\dfrac{1}{10}}$；$\boldsymbol{0}$；分布函数如上。
\end{solution}

\begin{exercise}[反正弦分布] \label{exer:2-38}
设连续型随机变量 $X$ 的分布函数为
\[
F(x)=
\begin{cases}
0, & x<-a, \\[6pt]
A+B\arcsin\dfrac{x}{a}, & -a\leqslant x\leqslant a, \\[6pt]
1, & x\geqslant a,
\end{cases}
\]
则 $A=\underline{\hspace{1cm}}$，$B=\underline{\hspace{1cm}}$，$X$ 的概率密度 $f(x)=\underline{\hspace{1cm}}$，$P\left\{-\dfrac{a}{2}<X<\dfrac{a}{2}\right\}=\underline{\hspace{1cm}}$。
\end{exercise}
\begin{solution}
由分布函数的连续性：
\[
\begin{cases}
\displaystyle\lim_{x\to(-a)^+}F(x) = A+B\arcsin(-1) = A-\frac{\pi}{2}B = 0, \\[10pt]
\displaystyle\lim_{x\to a^-}F(x) = A+B\arcsin(1) = A+\frac{\pi}{2}B = 1.
\end{cases}
\]
解得 $A=\dfrac{1}{2}$，$B=\dfrac{1}{\pi}$。\\
概率密度：
\[
f(x) = F'(x) = 
\begin{cases}
\dfrac{1}{\pi\sqrt{a^2-x^2}}, & -a<x<a, \\[6pt]
0, & \text{其他}.
\end{cases}
\]
概率计算：
\[
P\left\{-\frac{a}{2}<X<\frac{a}{2}\right\} = F\left(\frac{a}{2}\right)-F\left(-\frac{a}{2}\right) = \frac{1}{2}+\frac{1}{\pi}\arcsin\frac{1}{2} - \left[\frac{1}{2}+\frac{1}{\pi}\arcsin\left(-\frac{1}{2}\right)\right] = \frac{2}{\pi}\cdot\frac{\pi}{6} = \frac{1}{3}.
\]
\end{solution}


\begin{exercise}[同分布随机变量参数求解] \label{exer:2-42}
设随机变量 $X$ 和 $Y$ 同分布，$X$ 的概率密度为
\[
f(x)=
\begin{cases}
\dfrac{3}{8}x^2, & 0<x<2, \\[6pt]
0, & \text{其他}.
\end{cases}
\]
已知事件 $A=\{X>a\}$ 和 $B=\{Y>a\}$ 独立，且 $P(A\cup B)=\dfrac{3}{4}$，则常数 $a=\underline{\hspace{2cm}}$。
\end{exercise}
\begin{solution}
首先计算 $P(A)=P\{X>a\}$：
\[
P\{X>a\} = \int_a^2 \frac{3}{8}x^2\,\mathrm{d}x = \left[\frac{x^3}{8}\right]_a^2 = 1 - \frac{a^3}{8}.
\]
由于 $X$ 和 $Y$ 同分布，故 $P(A)=P(B)$。由独立性及加法公式：
\[
P(A\cup B) = P(A) + P(B) - P(AB) = 2P(A) - [P(A)]^2 = \frac{3}{4}.
\]
设 $p=P(A)$，解方程 $2p-p^2=\dfrac{3}{4}$ 得 $p=\dfrac{1}{2}$（舍去 $p=\dfrac{3}{2}$）。
于是
\[
1 - \frac{a^3}{8} = \frac{1}{2} \implies \frac{a^3}{8} = \frac{1}{2} \implies a^3 = 4 \implies a = \sqrt[3]{4}.
\]
\end{solution}


\begin{exercise}[分段概率密度的分布函数] \label{exer:2-47}
随机变量 $X$ 的概率密度为
\[
f(x)=
\begin{cases}
x, & 0\leqslant x<1, \\[6pt]
2-x, & 1\leqslant x<2, \\[6pt]
0, & \text{其他}.
\end{cases}
\]
求 $X$ 的分布函数 $F(x)$。
\end{exercise}
\begin{solution}
分段计算分布函数 $F(x) = \int_{-\infty}^{x} f(t)\,dt$：
\begin{enumerate}
    \item 当 $x<0$ 时：$F(x)=0$；
    \item 当 $0\leqslant x<1$ 时：
    \[
    F(x) = \int_0^x t\,dt = \frac{x^2}{2};
    \]
    \item 当 $1\leqslant x<2$ 时：
    \[
    F(x) = \int_0^1 t\,dt + \int_1^x (2-t)\,dt = \frac{1}{2} + \left[2t - \frac{t^2}{2}\right]_1^x = -\frac{x^2}{2} + 2x - 1;
    \]    
    \item 当 $x\geqslant2$ 时：$F(x)=1$。
\end{enumerate}
综上
\[
F(x)=
\begin{cases}
0, & x<0, \\[6pt]
\dfrac{x^2}{2}, & 0\leqslant x<1, \\[6pt]
-\dfrac{x^2}{2}+2x-1, & 1\leqslant x<2, \\[6pt]
1, & x\geqslant2.
\end{cases}
\]
\end{solution}

\begin{exercise}[分布函数性质判断]
下列函数可作为随机变量的分布函数的是（\quad）。
\[
\begin{aligned}
&\text{A. } F(x)=
\begin{cases}
0, & x<-1, \\[4pt]
\dfrac{1}{3}, & -1\leqslant x\leqslant2, \\[4pt]
1, & x>2,
\end{cases}
&&\text{B. } F(x)=
\begin{cases}
0, & x<0, \\[4pt]
\sin x, & 0\leqslant x<\pi, \\[4pt]
1, & x\geqslant\pi,
\end{cases} \\[10pt]
&\text{C. } F(x)=
\begin{cases}
0, & x<0, \\[4pt]
\dfrac{x+2}{5}, & 0\leqslant x<2, \\[4pt]
1, & x\geqslant2,
\end{cases}
&&\text{D. } F(x)=\dfrac{3}{4}+\dfrac{1}{2\pi}\arctan x.
\end{aligned}
\]
\end{exercise}
\begin{solution}
选项C正确。\textbf{分析}：
\begin{itemize}
    \item 选项A：在 $x=2$ 处，$F(2)=\dfrac{1}{3}$，但 $\displaystyle\lim_{x\to2^+}F(x)=1$，不满足右连续性；
    \item 选项B：$\sin x$ 在 $[0,\pi)$ 上不单调（先增后减），违反分布函数的单调不减性；
    \item 选项C：满足分布函数四条性质：单调不减、右连续、$F(-\infty)=0$、$F(+\infty)=1$；
    \item 选项D：$F(-\infty)=\dfrac{3}{4}-\dfrac{1}{4}=\dfrac{1}{2}\neq0$，不满足极限性质。
\end{itemize}
\end{solution}


\begin{exercise}[反正切分布函数]
已知随机变量 $X$ 的分布函数为
\[
F(x)=
\begin{cases}
A_1, & x<-a, \\[4pt]
A_2+A_3\arctan\dfrac{x}{a}, & x\in[-a,a], \\[6pt]
A_4, & x>a,
\end{cases}
\]
$a$ 为已知常数，求：
\begin{enumerate}
    \item[(1)] 使 $F(x)$ 连续的常数 $A_i\ (i=1,2,3,4)$；
    \item[(2)] 随机变量 $X$ 的概率密度函数；
    \item[(3)] 二次方程 $t^2+Xt+\dfrac{a^2}{16}=0$ 的两个根均为正根的概率。
\end{enumerate}
\end{exercise}
\begin{solution}
(1) 由分布函数极限性质：
\[
\lim_{x\to-\infty}F(x)=0,\quad \lim_{x\to+\infty}F(x)=1 \Rightarrow A_1=0,\ A_4=1.
\]
由连续性条件：
\[
\begin{cases}
F(-a)=A_2+A_3\arctan(-1)=A_2-\dfrac{\pi}{4}A_3=0, \\[6pt]
F(a)=A_2+A_3\arctan(1)=A_2+\dfrac{\pi}{4}A_3=1.
\end{cases}
\]
解得 $A_2=\dfrac{1}{2},\ A_3=\dfrac{2}{\pi}$。\\
(2) 概率密度为分布函数的导数：
\[
f(x)=F'(x)=
\begin{cases}
\dfrac{2a}{\pi(x^2+a^2)}, & -a\leqslant x\leqslant a, \\[6pt]
0, & \text{其他}.
\end{cases}
\]
(3) 方程有两个正根的条件：
\[
\begin{cases}
\Delta=X^2-\dfrac{a^2}{4}\geqslant0, \\[4pt]
X<0 \quad (\text{由韦达定理，根的和}=-X>0).
\end{cases}
\]
即 $X\leqslant-\dfrac{a}{2}$。所求概率：
\[
P\left\{X\leqslant-\frac{a}{2}\right\}=F\left(-\frac{a}{2}\right)=\frac{1}{2}-\frac{2}{\pi}\arctan\frac{1}{2}.
\]
\end{solution}


\begin{exercise}[分布函数线性组合]
设 $F_1(x),F_2(x)$ 分别为随机变量 $X_1,X_2$ 的分布函数，为使 $aF_1(x)-bF_2(x)$ 是某一随机变量的分布函数，在下列给定的各组数值中应取（\quad）。\\
(A)$a=\dfrac{3}{5}, b=-\dfrac{2}{5}$~~~(B)$a=\dfrac{2}{3}, b=\dfrac{2}{3}$~~~(C)$a=-\dfrac{1}{2}, b=\dfrac{3}{2}$~~~(D)$a=\dfrac{1}{2}, b=-\dfrac{3}{2}$。
\end{exercise}
\begin{solution}
答案选\textbf{(A)}。由分布函数性质，必须满足 $F(+\infty)=1$，即
\[
F(+\infty)=aF_1(+\infty)-bF_2(+\infty)=a-b=1.
\]
只有选项 (A) 满足 $a-b=\dfrac{3}{5}-\left(-\dfrac{2}{5}\right)=1$。其他选项：
\begin{itemize}
    \item (B): $a-b=0\neq1$；
    \item (C): $a-b=-2\neq1$；
    \item (D): $a-b=2\neq1$。
\end{itemize}
\end{solution}





\section{一维随机变量函数的分布} \label{sec:random-variable-function}

\subsection{基本概念} \label{subsec:function-concept}

设 $X$ 为随机变量，函数 $y = g(x)$ 在随机变量 $X$ 的取值范围内有定义，则以随机变量 $X$ 作为自变量的函数 $Y = g(X)$ 也是一个随机变量，称为随机变量 $X$ 的\textbf{函数}。例如：
\[
Y = aX^2 + bX + c,\quad Y = |X-a|,\quad Y = \begin{cases} X, & X \le 1, \\ 1, & X > 1 \end{cases}
\]
等都是随机变量的函数。本节主要研究如何由已知随机变量 $X$ 的分布来求其函数 $Y = g(X)$ 的分布。

\begin{note}
随机变量的函数在概率论与数理统计中具有重要应用。例如，在物理实验中，我们可能直接测量的是电压 $X$，但关心的却是功率 $Y = X^2/R$；在金融领域，资产的收益率经过某种变换后得到新的随机变量。因此，掌握随机变量函数的分布求法是解决实际问题的关键。
\end{note}

\subsection{离散型随机变量函数的分布} \label{subsec:discrete-to-discrete}

设 $X$ 为离散型随机变量，其概率分布为
\[
P\{X = x_i\} = p_i,\quad i=1,2,\dots
\]
则 $X$ 的函数 $Y = g(X)$ 也是离散型随机变量。若 $g(x_i)$ 的值互不相同，则 $Y$ 的分布律为
\[
P\{Y = g(x_i)\} = p_i,\quad i=1,2,\dots
\]
即
\[
Y \sim \begin{pmatrix} g(x_1) & g(x_2) & \cdots \\ p_1 & p_2 & \cdots \end{pmatrix}.
\]

如果有若干个 $g(x_i)$ 的值相同，则需要将这些相同值对应的概率相加合并。具体来说，若 $g(x_{i_1}) = g(x_{i_2}) = \cdots = g(x_{i_k}) = y_j$，则
\[
P\{Y = y_j\} = \sum_{t=1}^{k} P\{X = x_{i_t}\} = \sum_{t=1}^{k} p_{i_t}.
\]

\begin{property}\label{prop:discrete-function}
离散型随机变量函数的分布求解步骤：
\begin{enumerate}
    \item 确定 $Y = g(X)$ 的所有可能取值 $y_j$；
    \item 对每个 $y_j$，找出所有满足 $g(x_i) = y_j$ 的 $x_i$；
    \item 将相应的概率 $p_i$ 相加，得到 $P\{Y = y_j\}$。
\end{enumerate}
\end{property}

\subsection{连续型随机变量函数的分布} \label{subsec:continuous-function}

设 $X$ 为连续型随机变量，其概率密度函数为 $f_X(x)$，分布函数为 $F_X(x)$。对于连续型随机变量 $X$ 的函数 $Y = g(X)$，需要根据 $Y$ 的类型（离散型或连续型）采用不同的方法。

\subsubsection{当 \texorpdfstring{$Y$}{Y} 为离散型时}

若 $Y = g(X)$ 只能取有限个或可列个值，则 $Y$ 是离散型随机变量。此时应先确定 $Y$ 的所有可能取值，然后通过 $X$ 的概率分布计算 $Y$ 取各个值的概率。

具体来说，若 $Y$ 的可能取值为 $y_1, y_2, \dots, y_n, \dots$，则
\[
P\{Y = y_j\} = P\{g(X) = y_j\} = \int_{\{x: g(x) = y_j\}} f_X(x) \,\mathrm{d}x.
\]
这通常涉及到对 $X$ 的取值空间按 $g(x)$ 的取值进行划分。

\subsubsection{当 \texorpdfstring{$Y$}{Y} 为连续型时} 

若 $Y = g(X)$ 的取值充满一个区间（或若干区间的并），则 $Y$ 是连续型随机变量。此时主要有两种求解方法：分布函数法（定义法）和公式法。

\paragraph{方法一：分布函数法（通用方法）} \label{par:distribution-function-method}
根据分布函数的定义，先求 $Y$ 的分布函数 $F_Y(y)$，再对 $y$ 求导得到概率密度 $f_Y(y)$。

具体步骤如下：
\begin{enumerate}
    \item 对任意实数 $y$，求 $Y$ 的分布函数：
    \[
    F_Y(y) = P\{Y \le y\} = P\{g(X) \le y\} = \int_{\{x: g(x) \le y\}} f_X(x) \,\mathrm{d}x.
    \]
    这一步的关键是解不等式 $g(x) \le y$，确定积分区域 $\{x: g(x) \le y\}$。
    
    \item 对 $F_Y(y)$ 关于 $y$ 求导（在可导点处）：
    \[
    f_Y(y) = F_Y'(y).
    \]
    
    \item 检查 $f_Y(y)$ 是否满足概率密度函数的基本条件：非负性和归一性。
\end{enumerate}

\begin{note}
分布函数法是求连续型随机变量函数分布的最通用方法，无论 $g(x)$ 是单调还是非单调、是连续还是分段连续，都适用。虽然计算过程可能相对复杂，但其思想直观、逻辑严密，是理解随机变量函数分布的基础。
\end{note}

\paragraph{方法二：公式法（单调函数情形）} \label{par:formula-method}
当 $y = g(x)$ 是严格单调函数时，可以使用公式法更快捷地求得 $Y$ 的概率密度。

\begin{theorem}[随机变量函数的密度公式] \label{thm:function-density-formula}
设连续型随机变量 $X$ 的概率密度为 $f_X(x)$，函数 $y = g(x)$ 在 $X$ 的取值区间 $(a,b)$ 上严格单调且可导，其反函数为 $x = h(y)$，且 $h(y)$ 连续可导。则 $Y = g(X)$ 的概率密度为
\[
f_Y(y) = \begin{cases}
f_X[h(y)] \cdot |h'(y)|, & \alpha < y < \beta, \\
0, & \text{其他},
\end{cases}
\]
其中 $(\alpha, \beta)$ 是函数 $y = g(x)$ 的值域，即
\[
\alpha = \min\left\{\lim_{x \to a^+} g(x), \lim_{x \to b^-} g(x)\right\}, \quad \beta = \max\left\{\lim_{x \to a^+} g(x), \lim_{x \to b^-} g(x)\right\}.
\]
\end{theorem}

\begin{proof}
考虑 $g(x)$ 严格单调递增的情形（单调递减类似）：
\[
F_Y(y) = P\{Y \le y\} = P\{g(X) \le y\} = P\{X \le h(y)\} = F_X[h(y)].
\]
两边对 $y$ 求导，由链式法则得
\[
f_Y(y) = F_Y'(y) = f_X[h(y)] \cdot h'(y).
\]
由于 $h'(y) > 0$（递增情形），故 $f_Y(y) = f_X[h(y)] \cdot |h'(y)|$。对于单调递减情形，$h'(y) < 0$，此时
\[
F_Y(y) = P\{g(X) \le y\} = P\{X \ge h(y)\} = 1 - F_X[h(y)],
\]
求导后得 $f_Y(y) = -f_X[h(y)] \cdot h'(y) = f_X[h(y)] \cdot |h'(y)|$。综上，公式成立。
\end{proof}

\begin{remark}
公式法的适用条件：
\begin{enumerate}
    \item $g(x)$ 必须是严格单调函数（递增或递减）；
    \item $g(x)$ 在 $X$ 的取值区间内可导；
    \item 若 $g(x)$ 在若干区间上分段单调，则需要在每个单调区间上分别应用公式法，然后将结果相加（这实际上等价于分布函数法）。
\end{enumerate}
当 $g(x)$ 不满足严格单调条件时，必须使用分布函数法。
\end{remark}


\begin{note}若 $X$ 既非离散型也非连续型（或形式未知），则只能使用分布函数的定义来求 $Y$ 的分布：
\[
F_Y(y) = P\{Y \le y\} = P\{g(X) \le y\}.
\]
这种方法不依赖于 $X$ 的具体类型，具有完全的通用性。
\end{note}

\section{典型例题}

\begin{exercise}[线性变换与平方变换] \label{exer:2-18}
设连续型随机变量 $X$ 的概率密度函数为
\[
f_X(x) = \begin{cases}
1, & -\dfrac{1}{2} < x < \dfrac{1}{2}, \\[6pt]
0, & \text{其他},
\end{cases}
\]
求下列随机变量 $Y$ 的概率密度函数：
\begin{enumerate}
    \item[(1)] $Y = a + bX$, 其中 $b \neq 0$；
    \item[(2)] $Y = X^2$；
    \item[(3)] $Y = |X|$。
\end{enumerate}
\end{exercise}
\begin{solution}
(1) \textbf{线性变换}：使用分布函数法求 $Y=a+bX$ 的概率密度。设 $Y$ 的分布函数为 $F_Y(y)$，由于 $b \neq 0$，需分情况讨论：
\begin{itemize}
\item 当 $b > 0$ 时：
\begin{align*}
F_Y(y) &= P\{Y \le y\} = P\{a + bX \le y\} = P\left\{X \le \frac{y-a}{b}\right\} = F_X\left(\frac{y-a}{b}\right). \\
\intertext{对 $y$ 求导可得：}
f_Y(y) &= F_Y'(y) = f_X\left(\frac{y-a}{b}\right) \cdot \frac{1}{b}.
\end{align*}
因此，当 $b > 0$ 时，$Y$ 的取值范围是
\[
\left(a-\dfrac{b}{2},\ a+\dfrac{b}{2}\right).
\]
\item 当 $b < 0$ 时：
\begin{align*}
F_Y(y) &= P\{Y \le y\} = P\{a + bX \le y\} = P\left\{X \ge \frac{y-a}{b}\right\} = 1 - F_X\left(\frac{y-a}{b}\right). \\
f_Y(y) &= F_Y'(y) = -f_X\left(\frac{y-a}{b}\right) \cdot \frac{1}{b} = f_X\left(\frac{y-a}{b}\right) \cdot \frac{1}{-b}.
\end{align*}
由于 $-\dfrac{1}{2} < X < \dfrac{1}{2}$ 且 $b < 0$，此时 $Y$ 的取值范围是 $\left(a+\dfrac{b}{2}, a-\dfrac{b}{2}\right)$。
\end{itemize}
综合以上两种情况，无论是 $b>0$（此时 $b=|b|$）还是 $b<0$（此时 $-b=|b|$），都可以统一写为：
\[
f_Y(y) = f_X\left(\frac{y-a}{b}\right) \cdot \frac{1}{|b|}.
\]
取值区间统一为 $\left(a-\dfrac{|b|}{2}, a+\dfrac{|b|}{2}\right)$。在该区间内，$f_X(\cdot) = 1$。因此 $Y$ 的概率密度函数为：
\[
f_Y(y) = \begin{cases}
\dfrac{1}{|b|}, & a - \dfrac{|b|}{2} < y < a + \dfrac{|b|}{2}, \\[8pt]
0, & \text{其他}.
\end{cases}
\]
(2) \textbf{平方变换}：函数 $y = x^2$ 在 $(-\dfrac{1}{2}, \dfrac{1}{2})$ 上不是单调函数，必须使用分布函数法。由于 $Y = X^2 \ge 0$，当 $y \le 0$ 时，$F_Y(y) = 0$。当 $0 < y < \dfrac{1}{4}$ 时：
\[
F_Y(y) = P\{X^2 \le y\} = P\{-\sqrt{y} \le X \le \sqrt{y}\} = \int_{-\sqrt{y}}^{\sqrt{y}} 1 \, \mathrm{d}x = 2\sqrt{y}.
\]
当 $y \ge \dfrac{1}{4}$ 时，$F_Y(y) = 1$。
求导得概率密度：
\[
f_Y(y) = \begin{cases}
\dfrac{1}{\sqrt{y}}, & 0 < y < \dfrac{1}{4}, \\[6pt]
0, & \text{其他}.
\end{cases}
\]
(3) \textbf{绝对值变换}：同样使用分布函数法。由于 $Y = |X| \ge 0$，当 $y < 0$ 时，$F_Y(y) = 0$。当 $0 \le y < \dfrac{1}{2}$ 时：
\[
F_Y(y) = P\{|X| \le y\} = P\{-y \le X \le y\} = \int_{-y}^{y} 1 \, \mathrm{d}x = 2y.
\]
当 $y \ge \dfrac{1}{2}$ 时，$F_Y(y) = 1$。求导得：
\[
f_Y(y) = \begin{cases}
2, & 0 \le y < \dfrac{1}{2}, \\[6pt]
0, & \text{其他}.
\end{cases}
\]
\end{solution}

\begin{exercise}[离散型随机变量的函数] \label{exer:2-19}
设随机变量 $X$ 的分布律为
\[
P\{X=-1\} = \frac{1}{4},\quad P\{X=0\} = \frac{1}{2},\quad P\{X=1\} = \frac{1}{4},
\]
求 $Y = X^2$ 的分布。
\end{exercise}

\begin{solution}
$Y = X^2$ 的可能取值为 $0$ 和 $1$。

当 $Y=0$ 时：
\[
P\{Y=0\} = P\{X^2=0\} = P\{X=0\} = \frac{1}{2}.
\]

当 $Y=1$ 时：
\[
P\{Y=1\} = P\{X^2=1\} = P\{X=1 \text{或} X=-1\} = P\{X=1\} + P\{X=-1\} = \frac{1}{4} + \frac{1}{4} = \frac{1}{2}.
\]

因此，$Y$ 的分布律为
\[
Y \sim \begin{pmatrix} 0 & 1 \\[4pt] \dfrac{1}{2} & \dfrac{1}{2} \end{pmatrix}.
\]
\end{solution}

\begin{exercise}[概率积分变换] \label{exer:2-20}
设随机变量 $X$ 的分布函数 $F_X(x)$ 是严格单调增加函数，其反函数 $F_X^{-1}(y)$ 存在，$Y = F_X(X)$。证明：$Y$ 服从区间 $(0,1)$ 上的均匀分布。
\end{exercise}

\begin{solution}
要证明 $Y$ 服从均匀分布，最直接的方法就是求它的分布函数 $F_Y(y)=P\{Y\le y\}$。

由于 $F_X(x)\in[0,1]$，所以 $Y=F_X(X)$ 的取值只可能落在 $[0,1]$ 内。下面分段讨论。

当 $y < 0$ 时，$F_Y(y) = P\{Y \le y\} = 0$。

当 $0 \le y < 1$ 时，因为 $F_X$ 严格单调增加，
\[
F_Y(y)=P\{Y\le y\}=P\{F_X(X)\le y\}=P\{X\le F_X^{-1}(y)\}=F_X\!\bigl(F_X^{-1}(y)\bigr)=y.
\]
这一步正是题目要求 $F_X$ 严格单调、反函数存在的原因。

当 $y \ge 1$ 时，$F_Y(y) = 1$。

综上，$Y$ 的分布函数为
\[
F_Y(y) = \begin{cases}
0, & y < 0, \\
y, & 0 \le y < 1, \\
1, & y \ge 1,
\end{cases}
\]
这正是均匀分布 $U(0,1)$ 的分布函数，因此 \(Y\sim U(0,1)\)。

换句话说，\textbf{任何连续型随机变量经过自己的分布函数变换后，都会变成标准均匀分布}。
\end{solution}

\begin{remark}
概率积分变换是统计模拟和假设检验中的重要工具。它表明，任何连续型随机变量都可以通过其分布函数变换为标准均匀分布。反之，若 $U \sim U(0,1)$，则 $X = F_X^{-1}(U)$ 服从分布 $F_X(x)$。这为随机数生成提供了理论基础。
\end{remark}

\begin{exercise}[分段随机变量的函数] \label{exer:2-21}
在区间 $(0,2)$ 上随机取一点，将该区间分成两段，较短一段的长度记为 $X$，较长一段的长度记为 $Y$。令 $Z = \dfrac{Y}{X}$。(1) 求 $X$ 的概率密度；(2) 求 $Z$ 的概率密度。
\end{exercise}
\begin{solution}
(1) 设随机点的坐标为 $V \sim U(0,2)$，由于区间总长为 $2$，较短一段的长度 $X = \min\{V, 2-V\}$。显然 $X$ 的取值范围为 $(0,1)$。
先求 $X$ 的分布函数 $F_X(x) = P\{X \le x\}$：
当 $x \le 0$ 时，$F_X(x) = 0$；
当 $x \ge 1$ 时，$F_X(x) = 1$。
当 $0 < x < 1$ 时，利用对立事件有：
\[
\begin{aligned}
F_X(x) &= P\{\min\{V, 2-V\} \le x\} = 1 - P\{\min\{V, 2-V\} > x\} \\
&= 1 - P\{V > x \text{ 且 } 2-V > x\} = 1 - P\{x < V < 2-x\}
\end{aligned}
\]
因为 $V$ 在 $(0,2)$ 上服从均匀分布，其密度函数为 $\frac{1}{2}$，所以：
\[
P\{x < V < 2-x\} = \int_{x}^{2-x} \frac{1}{2} \,\mathrm{d}v = \frac{(2-x) - x}{2} = 1 - x
\]
从而 $F_X(x) = 1 - (1 - x) = x$。对分布函数求导，得到 $X$ 的概率密度为：
\[
f_X(x) = F'_X(x) = \begin{cases}
1, & 0 < x < 1, \\[4pt]
0, & \text{其他}.
\end{cases}
\]
（即 $X \sim U(0,1)$）

(2) 由于 $Y = 2 - X$，故 $Z = \dfrac{Y}{X} = \dfrac{2-X}{X} = \dfrac{2}{X} - 1$。
因为 $X \in (0,1)$，所以 $Z$ 的取值范围为 $(1, +\infty)$。
现在用分布函数法求 $Z$ 的概率密度。设 $Z$ 的分布函数为 $F_Z(z) = P\{Z \le z\}$：
当 $z \le 1$ 时，$F_Z(z) = 0$。
当 $z > 1$ 时：
\[
\begin{aligned}
F_Z(z) &= P\left\{\frac{2}{X} - 1 \le z\right\} = P\left\{\frac{2}{X} \le z + 1\right\} = P\left\{X \ge \frac{2}{z+1}\right\}
\end{aligned}
\]
由于 $X \sim U(0,1)$，当 $z > 1$ 时，$\dfrac{2}{z+1} \in (0, 1)$，因此：
\[
F_Z(z) = \int_{\frac{2}{z+1}}^{1} f_X(x) \,\mathrm{d}x = \int_{\frac{2}{z+1}}^{1} 1 \,\mathrm{d}x = 1 - \frac{2}{z+1} = \frac{z-1}{z+1}
\]
对分布函数求导，得到 $Z$ 的概率密度为：
\[
f_Z(z) = F'_Z(z) = \left( 1 - \frac{2}{z+1} \right)' = \frac{2}{(z+1)^2}, \quad z > 1.
\]
故 $Z$ 的概率密度为：
\[
f_Z(z) = \begin{cases}
\dfrac{2}{(1+z)^2}, & z > 1, \\[6pt]
0, & \text{其他}.
\end{cases}
\]
\end{solution}

\begin{exercise}[分段密度函数的变换] \label{exer:2-22}
设随机变量 $X$ 的概率密度为
\[
f_X(x) = \begin{cases}
\dfrac{1}{2}, & -1 < x < 0, \\[6pt]
\dfrac{1}{4}, & 0 \le x < 2, \\[6pt]
0, & \text{其他},
\end{cases}
\]
令 $Y = X^2$，求 $Y$ 的概率密度 $f_Y(y)$。
\end{exercise}

\begin{solution}
使用分布函数法。由于 $Y = X^2 \ge 0$，当 $y < 0$ 时，$F_Y(y) = 0$。

当 $0 \le y < 1$ 时：
\[
\begin{aligned}
F_Y(y) &= P\{-\sqrt{y} \le X \le \sqrt{y}\} \\
&= P\{-\sqrt{y} \le X < 0\} + P\{0 \le X \le \sqrt{y}\} \\
&= \frac{1}{2}\sqrt{y} + \frac{1}{4}\sqrt{y} = \frac{3}{4}\sqrt{y}.
\end{aligned}
\]

当 $1 \le y < 4$ 时（注意 $-1 < x < 0$ 对应 $0 < y < 1$，而 $0 \le x < 2$ 对应 $0 \le y < 4$）：
\[
F_Y(y) = P\{-1 < X < 0\} + P\{0 \le X \le \sqrt{y}\} = \frac{1}{2} + \frac{1}{4}\sqrt{y}.
\]

当 $y \ge 4$ 时，$F_Y(y) = 1$。

求导得概率密度：
\[
f_Y(y) = \begin{cases}
\dfrac{3}{8\sqrt{y}}, & 0 < y < 1, \\[6pt]
\dfrac{1}{8\sqrt{y}}, & 1 \le y < 4, \\[6pt]
0, & \text{其他}.
\end{cases}
\]

此题体现了当 $X$ 的密度函数分段定义时，$Y = g(X)$ 的分布需要分段讨论，特别注意变换后不同区间的合并与拆分。
\end{solution}


\begin{exercise}[离散型随机变量函数分布] \label{exer:2-43}
已知 $X$ 为随机变量，$Y=X^2+1$，已知 $X$ 的概率分布为 $P\{X=-1\}=P\{X=0\}=P\{X=1\}=\dfrac{1}{3}$，则 $Y$ 的分布函数 $F_Y(y)=\underline{\hspace{2cm}}$。
\end{exercise}
\begin{solution}
$Y$ 的可能取值为：
\begin{itemize}
    \item 当 $X=-1$ 或 $X=1$ 时，$Y=(-1)^2+1=1^2+1=2$
    \item 当 $X=0$ 时，$Y=0^2+1=1$
\end{itemize}
因此
\[
P\{Y=1\}=P\{X=0\}=\frac{1}{3}, \quad P\{Y=2\}=P\{X=-1\}+P\{X=1\}=\frac{2}{3}.
\]
分布函数：
\[
F_Y(y)=
\begin{cases}
0, & y<1, \\[6pt]
\dfrac{1}{3}, & 1\leqslant y<2, \\[6pt]
1, & y\geqslant2.
\end{cases}
\]
\end{solution}


\begin{exercise}\label{exer:2-30}
设随机变量 $X$ 服从参数为 2 的指数分布，求 $Y=1-e^{-2X}$ 的概率密度 $f_Y(y)$。
\end{exercise}

\begin{solution}
由题意，$X \sim E(2)$，概率密度为 $f_X(x) = \begin{cases} 2e^{-2x}, & x>0 \\ 0, & x \le 0 \end{cases}$。\\
函数 $y = 1-e^{-2x}$ 在 $x>0$ 上严格单调递增，值域为 $(0,1)$。反函数为 $x = -\dfrac{1}{2}\ln(1-y)$，导数为
\[
\frac{\mathrm{d}x}{\mathrm{d}y} = \frac{1}{2(1-y)}.
\]
由公式法：
\begin{align*}
f_Y(y) &= f_X\!\left(-\frac{1}{2}\ln(1-y)\right) \cdot \left|\frac{1}{2(1-y)}\right| \\
&= 2e^{-2 \cdot \left[-\frac{1}{2}\ln(1-y)\right]} \cdot \frac{1}{2(1-y)} \\
&= 2(1-y) \cdot \frac{1}{2(1-y)} = 1, \quad 0 < y < 1.
\end{align*}
故
\[
f_Y(y) = \begin{cases}
1, & 0 < y < 1, \\[4pt]
0, & \text{其他}.
\end{cases}
\]
即 $Y \sim U(0,1)$。
\end{solution}

\begin{exercise}\label{exer:2-31}
设随机变量 $X \sim N(4,16)$，求 $Y=|X-4|$ 的概率密度 $f_Y(y)$。
\end{exercise}

\begin{solution}
令 $Z = X-4$，则 $Z \sim N(0,16)$，概率密度为 $f_Z(z) = \dfrac{1}{\sqrt{32\pi}} e^{-\frac{z^2}{32}}$。

函数 $y = |z|$ 在 $(-\infty,+\infty)$ 上不是严格单调函数，使用分布函数法。

由于 $Y = |Z| \ge 0$，当 $y \le 0$ 时，$F_Y(y) = 0$。

当 $y > 0$ 时：
\begin{align*}
F_Y(y) &= P\{|Z| \le y\} = P\{-y \le Z \le y\} = F_Z(y) - F_Z(-y).
\end{align*}

求导得：
\[
f_Y(y) = F_Z'(y) - F_Z'(-y) \cdot (-1) = f_Z(y) + f_Z(-y) = \frac{2}{\sqrt{32\pi}} e^{-\frac{y^2}{32}}.
\]
故
\[
f_Y(y) = \begin{cases}
\dfrac{1}{\sqrt{8\pi}} e^{-\frac{y^2}{32}}, & y > 0, \\[6pt]
0, & y \le 0.
\end{cases}
\]
\end{solution}


\begin{exercise}[柯西分布变换] \label{exer:2-49}
设 $f(x)=\dfrac{k}{1+x^2}$，求：(1) $k$~~~(2)$Y=1-\sqrt[3]{X}$ 的概率密度。
\end{exercise}
\begin{solution}
(1) 由概率密度的归一性条件：
\[
\int_{-\infty}^{+\infty} \frac{k}{1+x^2}\,\mathrm{d}x = k\left[\arctan x\right]_{-\infty}^{+\infty} = k\left(\frac{\pi}{2} - \left(-\frac{\pi}{2}\right)\right) = k\pi = 1 \implies k = \frac{1}{\pi}.
\]
即 $X$ 服从标准柯西分布，概率密度为 $f_X(x) = \frac{1}{\pi(1+x^2)}$。

(2) 使用分布函数法。设 $Y$ 的分布函数为 $F_Y(y)$，由于 $X \in (-\infty, +\infty)$，显然 $Y \in (-\infty, +\infty)$。
\begin{align*}
F_Y(y) &= P\{Y \le y\} \\
&= P\{1-\sqrt[3]{X} \le y\} \\
&= P\{\sqrt[3]{X} \ge 1-y\} 
\end{align*}
由于函数 $g(t) = t^3$ 在定义域上单调递增，两边同时立方不等号方向不变：
\begin{align*}
F_Y(y) &= P\{X \ge (1-y)^3\} \\
&= 1 - P\{X \le (1-y)^3\} \quad \text{（因 $X$ 是连续型，等号处概率为 $0$）} \\
&= 1 - F_X\left((1-y)^3\right)
\end{align*}

对 $y$ 求导得到 $Y$ 的概率密度 $f_Y(y)$，利用复合函数求导法则：
\begin{align*}
f_Y(y) &= F_Y'(y) = \frac{\mathrm{d}}{\mathrm{d}y} \left[ 1 - F_X\left((1-y)^3\right) \right] \\
&= -f_X\left((1-y)^3\right) \cdot \left[3(1-y)^2 \cdot (-1)\right] \\
&= 3(1-y)^2 f_X\left((1-y)^3\right)
\end{align*}

将 $f_X(x) = \dfrac{1}{\pi(1+x^2)}$ 代入上式：
\[
f_Y(y) = 3(1-y)^2 \cdot \frac{1}{\pi\left[1+\left((1-y)^3\right)^2\right]} = \frac{3(1-y)^2}{\pi\left[1+(1-y)^6\right]}, \quad y\in\mathbb{R}.
\]

答案：(1) $\boldsymbol{k=\dfrac{1}{\pi}}$；(2) $\boldsymbol{f_Y(y)=\dfrac{3(1-y)^2}{\pi[1+(1-y)^6]}}$。
\end{solution}

\begin{exercise}[柯西分布尺度变换]
设随机变量 $X$ 的概率密度 $f(x)=\dfrac{1}{\pi(1+x^2)}$，则 $Y=2X$ 的分布密度为（\quad）。~~~~
（A） $\dfrac{1}{\pi(1+4y^2)}$；~~~~（B）$\dfrac{2}{\pi(4+y^2)}$；~~~~（C）$\dfrac{1}{\pi(1+y^2)}$；~~~~（D）$\dfrac{1}{\pi}\arctan y$。
\end{exercise}
\begin{solution}
答案选\textbf{(B)}。

使用分布函数法求解。设 $Y$ 的分布函数为 $F_Y(y)$，则：
\[
F_Y(y) = P\{Y \le y\} = P\{2X \le y\} = P\left\{X \le \frac{y}{2}\right\} = F_X\left(\frac{y}{2}\right)
\]
两边对 $y$ 求导，得到 $Y$ 的概率密度 $f_Y(y)$：
\[
f_Y(y) = F_Y'(y) = f_X\left(\frac{y}{2}\right) \cdot \frac{1}{2}
\]
将 $X$ 的概率密度 $f_X(x) = \dfrac{1}{\pi(1+x^2)}$ 代入：
\[
f_Y(y) = \frac{1}{\pi\left[1+\left(\frac{y}{2}\right)^2\right]} \cdot \frac{1}{2} = \frac{1}{\pi\left(1+\frac{y^2}{4}\right)} \cdot \frac{1}{2} = \frac{4}{\pi(4+y^2)} \cdot \frac{1}{2} = \frac{2}{\pi(4+y^2)}
\]
故选项 (B) 正确。
\end{solution}

\begin{exercise}[正态分布平方变换] 
设随机变量 $X\sim N(0,1)$，求 $Y=2X^2+1$ 的概率密度函数。
\end{exercise}
\begin{solution}
当 $y\leqslant1$ 时，$F_Y(y)=0$，故 $f_Y(y)=0$。\\
当 $y>1$ 时：
\[
\begin{aligned}
F_Y(y) &= P\{2X^2+1\leqslant y\} = P\left\{X^2\leqslant\frac{y-1}{2}\right\} = P\left\{-\sqrt{\frac{y-1}{2}}\leqslant X\leqslant\sqrt{\frac{y-1}{2}}\right\} \\
&= 2\Phi\left(\sqrt{\frac{y-1}{2}}\right)-1,
\end{aligned}
\]
其中 $\Phi$ 为标准正态分布函数。求导得：
\[
f_Y(y) = \frac{\mathrm{d}}{\mathrm{d}y}F_Y(y) = 2\varphi\left(\sqrt{\frac{y-1}{2}}\right)\cdot\frac{1}{2\sqrt{2(y-1)}},
\]
其中 $\varphi(x)=\dfrac{1}{\sqrt{2\pi}}e^{-x^2/2}$ 为标准正态密度。化简得：
\[
f_Y(y) =
\begin{cases}
\dfrac{1}{2\sqrt{\pi(y-1)}}\exp\left\{-\dfrac{y-1}{4}\right\}, & y>1, \\[8pt]
0, & y\leqslant1.
\end{cases}
\]
\end{solution}

\begin{exercise}[逻辑分布变换] \label{exer:2-50}
设 $f(x)=\dfrac{e^x}{(1+e^x)^2}$，$Y=e^X$，求$X$ 的分布函数和$Y$ 的概率密度。
\end{exercise}
\begin{solution}
(1) 求 $X$ 的分布函数 $F_X(x)$：
\begin{align*}
F_X(x) &= \int_{-\infty}^{x} \frac{e^t}{(1+e^t)^2}\,\mathrm{d}t \\
&= \left[-\frac{1}{1+e^t}\right]_{-\infty}^{x} \\
&= \left(-\frac{1}{1+e^x}\right) - \left(-\frac{1}{1+0}\right) \\
&= 1 - \frac{1}{1+e^x} = \frac{e^x}{1+e^x}, \quad x\in\mathbb{R}.
\end{align*}
这是标准的逻辑分布（Logistic distribution）。

(2) 使用分布函数法求 $Y$ 的概率密度。设 $Y$ 的分布函数为 $F_Y(y)$。
由于 $Y = e^X > 0$，所以当 $y \le 0$ 时，$F_Y(y) = P\{Y \le y\} = 0$。

当 $y > 0$ 时：
\[
F_Y(y) = P\{Y \le y\} = P\{e^X \le y\} = P\{X \le \ln y\} = F_X(\ln y).
\]
将第 (1) 问求得的 $F_X(x) = \dfrac{e^x}{1+e^x}$ 代入上式，可得：
\[
F_Y(y) = \frac{e^{\ln y}}{1+e^{\ln y}} = \frac{y}{1+y}.
\]
对 $y$ 求导即可得到 $Y$ 的概率密度 $f_Y(y)$。当 $y > 0$ 时：
\[
f_Y(y) = F_Y'(y) = \frac{\mathrm{d}}{\mathrm{d}y} \left( \frac{y}{1+y} \right) = \frac{1 \cdot (1+y) - y \cdot 1}{(1+y)^2} = \frac{1}{(1+y)^2}.
\]
当 $y \le 0$ 时，$f_Y(y) = 0$。

综合可得：
\[
f_Y(y)=
\begin{cases}
\dfrac{1}{(1+y)^2}, & y>0, \\[8pt]
0, & y\leqslant0.
\end{cases}
\]
答案：(1) $\boldsymbol{F_X(x)=\dfrac{e^x}{1+e^x}}$；(2) 密度函数如上。
\end{solution}


\begin{exercise}[均匀分布绝对值变换] \label{exer:2-51}
设 $X\sim U(-1,2)$，求 $Y=|X|$ 的概率密度。
\end{exercise}
\begin{solution}
$X$ 的概率密度为
\[
f_X(x)=
\begin{cases}
\dfrac{1}{3}, & -1<x<2, \\[6pt]
0, & \text{其他}.
\end{cases}
\]
函数 $y=|x|$ 在区间 $(-1,2)$ 上不是单调函数，使用分布函数法。由于 $Y=|X|\in[0,2)$：
\begin{itemize}
    \item 当 $0<y<1$ 时：
    \[
    F_Y(y) = P\{|X|\leqslant y\} = P\{-y\leqslant X\leqslant y\} = \frac{y-(-y)}{3} = \frac{2y}{3},
    \]
    故 $f_Y(y)=\dfrac{2}{3}$；
    \item 当 $1\leqslant y<2$ 时：
    \[
    F_Y(y) = P\{-1\leqslant X\leqslant y\} = \frac{y-(-1)}{3} = \frac{y+1}{3},
    \]
    故 $f_Y(y)=\dfrac{1}{3}$；
    \item 当 $y\leqslant0$ 时，$F_Y(y)=0$；当 $y\geqslant2$ 时，$F_Y(y)=1$。
\end{itemize}
综上，
\[
f_Y(y)=
\begin{cases}
\dfrac{2}{3}, & 0<y<1, \\[6pt]
\dfrac{1}{3}, & 1\leqslant y<2, \\[6pt]
0, & \text{其他}.
\end{cases}
\]
\end{solution}

\begin{table}[htbp]
\centering
\caption{常见分布速查表}
\label{tab:distribution_summary}
\resizebox{\textwidth}{!}{
\begin{tabular}{llllll}
\toprule
\textbf{分布} & \textbf{记号与参数} & \textbf{分布律/密度} & $EX$ & $DX$ & \textbf{典型场景} \\
\midrule
\multicolumn{6}{c}{\textbf{离散型}} \\
\midrule
两点分布 & $B(1,p)$ & $P(X=1)=p$ & $p$ & $p(1-p)$ & 单次试验成功/失败 \\
二项分布 & $B(n,p)$ & $\binom{n}{k}p^k(1-p)^{n-k}$ & $np$ & $np(1-p)$ & $n$次独立重复试验成功次数 \\
泊松分布 & $P(\lambda)$ & $\frac{\lambda^k}{k!}e^{-\lambda}$ & $\lambda$ & $\lambda$ & 稀有事件计数 \\
几何分布 & $G(p)$ & $p(1-p)^{k-1}$ & $\frac{1}{p}$ & $\frac{1-p}{p^2}$ & 首次成功的试验次数 \\
超几何分布 & $H(n,M,N)$ & $\frac{\binom{M}{k}\binom{N-M}{n-k}}{\binom{N}{n}}$ & $\frac{nM}{N}$ & $\frac{nM(N-M)(N-n)}{N^2(N-1)}$ & 不放回抽样 \\
\midrule
\multicolumn{6}{c}{\textbf{连续型}} \\
\midrule
均匀分布 & $U(a,b)$ & $\frac{1}{b-a},\ x\in(a,b)$ & $\frac{a+b}{2}$ & $\frac{(b-a)^2}{12}$ & 等可能取值 \\
指数分布 & $E(\lambda)$ & $\lambda e^{-\lambda x},\ x>0$ & $\frac{1}{\lambda}$ & $\frac{1}{\lambda^2}$ & 寿命/等待时间（无记忆性） \\
正态分布 & $N(\mu,\sigma^2)$ & $\frac{1}{\sqrt{2\pi}\sigma}e^{-\frac{(x-\mu)^2}{2\sigma^2}}$ & $\mu$ & $\sigma^2$ & 误差/自然现象 \\
$\chi^2$分布 & $\chi^2(n)$ & — & $n$ & $2n$ & $n$个标准正态平方和 \\
$t$分布 & $t(n)$ & — & $0\ (n>1)$ & $\frac{n}{n-2}\ (n>2)$ & 均值检验（$\sigma$未知） \\
$F$分布 & $F(n_1,n_2)$ & — & $\frac{n_2}{n_2-2}\ (n_2>2)$ & 复杂 & 方差比检验 \\
\bottomrule
\end{tabular}
}
\end{table}

\begin{note}
\textbf{记忆技巧：}
\begin{itemize}
    \item 泊松分布"期望=方差=$\lambda$"——只有一个参数，两个特征值相等。
    \item 指数分布"期望=标准差=$1/\lambda$"——也是一个参数决定一切。
    \item 二项分布 $B(n,p)$：$EX=np$（$n$次试验每次成功概率$p$），$DX=np(1-p)$（乘以失败概率）。
    \item $\chi^2(n)$：期望=自由度，方差=2倍自由度。
\end{itemize}
\end{note}

\begin{note}
\textbf{【高频题型模板】求 $Y=g(X)$ 的密度（CDF 法万能模板）：}

\textbf{适用场景：} 已知连续型 $X$ 的密度 $f_X(x)$，求 $Y=g(X)$ 的密度。当 $g$ 不单调或分段时，公式法不好用，CDF 法永远可行。

\textbf{机械化三步：}
\begin{enumerate}
    \item \textbf{写 CDF：} $F_Y(y)=P\{g(X)\leq y\}$，把 $\{g(X)\leq y\}$ 化成 $X$ 的事件（解不等式）；
    \item \textbf{用已知密度积分：} $F_Y(y)=\displaystyle\int_{\{x:\,g(x)\leq y\}} f_X(x)\,\mathrm{d}x$；
    \item \textbf{对 $y$ 求导：} $f_Y(y)=F_Y'(y)$。
\end{enumerate}
\textbf{易错提醒：} 第1步解不等式时要\textbf{分 $y$ 的范围讨论}（如 $y>0$ 和 $y\leq 0$），不要漏掉分段。
\end{note}


\chapter{二维随机变量及分布}

\begin{note}
\textbf{自学导读：从一维到二维的跨越。}

在第 1--2 章中，我们只用一个随机变量 $X$ 来描述试验结果。但现实中的问题往往涉及\textbf{两个或多个}随机量之间的关系：
\begin{itemize}
    \item 身高和体重是否"联动"？知道一个人身高 180cm，能否推断他体重大概率超过 70kg？
    \item 两门课的成绩之间有没有关联？数学好的人概率论是否也好？
    \item 两个零件的尺寸误差是否独立？如果不独立，装配后的总误差怎么算？
\end{itemize}
这些问题仅靠 $X$ 的分布和 $Y$ 的分布\textbf{分别}回答不了——你必须知道 $X$ 和 $Y$ 的\textbf{联合}行为。这就是本章的核心任务。

\textbf{学习路线图：}
\begin{enumerate}
    \item 先学联合分布（把两个变量"绑在一起"看）；
    \item 再学边缘分布（从联合分布中"提取"单个变量的信息——你会发现这就是第 2 章学过的东西）；
    \item 然后学条件分布（已知一个变量取某值时，另一个变量怎么分布——这和第 1 章的条件概率一脉相承）；
    \item 最后学独立性（什么时候联合分布可以"拆开"成两个边缘分布的乘积）。
\end{enumerate}
整章的逻辑是：\textbf{联合 $\to$ 边缘 $\to$ 条件 $\to$ 独立}。每一步都在回答"两个随机变量之间的关系能告诉我们什么"。
\end{note}

\[
\begin{cases}
\text{二维随机变量的概念：二维平面中随机点的坐标} \\
\text{二维随机变量的概率分布定义及性质}
\begin{cases}
(1) \text{ 单调不减性} \\
(2) \text{ 规范性} \\
(3) \text{ 右连续性}
\end{cases}
\text{ 及 }
\begin{cases}
\text{边缘分布} \\
\text{条件分布}
\end{cases} \\[6pt]

\text{二维离散型}
\begin{cases}
\text{联合分布律：} P\{X=x_i,Y=y_j\}=p_{ij};\ p_{ij}\geqslant0,\ \sum\limits_{i=1}^{\infty}\sum\limits_{j=1}^{\infty}p_{ij}=1 \\[8pt]
\text{边缘分布律：}
\begin{cases}
P\{X=x_i\}=\sum\limits_{j=1}^{\infty}p_{ij}=p_{i\cdot} \\[8pt]
P\{Y=y_j\}=\sum\limits_{i=1}^{\infty}p_{ij}=p_{\cdot j}
\end{cases} \\[12pt]
\text{条件分布律：}
\begin{cases}
P\{X=x_i|Y=y_j\}=\dfrac{p_{ij}}{p_{\cdot j}} \\[8pt]
P\{Y=y_j|X=x_i\}=\dfrac{p_{ij}}{p_{i\cdot}}
\end{cases}
\end{cases} \\[18pt]

\text{二维连续型}
\begin{cases}
\text{联合概率密度：} f(x,y)\geqslant0,\ \displaystyle\int_{-\infty}^{+\infty}\int_{-\infty}^{+\infty}f(u,v)dudv=1 \\[12pt]
\text{边缘分布：}
\begin{cases}
f_X(x)=\displaystyle\int_{-\infty}^{+\infty}f(x,y)dy,\ F_X(x)=\displaystyle\int_{-\infty}^{x}f_X(t)dt \\[12pt]
f_Y(y)=\displaystyle\int_{-\infty}^{+\infty}f(x,y)dx,\ F_Y(y)=\displaystyle\int_{-\infty}^{y}f_Y(s)ds
\end{cases} \\[18pt]
\text{条件概率密度：}
\begin{cases}
f_{X|Y}(x|y)=\dfrac{f(x,y)}{f_Y(y)} \\[8pt]
f_{Y|X}(y|x)=\dfrac{f(x,y)}{f_X(x)}
\end{cases}
\end{cases} \\[18pt]

\text{随机变量的独立性：}
\begin{cases}
\text{离散型：} p_{ij}=p_{i\cdot}p_{\cdot j} \\
\text{连续型：} f(x,y)=f_X(x)\cdot f_Y(y)
\end{cases}
\iff F(x,y)=F_X(x)\cdot F_Y(y) \\[12pt]

\text{二维随机变量函数的概率分布}
\begin{cases}
\text{离散型：逐点法或表格法} \\
\text{连续型：卷积公式法，分布函数法}
\end{cases}
\end{cases}
\]

\begin{introduction}
  \item 二维随机变量与联合分布函数~\ref{def:2d-random-variable}~\ref{def:joint-df}~\ref{prop:joint-df}
  \item 边缘分布函数~\ref{def:marginal-df}
  \item 二维离散型随机变量：联合分布律~\ref{def:2d-discrete}、边缘分布律~\ref{def:marginal-pm}、条件分布律~\ref{def:conditional-pm}
  \item 二维连续型随机变量：联合、边缘与条件概率密度~\ref{def:2d-continuous}
  \item 常见二维分布：均匀分布~\ref{def:2d-uniform}、正态分布~\ref{def:2d-normal}~\ref{prop:2d-normal}
  \item $n$ 维随机变量~\ref{def:n-d-random-variable}~\ref{def:n-joint-df}
  \item 随机变量的独立性：定义~\ref{def:2d-independence}~\ref{def:n-independence}、充要条件与性质~\ref{prop:independence-implications}
  \item 二维随机变量函数的分布：分布函数法、卷积公式~\ref{thm:convolution}、$\max$ 和 $\min$ 的分布、常见分布的可加性
\end{introduction}

\section{二维随机变量}

在许多实际问题中，随机试验的结果往往需要用两个或更多个随机变量来描述。例如，研究某地区的气象状况，需要同时考虑温度和湿度；检查产品的质量，需要同时考察长度和直径等。这时，我们不仅要研究每个随机变量的统计规律，还要研究它们之间的相互关系。

\begin{definition}[二维随机变量] \label{def:2d-random-variable}
设 \(E\) 是一个随机试验，其样本空间为 \(\Omega\)。若 \(X=X(\omega)\) 和 \(Y=Y(\omega)\) 是定义在同一样本空间 \(\Omega\) 上的两个随机变量，则由它们构成的向量 \((X,Y)\) 称为\textbf{二维随机向量}（或\textbf{二维随机变量}）。
\end{definition}

\begin{note}
可以把 \((X,Y)\) 看成平面上的一个随机点：一次试验先产生样本点 \(\omega\)，再映射成坐标 \((X(\omega),Y(\omega))\)。
研究二维随机变量时，除了关心 \(X\)、\(Y\) 各自怎么分布，还必须关心它们之间是否“联动”。例如，单独知道身高和体重各自的分布，并不能回答“高的人是否通常更重”这一问题；这正是联合分布要描述的内容。
\end{note}

\begin{note}
\textbf{自学追问：} 如果只知道 \(X\) 和 \(Y\) 的边缘分布，为什么还不够？因为二维随机变量真正要描述的是“联动关系”，而联动关系只有联合分布才能完整刻画。后面凡是问“能否独立”“能否求条件分布”“能否判断函数分布”，第一步都应先回到联合分布。
\end{note}

类似地，如果 $X_1,X_2,\dots,X_n$ 是定义在同一个样本空间 $\Omega$ 上的 $n$ 个随机变量，则称 $(X_1,X_2,\dots,X_n)$ 为\textbf{$n$维随机变量}或\textbf{$n$维随机向量}，$X_i$ ($i=1,2,\dots,n$) 称为第 $i$ 个\textbf{分量}。

\subsection{联合分布函数}

\begin{definition}[联合分布函数] \label{def:joint-df}
设 $(X,Y)$ 是二维随机变量，对于任意实数 $x,y$，称二元函数
\[
F(x,y) = P\{X \leqslant x, Y \leqslant y\}
\]
为二维随机变量 $(X,Y)$ 的\textbf{分布函数}，或称为随机变量 $X$ 和 $Y$ 的\textbf{联合分布函数}，记为 $(X,Y) \sim F(x,y)$。
\end{definition}

\begin{note}
几何上，$F(x_0,y_0)$ 表示随机点 $(X,Y)$ 落在“左下无穷矩形”
\[
\{(x,y)\mid x\le x_0,\ y\le y_0\}
\]
中的概率。要特别注意：它表示的是\textbf{一个区域}的概率，而不是点 $(x_0,y_0)$ 本身的概率；对连续型随机变量来说，单点概率通常为 $0$。
\end{note}

\begin{property}\label{prop:joint-df}
联合分布函数 $F(x,y)$ 具有以下基本性质：
\begin{enumerate}
    \item \textbf{单调性}：$F(x,y)$ 关于每个变量单调不减。即对任意固定的 $y$，当 $x_1 < x_2$ 时，$F(x_1,y) \leqslant F(x_2,y)$；对任意固定的 $x$，当 $y_1 < y_2$ 时，$F(x,y_1) \leqslant F(x,y_2)$。
    \item \textbf{有界性}：$0 \leqslant F(x,y) \leqslant 1$，且
    \[
    F(-\infty, y) = 0, \quad F(x, -\infty) = 0, \quad F(-\infty, -\infty) = 0, \quad F(+\infty, +\infty) = 1.
    \]
    \item \textbf{右连续性}：$F(x,y)$ 关于每个变量右连续，即
    \[
    F(x+0, y) = F(x,y), \quad F(x, y+0) = F(x,y).
    \]
    \item \textbf{非负性}：对于任意的 $x_1 < x_2$，$y_1 < y_2$，有
    \[
    P\{x_1 < X \leqslant x_2,\ y_1 < Y \leqslant y_2\} = F(x_2,y_2) - F(x_2,y_1) - F(x_1,y_2) + F(x_1,y_1) \geqslant 0.
    \]
\end{enumerate}
\end{property}

\begin{remark}
性质 (4) 是二维分布函数特有的性质，也是区别于一维情形的关键。仅满足前三个性质的函数不一定能成为某个二维随机变量的联合分布函数，必须同时满足性质 (4)。例如，设
\[
F(x,y) = \begin{cases} 1, & x+y \geqslant 0, \\ 0, & x+y < 0, \end{cases}
\]
满足性质 (1)--(3)，但不满足性质 (4)，因此不是联合分布函数。
\end{remark}

\begin{definition}[边缘分布函数] \label{def:marginal-df}
设二维随机变量 $(X,Y)$ 的联合分布函数为 $F(x,y)$，分别称
\begin{align*}
F_X(x) = P\{X \leqslant x\} = F(x, +\infty) = \lim_{y \to +\infty} F(x,y)\\
 F_Y(y) = P\{Y \leqslant y\} = F(+\infty, y) = \lim_{x \to +\infty} F(x,y)
\end{align*}
为 $(X,Y)$ 关于 $X$ 和关于 $Y$ 的\textbf{边缘分布函数}。
\end{definition}

\begin{note}
边缘分布的含义是：只看其中一个分量，把另一个分量“放开不管”。计算时体现为让另一个变量趋于 $+\infty$。但要注意，\textbf{知道两个边缘分布通常还不够}，一般不能反推出联合分布；除非额外知道 $X$ 与 $Y$ 独立，或者知道更多结构信息。
\end{note}


\begin{exercise}\label{exer:3-1}
设 $(X,Y)$ 的联合分布函数为 $F(x,y) = a\left(b + \arctan\dfrac{x}{2}\right)\left(c + \arctan\dfrac{y}{2}\right)$，求常数 $a, b, c$ 及 $(X,Y)$ 的概率密度。
\end{exercise}
\begin{solution}
由联合分布函数的性质：\\
(1) $F(-\infty, y) = a\left(b - \dfrac{\pi}{2}\right)\left(c + \arctan\dfrac{y}{2}\right) = 0$，由于此式对所有 $y$ 成立，故 $b = \dfrac{\pi}{2}$。\\
(2) $F(x, -\infty) = a\left(b + \arctan\dfrac{x}{2}\right)\left(c - \dfrac{\pi}{2}\right) = 0$，同理得 $c = \dfrac{\pi}{2}$。\\
(3) $F(+\infty, +\infty) = a\left(\dfrac{\pi}{2} + \dfrac{\pi}{2}\right)\left(\dfrac{\pi}{2} + \dfrac{\pi}{2}\right) = a\pi^2 = 1$，故 $a = \dfrac{1}{\pi^2}$。\\
概率密度为联合分布函数的二阶混合偏导数：
\[
f(x,y) = \frac{\partial^2 F(x,y)}{\partial x\,\partial y} = \frac{1}{\pi^2} \cdot \frac{1/2}{1 + (x/2)^2} \cdot \frac{1/2}{1 + (y/2)^2} = \frac{4}{\pi^2(4+x^2)(4+y^2)}.
\]
\end{solution}


\begin{exercise}[密度函数线性组合] \label{exer:3-c1}
设二维随机变量 $(X_1,Y_1)$ 和 $(X_2,Y_2)$ 的概率密度分别为 $f_1(x,y)$ 与 $f_2(x,y)$，令 $f(x,y) = af_1(x,y) + bf_2(x,y)$。若 $f(x,y)$ 是某二维连续型随机变量的概率密度，则 $a, b$ 满足条件（\quad）。
\begin{tabular}{ll}
(A) $a+b=1$ & (B) $a>0$ 且 $b>0$ \\
(C) $0\leqslant a\leqslant 1,\ 0\leqslant b\leqslant 1$ & (D) $a\geqslant 0,\ b\geqslant 0$，且 $a+b=1$
\end{tabular}
\end{exercise}
\begin{solution}
\textbf{题型识别：} 两个概率密度的线性组合何时仍是概率密度。

\textbf{依据：} $f(x,y)$ 为概率密度的充要条件是 (i) $f(x,y) \geqslant 0$，(ii) $\displaystyle\iint f(x,y)\,\mathrm{d}x\mathrm{d}y = 1$。

\textbf{逐条检验：}

条件 (ii)：由于 $f_1, f_2$ 各自归一，
\[
\iint f\,\mathrm{d}x\mathrm{d}y = a\underbrace{\iint f_1}_{=1} + b\underbrace{\iint f_2}_{=1} = a + b = 1.
\]

条件 (i)：$f_1, f_2 \geqslant 0$，要使 $af_1 + bf_2 \geqslant 0$ 对一切 $(x,y)$ 成立，需要 $a \geqslant 0$ 且 $b \geqslant 0$。

两个条件合起来就是 $a \geqslant 0$，$b \geqslant 0$，$a + b = 1$。答案选 \textbf{D}。
\end{solution}

\begin{exercise}[标准正态独立随机变量] \label{exer:3-c2}
设随机变量 $X, Y$ 相互独立，且均服从标准正态分布 $N(0,1)$，则（\quad）。
\begin{tabular}{ll}
(A) $P\{X+Y \geqslant 1\} = 1/2$ & (B) $P\{X-Y \geqslant 1\} = 1/2$ \\
(C) $P\{\max(X,Y) \geqslant 0\} = 1/4$ & (D) $P\{\min(X,Y) \geqslant 0\} = 1/4$
\end{tabular}
\end{exercise}
\begin{solution}
应选 D。

逐项判断最稳妥。

先看 A。由于 \(X,Y\) 独立且都服从 \(N(0,1)\)，所以 \(X+Y\sim N(0,2)\)。
它关于 0 对称，但题目问的是 \(P\{X+Y\ge 1\}\)，阈值是 1 而不是对称中心 0，因此这个概率不可能等于 \(\dfrac12\)。所以 A 错。

同理，\(X-Y\sim N(0,2)\)，于是 \(P\{X-Y\ge 1\}\ne \frac12\)，故 B 也错。

再看 C。事件 \(\{\max(X,Y)\ge 0\}\) 的对立事件是“两个都小于 0”，由独立性和标准正态的对称性，
\[
P\{\max(X,Y)\ge 0\}=1-P\{X<0,Y<0\}=1-\frac12\cdot\frac12=\frac34,
\]
并不是 \(\dfrac14\)，故 C 错。

最后看 D。事件 \(\{\min(X,Y)\ge 0\}\) 等价于“两个都不小于 0”，因此
\[
P\{\min(X,Y)\ge 0\}=P\{X\ge 0\}P\{Y\ge 0\}
=\frac12\cdot\frac12
=\frac14.
\]
所以 D 正确。
\end{solution}



\begin{exercise}[max 的分布函数] \label{exer:3-c3}
设随机变量 $X$ 和 $Y$ 独立同分布，且 $X$ 的分布函数为 $F(x)$，则 $Z = \max\{X,Y\}$ 的分布函数为（\quad）。
\begin{tabular}{ll}
(A) $F^2(z)$ & (B) $F(x)F(y)$ \\
(C) $1-[1-F(z)]^2$ & (D) $[1-F(x)][1-F(y)]$
\end{tabular}
\end{exercise}
\begin{solution}
求最大值的分布函数时，最关键的是先把事件 \(\{\max(X,Y)\le z\}\) 翻译成关于 \(X,Y\) 的同时事件。
\[
F_Z(z) = P\{\max(X,Y) \leqslant z\} = P\{X \leqslant z,\ Y \leqslant z\} = F(z) \cdot F(z) = F^2(z).
\]

上面第二步用到了独立同分布条件 \(P\{X\le z,\ Y\le z\}=P\{X\le z\}P\{Y\le z\}=F(z)\cdot F(z)\)，故应选 A。

顺便提醒：选项 C 的 \(1-[1-F(z)]^2\) 对应的是事件“至少有一个不大于 \(z\)”的形式，更常在处理最小值时出现，不能和最大值公式混淆。
\end{solution}



\begin{exercise}[边缘分布函数] \label{exer:3-c4}
如果二维随机变量 $(X,Y)$ 的分布函数为
\[
F(x,y) = \begin{cases} (1-e^{-\lambda_1 x})(1-e^{-\lambda_2 y} + e^{-\lambda_3 x-\lambda_2 y}), & x > 0,\ y > 0, \\ 0, & \text{其他}, \end{cases}
\]
其中 $\lambda_1, \lambda_2, \lambda_3 > 0$，求 $X$ 和 $Y$ 的边缘分布函数。
\end{exercise}
\begin{solution}
\[
F_X(x) = \lim_{y\to+\infty}F(x,y) = \begin{cases} (1-e^{-\lambda_1 x})\cdot 1 = 1-e^{-\lambda_1 x}, & x > 0, \\ 0, & \text{其他}. \end{cases}
\]
\[
F_Y(y) = \lim_{x\to+\infty}F(x,y) = \begin{cases} 1\cdot(1-e^{-\lambda_2 y}) = 1-e^{-\lambda_2 y}, & y > 0, \\ 0, & \text{其他}. \end{cases}
\]
\begin{note}
$X \sim E(\lambda_1)$，$Y \sim E(\lambda_2)$。注意 $F(x,y) \neq F_X(x)F_Y(y)$，故 $X$ 与 $Y$ 不独立（除非 $\lambda_3 = 0$）。
\end{note}
\end{solution}



\begin{exercise}[边缘密度不保证联合为正态] \label{exer:3-c5}
随机变量 $(X,Y)$ 的概率密度为
\[
f(x,y) = \frac{1+\sin x\sin y}{2\pi}\,e^{-\frac{x^2+y^2}{2}}, \quad -\infty < x, y < +\infty,
\]
求边缘概率密度 $f_X(x)$ 和 $f_Y(y)$。
\end{exercise}
\begin{solution}
$\sin y \cdot e^{-y^2/2}$ 是奇函数，在 $\mathbb{R}$ 上积分为 $0$，故
\[
f_X(x) = \int_{-\infty}^{+\infty}\frac{1+\sin x\sin y}{2\pi}\,e^{-\frac{x^2+y^2}{2}}\,\mathrm{d}y = \frac{e^{-x^2/2}}{2\pi}\left[\int_{-\infty}^{+\infty}e^{-y^2/2}\,\mathrm{d}y + \sin x \cdot 0\right] = \frac{1}{\sqrt{2\pi}}\,e^{-x^2/2}.
\]
同理 $f_Y(y) = \dfrac{1}{\sqrt{2\pi}}\,e^{-y^2/2}$。
\begin{note}
$X \sim N(0,1)$，$Y \sim N(0,1)$，但 $(X,Y)$ \textbf{不服从}二维正态分布（密度中含 $\sin x\sin y$ 项）。这说明：\textbf{边缘分布为正态不能推出联合分布为二维正态}。
\end{note}
\end{solution}


\section{n维随机变量}

前面主要讨论了二维随机变量，本节将相关概念推广到 $n$ 维情形。在实际问题中，有时需要同时研究多个随机变量之间的关系。

\begin{definition}[$n$维随机变量] \label{def:n-d-random-variable}
如果 $X_1, X_2, \dots, X_n$ 是定义在同一个样本空间 $\Omega$ 上的 $n$ 个随机变量，则称 $(X_1, X_2, \dots, X_n)$ 为\textbf{$n$维随机变量}（或\textbf{$n$维随机向量}），$X_i$（$i=1,2,\dots,n$）称为第 $i$ 个\textbf{分量}。
\end{definition}

\begin{definition}[$n$维联合分布函数] \label{def:n-joint-df}
对任意的 $n$ 个实数 $x_1, x_2, \dots, x_n$，称 $n$ 元函数
\[
F(x_1, x_2, \dots, x_n) = P\{X_1 \leqslant x_1, X_2 \leqslant x_2, \dots, X_n \leqslant x_n\}
\]
为 $n$ 维随机变量 $(X_1, X_2, \dots, X_n)$ 的\textbf{联合分布函数}（或\textbf{分布函数}）。
\end{definition}

\begin{note}
把二维情形推广到 $n$ 维时，思想并没有变化：联合分布函数仍然是在问“所有分量同时落在给定阈值以下的概率是多少”。只是维数增大后，记号更长、几何图形不再直观，但定义和用法与二维完全平行。
\end{note}

\begin{property}[$n$维联合分布函数的性质]\label{prop:n-joint-df-property}
$n$ 维联合分布函数 $F(x_1, x_2, \dots, x_n)$ 具有以下性质：
\begin{enumerate}
    \item \textbf{单调性}：$F$ 关于每个变量 $x_i$ 单调不减；
    \item \textbf{右连续性}：$F$ 关于每个变量 $x_i$ 右连续；
    \item \textbf{有界性}：
    \[
    F(\underbrace{-\infty, \dots, -\infty}_{k\text{个}}, x_{k+1}, \dots, x_n) = 0 \quad (k=1,2,\dots,n-1), \quad F(+\infty, +\infty, \dots, +\infty) = 1;
    \]
    \item \textbf{非负性}：对任意的 $a_i < b_i$（$i=1,2,\dots,n$），以 $F$ 计算的 $n$ 维矩形区域的概率
    \[
    P\{a_1 < X_1 \leqslant b_1, \dots, a_n < X_n \leqslant b_n\} \geqslant 0.
    \]
\end{enumerate}
\end{property}

\begin{definition}[$n$维边缘分布函数]\label{def:n-marginal-df}
$n$ 维随机变量 $(X_1, X_2, \dots, X_n)$ 中任意 $k$ 个（$1 \leqslant k < n$）分量所构成的 $k$ 维随机变量的分布函数，称为原 $n$ 维随机变量的\textbf{$k$维边缘分布函数}。特别地，关于第 $i$ 个分量 $X_i$ 的边缘分布函数为
\[
F_{X_i}(x_i) = F(+\infty, \dots, +\infty, x_i, +\infty, \dots, +\infty).
\]
\end{definition}

\begin{note}
$n$ 维边缘分布就是“只保留感兴趣的那几个分量，其余分量全部积分掉或放到 $+\infty$”。因此边缘分布反映的是局部信息，而联合分布保存的是整体联动信息。
\end{note}

类似地，对于 $n$ 维离散型随机变量，可以定义\textbf{联合分布律} \[p_{i_1 i_2 \cdots i_n} = P\{X_1 = x_{i_1}, X_2 = x_{i_2}, \dots, X_n = x_{i_n}\}\] 及相应的边缘分布律和条件分布律。

对于 $n$ 维连续型随机变量，可以定义\textbf{联合概率密度} $f(x_1, x_2, \dots, x_n)$，满足
\[
F(x_1, x_2, \dots, x_n) = \int_{-\infty}^{x_n}\cdots\int_{-\infty}^{x_1}f(t_1, t_2, \dots, t_n)\,\mathrm{d}t_1\cdots\mathrm{d}t_n,
\]
以及边缘概率密度、条件概率密度和概率密度乘法公式等，这些概念与二维情形完全类似。

\section{随机变量的独立性}

\begin{note}
\textbf{自学导读：为什么独立性如此重要？}

回忆第 1 章：事件 $A$、$B$ 独立意味着 $P(AB)=P(A)P(B)$。现在把这个想法从"事件"升级到"随机变量"。

独立性的实际意义：如果 $X$ 和 $Y$ 独立，那么知道 $X$ 的取值\textbf{不会改变}你对 $Y$ 的任何判断。这意味着：
\begin{itemize}
    \item 联合分布可以直接拆成边缘分布的乘积——计算量大幅降低；
    \item $E(XY)=E(X)E(Y)$，$D(X+Y)=D(X)+D(Y)$——后面第 4 章会反复用到；
    \item 卷积公式、$\max/\min$ 分布等工具都以独立性为前提。
\end{itemize}
所以本节不是一个孤立的定义，而是后续所有章节的\textbf{基础假设}。考试中凡是看到"$X,Y$ 相互独立"这个条件，就意味着你可以把联合问题拆成两个一维问题分别处理。
\end{note}

在第一章中我们讨论了事件的独立性，现在将这一概念推广到随机变量的情形。独立性是概率论中最重要的概念之一，它能够大大简化联合分布的研究。

\subsection{独立性的定义}

\begin{definition}[两个随机变量的独立性] \label{def:2d-independence}
设 $F(x,y)$ 及 $F_X(x)$、$F_Y(y)$ 分别是二维随机变量 $(X,Y)$ 的联合分布函数及边缘分布函数。若对任意实数 $x, y$，都有
\[
P\{X \leqslant x, Y \leqslant y\} = P\{X \leqslant x\} \cdot P\{Y \leqslant y\},
\]
即
\[
F(x,y) = F_X(x) \cdot F_Y(y),
\]
则称随机变量 $X$ 和 $Y$ 是\textbf{相互独立}的，否则称 $X$ 与 $Y$ 不相互独立。
\end{definition}

\begin{note}
独立性的意思不是“某一个特殊事件能拆开”，而是\textbf{对任意阈值 $x,y$，相应矩形事件都能拆成乘积}。因此：
\begin{enumerate}
    \item 证明独立时，通常要证明联合分布函数、联合分布律或联合密度能够整体分解；
    \item 证明不独立时，只需找到一个反例，使 $F(x,y)\neq F_X(x)F_Y(y)$ 即可。
\end{enumerate}
\end{note}

将独立性推广到 $n$ 个随机变量的情形：

\begin{definition}[$n$ 个随机变量的相互独立性] \label{def:n-independence}
设 $n$ 维随机变量 $(X_1, X_2, \dots, X_n)$ 的联合分布函数为 $F(x_1, x_2, \dots, x_n)$，$X_i$ 的边缘分布函数为 $F_{X_i}(x_i)$。若对任意实数 $x_1, x_2, \dots, x_n$，都有
\[
F(x_1, x_2, \dots, x_n) = \prod_{i=1}^{n}F_{X_i}(x_i),
\]
则称 $X_1, X_2, \dots, X_n$ 是\textbf{相互独立}的。

进一步地，两个多维随机变量 $(X_1, \dots, X_n)$ 与 $(Y_1, \dots, Y_m)$ 相互独立，是指对任意实数 $x_i$（$i=1,\dots,n$）与 $y_j$（$j=1,\dots,m$），都有
\[
F(x_1, \dots, x_n, y_1, \dots, y_m) = F_X(x_1, \dots, x_n) \cdot F_Y(y_1, \dots, y_m),
\]
其中 $F_X$ 和 $F_Y$ 分别是 $(X_1, \dots, X_n)$ 和 $(Y_1, \dots, Y_m)$ 的分布函数。
\end{definition}

\begin{note}
这里的“相互独立”比“两两独立”强得多。两两独立只要求任意两个变量独立；而相互独立要求任意有限个变量联合起来都能分解成边缘分布的乘积。以后做选择题时，这两个概念非常容易混淆。
\end{note}

\begin{definition}[随机变量序列的独立性]\label{def:sequence-independence} 
若对任意正整数 $k$ 和任意 $1 \leqslant i_1 < i_2 < \dots < i_k$，随机变量 $X_{i_1}, X_{i_2}, \dots, X_{i_k}$ 都相互独立，则称随机变量序列 $X_1, X_2, \dots$ 是\textbf{相互独立}的。
\end{definition}

\subsection{独立的充要条件}

$n$ 个随机变量 $X_1, X_2, \dots, X_n$ 相互独立
$\iff$ 对任意的 $n$ 个实数 $x_i$（$i=1,2,\dots,n$），$n$ 个事件 $\{X_1 \leqslant x_1\}, \{X_2 \leqslant x_2\}, \dots, \{X_n \leqslant x_n\}$ 相互独立。

\subsubsection{离散型随机变量的独立性}

\begin{theorem}[离散型随机变量独立的判定]\label{thm:discrete-independence}
设 $(X,Y)$ 为二维离散型随机变量，其联合分布律为 $P\{X = x_i, Y = y_j\} = p_{ij}$，则 $X$ 与 $Y$ 相互独立的\textbf{充分必要条件}是对一切 $i, j$，都有
\[
p_{ij} = p_{i\cdot} \cdot p_{\cdot j}.
\]
更一般地，$n$ 个离散型随机变量 $X_1, X_2, \dots, X_n$ 相互独立的充分必要条件是对任意的 $x_i \in D_i$（$X_i$ 的一切可能取值），有
\[
P\{X_1 = x_1, \dots, X_n = x_n\} = \prod_{i=1}^{n}P\{X_i = x_i\}.
\]
\end{theorem}

\begin{proof}
必要性显然。充分性：若 $p_{ij} = p_{i\cdot} \cdot p_{\cdot j}$ 对一切 $i,j$ 成立，则
\[
F(x,y) = \sum_{x_i \leqslant x}\sum_{y_j \leqslant y}p_{ij} = \sum_{x_i \leqslant x}\sum_{y_j \leqslant y}p_{i\cdot}\,p_{\cdot j} = \left(\sum_{x_i \leqslant x}p_{i\cdot}\right)\left(\sum_{y_j \leqslant y}p_{\cdot j}\right) = F_X(x)\,F_Y(y).
\]
\end{proof}

\subsubsection{连续型随机变量的独立性}

\begin{theorem}[连续型随机变量独立的判定]\label{thm:continuous-independence}
设 $(X,Y)$ 为二维连续型随机变量，联合概率密度为 $f(x,y)$，边缘概率密度为 $f_X(x)$ 和 $f_Y(y)$，则 $X$ 与 $Y$ 相互独立的\textbf{充分必要条件}是
\[
f(x,y) = f_X(x) \cdot f_Y(y)
\]
几乎处处成立（即除了面积为零的集合外处处成立）。

更一般地，$n$ 维连续型随机变量 $(X_1, \dots, X_n)$ 相互独立的充分必要条件是
\[
f(x_1, x_2, \dots, x_n) = \prod_{i=1}^{n}f_i(x_i),
\]
其中 $f_i(x_i)$ 为 $X_i$ 的边缘概率密度。
\end{theorem}

\begin{proof}
充分性：若 $f(x,y) = f_X(x)\,f_Y(y)$，则
\[
F(x,y) = \int_{-\infty}^{y}\int_{-\infty}^{x}f(u,v)\,\mathrm{d}u\,\mathrm{d}v = \int_{-\infty}^{x}f_X(u)\,\mathrm{d}u \cdot \int_{-\infty}^{y}f_Y(v)\,\mathrm{d}v = F_X(x)\,F_Y(y).
\]

必要性：若 $F(x,y) = F_X(x)\,F_Y(y)$，求二阶混合偏导数即得 $f(x,y) = f_X(x)\,f_Y(y)$。
\end{proof}

\begin{note}
利用独立性，可以大大简化联合分布的确定和概率计算。例如，当 $X$ 与 $Y$ 独立时，
\[
P\{a < X \leqslant b,\ c < Y \leqslant d\} = P\{a < X \leqslant b\} \cdot P\{c < Y \leqslant d\}.
\]
此外，常数与任何随机变量相互独立。
\end{note}

\subsection{独立性的性质}

\begin{property}\label{prop:independence-implications}
随机变量的独立性具有以下重要性质：
\begin{enumerate}
    \item 若 $X_1, X_2, \dots, X_n$ 相互独立，则其中任意 $k$ 个（$2 \leqslant k \leqslant n$）随机变量也相互独立。
    \item 设 $(X,Y)$ 为二维离散型随机变量，$X$ 与 $Y$ 独立，则条件分布等于边缘分布：
    \begin{align*}
    P\{X = x_i \mid Y = y_j\} = P\{X = x_i\}\ (P\{Y = y_j\} > 0)\\
     P\{Y = y_j \mid X = x_i\} = P\{Y = y_j\}\ (P\{X = x_i\} > 0)
    \end{align*}
    \item 设 $(X,Y)$ 为二维连续型随机变量，$X$ 与 $Y$ 独立，则条件密度等于边缘密度：
    \[
    f_{X|Y}(x|y) = \frac{f(x,y)}{f_Y(y)} = f_X(x)\ (f_Y(y) > 0), \quad f_{Y|X}(y|x) = f_Y(y)\ (f_X(x) > 0).
    \]
    \item 若 $X_1, X_2, \dots, X_n$ 相互独立，$g_1, g_2, \dots, g_n$ 为一元连续函数，则 $g_1(X_1), g_2(X_2), \dots, g_n(X_n)$ 也相互独立。更一般地，若 $X_{11}, \dots, X_{1i_1}, X_{21}, \dots, X_{2i_2}, \dots, X_{n1}, \dots, X_{ni_n}$ 相互独立，$g_k$ 是 $i_k$（$k=1,2,\dots,n$）元连续函数，则 $g_1(X_{11}, \dots, X_{1i_1}), g_2(X_{21}, \dots, X_{2i_2}), \dots, g_n(X_{n1}, \dots, X_{ni_n})$ 也相互独立。
    \item 若 $(X,Y)$ 服从二维正态分布，则 $X$ 与 $Y$ 相互独立当且仅当 $\rho = 0$。但一般而言，$X$ 与 $Y$ 不相关（$\mathrm{Cov}(X,Y) = 0$）\textbf{不能}推出 $X$ 与 $Y$ 独立。
    \item 独立一定不相关，但不相关不一定独立（除非联合分布为正态分布）。
\end{enumerate}
\end{property}

\subsection{不独立的判断与证明}

\begin{property}[不独立的判定]\label{prop:dependence-criterion}
$X$ 与 $Y$ 不独立 $\iff$ 存在 $x_0, y_0$，使得
\[
F(x_0, y_0) \neq F_X(x_0) \cdot F_Y(y_0).
\]
也可通过取特定值使 $P\{X \leqslant x_0, Y \leqslant y_0\} \neq P\{X \leqslant x_0\} \cdot P\{Y \leqslant y_0\}$ 来证明不独立。对于离散型，寻找一对 $(i,j)$ 使 $p_{ij} \neq p_{i\cdot} \cdot p_{\cdot j}$；对于连续型，验证 $f(x,y) \neq f_X(x)\,f_Y(y)$。
\end{property}


\begin{exercise}[独立性选择题] \label{exer:3-12d}
设随机变量 $X, Y$ 相互独立，分布函数分别为 $F_X(x), F_Y(y)$，则对任意实数 $z$，$P\{\max(X,Y) \leqslant z\}$ 等于（\quad）。\\
\begin{tabular}{ll}
A. $F_X(z)F_Y(z)$ & B. $(1-F_X(z))F_Y(z)$ \\
C. $1-(1-F_X(z))(1-F_Y(z))$ & D. $(1-F_X(z))F_Y(z)$
\end{tabular}
\end{exercise}
\begin{solution}
由独立性：
\[
P\{\max(X,Y) \leqslant z\} = P\{X \leqslant z,\ Y \leqslant z\} = P\{X \leqslant z\}\,P\{Y \leqslant z\} = F_X(z)\,F_Y(z).
\]
答案选 \textbf{A}。
\begin{note}
选项 C 为 $P\{\min(X,Y) \leqslant z\}$ 的表达式。不要混淆 max 和 min 的公式。
\end{note}
\end{solution}

\begin{exercise}[证明 $X$ 与 $|X|$ 不独立] \label{exer:3-5c}
设随机变量 $X$ 的概率密度为
\[
f(x) = \frac{1}{2}e^{-|x|}, \quad x \in (-\infty, +\infty),
\]
证明：$X$ 与 $|X|$ 不独立。
\end{exercise}
\begin{solution}
要证明两个随机变量不独立，只需找到一个具体的事件，使“联合概率 $\neq$ 边缘概率之积”即可。这里取最方便的一组事件：
\[
\{X\le 1\},\qquad \{|X|\le 1\}.
\]

先分别计算这两个事件的概率。
\[
\begin{aligned}
P\{X\le 1\}
&=\int_{-\infty}^1 \frac{1}{2}e^{-|x|}\,\mathrm{d}x \\
&=\frac{1}{2}\int_{-\infty}^{0} e^x\,\mathrm{d}x+\frac{1}{2}\int_{0}^{1} e^{-x}\,\mathrm{d}x \\
&=\frac{1}{2}+\frac{1}{2}(1-e^{-1})
=1-\frac{1}{2}e^{-1},
\end{aligned}
\]
\[
\begin{aligned}
P\{|X|\le 1\}
&=\frac{1}{2}\int_{-1}^{1}e^{-|x|}\,\mathrm{d}x
=\int_{0}^{1}e^{-x}\,\mathrm{d}x
=1-e^{-1}.
\end{aligned}
\]

再看联合事件。因为 $\{|X|\le 1\}$ 已经蕴含 $X\le 1$，所以
\[
P\{X\le 1,\ |X|\le 1\}=P\{|X|\le 1\}=1-e^{-1}.
\]

如果 $X$ 与 $|X|$ 独立，就应当有
\[
P\{X\le 1,\ |X|\le 1\}=P\{X\le 1\}\,P\{|X|\le 1\}.
\]
但实际上
\[
\left(1-\frac{1}{2}e^{-1}\right)(1-e^{-1})
=1-\frac{3}{2}e^{-1}+\frac{1}{2}e^{-2}
\neq 1-e^{-1}.
\]
于是
\[
P\{X\le 1,\ |X|\le 1\}\neq P\{X\le 1\}\,P\{|X|\le 1\},
\]
所以 $X$ 与 $|X|$ 不独立。

这个例子也说明：\textbf{一个随机变量与它的非平凡函数通常不会独立}，因为函数值已经携带了原变量的信息。
\end{solution}

\begin{note}
\textbf{自学检查点——独立性：} 学完本节后，你应该能回答以下问题：
\begin{enumerate}
    \item 判断独立性时，离散型看什么？连续型看什么？（答：离散型看是否对所有 $i,j$ 都有 $p_{ij}=p_{i\cdot}\cdot p_{\cdot j}$；连续型看是否 $f(x,y)=f_X(x)\cdot f_Y(y)$）
    \item $X,Y$ 独立能推出不相关吗？不相关能推出独立吗？什么情况下两者等价？（答：独立 $\Rightarrow$ 不相关；反之一般不成立；二维正态时等价）
    \item 证明不独立需要验证所有点吗？（答：不需要，只需找到一个反例）
\end{enumerate}
\end{note}

\section{二维离散型随机变量}

\begin{note}
\textbf{自学导读：} 本节是第 2 章"离散型随机变量"的二维推广。核心操作只有一个：\textbf{列表格}。把 $(X,Y)$ 所有可能的取值对 $(x_i, y_j)$ 及其概率 $p_{ij}$ 写成一张二维表格，就是联合分布律。

你在第 2 章学过的"列出所有取值和概率"的方法，现在从一行变成了一张矩阵——行对应 $X$ 的取值，列对应 $Y$ 的取值，格子里填概率。掌握这张表格的读法，后面的边缘分布（按行/列求和）和条件分布（固定某行/列再归一化）就水到渠成。
\end{note}

\begin{definition}[二维离散型随机变量] \label{def:2d-discrete}
如果二维随机变量 $(X,Y)$ 全部可能取到的值是有限对或可列无限对，则称 $(X,Y)$ 为\textbf{二维离散型随机变量}。
\end{definition}

\subsection{联合分布律}

\begin{definition}[联合分布律]  \label{def:joint-pm}
设二维离散型随机变量 $(X,Y)$ 所有可能取的值为 $(x_i, y_j)$ ($i,j = 1,2,\dots$)，则称
\[
p_{ij} = P\{X = x_i, Y = y_j\}, \quad i,j = 1,2,\dots
\]
为 $(X,Y)$ 的\textbf{联合分布律}，或称为随机变量 $X$ 和 $Y$ 的\textbf{联合概率分布}。
\end{definition}

联合分布律也可用表格形式表示：

\begin{center}
\begin{tabular}{c|cccc|c}
\toprule
$X \backslash Y$ & $y_1$ & $y_2$ & $\cdots$ & $y_j$ & $\cdots$ \\
\midrule
$x_1$ & $p_{11}$ & $p_{12}$ & $\cdots$ & $p_{1j}$ & $\cdots$ \\
$x_2$ & $p_{21}$ & $p_{22}$ & $\cdots$ & $p_{2j}$ & $\cdots$ \\
$\vdots$ & $\vdots$ & $\vdots$ & & $\vdots$ & \\
$x_i$ & $p_{i1}$ & $p_{i2}$ & $\cdots$ & $p_{ij}$ & $\cdots$ \\
$\vdots$ & $\vdots$ & $\vdots$ & & $\vdots$ & \\
\bottomrule
\end{tabular}
\end{center}

\begin{property}\label{prop:joint-pm-condition}
数列 $\{p_{ij}\}$ ($i,j = 1,2,\dots$) 是某一二维离散型随机变量的联合分布律的充分必要条件为：
\begin{enumerate}
    \item \textbf{非负性}：$p_{ij} \geqslant 0$；
    \item \textbf{归一性}：$\displaystyle\sum_{i=1}^{\infty}\sum_{j=1}^{\infty}p_{ij} = 1$。
\end{enumerate}
\end{property}

由联合分布律可以求得联合分布函数：
\[
F(x,y) = \sum_{x_i \leqslant x}\sum_{y_j \leqslant y}p_{ij}.
\]

更一般地，设 $G$ 为平面上的某个区域，则
\[
P\{(X,Y) \in G\} = \sum_{(x_i,y_j) \in G}p_{ij},
\]
即 $(X,Y)$ 落入区域 $G$ 的概率等于所有落在 $G$ 内的可能值 $(x_i, y_j)$ 对应的概率之和。

\subsection{边缘分布律}

\begin{definition}[边缘分布律] \label{def:marginal-pm}
设 $(X,Y)$ 的联合分布律为 $P\{X = x_i, Y = y_j\} = p_{ij}$，则分别称
\[
p_{i\cdot} = P\{X = x_i\} = \sum_{j=1}^{\infty}p_{ij}, \quad i = 1,2,\dots
\]
\[
p_{\cdot j} = P\{Y = y_j\} = \sum_{i=1}^{\infty}p_{ij}, \quad j = 1,2,\dots
\]
为 $(X,Y)$ 关于 $X$ 和关于 $Y$ 的\textbf{边缘分布律}。
\end{definition}

\begin{note}
边缘分布律写在联合分布律表格的最后一列和最后一行。需要注意的是，联合分布律可以完全确定边缘分布律，但反过来由边缘分布律一般\textbf{不能}确定联合分布律，因为边缘分布律不包含两个随机变量之间相互关系的信息。
\end{note}

\begin{note}
\textbf{直觉理解：} 边缘分布就是"忽略另一个变量，只看自己"。操作上：
\begin{itemize}
    \item 求 $X$ 的边缘分布 $\to$ 把联合表格\textbf{按行求和}（把 $Y$ 的所有可能"加在一起"）；
    \item 求 $Y$ 的边缘分布 $\to$ 把联合表格\textbf{按列求和}。
\end{itemize}
这和第 1 章全概率公式的思想一致：$P(X=x_i) = \sum_j P(X=x_i, Y=y_j)$，就是对 $Y$ 的所有取值做"穷举求和"。
\end{note}

\subsection{条件分布律}

\begin{definition}[条件分布律] \label{def:conditional-pm}
设 $(X,Y)$ 为二维离散型随机变量，其联合分布律为 $P\{X = x_i, Y = y_j\} = p_{ij}$ ($i,j = 1,2,\dots$)。

若对固定的 $j$，$p_{\cdot j} = P\{Y = y_j\} > 0$，则称
\[
P\{X = x_i \mid Y = y_j\} = \frac{P\{X = x_i, Y = y_j\}}{P\{Y = y_j\}} = \frac{p_{ij}}{p_{\cdot j}}, \quad i = 1,2,\dots
\]
为在 $\{Y = y_j\}$ 条件下 $X$ 的\textbf{条件分布律}。

同理，若对固定的 $i$，$p_{i\cdot} = P\{X = x_i\} > 0$，则称
\[
P\{Y = y_j \mid X = x_i\} = \frac{p_{ij}}{p_{i\cdot}}, \quad j = 1,2,\dots
\]
为在 $\{X = x_i\}$ 条件下 $Y$ 的\textbf{条件分布律}。
\end{definition}

\begin{note}
\textbf{自学追问：} 条件分布律和第一章的条件概率有什么关系？本质上是同一件事：分母是"已知条件"的概率，分子是"条件与目标同时发生"的概率。区别在于，第一章处理的是事件层面的 $P(B|A)$，这里把它推广到了随机变量的每一个取值上，形成一张完整的"在某条件下的分布表"。

\textbf{易错点：} 条件分布律的分母是边缘概率 $p_{\cdot j}$（或 $p_{i\cdot}$），不是联合概率 $p_{ij}$。另外，条件分布律本身也满足概率公理（非负、归一），可以当作一个"新的分布"来使用——求条件期望、条件方差时都以它为基础。
\end{note}

\begin{exercise}\label{exer:3-2}
将两封信投入 3 个信箱，设 $X_1, X_2$ 分别表示第一个和第二个信箱投进的信的数量。求：
\begin{enumerate}
    \item[(1)] $(X_1, X_2)$ 的联合分布律及边缘分布律；
    \item[(2)] 在条件 $X_2 = 1$ 下，$X_1$ 的条件分布。
\end{enumerate}
\end{exercise}

\begin{solution}
(1) 每封信有 3 种投法，两封信共有 $3^2 = 9$ 种等可能结果。$X_1$ 与 $X_2$ 的可能取值均为 $0, 1, 2$。

\[
\begin{aligned}
p_{00} &= P\{X_1=0, X_2=0\} = \frac{1}{9}, &\quad p_{01} &= P\{X_1=0, X_2=1\} = \frac{2}{9}, \\
p_{02} &= P\{X_1=0, X_2=2\} = \frac{1}{9}, &\quad p_{10} &= P\{X_1=1, X_2=0\} = \frac{2}{9}, \\
p_{11} &= P\{X_1=1, X_2=1\} = \frac{2}{9}, &\quad p_{12} &= P\{X_1=1, X_2=2\} = 0, \\
p_{20} &= P\{X_1=2, X_2=0\} = \frac{1}{9}, &\quad p_{21} &= p_{22} = 0.
\end{aligned}
\]

联合分布律与边缘分布律如下表所示：

\begin{center}
\begin{tabular}{c|ccc|c}
\toprule
$X_1 \backslash X_2$ & $0$ & $1$ & $2$ & $p_{i\cdot}$ \\
\midrule
$0$ & $\dfrac{1}{9}$ & $\dfrac{2}{9}$ & $\dfrac{1}{9}$ & $\dfrac{4}{9}$ \\[6pt]
$1$ & $\dfrac{2}{9}$ & $\dfrac{2}{9}$ & $0$ & $\dfrac{4}{9}$ \\[6pt]
$2$ & $\dfrac{1}{9}$ & $0$ & $0$ & $\dfrac{1}{9}$ \\
\midrule
$p_{\cdot j}$ & $\dfrac{4}{9}$ & $\dfrac{4}{9}$ & $\dfrac{1}{9}$ & $1$ \\
\bottomrule
\end{tabular}
\end{center}

(2) $P\{X_2=1\} = p_{\cdot 1} = \dfrac{4}{9}$。在 $X_2 = 1$ 的条件下，$X_1$ 的条件分布为：
\begin{align*}
P\{X_1=0 \mid X_2=1\} &= \frac{p_{01}}{p_{\cdot 1}} = \frac{2/9}{4/9} = \frac{1}{2}\\
P\{X_1=1 \mid X_2=1\} &= \frac{p_{11}}{p_{\cdot 1}} = \frac{2/9}{4/9} = \frac{1}{2}\\
P\{X_1=2 \mid X_2=1\} &= 0
\end{align*}
即 $X_1 \mid (X_2=1) \sim \begin{pmatrix} 0 & 1 \\ 1/2 & 1/2 \end{pmatrix}$。
\end{solution}

\begin{exercise}\label{exer:3-3}
已知二维随机变量 $(X,Y)$ 的联合概率分布如下表，若随机事件 $\{X=0\}$ 与 $\{X+Y=1\}$ 相互独立，则（\quad）。
\begin{center}
\begin{tabular}{c|cc}
\toprule
$X \backslash Y$ & $0$ & $1$ \\
\midrule
$0$ & $0.4$ & $a$ \\
$1$ & $b$ & $0.1$ \\
\bottomrule
\end{tabular}
\end{center}
\begin{tabular}{llll}
A. $a=0.2, b=0.3$ & B. $a=0.1, b=0.4$ & C. $a=0.3, b=0.2$ & D. $a=0.4, b=0.1$
\end{tabular}
\end{exercise}
\begin{solution}
答案选 \textbf{D}。由联合分布的性质，先由归一性得：
\[
0.4 + a + b + 0.1 = 1 \implies a + b = 0.5.
\]
分别计算三个事件的概率：
\[
P\{X=0\} = 0.4 + a, \quad P\{X+Y=1\} = a + b = 0.5, \quad P\{X=0, X+Y=1\} = P\{X=0, Y=1\} = a.
\]
由 $\{X=0\}$ 与 $\{X+Y=1\}$ 相互独立，有 $P\{X=0\} \cdot P\{X+Y=1\} = P\{X=0, X+Y=1\}$，即
\[
(0.4 + a) \times 0.5 = a \implies 0.2 + 0.5a = a \implies 0.5a = 0.2 \implies a = 0.4.
\]
从而 $b = 0.5 - a = 0.1$。
\end{solution}

\begin{exercise}[联合分布的确定与独立性] \label{exer:3-3b}
已知 $X_1$ 和 $X_2$ 的分布律分别为
\begin{center}
\begin{tabular}{c|ccc}
$X_1$ & $-1$ & $0$ & $1$ \\
\midrule
$P$ & $\dfrac{1}{4}$ & $\dfrac{1}{2}$ & $\dfrac{1}{4}$
\end{tabular}
\quad
\begin{tabular}{c|cc}
$X_2$ & $0$ & $1$ \\
\midrule
$P$ & $\dfrac{1}{2}$ & $\dfrac{1}{2}$
\end{tabular}
\end{center}
且 $P\{X_1 X_2 = 0\} = 1$。求：
\begin{enumerate}
    \item[(1)] $X_1$ 和 $X_2$ 的联合分布律；
    \item[(2)] $X_1$ 和 $X_2$ 是否独立？为什么？
    \item[(3)] $Z = X_1 - X_2$ 的分布律。
\end{enumerate}
\end{exercise}
\begin{solution}
(1) 先利用题目给出的特殊条件 $P\{X_1X_2=0\}=1$。这说明 $X_1X_2\neq 0$ 的情形根本不会发生，因此
\[
P\{X_1 = -1, X_2 = 1\} = 0, \quad P\{X_1 = 1, X_2 = 1\} = 0.
\]

接着按边缘分布补齐其余格子。由 $P\{X_2=1\}=1/2$，
\begin{align*}
    P\{X_2 = 1\} &= \frac{1}{2} = P\{X_1=-1,X_2=1\} + P\{X_1=0,X_2=1\} + P\{X_1=1,X_2=1\} \\
    &= 0 + P\{X_1=0,X_2=1\} + 0
\end{align*}
故 $P\{X_1=0, X_2=1\} = 1/2$。

再由 $X_1$ 的边缘分布，
\[
P\{X_1=-1,X_2=0\}=P\{X_1=-1\}-P\{X_1=-1,X_2=1\}=\frac{1}{4},
\]
\[
P\{X_1=1,X_2=0\}=P\{X_1=1\}-P\{X_1=1,X_2=1\}=\frac{1}{4},
\]
\[
P\{X_1=0,X_2=0\}=P\{X_1=0\}-P\{X_1=0,X_2=1\}=\frac{1}{2}-\frac{1}{2}=0.
\]
因此联合分布律为
\begin{center}
\begin{tabular}{c|ccc|c}
\toprule
$X_2 \backslash X_1$ & $-1$ & $0$ & $1$ & $p_{\cdot j}$ \\
\midrule
$0$ & $\dfrac{1}{4}$ & $0$ & $\dfrac{1}{4}$ & $\dfrac{1}{2}$ \\[6pt]
$1$ & $0$ & $\dfrac{1}{2}$ & $0$ & $\dfrac{1}{2}$ \\
\midrule
$p_{i\cdot}$ & $\dfrac{1}{4}$ & $\dfrac{1}{2}$ & $\dfrac{1}{4}$ & $1$ \\
\bottomrule
\end{tabular}
\end{center}

(2) 判断独立性时，只需找一组反例即可。例如
\[
P\{X_2=0,\ X_1=-1\}=\frac{1}{4},
\]
而
\[
P\{X_2=0\}\,P\{X_1=-1\}
=\frac{1}{2}\times\frac{1}{4}
=\frac{1}{8}.
\]
两者不相等，所以 $X_1$ 与 $X_2$ 不独立。

(3) 现在求 $Z=X_1-X_2$ 的分布。先列出可能取值：
\[
Z\in\{-2,-1,0,1\}.
\]
再逐个由联合分布表求概率：
\begin{align*}
    P\{Z = -2\} &= P\{X_1 = -1, X_2 = 1\} = 0, \\
    P\{Z = -1\} &= P\{X_1 = -1, X_2 = 0\} + P\{X_1 = 0, X_2 = 1\} = \frac{1}{4} + \frac{1}{2} = \frac{3}{4}, \\
    P\{Z = 0\} &= P\{X_1 = 0, X_2 = 0\} + P\{X_1 = 1, X_2 = 1\} = 0 + 0 = 0, \\
P\{Z = 1\} &= P\{X_1 = 1, X_2 = 0\} = \frac{1}{4}.
\end{align*}
\[
Z \sim \begin{pmatrix} -2 & -1 & 0 & 1 \\ 0 & 3/4 & 0 & 1/4 \end{pmatrix}.
\]
\end{solution}



\begin{exercise}[由独立性求未知参数] \label{exer:3-3c}
设二维随机变量 $(X,Y)$ 的联合分布律为
\begin{center}
\begin{tabular}{c|cccccc}
\toprule
$(X,Y)$ & $(1,1)$ & $(1,2)$ & $(1,3)$ & $(2,1)$ & $(2,2)$ & $(2,3)$ \\
\midrule
$P$ & $\dfrac{1}{6}$ & $\dfrac{1}{9}$ & $\dfrac{1}{18}$ & $\dfrac{1}{3}$ & $\alpha$ & $\beta$ \\
\bottomrule
\end{tabular}
\end{center}
若 $X$ 与 $Y$ 相互独立，求 $\alpha$ 和 $\beta$。
\end{exercise}
\begin{solution}
列表格（联合分布律）：
\begin{center}
\begin{tabular}{c|ccc|c}
\toprule
$X \backslash Y$ & $1$ & $2$ & $3$ & $p_{\cdot j}$ \\
\midrule
$1$ & $\dfrac{1}{6}$ & $\dfrac19$ & $\dfrac{1}{18}$ & $\dfrac{1}{3}$ \\[6pt]
$2$ & $\dfrac13$ & $\alpha$ & $\beta$ & $\dfrac{2}{3}$ \\
\midrule
$p_{i\cdot}$ & $\dfrac{1}{2}$ & $\dfrac{1}{9}+\alpha$ & $\dfrac{1}{18}+\beta$ & $1$ \\
\bottomrule
\end{tabular}
\end{center}
显然$\alpha+\beta=\frac13$，然后随便选一行一列的边缘概率相乘等于单元格内部的概率：
\[\frac13(\frac19+\alpha)=\frac19\implies \alpha=\frac29\implies \beta=\frac19\]
\begin{note}
本题的解题技巧是利用独立性的比例关系 $\dfrac{a_1}{a_2} = \dfrac{a_1 b_j}{a_2 b_j}$（对任意 $j$），由已知项快速求出未知项。
\end{note}
\end{solution}

\begin{exercise}[行列式分布] \label{exer:3-c6}
设随机变量 $X_1, X_2, X_3, X_4$ 相互独立且同分布，$P\{X_i=0\} = 0.6$，$P\{X_i=1\} = 0.4$（$i=1,2,3,4$），求行列式 $Z = \begin{vmatrix} X_1 & X_2 \\ X_3 & X_4 \end{vmatrix}$ 的概率分布。
\end{exercise}
\begin{solution}
$Z = X_1X_4 - X_2X_3$，可能取值为 $\{-1, 0, 1\}$。
\begin{align*}
P\{Z=1\} &= P\{X_1X_4=1\}\cdot P\{X_2X_3=0\}\\
& = P\{X_1=1,X_4=1\}\cdot P\{(X_2=0)\cup(X_3=0)\}\\
&= 0.4^2 \cdot (1-0.4^2) = 0.16 \times 0.84 = 0.1344.\\
P\{Z=-1\} &= P\{X_1X_4=0\}\cdot P\{X_2X_3=1\} \\
&= (1-0.4^2)\cdot 0.4^2 = 0.84 \times 0.16 = 0.1344\\
P\{Z=0\} &= 1 - 0.1344 - 0.1344 = 0.7312.
\end{align*}
所以
\[
Z \sim \begin{pmatrix} -1 & 0 & 1 \\ 0.1344 & 0.7312 & 0.1344 \end{pmatrix}.
\]
\end{solution}

\begin{exercise}[独立性填表] \label{exer:3-c7}
设 $X$ 与 $Y$ 相互独立，下表是 $(X,Y)$ 的联合分布律以及边缘分布律中的部分数值，试将其余数值填入空白处：
\begin{center}
\begin{tabular}{c|ccc|c}
\toprule
$X \backslash Y$ & $y_1$ & $y_2$ & $y_3$ & $P\{X=x_i\}$ \\
\midrule
$x_1$ & & $1/8$ & & \\
$x_2$ & $1/8$ & & & \\
\midrule
$P\{Y=y_j\}$ & $1/6$ & & & $1$ \\
\bottomrule
\end{tabular}
\end{center}
\end{exercise}
\begin{solution}
独立性给出的核心关系是
\[
p_{ij}=p_{i\cdot}p_{\cdot j}.
\]
因此这题的做法不是“凭感觉填表”，而是先用已知单元格反推出边缘分布，再把整张表补齐。

由
\[
p_{21}=p_{2\cdot}p_{\cdot 1}=\frac{1}{8},
\qquad
p_{\cdot 1}=\frac{1}{6},
\]
得
\[
p_{2\cdot}=\frac{1/8}{1/6}=\frac{3}{4}.
\]
于是
\[
p_{1\cdot}=1-p_{2\cdot}=1-\frac{3}{4}=\frac{1}{4}.
\]

再由
\[
p_{12}=p_{1\cdot}p_{\cdot 2}=\frac{1}{8},
\]
得
\[
p_{\cdot 2}=\frac{1/8}{1/4}=\frac{1}{2}.
\]
最后利用边缘概率之和为 $1$，
\[
p_{\cdot 3}=1-p_{\cdot 1}-p_{\cdot 2}=1-\frac{1}{6}-\frac{1}{2}=\frac{1}{3}.
\]

这样三列、两行的边缘分布都已确定，剩余单元格直接用“行概率 $\times$ 列概率”即可。填满后得到
\begin{center}
\begin{tabular}{c|ccc|c}
\toprule
$X \backslash Y$ & $y_1$ & $y_2$ & $y_3$ & $P\{X=x_i\}$ \\
\midrule
$x_1$ & $1/24$ & $1/8$ & $1/12$ & $1/4$ \\[4pt]
$x_2$ & $1/8$ & $3/8$ & $1/4$ & $3/4$ \\
\midrule
$P\{Y=y_j\}$ & $1/6$ & $1/2$ & $1/3$ & $1$ \\
\bottomrule
\end{tabular}
\end{center}
最后检验一下更稳妥：行和分别是
\[
\frac{1}{24}+\frac{1}{8}+\frac{1}{12}=\frac{1}{4},\qquad
\frac{1}{8}+\frac{3}{8}+\frac{1}{4}=\frac{3}{4},
\]
列和分别是
\[
\frac{1}{24}+\frac{1}{8}=\frac{1}{6},\qquad
\frac{1}{8}+\frac{3}{8}=\frac{1}{2},\qquad
\frac{1}{12}+\frac{1}{4}=\frac{1}{3}.
\]
与边缘分布完全一致，所以填表正确。
\end{solution}



\begin{exercise}[无放回取球] \label{exer:3-c8}
袋中有编号为 $1,1,2,3$ 的四个球，从中无放回地取两次，每次任取一个。设 $X_1, X_2$ 分别为第一次、第二次取得的球的号码。求：
\begin{enumerate}
    \item[(1)] $(X_1,X_2)$ 的联合分布律，并判断 $X_1$ 与 $X_2$ 的独立性；
    \item[(2)] 在 $X_2=2$ 的条件下，$X_1$ 的条件分布律；
    \item[(3)] $Y = X_1X_2$ 的分布。
\end{enumerate}
\end{exercise}
\begin{solution}
(1) 四个球编号 $\{1,1,2,3\}$，无放回取两次，共 $4\times 3 = 12$ 种等可能结果。
\begin{center}
\begin{tabular}{c|ccc|c}
\toprule
$X_1 \backslash X_2$ & $1$ & $2$ & $3$ & $p_{i\cdot}$ \\
\midrule
$1$ & $1/6$ & $1/6$ & $1/6$ & $1/2$ \\[4pt]
$2$ & $1/6$ & $0$ & $1/12$ & $1/4$ \\[4pt]
$3$ & $1/6$ & $1/12$ & $0$ & $1/4$ \\
\midrule
$p_{\cdot j}$ & $1/2$ & $1/4$ & $1/4$ & $1$ \\
\bottomrule
\end{tabular}
\end{center}
独立性：$p_{11} = 1/6$，$p_{1\cdot}\cdot p_{\cdot 1} = (1/2)(1/2) = 1/4 \neq 1/6$，故 $X_1$ 与 $X_2$ \textbf{不独立}。\\
(2) $P\{X_2=2\} = p_{\cdot 2} = 1/4$。
\begin{align*}
P\{X_1=1 \mid X_2=2\} =& \frac{p_{12}}{p_{\cdot 2}} = \frac{1/6}{1/4} = \frac{2}{3}\\
P\{X_1=2 \mid X_2=2\} = &\frac{p_{22}}{p_{\cdot 2}} = \frac{0}{1/4} = 0\\
 P\{X_1=3 \mid X_2=2\} = &\frac{p_{32}}{p_{\cdot 2}} = \frac{1/12}{1/4} = \frac{1}{3}.
\end{align*}
(3) $Y = X_1X_2$ 的可能取值及概率：
\begin{align*}
P\{Y=1\} &= P\{X_1=1,X_2=1\} = 1/6, \\
P\{Y=2\} &= P\{X_1=1,X_2=2\} + P\{X_1=2,X_2=1\} = 1/6 + 1/6 = 1/3, \\
P\{Y=3\} &= P\{X_1=1,X_2=3\} + P\{X_1=3,X_2=1\} = 1/6 + 1/6 = 1/3, \\
P\{Y=6\} &= P\{X_1=2,X_2=3\} + P\{X_1=3,X_2=2\} = 1/12 + 1/12 = 1/6.
\end{align*}
\[
Y \sim \begin{pmatrix} 1 & 2 & 3 & 6 \\ 1/6 & 1/3 & 1/3 & 1/6 \end{pmatrix}.
\]
\end{solution}

\begin{exercise}[离散型函数的分布] \label{exer:3-c9}
已知 $(X,Y)$ 的概率分布为
\begin{center}
\begin{tabular}{c|cc}
\toprule
$X \backslash Y$ & $-1$ & $1$ \\
\midrule
$-1$ & $1/3$ & $1/6$ \\[4pt]
$0$ & $1/3$ & $1/6$ \\
\bottomrule
\end{tabular}
\end{center}
求：(1) $Z = X - Y$ 的概率分布；(2) 设 $U = \max\{X,Y\}$，$V = \min\{X,Y\}$，求 $(U,V)$ 及 $UV$ 的分布。
\end{exercise}
\begin{solution}
(1) $Z$ 的可能取值：$Z(-1,-1)=0$，$Z(-1,1)=-2$，$Z(0,-1)=1$，$Z(0,1)=-1$。
\[
P\{Z=-2\} = 1/6,\quad P\{Z=-1\} = 1/6,\quad P\{Z=0\} = 1/3,\quad P\{Z=1\} = 1/3.
\]
\[
Z \sim \begin{pmatrix} -2 & -1 & 0 & 1 \\ 1/6 & 1/6 & 1/3 & 1/3 \end{pmatrix}.
\]
(2) $(U,V)$ 的可能取值：
\begin{align*}
(X,Y)=(-1,-1) &: U=-1,\ V=-1, \quad P = 1/3. \\
(X,Y)=(-1,1) &: U=1,\ V=-1, \quad P = 1/6. \\
(X,Y)=(0,-1) &: U=0,\ V=-1, \quad P = 1/3. \\
(X,Y)=(0,1) &: U=1,\ V=0, \quad P = 1/6.
\end{align*}
$(U,V)$ 联合分布律：
\begin{center}
\begin{tabular}{c|ccc}
\toprule
$V \backslash U$ & $-1$ & $0$ & $1$ \\
\midrule
$-1$ & $1/3$ & $1/3$ & $1/6$ \\[4pt]
$0$ & $0$ & $0$ & $1/6$ \\
\bottomrule
\end{tabular}
\end{center}
$UV$ 的分布：$UV \in \{-1, 0, 1\}$。
\[
UV \sim \begin{pmatrix} -1 & 0 & 1 \\ 1/6 & 1/2 & 1/3 \end{pmatrix}.
\]
\end{solution}


\begin{exercise}[事件的示性函数] \label{exer:3-d1}
设 $A,B$ 为两个随机事件，$P(A) = 1/4$，$P(B) = 1/6$，$P(AB) = 1/12$。令
\[
X = \begin{cases} 1, & A \text{发生}, \\ 0, & A \text{不发生}, \end{cases} \quad Y = \begin{cases} 1, & B \text{发生}, \\ 0, & B \text{不发生}. \end{cases}
\]
求：(1) $(X,Y)$ 的联合分布律；(2) $Z = X^2+Y^2$ 的分布律；(3) $W = \min(X,Y)$ 的分布律。
\end{exercise}
\begin{solution}
(1) 由事件的概率关系：
\begin{align*}
P(X=0,Y=0) &= P(\bar{A}\bar{B}) = 1 - P(A\cup B) = 1 - (P(A)+P(B)-P(AB)) = \frac{2}{3}, \\
P(X=1,Y=0) &= P(A\bar{B}) = P(A) - P(AB) = \frac{1}{4}-\frac{1}{12} = \frac{1}{6}, \\
P(X=0,Y=1) &= P(\bar{A}B) = P(B) - P(AB) = \frac{1}{6}-\frac{1}{12} = \frac{1}{12}, \\
P(X=1,Y=1) &= P(AB) = \frac{1}{12}.
\end{align*}
\begin{center}
\begin{tabular}{c|cc|c}
\toprule
$X \backslash Y$ & $0$ & $1$ & $P\{X=x_i\}$ \\
\midrule
$0$ & $2/3$ & $1/12$ & $3/4$ \\[4pt]
$1$ & $1/6$ & $1/12$ & $1/4$ \\
\midrule
$P\{Y=y_j\}$ & $5/6$ & $1/6$ & $1$ \\
\bottomrule
\end{tabular}
\end{center}
(2) $Z = X^2+Y^2$：$(0,0)\to 0$，$(1,0)$ 或 $(0,1)\to 1$，$(1,1)\to 2$。
\[
P(Z=0) = 2/3,\quad P(Z=1) = 1/6+1/12 = 1/4,\quad P(Z=2) = 1/12.
\]
\[
Z \sim \begin{pmatrix} 0 & 1 & 2 \\ 2/3 & 1/4 & 1/12 \end{pmatrix}.
\]
(3) $W = \min(X,Y)$：$W=0$ 当至少一个为 $0$，$W=1$ 仅当 $(1,1)$。
\[
P(W=0) = 2/3+1/6+1/12 = 11/12,\quad P(W=1) = 1/12.
\]
\end{solution}



\begin{exercise}[独立性条件求参数] \label{exer:3-d2}
设 $X,Y$ 相互独立，$P(X=1)=P(Y=1)=p$，$P(X=0)=P(Y=0)=1-p$（$0<p<1$）。令
\[
Z = \begin{cases} 1, & X+Y \text{为偶数}, \\ 0, & X+Y \text{为奇数}. \end{cases}
\]
要使 $X$ 与 $Z$ 独立，则 $p = $（\quad）
\end{exercise}
\begin{solution}
$X+Y$ 为偶数当且仅当 $(X,Y) = (0,0)$ 或 $(1,1)$，故
\[
P(Z=1) = (1-p)^2+p^2,\quad P(Z=0) = 2p(1-p).
\]
又 $P(X=1,Z=1) = P(X=1,Y=1) = p^2$。由独立性 $P(X=1,Z=1) = P(X=1)\,P(Z=1)$：
\[
p^2 = p\left[(1-p)^2+p^2\right].
\]
若 $p=0$ 或 $p=1$，$Z$ 退化为常数，平凡独立。对 $0<p<1$，两边除以 $p$：
\[
p = 1-2p+2p^2 \implies 2p^2-3p+1 = 0 \implies (2p-1)(p-1) = 0 \implies p = 1/2.
\]
\begin{note}
验证 $p=1/2$：$P(Z=1)=1/2$，$P(X=1,Z=1)=1/4=P(X=1)P(Z=1)$；$P(X=1,Z=0)=1/4=P(X=1)P(Z=0)$ 。
\end{note}
\end{solution}



\begin{exercise}[离散型函数分布：和、max、min] \label{exer:3-d3}
$X,Y$ 独立同分布，$P(X=0)=1/2$，$P(X=1)=1/4$，$P(X=2)=1/4$。求：
(1) $Z=X+Y$；(2) $U=\max(X,Y)$；(3) $W=\min(X,Y)$ 的分布律。
\end{exercise}
\begin{solution}
(1) $Z=X+Y$ 的可能取值为 $0,1,2,3,4$。由独立性逐项计算：
\begin{align*}
P(Z=0) &= P(X=0)P(Y=0) = \frac{1}{2}\cdot\frac{1}{2} = \frac{1}{4}, \\
P(Z=1) &= P(X=0)P(Y=1)+P(X=1)P(Y=0) = 2\cdot\frac{1}{2}\cdot\frac{1}{4} = \frac{1}{4}, \\
P(Z=2) &= P(0,2)+P(2,0)+P(1,1) = 2\cdot\frac{1}{2}\cdot\frac{1}{4}+\frac{1}{4}\cdot\frac{1}{4} = \frac{5}{16}, \\
P(Z=3) &= P(1,2)+P(2,1) = 2\cdot\frac{1}{4}\cdot\frac{1}{4} = \frac{1}{8}, \\
P(Z=4) &= P(2,2) = \frac{1}{16}.
\end{align*}
(2) $U=\max(X,Y)$：利用 $P(U \leqslant k) = P(X \leqslant k)^2$：
\[
P(U=0) = \frac{1}{4},\quad P(U=1) = \left(\frac{3}{4}\right)^2 - \left(\frac{1}{2}\right)^2 = \frac{5}{16},\quad P(U=2) = 1-\frac{1}{4}-\frac{5}{16} = \frac{7}{16}.
\]
(3) $W=\min(X,Y)$：利用 $P(W \geqslant k) = P(X \geqslant k)^2$：
\[
P(W=0) = 1-\left(\frac{1}{2}\right)^2 = \frac{3}{4},\quad P(W=1) = \left(\frac{1}{2}\right)^2-\left(\frac{1}{4}\right)^2 = \frac{3}{16},\quad P(W=2) = \frac{1}{16}.
\]
\begin{center}
\begin{tabular}{c|ccccc}
\toprule
$Z$ & $0$ & $1$ & $2$ & $3$ & $4$ \\
\midrule
$P$ & $1/4$ & $1/4$ & $5/16$ & $1/8$ & $1/16$ \\
\bottomrule
\end{tabular}
\quad
\begin{tabular}{c|ccc}
\toprule
$U$ & $0$ & $1$ & $2$ \\
\midrule
$P$ & $1/4$ & $5/16$ & $7/16$ \\
\bottomrule
\end{tabular}
\quad
\begin{tabular}{c|ccc}
\toprule
$W$ & $0$ & $1$ & $2$ \\
\midrule
$P$ & $3/4$ & $3/16$ & $1/16$ \\
\bottomrule
\end{tabular}
\end{center}
\end{solution}



\begin{exercise}[泊松变量的复合函数] \label{exer:3-d4}
设 $X_1,X_2,X_3$ 相互独立且同服从 $P(1)$，记
\[
X = \begin{cases} 1, & X_1+X_2=1, \\ 0, & X_1+X_2\neq 1, \end{cases} \quad Y = \begin{cases} 1, & X_2+X_3=1, \\ 0, & X_2+X_3\neq 1. \end{cases}
\]
求：(1) $(X,Y)$ 的联合分布律；(2) 边缘分布律；(3) $P(X+Y \leqslant 1)$。
\end{exercise}
\begin{solution}
(1) $\{X=1,Y=1\}$ 要求 $X_1+X_2=1$ 且 $X_2+X_3=1$。分两种情况：\\
\textbf{情形 $X_2=0$}：需 $X_1=1$，$X_3=1$。$P = e^{-1}\cdot e^{-1}\cdot e^{-1} = e^{-3}$。\\
\textbf{情形 $X_2=1$}：需 $X_1=0$，$X_3=0$。$P = e^{-1}\cdot e^{-1}\cdot e^{-1} = e^{-3}$。\\
当 $X_2 \geqslant 2$ 时 $X_1+X_2 \geqslant 2 \neq 1$，不贡献。
\[
P(X=1,Y=1) = 2e^{-3}.
\]
边缘概率：
\[
P(X=1) = P(X_1+X_2=1) = P(1,0)+P(0,1) = 2e^{-2}, \qquad P(Y=1) = P(X_2+X_3=1) = 2e^{-2}.
\]
\begin{align*}
P(X=0,Y=1) &= P(Y=1)-P(X=1,Y=1) = 2e^{-2}-2e^{-3}, \\
P(X=1,Y=0) &= P(X=1)-P(X=1,Y=1) = 2e^{-2}-2e^{-3}, \\
P(X=0,Y=0) &= 1-2e^{-3}-(2e^{-2}-2e^{-3})-(2e^{-2}-2e^{-3}) = 1-4e^{-2}+2e^{-3}.
\end{align*}
\begin{center}
\begin{tabular}{c|cc|c}
\toprule
$X \backslash Y$ & $1$ & $0$ & $P\{X=x_i\}$ \\
\midrule
$1$ & $2e^{-3}$ & $2e^{-2}-2e^{-3}$ & $2e^{-2}$ \\[4pt]
$0$ & $2e^{-2}-2e^{-3}$ & $1-4e^{-2}+2e^{-3}$ & $1-2e^{-2}$ \\
\midrule
$P\{Y=y_j\}$ & $2e^{-2}$ & $1-2e^{-2}$ & $1$ \\
\bottomrule
\end{tabular}
\end{center}
(2) 边缘分布律见上表右侧和下方。\\
(3) $P(X+Y \leqslant 1) = 1-P(X=1,Y=1) = 1-2e^{-3}$。
\begin{note}
$X$ 与 $Y$ 不独立：$P(X=1,Y=1)=2e^{-3}$，而 $P(X=1)P(Y=1)=4e^{-4} \neq 2e^{-3}$。这是因为 $X$ 和 $Y$ 共同依赖 $X_2$。
\end{note}
\end{solution}



\begin{exercise}[离散条件均匀分布] \label{exer:3-d5}
从 $\{1,2,3\}$ 中等可能地取一个数记为 $X$，再从 $\{1,\dots,X\}$ 中等可能地取一个整数记为 $Y$。求 $(X,Y)$ 的联合分布律和边缘分布律，并计算 $P(X+Y=4)$。
\end{exercise}
\begin{solution}
$P(X=a,Y=b) = P(X=a)\,P(Y=b|X=a) = \dfrac{1}{3}\cdot\dfrac{1}{a}$（$1 \leqslant b \leqslant a$）。
逐项列出联合分布律：
\begin{center}
\begin{tabular}{c|ccc|c}
\toprule
$X \backslash Y$ & $1$ & $2$ & $3$ & $P\{X=x_i\}$ \\
\midrule
$1$ & $1/3$ & $0$ & $0$ & $1/3$ \\[4pt]
$2$ & $1/6$ & $1/6$ & $0$ & $1/3$ \\[4pt]
$3$ & $1/9$ & $1/9$ & $1/9$ & $1/3$ \\
\midrule
$P\{Y=y_j\}$ & $1/3+1/6+1/9$ & $1/6+1/9$ & $1/9$ & $1$ \\
& $= 11/18$ & $= 5/18$ & $= 1/9$ & \\
\bottomrule
\end{tabular}
\end{center}
$P(X+Y=4)$：满足 $a+b=4$ 且 $1\leqslant b\leqslant a\leqslant 3$ 的对为 $(2,2)$ 和 $(3,1)$。
\[
P(X+Y=4) = P(X=2,Y=2)+P(X=3,Y=1) = \frac{1}{6}+\frac{1}{9} = \frac{5}{18}.
\]
\end{solution}


\section{二维连续型随机变量}

\begin{note}
\textbf{自学导读：} 本节是第 2 章"连续型随机变量"的二维推广。回忆一维时：概率密度 $f(x)\geqslant 0$，积分为 1，概率 = 面积。现在推广到二维：概率密度 $f(x,y)\geqslant 0$，二重积分为 1，概率 = \textbf{体积}（密度曲面下方的体积）。

操作上的变化：一维求 $P(a<X<b)$ 是定积分；二维求 $P((X,Y)\in G)$ 是二重积分。如果你对二重积分生疏，建议先复习"先对 $x$ 积分再对 $y$ 积分"的累次积分写法，本节所有计算都依赖这一技巧。
\end{note}

\begin{definition}[二维连续型随机变量] \label{def:2d-continuous}
设 $(X,Y)$ 为二维随机变量，若存在非负可积函数 $f(x,y)$，使得对于任意实数 $x,y$，都有
\[
F(x,y) = \int_{-\infty}^{y}\int_{-\infty}^{x}f(u,v)\,\mathrm{d}u\,\mathrm{d}v,
\]
则称 $(X,Y)$ 为\textbf{二维连续型随机变量}，称 $f(x,y)$ 为 $(X,Y)$ 的\textbf{联合概率密度函数}（简称\textbf{概率密度}），记为 $(X,Y) \sim f(x,y)$。
\end{definition}

\begin{property}\label{prop:joint-pdf}
联合概率密度 $f(x,y)$ 具有以下性质：
\begin{enumerate}
    \item \textbf{非负性}：$f(x,y) \geqslant 0$；
    \item \textbf{归一性}：$\displaystyle\int_{-\infty}^{+\infty}\int_{-\infty}^{+\infty}f(x,y)\,\mathrm{d}x\,\mathrm{d}y = 1$；
    \item 设 $G$ 为 $xOy$ 平面上的某个区域，则
    \[
    P\{(X,Y) \in G\} = \iint_G f(x,y)\,\mathrm{d}x\,\mathrm{d}y;
    \]
    \item 若 $f(x,y)$ 在点 $(x,y)$ 处连续，则
    \[
    f(x,y) = \frac{\partial^2 F(x,y)}{\partial x\,\partial y}.
    \]
\end{enumerate}
\end{property}

\begin{note}
几何意义：$z = f(x,y)$ 表示空间中的一个曲面，该曲面与 $xOy$ 平面之间所围成的空间区域的体积为 $1$。$P\{(X,Y) \in G\}$ 的值等于以 $G$ 为底、以曲面 $z = f(x,y)$ 为顶面的曲顶柱体的体积。改变 $f(x,y)$ 在有限个点（甚至一条曲线上）的值（仍取非负值），$f(x,y)$ 仍然是 $(X,Y)$ 的概率密度。
\end{note}


\begin{exercise}[概率计算] \label{exer:3-c12}
$f(x,y) = \begin{cases} Ae^{-y}, & 0 < x < y, \\ 0, & \text{其他}. \end{cases}$，求 $A$，$P\{X+Y \geqslant 1\}$，$P\{X/Y \leqslant 1/2\}$。
\end{exercise}
\begin{solution}
(1) \textbf{求解常数 $\boldsymbol{A}$}：由联合概率密度函数的归一性（即全空间积分为 1）可知：
\[
\iint\limits_{\mathbb{R}^2} f(x,y) \,\mathrm{d}x\,\mathrm{d}y = 1
\]
根据题目给定的非零区域 $0 < x < y < +\infty$，我们先对 $x$ 积分，再对 $y$ 积分，确定上下限如下：
\[
\int_{0}^{+\infty} \left( \int_{0}^{y} A e^{-y} \,\mathrm{d}x \right) \mathrm{d}y = 1
\]
计算内层积分（此时 $y$ 视为常数）：
\[
A \int_{0}^{+\infty} e^{-y} \Big[ x \Big]_{0}^{y} \,\mathrm{d}y = A \int_{0}^{+\infty} y e^{-y} \,\mathrm{d}y = 1
\]
利用分部积分法或 Gamma 函数性质 $\Gamma(2) = \int_0^\infty y e^{-y} \mathrm{d}y = 1! = 1$，易得：
\[
A \times 1 = 1 \implies A = 1
\]
(2) \textbf{求解 $\boldsymbol{P\{X+Y \geqslant 1\}}$}：若直接求解 $X+Y \geqslant 1$，积分区域是一个向右上角无限延伸的无界区域，计算较为繁琐。因此，我们利用对立事件（补集）来简化计算：计算图形封闭的 $P\{X+Y < 1\}$。
\[
P\{X+Y \geqslant 1\} = 1 - P\{X+Y < 1\}
\]
事件 $\{X+Y < 1\}$ 对应的概率本质上是联合密度函数在特定二维区域 $D$ 上的二重积分。该区域 $D$ 由以下三个不等式共同围成：
\[
\begin{cases}
0 < x < y & \text{(密度函数非零的固有条件)} \\
x + y < 1 & \text{(题目要求的事件条件)}
\end{cases}
\]
在平面直角坐标系中，直线 $y = x$ 与直线 $x + y = 1$ 的交点为 $(1/2, 1/2)$。观察该三角形区域 $D$，我们可以采用 $X$-型区域（先对 $y$ 积分，后对 $x$ 积分）来设置上下限：
外层 $x$ 的投影范围是 $(0, 1/2)$；对于固定的 $x$，内层 $y$ 的下界是直线 $y=x$，上界是直线 $y=1-x$。因此：
\begin{align*}
P\{X+Y < 1\} &= \iint\limits_{D} f(x,y) \,\mathrm{d}x\,\mathrm{d}y = \int_{0}^{1/2} \mathrm{d}x \int_{x}^{1-x} e^{-y} \,\mathrm{d}y \\
&= \int_{0}^{1/2} \Big[ -e^{-y} \Big]_{x}^{1-x} \mathrm{d}x = \int_{0}^{1/2} \left( e^{-x} - e^{-(1-x)} \right) \mathrm{d}x \\
&= \Big[ -e^{-x} - e^{x-1} \Big]_{0}^{1/2} = \left( -e^{-1/2} - e^{-1/2} \right) - \left( -1 - e^{-1} \right) \\
&= 1 + e^{-1} - 2e^{-1/2}
\end{align*}
最后，代回对立事件公式中：
\[
P\{X+Y \geqslant 1\} = 1 - \left( 1 + e^{-1} - 2e^{-1/2} \right) = 2e^{-1/2} - e^{-1}
\]
(3) \textbf{求解 $\boldsymbol{P\{X/Y \leqslant 1/2\}}$}：由于在非零区域内 $y > 0$，不等式 $X/Y \leqslant 1/2$ 可直接等价变形为 $X \leqslant Y/2$。
同样地，我们需要寻找联合密度函数非零区域 $\{0 < x < y\}$ 与目标事件区域 $\{x \leqslant y/2\}$ 的交集。综合这两个条件，得到真正的积分区域：$0 < x \leqslant y/2 < +\infty$。
这次我们采用 $Y$-型区域（先对 $x$ 积分，后对 $y$ 积分）来设置上下限更为简便：
外层 $y$ 的范围是 $(0, +\infty)$；对于固定的 $y$，内层 $x$ 的下界是 $0$，上界是 $y/2$。
\begin{align*}
P\{X \leqslant Y/2\} &= \int_{0}^{+\infty} \mathrm{d}y \int_{0}^{y/2} e^{-y} \,\mathrm{d}x = \int_{0}^{+\infty} e^{-y} \Big[ x \Big]_{0}^{y/2} \mathrm{d}y \\
&= \int_{0}^{+\infty} \frac{y}{2} e^{-y} \,\mathrm{d}y = \frac{1}{2} \int_{0}^{+\infty} y e^{-y} \,\mathrm{d}y
\end{align*}
由第(1)问已知的结论 $\int_{0}^{+\infty} y e^{-y} \,\mathrm{d}y = 1$，故：
\[
P\{X/Y \leqslant 1/2\} = \frac{1}{2} \times 1 = \frac{1}{2}
\]
\end{solution}

\begin{exercise}[混合型：连续+离散] \label{exer:3-c13}
$X \sim E(4)$，$Y$ 为离散型随机变量：$P\{Y=-1\} = 1/4$，$P\{Y=1\} = 3/4$，且 $X$ 与 $Y$ 相互独立。求 $P\{X-Y \leqslant 1\}$ 和 $P\{XY \leqslant 2\}$。
\end{exercise}
\begin{solution}
已知 $X \sim E(4)$，即 $X$ 服从参数 $\lambda = 4$ 的指数分布。根据指数分布的定义，其概率密度函数 $f_X(x)$ 和累积分布函数 $F_X(x)$ 分别为：
\[
f_X(x) = \begin{cases} 4e^{-4x}, & x > 0 \\ 0, & x \leqslant 0 \end{cases}
\quad \text{与} \quad
F_X(x) = P\{X \leqslant x\} = \begin{cases} 1 - e^{-4x}, & x > 0 \\ 0, & x \leqslant 0 \end{cases}
\]
已知 $Y$ 的分布律为：$P\{Y=-1\} = \frac{1}{4}$，$P\{Y=1\} = \frac{3}{4}$。
最关键的条件是：\textbf{$\boldsymbol{X}$ 与 $\boldsymbol{Y}$ 相互独立}，且本题中的 \(Y\) 只取 \(-1,1\) 两个离散值。因此对任意事件 \(A\)，只要 \(P\{Y=y\}>0\)，就有
\[
P\{X\in A\mid Y=y\}=P\{X\in A\},
\]
等价地，
\[
P(\{X\in A\}\cap\{Y=y\})=P\{X\in A\}\cdot P\{Y=y\}.
\]
基于这种“离散 + 独立”的结构，我们可按 \(Y\) 的可能取值把复杂事件拆成两个情形，再用全概率公式逐项计算。\\
\textbf{（1）求解 $\boldsymbol{P\{X-Y \leqslant 1\}}$}:根据 $Y$ 的取值将目标事件 $\{X-Y \leqslant 1\}$ 分解为两个互斥子事件的和：
\begin{align*}
P\{X-Y \leqslant 1\} &= P\big( \{X-Y \leqslant 1\} \cap \{Y=-1\} \big) + P\big( \{X-Y \leqslant 1\} \cap \{Y=1\} \big) \\
&= P\big( \{X-(-1) \leqslant 1\} \cap \{Y=-1\} \big) + P\big( \{X-1 \leqslant 1\} \cap \{Y=1\} \big) \\
&= P\big( \{X \leqslant 0\} \cap \{Y=-1\} \big) + P\big( \{X \leqslant 2\} \cap \{Y=1\} \big)
\end{align*}
利用 $X, Y$ 的相互独立性，将联合概率拆分为边缘概率的乘积：
\begin{align*}
P\{X-Y \leqslant 1\} &= P\{X \leqslant 0\} \cdot P\{Y=-1\} + P\{X \leqslant 2\} \cdot P\{Y=1\}\\
P\{X \leqslant 0\} &= F_X(0) = 0 \quad (\text{因为 } X \text{ 恒大于 } 0) \\
P\{X \leqslant 2\} &= F_X(2) = 1 - e^{-4 \times 2} = 1 - e^{-8}
\end{align*}
将其代回主式：
\begin{align*}
P\{X-Y \leqslant 1\} &= 0 \times \frac{1}{4} + (1 - e^{-8}) \times \frac{3}{4} = \frac{3}{4}(1 - e^{-8})
\end{align*}
\textbf{（2）求解 $\boldsymbol{P\{XY \leqslant 2\}}$}:采用相同的全概率分解逻辑，按 $Y$ 的取值将事件 $\{XY \leqslant 2\}$ 拆分：
\begin{align*}
P\{XY \leqslant 2\} &= P\big( \{XY \leqslant 2\} \cap \{Y=-1\} \big) + P\big( \{XY \leqslant 2\} \cap \{Y=1\} \big) \\
&= P\big( \{X(-1) \leqslant 2\} \cap \{Y=-1\} \big) + P\big( \{X(1) \leqslant 2\} \cap \{Y=1\} \big) \\
&= P\big( \{-X \leqslant 2\} \cap \{Y=-1\} \big) + P\big( \{X \leqslant 2\} \cap \{Y=1\} \big)
\end{align*}
由于 $X, Y$ 独立，原式转化为：
\begin{align*}
P\{XY \leqslant 2\} &= P\{-X \leqslant 2\} \cdot P\{Y=-1\} + P\{X \leqslant 2\} \cdot P\{Y=1\} \\
&= P\{X \geqslant -2\} \cdot \frac{1}{4} + P\{X \leqslant 2\} \cdot \frac{3}{4}
\end{align*}
由于 $X \sim E(4)$，其取值范围严格为 $X > 0$。因此，事件 $\{X \geqslant -2\}$ 是一个\textbf{必然事件}（即 $X$ 的取值永远大于 0，当然也就大于 -2），故：
\[
P\{X \geqslant -2\} = 1
\]
再代入前面算出的 $P\{X \leqslant 2\} = 1 - e^{-8}$，得到：
\begin{align*}
P\{XY \leqslant 2\} &= 1 \times \frac{1}{4} + (1 - e^{-8}) \times \frac{3}{4} = \frac{1}{4} + \frac{3}{4} - \frac{3}{4} e^{-8} = 1 - \frac{3}{4} e^{-8}
\end{align*}
\end{solution}

\subsection{边缘概率密度}

\begin{definition}[边缘概率密度] \label{def:marginal-pdf}
设 $(X,Y) \sim f(x,y)$，则分别称
\[
f_X(x) = \int_{-\infty}^{+\infty}f(x,y)\,\mathrm{d}y, \quad f_Y(y) = \int_{-\infty}^{+\infty}f(x,y)\,\mathrm{d}x
\]
为 $(X,Y)$ 关于 $X$ 和关于 $Y$ 的\textbf{边缘概率密度}。
\end{definition}

相应的边缘分布函数分别为
\[
F_X(x) = \int_{-\infty}^{x}f_X(t)\,\mathrm{d}t = \int_{-\infty}^{x}\left[\int_{-\infty}^{+\infty}f(t,y)\,\mathrm{d}y\right]\mathrm{d}t,
\]
\[
F_Y(y) = \int_{-\infty}^{y}f_Y(t)\,\mathrm{d}t = \int_{-\infty}^{y}\left[\int_{-\infty}^{+\infty}f(x,t)\,\mathrm{d}x\right]\mathrm{d}t.
\]

\begin{exercise}[两种边缘密度计算] \label{exer:3-c10}
求下列联合密度的边缘概率密度：\\
（1）$f(x,y) = \begin{cases} e^{-y}, & 0 < x < y, \\ 0, & \text{其他}; \end{cases}$~~（2）$f_2(x,y) = \begin{cases} 5(x^2+y), & 0 < y < 1-x^2, \\ 0, & \text{其他}. \end{cases}$
\end{exercise}
\begin{solution}
(1)
\[
f_X(x) = \int_x^{+\infty}e^{-y}\,\mathrm{d}y = e^{-x}\ (x > 0), \qquad f_Y(y) = \int_0^y e^{-y}\,\mathrm{d}x = ye^{-y}\ (y > 0).
\]
即 $X \sim E(1)$，$Y \sim \Gamma(2,1)$。\\
(2) $f_2(x,y) \neq 0$ 要求 $0 < y < 1-x^2$，即 $|x| < 1$。
\[
f_X(x) = \int_0^{1-x^2}5(x^2+y)\,\mathrm{d}y = 5\left[x^2y + \frac{y^2}{2}\right]_0^{1-x^2} = 5\left[x^2(1-x^2) + \frac{(1-x^2)^2}{2}\right] = \frac{5}{2}(1-x^4),\quad |x|<1.
\]
\[
f_Y(y) = \int_{-\sqrt{1-y}}^{\sqrt{1-y}}5(x^2+y)\,\mathrm{d}x = 10\sqrt{1-y}\cdot\frac{1+2y}{3} = \frac{10}{3}(1+2y)\sqrt{1-y},\quad 0<y<1.
\]
\end{solution}

\begin{exercise}\label{exer:3-4}
设 $(X,Y)$ 是二维随机变量，$X$ 的边缘概率密度为
\[
f_X(x) = \begin{cases} 3x^2, & 0 < x < 1, \\ 0, & \text{其他}, \end{cases}
\]
在给定 $X = x$（$0 < x < 1$）的条件下，$Y$ 的条件概率密度为
\[
f_{Y|X}(y|x) = \begin{cases} \dfrac{3y^2}{x^3}, & 0 < y < x, \\ 0, & \text{其他}. \end{cases}
\]
\begin{enumerate}
    \item[(1)] 求 $(X,Y)$ 的联合概率密度 $f(x,y)$；
    \item[(2)] 求 $Y$ 的边缘概率密度 $f_Y(y)$；
    \item[(3)] 求 $P\{X > 2Y\}$。
\end{enumerate}
\end{exercise}
\begin{solution}
(1) 由概率密度乘法公式：
\[
f(x,y) = f_{Y|X}(y|x)\,f_X(x) = \begin{cases} \dfrac{3y^2}{x^3} \cdot 3x^2 = \dfrac{9y^2}{x}, & 0 < y < x < 1, \\ 0, & \text{其他}. \end{cases}
\]
(2) 当 $0 < y \leqslant 1$ 时，$f_Y(y) = \displaystyle\int_y^1 \frac{9y^2}{x}\,\mathrm{d}x = 9y^2\bigl[-\ln y\bigr]_y^1 = -9y^2\ln y$；其余范围 $f_Y(y) = 0$。\\
(3) $\{X > 2Y\}$ 等价于 $\{Y < X/2\}$，积分区域为 $\{(x,y): 0 < x < 1,\ 0 < y < x/2\}$：
\[
P\{X > 2Y\} = \int_0^1\mathrm{d}x\int_0^{x/2}\frac{9y^2}{x}\,\mathrm{d}y = \int_0^1 \frac{3x^2}{8}\,\mathrm{d}x = \left[\frac{x^3}{8}\right]_0^1 = \frac{1}{8}.
\]
\end{solution}

\subsection{条件概率密度}

\begin{definition}[条件概率密度] \label{def:conditional-pdf}
设 $(X,Y) \sim f(x,y)$，若边缘概率密度 $f_X(x) > 0$，则称
\[
f_{Y|X}(y|x) = \frac{f(x,y)}{f_X(x)}
\]
为在 $\{X = x\}$ 条件下 $Y$ 的\textbf{条件概率密度}。

同理，若 $f_Y(y) > 0$，则称
\[
f_{X|Y}(x|y) = \frac{f(x,y)}{f_Y(y)}
\]
为在 $\{Y = y\}$ 条件下 $X$ 的\textbf{条件概率密度}。
\end{definition}

\begin{property}[概率密度乘法公式]\label{prop:pdf-multiplication}
若 $f_X(x) > 0$，$f_Y(y) > 0$，则有
\[
f(x,y) = f_X(x)\,f_{Y|X}(y|x) = f_Y(y)\,f_{X|Y}(x|y).
\]
\end{property}

\begin{note}
\textbf{自学追问：} 条件概率密度和离散型的条件分布律在形式上完全平行——都是"联合除以边缘"。不同之处在于：离散型直接用概率比值 $p_{ij}/p_{\cdot j}$，连续型用密度比值 $f(x,y)/f_Y(y)$。后面做题碰到"在 $X=x$ 条件下求 $Y$ 的分布"，第一步永远是算出边缘密度 $f_X(x)$，然后做除法。

\textbf{易错点：} 条件概率密度 $f_{Y|X}(y|x)$ 中的 $x$ 是固定常数，不再是积分变量。写积分上下限时只对 $y$ 积分。另外，只有当 $f_X(x)>0$ 时条件密度才有定义；若 $x$ 落在边缘密度为 $0$ 的区域，条件分布无意义。
\end{note}

\begin{note}
\textbf{求条件密度的机械化四步（连续型）：}
\begin{enumerate}
    \item \textbf{画支撑域。} 把 $f(x,y)>0$ 的区域画在 $xOy$ 平面上。
    \item \textbf{求边缘密度。} 例如求 $f_{Y|X}(y|x)$ 时，先对 $y$ 积分得 $f_X(x)$：
    \[
    f_X(x)=\int_{\text{$y$的范围（由支撑域读出）}} f(x,y)\,\mathrm{d}y.
    \]
    \item \textbf{做除法。} $f_{Y|X}(y|x)=\dfrac{f(x,y)}{f_X(x)}$，此时 $x$ 视为常数。
    \item \textbf{确定条件密度的 $y$ 范围。} 在支撑域图上，固定 $x$ 画一条竖线，竖线截出的 $y$ 区间就是条件密度的定义域。
\end{enumerate}
\textbf{最常见错误：} 第4步中把 $y$ 的范围写成了不含 $x$ 的固定区间。实际上条件密度的 $y$ 范围通常依赖于 $x$ 的取值——必须回到支撑域图上读。
\end{note}

\begin{exercise}[判断独立性与条件概率] \label{exer:3-5d}
设随机变量 $X, Y$ 的联合概率密度函数为
\[
f(x,y) = \begin{cases} e^{-y}, & x > 0,\ y > x, \\ 0, & \text{其他}, \end{cases}
\]
(1) $X, Y$ 是否相互独立？(2) 求 $P\{X > 2 \mid Y < 4\}$。
\end{exercise}
\begin{solution}
(1) 求边缘概率密度：
\[
f_X(x) = \int_x^{+\infty} e^{-y}\,\mathrm{d}y = e^{-x}\ (x > 0), \qquad f_Y(y) = \int_0^y e^{-y}\,\mathrm{d}x = ye^{-y}\ (y > 0).
\]
由于 $f_X(x)\,f_Y(y) = e^{-x} \cdot ye^{-y} = ye^{-(x+y)} \neq e^{-y} = f(x,y)$（当 $y > x > 0$ 时），故 $X$ 与 $Y$ 不相互独立。

(2)
\[
P\{X > 2,\ Y < 4\} = \int_2^4\mathrm{d}y\int_2^y e^{-y}\,\mathrm{d}x = \int_2^4(y-2)e^{-y}\,\mathrm{d}y = e^{-2} - 3e^{-4},
\]
\[
P\{Y < 4\} = \int_0^4 ye^{-y}\,\mathrm{d}y = 1 - 5e^{-4},
\]
故
\[
P\{X > 2 \mid Y < 4\} = \frac{P\{X > 2,\ Y < 4\}}{P\{Y < 4\}} = \frac{e^{-2} - 3e^{-4}}{1 - 5e^{-4}}.
\]
\end{solution}


\begin{exercise}[求系数、边缘密度与条件密度] \label{exer:3-4b}
二维随机变量 $(X,Y)$ 的联合概率密度为
\[
f(x,y) = \begin{cases} A(6-x-y), & 0 \leqslant x \leqslant 2,\ 2 \leqslant y \leqslant 4, \\ 0, & \text{其他}. \end{cases}
\]
(1) 求系数 $A$；(2) 求 $X, Y$ 的边缘概率密度；(3) $X, Y$ 是否独立？(4) 求条件概率密度 $f_{X|Y}(x|y)$ 和 $f_{Y|X}(y|x)$。
\end{exercise}
\begin{solution}
(1) 由归一性：
\[
\int_2^4\mathrm{d}y\int_0^2 A(6-x-y)\,\mathrm{d}x = A\int_2^4\bigl[6x - \tfrac{x^2}{2} - xy\bigr]_0^2\,\mathrm{d}y = A\int_2^4(10 - 2y)\,\mathrm{d}y = A \cdot 8 = 1,~~A = \frac18
\]
(2) 边缘概率密度：
\[
f_X(x) = \int_2^4 \frac{1}{8}(6-x-y)\,\mathrm{d}y = \frac{3-x}{4}, \quad 0 < x < 2.
\]
\[
f_Y(y) = \int_0^2 \frac{1}{8}(6-x-y)\,\mathrm{d}x = \frac{5-y}{4}, \quad 2 < y < 4.
\]
(3) 取 $x=1, y=3$：$f(1,3) = \frac{1}{8}(6-1-3) = \frac{1}{4}$，而 $f_X(1) \cdot f_Y(3) = \frac{2}{4} \cdot \frac{2}{4} = \frac{1}{4}$。看似相等，但一般点检查不充分。
实际上 $f(x,y) = \dfrac{1}{8}(6-x-y)$，$f_X(x)\,f_Y(y) = \dfrac{(3-x)(5-y)}{16} = \dfrac{15 - 8x - 3y + xy}{16}$，而 $\dfrac{6-x-y}{8} = \dfrac{12 - 2x - 2y}{16}$。两者不相等（例如取 $x=0, y=2$：前者 $= 12/16 = 3/4$，后者 $= 15/16$），故 $X$ 与 $Y$ \textbf{不独立}。\\
(4)
\[
f_{X|Y}(x|y) = \frac{f(x,y)}{f_Y(y)} = \frac{(6-x-y)/8}{(5-y)/4} = \frac{6-x-y}{2(5-y)}, \quad 0 < x < 2.
\]
\[
f_{Y|X}(y|x) = \frac{f(x,y)}{f_X(x)} = \frac{(6-x-y)/8}{(3-x)/4} = \frac{6-x-y}{2(3-x)}, \quad 2 < y < 4.
\]
\end{solution}


\begin{definition}[条件分布函数]\label{def:conditional-df} 
在 $\{X=x\}$ 条件下 $Y$ 的\textbf{条件分布函数}为
\[
F_{Y|X}(y|x) = \int_{-\infty}^{y}f_{Y|X}(v|x)\,\mathrm{d}v = \int_{-\infty}^{y}\frac{f(x,v)}{f_X(x)}\,\mathrm{d}v.
\]
同理，在 $\{Y=y\}$ 条件下 $X$ 的\textbf{条件分布函数}为
\[
F_{X|Y}(x|y) = \int_{-\infty}^{x}f_{X|Y}(u|y)\,\mathrm{d}u = \int_{-\infty}^{x}\frac{f(u,y)}{f_Y(y)}\,\mathrm{d}u.
\]
\end{definition}

\begin{note}
在研究二维连续型随机变量时，通常需要先确定联合密度函数 $f(x,y)$，一般有三种情形：
\begin{enumerate}
    \item 题目中直接给出 $f(x,y)$；
    \item 通过边缘密度函数及随机变量的独立性来确定；
    \item 通过条件密度及边缘密度利用乘法公式求出。
\end{enumerate}
\end{note}


\begin{property}[区域概率公式]\label{prop:region-prob}
对于联合分布函数 $F(x,y)$ 及边缘分布函数 $F_X(x), F_Y(y)$，有
\[
P\{X > x_0, Y > y_0\} = 1 - F_X(x_0) - F_Y(y_0) + F(x_0, y_0).
\]
\textbf{推导}：将全空间分为四个区域，
\[
P\{X > x_0, Y > y_0\} = 1 - P\{X \leqslant x_0\} - P\{Y \leqslant y_0\} + P\{X \leqslant x_0, Y \leqslant y_0\}.
\]
注意：只有当 $X$ 与 $Y$ 独立时，$P\{X > x_0, Y > y_0\} = [1-F_X(x_0)][1-F_Y(y_0)]$。
\end{property}

\begin{exercise}[混合条件分布求 $F_Y$] \label{exer:3-c17}
$P\{X=1\} = P\{X=2\} = 1/2$，在 $X=i$ 的条件下 $Y \sim U(0,i)$。求 $Y$ 的分布函数 $F_Y(y)$ 和密度 $f_Y(y)$。
\end{exercise}
\begin{solution}
由全概率公式：$F_Y(y) = P\{X=1\}\,F_{Y|X=1}(y) + P\{X=2\}\,F_{Y|X=2}(y) = \frac{1}{2}F_1(y) + \frac{1}{2}F_2(y)$。\\
$Y|X=1 \sim U(0,1)$：$F_1(y) = \begin{cases} 0, & y \leqslant 0, \\ y, & 0 < y < 1, \\ 1, & y \geqslant 1. \end{cases}$\\
$Y|X=2 \sim U(0,2)$：$F_2(y) = \begin{cases} 0, & y \leqslant 0, \\ y/2, & 0 < y < 2, \\ 1, & y \geqslant 2. \end{cases}$
分区间计算 $F_Y(y)$：
\begin{itemize}
    \item $y \leqslant 0$：$F_Y(y) = 0$。
    \item $0 < y \leqslant 1$：$F_Y(y) = \dfrac{1}{2}y + \dfrac{1}{2}\cdot\dfrac{y}{2} = \dfrac{3y}{4}$。
    \item $1 < y \leqslant 2$：$F_Y(y) = \dfrac{1}{2}\cdot 1 + \dfrac{1}{2}\cdot\dfrac{y}{2} = \dfrac{1}{2} + \dfrac{y}{4}$。
    \item $y \geqslant 2$：$F_Y(y) = 1$。
\end{itemize}
求导得密度：
\[
F_Y(y) = \begin{cases} 0, & y \leqslant 0, \\ \dfrac{3y}{4}, & 0 < y \leqslant 1, \\ \dfrac{1}{2} + \dfrac{y}{4}, & 1 < y \leqslant 2, \\ 1, & y \geqslant 2. \end{cases}
\qquad
f_Y(y) = \begin{cases} \dfrac{3}{4}, & 0 < y \leqslant 1, \\ \dfrac{1}{4}, & 1 < y \leqslant 2, \\ 0, & \text{其他}. \end{cases}
\]
\end{solution}


\section{常见的二维分布}

\subsection{二维均匀分布}

\begin{definition}[二维均匀分布] \label{def:2d-uniform}
设 $D$ 为 $xOy$ 平面上的有界区域，其面积为 $S_D$。如果 $(X,Y)$ 的联合概率密度为
\[
f(x,y) = \begin{cases} \dfrac{1}{S_D}, & (x,y) \in D, \\ 0, & \text{其他}, \end{cases}
\]
则称 $(X,Y)$ 在区域 $D$ 上服从\textbf{二维均匀分布}，记为 $(X,Y) \sim U(D)$。
\end{definition}

二维均匀分布是区间上均匀分布在二维情形的推广。若 $(X,Y) \sim U(D)$，则对 $D$ 的任何子区域 $G$（其面积为 $S_G$），有
\[
P\{(X,Y) \in G\} = \frac{S_G}{S_D}.
\]


\begin{exercise}[抛物线区域上的均匀分布] \label{exer:3-c7b}
设 $(X,Y)$ 服从区域 $G=\{(x,y): 0 \leqslant y \leqslant 1-x^2,-1<x<1\}$ 上的均匀分布。求：\\
(1) 联合概率密度 $f(x,y)$；\\
(2) 边缘密度 $f_X(x)$，$f_Y(y)$ 及条件密度 $f_{Y|X}(y|x)$；\\
(3) $X$ 与 $Y$ 是否独立？\\
(4) $P\{Y \geqslant X^2\}$。
\end{exercise}
\begin{solution}
(1) 区域 $G$ 的面积：
\[
S_G = \int_{-1}^1(1-x^2)\,\mathrm{d}x = \left[x-\frac{x^3}{3}\right]_{-1}^1 = \frac{4}{3}.
\]
故
\[
f(x,y) = \begin{cases} 3/4, & -1 < x < 1,\ 0 \leqslant y \leqslant 1-x^2, \\ 0, & \text{其他}. \end{cases}
\]
(2) 边缘密度：
\[
f_X(x) = \int_0^{1-x^2}\frac{3}{4}\,\mathrm{d}y = \frac{3}{4}(1-x^2),\quad |x|<1.
\]
\[
f_Y(y) = \int_{-\sqrt{1-y}}^{\sqrt{1-y}}\frac{3}{4}\,\mathrm{d}x = \frac{3}{2}\sqrt{1-y},\quad 0 \leqslant y \leqslant 1.
\]
条件密度（$|x|<1$ 时）：
\[
f_{Y|X}(y|x) = \frac{f(x,y)}{f_X(x)} = \begin{cases} \dfrac{1}{1-x^2}, & 0 \leqslant y \leqslant 1-x^2, \\ 0, & \text{其他}. \end{cases}
\]
(3) $f_X(x)\,f_Y(y) = \dfrac{3}{4}(1-x^2)\cdot\dfrac{3}{2}\sqrt{1-y} = \dfrac{9}{8}(1-x^2)\sqrt{1-y} \neq \dfrac{3}{4} = f(x,y)$，故不独立。\\
(4) $\{Y \geqslant X^2\}$ 要求 $x^2 \leqslant y \leqslant 1-x^2$，即 $2x^2 \leqslant 1$，$|x| \leqslant 1/\sqrt{2}$：
\[
P\{Y \geqslant X^2\} = \int_{-1/\sqrt{2}}^{1/\sqrt{2}}\int_{x^2}^{1-x^2}\frac{3}{4}\,\mathrm{d}y\,\mathrm{d}x = \frac{3}{2}\int_0^{1/\sqrt{2}}(1-2x^2)\,\mathrm{d}x
\]
\[
= \frac{3}{2}\left[x-\frac{2x^3}{3}\right]_0^{1/\sqrt{2}} = \frac{3}{2}\left(\frac{1}{\sqrt{2}}-\frac{2}{3\sqrt{8}}\right) = \frac{3}{2}\cdot\frac{2}{3\sqrt{2}} = \frac{\sqrt{2}}{2}.
\]
\end{solution}


\begin{exercise}[联合均匀分布的判断] \label{exer:3-c12b}
$X,Y$ 独立且均服从 $U[0,1]$，则服从均匀分布的随机变量是（\quad）\\
A. $(X,Y)$ \quad B. $X+Y$ \quad C. $X-Y$ \quad D. $X^2$
\end{exercise}
\begin{solution}
\textbf{题型识别：} 判断哪个随机变量（或向量）服从均匀分布。

\textbf{依据：} 均匀分布的密度在支撑区域上为常数。逐项检查各选项的密度是否为常数。

\textbf{逐项分析：}

$(X,Y)$ 的联合密度 $f(x,y)=1$（$0\leqslant x,y\leqslant 1$），在正方形上为常数，是均匀分布。

$X+Y$ 的密度为三角形分布 $f_Z(z) = z$（$0<z<1$）或 $2-z$（$1\leqslant z<2$），密度不是常数，非均匀。

$X-Y$ 的密度为 $f_W(w) = 1-|w|$（$|w|<1$），密度不是常数，非均匀。

$X^2$ 的密度为 $f(u) = 1/(2\sqrt{u})$（$0<u<1$），即 $\text{Beta}(1/2,1)$ 分布，非均匀。

答案选 \textbf{A}。
\end{solution}


\begin{exercise}[抛物线区域上的均匀分布] \label{exer:3-c11}
\[f(x,y) = \begin{cases} kx^2, & 0 < x < 1,\ x^2 < y < 1, \\ 0, & \text{其他}. \end{cases}\]
求：(1) $k$；(2) $P\{X > 1/2\}$，$P\{Y < 1/2\}$。
\end{exercise}
\begin{solution}
(1) 由归一性：
\[
\int_0^1\int_{x^2}^1 kx^2\,\mathrm{d}y\,\mathrm{d}x = k\int_0^1 x^2(1-x^2)\,\mathrm{d}x = k\left[\frac{x^3}{3} - \frac{x^5}{5}\right]_0^1 = k\cdot\frac{2}{15} = 1,
\]
故 $\boldsymbol{k = 15/2}$。\\
(2)
\[
P\left\{X > \frac{1}{2}\right\} = \frac{15}{2}\int_{1/2}^1\int_{x^2}^1 x^2\,\mathrm{d}y\,\mathrm{d}x = \frac{15}{2}\int_{1/2}^1 x^2(1-x^2)\,\mathrm{d}x
\]
\[
= \frac{15}{2}\left[\frac{x^3}{3} - \frac{x^5}{5}\right]_{1/2}^1 = \frac{15}{2}\cdot\frac{47}{480} = \frac{47}{64}.
\]
\[
P\left\{Y < \frac{1}{2}\right\} = \frac{15}{2}\int_0^{1/\sqrt{2}}\int_{x^2}^{1/2} x^2\,\mathrm{d}y\,\mathrm{d}x = \frac{15}{2}\int_0^{1/\sqrt{2}}x^2\!\left(\frac{1}{2}-x^2\right)\mathrm{d}x
\]
\[
= \frac{15}{2}\left[\frac{x^3}{6} - \frac{x^5}{5}\right]_0^{1/\sqrt{2}} = \frac{15}{2}\cdot\frac{1}{15\sqrt{2}} = \frac{1}{2\sqrt{2}} = \frac{\sqrt{2}}{4}.
\]
\end{solution}


\begin{exercise}[三角形均匀分布求概率] \label{exer:3-c14}
\[f(x,y) = \begin{cases} 2-x-y, & 0 < x < 1,\ 0 < y < 1, \\ 0, & \text{其他}. \end{cases}\]
求：(1) $P\{X > 2Y\}$；(2) $Z = X+Y$ 的密度。
\end{exercise}
\begin{solution}
验证归一性：$\int_0^1\int_0^1(2-x-y)\,\mathrm{d}y\,\mathrm{d}x = 1$。\\
(1) $\{X > 2Y\}$ 等价于 $\{0 < y < x/2\}$（$0 < x < 1$）：
\[
P\{X > 2Y\} = \int_0^1\int_0^{x/2}(2-x-y)\,\mathrm{d}y\,\mathrm{d}x = \int_0^1\left[2y-xy-\frac{y^2}{2}\right]_0^{x/2}\mathrm{d}x
\]
\[
= \int_0^1\left(x - \frac{x^2}{2} - \frac{x^2}{8}\right)\mathrm{d}x = \int_0^1\left(x - \frac{5x^2}{8}\right)\mathrm{d}x = \left[\frac{x^2}{2} - \frac{5x^3}{24}\right]_0^1 = \frac{7}{24}.
\]
(2) 当 $0 < z < 1$ 时（$x \in (0,z)$，$y \in (0, z-x)$）：
\[
f_Z(z) = \int_0^z(2-z)\,\mathrm{d}x = z(2-z).
\]
当 $1 \leqslant z < 2$ 时（$x \in (z-1, 1)$，$y \in (0, z-x)$）：
\[
f_Z(z) = \int_{z-1}^1(2-z)\,\mathrm{d}x = (2-z)(2-z) = (2-z)^2.
\]
综上：
\[
f_Z(z) = \begin{cases} z(2-z), & 0 < z < 1, \\ (2-z)^2, & 1 \leqslant z < 2, \\ 0, & \text{其他}. \end{cases}
\]
\end{solution}


\begin{exercise}[三角形区域上的均匀分布] \label{exer:3-13}
已知二维随机变量 $(X,Y)$ 在区域 $G$ 上服从均匀分布，$G$ 由直线 $x-y=0$、$x+y=2$ 与 $y=0$ 围成。求：\\
(1)边缘概率密度 $f_X(x)$ 和 $f_Y(y)$；~~(2)条件概率密度 $f_{X|Y}(x|y)$。
\end{exercise}
\begin{solution}
$(X,Y)$ 的正概率密度区域 $G$ 是以 $(0,0)$、$(2,0)$、$(1,1)$ 为顶点的三角形，面积 $S_G = \dfrac{1}{2} \times 2 \times 1 = 1$。
由于面积恰好为 $1$，联合概率密度为
\[
f(x,y) = \begin{cases} 1, & 0 \leqslant y \leqslant 1,\ y \leqslant x \leqslant 2-y, \\ 0, & \text{其他}. \end{cases}
\]
(1) 边缘概率密度。$f_X(x)$：固定 $x$，对 $y$ 积分。
\begin{itemize}
    \item 当 $0 < x < 1$ 时，$y$ 从 $0$ 到 $x$：$f_X(x) = \displaystyle\int_0^x 1\,\mathrm{d}y = x$；
    \item 当 $1 < x < 2$ 时，$y$ 从 $0$ 到 $2-x$：$f_X(x) = \displaystyle\int_0^{2-x} 1\,\mathrm{d}y = 2-x$。
\end{itemize}
\[
f_X(x) = \begin{cases} x, & 0 < x < 1, \\ 2-x, & 1 < x < 2, \\ 0, & \text{其他}. \end{cases}
\]
$f_Y(y)$：固定 $y$，对 $x$ 积分。
\begin{itemize}
    \item 当 $0 < y < 1$ 时，$x$ 从 $y$ 到 $2-y$：$f_Y(y) = \displaystyle\int_y^{2-y} 1\,\mathrm{d}x = 2-2y$。
\end{itemize}
\[
f_Y(y) = \begin{cases} 2-2y, & 0 < y < 1, \\ 0, & \text{其他}. \end{cases}
\]
(2) 在 $0 < y < 1$ 的条件下：
\[
f_{X|Y}(x|y) = \frac{f(x,y)}{f_Y(y)} = \frac{1}{2-2y}, \quad y \leqslant x \leqslant 2-y.
\]
即
\[
f_{X|Y}(x|y) = \begin{cases} \dfrac{1}{2-2y}, & y \leqslant x \leqslant 2-y, \\ 0, & \text{其他}. \end{cases}
\]
\end{solution}
\begin{note}
求边缘概率密度的技巧（口诀）：\textbf{求谁不积谁，不积先定限，限内画条线，先交写下限，后交写上限。}即对 $f(x,y)$ 求关于 $y$ 的积分得 $f_X(x)$，对 $f(x,y)$ 求关于 $x$ 的积分得 $f_Y(y)$。定积分限时，画出正概率密度区域的图形，用平行于积分轴的直线穿过区域，先交到的位置为下限，后交到的位置为上限。
\end{note}

\begin{exercise}[条件均匀分布] \label{exer:3-14}
已知随机变量 $X$ 在区间 $[0,1]$ 上服从均匀分布，在 $X=x$（$0<x<1$）的条件下，随机变量 $Y$ 在区间 $[0,x]$ 上服从均匀分布。求：\\
(1) $(X,Y)$ 的联合概率密度；~~(2)$Y$ 的概率密度；~~(3)概率 $P\{X+Y>1\}$。
\end{exercise}
\begin{solution}
(1) 由题设，$X$ 的边缘概率密度为
\[
f_X(x) = \begin{cases} 1, & 0 < x < 1, \\ 0, & \text{其他}. \end{cases}
\]
在 $X=x$（$0<x<1$）的条件下，$Y$ 在 $[0,x]$ 上的条件概率密度为
\[
f_{Y|X}(y|x) = \begin{cases} \dfrac{1}{x}, & 0 < y < x, \\ 0, & \text{其他}. \end{cases}
\]
由乘法公式 $f(x,y) = f_X(x)\,f_{Y|X}(y|x)$，当 $0<y<x<1$ 时 $f(x,y) = 1\cdot\dfrac{1}{x} = \dfrac{1}{x}$。故联合概率密度为
\[
f(x,y) = \begin{cases} \dfrac{1}{x}, & 0 < y < x < 1, \\ 0, & \text{其他}. \end{cases}
\]
(2) 当 $0<y<1$ 时，固定 $y$ 对 $x$ 积分（$x$ 从 $y$ 到 $1$）：
\[
f_Y(y) = \int_y^1\frac{1}{x}\,\mathrm{d}x = \bigl[-\ln y\bigr]_y^1 = -\ln y.
\]
当 $y \leqslant 0$ 或 $y \geqslant 1$ 时，$f_Y(y) = 0$。因此
\[
f_Y(y) = \begin{cases} -\ln y, & 0 < y < 1, \\ 0, & \text{其他}. \end{cases}
\]
(3) $\{X+Y>1\}$ 的积分区域为 $\{(x,y): 0<y<x<1,\ x+y>1\}$，即 $x>\dfrac{1}{2}$ 且 $1-x<y<x$：
\[
P\{X+Y>1\} = \int_{1/2}^1\mathrm{d}x\int_{1-x}^x\frac{1}{x}\,\mathrm{d}y = \int_{1/2}^1\frac{x-(1-x)}{x}\,\mathrm{d}x = \int_{1/2}^1\frac{2x-1}{x}\,\mathrm{d}x = \int_{1/2}^1\left(2-\frac{1}{x}\right)\mathrm{d}x
\]
\[
= \bigl[2x - \ln x\bigr]_{1/2}^1 = (2-0) - (1-\ln 2) = 1 - \ln 2.
\]
\end{solution}
\begin{note}
第 (1) 问中，乘法公式 $f(x,y) = f_X(x)\,f_{Y|X}(y|x)$ 仅在 $f_X(x)>0$ 时成立，因此最终结果需要加上 $f_X(x)=0$ 的部分，在分段表达式的``其他''中自然已包含。
\end{note}


\subsection{二维正态分布}

\begin{definition}[二维正态分布] \label{def:2d-normal}
如果 $(X,Y)$ 的联合概率密度为
\[
f(x,y) = \frac{1}{2\pi\sigma_1\sigma_2\sqrt{1-\rho^2}}\exp\left\{-\frac{1}{2(1-\rho^2)}\left[\frac{(x-\mu_1)^2}{\sigma_1^2} - 2\rho\frac{(x-\mu_1)(y-\mu_2)}{\sigma_1\sigma_2} + \frac{(y-\mu_2)^2}{\sigma_2^2}\right]\right\},
\]
其中 $\mu_1, \mu_2 \in \mathbb{R}$，$\sigma_1 > 0$，$\sigma_2 > 0$，$-1 < \rho < 1$，则称 $(X,Y)$ 服从参数为 $\mu_1, \mu_2, \sigma_1^2, \sigma_2^2, \rho$ 的\textbf{二维正态分布}，记为 $(X,Y) \sim N(\mu_1, \mu_2, \sigma_1^2, \sigma_2^2; \rho)$。
\end{definition}

\begin{property}\label{prop:2d-normal}
二维正态分布具有以下重要性质：
\begin{enumerate}
    \item \textbf{边缘分布仍为正态分布}：若 $(X,Y) \sim N(\mu_1, \mu_2, \sigma_1^2, \sigma_2^2; \rho)$，则
    \[
    X \sim N(\mu_1, \sigma_1^2), \quad Y \sim N(\mu_2, \sigma_2^2).
    \]
    \item \textbf{独立性的充要条件}：若 $(X,Y)$ 服从二维正态分布，则 $X$ 与 $Y$ 相互独立的充分必要条件是 $\rho = 0$。
    \item \textbf{独立正态变量联合分布}：若 $X \sim N(\mu_1, \sigma_1^2)$，$Y \sim N(\mu_2, \sigma_2^2)$，且 $X$ 与 $Y$ 相互独立，则
    \[
    (X,Y) \sim N(\mu_1, \mu_2, \sigma_1^2, \sigma_2^2; 0).
    \]
\end{enumerate}
\end{property}

\begin{note}
需要注意以下几点：
\begin{enumerate}
    \item 边缘分布为正态分布\textbf{不能}推出联合分布为二维正态分布。例如，设 $(X,Y)$ 的联合密度为
    \[
    f(x,y) = \frac{1}{2\pi}e^{-(x^2+y^2)/2}\bigl[1 + \sin x\sin y\bigr],
    \]
    则 $X$ 和 $Y$ 都服从标准正态分布，但 $(X,Y)$ 并不服从二维正态分布。
    \item 性质 (2) 仅对二维正态分布成立。一般地，由 $X$ 与 $Y$ 的边缘分布和 $\rho_{XY} = 0$ 不能推出 $X$ 与 $Y$ 独立。
    \item 二维正态分布的条件分布仍为正态分布：在 $X = x$ 的条件下，$Y$ 的条件分布为
    \[
    Y \mid X=x \sim N\!\left(\mu_2 + \rho\frac{\sigma_2}{\sigma_1}(x-\mu_1),\ \sigma_2^2(1-\rho^2)\right).
    \]
\end{enumerate}
\end{note}

\begin{exercise}[二维分布的性质判断] \label{exer:3-c5b}
下列叙述中正确的是（\quad）
\begin{enumerate}
    \item[(1)] 二维正态分布的边缘分布是正态分布，条件分布也是正态分布；
    \item[(2)] 二维正态分布的边缘分布是正态分布，但条件分布不是正态分布；
    \item[(3)] 二维均匀分布的边缘分布是均匀分布，条件分布是均匀分布；
    \item[(4)] 二维均匀分布的边缘分布不一定是均匀分布，条件分布也不一定为均匀分布。
\end{enumerate}
A. (1)(3) \quad B. (2)(3) \quad C. (1)(4) \quad D. (2)(4)
\end{exercise}
\begin{solution}
\textbf{答案选 C。}
(1) 正确：二维正态分布的边缘分布和条件分布均为正态分布。\\
(4) 正确：二维均匀分布的边缘不一定均匀。例如
\[
f(x,y) = \begin{cases} 2, & 0 < x < 1,\ 0 < y < 1-x, \\ 0, & \text{其他}, \end{cases}
\]
则 $f_X(x) = 2(1-x)$（$0<x<1$），$f_Y(y) = 2(1-y)$（$0<y<1$），均非均匀分布；条件密度 $f_{Y|X}(y|x) = 1/(1-x)$（$0<y<1-x$），在 $[0,1]$ 上也不均匀。
\end{solution}

\begin{exercise}\label{exer:3-5}
设随机变量 $X$ 与 $Y$ 相互独立，且均服从参数为 $1$ 的指数分布，即 $X, Y \sim E(1)$，求 $P\{X < Y\}$。
\end{exercise}

\begin{solution}
由独立性，联合概率密度为
\[
f(x,y) = f_X(x)\,f_Y(y) = \begin{cases} e^{-(x+y)}, & x > 0,\ y > 0, \\ 0, & \text{其他}. \end{cases}
\]
\[
P\{X < Y\} = \iint_{x < y}f(x,y)\,\mathrm{d}x\,\mathrm{d}y = \int_0^{+\infty}\mathrm{d}y\int_0^y e^{-(x+y)}\,\mathrm{d}x = \int_0^{+\infty}e^{-y}(1-e^{-y})\,\mathrm{d}y
\]
\[
= \int_0^{+\infty}(e^{-y} - e^{-2y})\,\mathrm{d}y = \left[-e^{-y} + \frac{1}{2}e^{-2y}\right]_0^{+\infty} = 1 - \frac{1}{2} = \frac{1}{2}.
\]
\end{solution}

\begin{exercise}[条件正态联合密度] \label{exer:3-5b}
设随机变量 $X \sim N(0,1)$，在 $X = x$ 的条件下，随机变量 $Y \sim N(x, 1)$，求 $(X,Y)$ 的概率密度 $f(x,y)$。
\end{exercise}

\begin{solution}
由题意知
\[
f_X(x) = \frac{1}{\sqrt{2\pi}}e^{-x^2/2}, \quad f_{Y|X}(y|x) = \frac{1}{\sqrt{2\pi}}e^{-(y-x)^2/2},
\]
故由乘法公式
\[
f(x,y) = f_X(x)\,f_{Y|X}(y|x) = \frac{1}{2\pi}e^{-\frac{x^2 + (y-x)^2}{2}} = \frac{1}{2\pi}e^{-\frac{2x^2 - 2xy + y^2}{2}}.
\]
\end{solution}


\section{二维随机变量函数的分布}

\begin{note}
\textbf{自学导读：本节解决"已知两个变量的联合分布，求它们的函数的分布"。}

这是第 2 章"一维随机变量函数的分布"的二维推广。核心方法有两种：
\begin{itemize}
    \item \textbf{分布函数法}（万能法）：先写 $F_Z(z)=P\{g(X,Y)\leqslant z\}$，把它转化为对联合密度的二重积分，再对 $z$ 求导得 $f_Z(z)$。
    \item \textbf{卷积公式}（专用于 $Z=X+Y$ 且 $X\perp Y$）：$f_Z(z)=\int_{-\infty}^{+\infty}f_X(t)\,f_Y(z-t)\,\mathrm{d}t$。
\end{itemize}

\textbf{自学提醒：} 做题时最容易出错的地方是\textbf{积分限}。建议每次都画出 $(x,y)$ 平面上的积分区域，标清边界线，再确定积分上下限。
\end{note}

设 $X, Y$ 为随机变量，$g(x,y)$ 是二元函数，则 $Z = g(X,Y)$ 也是随机变量，称为随机变量 $X, Y$ 的函数。本节讨论已知 $(X,Y)$ 的分布，求 $Z = g(X,Y)$ 的分布的方法。

常考问题包括：已知 $(X,Y)$ 的分布，求 $Z = X + Y$、$Z = XY$、$Z = X/Y$、$Z = \max\{X,Y\}$、$Z = \min\{X,Y\}$ 等的分布。

\subsection{离散型随机变量函数的分布}

对于离散型随机变量 $(X,Y)$，函数 $Z = g(X,Y)$ 的分布可以通过\textbf{逐点法}（也称\textbf{表格法}）来确定。

\begin{property}[逐点法]\label{prop:pointwise-method} 
设 $(X,Y)$ 的联合分布律为 $P\{X = x_i, Y = y_j\} = p_{ij}$，则 $Z = g(X,Y)$ 的分布律为
\[
P\{Z = z_k\} = \sum_{g(x_i,y_j) = z_k}p_{ij},
\]
即将使得 $g(x_i, y_j) = z_k$ 的所有 $(i,j)$ 对应的 $p_{ij}$ 求和。若不同的 $(x_i, y_j)$ 对应相同的 $z_k$ 值，则需要合并概率。
\end{property}

\begin{exercise}\label{exer:3-6}
设随机变量 $X$ 和 $Y$ 相互独立，且 $X \sim B(2, p)$，$Y \sim B(4, p)$，求 $Z = X + Y$ 的分布律。
\end{exercise}

\begin{solution}
由于 $X$ 和 $Y$ 相互独立，利用二项分布的可加性（独立二项分布之和仍为二项分布）：
\[
Z = X + Y \sim B(2+4, p) = B(6, p).
\]
即 $P\{Z = k\} = \mathrm{C}_6^k\,p^k(1-p)^{6-k}$，$k = 0, 1, \dots, 6$。

\begin{note}
更一般地，若 $X \sim B(n_1, p)$，$Y \sim B(n_2, p)$，且 $X$ 与 $Y$ 独立，则 $X + Y \sim B(n_1 + n_2, p)$。类似地，若 $X \sim P(\lambda_1)$，$Y \sim P(\lambda_2)$，且 $X$ 与 $Y$ 独立，则 $X + Y \sim P(\lambda_1 + \lambda_2)$。
\end{note}
\end{solution}

\begin{exercise}[和与差的分布] \label{exer:3-6b}
将两封信投入 3 个信箱，设 $X_1, X_2$ 分别表示第一个和第二个信箱投进的信的数量，求随机变量 $Y_1 = X_1 + X_2$ 与 $Y_2 = X_1 - X_2$ 的分布。
\end{exercise}

\begin{solution}
由 \ref{exer:3-2} 知 $(X_1, X_2)$ 的联合分布律。$Y_1 = X_1 + X_2$ 的可能取值为 $0, 1, 2$：
\[
P\{Y_1 = 0\} = p_{00} = \frac{1}{9}, \quad P\{Y_1 = 1\} = p_{01} + p_{10} = \frac{4}{9}, \quad P\{Y_1 = 2\} = p_{02} + p_{11} + p_{20} = \frac{4}{9}.
\]

$Y_2 = X_1 - X_2$ 的可能取值为 $-2, -1, 0, 1, 2$：
\[
P\{Y_2 = -2\} = p_{02} = \frac{1}{9}, \quad P\{Y_2 = -1\} = p_{01} = \frac{2}{9}, \quad P\{Y_2 = 0\} = p_{00} + p_{11} = \frac{3}{9} = \frac{1}{3},
\]
\[
P\{Y_2 = 1\} = p_{10} = \frac{2}{9}, \quad P\{Y_2 = 2\} = p_{20} = \frac{1}{9}.
\]

综上：
\[
Y_1 \sim \begin{pmatrix} 0 & 1 & 2 \\ 1/9 & 4/9 & 4/9 \end{pmatrix}, \qquad Y_2 \sim \begin{pmatrix} -2 & -1 & 0 & 1 & 2 \\ 1/9 & 2/9 & 1/3 & 2/9 & 1/9 \end{pmatrix}.
\]
\end{solution}

\subsection{连续型随机变量函数的分布}

对于连续型随机变量 $(X,Y)$，求 $Z = g(X,Y)$ 的分布主要有\textbf{分布函数法}和\textbf{卷积公式法}两种方法。

如果 $(X,Y)$ 中一个是离散型、一个是连续型，可以将事件对离散型的一切可能值进行\textbf{全概率公式分解}，逐项求分布（见 \ref{exer:3-8d}）。

\begin{note}
\textbf{自学追问：} 本章的学习路线可以概括为"联合分布 $\to$ 边缘分布 $\to$ 条件分布 $\to$ 函数分布"。前三步回答的是"$(X,Y)$ 本身怎么分布"，这一步回答的是"$Z=g(X,Y)$ 怎么分布"。方法选择的口诀：\textbf{万能用分布函数法}（先求 $F_Z(z)$ 再求导）；\textbf{求和用卷积公式}（$Z=X+Y$ 时直接套积分）；\textbf{求极值用独立乘积}（$\max$/$\min$ 拆成各自分布函数的乘积）。

\textbf{易错点：} 分布函数法中积分区域随 $z$ 变化，画图定上下限是关键——积分限写错是此类题最常见的扣分原因。卷积公式要求 $X,Y$ 独立，如果题目没有给独立条件，只能回到分布函数法。
\end{note}

\subsubsection{分布函数法}

分布函数法是最基本的方法，适用于任何函数 $g(X,Y)$。

\begin{property}[分布函数法]\label{prop:df-method}
求 $Z = g(X,Y)$ 的分布，步骤如下：
\begin{enumerate}
    \item 先求 $Z$ 的分布函数：
    \[
    F_Z(z) = P\{Z \leqslant z\} = P\{g(X,Y) \leqslant z\} = \iint_{\{(x,y): g(x,y) \leqslant z\}}f(x,y)\,\mathrm{d}x\,\mathrm{d}y;
    \]
    \item 再对 $F_Z(z)$ 求导得概率密度：$f_Z(z) = F_Z'(z)$。
\end{enumerate}
\end{property}

\begin{exercise}\label{exer:3-7}
设随机变量 $X$ 和 $Y$ 相互独立，$X \sim U(0, 1)$，$Y \sim E(1)$，求 $Z = \max\{X, Y\}$ 的分布函数 $F_Z(z)$。
\end{exercise}

\begin{solution}
$X$ 的分布函数为 $F_X(x) = x$（$0 < x < 1$），$Y$ 的分布函数为 $F_Y(y) = 1 - e^{-y}$（$y > 0$）。

对于独立的 $X$ 和 $Y$：
\[
\max\{X, Y\} \leqslant z \iff X \leqslant z \text{ 且 } Y \leqslant z,
\]
故
\[
F_Z(z) = P\{Z \leqslant z\} = P\{X \leqslant z\}\,P\{Y \leqslant z\} = F_X(z)\,F_Y(z).
\]

当 $z \leqslant 0$ 时，$F_Z(z) = 0$；

当 $0 < z < 1$ 时，$F_Z(z) = z(1 - e^{-z})$；

当 $z \geqslant 1$ 时，$F_Z(z) = 1 \cdot (1-e^{-z}) = 1 - e^{-z}$。

综上：
\[
F_Z(z) = \begin{cases} 0, & z \leqslant 0, \\ z(1 - e^{-z}), & 0 < z < 1, \\ 1 - e^{-z}, & z \geqslant 1. \end{cases}
\]
\end{solution}

\begin{exercise}\label{exer:3-8}
设 $(X,Y)$ 的联合概率密度为
\[
f(x,y) = \begin{cases} e^{-(x+y)}, & x > 0,\ y > 0, \\ 0, & \text{其他}, \end{cases}
\]
求 $Z = X + Y$ 的概率密度。
\end{exercise}

\begin{solution}
用分布函数法。当 $z > 0$ 时：
\[
F_Z(z) = P\{X+Y \leqslant z\} = \iint_{x+y \leqslant z}e^{-(x+y)}\,\mathrm{d}x\,\mathrm{d}y = \int_0^z\mathrm{d}x\int_0^{z-x}e^{-(x+y)}\,\mathrm{d}y.
\]

内层积分：
\[
\int_0^{z-x}e^{-(x+y)}\,\mathrm{d}y = e^{-x}(1-e^{-(z-x)}) = e^{-x} - e^{-z}.
\]

外层积分：
\[
F_Z(z) = \int_0^z(e^{-x} - e^{-z})\,\mathrm{d}x = 1 - e^{-z} - ze^{-z} = 1 - (1+z)e^{-z}.
\]

当 $z \leqslant 0$ 时，$F_Z(z) = 0$。求导得：
\[
f_Z(z) = \begin{cases} ze^{-z}, & z > 0, \\ 0, & z \leqslant 0. \end{cases}
\]
即 $Z \sim \Gamma(2, 1)$（形状参数为2、尺度参数为1的Gamma分布）。
\end{solution}

\begin{exercise}[$Z=AX+BY$型:] 已知 $X, Y$ 独立，同时满足参数为2的指数分布$\text{Exp}(2)$，求$Z=2X+3Y$的分布密度函数。
\end{exercise}
\begin{solution}
\[f_X(x)=\begin{cases}2\mathrm{e}^{-2x},&x\geqslant0\\0,&x<0\end{cases},~f_Y(y)=\begin{cases}2\mathrm{e}^{-2y},&y\geqslant0\\0,&y<0\end{cases}\]
分布函数为：\[F_T(t) =\begin{cases}1 - e^{-2t}, & t>0 \\0, & t\le 0\end{cases}\quad (T=X,Y)\]
下求$Z$的分布函数，首先，$F_{Z}(z)=P(Z\leqslant z)=P(2X+3Y\leqslant z)$
分两种情况讨论：
\begin{enumerate}
    \item $z<0$时，$F_{Z}(z)=0$
    \item $z\geqslant 0$时，首先搞清楚积分区域为 \( x>0,y>0, 2x+3y \le z \)，即 \( 0 < x < \frac{z}{2}, 0 < y < \frac{z-2x}{3} \)。
\begin{align*}
    F_{Z}(z)
    &= \iint_{2x+3y\leqslant z} f_X(x) f_Y(y) \,\mathrm{d}x\,\mathrm{d}y
    = \int_{0}^{\frac{z}{2}} \int_{0}^{\frac{z-2x}{3}} 2e^{-2x} \cdot 2e^{-2y}\,\mathrm{d}y\,\mathrm{d}x \\
    &= 4 \int_{0}^{\frac{z}{2}} e^{-2x} \left( \int_{0}^{\frac{z-2x}{3}} e^{-2y} \mathrm{d}y \right) \mathrm{d}x
    = 4 \int_{0}^{\frac{z}{2}} e^{-2x} \left( \left. -\frac{1}{2}e^{-2y} \right|_{0}^{\frac{z-2x}{3}} \right) \mathrm{d}x \\
    &= 4 \int_{0}^{\frac{z}{2}} e^{-2x} \cdot \frac{1}{2}\left( 1 - e^{-\frac{2(z-2x)}{3}} \right) \mathrm{d}x
    = 2 \int_{0}^{\frac{z}{2}} \left( e^{-2x} - e^{-\frac{2z}{3} - \frac{2x}{3}} \right) \mathrm{d}x \\
    &= 2\left[ \int_{0}^{\frac{z}{2}} e^{-2x} \mathrm{d}x - e^{-\frac{2z}{3}} \int_{0}^{\frac{z}{2}} e^{-\frac{2x}{3}} \mathrm{d}x \right] = 2\left[ \left. \left( -\frac{1}{2}e^{-2x} \right) \right|_{0}^{\frac{z}{2}} - e^{-\frac{2z}{3}} \cdot \left. \left( -\frac{3}{2}e^{-\frac{2x}{3}} \right) \right|_{0}^{\frac{z}{2}} \right] \\
    &= 2\left[ \frac{1}{2}\left( 1 - e^{-z} \right) - e^{-\frac{2z}{3}} \cdot \frac{3}{2}\left( 1 - e^{-\frac{z}{3}} \right) \right] = (1 - e^{-z}) - 3e^{-\frac{2z}{3}}\left( 1 - e^{-\frac{z}{3}} \right) \\
    &= 1 - e^{-z} - 3e^{-\frac{2z}{3}} + 3e^{-z} = 1 - 3e^{-\frac{2z}{3}} + 2e^{-z}
\end{align*}
求导：
\begin{align*}
    f_Z(z) &= \frac{\mathrm{d}}{\mathrm{d}z}F_Z(z) = \frac{\mathrm{d}}{\mathrm{d}z}\left( 1 - 3e^{-\frac{2z}{3}} + 2e^{-z} \right) = -3 \cdot \left( -\frac{2}{3} \right)e^{-\frac{2z}{3}} + 2 \cdot (-1)e^{-z} \\
    &= 2e^{-\frac{2z}{3}} - 2e^{-z} = 2\left( e^{-\frac{2z}{3}} - e^{-z} \right), \quad z>0
\end{align*}
\end{enumerate}
\( Z = 2X + 3Y \) 的分布密度函数为：
\[f_Z(z) =\begin{cases}2\left( e^{-\frac{2z}{3}} - e^{-z} \right), & z > 0 \\0, & z \le 0\end{cases}\]
\end{solution}



\begin{exercise}[Gamma型联合密度] \label{exer:3-c9b}
$(X,Y)$ 的概率密度为
\[
f(x,y) = \begin{cases} \dfrac{1}{2}(x+y)e^{-(x+y)}, & x>0,\ y>0, \\ 0, & \text{其他}. \end{cases}
\]
求：(1) 边缘密度；(2) 独立性；(3) $Z=X+Y$ 的密度。
\end{exercise}
\begin{solution}
(1) 利用 $\displaystyle\int_0^{+\infty}ye^{-y}\,\mathrm{d}y = 1$ 和 $\displaystyle\int_0^{+\infty}y^2 e^{-y}\,\mathrm{d}y = 2$：
\[
f_X(x) = \int_0^{+\infty} \frac{1}{2}(x+y)e^{-(x+y)} \,\mathrm{d}y = \frac{e^{-x}}{2}\left(x\int_0^{+\infty}e^{-y}\,\mathrm{d}y+\int_0^{+\infty}ye^{-y}\,\mathrm{d}y\right) = \frac{x+1}{2}e^{-x},\quad x>0.
\]
同理可得边缘密度：\[f_Y(y) = \frac{y+1}{2}e^{-y},\quad y>0.\]

(2) 由于 $f_X(x)\,f_Y(y) = \dfrac{(x+1)(y+1)}{4}e^{-(x+y)} \neq \dfrac{x+y}{2}e^{-(x+y)} = f(x,y)$，故 $X$ 与 $Y$ 不独立。

(3) 下求 $Z=X+Y$ 的分布函数 $F_Z(z) = P(X+Y \leqslant z)$。
分两种情况讨论：
\begin{enumerate}
    \item $z \leqslant 0$ 时，$F_Z(z) = 0$
    \item $z > 0$ 时，积分区域为 $x>0, y>0, x+y \leqslant z$，即 $0 < x < z, 0 < y < z-x$。
\begin{align*}
    F_Z(z) &= \iint_{x+y\leqslant z} f(x,y) \,\mathrm{d}x\,\mathrm{d}y = \int_0^z \int_0^{z-x} \frac{1}{2}(x+y)e^{-(x+y)} \,\mathrm{d}y\,\mathrm{d}x \\
    &= \frac{1}{2} \int_0^z \left( \int_x^z u e^{-u} \,\mathrm{d}u \right) \,\mathrm{d}x \quad \text{(令 $u=x+y$)} \\
    &= \frac{1}{2} \int_0^z \left( \left. -(u+1)e^{-u} \right|_x^z \right) \,\mathrm{d}x \\
    &= \frac{1}{2} \int_0^z \left( (x+1)e^{-x} - (z+1)e^{-z} \right) \,\mathrm{d}x \\
    &= \frac{1}{2} \left[ \left. -(x+2)e^{-x} \right|_0^z - (z+1)e^{-z} \cdot z \right] \\
    &= \frac{1}{2} \left[ 2 - (z+2)e^{-z} - (z^2+z)e^{-z} \right] \\
    &= 1 - \left( 1 + z + \frac{z^2}{2} \right) e^{-z}
\end{align*}
求导得 $Z$ 的概率密度：
\begin{align*}
    f_Z(z) &= \frac{\mathrm{d}}{\mathrm{d}z} F_Z(z) = - \left( 1 + z \right) e^{-z} + \left( 1 + z + \frac{z^2}{2} \right) e^{-z} = \frac{z^2}{2}e^{-z}, \quad z>0
\end{align*}
\end{enumerate}
$Z = X + Y$ 的分布密度函数为：
\[
f_Z(z) = \begin{cases} \dfrac{z^2}{2}e^{-z}, & z > 0 \\ 0, & z \leqslant 0 \end{cases}
\]
\begin{note}
验证：$\displaystyle\int_0^{+\infty}\frac{z^2}{2}e^{-z}\,\mathrm{d}z = \frac{1}{2}\Gamma(3) = 1$。即 $Z \sim \Gamma(3,1)$。
\end{note}
\end{solution}


\begin{note}
\textbf{难点拆解：分布函数法中如何确定积分限（画图法）。}

当 $X$ 和 $Y$ 的支撑集有限（如 $0<x<a$，$0<y<b$）时，积分区域 $\{(x,y): x+y\leqslant z\}$ 与支撑集的交集会随 $z$ 变化。操作步骤：
\begin{enumerate}
    \item 在 $(x,y)$ 平面上画出支撑集（通常是矩形或三角形）；
    \item 画直线 $x+y=z$（斜率 $-1$，截距 $z$）；
    \item 随着 $z$ 从小到大移动这条直线，观察它与支撑集的交集形状何时发生变化——每个"变化点"就是一个分段边界；
    \item 在每一段内，分别写出积分上下限。
\end{enumerate}
\textbf{典型分段：} 若 $X\in(0,\infty)$，$Y\in(0,1)$，则 $z=1$ 是分段点——当 $0<z<1$ 时直线只切到 $Y$ 支撑集的一部分；当 $z\geqslant 1$ 时 $Y$ 的整个范围都被覆盖。
\end{note}

\begin{exercise}[卷积公式分段积分] \label{exer:3-8b}
设随机变量 $X, Y$ 相互独立，其概率密度分别为
\[
f_X(x) = \begin{cases} 2e^{-2x}, & x > 0, \\ 0, & \text{其他}, \end{cases} \quad f_Y(y) = \begin{cases} 3y^2, & 0 < y < 1, \\ 0, & \text{其他}. \end{cases}
\]
(1) 求 $(X,Y)$ 的概率密度；(2) 求 $Z = X + Y$ 的概率密度。
\end{exercise}

\begin{solution}
(1) 由独立性，联合概率密度为：
\[
f(x,y) = f_X(x)\,f_Y(y) = \begin{cases} 6y^2 e^{-2x}, & x > 0,\ 0 < y < 1 \\ 0, & \text{其他} \end{cases}
\]

(2) 下求 $Z=X+Y$ 的分布函数 $F_Z(z) = P(X+Y \leqslant z)$。
分三种情况讨论：
\begin{enumerate}
    \item $z \leqslant 0$ 时，$F_Z(z) = 0$
    \item $0 < z < 1$ 时，积分区域为 $x>0, 0<y<1, x+y \leqslant z$，即 $0 < y < z, 0 < x < z-y$。
\begin{align*}
    F_Z(z) &= \iint_{x+y\leqslant z} f(x,y) \,\mathrm{d}x\,\mathrm{d}y = \int_0^z \int_0^{z-y} 6y^2 e^{-2x} \,\mathrm{d}x\,\mathrm{d}y \\
    &= 3 \int_0^z y^2 \left( \left. -e^{-2x} \right|_0^{z-y} \right) \,\mathrm{d}y = 3 \int_0^z y^2 \left( 1 - e^{-2(z-y)} \right) \,\mathrm{d}y \\
    &= 3 \int_0^z y^2 \,\mathrm{d}y - 3e^{-2z} \int_0^z y^2 e^{2y} \,\mathrm{d}y \\
    &= z^3 - 3e^{-2z} \left[ \left( \frac{1}{2}y^2 - \frac{1}{2}y + \frac{1}{4} \right) e^{2y} \right]_0^z \\
    &= z^3 - 3e^{-2z} \left( \left( \frac{1}{2}z^2 - \frac{1}{2}z + \frac{1}{4} \right) e^{2z} - \frac{1}{4} \right) \\
    &= z^3 - \frac{3}{2}z^2 + \frac{3}{2}z - \frac{3}{4} + \frac{3}{4}e^{-2z}
\end{align*}
求导：
\begin{align*}
    f_Z(z) &= \frac{\mathrm{d}}{\mathrm{d}z} F_Z(z) = 3z^2 - 3z + \frac{3}{2} - \frac{3}{2}e^{-2z}
\end{align*}

    \item $z \geqslant 1$ 时，积分区域为 $x>0, 0<y<1, x+y \leqslant z$，即 $0 < y < 1, 0 < x < z-y$。
\begin{align*}
    F_Z(z) &= \iint_{x+y\leqslant z} f(x,y) \,\mathrm{d}x\,\mathrm{d}y = \int_0^1 \int_0^{z-y} 6y^2 e^{-2x} \,\mathrm{d}x\,\mathrm{d}y \\
    &= 3 \int_0^1 y^2 \left( 1 - e^{-2(z-y)} \right) \,\mathrm{d}y = 3 \int_0^1 y^2 \,\mathrm{d}y - 3e^{-2z} \int_0^1 y^2 e^{2y} \,\mathrm{d}y \\
    &= 1 - 3e^{-2z} \left[ \left( \frac{1}{2}y^2 - \frac{1}{2}y + \frac{1}{4} \right) e^{2y} \right]_0^1 \\
    &= 1 - 3e^{-2z} \left( \frac{1}{4}e^2 - \frac{1}{4} \right) = 1 - \frac{3(e^2-1)}{4}e^{-2z}
\end{align*}
求导：
\begin{align*}
    f_Z(z) &= \frac{\mathrm{d}}{\mathrm{d}z} F_Z(z) = \frac{3(e^2-1)}{2}e^{-2z}
\end{align*}
\end{enumerate}
综上，$Z = X + Y$ 的分布密度函数为：
\[
f_Z(z) = \begin{cases} 3z^2 - 3z + \dfrac{3}{2} - \dfrac{3}{2}e^{-2z}, & 0 < z < 1 \\ \dfrac{3(e^2-1)}{2}e^{-2z}, & z \geqslant 1 \\ 0, & \text{其他} \end{cases}
\]
\end{solution}

\begin{exercise}[积的分布：矩形面积] \label{exer:3-8c}
设二维随机变量 $(X,Y)$ 在矩形区域 $D = \{(x,y) : 0 \leqslant x \leqslant 2,\ 0 \leqslant y \leqslant 1\}$ 上服从均匀分布，求边长为 $X$ 和 $Y$ 的矩形面积 $Z = XY$ 的概率密度。
\end{exercise}

\begin{solution}
由题设，联合概率密度为：
\[
f(x,y) = \begin{cases} \dfrac{1}{2}, & 0 \leqslant x \leqslant 2,\ 0 \leqslant y \leqslant 1 \\ 0, & \text{其他} \end{cases}
\]
下求 $Z=XY$ 的分布函数 $F_Z(z) = P(XY \leqslant z)$。
分三种情况讨论：
\begin{enumerate}
    \item $z \leqslant 0$ 时，$F_Z(z) = 0$
    \item $0 < z < 2$ 时，积分区域为矩形 $D$ 内满足 $xy \leqslant z$ 的部分。边界曲线为 $y = \frac{z}{x}$，与 $y=1$ 交于 $x=z$。我们将积分区域按 $x$ 划分：
\begin{align*}
    F_Z(z) &= \iint_{xy \leqslant z} f(x,y) \,\mathrm{d}x\,\mathrm{d}y = \int_0^z \left( \int_0^1 \frac{1}{2} \,\mathrm{d}y \right) \mathrm{d}x + \int_z^2 \left( \int_0^{\frac{z}{x}} \frac{1}{2} \,\mathrm{d}y \right) \mathrm{d}x \\
    &= \frac{1}{2} \int_0^z 1 \,\mathrm{d}x + \frac{1}{2} \int_z^2 \frac{z}{x} \,\mathrm{d}x \\
    &= \frac{1}{2}z + \frac{z}{2} \left( \left. \ln x \right|_z^2 \right) = \frac{z}{2} + \frac{z}{2} (\ln 2 - \ln z) \\
    &= \frac{z}{2}(1 + \ln 2 - \ln z)
\end{align*}
求导：
\begin{align*}
    f_Z(z) &= \frac{\mathrm{d}}{\mathrm{d}z} F_Z(z) = \frac{\mathrm{d}}{\mathrm{d}z} \left[ \frac{z}{2}(1 + \ln 2 - \ln z) \right] \\
    &= \frac{1}{2}(1 + \ln 2 - \ln z) + \frac{z}{2} \left( -\frac{1}{z} \right) = \frac{1}{2}(\ln 2 - \ln z)
\end{align*}
    \item $z \geqslant 2$ 时，矩形区域完全包含在 $xy \leqslant z$ 内，$F_Z(z) = 1$。求导得 $f_Z(z) = 0$
\end{enumerate}
综上，$Z = XY$ 的分布密度函数为：
\[
f_Z(z) = \begin{cases} \dfrac{1}{2}(\ln 2 - \ln z), & 0 < z < 2 \\ 0, & \text{其他} \end{cases}
\]
\end{solution}

\begin{exercise}[混合型：全概率公式] \label{exer:3-8d}
设随机变量 $X_1, X_2, X_3$ 相互独立，其中 $X_1$ 与 $X_2$ 均服从标准正态分布，$X_3$ 的概率分布为 $P\{X_3 = 0\} = P\{X_3 = 1\} = 1/2$。令  $Y = X_1 X_3 + X_2 (1-X_3)$。(1) 求 $(X_1, Y)$ 的分布函数（结果用 $\Phi(x)$ 表示）；(2) 证明 $Y$ 服从标准正态分布。
\end{exercise}
\begin{solution}
(1) 记 $(X_1, Y)$ 的联合分布函数为 $F(x,y) = P(X_1 \leqslant x, Y \leqslant y)$。
由全概率公式，按 $X_3$ 的取值进行分解：
\begin{align*}
    F(x,y) &= P(X_3 = 0) P(X_1 \leqslant x, Y \leqslant y \mid X_3 = 0) + P(X_3 = 1) P(X_1 \leqslant x, Y \leqslant y \mid X_3 = 1) \\
    &= \frac{1}{2} P(X_1 \leqslant x, X_2 \leqslant y \mid X_3 = 0) + \frac{1}{2} P(X_1 \leqslant x, X_1 \leqslant y \mid X_3 = 1)
\end{align*}
由于 $X_1, X_2, X_3$ 相互独立，条件概率与 $X_3$ 无关，可直接转化为无条件概率：
\begin{align*}
    F(x,y) &= \frac{1}{2} P(X_1 \leqslant x, X_2 \leqslant y) + \frac{1}{2} P(X_1 \leqslant x, X_1 \leqslant y) \\
    &= \frac{1}{2} P(X_1 \leqslant x)P(X_2 \leqslant y) + \frac{1}{2} P(X_1 \leqslant \min(x,y)) \\
    &= \frac{1}{2} \Phi(x)\Phi(y) + \frac{1}{2} \Phi(\min(x,y))
\end{align*}
(2) 下求 $Y$ 的边缘分布函数 $F_Y(y)$：
\begin{align*}
    F_Y(y) &= \lim_{x \to +\infty} F(x,y) = \lim_{x \to +\infty} \left[ \frac{1}{2} \Phi(x)\Phi(y) + \frac{1}{2} \Phi(\min(x,y)) \right] \\
    &= \frac{1}{2} \cdot 1 \cdot \Phi(y) + \frac{1}{2} \Phi(y) \\
    &= \Phi(y)
\end{align*}
由于 $F_Y(y) = \Phi(y)$，因此严格证明了 $Y$ 服从标准正态分布 $N(0,1)$。
\end{solution}

\subsubsection{卷积公式与常用分布函数公式}

\begin{theorem}[卷积公式] \label{thm:convolution}
设 $(X,Y) \sim f(x,y)$，则 $Z = X + Y$ 的概率密度为
\[
f_Z(z) = \int_{-\infty}^{+\infty}f(x, z-x)\,\mathrm{d}x = \int_{-\infty}^{+\infty}f(z-y, y)\,\mathrm{d}y.
\]
当 $X$ 与 $Y$ 相互独立时，
\[
f_Z(z) = \int_{-\infty}^{+\infty}f_X(x)\,f_Y(z-x)\,\mathrm{d}x = \int_{-\infty}^{+\infty}f_X(z-y)\,f_Y(y)\,\mathrm{d}y.
\]
\end{theorem}
\begin{proof}
由分布函数法，\(F_Z(z) = P\{X+Y \leqslant z\} = \iint_{x+y \leqslant z} f(x,y)\,\mathrm{d}x\,\mathrm{d}y\)。对内层积分做变量代换 \(u = x\)，\(y = v - u\)（即固定 \(u\)，令 \(v = u+y\)），则
\[
F_Z(z) = \int_{-\infty}^{z}\left[\int_{-\infty}^{+\infty} f(u, v-u)\,\mathrm{d}u\right]\mathrm{d}v.
\]
对 \(z\) 求导即得 \(f_Z(z) = \int_{-\infty}^{+\infty} f(x, z-x)\,\mathrm{d}x\)。独立时 \(f(x,y)=f_X(x)f_Y(y)\)，即得卷积公式。
\end{proof}

\begin{note}
\textbf{直觉理解卷积公式：} 要求 $Z=X+Y$ 在 $z$ 处的密度，就是问"所有使得 $X+Y=z$ 的 $(x,y)$ 对，它们的联合密度加起来是多少"。满足 $x+y=z$ 的点构成一条直线 $y=z-x$，所以你沿着这条直线"扫描"所有 $x$ 值，把每个点的联合密度 $f(x,z-x)$ 积起来。

\textbf{实际操作要点：}
\begin{enumerate}
    \item 写出 $f_X(x)$ 和 $f_Y(z-x)$ 的表达式；
    \item 确定积分限：$f_X(x)\neq 0$ 且 $f_Y(z-x)\neq 0$ 的 $x$ 范围（这一步最容易出错！）；
    \item 积分限通常随 $z$ 的取值范围而变化，所以往往需要\textbf{分段讨论}。
\end{enumerate}
\end{note}

类似地，对于其他形式的函数，有如下公式：

\begin{property}[差、积、商的分布]\label{prop:diff-prod-quot-dist}
设 $(X,Y) \sim f(x,y)$，则
\begin{enumerate}
    \item \textbf{差的分布}（$Z = X - Y$）：
    \[
    f_Z(z) = \int_{-\infty}^{+\infty}f(x, x-z)\,\mathrm{d}x = \int_{-\infty}^{+\infty}f(z+y, y)\,\mathrm{d}y.
    \]
    当 $X$ 与 $Y$ 独立时，$f_Z(z) = \displaystyle\int_{-\infty}^{+\infty}f_X(x)\,f_Y(x-z)\,\mathrm{d}x$。

    \item \textbf{积的分布}（$Z = XY$）：
    \[
    f_Z(z) = \int_{-\infty}^{+\infty}\frac{1}{|x|}\,f\!\left(x, \frac{z}{x}\right)\,\mathrm{d}x = \int_{-\infty}^{+\infty}\frac{1}{|y|}\,f\!\left(\frac{z}{y}, y\right)\,\mathrm{d}y.
    \]
    当 $X$ 与 $Y$ 独立时，$f_Z(z) = \displaystyle\int_{-\infty}^{+\infty}\frac{1}{|x|}\,f_X(x)\,f_Y\!\left(\frac{z}{x}\right)\,\mathrm{d}x$。

    \item \textbf{商的分布}（$Z = X / Y$）：
    \[
    f_Z(z) = \int_{-\infty}^{+\infty}|y|\,f(yz, y)\,\mathrm{d}y.
    \]
    当 $X$ 与 $Y$ 独立时，$f_Z(z) = \displaystyle\int_{-\infty}^{+\infty}|y|\,f_X(yz)\,f_Y(y)\,\mathrm{d}y$。
\end{enumerate}
\end{property}

\begin{note}
上述卷积公式可用\textbf{口诀}记忆：\textbf{积谁不换谁，换完求偏导}。

例如，$Z = X - Y$ 中对 $x$ 积分（``积谁''= 积 $x$），则 $x$ 不换，$y$ 换为 $y = x - z$，写成 $f(x, x-z)$；然后``换完求偏导''：$\dfrac{\partial(x-z)}{\partial z} = -1$，取绝对值为 $|-1| = 1$，即得公式 $f_Z(z) = \displaystyle\int_{-\infty}^{+\infty}f(x, x-z)\,\mathrm{d}x$。

更一般地，若 $Z = g(X,Y)$，由 $u = g(x,y)$ 解出 $y = h(x,u)$（要求存在且可导），则
\[
f_Z(z) = \int_{-\infty}^{+\infty}f(x, h(x,z))\left|\frac{\partial h(x,z)}{\partial z}\right|\,\mathrm{d}x.
\]
\end{note}

\begin{proof}[差的分布 $Z=X-Y$ 的证明]
\begin{enumerate}
    \item \textbf{分布函数法}：
    先求 $Z$ 的分布函数 $F_Z(z) = P(X-Y \leqslant z) = P(X \leqslant z+Y)$。
    由联合密度积分可得：
    \[
    F_Z(z) = \iint_{x-y \leqslant z} f(x,y)\,\mathrm{d}x\,\mathrm{d}y = \int_{-\infty}^{+\infty} \left( \int_{-\infty}^{z+y} f(x,y)\,\mathrm{d}x \right) \mathrm{d}y.
    \]
    对内层积分做变量代换，令 $u = x-y$（即固定 $y$，有 $x = u+y, \mathrm{d}x = \mathrm{d}u$），当 $x$ 从 $-\infty$ 积到 $z+y$ 时，$u$ 从 $-\infty$ 积到 $z$：
    \[
    F_Z(z) = \int_{-\infty}^{+\infty} \left( \int_{-\infty}^{z} f(u+y, y)\,\mathrm{d}u \right) \mathrm{d}y = \int_{-\infty}^{z} \left[ \int_{-\infty}^{+\infty} f(u+y, y)\,\mathrm{d}y \right] \mathrm{d}u.
    \]
    两边对 $z$ 求导，即得 $f_Z(z) = \displaystyle\int_{-\infty}^{+\infty}f(z+y, y)\,\mathrm{d}y$。
    同理，若按 $\mathrm{d}x$ 在外层积分，可证得另一种形式 $f_Z(z) = \displaystyle\int_{-\infty}^{+\infty}f(x, x-z)\,\mathrm{d}x$。当 $X, Y$ 独立时，$f(x,y)=f_X(x)f_Y(y)$，代入即得结论。

    \item \textbf{卷积公式法（利用和的分布变形）}：
    令 $Y' = -Y$，则 $(X, Y')$ 的联合密度为 $g(x, y') = f(x, -y')$。
    由于 $Z = X - Y = X + Y'$，直接应用定理 \ref{thm:convolution} 中和的卷积公式：
    \[
    f_Z(z) = \int_{-\infty}^{+\infty} g(x, z-x)\,\mathrm{d}x = \int_{-\infty}^{+\infty} f(x, -(z-x))\,\mathrm{d}x = \int_{-\infty}^{+\infty} f(x, x-z)\,\mathrm{d}x.
    \]
    若 $X, Y$ 独立，则 $X, Y'$ 亦独立，此时 $g(x, z-x) = f_X(x)f_{Y'}(z-x) = f_X(x)f_Y(x-z)$，得证。
\end{enumerate}
\end{proof}


\begin{proof}[积的分布 $Z=XY$ 的证明]
\begin{enumerate}
    \item \textbf{分布函数法}：
    求 $Z$ 的分布函数 $F_Z(z) = P(XY \leqslant z)$，积分区域按 $x$ 的正负分为两部分：
    \begin{align*}
        F_Z(z) &= \iint_{xy \leqslant z} f(x,y)\,\mathrm{d}x\,\mathrm{d}y \\
        &= \int_{0}^{+\infty} \left( \int_{-\infty}^{\frac{z}{x}} f(x,y)\,\mathrm{d}y \right) \mathrm{d}x + \int_{-\infty}^{0} \left( \int_{\frac{z}{x}}^{+\infty} f(x,y)\,\mathrm{d}y \right) \mathrm{d}x.
    \end{align*}
    对内层积分做代换，令 $u = xy$（固定 $x$，则 $y = \frac{u}{x}$，$\mathrm{d}y = \frac{1}{x}\mathrm{d}u$）。
    当 $x > 0$ 时，$y$ 的区间 $(-\infty, \frac{z}{x}]$ 对应 $u$ 的区间 $(-\infty, z]$；
    当 $x < 0$ 时，$y$ 的区间 $[\frac{z}{x}, +\infty)$ 对应 $u$ 的区间 $[z, -\infty)$，此时积分限反转带来负号，即 $\mathrm{d}y = \frac{1}{x}\mathrm{d}u = -\frac{1}{|x|}\mathrm{d}u$，吸收反转积分限的负号后，两部分合并为：
    \begin{align*}
    F_Z(z) &= \int_{0}^{+\infty} \left( \int_{-\infty}^{z} f\!\left(x,\frac{u}{x}\right)\frac{1}{|x|}\,\mathrm{d}u \right) \mathrm{d}x + \int_{-\infty}^{0} \left( \int_{-\infty}^{z} f\!\left(x,\frac{u}{x}\right)\frac{1}{|x|}\,\mathrm{d}u \right) \mathrm{d}x \\
    &= \int_{-\infty}^{z} \left[ \int_{-\infty}^{+\infty} \frac{1}{|x|}f\!\left(x,\frac{u}{x}\right)\,\mathrm{d}x \right] \mathrm{d}u.
    \end{align*}
    对 $z$ 求导即得 $f_Z(z) = \displaystyle\int_{-\infty}^{+\infty}\frac{1}{|x|}\,f\!\left(x, \frac{z}{x}\right)\,\mathrm{d}x$。独立时展开联合密度即得特例。

    \item \textbf{二维变量代换法（雅可比行列式法）}：
    作变换 $Z = XY, W = X$。该变换的逆变换为 $x = w, y = \frac{z}{w}$。
    雅可比行列式为 $J = \frac{\partial(x,y)}{\partial(z,w)} = \begin{vmatrix} 0 & 1 \\ \frac{1}{w} & -\frac{z}{w^2} \end{vmatrix} = -\frac{1}{w}$，故 $|J| = \frac{1}{|w|}$。
    因此 $(Z,W)$ 的联合概率密度为 $g(z,w) = f\!\left(w, \frac{z}{w}\right) |J| = \frac{1}{|w|}f\!\left(w, \frac{z}{w}\right)$。
    求边缘密度，将 $w$ 积分掉（用 $x$ 替换积分变量 $w$）：
    \[
    f_Z(z) = \int_{-\infty}^{+\infty} g(z,w)\,\mathrm{d}w = \int_{-\infty}^{+\infty} \frac{1}{|x|}f\!\left(x, \frac{z}{x}\right)\,\mathrm{d}x.
    \]
\end{enumerate}
\end{proof}


\begin{proof}[商的分布 $Z=X/Y$ 的证明]
\begin{enumerate}
    \item \textbf{分布函数法}：
    求 $F_Z(z) = P(X/Y \leqslant z)$，按 $y$ 的正负分为两部分：
    \begin{align*}
        F_Z(z) &= P(X \leqslant zY, Y > 0) + P(X \geqslant zY, Y < 0) \\
        &= \int_{0}^{+\infty} \left( \int_{-\infty}^{zy} f(x,y)\,\mathrm{d}x \right) \mathrm{d}y + \int_{-\infty}^{0} \left( \int_{zy}^{+\infty} f(x,y)\,\mathrm{d}x \right) \mathrm{d}y.
    \end{align*}
    对内层积分做代换，令 $u = x/y$（固定 $y$，则 $x = uy$，$\mathrm{d}x = y\,\mathrm{d}u$）。
    当 $y > 0$ 时，$x$ 的区间 $(-\infty, zy]$ 对应 $u$ 的区间 $(-\infty, z]$；
    当 $y < 0$ 时，$x$ 的区间 $[zy, +\infty)$ 对应 $u$ 的区间 $[z, -\infty)$，此时 $y = -|y|$，积分限反转吸收负号后得：
    \begin{align*}
    F_Z(z)& = \int_{0}^{+\infty} \left( \int_{-\infty}^{z} f(uy, y)y\,\mathrm{d}u \right) \mathrm{d}y + \int_{-\infty}^{0} \left( \int_{-\infty}^{z} f(uy, y)(-y)\,\mathrm{d}u \right) \mathrm{d}y\\
    & = \int_{-\infty}^{z} \left[ \int_{-\infty}^{+\infty} |y|f(uy, y)\,\mathrm{d}y \right] \mathrm{d}u.
    \end{align*}
    对 $z$ 求导即得 $f_Z(z) = \displaystyle\int_{-\infty}^{+\infty}|y|\,f(yz, y)\,\mathrm{d}y$。独立时展开联合密度即得特例。

    \item \textbf{二维变量代换法（雅可比行列式法）}：
    作变换 $Z = X/Y, W = Y$。逆变换为 $x = zw, y = w$。
    雅可比行列式为 $J = \frac{\partial(x,y)}{\partial(z,w)} = \begin{vmatrix} w & z \\ 0 & 1 \end{vmatrix} = w$，故 $|J| = |w|$。
    $(Z,W)$ 的联合密度为 $g(z,w) = f(zw, w) |w|$。
    求边缘密度，将 $w$ 积分掉（用 $y$ 替换积分变量 $w$）：
    \[
    f_Z(z) = \int_{-\infty}^{+\infty} g(z,w)\,\mathrm{d}w = \int_{-\infty}^{+\infty} |y|f(zy, y)\,\mathrm{d}y.
    \]
\end{enumerate}
\end{proof}

\begin{property}[常见分布的可加性]\label{prop:additivity}
有些相互独立且服从同类型分布的随机变量，其和的分布也是同类型的。设 $X$ 与 $Y$ 相互独立，则
\begin{enumerate}
    \item 二项分布：若 $X \sim B(n,p)$，$Y \sim B(m,p)$，则 $X + Y \sim B(n+m,p)$（注意 $p$ 必须相同）。
    \item 泊松分布：若 $X \sim P(\lambda_1)$，$Y \sim P(\lambda_2)$，则 $X + Y \sim P(\lambda_1 + \lambda_2)$。
    \item 正态分布：若 $X \sim N(\mu_1, \sigma_1^2)$，$Y \sim N(\mu_2, \sigma_2^2)$，则 $X + Y \sim N(\mu_1 + \mu_2, \sigma_1^2 + \sigma_2^2)$。更一般地，独立正态变量的线性组合
    \[
    \sum_{i=1}^{n}c_i X_i \sim N\!\left(\sum_{i=1}^{n}c_i\mu_i,\ \sum_{i=1}^{n}c_i^2\sigma_i^2\right).
    \]
    \item $\chi^2$ 分布：若 $X \sim \chi^2(n_1)$，$Y \sim \chi^2(n_2)$，则 $X + Y \sim \chi^2(n_1 + n_2)$（$\chi^2$ 分布详见第六章）。
    \item Gamma分布：若 $X \sim \Gamma(\alpha_1, \beta)$，$Y \sim \Gamma(\alpha_2, \beta)$，则 $X + Y \sim \Gamma(\alpha_1 + \alpha_2, \beta)$。特别地，$X_i \sim E(\lambda)$ 独立时，$\sum_{i=1}^{n}X_i \sim \Gamma(n, \lambda)$。
\end{enumerate}
以上结果对 $n$ 个相互独立的随机变量也成立。
\end{property}
\begin{proof}[二项分布可加性的证明]
利用离散型随机变量的卷积公式（也可视为全概率公式）。

对于任意 $k \in \{0, 1, \dots, n+m\}$，有
\begin{align*}
    P(X+Y=k) &= \sum_{i=0}^{k} P(X=i)P(Y=k-i) \\
    &= \sum_{i=0}^{k} \left[ \binom{n}{i} p^i (1-p)^{n-i} \right] \left[ \binom{m}{k-i} p^{k-i} (1-p)^{m-(k-i)} \right] \\
    &= p^k (1-p)^{n+m-k} \sum_{i=0}^{k} \binom{n}{i} \binom{m}{k-i}.
\end{align*}
由组合恒等式（范德蒙恒等式）可知 $\sum_{i=0}^{k} \binom{n}{i} \binom{m}{k-i} = \binom{n+m}{k}$。因此，
\[
    P(X+Y=k) = \binom{n+m}{k} p^k (1-p)^{n+m-k},
\]
即 $X+Y \sim B(n+m, p)$。
\end{proof}

\begin{proof}[泊松分布可加性的证明]
利用离散型随机变量的卷积公式。对于任意 $k \in \{0, 1, 2, \dots\}$，有
\begin{align*}
    P(X+Y=k) &= \sum_{i=0}^{k} P(X=i)P(Y=k-i) \\
    &= \sum_{i=0}^{k} \left( \frac{\lambda_1^i}{i!} e^{-\lambda_1} \right) \left( \frac{\lambda_2^{k-i}}{(k-i)!} e^{-\lambda_2} \right) \\
    &= e^{-(\lambda_1+\lambda_2)} \sum_{i=0}^{k} \frac{\lambda_1^i \lambda_2^{k-i}}{i!(k-i)!} \\
    &= \frac{e^{-(\lambda_1+\lambda_2)}}{k!} \sum_{i=0}^{k} \frac{k!}{i!(k-i)!} \lambda_1^i \lambda_2^{k-i}.
\end{align*}
利用二项式定理，$\sum_{i=0}^{k} \binom{k}{i} \lambda_1^i \lambda_2^{k-i} = (\lambda_1 + \lambda_2)^k$。因此，
\[
    P(X+Y=k) = \frac{(\lambda_1+\lambda_2)^k}{k!} e^{-(\lambda_1+\lambda_2)},
\]
即 $X+Y \sim P(\lambda_1+\lambda_2)$。
\end{proof}

\begin{proof}[正态分布可加性的证明]
利用特征函数证明。若 $X \sim N(\mu, \sigma^2)$，其特征函数为
\[
\varphi_X(t) = \exp\left(i\mu t - \frac{1}{2}\sigma^2 t^2\right).
\]
由于 $X_i$ 相互独立，对于任意常数 $c_j$，$Z = \sum_{j=1}^{n}c_j X_j$ 的特征函数为
\begin{align*}
    \varphi_Z(t) &= E\left[\exp\left(it\sum_{j=1}^{n}c_j X_j\right)\right] = \prod_{j=1}^{n} E\left[\exp(it c_j X_j)\right] \\
    &= \prod_{j=1}^{n} \varphi_{X_j}(c_j t) = \prod_{j=1}^{n} \exp\left(i\mu_j c_j t - \frac{1}{2}\sigma_j^2 c_j^2 t^2\right) \\
    &= \exp\left(it\sum_{j=1}^{n}c_j\mu_j - \frac{1}{2}t^2\sum_{j=1}^{n}c_j^2\sigma_j^2\right).
\end{align*}
此形式恰好是均值为 $\sum_{j=1}^{n}c_j\mu_j$，方差为 $\sum_{j=1}^{n}c_j^2\sigma_j^2$ 的正态分布的特征函数。由特征函数唯一决定分布定理，结论得证。
\end{proof}

\begin{proof}[$\chi^2$ 分布可加性的证明]
利用矩母函数（或特征函数）证明。$\chi^2(n)$ 分布的矩母函数为 $M(t) = (1-2t)^{-n/2}$（其中 $t < 1/2$）。
由于 $X \sim \chi^2(n_1)$ 和 $Y \sim \chi^2(n_2)$ 相互独立，其和的矩母函数等于各自独立矩母函数的乘积：
\[
    M_{X+Y}(t) = M_X(t) M_Y(t) = (1-2t)^{-n_1/2} (1-2t)^{-n_2/2} = (1-2t)^{-(n_1+n_2)/2}.
\]
这恰为自由度为 $n_1+n_2$ 的 $\chi^2$ 分布的矩母函数。由矩母函数的唯一性定理，结论得证。
\end{proof}

\begin{proof}[Gamma分布可加性的证明]
利用矩母函数证明。根据题设中指数分布 $X_i \sim E(\lambda)$ 为特例的参数化约定，此处的 $\beta$ 为率参数。Gamma 分布 $\Gamma(\alpha, \beta)$ 的矩母函数为 $M(t) = \left( \frac{\beta}{\beta-t} \right)^\alpha$（其中 $t < \beta$）。
由于 $X \sim \Gamma(\alpha_1, \beta)$ 和 $Y \sim \Gamma(\alpha_2, \beta)$ 相互独立，有：
\[
    M_{X+Y}(t) = M_X(t) M_Y(t) = \left( \frac{\beta}{\beta-t} \right)^{\alpha_1} \left( \frac{\beta}{\beta-t} \right)^{\alpha_2} = \left( \frac{\beta}{\beta-t} \right)^{\alpha_1+\alpha_2}.
\]
这正是参数为 $\alpha_1+\alpha_2$ 和 $\beta$ 的 Gamma 分布的矩母函数。由唯一性定理，结论得证。

特别地，指数分布 $E(\lambda)$ 可视为 $\Gamma(1, \lambda)$。若有 $n$ 个相互独立的 $X_i \sim E(\lambda)$，将其矩母函数累乘即可得到 $M_{\sum X_i}(t) = \left( \frac{\lambda}{\lambda-t} \right)^n$，因此 $\sum_{i=1}^{n}X_i \sim \Gamma(n, \lambda)$。
\end{proof}


\subsubsection{\texorpdfstring{$\max$ 和 $\min$ 的分布}{max 和 min 的分布}}

在可靠性理论和排序问题中，$\max\{X,Y\}$ 和 $\min\{X,Y\}$ 的分布具有重要的应用价值。

\begin{property}[$\max$ 和 $\min$ 的分布]\label{prop:max-min-dist}
设 $(X,Y) \sim f(x,y)$，分布函数为 $F(x,y)$。

\textbf{(1) 一般情形}：
\begin{enumerate}
    \item $Z = \max\{X, Y\}$ 的分布函数为
    \[
    F_{\max}(z) = P\{X \leqslant z, Y \leqslant z\} = F(z, z).
    \]
    \item $Z = \min\{X, Y\}$ 的分布函数为
    \[
    F_{\min}(z) = P\{X \leqslant z\} + P\{Y \leqslant z\} - P\{X \leqslant z, Y \leqslant z\} = F_X(z) + F_Y(z) - F(z, z).
    \]
\end{enumerate}

\textbf{(2) 独立情形}：
\begin{enumerate}
    \item 当 $X$ 与 $Y$ 独立时，$F_{\max}(z) = F_X(z) \cdot F_Y(z)$，$F_{\min}(z) = 1 - [1 - F_X(z)][1 - F_Y(z)]$。
    \item 推广到 $n$ 个相互独立的随机变量 $X_1, X_2, \dots, X_n$（分布函数为 $F_{X_i}(x)$）：
    \[
    F_{\max}(z) = \prod_{i=1}^{n}F_{X_i}(z), \qquad F_{\min}(z) = 1 - \prod_{i=1}^{n}[1 - F_{X_i}(z)].
    \]
    \item 当 $X_1, \dots, X_n$ 独立同分布（分布函数为 $F(x)$，概率密度为 $f(x)$）时：
    \[
    F_{\max}(z) = [F(z)]^n, \quad f_{\max}(z) = n[F(z)]^{n-1}f(z),
    \]
    \[
    F_{\min}(z) = 1 - [1 - F(z)]^n, \quad f_{\min}(z) = n[1 - F(z)]^{n-1}f(z).
    \]
    这些结果在数理统计部分有重要应用。
\end{enumerate}
\end{property}

\begin{remark}
\begin{enumerate}
    \item "$\max\{X_1, \dots, X_n\} \leqslant z$" 等价于"所有 $X_i \leqslant z$"，因此取概率时用乘积。
    \item "$\min\{X_1, \dots, X_n\} > z$" 等价于"所有 $X_i > z$"，因此取逆事件后用乘积。
    \item 推导 $F_{\min}(z)$ 的关键步骤：
    \[
    P\{\min(X,Y) \leqslant z\} = 1 - P\{\min(X,Y) > z\} = 1 - P\{X > z, Y > z\} = 1 - [1-F_X(z)][1-F_Y(z)].
    \]
\end{enumerate}
\end{remark}

\begin{exercise}\label{exer:3-9}
一辆新车装有 4 个轮胎并配有 1 个备胎，当使用的轮胎损坏时就用备胎更换。设所有轮胎的寿命均服从参数为 $\lambda = 0.2$（单位：万小时）的指数分布且相互独立，记 $T$ 为开始使用备胎的时间，求 $T$ 的分布函数。
\end{exercise}

\begin{solution}
开始使用备胎的时刻 $T$ 等于 4 个轮胎中寿命最短的那个的寿命，即 $T = \min\{X_1, X_2, X_3, X_4\}$。

每个轮胎的寿命 $X_i \sim E(0.2)$，分布函数为 $F(x) = 1 - e^{-0.2x}$（$x \geqslant 0$）。

由最小值的分布公式：
\[
F_T(t) = 1 - [1 - F(t)]^4 = 1 - [e^{-0.2t}]^4 = 1 - e^{-0.8t}, \quad t \geqslant 0.
\]

即 $T \sim E(0.8)$，备胎在参数为 $0.8$ 的指数分布时刻开始被使用。

\begin{note}
更一般地，若 $X_1, X_2, \dots, X_n$ 独立同分布于 $E(\lambda)$，则 $\min\{X_1, \dots, X_n\} \sim E(n\lambda)$。
\end{note}
\end{solution}

\begin{exercise}[并联系统最大值分布] \label{exer:3-10}
设系统由两个元件并联组成，两个元件的寿命分别为 $X$ 和 $Y$，且 $X, Y$ 相互独立，均服从参数为 $\lambda = 1/4$（万小时）的指数分布。求系统寿命 $S = \max\{X, Y\}$ 大于 5 万小时的概率。
\end{exercise}

\begin{solution}
$X, Y \sim E(1/4)$，分布函数 $F(x) = 1 - e^{-x/4}$（$x \geqslant 0$）。

$S = \max\{X, Y\}$ 的分布函数为 $F_S(s) = [F(s)]^2 = (1-e^{-s/4})^2$。

所求概率：
\[
P\{S > 5\} = 1 - F_S(5) = 1 - (1 - e^{-5/4})^2 = 2e^{-5/4} - e^{-5/2}.
\]
\end{solution}

\begin{exercise}[均匀分布最大值密度] \label{exer:3-10b}
设随机变量 $X, Y$ 相互独立，且都在区间 $[a, b]$ 上服从均匀分布，求 $Z = \max\{X, Y\}$ 的概率密度。
\end{exercise}

\begin{solution}
由独立性，$(X,Y)$ 的概率密度为
\[
f(x,y) = \begin{cases} \dfrac{1}{(b-a)^2}, & a \leqslant x \leqslant b,\ a \leqslant y \leqslant b, \\ 0, & \text{其他}. \end{cases}
\]

$Z = \max\{X, Y\}$ 的取值范围为 $[a, b]$。当 $a \leqslant z \leqslant b$ 时，
\[
F_Z(z) = P\{Z \leqslant z\} = P\{X \leqslant z,\ Y \leqslant z\} = \int_a^z\mathrm{d}x\int_a^z \frac{1}{(b-a)^2}\,\mathrm{d}y = \frac{(z-a)^2}{(b-a)^2}.
\]

求导得概率密度：
\[
f_Z(z) = F_Z'(z) = \begin{cases} \dfrac{2(z-a)}{(b-a)^2}, & a \leqslant z \leqslant b, \\ 0, & \text{其他}. \end{cases}
\]

\begin{note}
这与独立同分布时 $f_{\max}(z) = n[F(z)]^{n-1}f(z)$ 的公式一致：$n=2$，$F(z) = \dfrac{z-a}{b-a}$，$f(z) = \dfrac{1}{b-a}$，故 $f_{\max}(z) = 2 \cdot \dfrac{z-a}{b-a} \cdot \dfrac{1}{b-a} = \dfrac{2(z-a)}{(b-a)^2}$。
\end{note}
\end{solution}

\begin{exercise}[柯西分布线性变换] \label{exer:3-11}
设随机变量 $X$ 的概率密度为 $f(x) = \dfrac{1}{\pi(1+x^2)}$（柯西分布），$Y = 2X$，求 $Y$ 的概率密度。
\end{exercise}

\begin{solution}
$y = g(x) = 2x$ 是严格单调函数，反函数 $x = h(y) = y/2$，$h'(y) = 1/2$。

由公式法：
\[
f_Y(y) = f_X\!\left(\frac{y}{2}\right) \cdot \left|\frac{1}{2}\right| = \frac{1}{\pi\left[1 + (y/2)^2\right]} \cdot \frac{1}{2} = \frac{2}{\pi(4+y^2)}.
\]
\end{solution}

\begin{exercise}[正态平方变换] \label{exer:3-12}
设 $X \sim N(0,1)$，求 $Y = 2X^2 + 1$ 的概率密度函数。
\end{exercise}
\begin{solution}
当 $y \leqslant 1$ 时，$2X^2 + 1 \leqslant y$ 无解（$X^2 \geqslant 0$），故 $F_Y(y) = 0$，$f_Y(y) = 0$。\\
当 $y > 1$ 时：
\[
F_Y(y) = P\{2X^2 + 1 \leqslant y\} = P\!\left\{-\sqrt{\frac{y-1}{2}} \leqslant X \leqslant \sqrt{\frac{y-1}{2}}\right\} = 2\Phi\!\left(\sqrt{\frac{y-1}{2}}\right) - 1.
\]
求导，令 $\varphi(x) = \dfrac{1}{\sqrt{2\pi}}e^{-x^2/2}$：
\[
f_Y(y) = 2\varphi\!\left(\sqrt{\frac{y-1}{2}}\right) \cdot \frac{1}{2\sqrt{2(y-1)}} = \frac{1}{2\sqrt{\pi(y-1)}}\,e^{-(y-1)/4}.
\]
综上：
\[
f_Y(y) = \begin{cases} \dfrac{1}{2\sqrt{\pi(y-1)}}\,\exp\!\left\{-\dfrac{y-1}{4}\right\}, & y > 1, \\ 0, & y \leqslant 1. \end{cases}
\]
\end{solution}



\begin{exercise}[指数分布条件区域] \label{exer:3-12b}
设二维随机变量 $(X,Y)$ 的概率密度函数为
\[
f(x,y) = \begin{cases} \lambda e^{-\lambda y}, & 0 < x < y, \\ 0, & \text{其他}, \end{cases}
\]
其中 $\lambda > 0$。求：(1) $Z = X + Y$ 的概率密度；(2) $N = \min(X,Y)$ 的概率密度。
\end{exercise}

\begin{solution}
验证归一性：$\displaystyle\int_0^{+\infty}\mathrm{d}y\int_0^y \lambda e^{-\lambda y}\,\mathrm{d}x = \int_0^{+\infty}y\lambda e^{-\lambda y}\,\mathrm{d}y = 1$（$Y \sim \Gamma(2,\lambda)$）。

边缘密度：
\[
f_X(x) = \int_x^{+\infty}\lambda e^{-\lambda y}\,\mathrm{d}y = e^{-\lambda x}, \quad x > 0. \qquad \text{即 } X \sim E(\lambda).
\]
\[
f_Y(y) = \int_0^y \lambda e^{-\lambda y}\,\mathrm{d}x = \lambda y\,e^{-\lambda y}, \quad y > 0. \qquad \text{即 } Y \sim \Gamma(2,\lambda).
\]

(1) 当 $z > 0$ 时，$\{X+Y \leqslant z\}$ 等价于 $\{0 < x < y,\ x+y \leqslant z\}$，即 $0 < x < z/2$，$x < y \leqslant z-x$：
\begin{align*}
F_Z(z) &= \int_0^{z/2}\mathrm{d}x\int_x^{z-x}\lambda e^{-\lambda y}\,\mathrm{d}y = \int_0^{z/2}(e^{-\lambda x} - e^{-\lambda(z-x)})\,\mathrm{d}x\\
&= \left[-\frac{1}{\lambda}e^{-\lambda x} - \frac{1}{\lambda}e^{-\lambda z}e^{\lambda x}\right]_0^{z/2} = \frac{1}{\lambda}\bigl(1 - e^{-\lambda z/2}\bigr) - \frac{1}{\lambda}e^{-\lambda z}(e^{\lambda z/2} - 1)\\
&= \frac{1}{\lambda}\bigl(1 - e^{-\lambda z/2} - e^{-\lambda z/2} + e^{-\lambda z}\bigr) = \frac{1}{\lambda}\bigl(1 - 2e^{-\lambda z/2} + e^{-\lambda z}\bigr).
\end{align*}
当 $z \leqslant 0$ 时，$F_Z(z) = 0$。求导：
\[
f_Z(z) = \frac{1}{\lambda}\bigl(\lambda e^{-\lambda z/2} - \lambda e^{-\lambda z}\bigr) = e^{-\lambda z/2} - e^{-\lambda z}, \quad z > 0.
\]
所以
\[
f_Z(z) = \begin{cases} e^{-\lambda z/2} - e^{-\lambda z}, & z > 0, \\ 0, & z \leqslant 0. \end{cases}
\]
(2) 在区域 $\{0 < x < y\}$ 中，$\min(X,Y) = X$。因此 $N = X$，其密度为
\[
f_N(n) = f_X(n) = \begin{cases} e^{-\lambda n}, & n > 0, \\ 0, & n \leqslant 0. \end{cases}
\]
即 $N \sim E(\lambda)$。
\end{solution}


\begin{exercise}[卷积公式：公式法与分布函数法对照] \label{exer:3-12c}
随机变量 $X, Y$ 相互独立，概率密度分别为
\[
f_X(x) = \begin{cases} e^{-x}, & x > 0, \\ 0, & x \leqslant 0, \end{cases} \quad f_Y(y) = \begin{cases} \dfrac{1}{2}e^{-y/2}, & y > 0, \\ 0, & y \leqslant 0. \end{cases}
\]
求 $Z = X + Y$ 的概率密度。
\end{exercise}
\begin{solution}
\textbf{方法一（公式法）}：由 $Z = X+Y$，故 $Z > 0$。$f_Z(z) = \displaystyle\int_{-\infty}^{+\infty}f_X(x)\,f_Y(z-x)\,\mathrm{d}x$。$f_Y(z-x) \neq 0$ 要求 $z-x > 0$，即 $x < z$；$f_X(x) \neq 0$ 要求 $x > 0$。当 $z > 0$ 时：
\[
f_Z(z) = \int_0^z e^{-x} \cdot \frac{1}{2}e^{-(z-x)/2}\,\mathrm{d}x = \frac{1}{2}e^{-z/2}\int_0^z e^{-x/2}\,\mathrm{d}x = \frac{1}{2}e^{-z/2}\bigl[-2e^{-x/2}\bigr]_0^z = e^{-z/2} - e^{-z}.
\]
\textbf{方法二（分布函数法）}：当 $z > 0$ 时，
\begin{align*}
F_Z(z)& = P\{X+Y \leqslant z\} = \int_0^z e^{-x}\,\mathrm{d}x\int_0^{z-x}\frac{1}{2}e^{-y/2}\,\mathrm{d}y = \int_0^z e^{-x}\bigl(1-e^{-(z-x)/2}\bigr)\,\mathrm{d}x\\
&= \int_0^z(e^{-x} - e^{-z/2}e^{-x/2})\,\mathrm{d}x = \bigl[-e^{-x} + 2e^{-z/2}e^{-x/2}\bigr]_0^z \\
&= 1 - e^{-z} - 2e^{-z/2} + 2e^{-z} = 1 - 2e^{-z/2} + e^{-z}.
\end{align*}
求导：$f_Z(z) = e^{-z/2} - e^{-z}$（$z > 0$）。综上：
\[
f_Z(z) = \begin{cases} e^{-z/2} - e^{-z}, & z > 0, \\ 0, & z \leqslant 0. \end{cases}
\]
\begin{note}
本题用两种方法得到一致结果。当 $X \sim E(1)$，$Y \sim E(1/2)$ 时，$Z = X+Y$ 的密度为 $e^{-z/2} - e^{-z}$。这恰好是两个不同参数指数分布之和，不服从简单的指数分布（因为参数不同，不可直接用Gamma可加性公式 $\Gamma(\alpha_1,\beta) + \Gamma(\alpha_2,\beta) = \Gamma(\alpha_1+\alpha_2,\beta)$，这里 $\beta$ 不同）。
\end{note}
\end{solution}

\begin{exercise}[综合性：常数、独立性、和的分布、min的概率] \label{exer:3-12e}
设二维随机变量 $(X,Y)$ 具有联合密度函数
\[
f(x,y) = \begin{cases} cxe^{-y}, & 0 < x < y < +\infty, \\ 0, & \text{其他}. \end{cases}
\]
求：
\begin{enumerate}
    \item[(1)] 常数 $c$；
    \item[(2)] $X$ 与 $Y$ 是否独立？
    \item[(3)] $Z = X + Y$ 的密度函数；
    \item[(4)] $P\{\min(X,Y) < 1\}$。
\end{enumerate}
\end{exercise}
\begin{solution}
(1) 由概率密度的归一性：
\[
1 = \int_{-\infty}^{+\infty}\int_{-\infty}^{+\infty}f(x,y)\,\mathrm{d}x\,\mathrm{d}y = \int_0^{+\infty}\int_0^y cxe^{-y}\,\mathrm{d}x\,\mathrm{d}y.
\]
内层积分 $\int_0^y x\,\mathrm{d}x = \dfrac{y^2}{2}$，因此
\[
\int_0^{+\infty}c \cdot \frac{y^2}{2}e^{-y}\,\mathrm{d}y = \frac{c}{2}\Gamma(3) = \frac{c}{2} \cdot 2! = c = 1.
\]
(2) 求边缘概率密度：
\[
f_X(x) = \int_x^{+\infty}xe^{-y}\,\mathrm{d}y = xe^{-x}, \quad x > 0. \qquad f_Y(y) = \int_0^y xe^{-y}\,\mathrm{d}x = \frac{y^2}{2}e^{-y}, \quad y > 0.
\]
由于 $f_X(x)\,f_Y(y) = \dfrac{xy^2}{2}e^{-(x+y)} \neq xe^{-y} = f(x,y)$（取 $x=1, y=2$ 可直接验证），故 $X$ 与 $Y$ \textbf{不独立}。\\
(3) \textbf{方法一（卷积公式）}：$f_Z(z) = \displaystyle\int_{-\infty}^{+\infty}f(x, z-x)\,\mathrm{d}x$。
$f(x, z-x) \neq 0$ 要求 $0 < x < z-x$，即 $0 < x < \dfrac{z}{2}$。当 $z > 0$ 时：
\[
f_Z(z) = \int_0^{z/2}xe^{-(z-x)}\,\mathrm{d}x = e^{-z}\int_0^{z/2}xe^x\,\mathrm{d}x = e^{-z}\left[\left(\frac{z}{2}-1\right)e^{z/2} + 1\right] = \left(\frac{z}{2}-1\right)e^{-z/2} + e^{-z}.
\]
\textbf{方法二（分布函数法）}：当 $z < 0$ 时，$F_Z(z) = 0$；当 $z \geqslant 0$ 时，积分区域为 $\{0 < x < y,\ x+y \leqslant z\}$，即 $0 < x < z/2$，$x < y \leqslant z-x$：
\[
F_Z(z) = \int_0^{z/2}\mathrm{d}x\int_x^{z-x}xe^{-y}\,\mathrm{d}y = \int_0^{z/2}x(e^{-x} - e^{-(z-x)})\,\mathrm{d}x
\]
\[
= \left[1 - \left(\frac{z}{2}+1\right)e^{-z/2}\right] - e^{-z}\left[\left(\frac{z}{2}-1\right)e^{z/2} + 1\right] = 1 - ze^{-z/2} - e^{-z}.
\]
求导验证：
\[
f_Z(z) = F_Z'(z) = \left(\frac{z}{2}-1\right)e^{-z/2} + e^{-z}, \quad z \geqslant 0,
\]
与方法一结果一致。综上：
\[
f_Z(z) = \begin{cases} \left(\dfrac{z}{2}-1\right)e^{-z/2} + e^{-z}, & z \geqslant 0, \\ 0, & z < 0. \end{cases}
\]
(4) 在区域 $\{0 < x < y\}$ 中，$\min(X,Y) = X$。利用逆事件：
\[
P\{\min(X,Y) < 1\} = 1 - P\{\min(X,Y) \geqslant 1\} = 1 - P\{X \geqslant 1, Y \geqslant 1\}
\]
\[
= 1 - \int_1^{+\infty}\mathrm{d}y\int_1^y xe^{-y}\,\mathrm{d}x = 1 - \int_1^{+\infty}\frac{y^2-1}{2}e^{-y}\,\mathrm{d}y
\]
\[
= 1 - \frac{1}{2}\left[\int_1^{+\infty}y^2 e^{-y}\,\mathrm{d}y - \int_1^{+\infty}e^{-y}\,\mathrm{d}y\right] = 1 - \frac{1}{2}(5e^{-1} - e^{-1}) = 1 - 2e^{-1}.
\]
\begin{note}
\begin{enumerate}
    \item 第 (3) 问分布函数法中，$F_Z(z) = 1 - ze^{-z/2} \boldsymbol{-} e^{-z}$（注意 $e^{-z}$ 前为负号）。可通过取 $z = 2$ 验证：直接计算得 $1 - 2e^{-1} - e^{-2}$，而 $1 - 2e^{-1} + e^{-2}$ 与之不符。
    \item 第 (4) 问的技巧：在 $\{0 < x < y\}$ 的条件下 $\min(X,Y) = X$，因此 $P\{X \geqslant 1, Y \geqslant 1\}$ 的积分区域为 $\{1 \leqslant x < y\}$。
\end{enumerate}
\end{note}
\end{solution}


\begin{exercise}[商的分布] \label{exer:3-c15}
$X, Y$ 独立同服从指数分布 $E(\lambda)$，求 $Z = X/Y$ 的概率密度。
\end{exercise}

\begin{solution}
由商的分布公式，$Z > 0$ 时：
\[
f_Z(z) = \int_0^{+\infty}|y|\,f_X(zy)\,f_Y(y)\,\mathrm{d}y = \int_0^{+\infty}y\cdot\lambda e^{-\lambda zy}\cdot\lambda e^{-\lambda y}\,\mathrm{d}y
\]
\[
= \lambda^2\int_0^{+\infty}ye^{-\lambda(z+1)y}\,\mathrm{d}y = \lambda^2\cdot\frac{1}{\lambda^2(z+1)^2} = \frac{1}{(z+1)^2}.
\]

\[
f_Z(z) = \begin{cases} \dfrac{1}{(z+1)^2}, & z > 0, \\ 0, & z \leqslant 0. \end{cases}
\]

\begin{note}
$Z = X/Y$ 服从 \textbf{Pareto 型分布}，其期望 $E[Z] = \displaystyle\int_0^{+\infty}\frac{z}{(z+1)^2}\,\mathrm{d}z$ 发散（令 $u = z+1$ 可验证）。
\end{note}
\end{solution}

\begin{exercise}[均匀分布 min 的密度] \label{exer:3-c16}
$X, Y$ 独立同服从均匀分布 $U[a,b]$，求 $Z = \min(X,Y)$ 的概率密度。
\end{exercise}
\begin{solution}
$X, Y$ 的分布函数 $F(x) = \dfrac{x-a}{b-a}$（$a \leqslant x \leqslant b$），密度 $f(x) = \dfrac{1}{b-a}$。
由 $\min$ 分布公式：
\[
F_Z(z) = 1 - [1-F(z)]^2 = 1 - \left(1 - \frac{z-a}{b-a}\right)^2 = 1 - \left(\frac{b-z}{b-a}\right)^2.
\]
求导得密度：
\[
f_Z(z) = \frac{2(b-z)}{(b-a)^2}\cdot\frac{1}{b-a} = \frac{2(b-z)}{(b-a)^2}, \quad a \leqslant z \leqslant b.
\]
所以
\[
f_Z(z) = \begin{cases} \dfrac{2(b-z)}{(b-a)^2}, & a \leqslant z \leqslant b, \\ 0, & \text{其他}. \end{cases}
\]
\begin{note}
对照 \ref{exer:3-10b} 中 $\max$ 的密度 $\dfrac{2(z-a)}{(b-a)^2}$，可见 $\min$ 和 $\max$ 的密度关于区间中点对称。
\end{note}
\end{solution}

\begin{exercise}[泊松可加性] \label{exer:3-c6b}
$X\sim P(5)$，$Y\sim P(3)$，且 $X$ 与 $Y$ 独立，则（\quad）\\
A. $X+Y\sim P(2)$ \quad B. $X-Y\sim P(2)$ \quad C. $X+Y\sim P(8)$ \quad D. $X-Y\sim P(8)$
\end{exercise}
\begin{solution}
应选 \textbf{C}。

本题直接考查泊松分布的可加性。若
\[
X\sim P(\lambda_1),\qquad Y\sim P(\lambda_2),
\]
并且 \(X,Y\) 相互独立，那么有
\[
X+Y\sim P(\lambda_1+\lambda_2).
\]

代入本题 \(\lambda_1=5,\ \lambda_2=3\)，得到 \(X+Y\sim P(5+3)=P(8)\)。

再看减法为什么不对。随机变量 \(X-Y\) 可能取负值，例如 \(X=1,\ Y=3\) 时就有 \(X-Y=-2\)。而泊松分布的取值只能是 \(0,1,2,\dots\) 这些非负整数，因此 \(X-Y\) 不可能服从泊松分布。

所以只有 C 正确。
\end{solution}


\begin{exercise}[次序统计量的分布] \label{exer:3-c8b}
$X_1,X_2,\dots,X_5$ 为总体 $N(1,4)$ 的简单随机样本，$X_{(5)} = \max(X_1,\dots,X_5)$，求 $P(X_{(5)}<3)$。
\end{exercise}
\begin{solution}
$P(X<3) = \Phi\!\left(\dfrac{3-1}{2}\right) = \Phi(1) \approx 0.8413$。\\
由独立性，$X_{(5)} = \max$ 的分布函数 $F_{(5)}(z) = [F(z)]^5$，故
\[
P(X_{(5)}<3) = \Phi^5(1) = 0.8413^5 \approx 0.4214.
\]
\end{solution}


\begin{exercise}[min+max=X+Y] \label{exer:3-c10b}
$(X,Y)$ 服从区域 $G=\{(x,y):0\leqslant x\leqslant 1,\ 0\leqslant y\leqslant 2\}$ 上的均匀分布，令 $Z=\min(X,Y)$，$W=\max(X,Y)$，求 $P(Z+W \geqslant 1)$。
\end{exercise}
\begin{solution}
先注意到
\[
Z=\min(X,Y),\qquad W=\max(X,Y),
\]
无论 \(X\) 与 \(Y\) 谁大谁小，总有
\[
Z+W=X+Y.
\]
因此题目要求的概率就是
\[
P(Z+W\geqslant 1)=P(X+Y\geqslant 1).
\]

又因为 \((X,Y)\) 在区域
\[
G=\{(x,y):0\le x\le 1,\ 0\le y\le 2\}
\]
上服从均匀分布，而该矩形面积为 \(1\times 2=2\)，所以联合密度为
\[
f(x,y)=\frac12\qquad ((x,y)\in G).
\]

直接算补事件更方便：
\[
P(X+Y \geqslant 1) = 1 - P(X+Y < 1) = 1 - \int_0^1\mathrm{d}x\int_0^{1-x}\frac{1}{2}\,\mathrm{d}y = 1 - \frac{1}{2}\int_0^1(1-x)\,\mathrm{d}x = 1-\frac{1}{4} = \frac{3}{4}.
\]

这里的积分区域 \(X+Y<1\) 在矩形 \(G\) 内正好是一个直角三角形，所以结果也可以理解为“用三角形面积除以矩形总面积”。

故
\[
P(Z+W\geqslant 1)=\frac34.
\]
\end{solution}


\begin{exercise}[半正态联合密度与极坐标变换] \label{exer:3-c13b}
$(X,Y)$ 的联合概率密度为
\[
f(x,y) = \begin{cases} \dfrac{2}{\pi}\,e^{-\frac{x^2+y^2}{2}}, & x>0,\ y>0, \\ 0, & \text{其他}. \end{cases}
\]
求：(1) 边缘密度；(2) 独立性；(3) $Z=X^2+Y^2$ 的密度。
\end{exercise}
\begin{solution}
(1) 利用 $\displaystyle\int_0^{+\infty}e^{-y^2/2}\,\mathrm{d}y = \sqrt{\frac{\pi}{2}}$：
\[
f_X(x) = \frac{2}{\pi}\,e^{-x^2/2}\int_0^{+\infty}e^{-y^2/2}\,\mathrm{d}y = \frac{2}{\pi}\,e^{-x^2/2}\cdot\sqrt{\frac{\pi}{2}} = \sqrt{\frac{2}{\pi}}\,e^{-x^2/2},\quad x>0.
\]
同理 $f_Y(y) = \sqrt{2/\pi}\,e^{-y^2/2}$（$y>0$）。此为\textbf{半正态分布}（折叠正态分布）。\\
(2) $f(x,y) = \dfrac{2}{\pi}\,e^{-(x^2+y^2)/2} = f_X(x)\,f_Y(y)$，故 $X$ 与 $Y$ \textbf{独立}。\\
(3) $Z=X^2+Y^2 > 0$。用极坐标变换 $x=r\cos\theta$，$y=r\sin\theta$（$0<\theta<\pi/2$，$r>0$）：
\[
F_Z(z) = \iint_{x^2+y^2 \leqslant z,\ x>0,\ y>0}\frac{2}{\pi}\,e^{-(x^2+y^2)/2}\,\mathrm{d}x\,\mathrm{d}y = \frac{2}{\pi}\int_0^{\pi/2}\mathrm{d}\theta\int_0^{\sqrt{z}}r\,e^{-r^2/2}\,\mathrm{d}r
\]
\[
= \frac{2}{\pi}\cdot\frac{\pi}{2}\left[-e^{-r^2/2}\right]_0^{\sqrt{z}} = 1-e^{-z/2},\quad z>0.
\]
求导得 $f_Z(z) = \dfrac{1}{2}e^{-z/2}$（$z>0$），即 $Z \sim E(1/2)$（或等价地 $Z\sim\chi^2(2)$）。
\begin{note}
本题说明：\textbf{两个独立半正态变量的平方和服从指数分布（自由度为2的$\chi^2$分布）}。更一般地，若 $X_i$ 独立同分布于 $N(0,\sigma^2)$，则 $\sum_{i=1}^n X_i^2 \sim \sigma^2\chi^2(n)$。
\end{note}
\end{solution}

\begin{note}
\textbf{【高频题型模板】二维连续型综合题的解题管线：}

考试中二维连续型大题通常按以下顺序出小问，每步都有固定套路：

\textbf{第一问：求边缘密度}
\[
f_X(x)=\int_{-\infty}^{+\infty}f(x,y)\,\mathrm{d}y\quad\text{（积分限由支撑域中 $y$ 的范围确定）}
\]

\textbf{第二问：判断独立性}
\[
\text{验证 } f(x,y)\overset{?}{=}f_X(x)\cdot f_Y(y) \text{ 在整个支撑域上是否恒成立}
\]
快速判断：若支撑域不是矩形（如三角形、圆形），则\textbf{几乎一定不独立}。

\textbf{第三问：求概率 $P\{(X,Y)\in D\}$}
\[
\iint_D f(x,y)\,\mathrm{d}x\,\mathrm{d}y\quad\text{（画图确定积分区域，选合适的积分次序）}
\]

\textbf{第四问：求 $Z=g(X,Y)$ 的分布}

用 CDF 法：$F_Z(z)=P\{g(X,Y)\leq z\}=\displaystyle\iint_{\{(x,y):g(x,y)\leq z\}}f(x,y)\,\mathrm{d}x\,\mathrm{d}y$，再对 $z$ 求导。
\end{note}


\chapter{随机变量的数字特征}

\begin{note}
\textbf{自学导读：为什么需要"数字特征"？}

前面三章你学会了用分布函数、分布律或密度函数\textbf{完整}描述一个随机变量。但完整描述有时太"重"了——实际中我们经常只需要几个关键数字就能做决策：
\begin{itemize}
    \item "这门课平均分多少？"——需要\textbf{期望}（中心位置）
    \item "成绩波动大不大？"——需要\textbf{方差}（离散程度）
    \item "数学成绩好的人，物理成绩是否也好？"——需要\textbf{协方差/相关系数}（线性关联强度）
\end{itemize}

\textbf{本章学习路线：}
\begin{enumerate}
    \item 期望 $EX$：用第 2 章的分布律或密度函数，加权平均得到"重心"；
    \item 方差 $DX$：衡量偏离重心的平均程度，计算公式 $DX=EX^2-(EX)^2$ 是核心工具；
    \item 协方差 $\text{Cov}(X,Y)$ 和相关系数 $\rho_{XY}$：用第 3 章的联合分布，衡量两个变量的线性联动。
\end{enumerate}

\textbf{关键提醒：} 期望和方差的\textbf{性质}（线性性、独立时可加性）比计算公式更重要——后面第 5--9 章几乎每一步推导都在调用这些性质。
\end{note}

{\small
\[
\begin{cases}
\text{数学期望}
\begin{cases}
\text{计算公式}
\begin{cases}
\text{离散型}
\begin{cases}
\text{一维:}
\begin{cases}
EX = \sum\limits_{i} x_i p_i \\
E\varphi(X) = \sum\limits_{i} \varphi(x_i) p_i
\end{cases} \\
\text{二维: } E[\varphi(X,Y)] = \sum\limits_{i=1}^{m} \sum\limits_{j=1}^{n} \varphi(x_i,y_j) p_{ij}
\end{cases} \\
\text{连续型}
\begin{cases}
\text{一维:}
\begin{cases}
EX = \int_{-\infty}^{+\infty} x f(x) dx \\
E\varphi(X) = \int_{-\infty}^{+\infty} \varphi(x) f(x) dx
\end{cases} \\
\text{二维: } E[\varphi(X,Y)] = \int_{-\infty}^{+\infty} \int_{-\infty}^{+\infty} \varphi(x,y) f(x,y) dxdy
\end{cases}
\end{cases} \\
\text{性质}
\begin{cases}
(1)\ \text{常数的期望不变: } EC=C,\ \text{特别地: } E(EX)=EX \\
(2)\ \text{线性性质: 如果 } EX \text{ 和 } EY \text{ 都存在, 则 } E(aX+bY)=aEX+bEY \\
(3)\ \text{如果 } EX \text{ 存在, 则 } E(kX)=kEX \\
(4)\ \text{若 } X,Y \text{ 相互独立, 且期望都存在, 则 } E(XY)=EXEY
\end{cases}
\end{cases} \\
\text{方差}
\begin{cases}
\text{计算公式: } DX=E(X-EX)^2=E(X^2)-(EX)^2 \\
\text{性质}
\begin{cases}
(1)\ \text{常数的方差为零: } DC=0,~~~~ (2) DX\geqslant0 \Rightarrow E(X^2)\geqslant (EX)^2 \\
(3)\ D(aX+b)=a^2 DX,~~~ (4)\ D(X\pm Y)=DX+DY\pm 2\text{Cov}(X,Y) \\
(5)\ \text{若 } X,Y \text{ 相互独立, 则 } D(X\pm Y)=DX+DY \\
(6)\ DX=E(X-EX)^2\leqslant E(X-c)^2
\end{cases}
\end{cases} \\
\text{协方差}
\begin{cases}
\text{计算公式: } \text{Cov}(X,Y)=E(X-EX)(Y-EY)=E(XY)-EXEY \\
\text{其中 } E(XY)=
\begin{cases}
\sum\limits_{i}\sum\limits_{j} x_i y_j P\{X=x_i,Y=y_j\} \quad (\text{离散型}) \\
\int_{-\infty}^{+\infty}\int_{-\infty}^{+\infty} xy f(x,y) dxdy \quad (\text{连续型})
\end{cases} \\
\text{性质}
\begin{cases}
(1)\ \text{Cov}(X,X)=DX ,~~~~~(2)\ \text{Cov}(X,Y)=\text{Cov}(Y,X) \\
(3)\ \text{Cov}(X,C)=0 \ (C \text{是常数}) \\
(4)\ \text{Cov}(aX+bY,Z)=a\text{Cov}(X,Z)+b\text{Cov}(Y,Z)
\end{cases}
\end{cases} \\
\text{相关系数}
\begin{cases}
\text{计算公式: } \rho_{XY}=\dfrac{\text{Cov}(X,Y)}{\sqrt{DX}\sqrt{DY}},\ |\rho_{XY}|\leqslant1 \\
\text{性质}
\begin{cases}
(1)\ |\rho_{XY}|=1 \Leftrightarrow P\{Y=aX+b\}=1 \\
(2)\ \rho_{XY}=0 \Leftrightarrow X,Y \text{ 不相关} \Leftrightarrow
\begin{cases}
(1)\ \text{Cov}(X,Y)=0 \\
(2)\ E(XY)=EXEY \\
(3)\ D(X\pm Y)=DX+DY
\end{cases} \\
(3)\ \text{若 } (X,Y)\sim N(\mu_1,\mu_2,\sigma_1^2,\sigma_2^2,\rho), \\
\qquad \text{则 } X,Y \text{ 独立与 } X,Y \text{ 不相关等价}
\end{cases}
\end{cases}
\end{cases}
\]
}

\begin{introduction}
    \item 数学期望的定义与性质~\ref{def:expectation}~\ref{prop:expectation}
    \item 方差的定义与性质~\ref{def:variance}~\ref{prop:variance}
    \item 常见分布的数学期望与方差~\ref{tab:common-ev}
    \item 协方差~\ref{def:covariance}~\ref{prop:covariance}
    \item 相关系数~\ref{def:correlation}~\ref{prop:correlation}
    \item 独立性与不相关性~\ref{prop:independence-uncorrelated}
    \item 矩（原点矩、中心矩、混合矩）~\ref{def:moments}
    \item 切比雪夫不等式~\ref{thm:chebyshev}
\end{introduction}

\begin{note}
\textbf{使用提醒：} 本章最稳妥的学习顺序是“先看中心位置，再看离散程度，最后看变量之间是否联动”。具体地，数学期望回答“平均落在哪里”，方差回答“围绕中心散得多不多”，协方差与相关系数回答“两个变量是否同涨同跌、线性联系有多强”。做题时若先分清“单变量还是双变量、中心还是波动、精确计算还是估计界”，后面的公式就不容易混用。
\end{note}

\section{数学期望}

\subsection{数学期望的定义}

\begin{definition}[数学期望] \label{def:expectation}
设 \(X\) 是一个随机变量。
\begin{enumerate}
    \item \textbf{离散型}：若 \(X\) 的分布律为 \(P\{X=x_i\} = p_i\ (i=1,2,\dots)\)，且级数 \(\sum\limits_{i=1}^{\infty} x_i p_i\) 绝对收敛，则称 \(X\) 的\textbf{数学期望}存在，其值为
    \[
    EX = \sum_{i=1}^{\infty} x_i p_i.
    \]
    若级数不绝对收敛，则称 \(X\) 的数学期望不存在。
    \item \textbf{连续型}：若 \(X\) 的概率密度为 \(f(x)\)，且积分 \(\int_{-\infty}^{+\infty} x f(x)\,\mathrm{d}x\) 绝对收敛，则称 \(X\) 的\textbf{数学期望}存在，其值为
    \[
    EX = \int_{-\infty}^{+\infty} x f(x)\,\mathrm{d}x.
    \]
\end{enumerate}
\end{definition}

\begin{exercise}\label{exer:4-3-cauchy}
设连续型随机变量 \(X\) 的概率密度为
\[
f(x) = \frac{1}{\pi(1+x^2)}, \quad -\infty < x < +\infty,
\]
则 \(EX\)（\quad）。~~A. 等于 0\quad B. 等于 1\quad C. 等于 \(\pi\)\quad D. 不存在
\end{exercise}
\begin{solution}
\textbf{题型识别：} 这是标准的柯西分布 $X\sim C(0,1)$，需要判断期望是否存在。

\textbf{判断依据：} 期望存在的前提是 $\displaystyle\int_{-\infty}^{+\infty} |x|\,f(x)\,\mathrm{d}x$ 收敛（绝对收敛）。

\textbf{验证：}
\[
\int_{-\infty}^{+\infty} \frac{|x|}{\pi(1+x^2)}\,\mathrm{d}x = \frac{2}{\pi}\int_0^{+\infty}\frac{x}{1+x^2}\,\mathrm{d}x = \frac{2}{\pi}\cdot\left.\frac{1}{2}\ln(1+x^2)\right|_0^{+\infty} = +\infty.
\]
绝对值积分发散，故 $EX$ 不存在。注意不能因为被积函数 $\dfrac{x}{\pi(1+x^2)}$ 是奇函数就断言"正负抵消等于 0"——期望的定义要求绝对收敛，而这里绝对值积分已经发散。

答案选 \textbf{D}。
\end{solution}

\begin{exercise}\label{exer:4-s1-linear}
已知连续型随机变量 \(X\) 与 \(Y\) 有相同的概率密度，且
\[
X \sim f(x) = \begin{cases} 2x\theta^2, & 0 < x < \dfrac{1}{\theta}\ (\theta > 0), \\[4pt] 0, & \text{其他}, \end{cases}
\quad E[a(X+2Y)] = \frac{1}{\theta},
\]
则 \(a =\,\underline{\hspace{2cm}}\)。~~~A. \(\dfrac{2}{3}\)\quad B. \(\dfrac{1}{2}\)\quad C. \(\dfrac{1}{3}\)\quad D. \(\dfrac{1}{6}\)
\end{exercise}
\begin{solution}
\textbf{题型识别：} 同分布随机变量的期望线性组合问题。

\textbf{依据：} $X$ 与 $Y$ 同分布 $\Rightarrow$ $EX = EY$，再利用期望的线性性 $E[a(X+2Y)] = a(EX + 2EY) = 3a\,EX$。

\textbf{计算 $EX$：}
\[
EX = \int_0^{1/\theta} x \cdot 2x\theta^2\,\mathrm{d}x = 2\theta^2\int_0^{1/\theta} x^2\,\mathrm{d}x = 2\theta^2\cdot\frac{1}{3\theta^3} = \frac{2}{3\theta}.
\]

\textbf{代入求 $a$：}
\[
E[a(X+2Y)] = 3a\cdot\frac{2}{3\theta} = \frac{2a}{\theta} = \frac{1}{\theta},
\]
解得 $a = \dfrac{1}{2}$。答案选 \textbf{B}。
\end{solution}


\begin{exercise}\label{exer:4-s8-poisson}
设随机变量 \(X\) 服从参数为 \(\lambda\) 的泊松分布，且已知 \(E[(X-1)(X-2)] = 1\)，则 \(\lambda =\,\underline{\hspace{2cm}}\)。
\end{exercise}
\begin{solution}
\textbf{题型识别：} 泊松分布的高阶矩计算，利用 $EX = DX = \lambda$ 将展开式化为关于 $\lambda$ 的方程。

\textbf{依据：} $E(X^2) = DX + (EX)^2 = \lambda + \lambda^2$。

\textbf{展开并代入：}
\[
E[(X-1)(X-2)] = E(X^2 - 3X + 2) = E(X^2) - 3EX + 2 = (\lambda + \lambda^2) - 3\lambda + 2 = \lambda^2 - 2\lambda + 2.
\]
令其等于 1：
\[
\lambda^2 - 2\lambda + 2 = 1 \implies (\lambda - 1)^2 = 0,
\]
解得 $\lambda = 1$。
\end{solution}

\begin{exercise}\label{exer:4-5-poisson}
设随机变量 \(X\) 的概率分布为 \(P\{X=k\} = \dfrac{C}{k!},\ k=0,1,2,\dots\)，则 \(E(|X - EX|) =\,\underline{\hspace{2cm}}\)。
\end{exercise}
\begin{solution}
应填 2。由 \(\sum\limits_{k=0}^{\infty} P\{X=k\} = \sum\limits_{k=0}^{\infty} \dfrac{C}{k!} = Ce = 1\)，得 \(C = e^{-1}\)，故 \(X \sim P(1)\)，\(EX = 1\)。
\[
\begin{aligned}
E(|X - EX|) &= E(|X - 1|) = \sum_{k=0}^{\infty} P(X=k)|k-1| \\
&= P(X=0) \cdot 1 + \sum_{k=2}^{\infty} P(X=k)(k-1) \\
&= e^{-1} + \sum_{k=2}^{\infty} \frac{e^{-1}}{k!}(k-1) = e^{-1} + e^{-1}\sum_{j=1}^{\infty}\frac{j}{(j+1)!} \\
&= e^{-1} + e^{-1}\sum_{j=1}^{\infty}\left(\frac{1}{j!} - \frac{1}{(j+1)!}\right) = e^{-1} + e^{-1}\cdot 1 = 2.
\end{aligned}
\]
\end{solution}

\begin{note}
由泊松分布的分布列 \(P\{X=k\} = \dfrac{\lambda^k e^{-\lambda}}{k!}\) 与本题的 \(P\{X=k\} = \dfrac{C}{k!}\) 比较可知，\(X\) 服从参数 \(\lambda = 1\) 的泊松分布。直接利用 \(EX = \lambda = 1\) 可简化计算。
\end{note}


\begin{exercise}\label{exer:4-s11-machine}
一部机器一天内发生故障的概率为 0.2，机器发生故障时全天停产。若一周 5 个工作日里：无故障获利 10 万元，发生一次故障获利 5 万元，发生两次故障获利 0 元，发生三次及以上故障亏损 2 万元。求一周内利润的期望。
\end{exercise}

\begin{solution}
设一周内故障次数为 \(X\)，则 \(X \sim B(5, 0.2)\)。

\[
P\{X=0\} = 0.8^5 = 0.32768,\quad P\{X=1\} = \mathrm{C}_5^1 \times 0.2 \times 0.8^4 = 0.4096,
\]
\[
P\{X=2\} = \mathrm{C}_5^2 \times 0.04 \times 0.512 = 0.2048,\quad P\{X \geqslant 3\} = 1 - \sum_{k=0}^{2}P\{X=k\} = 0.05792.
\]

设利润为 \(Y\)（万元），则
\[
EY = 10 \times 0.32768 + 5 \times 0.4096 + 0 \times 0.2048 + (-2) \times 0.05792 = 5.20896 \text{（万元）}.
\]
\end{solution}

\begin{exercise}\label{exer:4-20-shop}
某商店销售某商品，销量 \(X\) 在 \((50,100)\) 上服从均匀分布。若每销售一单位获利 500 元；如果需求量大于进货量 \(n\)，可以从其他部门调剂，此时每单位获利 300 元；如果有积压，则每单位亏损 200 元。进货量 \(n\) 为多少时平均获利最大？最大为多少？
\end{exercise}
\begin{solution}
设进货量为 \(n\)（\(n \in [50,100]\)），获利为
\[
Y = \begin{cases} 500n + 300(X-n) = 300X + 200n, & X \geqslant n, \\[4pt] 500X - 200(n-X) = 700X - 200n, & X < n. \end{cases}
\]
\[
\begin{aligned}
EY &= \frac{1}{50}\left[\int_{50}^{n}(700x - 200n)\,\mathrm{d}x + \int_{n}^{100}(300x + 200n)\,\mathrm{d}x\right] \\
&= \frac{1}{50}\left[350(n^2-2500) - 200n(n-50) + 150(10000-n^2) + 200n(100-n)\right] \\
&= \frac{1}{50}\left[-200n^2 + 30000n + 625000\right] \\
&= -4n^2 + 600n + 12500.
\end{aligned}
\]
令 \((EY)' = -8n + 600 = 0\)，得 \(n = 75\)。此时
\[
EY_{\max} = -4 \times 75^2 + 600 \times 75 + 12500 = 35000.
\]
\end{solution}

\begin{exercise}\label{exer:4-trade}
假定国际市场每年对我国某种商品的需求量 \(X\)（单位：吨）服从 \([2000, 4000]\) 上的均匀分布。每售出一吨获利 3 万美元，若销售不出则每吨需仓储费 1 万美元。问：每年组织多少资源，才能使期望收益最大？
\end{exercise}
\begin{solution}
设进货量为 \(y\)，收益为
\[
L(y,X) = \begin{cases} 3y, & X \geqslant y, \\[4pt] 3X - (y - X) = 4X - y, & X < y. \end{cases}
\]
\[
E[L] = \frac{1}{2000}\left[\int_{2000}^{y}(4x - y)\,\mathrm{d}x + \int_{y}^{4000}3y\,\mathrm{d}x\right]
= \frac{1}{2000}(-2y^2 + 14000y - 8000000).
\]
令 \(\dfrac{\mathrm{d}E[L]}{\mathrm{d}y} = \dfrac{1}{2000}(-4y + 14000) = 0\)，得 \(y = 3500\) 吨。
\end{solution}

\begin{exercise}\label{exer:4-basketball}
某决赛采取五局三胜制（三局先胜后比赛终止），双方每局胜率相同。目前比分 1:0（对方领先）。若每场比赛总收入 160 万元，胜方分得 120 万，败方分得 40 万，求对方收入的期望。
\end{exercise}
\begin{solution}
先看状态：对方已经 1 胜，还差 2 胜；我们还差 3 胜。把后面最多 4 局补满以后，对方在 4 局中至少赢 2 局，就一定能先达到 3 胜。因此
\[
P\{\text{对方夺冠}\} = \mathrm{C}_4^2\!\left(\frac{1}{2}\right)^{\!4} + \mathrm{C}_4^3\!\left(\frac{1}{2}\right)^{\!4} + \mathrm{C}_4^4\!\left(\frac{1}{2}\right)^{\!4} = \frac{11}{16}.
\]
但题目问的是“收入”，不能只看最终谁夺冠，而要把每一局的 120 万/40 万累计起来。

设 \(E(a,b)\) 表示从当前状态出发、对方还差 \(a\) 胜、我方还差 \(b\) 胜时，对方在\textbf{后续比赛}中的收入期望。则有递推关系
\[
E(a,b)=\frac12\bigl(120+E(a-1,b)\bigr)+\frac12\bigl(40+E(a,b-1)\bigr),\qquad E(0,b)=E(a,0)=0.
\]
顺次计算：
\[
E(1,1)=80,\quad E(1,2)=120,\quad E(2,1)=120,\quad E(1,3)=140,\quad E(2,2)=200.
\]
因此
\[
E(2,3)=\frac12(120+140)+\frac12(40+200)=250.
\]
又因为第一局已经打完且对方获胜，所以这 120 万元是确定收入。故对方总收入的期望为
\[
120+250=370\text{（万元）}.
\]
\end{solution}

\begin{definition}[随机变量函数的期望] \label{def:function-expectation}
设 \(Y = g(X)\)（\(g\) 为连续函数）。
\begin{enumerate}
    \item \textbf{离散型}：若 \(\sum\limits_{i=1}^{\infty} g(x_i) p_i\) 绝对收敛，则
    \[
    E[g(X)] = \sum_{i=1}^{\infty} g(x_i) p_i.
    \]
    \item \textbf{连续型}：若 \(\int_{-\infty}^{+\infty} g(x) f(x)\,\mathrm{d}x\) 绝对收敛，则
    \[
    E[g(X)] = \int_{-\infty}^{+\infty} g(x) f(x)\,\mathrm{d}x.
    \]
\end{enumerate}
\end{definition}

\begin{note}
(1) 数学期望又称为概率平均值，常简称期望或均值。它是描述随机变量平均取值状况的数字特征，刻画了随机变量一切可能值的集中位置。\par
(2) 定义中要求级数（或积分）\textbf{绝对收敛}，这是为了保证期望的值不依赖于求和（或积分）的顺序。\par
(3) 并非所有随机变量都存在数学期望。例如柯西分布 \(X\sim f(x)=\dfrac{1}{\pi(1+x^2)}\)，由于
\[
\int_{-\infty}^{+\infty}\frac{|x|}{\pi(1+x^2)}\,\mathrm{d}x = +\infty,
\]
其数学期望不存在。
\end{note}



\begin{exercise}\label{exer:4-1-2x}
设随机变量 \(X \sim B(2, 0.1)\)，则 \(E(2^X) =\,\underline{\hspace{2cm}}\)。
\end{exercise}

\begin{solution}
由于 \(X\sim B(2,0.1)\)，所以 \(X\) 只能取 \(0,1,2\) 三个值，并且
\[
P(X=0)=0.9^2,\qquad
P(X=1)=\mathrm{C}_2^1\cdot 0.1\cdot 0.9,\qquad
P(X=2)=0.1^2.
\]

按离散型随机变量期望的定义，
\[
E(2^X)
=\sum_x 2^xP(X=x)
=0.9^2 \times 2^0+\mathrm{C}_2^1 \times 0.1 \times 0.9 \times 2^1+0.1^2 \times 2^2.
\]

逐项计算得
\[
E(2^X)=0.81+0.36+0.04=1.21.
\]

因此应填 \(1.21\)。
\end{solution}

\begin{exercise}\label{exer:4-s14-cauchymin}
设连续型随机变量 \(X\) 的概率密度为 \(f(x) = \dfrac{1}{\pi(1+x^2)}\ (-\infty < x < +\infty)\)，求 \(E(\min(|X|,1))\)。
\end{exercise}

\begin{solution}
由对称性，
\[
\begin{aligned}
E(\min(|X|,1)) &= 2\int_0^{+\infty}\min(x,1)\,f(x)\,\mathrm{d}x = 2\left[\int_0^1 x f(x)\,\mathrm{d}x + \int_1^{+\infty} f(x)\,\mathrm{d}x\right] \\
&= 2\left(\int_0^1 \frac{x}{\pi(1+x^2)}\,\mathrm{d}x + \int_1^{+\infty}\frac{1}{\pi(1+x^2)}\,\mathrm{d}x\right) \\
&= \frac{1}{\pi}\ln(1+x^2)\Big|_0^1 + \frac{2}{\pi}\arctan x\Big|_1^{+\infty} = \frac{\ln 2}{\pi} + \frac{1}{2}.
\end{aligned}
\]
\end{solution}

\begin{exercise}\label{exer:4-s7-rayleigh}
已知连续型随机变量 \(Y\) 的概率密度为
\[
f(y) = \begin{cases} \dfrac{y}{a^2}\,e^{-\frac{y^2}{2a^2}}, & y > 0,\ a > 0, \\[4pt] 0, & y \leqslant 0, \end{cases}
\]
求 \(Z = \dfrac{1}{Y}\) 的数学期望 \(EZ\)。
\end{exercise}

\begin{solution}
\[
\begin{aligned}
EZ &= E\!\left(\frac{1}{Y}\right) =0+ \int_0^{+\infty}\frac{1}{y}\cdot\frac{y}{a^2}\,e^{-\frac{y^2}{2a^2}}\,\mathrm{d}y = \frac{1}{a^2}\int_0^{+\infty}e^{-\frac{y^2}{2a^2}}\,\mathrm{d}y.
\end{aligned}
\]
利用正态分布密度的归一性：\(N(0,a^2)\) 的密度 \(\varphi(y) = \dfrac{1}{\sqrt{2\pi}\,a}\,e^{-\frac{y^2}{2a^2}}\) 满足 \(\displaystyle\int_{-\infty}^{+\infty}\varphi(y)\,\mathrm{d}y = 1\)，由对称性
\[
\int_0^{+\infty}\frac{1}{\sqrt{2\pi}\,a}\,e^{-\frac{y^2}{2a^2}}\,\mathrm{d}y = \frac{1}{2},
\]
故 \(\displaystyle\int_0^{+\infty}e^{-\frac{y^2}{2a^2}}\,\mathrm{d}y = \frac{\sqrt{2\pi}\,a}{2}\)，于是
\[
EZ = \frac{1}{a^2}\cdot\frac{\sqrt{2\pi}\,a}{2} = \frac{\sqrt{2\pi}}{2a}.
\]
\end{solution}


\begin{exercise}\label{exer:4-dy-seg}
设随机变量 \(X\) 在 \([-1,2]\) 上服从均匀分布，令
\[
Y = \begin{cases} 1, & X > 0, \\ 0, & X = 0, \\ -1, & X < 0, \end{cases}
\]
则 \(DY =\,\underline{\hspace{2cm}}\)。
\end{exercise}
\begin{solution}
\(X \sim U(-1,2)\)，\(f(x) = \dfrac{1}{3}\ (-1 < x < 2)\)。
\begin{align*}
P\{Y=1\} &= P\{X>0\} = \int_0^2 \frac{1}{3}\,\mathrm{d}x = \frac{2}{3}\\
P\{Y=0\} &= P\{X=0\} = 0\\
P\{Y=-1\} &= P\{X<0\} = \int_{-1}^0 \frac{1}{3}\,\mathrm{d}x = \frac{1}{3}\\
EY &= 1 \times \frac{2}{3} + (-1) \times \frac{1}{3} = \frac{1}{3},\quad E(Y^2) = 1^2 \times \frac{2}{3} + (-1)^2 \times \frac{1}{3} = 1\\
DY &= E(Y^2) - (EY)^2 = 1 - \frac{1}{9} = \frac{8}{9}.
\end{align*}
\begin{note}
\textbf{易错点：} \(X=0\) 的概率为 0，所以 \(Y=0\) 也是零概率事件；算期望和方差时真正起作用的是 \(X>0\) 与 \(X<0\) 两块区域。
\end{note}
\end{solution}

\begin{exercise}\label{exer:4-s9-maxbvn}
已知随机变量 \((X,Y) \sim N(0,0;\,1,1;\,0.5)\)，\(U = \max\{X,Y\}\)，\(V = \min\{X,Y\}\)，求 \(E(U+V)\) 和 \(E(UV)\)。
\end{exercise}
\begin{solution}
根据二维正态分布的参数记号 \((X,Y) \sim N(\mu_1, \mu_2;\, \sigma_1^2, \sigma_2^2;\, \rho)\)，由题意可知：
\begin{itemize}
    \item 期望：$E(X) = 0,\quad E(Y) = 0$
    \item 方差：$D(X) = 1,\quad D(Y) = 1$
    \item 相关系数：$\rho = 0.5$
\end{itemize}
因为 $U = \max\{X,Y\}$，$V = \min\{X,Y\}$，对于任意一对实数取值 $X$ 和 $Y$，无论哪个数较大，$U$ 和 $V$ 必然一个是 $X$ 另一个是 $Y$。因此它们的和与乘积在变换前后保持不变，即有如下代数恒等式：
\[
U + V = X + Y, \quad U \cdot V = X \cdot Y.
\]
(1) \textbf{求 $\boldsymbol{E(U+V)}$}：\\
利用数学期望的线性性质，直接代入即可：
\[
E(U+V) = E(X+Y) = E(X) + E(Y) = 0 + 0 = 0.
\]
(2) \textbf{求 $\boldsymbol{E(UV)}$}：\\
要求 $E(UV)$，即求 $E(XY)$。首先，根据相关系数与协方差的关系公式计算协方差：
\[
\mathrm{Cov}(X,Y) = \rho \cdot \sqrt{D(X)} \sqrt{D(Y)} = 0.5 \times \sqrt{1} \times \sqrt{1} = 0.5.
\]
又根据协方差的计算公式 $\mathrm{Cov}(X,Y) = E(XY) - E(X)E(Y)$，可得：
\[
E(UV) = E(XY) = \mathrm{Cov}(X,Y) + E(X)E(Y) = 0.5 + 0 \times 0 = 0.5.
\]
综上所述，$E(U+V) = 0$，$E(UV) = 0.5$。
\end{solution}


\begin{exercise}
设随机变量 $X$ 服从参数为 $0.5$ 的泊松分布，试求随机变量 $Y = X/(1+X)$ 的数学期望 $E[Y]$.
\end{exercise}

\begin{solution}
由于 $X \sim P(\lambda)$，随机变量 $X$ 的概率质量函数为：
$$P(X = k) = \frac{\lambda^k e^{-\lambda}}{k!}, \quad k = 0, 1, 2, \dots$$

根据离散型随机变量函数的期望公式，展开 $E[Y]$：
$$E[Y] = E\left[\frac{X}{1+X}\right] = \sum_{k=0}^{\infty} \frac{k}{k+1} \frac{\lambda^k e^{-\lambda}}{k!} = \sum_{k=0}^{\infty} \left(1 - \frac{1}{k+1}\right) \frac{\lambda^k e^{-\lambda}}{k!}$$

将其拆分为两部分：
$$E[Y] = \sum_{k=0}^{\infty} \frac{\lambda^k e^{-\lambda}}{k!} - \sum_{k=0}^{\infty} \frac{1}{k+1} \frac{\lambda^k e^{-\lambda}}{k!}$$

第一部分为所有概率之和，由全概率公式可知等于 1。第二部分化简如下：
$$\sum_{k=0}^{\infty} \frac{1}{k+1} \frac{\lambda^k e^{-\lambda}}{k!} = \sum_{k=0}^{\infty} \frac{\lambda^k e^{-\lambda}}{(k+1)!} = \frac{e^{-\lambda}}{\lambda} \sum_{k=0}^{\infty} \frac{\lambda^{k+1}}{(k+1)!}$$

令 $j = k+1$，当 $k$ 从 $0$ 到 $\infty$ 时，$j$ 从 $1$ 到 $\infty$：
$$\frac{e^{-\lambda}}{\lambda} \sum_{j=1}^{\infty} \frac{\lambda^j}{j!} = \frac{e^{-\lambda}}{\lambda} (e^\lambda - 1) = \frac{1 - e^{-\lambda}}{\lambda}$$

因此，$E[Y]$ 的通用表达式为：
$$E[Y] = 1 - \frac{1 - e^{-\lambda}}{\lambda}$$

代入给定的参数 $\lambda = 0.5$：
$$E[Y] = 1 - \frac{1 - e^{-0.5}}{0.5} = 1 - 2(1 - e^{-0.5}) = 2e^{-0.5} - 1$$

故随机变量 $Y$ 的数学期望为 $2e^{-0.5} - 1$。
\end{solution}

\vspace{2em}

\begin{exercise}
现有 3 个袋子，各装有 $a$ 个白球和 $b$ 个黑球。先从第 1 个袋子中摸出一球，记下颜色后把它放入第 2 个袋子中；再从第 2 个袋子中摸出一球，记下颜色后把它放入第 3 个袋子中；最后从第 3 个袋子中摸出一球，记下颜色。记这 3 次摸球中所得的白球总数为 $X$，求 $E[X]$.
\end{exercise}
\begin{solution}
引入指示变量 $X_i$（$i = 1, 2, 3$）：
若第 $i$ 次摸出的是白球，则 $X_i = 1$；若第 $i$ 次摸出的是黑球，则 $X_i = 0$。
这 3 次摸球中摸出的白球总数 $X$ 可以表示为：
$$X = X_1 + X_2 + X_3$$
根据数学期望的线性性质：
$$E[X] = E[X_1 + X_2 + X_3] = E[X_1] + E[X_2] + E[X_3]$$
因为 $X_i$ 只能取 0 或 1，所以 $E[X_i] = 1 \cdot P(X_i = 1) + 0 \cdot P(X_i = 0) = P(X_i = 1)$。
第一次摸球（从第 1 个袋子）：
袋中有 $a$ 个白球，$b$ 个黑球。
$$P(X_1 = 1) = \frac{a}{a+b}$$
第二次摸球（从第 2 个袋子）：
利用全概率公式，考虑第一次摸出的是白球还是黑球：
\begin{align*}
P(X_2 = 1) &= P(X_2 = 1 \mid X_1 = 1)P(X_1 = 1) + P(X_2 = 1 \mid X_1 = 0)P(X_1 = 0) \\
&= \left(\frac{a+1}{a+b+1}\right)\left(\frac{a}{a+b}\right) + \left(\frac{a}{a+b+1}\right)\left(\frac{b}{a+b}\right) \\
&= \frac{a(a+1) + ab}{(a+b+1)(a+b)} = \frac{a(a+b+1)}{(a+b+1)(a+b)} = \frac{a}{a+b}
\end{align*}
第三次摸球（从第 3 个袋子）：
同理，根据全概率公式，由于从第 2 个袋子摸出白球的无条件概率仍为 $\frac{a}{a+b}$，过程具有完全的对称性。因此，第三次摸出白球的概率依然不变：
$$P(X_3 = 1) = \frac{a}{a+b}$$
将上述三次的期望加总：
$$E[X] = E[X_1] + E[X_2] + E[X_3] = \frac{a}{a+b} + \frac{a}{a+b} + \frac{a}{a+b} = \frac{3a}{a+b}$$
故这 3 次摸球中所得的白球总数 $X$ 的期望 $E[X]$ 为 $\frac{3a}{a+b}$。
\end{solution}


\subsection{二维随机变量函数的数学期望}

\begin{definition}[二维随机变量函数的期望] \label{def:2d-expectation}
设 \((X,Y)\) 为二维随机变量，\(g(X,Y)\) 为连续函数。
\begin{enumerate}
    \item \textbf{离散型}：若 \((X,Y)\) 的联合分布律为 \(p_{ij} = P\{X=x_i, Y=y_j\}\ (i,j=1,2,\dots)\)，且级数 \(\sum\limits_{i}\sum\limits_{j} |g(x_i,y_j)|\,p_{ij}\) 收敛，则
    \[
    E[g(X,Y)] = \sum_{i=1}^{\infty}\sum_{j=1}^{\infty} g(x_i, y_j)\,p_{ij}.
    \]
    \item \textbf{连续型}：若 \((X,Y)\) 的联合概率密度为 \(f(x,y)\)，且积分 \(\int_{-\infty}^{+\infty}\int_{-\infty}^{+\infty} |g(x,y)|\,f(x,y)\,\mathrm{d}x\mathrm{d}y\) 收敛，则
    \[
    E[g(X,Y)] = \int_{-\infty}^{+\infty}\int_{-\infty}^{+\infty} g(x,y)\,f(x,y)\,\mathrm{d}x\mathrm{d}y.
    \]
\end{enumerate}
\end{definition}

\subsection{数学期望的性质}

\begin{property}\label{prop:expectation}
数学期望具有以下性质（假设所涉及的期望均存在）：
\begin{enumerate}
    \item \(E(C) = C\)，其中 \(C\) 为任意常数。特别地，\(E(EX) = EX\)。
    \item \(E(kX) = kEX\)，其中 \(k\) 为任意常数。
    \item \textbf{线性性质}：\(E(aX + bY) = aEX + bEY\)。更一般地，
    \[
    E\!\left(\sum_{i=1}^{n} a_i X_i\right) = \sum_{i=1}^{n} a_i EX_i.
    \]
    \item 若 \(X, Y\) 相互独立，则 \(E(XY) = EX \cdot EY\)。
\end{enumerate}
\end{property}
\begin{proof}[性质 (4) 的证明]
以连续型为例。\(X, Y\) 独立时联合密度 \(f(x,y) = f_X(x)\,f_Y(y)\)，故
\begin{align*}
E(XY) &= \iint_{\mathbb{R}^2} xy\,f(x,y)\,\mathrm{d}x\,\mathrm{d}y
= \int_{-\infty}^{+\infty} x\,f_X(x)\,\mathrm{d}x \cdot \int_{-\infty}^{+\infty} y\,f_Y(y)\,\mathrm{d}y \\
&= EX \cdot EY.
\end{align*}
离散型类似（将积分改为求和）。
\end{proof}

\begin{note}
(1) 性质 4 的逆命题不成立，即由 \(E(XY) = EX \cdot EY\) 不能推出 \(X\) 与 \(Y\) 独立。\par
(2) 在计算期望时，如果积分区间关于原点对称，应当考虑被积函数的奇偶性，这往往能简化计算。\par
(3) 实数域上的 Gamma 函数定义为 \(\Gamma(\alpha) = \int_0^{+\infty} t^{\alpha-1}e^{-t}\,\mathrm{d}t\)，其有以下性质：
\[
\Gamma\!\left(\frac{1}{2}\right) = \sqrt{\pi},\quad \Gamma(x+1) = x\Gamma(x),\quad \Gamma(n) = (n-1)!.
\]
在计算正态分布的各阶矩时，利用 Gamma 函数能够降低计算量。
\end{note}


\begin{exercise}\label{exer:4-6-fx}
设随机变量 \(X\) 的概率密度为 \(f(x) = \begin{cases} \frac12x, & 0 < x < 2, \\ 0, & \text{其他}, \end{cases}\) \(F(x)\) 为 \(X\) 的分布函数，则 \(P\{F(X) > EX - 1\} =\,\underline{\hspace{2cm}}\)。
\end{exercise}

\begin{solution}
应填 \(\dfrac23\)。

先算右端的阈值 \(EX-1\)。由密度
\[
f(x)=\frac12x,\qquad 0<x<2,
\]
有
\[
EX=\int_0^2 x\cdot \frac12x\,\mathrm dx
=\frac12\int_0^2 x^2\,\mathrm dx
=\frac12\cdot\frac{8}{3}
=\frac43.
\]
因此 \(EX-1=\frac43-1=\frac13\)。

接下来处理 \(F(X)\)。对于连续型随机变量，分布函数变换 \(F(X)\) 服从 \(U(0,1)\)，所以令 \(Y=F(X)\)，则 \(Y\sim U(0,1)\)。

于是原概率化为
\[
P\{F(X)>EX-1\}=P\left(Y>\frac13\right).
\]
而 \(Y\sim U(0,1)\)，故
\[
P\left(Y>\frac13\right)=1-\frac13=\frac23.
\]

所以答案是 \(\dfrac23\)。
\end{solution}
\begin{note}
本题若不知 \(Y = F(X) \sim U(0,1)\) 的结论，可直接计算：
\[
F(x) = \begin{cases} 0, & x \leqslant 0, \\ \dfrac{x^2}{4}, & 0 < x < 2, \\ 1, & x \geqslant 2. \end{cases}
\]
于是 \(P\!\left(F(X) > \dfrac{1}{3}\right) = P\!\left(\dfrac{X^2}{4} > \dfrac{1}{3},\ 0 \leqslant X < 2\right) = P\!\left(X > \dfrac{2}{\sqrt{3}}\right) = \displaystyle\int_{\frac{2}{\sqrt{3}}}^{2} x\,\mathrm{d}x = \dfrac{2}{3}\)。
\end{note}


\begin{exercise}\label{exer:4-6-exy}
设随机变量 \(X, Y\) 相互独立，且 \(X\) 服从区间 \([0,2]\) 上的均匀分布，\(Y\) 服从参数为 3 的指数分布，则 \(E(XY) =\,\underline{\hspace{2cm}}\)。
\end{exercise}
\begin{solution}
由于 \(X,Y\) 相互独立，所以可以直接使用 \(E(XY)=EX\cdot EY\)。又 \(X\sim U(0,2)\)，故 \(EX=\frac{0+2}{2}=1\)；\(Y\) 服从参数为 3 的指数分布，因此 \(EY=\frac{1}{3}\)。

于是
\[
E(XY) = EX \cdot EY = 1 \times \frac{1}{3} = \frac{1}{3}.
\]

所以应填 \(\dfrac13\)。
\end{solution}


\begin{exercise}\label{exer:4-8-wait}
甲、乙两人相约于 12:00\textasciitilde 13:00 会面。设 \(X, Y\) 分别是甲、乙到达超过 12:00 的时间（单位：小时），且相互独立，已知
\[
f_X(x) = \begin{cases} 3x^2, & 0 < x < 1, \\ 0, & \text{其他}, \end{cases}\quad
f_Y(y) = \begin{cases} 2y, & 0 < y < 1, \\ 0, & \text{其他}, \end{cases}
\]
求先到达者需要等待的时间的数学期望。
\end{exercise}
\begin{solution}
所求为 \(E(|X-Y|)\)。由独立性，\(f(x,y) = 6x^2y\ (0 < x,y < 1)\)，故
\[
\begin{aligned}
E(|X-Y|) &= \int_0^1\!\!\int_0^1 |x-y| \cdot 6x^2 y\,\mathrm{d}x\mathrm{d}y \\
&= \iint_{x>y} (x-y) \cdot 6x^2 y\,\mathrm{d}x\mathrm{d}y + \iint_{x<y} (y-x) \cdot 6x^2 y\,\mathrm{d}x\mathrm{d}y \\
&= \int_0^1\!\!\int_0^x 6x^3y - 6x^2y^2\,\mathrm{d}y\,\mathrm{d}x + \int_0^1\!\!\int_0^y 6x^2y^2 - 6x^3y\,\mathrm{d}x\,\mathrm{d}y \\
&= \int_0^1 \!\left(3x^5 - 2x^5\right)\mathrm{d}x + \int_0^1 \!\left(2y^5 - \frac{3}{2}y^5\right)\mathrm{d}y = \frac{1}{6} + \frac{1}{12} = \frac{1}{4}.
\end{aligned}
\]
\end{solution}


\begin{exercise}\label{exer:4-two-points}
从长度为 3 的线段上任取两个点，求两点间距离 \(D\) 的数学期望和方差。
\end{exercise}
\begin{solution}
设两个点的位置分别为 \(X_1, X_2\)，独立且均服从 \(U(0,3)\)，则距离 \(D = |X_1 - X_2|\)。单个均匀分布 \( U(0,3) \) 的概率密度为：
\[f_{X_i}(x_i) =\begin{cases}\displaystyle \frac{1}{3}, & 0 < x_i < 3 \\0, & \text{其他}\end{cases}\quad (i=1,2)\]
由独立性，联合密度 = 边缘密度的乘积：
  \[
  f(x_1,x_2) = f_{X_1}(x_1) \cdot f_{X_2}(x_2) =
  \begin{cases}
  \displaystyle \frac{1}{3} \cdot \frac{1}{3} = \frac{1}{9}, & 0 < x_1 < 3,\ 0 < x_2 < 3 \\
  0, & \text{其他}
  \end{cases}
  \]
由对称性，只需计算 \(x_1 > x_2\) 的区域再乘以 2：
\[
ED = 2\int_0^3 \int_0^{x_1} (x_1 - x_2) \cdot \frac{1}{9}\,\mathrm{d}x_2\,\mathrm{d}x_1
= \frac{2}{9}\int_0^3 \frac{x_1^2}{2}\,\mathrm{d}x_1 = 1.
\]
\[
E(X_i) = \frac{3}{2},\quad E(X_i^2) = DX_i + (EX_i)^2 = \frac{3}{4} + \frac{9}{4} = 3.
\]
\[
E(D^2) = E[(X_1 - X_2)^2] = E(X_1^2) - 2EX_1 EX_2 + E(X_2^2) = 3 - 2 \times \frac{9}{4} + 3 = \frac{3}{2},
\]
\[
DD = ED^2 - (ED)^2 = \frac{3}{2} - 1 = \frac{1}{2}.
\]
\end{solution}

\begin{proposition}[均匀分布次序统计量的通用结论]\label{prop:order-stat-uniform}
对于 \( U(0,1) \) 上的第 \( k \) 个次序统计量 \( X_{(k)} \)，其期望为：
\[
\boldsymbol{E(X_{(k)}) = \frac{k}{n+1}}
\]
当 \( k=1 \)（最小值）：\( E(X_{(1)}) = \frac{1}{n+1} \)； 当 \( k=n \)（最大值）：\( E(X_{(n)}) = \frac{n}{n+1} \)
\end{proposition}
\begin{proof}
对于 \( U(0,1) \)，单个变量的分布函数 \( F(x)=x \ (0\leq x\leq1) \)，密度 \( f(x)=1 \)。
第 \( k \) 个次序统计量 \( X_{(k)} \) 的密度公式为：
\[
f_{X_{(k)}}(x) = \frac{n!}{(k-1)!(n-k)!} [F(x)]^{k-1} [1-F(x)]^{n-k} f(x)
\]
代入 \( F(x)=x, f(x)=1 \)，得：
\[
f_{X_{(k)}}(x) = \frac{n!}{(k-1)!(n-k)!} x^{k-1} (1-x)^{n-k}, \quad 0<x<1
\]
这正是Beta分布 \( \text{Beta}(k, n-k+1) \) 的密度。Beta分布 \( \text{Beta}(\alpha,\beta) \) 的期望为 \( \frac{\alpha}{\alpha+\beta} \)，这里 \( \alpha=k, \beta=n-k+1 \)，因此：
\[
E(X_{(k)}) = \frac{k}{k + (n-k+1)} = \frac{k}{n+1}
\]
\end{proof}

\begin{exercise}\label{exer:4-order-stat}
设相互独立的随机变量 \(X_1, X_2, \dots, X_n\) 均服从 \((0, \theta)\) 上的均匀分布。求
\[
M = \max\{X_1, \dots, X_n\},\qquad
N = \min\{X_1, \dots, X_n\}
\]
的数学期望。
\end{exercise}
\begin{solution}
\(F(x) = \dfrac{x}{\theta}\)，\(f(x) = \dfrac{1}{\theta}\ (0 < x < \theta)\)。\\
\textbf{最大值}：\(F_M(x) = [F(x)]^n = \left(\dfrac{x}{\theta}\right)^n\)，\(f_M(x) = \dfrac{n\,x^{n-1}}{\theta^n}\)，
\[
EM = \int_0^\theta x \cdot \frac{n\,x^{n-1}}{\theta^n}\,\mathrm{d}x = \frac{n}{n+1}\,\theta.
\]
\textbf{最小值}：\(1 - F_N(x) = [1 - F(x)]^n = \left(1 - \dfrac{x}{\theta}\right)^n\)，\(f_N(x) = \dfrac{n}{\theta}\!\left(1 - \dfrac{x}{\theta}\right)^{n-1}\)，
\[
EN = \int_0^\theta x \cdot \frac{n}{\theta}\!\left(1 - \frac{x}{\theta}\right)^{n-1}\mathrm{d}x = \frac{1}{n+1}\,\theta.
\]
\end{solution}



\begin{exercise}\label{exer:4-range}
在区间 \([0,1]\) 上任取 \(n\) 个点 \(X_1, X_2, \dots, X_n\)，记 \[X_{(1)} = \min\{X_1, \dots, X_n\},~~X_{(n)} = \max\{X_1, \dots, X_n\},~~R = X_{(n)} - X_{(1)}\text{（极差）}\]求 \(ER\)。
\end{exercise}
\begin{solution}
设 \( X_1,X_2,\dots,X_n \) 是独立同分布的随机变量，都服从区间 \([0,1]\) 上的均匀分布 \( U(0,1) \)。 定义次序统计量：
\begin{itemize}
    \item 最小值 \( X_{(1)} = \min\{X_1,X_2,\dots,X_n\} \)（第1个次序统计量）
    \item 最大值 \( X_{(n)} = \max\{X_1,X_2,\dots,X_n\} \)（第n个次序统计量）
\end{itemize}
极差 \( R = X_{(n)} - X_{(1)} \)，表示n个点覆盖的区间长度，我们要求它的数学期望 \( ER \)。由次序统计量的期望公式（\(\theta = 1\)）：
\[
E(X_{(n)}) = \frac{n}{n+1},\quad E(X_{(1)}) = \frac{1}{n+1},\quad ER = \frac{n}{n+1} - \frac{1}{n+1} = \frac{n-1}{n+1}.
\]
\end{solution}



\begin{exercise}
设随机变量 $X_1, X_2, X_3$ 独立同分布，且 $X_i (i=1,2,3)$ 的分布列为：$P(X_i = k) = \frac{1}{3} \quad (k=1,2,3)$，求 $Y = \max\{X_1, X_2, X_3\}$ 的数学期望.
\end{exercise}
\begin{solution}
因为 $Y = \max\{X_1, X_2, X_3\}$，所以事件 $\{Y \le k\}$ 等价于所有三个随机变量都小于等于 $k$，即 $\{X_1 \le k, X_2 \le k, X_3 \le k\}$。
由于 $X_1, X_2, X_3$ 独立同分布，我们可以先求出 $Y$ 的分布函数（即累积概率）：
$$P(Y \le k) = P(X_1 \le k, X_2 \le k, X_3 \le k) = P(X_1 \le k)P(X_2 \le k)P(X_3 \le k) = [P(X_1 \le k)]^3$$
根据已知条件，$X_i$ 的分布列为 $P(X_i = k) = \frac{1}{3}$（$k=1,2,3$），我们可以算出 $X_i$ 的累积概率：
\begin{itemize}
\item 当 $k=1$ 时，$P(X_1 \le 1) = P(X_1 = 1) = \frac{1}{3}$
\item 当 $k=2$ 时，$P(X_1 \le 2) = P(X_1 = 1) + P(X_1 = 2) = \frac{2}{3}$
\item 当 $k=3$ 时，$P(X_1 \le 3) = P(X_1 = 1) + P(X_1 = 2) + P(X_1 = 3) = 1$
\end{itemize}
代入前面的公式，得到 $Y$ 的累积概率：
$$P(Y \le 1) = \left(\frac{1}{3}\right)^3 = \frac{1}{27}$$
$$P(Y \le 2) = \left(\frac{2}{3}\right)^3 = \frac{8}{27}$$
$$P(Y \le 3) = 1^3 = 1$$
接下来，我们利用 $P(Y = k) = P(Y \le k) - P(Y \le k-1)$ 来求 $Y$ 的分布列：
$$P(Y = 1) = P(Y \le 1) = \frac{1}{27}$$
$$P(Y = 2) = P(Y \le 2) - P(Y \le 1) = \frac{8}{27} - \frac{1}{27} = \frac{7}{27}$$
$$P(Y = 3) = P(Y \le 3) - P(Y \le 2) = 1 - \frac{8}{27} = \frac{19}{27}$$
最后，根据离散型随机变量的期望公式计算 $Y$ 的数学期望 $E[Y]$：
\begin{align*}
E[Y] &= \sum_{k=1}^{3} k \cdot P(Y = k) = 1 \cdot \frac{1}{27} + 2 \cdot \frac{7}{27} + 3 \cdot \frac{19}{27} = \frac{8}{3}
\end{align*}
\end{solution}

\section{方差}

\begin{note}
\textbf{自学导读：} 期望告诉你"平均落在哪里"，但两个分布可以有相同的期望却差别巨大——比如"稳定月薪 1 万"和"50\% 概率 0 元、50\% 概率 2 万"期望都是 1 万，但风险完全不同。方差就是量化这种"风险/波动"的工具。

\textbf{计算公式的直觉：} $DX = E(X-EX)^2$ 是"偏差的平方的平均"。展开后得到 $DX = EX^2 - (EX)^2$，这个公式几乎是做题的唯一入口——先算 $EX$ 和 $EX^2$，再相减。
\end{note}

\subsection{方差的定义}

\begin{definition}[方差] \label{def:variance}
设 \(X\) 是随机变量，若 \(E[(X - EX)^2]\) 存在，则称其为 \(X\) 的\textbf{方差}，记为 \(DX\)（或 \(\mathrm{Var}(X)\)），即
\[
DX = E[(X - EX)^2] = E(X^2) - (EX)^2.
\]
称 \(\sqrt{DX}\) 为 \(X\) 的\textbf{标准差}（或均方差），记为 \(\sigma(X)\)。

若 \(DX > 0\)，称 \(X^* = \dfrac{X - EX}{\sqrt{DX}}\) 为 \(X\) 的\textbf{标准化随机变量}，此时 \(EX^* = 0\)，\(DX^* = 1\)。
\end{definition}

\begin{note}
从方差的定义可以看出，只有当随机变量的期望存在时，其方差才可能存在。但期望存在的随机变量，其方差也不一定存在。
\end{note}



\begin{exercise}\label{exer:4-s5-dx2}
设随机变量 \(X\) 的分布函数为
\[
F(x) = \begin{cases} 0, & x < -1, \\ 0.2, & -1 \leqslant x < 0, \\ 0.8, & 0 \leqslant x < 1, \\ 1, & x \geqslant 1, \end{cases}
\]
则 \(D(X^2) =\,\underline{\hspace{2cm}}\)。
\end{exercise}

\begin{solution}
由分布函数知 \(X\) 的分布列为
\[
X \sim \begin{pmatrix} -1 & 0 & 1 \\ 0.2 & 0.6 & 0.2 \end{pmatrix}.
\]
于是
\[
E(X^2) = (-1)^2 \times 0.2 + 0^2 \times 0.6 + 1^2 \times 0.2 = 0.4,
\]
\[
E(X^4) = (-1)^4 \times 0.2 + 0^4 \times 0.6 + 1^4 \times 0.2 = 0.4,
\]
\[
D(X^2) = E(X^4) - [E(X^2)]^2 = 0.4 - 0.16 = 0.24.
\]
\end{solution}

\begin{exercise}
设随机变量 $X$ 的分布密度函数为
$$f(x) = \begin{cases} 2x, & 0 < x < 1, \\ 0, & \text{其他.} \end{cases}$$
求随机变量 $X$ 的方差.
\end{exercise}

\begin{solution}
根据连续型随机变量方差的计算公式 $D(X) = E[X^2] - (E[X])^2$，我们需要先分别求出 $X$ 的数学期望 $E[X]$ 和 $X^2$ 的数学期望 $E[X^2]$。
首先，计算随机变量 $X$ 的数学期望 $E[X]$：
\begin{align*}
E[X] &= \int_{-\infty}^{+\infty} x f(x) dx = \int_{0}^{1} x \cdot (2x) dx = \int_{0}^{1} 2x^2 dx = \left[ \frac{2}{3}x^3 \right]_0^1 = \frac{2}{3}
\end{align*}
接着，计算 $E[X^2]$：
\begin{align*}
E[X^2] &= \int_{-\infty}^{+\infty} x^2 f(x) dx = \int_{0}^{1} x^2 \cdot (2x) dx = \int_{0}^{1} 2x^3 dx = \left[ \frac{2}{4}x^4 \right]_0^1 = \frac{1}{2}
\end{align*}
最后，代入方差公式计算 $D(X)$：
\begin{align*}
D(X) &= E[X^2] - (E[X])^2 = \frac{1}{2} - \left(\frac{2}{3}\right)^2 = \frac{1}{2} - \frac{4}{9} = \frac{9}{18} - \frac{8}{18} = \frac{1}{18}
\end{align*}
\end{solution}

\begin{exercise}\label{exer:4-s6-dexp}
设连续型随机变量 \(X\) 服从参数为 1 的指数分布，则 \(D(e^{-2X}) =\,\underline{\hspace{2cm}}\)。
\end{exercise}
\begin{solution}
\textbf{题型识别：} 求随机变量函数 $g(X) = e^{-2X}$ 的方差，其中 $X \sim E(1)$。

\textbf{依据：} $D[g(X)] = E[g(X)^2] - \{E[g(X)]\}^2 = E(e^{-4X}) - [E(e^{-2X})]^2$。

\textbf{分步计算：} $X \sim E(1)$ 的密度为 $f(x) = e^{-x}$（$x > 0$），故
\[
E(e^{-2X}) = \int_0^{+\infty} e^{-2x}\,e^{-x}\,\mathrm{d}x = \int_0^{+\infty} e^{-3x}\,\mathrm{d}x = \frac{1}{3},
\]
\[
E(e^{-4X}) = \int_0^{+\infty} e^{-4x}\,e^{-x}\,\mathrm{d}x = \int_0^{+\infty} e^{-5x}\,\mathrm{d}x = \frac{1}{5}.
\]

\textbf{代入方差公式：}
\[
D(e^{-2X}) = \frac{1}{5} - \left(\frac{1}{3}\right)^2 = \frac{1}{5} - \frac{1}{9} = \frac{4}{45}.
\]
\end{solution}


\begin{exercise}\label{exer:4-1-binomial}
设一次试验成功的概率为 \(p\)，进行 100 次独立重复试验，则成功次数的标准差的最大值为\,\underline{\hspace{2cm}}。
\end{exercise}
\begin{solution}
\textbf{题型识别：} 二项分布标准差的极值问题。

\textbf{依据：} 成功次数 $X \sim B(100,p)$，则 $DX = 100p(1-p)$，标准差 $\sqrt{DX} = 10\sqrt{p(1-p)}$。

\textbf{求最大值：} 由均值不等式，$p(1-p) \leqslant \left(\dfrac{p + (1-p)}{2}\right)^2 = \dfrac{1}{4}$，当且仅当 $p = \dfrac{1}{2}$ 时等号成立。故标准差最大值为
\[
10\sqrt{\frac{1}{4}} = 5.
\]
\end{solution}

\begin{exercise}\label{exer:4-2-cos}
设随机变量 \(X\) 的概率密度为
\[
f(x) = \begin{cases} \dfrac{1}{2}\sin x, & 0 \leqslant x \leqslant \pi, \\[4pt] 0, & \text{其他}, \end{cases}
\]
对 \(X\) 独立观察 4 次，用 \(Y\) 表示观察值大于 \(\dfrac{\pi}{3}\) 的次数，求 \(E(Y^2)\)。
\end{exercise}
\begin{solution}
先验证单次观察中“\(X>\dfrac{\pi}{3}\)”发生的概率。由题给密度，
\[
P\!\left(X > \frac{\pi}{3}\right) = \int_{\frac{\pi}{3}}^{\pi} \frac{1}{2}\sin x\,\mathrm{d}x
= \frac12\bigl[-\cos x\bigr]_{\pi/3}^{\pi}
= \frac12\left(1+\frac12\right)
= \frac34.
\]

因此每次观察都可视为一次成功概率为 \(\dfrac34\) 的伯努利试验；又因为 4 次观察相互独立，所以
\[
Y\sim B\!\left(4,\frac34\right).
\]
于是
\[
EY=np=4\times \frac34=3,\qquad DY=np(1-p)=4\times \frac34\times \frac14=\frac34.
\]
再利用恒等式 \(E(Y^2)=DY+[EY]^2\)，得到
\[
E(Y^2)=\frac34+3^2=\frac34+9=\frac{39}{4}.
\]

所以 \(E(Y^2)=\frac{39}{4}\)。
\end{solution}



\subsection{方差的性质}

\begin{property}\label{prop:variance}
方差具有以下性质（假设所涉及的方差均存在）：
\begin{enumerate}
    \item \(DX \geqslant 0\)，从而 \(E(X^2) = DX + (EX)^2 \geqslant (EX)^2\)。
    \item \(D(C) = 0\)，其中 \(C\) 为常数。特别地，\(D(EX) = 0\)。
    \item \(DX = 0 \Leftrightarrow\) 存在常数 \(a\)，使得 \(P\{X = a\} = 1\)。
    \item \(D(aX + b) = a^2 DX\)，其中 \(a, b\) 为常数。
    \item 对任意两个随机变量 \(X, Y\)，有
    \[
    D(X \pm Y) = DX + DY \pm 2\,\mathrm{Cov}(X,Y).
    \]
    更一般地，
    \[
    D\!\left(\sum_{i=1}^{n} a_i X_i\right) = \sum_{i=1}^{n} a_i^2 DX_i + 2\sum_{1 \leqslant i < j \leqslant n} a_i a_j\,\mathrm{Cov}(X_i,X_j).
    \]
    \item 若 \(X, Y\) 相互独立，则 \(D(X \pm Y) = DX + DY\)，且
    \[
    D(XY) = DX \cdot DY + DX \cdot (EY)^2 + DY \cdot (EX)^2 \geqslant DX \cdot DY.
    \]
    一般地，若 \(X_1, X_2, \dots, X_n\) 相互独立，\(g_i\) 为连续函数，则
    \[
    D\!\left(\sum_{i=1}^{n} g_i(X_i)\right) = \sum_{i=1}^{n} D[g_i(X_i)].
    \]
    \item 对任意常数 \(c\)，有 \(DX = E[(X - EX)^2] \leqslant E[(X - c)^2]\)。
\end{enumerate}
\end{property}

\begin{note}
\textbf{自学追问：} 方差的核心计算公式是哪个？是 $DX = E(X^2) - (EX)^2$。几乎所有方差题的第一步都是先算 $EX$ 和 $E(X^2)$，再用这个公式做减法。如果题目给的是分布律或密度，就套定义求期望；如果给的是已知分布（如二项、泊松、正态），直接查表代入 $DX$ 的公式即可。

\textbf{使用提醒：} 性质 5 中 $D(X \pm Y) = DX + DY \pm 2\,\mathrm{Cov}(X,Y)$ 是最常考的方差公式。若 $X,Y$ 独立，交叉项为 $0$，方差直接相加；若不独立，必须先算协方差。做题时一定先确认"是否独立"，再决定要不要加交叉项。

\textbf{易错点：} $D(aX+b) = a^2 DX$，常数 $b$ 对方差无贡献，但系数 $a$ 要平方。另外，$D(X-Y) = DX + DY - 2\,\mathrm{Cov}(X,Y)$，注意减号在协方差前面，不是方差前面——方差永远是加号。
\end{note}

\begin{proof}[性质 6 中 \(D(XY)\) 的证明]
因为 \(X, Y\) 相互独立，故 \(X^2, Y^2\) 也相互独立，即有
\[
E(XY) = EX \cdot EY, \quad E(X^2 Y^2) = E(X^2) \cdot E(Y^2),
\]
于是
\[
\begin{aligned}
D(XY) &= E(X^2 Y^2) - [E(XY)]^2 = E(X^2)E(Y^2) - (EX)^2(EY)^2 \\
&= [DX + (EX)^2][DY + (EY)^2] - (EX)^2(EY)^2 \\
&= DX \cdot DY + (EX)^2 \cdot DY + (EY)^2 \cdot DX.
\end{aligned}
\]
又 \((EX)^2 \cdot DY + (EY)^2 \cdot DX \geqslant 0\)，所以 \(D(XY) \geqslant DX \cdot DY\)。
\end{proof}
\begin{proof}[性质 (1)–(5) 的证明]
(1) \(DX = E[(X-EX)^2] \geqslant 0\)。由 \(DX \geqslant 0\)，\(E(X^2) = DX + (EX)^2 \geqslant (EX)^2\)。\par
(2) \(D(C) = E[(C-C)^2] = 0\)。\par
(4) \(D(aX+b) = E[(aX+b-aEX-b)^2] = E[a^2(X-EX)^2] = a^2 DX\)。\par
(5) \(D(X\pm Y) = E[(X\pm Y - EX \mp EY)^2] = E[(X-EX)^2] + E[(Y-EY)^2] \pm 2E[(X-EX)(Y-EY)] = DX + DY \pm 2\mathrm{Cov}(X,Y)\)。特别地，若 \(X,Y\) 独立，则 \(\mathrm{Cov}(X,Y)=0\)，故 \(D(X\pm Y)=DX+DY\)。
\end{proof}

\begin{proof}[性质 7 的证明]
令 \(g(c) = E[(X - c)^2] = E(X^2) - 2cEX + c^2\)。由
\[
g'(c) = -2EX + 2c \stackrel{\text{令}}{=} 0,
\]
得 \(c = EX\)。又 \(g''(c) = 2 > 0\)，故 \(c = EX\) 是 \(g(c)\) 的最小值点。
\end{proof}

\begin{proposition}[均值的期望与方差]\label{prop:sample-mean-moment}
    若总体 $X$ 的期望为 $\mu$，方差为 $\sigma^2$，$X_1,\dots,X_n$ 是简单随机样本，样本均值为 $\overline{X}$，则：$$\boldsymbol{E(\overline{X}) = \mu}, \quad \boldsymbol{D(\overline{X}) = \frac{\sigma^2}{n}}$$
\end{proposition}
\begin{note}
这条结论只要求样本是独立同分布、总体具有有限方差，并不要求总体正态。它常用来说明样本均值是 \(\mu\) 的无偏估计，且样本量越大，样本均值的波动越小；但若题目进一步要求 \(\overline{X}\) 的精确分布，就还需要额外条件，不能只凭这两个矩结论直接下结论。
\end{note}
\begin{proof} 
    \[ E(\overline{X}) = E\!\left(\frac{1}{n}\sum_{i=1}^{n}X_i\right) = \frac{1}{n}\sum_{i=1}^{n}EX_i = \frac{1}{n}\cdot n\mu = \mu. \] 
    由于简单随机样本中 \(X_1, \dots, X_n\) 相互独立，故 \[ D(\overline{X}) = D\!\left(\frac{1}{n}\sum_{i=1}^{n}X_i\right) = \frac{1}{n^2}\sum_{i=1}^{n}DX_i = \frac{1}{n^2}\cdot n\sigma^2 = \frac{\sigma^2}{n}. \] 
\end{proof}

\begin{exercise}\label{exer:4-8-exye}
对于任意两个随机变量 \(X\) 和 \(Y\)，若 \(E(XY) = EXEY\)，则（\quad）。\par
A. \(X,Y\) 相互独立\quad B. \(X,Y\) 不相互独立\quad C. \(D(XY) = DX \cdot DY\)\quad D. \(D(X+Y) = DX + DY\)
\end{exercise}
\begin{solution}
应选 D。

题设
\[
E(XY)=EX\cdot EY
\]
等价于
\[
\mathrm{Cov}(X,Y)=E(XY)-EX\cdot EY=0.
\]
这说明 \(X,Y\) \textbf{不相关}，但并不能推出它们独立。

现在逐项判断：

A “\(X,Y\) 相互独立”不一定成立，因为不相关不等于独立，故 A 错。

B “\(X,Y\) 不相互独立”也不能必然成立，因为独立时当然也可能满足 \(E(XY)=EXEY\)，所以 B 也错。

C “\(D(XY)=DX\cdot DY\)”没有这样的普遍结论。即便 \(X,Y\) 独立，也应有
\[
D(XY)=DX\cdot DY+DX(EY)^2+DY(EX)^2,
\]
故 C 错。

D 由方差公式
\[
D(X+Y)=DX+DY+2\mathrm{Cov}(X,Y)
\]
立刻得到 \(D(X+Y)=DX+DY\)。
因此 D 正确。
\end{solution}

\begin{exercise}\label{exer:4-9-dxy}
已知 \(D(X) = 4\)，\(D(Y) = 1\)，\(D(X+Y) = 4\)，则 \(D(X-Y) =\,\underline{\hspace{2cm}}\)。
\end{exercise}
\begin{solution}
应填 \(6\)。

先由已知的 \(D(X+Y)\) 反求协方差。根据公式
\[
D(X+Y)=DX+DY+2\mathrm{Cov}(X,Y),
\]
代入题设 \(DX=4\)、\(DY=1\)、\(D(X+Y)=4\)，得
\[
4=4+1+2\mathrm{Cov}(X,Y)\implies 2\mathrm{Cov}(X,Y)=-1,\qquad \mathrm{Cov}(X,Y)=-\frac12.
\]

再用差的方差公式
\[
D(X-Y)=DX+DY-2\mathrm{Cov}(X,Y),
\]
于是
\[
D(X-Y)=4+1-2\left(-\frac12\right)=5+1=6.
\]

因此应填 \(6\)。
\end{solution}

\begin{exercise}\label{exer:4-10-suff}
设随机变量 \(X\) 和 \(Y\) 的方差存在且不等于 0，则 \(D(X+Y) = D(X) + D(Y)\) 是 $X$ 与 $Y$（\quad）。\par
A. 不相关的充分不必要条件\quad B. 不相关的充分必要条件\quad\\
C. 独立的充分不必要条件\quad D. 独立的充分必要条件
\end{exercise}
\begin{solution}
应选 B。

这道题要判断的是“\(D(X+Y)=DX+DY\)”与“不相关”之间到底是什么关系。最核心的公式是
\[
D(X+Y)=DX+DY+2\mathrm{Cov}(X,Y).
\]

\textbf{先证充分性}：若 \(D(X+Y)=DX+DY\)，代入上式得
\[
2\mathrm{Cov}(X,Y)=0 \implies \mathrm{Cov}(X,Y)=0.
\]
由于题目已说明 \(DX,DY\) 都存在且不为 0，所以这就等价于 \(X,Y\) 不相关。

\textbf{再证必要性}：若 \(X,Y\) 不相关，则 \(\mathrm{Cov}(X,Y)=0\)，再代回方差公式，立即得到 \(D(X+Y)=DX+DY\)。

因此，“\(D(X+Y)=DX+DY\)”既能推出“不相关”，反过来“不相关”也能推出它，所以这是\textbf{不相关的充分必要条件}。

故答案选 B。
\end{solution}


\begin{exercise}\label{exer:4-4-normal}
设连续型随机变量 \(X\) 的概率密度为
\[
f(x) = \frac{1}{\sqrt{\pi}}\,e^{-x^2 + 2x - 1}, \quad -\infty < x < +\infty,
\]
则 \(EX \cdot \sqrt{DX} =\,\underline{\hspace{2cm}}\)。
\end{exercise}
\begin{solution}
\textbf{题型识别：} 从密度函数中识别正态分布的参数。

\textbf{方法：} 把指数部分凑成 $-\dfrac{(x-\mu)^2}{2\sigma^2}$ 的标准形式，再与系数 $\dfrac{1}{\sqrt{2\pi}\,\sigma}$ 比对。

\textbf{配方：}
\[
f(x) = \frac{1}{\sqrt{\pi}}\,e^{-x^2+2x-1} = \frac{1}{\sqrt{\pi}}\,e^{-(x-1)^2} = \frac{1}{\sqrt{2\pi}\cdot\frac{1}{\sqrt{2}}}\exp\!\left\{-\frac{(x-1)^2}{2\cdot\left(\frac{1}{\sqrt{2}}\right)^2}\right\}.
\]

\textbf{读参数：} 与 $N(\mu,\sigma^2)$ 的密度比对，得 $\mu = 1$，$\sigma^2 = \dfrac{1}{2}$，即 $X \sim N\!\left(1,\dfrac{1}{2}\right)$。

\textbf{结论：} $EX = 1$，$\sqrt{DX} = \dfrac{1}{\sqrt{2}}$，故 $EX \cdot \sqrt{DX} = \dfrac{1}{\sqrt{2}}$。
\end{solution}

\begin{exercise}\label{exer:4-7-square}
设 \(X, Y\) 是随机变量，相关系数为 0.4，\(EX = 2\)，\(DX = 25\)，\(EY = 1\)，\(DY = 16\)，则 \(E(2X - 3Y + 4)^2 =\,\underline{\hspace{2cm}}\)。
\end{exercise}

\begin{solution}
先计算各分量：
\[
EX^2 = DX + (EX)^2 = 29,\quad EY^2 = DY + (EY)^2 = 17,
\]
\[
EXY = EX \cdot EY + \mathrm{Cov}(X,Y) = 2 + 0.4 \times 5 \times 4 = 10.
\]
于是
\[
\begin{aligned}
E(2X-3Y+4)^2 &= E(4X^2 + 9Y^2 - 12XY + 16 + 16X - 24Y) \\
&= 4 \times 29 + 9 \times 17 - 12 \times 10 + 16 + 16 \times 2 - 24 \times 1 \\
&= 116 + 153 - 120 + 16 + 32 - 24 = 173.
\end{aligned}
\]
\end{solution}


\begin{exercise}\label{exer:4-3-abs}
设随机变量 \((X,Y)\) 服从区域 \(G = \{(x,y) \mid 0 \leqslant x \leqslant 2,\ 0 \leqslant y \leqslant 2\}\) 上的均匀分布，\(U = |X-Y|\)。\par
(1) 求 \(U\) 的概率密度 \(f_U(u)\)；\par
(2) 求 \(EU\) 和 \(DU\)。
\end{exercise}
\begin{solution}
(1) \(F_U(u) = P(|X-Y| \leqslant u)\)。当 \(u \leqslant 0\) 时，\(F_U(u) = 0\)；当 \(u \geqslant 2\) 时，\(F_U(u) = 1\)。\\
当 \(0 < u < 2\) 时，区域 \(|x-y| \leqslant u\) 在 \([0,2]^2\) 内的面积为 \(4 - 2 \times \dfrac{1}{2}(2-u)^2 = 4u - u^2\)，故
\[
F_U(u) = \frac{4u - u^2}{4} = \frac{(4-u)u}{4}.
\]
因此
\[
f_U(u) = F_U'(u) = \begin{cases} \dfrac{2-u}{2}, & 0 < u < 2, \\[4pt] 0, & \text{其他}. \end{cases}
\]
(2) \(EU = \displaystyle\int_0^2 u \cdot \frac{2-u}{2}\,\mathrm{d}u = \frac{2}{3}\)；
\(EU^2 = \displaystyle\int_0^2 u^2 \cdot \frac{2-u}{2}\,\mathrm{d}u = \frac{2}{3}\)；
\(DU = EU^2 - (EU)^2 = \dfrac{2}{3} - \dfrac{4}{9} = \dfrac{2}{9}\)。
\end{solution}

\section{常见分布的数学期望与方差}

\begin{table}[htbp]
    \centering
    \caption{常见分布的数学期望与方差} \label{tab:common-ev}
    \begin{tabular}{cccc}
        \toprule
        分布 & 分布律或概率密度 & 数学期望 & 方差 \\
        \midrule
        \(0\text{-}1\) 分布 & \(P\{X=k\} = p^k(1-p)^{1-k},\ k=0,1\) & \(p\) & \(p(1-p)\) \\[6pt]
        \(B(n,p)\) & \(P\{X=k\} = \mathrm{C}_n^k p^k(1-p)^{n-k},\ k=0,1,\dots,n\) & \(np\) & \(np(1-p)\) \\[6pt]
        \(P(\lambda)\) & \(P\{X=k\} = \dfrac{\lambda^k e^{-\lambda}}{k!},\ k=0,1,\dots\) & \(\lambda\) & \(\lambda\) \\[6pt]
        \(G(p)\) & \(P\{X=k\} = (1-p)^{k-1}p,\ k=1,2,\dots\) & \(\dfrac{1}{p}\) & \(\dfrac{1-p}{p^2}\) \\[6pt]
        \(U(a,b)\) & \(f(x) = \dfrac{1}{b-a},\ a < x < b\) & \(\dfrac{a+b}{2}\) & \(\dfrac{(b-a)^2}{12}\) \\[6pt]
        \(E(\lambda)\) & \(f(x) = \lambda e^{-\lambda x},\ x > 0\) & \(\dfrac{1}{\lambda}\) & \(\dfrac{1}{\lambda^2}\) \\[6pt]
        \(N(\mu,\sigma^2)\) & \(f(x) = \dfrac{1}{\sqrt{2\pi}\,\sigma}\exp\!\left\{-\dfrac{(x-\mu)^2}{2\sigma^2}\right\}\) & \(\mu\) & \(\sigma^2\) \\[6pt]
        $\chi^2(n)$ & （概率密度不作要求） & \(n\) & \(2n\) \\
        \bottomrule
    \end{tabular}
\end{table}

\begin{note}
表中仅列出各分布概率密度（或分布律）的非零区域。
\end{note}

\begin{proof}[\(0\text{-}1\) 分布]
\(P\{X=1\}=p,\ P\{X=0\}=1-p\)。
\[EX = 1\cdot p + 0\cdot(1-p) = p.\]
\[EX^2 = 1^2\cdot p + 0^2\cdot(1-p) = p.\]
\[DX = EX^2 - (EX)^2 = p - p^2 = p(1-p).\]
\end{proof}

\begin{proof}[\(B(n,p)\)]
设 \(X_i\ (i=1,\dots,n)\) 相互独立，均服从 \(0\text{-}1\) 分布，则 \(X = X_1+X_2+\cdots+X_n \sim B(n,p)\)。
\[EX = \sum_{i=1}^{n}EX_i = np,\qquad DX = \sum_{i=1}^{n}DX_i = np(1-p).\]
\end{proof}

\begin{proof}[\(P(\lambda)\)]
\[EX = \sum_{k=0}^{\infty}k\cdot\frac{\lambda^k e^{-\lambda}}{k!} = \sum_{k=1}^{\infty}\frac{\lambda^k e^{-\lambda}}{(k-1)!} = \lambda e^{-\lambda}\sum_{k=0}^{\infty}\frac{\lambda^k}{k!} = \lambda e^{-\lambda}\cdot e^{\lambda} = \lambda.\]
\[EX^2 = \sum_{k=0}^{\infty}k^2\cdot\frac{\lambda^k e^{-\lambda}}{k!} = \sum_{k=1}^{\infty}k\cdot\frac{\lambda^k e^{-\lambda}}{(k-1)!} = \sum_{k=1}^{\infty}[(k-1)+1]\frac{\lambda^k e^{-\lambda}}{(k-1)!}\]
\[= \sum_{k=2}^{\infty}(k-1)\frac{\lambda^k e^{-\lambda}}{(k-1)!} + \sum_{k=1}^{\infty}\frac{\lambda^k e^{-\lambda}}{(k-1)!} = \lambda^2 e^{-\lambda}\sum_{k=2}^{\infty}\frac{\lambda^{k-2}}{(k-2)!} + \lambda = \lambda^2 + \lambda.\]
\[DX = EX^2 - (EX)^2 = \lambda^2 + \lambda - \lambda^2 = \lambda.\]
\end{proof}

\begin{proof}[\(G(p)\)]
令 \(q=1-p\)，利用幂级数求导技巧：
\[\sum_{k=1}^{\infty}kq^{k-1} = \frac{\mathrm{d}}{\mathrm{d}q}\!\left(\sum_{k=0}^{\infty}q^k\right) = \frac{\mathrm{d}}{\mathrm{d}q}\!\left(\frac{1}{1-q}\right) = \frac{1}{(1-q)^2} = \frac{1}{p^2},
\]
\[
\sum_{k=1}^{\infty}k^2 q^{k-1} = \frac{\mathrm{d}}{\mathrm{d}q}\!\left(\sum_{k=1}^{\infty}kq^k\right) = \frac{\mathrm{d}}{\mathrm{d}q}\!\left(\frac{q}{(1-q)^2}\right) = \frac{1+q}{(1-q)^3} = \frac{2-p}{p^3}.
\]
\[
EX = \sum_{k=1}^{\infty}k\,p\,q^{k-1} = p\cdot\frac{1}{p^2} = \frac{1}{p}.
\]
\[
EX^2 = \sum_{k=1}^{\infty}k^2\,p\,q^{k-1} = p\cdot\frac{2-p}{p^3} = \frac{2-p}{p^2}.
\]
\[
DX = EX^2 - (EX)^2 = \frac{2-p}{p^2} - \frac{1}{p^2} = \frac{1-p}{p^2}.
\]
\end{proof}

\begin{proof}[\(U(a,b)\)]
\[
EX = \int_a^b x\cdot\frac{1}{b-a}\,\mathrm{d}x = \frac{1}{b-a}\cdot\frac{b^2-a^2}{2} = \frac{a+b}{2}.
\]
\[
EX^2 = \int_a^b x^2\cdot\frac{1}{b-a}\,\mathrm{d}x = \frac{1}{b-a}\cdot\frac{b^3-a^3}{3} = \frac{a^2+ab+b^2}{3}.
\]
\[
DX = EX^2 - (EX)^2 = \frac{a^2+ab+b^2}{3} - \frac{(a+b)^2}{4} = \frac{4(a^2+ab+b^2)-3(a^2+2ab+b^2)}{12} = \frac{(b-a)^2}{12}.
\]
\end{proof}

\begin{proof}[\(E(\lambda)\)]
利用分部积分（\(\int_0^\infty x^n e^{-\lambda x}\,\mathrm{d}x = n!/\lambda^{n+1}\)）：
\[
EX = \int_0^\infty x\cdot\lambda e^{-\lambda x}\,\mathrm{d}x = \lambda\cdot\frac{1}{\lambda^2} = \frac{1}{\lambda}.
\]
\[
EX^2 = \int_0^\infty x^2\cdot\lambda e^{-\lambda x}\,\mathrm{d}x = \lambda\cdot\frac{2}{\lambda^3} = \frac{2}{\lambda^2}.
\]
\[
DX = EX^2 - (EX)^2 = \frac{2}{\lambda^2} - \frac{1}{\lambda^2} = \frac{1}{\lambda^2}.
\]
\end{proof}

\begin{proof}[\(\chi^2(n)\)]
由定义，\(\chi^2 = X_1^2+X_2^2+\cdots+X_n^2\)，其中 \(X_i \sim N(0,1)\) 且相互独立。
\[
E(\chi^2) = \sum_{i=1}^{n}E(X_i^2) = n\cdot 1 = n.
\]
\[
D(\chi^2) = \sum_{i=1}^{n}D(X_i^2).
\]
利用 \(X_i \sim N(0,1)\)：\(E(X_i^2)=1,\ E(X_i^4)=3\)（由正态分布的四阶矩公式），故
\[
D(X_i^2) = E(X_i^4) - [E(X_i^2)]^2 = 3 - 1 = 2.
\]
\[
D(\chi^2) = 2n.
\]
\end{proof}

\begin{note}
\textbf{自学检查点——期望与方差：} 学完本节后，快速自测：
\begin{enumerate}
    \item $E(2X+3)=?$（答：$2EX+3$）\quad $D(2X+3)=?$（答：$4DX$）
    \item $X\sim B(n,p)$ 的期望和方差？（答：$np$，$np(1-p)$）
    \item $DX=EX^2-(EX)^2$，如果 $EX=3$，$EX^2=10$，则 $DX=?$（答：$1$）
    \item $X,Y$ 独立时，$D(X-Y)=?$（答：$DX+DY$，注意是加号！）
    \item 为什么样本方差除以 $n-1$ 而不是 $n$？（答：为了使 $E(S^2)=\sigma^2$，即无偏性）
\end{enumerate}
\end{note}


\section{协方差与相关系数}

\begin{note}
\textbf{自学导读：从"单个变量的波动"到"两个变量的联动"。}

期望和方差只描述\textbf{一个}随机变量。但第 3 章告诉我们，两个变量之间可能存在联动关系。协方差和相关系数就是用\textbf{一个数字}来量化这种联动的强弱和方向：
\begin{itemize}
    \item $\text{Cov}(X,Y)>0$：$X$ 大时 $Y$ 倾向于也大（正相关）；
    \item $\text{Cov}(X,Y)<0$：$X$ 大时 $Y$ 倾向于小（负相关）；
    \item $\text{Cov}(X,Y)=0$：$X$ 和 $Y$ 不线性相关（但未必独立！）。
\end{itemize}

\textbf{计算捷径：} 实际做题时，几乎永远用 $\text{Cov}(X,Y)=E(XY)-EX\cdot EY$，而不是定义式。所以关键是会算 $E(XY)$——这需要第 3 章的联合分布。

\textbf{易错点：} 不相关（$\rho=0$）$\not\Rightarrow$ 独立。只有在二维正态分布下，不相关才等价于独立。
\end{note}

\subsection{协方差}

\begin{definition}[协方差] \label{def:covariance}
设随机变量 \(X\) 与 \(Y\) 的方差都存在，称
\[
\mathrm{Cov}(X,Y) = E[(X - EX)(Y - EY)] = E(XY) - EX \cdot EY
\]
为 \(X\) 与 \(Y\) 的\textbf{协方差}。其中
\[
E(XY) = \begin{cases}
\sum\limits_{i}\sum\limits_{j} x_i y_j P\{X=x_i,Y=y_j\} & \text{（离散型）}, \\[8pt]
\int_{-\infty}^{+\infty}\int_{-\infty}^{+\infty} xy\,f(x,y)\,\mathrm{d}x\mathrm{d}y & \text{（连续型）}.
\end{cases}
\]
\end{definition}

\begin{property}\label{prop:covariance}
协方差具有以下性质：
\begin{enumerate}
    \item \(\mathrm{Cov}(X,X) = DX\)。
    \item \(\mathrm{Cov}(X,Y) = \mathrm{Cov}(Y,X)\)（对称性）。
    \item \(\mathrm{Cov}(X,C) = 0\)，其中 \(C\) 为常数。
    \item \(\mathrm{Cov}(aX, bY) = ab\,\mathrm{Cov}(X,Y)\)。
    \item \(\mathrm{Cov}(X_1 + X_2, Y) = \mathrm{Cov}(X_1, Y) + \mathrm{Cov}(X_2, Y)\)（双线性）。更一般地，
    \[
    \mathrm{Cov}\!\left(\sum_{i=1}^{n} a_i X_i,\ \sum_{j=1}^{m} b_j Y_j\right) = \sum_{i=1}^{n}\sum_{j=1}^{m} a_i b_j\,\mathrm{Cov}(X_i, Y_j).
    \]
    \item 若 \(\mathrm{Cov}(X,Y) = 0\)，则 \(D(aX + bY) = a^2 DX + b^2 DY\)。更一般地，\(D(aX + bY) = a^2 DX + b^2 DY+2ab\mathrm{Cov}(X,Y) \)
\end{enumerate}
\end{property}
\begin{proof}
(1)--(4) 由定义 \(\mathrm{Cov}(X,Y)=E(XY)-EX\cdot EY\) 直接展开即可验证。\par
(5) 双线性：利用协方差定义和期望的线性性质：
\[
\mathrm{Cov}(X_1+X_2,Y) = E[(X_1+X_2)Y] - E(X_1+X_2)\,EY = \mathrm{Cov}(X_1,Y)+\mathrm{Cov}(X_2,Y).
\]
反复应用此式即得一般形式。
\end{proof}

\begin{note}
协方差 \(\mathrm{Cov}(X,Y)\) 研究的是 \(X\) 与 \(Y\) 偏离各自均值时，是否倾向于同涨同跌：若二者常常同时偏大或同时偏小，协方差往往为正；若一个偏大时另一个偏小，协方差往往为负。

\textbf{使用提醒：} 协方差更适合判断共同波动的方向，后面在计算 \(D(X+Y)\)、\(D(aX+bY)\) 时会频繁出现。

\textbf{易错点：} 协方差带有量纲，数值大小会受单位影响，不能直接跨题比较“相关强弱”；若要比较线性相关程度，更适合看相关系数。
\end{note}

\begin{exercise}\label{exer:4-bvn-dz}
二维随机变量 \((X,Y) \sim N(1,1;\,4,9;\,-0.5)\)，设 \(Z = \dfrac{1}{2}X - \dfrac{1}{3}Y\)，则 \(DZ =\,\underline{\hspace{2cm}}\)。
\end{exercise}
\begin{solution}
应填 3。已知 \(DX = 4\)，\(DY = 9\)，\(\rho = -0.5\)，则
\[
\mathrm{Cov}(X,Y) = \rho\sqrt{DX}\sqrt{DY} = -0.5 \times 2 \times 3 = -3.
\]
\[
DZ = \frac{1}{4}\,DX + \frac{1}{9}\,DY + 2 \times \frac{1}{2} \times \left(-\frac{1}{3}\right) \mathrm{Cov}(X,Y)
= 1 + 1 + 1 = 3.
\]
\end{solution}

\begin{exercise}\label{exer:4-sample-cov}
设 \(X_1, X_2, \dots, X_5\) 为总体 \(N(1,4)\) 的简单随机样本，\(X_{(5)} = \max\{X_1, X_2, \dots, X_5\}\)。求：(1) \(P(X_{(5)} < 3)\)；(2) \(\mathrm{Cov}(X_1, \bar{X})\)，其中 \(\bar{X} = \dfrac{1}{5}\sum\limits_{i=1}^{5}X_i\)。
\end{exercise}
\begin{solution}
(1) \(P(X_i < 3) = \Phi\!\left(\dfrac{3-1}{2}\right) = \Phi(1)\)，由样本独立性
\[
P(X_{(5)} < 3) = [P(X_1 < 3)]^5 = [\Phi(1)]^5.
\]
(2) 先把样本均值展开：
\[
\bar X=\frac15(X_1+X_2+X_3+X_4+X_5).
\]
于是
\[
\mathrm{Cov}(X_1,\bar X)
=\mathrm{Cov}\!\left(X_1,\frac15\sum_{i=1}^{5}X_i\right)
=\frac15\sum_{i=1}^{5}\mathrm{Cov}(X_1,X_i).
\]

这里最关键的一步是：由于简单随机样本中的 \(X_1,\dots,X_5\) 相互独立，所以
\[
\mathrm{Cov}(X_1,X_i)=0\qquad (i=2,3,4,5),
\]
只剩下
\[
\mathrm{Cov}(X_1,X_1)=DX_1.
\]
因此
\[
\mathrm{Cov}(X_1,\bar X)=\frac15\,DX_1.
\]

又因为总体 \(N(1,4)\) 的方差为 4，所以 \(DX_1=4\)。故
\[
\mathrm{Cov}(X_1,\bar X)=\frac45.
\]
\end{solution}



\begin{exercise}\label{exer:4-joint-xey}
设二维随机变量 \((X,Y)\) 的概率密度函数为
\[
f(x,y) = \begin{cases} xe^{-y}, & 0 < x < y, \\ 0, & \text{其他}, \end{cases}
\]
求：(1) \(Z = X + Y\) 的概率密度；(2) \(N = \min(X,Y)\) 的概率密度；(3) \(EZ\) 和 \(EN\)。
\end{exercise}
\begin{solution}
(1) 由卷积公式，约束 \(0 < x < y\) 且 \(x + y = z\)，即 \(0 < x < z - x\)，故 \(0 < x < z/2\)（\(z > 0\)）：
\begin{align*}
f_Z(z) &= \int_0^{z/2} x\,e^{-(z-x)}\,\mathrm{d}x= e^{-z}\!\int_0^{z/2} x\,e^x\,\mathrm{d}x\\
&= e^{-z}\bigl[e^{z/2}(z/2 - 1) + 1\bigr]= e^{-z/2}\!\left(\frac{z}{2}-1\right) + e^{-z}\quad(z>0).
\end{align*}
(2) 在定义域内始终有 \(0 < X < Y\)，故 \(\min(X,Y) \equiv X\)。
\[f_X(x) = \displaystyle\int_x^{+\infty} x\,e^{-y}\,\mathrm{d}y = x\,e^{-x}\quad(x > 0)\]
所以 \(f_N(n) = n\,e^{-n}\ (n > 0)\)\\
(3) \(EN = EX = \displaystyle\int_0^{+\infty} x^2\,e^{-x}\,\mathrm{d}x = \Gamma(3) = 2\)。\\
\(f_Y(y) = \displaystyle\int_0^y x\,e^{-y}\,\mathrm{d}x = \frac{y^2}{2}\,e^{-y}\ (y > 0)\)，\(EY = \dfrac{1}{2}\,\Gamma(4) = 3\)。\\
\(EZ = EX + EY = 2 + 3 = 5\)。
\end{solution}


\begin{exercise}\label{exer:4-s3-maxmin}
设二维随机变量 \((X,Y)\) 的概率分布为
\[
\begin{array}{c|ccc}
X \backslash Y & 0 & 1 & 2 \\
\hline
-1 & 0.1 & 0.1 & b \\
1 & a & 0.1 & 0.1
\end{array}
\]
若事件 \(\{\max(X,Y)=2\}\) 与事件 \(\{\min(X,Y)=1\}\) 相互独立，则 \(\mathrm{Cov}(X,Y) =\,\underline{\hspace{2cm}}\)。
\end{exercise}

\begin{solution}
首先确定未知参数：
\begin{align*}
P\{\max(X,Y)=2\} &= P\{Y=2\} = b + 0.1,\\
P\{\min(X,Y)=1\} &= P\{X=1,Y=1\} + P\{X=1,Y=2\} = 0.2.
\end{align*}
由独立性：\(P\{X=1,Y=2\} = P\{\max=2\}\,P\{\min=1\}\)，即 \(0.1 = (b+0.1) \times 0.2\)，解得 \(b = 0.4\)。

由分布律之和为 1：\(0.1+0.1+0.4+a+0.1+0.1 = 1\)，解得 \(a = 0.2\)。

计算协方差：
\[
EX = -1 \times 0.6 + 1 \times 0.4 = -0.2,\quad EY = 0 \times 0.3 + 1 \times 0.2 + 2 \times 0.5 = 1.2,
\]
\[
E(XY) = (-1)(2)(0.4) + 1(2)(0.1) = -0.6,
\]
\[
\mathrm{Cov}(X,Y) = -0.6 - (-0.2)(1.2) = -0.36.
\]
\end{solution}

\begin{exercise}\label{exer:4-s12-cov2joint}
设随机变量
\[
X \sim \begin{pmatrix} 0 & 1 \\ \frac{1}{4} & \frac{3}{4} \end{pmatrix},\quad
Y \sim \begin{pmatrix} 0 & 1 \\ \frac{1}{2} & \frac{1}{2} \end{pmatrix},
\]
且 \(\mathrm{Cov}(X,Y) = \dfrac{1}{8}\)，求 \(X\) 与 \(Y\) 的联合分布律。
\end{exercise}

\begin{solution}
\(EX = \dfrac{3}{4}\)，\(EY = \dfrac{1}{2}\)，由
\[
\mathrm{Cov}(X,Y) = E(XY) - EX \cdot EY = E(XY) - \frac{3}{8} = \frac{1}{8},
\]
得 \(E(XY) = \dfrac{1}{2}\)。而 \(E(XY) = P\{X=1,Y=1\}\)，故 \(P\{X=1,Y=1\} = \dfrac{1}{2}\)。
由此：
\begin{align*}
P\{X=1,Y=0\} &= P\{X=1\} - P\{X=1,Y=1\} = \frac{1}{4},\\
P\{X=0,Y=1\} &= P\{Y=1\} - P\{X=1,Y=1\} = 0,\\
P\{X=0,Y=0\} &= P\{X=0\} = \frac{1}{4}.
\end{align*}
联合分布律为
\[
\begin{array}{c|cc|c}
X \backslash Y & 0 & 1 & P_i \\
\hline
0 & \frac{1}{4} & 0 & \frac{1}{4} \\[4pt]
1 & \frac{1}{4} & \frac{1}{2} & \frac{3}{4} \\
\hline
P_j & \frac{1}{2} & \frac{1}{2} & 1
\end{array}
\]
\end{solution}


\subsection{相关系数}

\begin{definition}[相关系数] \label{def:correlation}
设 \(DX > 0\)，\(DY > 0\)，称
\[
\rho_{XY} = \frac{\mathrm{Cov}(X,Y)}{\sqrt{DX \cdot DY}}
\]
为 \(X\) 与 \(Y\) 的\textbf{相关系数}。

若 \(\rho_{XY} = 0\)，则称 \(X\) 与 \(Y\)\textbf{不相关}；若 \(\rho_{XY} > 0\)，称 \(X\) 与 \(Y\)\textbf{正相关}；若 \(\rho_{XY} < 0\)，称 \(X\) 与 \(Y\)\textbf{负相关}。
\end{definition}

\begin{property}\label{prop:correlation}
相关系数具有以下性质：
\begin{enumerate}
    \item \(|\rho_{XY}| \leqslant 1\)。
    \item \(|\rho_{XY}| = 1 \Leftrightarrow\) 存在常数 \(a \neq 0, b\)，使得 \(P\{Y = aX + b\} = 1\)。具体地：
    \begin{itemize}
        \item \(\rho_{XY} = 1 \Leftrightarrow P\{Y = aX + b\} = 1\ (a > 0)\)；
        \item \(\rho_{XY} = -1 \Leftrightarrow P\{Y = aX + b\} = 1\ (a < 0)\)。
    \end{itemize}
    \item \(\rho_{XY} = 0\)（即不相关）的等价条件：
    \begin{itemize}
        \item \(\mathrm{Cov}(X,Y) = 0\)；
        \item \(E(XY) = EX \cdot EY\)；
        \item \(D(X + Y) = DX + DY\)。
    \end{itemize}
    \item 若 \((X,Y) \sim N(\mu_1, \sigma_1^2;\, \mu_2, \sigma_2^2;\, \rho)\)，则 \(X\) 与 \(Y\) 独立等价于 \(X\) 与 \(Y\) 不相关。
\end{enumerate}
\end{property}
\begin{proof}[性质 (1)(2) 的证明]
考虑方差非负性：对任意实数 \(t\)，有
\[
D(Y - tX) = E[(Y-tX - E(Y-tX))^2] \geqslant 0.
\]
展开：
\[
DY - 2t\,\mathrm{Cov}(X,Y) + t^2 DX \geqslant 0, \quad \forall\, t \in \mathbb{R}.
\]
这是一个关于 \(t\) 的非负二次多项式，故判别式 \(\Delta \leqslant 0\)：
\[
4[\mathrm{Cov}(X,Y)]^2 - 4\,DX\cdot DY \leqslant 0 \implies \frac{[\mathrm{Cov}(X,Y)]^2}{DX \cdot DY} \leqslant 1 \implies |\rho_{XY}| \leqslant 1.
\]
等号成立当且仅当 \(\Delta=0\)，即存在 \(t_0\) 使 \(D(Y - t_0 X)=0\)，亦即 \(P\{Y=t_0 X + b\}=1\)（其中 \(b = EY - t_0 EX\)）。取 \(a=t_0\)，即可看出 \(\rho=1\) 时 \(a>0\)，\(\rho=-1\) 时 \(a<0\)。
\end{proof}

\begin{note}
\textbf{直觉理解 $\rho$ 的几何含义：} 把 $(X,Y)$ 的散点图想象成一团点云。
\begin{itemize}
    \item $\rho=+1$：所有点精确落在一条正斜率直线上；
    \item $\rho=-1$：所有点精确落在一条负斜率直线上；
    \item $\rho=0$：点云没有"倾斜"趋势（但可能呈圆形、X 形等非线性图案）；
    \item $0<\rho<1$：点云呈椭圆形，长轴向右上倾斜，$\rho$ 越接近 1 椭圆越"瘦"。
\end{itemize}
\end{note}

\begin{note}
相关系数 \(\rho_{XY}\) 是刻画 \(X\) 与 \(Y\) 之间\textbf{线性相关程度}的度量，本质上就是把协方差做了标准化。\(\rho_{XY} = 0\) 只表示不存在\textbf{线性}关系，但它们之间仍可能存在非线性关系。因此，\textbf{不相关不等于独立}。\footnote{不相关：没有直线关系。独立：什么关系都没有。}当 \(X\) 与 \(Y\) 相互独立时，可推出 \(X\) 与 \(Y\) 不相关；但当 \(X\) 与 \(Y\) 不相关时，\(X\) 与 \(Y\) 却不一定相互独立。

\textbf{易错点：} \(\rho_{XY}=0\) 只能说明“没有线性关系”，不能说明“没有关系”；而 \(|\rho_{XY}|=1\) 则意味着二者几乎处处落在一条直线上。
\end{note}

\begin{property}\label{prop:independence-uncorrelated}
独立性与不相关性的关系：
\begin{enumerate}
    \item \(X\) 与 \(Y\) 相互独立 \(\Rightarrow\) \(X\) 与 \(Y\) 不相关，但反之不一定成立。
    \item 在以下特殊情形中，独立性与不相关性等价：
    \begin{itemize}
        \item \((X,Y)\) 服从二维正态分布；
        \item \(X\) 与 \(Y\) 均服从两点分布。
    \end{itemize}
\end{enumerate}
\end{property}

\begin{note}
\textbf{使用提醒：} 做题时只有在”已知独立”时，才能直接推出协方差为 0 或相关系数为 0；若题目只给出”不相关”，一般不能反推独立，除非额外说明是二维正态等特殊情形。

\textbf{自学追问：} “独立”和”不相关”到底差在哪里？独立是说 $X$ 和 $Y$ 之间\textbf{什么关系都没有}（联合分布完全可拆）；不相关只是说\textbf{没有线性关系}（$\mathrm{Cov}(X,Y)=0$）。独立 $\Rightarrow$ 不相关永远成立，但不相关 $\nRightarrow$ 独立（除非是二维正态或两点分布等特殊情形）。考试中一旦看到”不相关”，脑中应立刻追问”题目是否还给了正态条件？”——有则可升级为独立，无则只能用不相关。
\end{note}

\begin{note}
\textbf{判断题速查口诀：}
\begin{itemize}
    \item 独立 $\Rightarrow$ 不相关 \checkmark（永远成立）
    \item 不相关 $\Rightarrow$ 独立 $\times$（一般不成立）
    \item 不相关 + 二维正态 $\Rightarrow$ 独立 \checkmark
    \item 不相关 + 各自正态（但未说联合正态）$\Rightarrow$ 独立 $\times$
\end{itemize}
\textbf{经典反例（必记）：} 设 $X\sim N(0,1)$，令 $Y=X^2$。则 $\mathrm{Cov}(X,Y)=EX^3-(EX)(EX^2)=0-0=0$，即 $X,Y$ 不相关；但 $Y$ 完全由 $X$ 决定，显然不独立。

\textbf{做题决策树：} 题目问”$X,Y$ 是否独立”时——
\begin{enumerate}
    \item 若已知联合分布 $\Rightarrow$ 直接验证 $f(x,y)=f_X(x)f_Y(y)$；
    \item 若只知不相关 $\Rightarrow$ 检查是否有”二维正态”条件；有则独立，无则不能判断。
\end{enumerate}
\end{note}

\begin{exercise}[多项分布的相关系数] 
袋中有红、白、黑三种颜色球若干，若从袋中任摸一球，摸出的球为红球的概率为 $p_1$，摸出的球为白球的概率为 $p_2$。现从袋中有放回地摸球 $n$ 次，共摸出红球 $X$ 次，摸出白球 $Y$ 次，试求 $X$ 与 $Y$ 的相关系数 $r(X,Y)$。
\end{exercise}

\begin{solution}
由于是有放回抽样，每次摸球是相互独立的，且 $X$ 和 $Y$ 分别服从二项分布：
\[
X \sim B(n, p_1), \quad Y \sim B(n, p_2)
\]
因此它们的方差分别为：
\[
D(X) = np_1(1-p_1), \quad D(Y) = np_2(1-p_2)
\]

为了求 $X$ 与 $Y$ 的协方差 $\text{Cov}(X,Y)$，我们引入示性变量（0-1 变量）：
设 
\[
X_i = \begin{cases}
1, & \text{第 } i \text{ 次摸出红球}, \\
0, & \text{其他},
\end{cases}
\qquad
Y_i = \begin{cases}
1, & \text{第 } i \text{ 次摸出白球}, \\
0, & \text{其他},
\end{cases}
\]
其中 $i = 1, 2, \cdots, n$。则 $X = \sum\limits_{i=1}^n X_i$，$Y = \sum\limits_{j=1}^n Y_j$。

利用协方差的双线性性质：
\[
\text{Cov}(X, Y) = \text{Cov}\left(\sum_{i=1}^n X_i, \sum_{j=1}^n Y_j\right) = \sum_{i=1}^n \sum_{j=1}^n \text{Cov}(X_i, Y_j)
\]

分两种情况讨论 $\text{Cov}(X_i, Y_j)$：
\begin{enumerate}
    \item[(1)] 当 $i \neq j$ 时，第 $i$ 次与第 $j$ 次摸球相互独立，因此 $X_i$ 与 $Y_j$ 独立，$\text{Cov}(X_i, Y_j) = 0$。
    \item[(2)] 当 $i = j$ 时，由于同一次摸球不可能既是红球又是白球，故恒有 $X_i Y_i = 0$，从而 $E(X_i Y_i) = 0$。
    此时：
    \[
    \text{Cov}(X_i, Y_i) = E(X_i Y_i) - E(X_i)E(Y_i) = 0 - p_1 p_2 = -p_1 p_2
    \]
\end{enumerate}

因此，协方差为：
\[
\text{Cov}(X,Y) = \sum_{i=1}^n \text{Cov}(X_i, Y_i) = \sum_{i=1}^n (-p_1 p_2) = -n p_1 p_2
\]

最后，求 $X$ 与 $Y$ 的相关系数：
\begin{align*}
r(X,Y) &= \frac{\text{Cov}(X,Y)}{\sqrt{D(X)} \sqrt{D(Y)}} \\[6pt]
&= \frac{-n p_1 p_2}{\sqrt{np_1(1-p_1)} \sqrt{np_2(1-p_2)}} \\[6pt]
&= \frac{-n p_1 p_2}{n \sqrt{p_1 p_2 (1-p_1)(1-p_2)}} \\[6pt]
&= -\sqrt{\frac{p_1 p_2}{(1-p_1)(1-p_2)}}
\end{align*}

\textbf{另解（方差公式法）}：
令 $Z = X + Y$，表示 $n$ 次摸球中摸出红球或白球的总次数。显然单次摸出红或白的概率为 $p_1 + p_2$，故 $Z \sim B(n, p_1+p_2)$。
其方差为 $D(X+Y) = n(p_1+p_2)(1 - p_1 - p_2)$。
又根据方差的性质：$D(X+Y) = D(X) + D(Y) + 2\text{Cov}(X,Y)$，代入展开：
\[
n(p_1+p_2)(1 - p_1 - p_2) = np_1(1-p_1) + np_2(1-p_2) + 2\text{Cov}(X,Y)
\]
化简同样可得 $\text{Cov}(X,Y) = -n p_1 p_2$，代入相关系数公式即可得相同结果。
\end{solution}


\begin{exercise}\label{exer:4-dens-24y}
设二维随机变量 \((X,Y)\) 的概率密度函数为
\[
f(x,y) = \begin{cases} 24y(1-x), & 0 \leqslant x \leqslant 1,\ 0 \leqslant y \leqslant x, \\ 0, & \text{其他}, \end{cases}
\]
求：(1) \(X\) 的边缘密度 \(f_X(x)\)；(2) \(P(X + Y < 1)\)；(3) \(\mathrm{Cov}(X,Y)\)。
\end{exercise}
\begin{solution}
(1) \(f_X(x) = \displaystyle\int_0^x 24y(1-x)\,\mathrm{d}y = 12x^2(1-x)\quad(0 \leqslant x \leqslant 1)\)。\\
(2) 积分区域为 \(\{0 \leqslant y \leqslant x \leqslant 1,\ x + y < 1\}\)。以 \(x = 1/2\) 为分界：
\begin{itemize}
    \item 当 \(0 \leqslant x \leqslant \dfrac{1}{2}\) 时，\(y\) 从 0 到 \(x\)（此时 \(x + y \leqslant 2x \leqslant 1\) 自动满足）；
    \item 当 \(\dfrac{1}{2} < x \leqslant 1\) 时，\(y\) 从 0 到 \(1 - x\)。
\end{itemize}
\[
P = \int_0^{1/2} 12x^2(1-x)\,\mathrm{d}x + \int_{1/2}^1 12(1-x)^3\,\mathrm{d}x
= \frac{5}{16} + \frac{3}{16} = \frac{1}{2}.
\]
(3)
\begin{align*}
EX &= \int_0^1 12x^3(1-x)\,\mathrm{d}x = 12\cdot B(4,2) = 12\cdot\frac{6}{60} = \frac{3}{5}\\
EY &= \int_0^1\!\!\int_0^x 24y^2(1-x)\,\mathrm{d}y\,\mathrm{d}x = \int_0^1 8x^3(1-x)\,\mathrm{d}x = 8\cdot B(4,2) = \frac{2}{5}\\
EXY &= \int_0^1\!\!\int_0^x 24xy^2(1-x)\,\mathrm{d}y\,\mathrm{d}x = \int_0^1 8x^4(1-x)\,\mathrm{d}x = 8\cdot B(5,2) = 8\cdot\frac{24}{120} = \frac{4}{15}\\
\mathrm{Cov}(X,Y) &= \frac{4}{15} - \frac{3}{5} \times \frac{2}{5} = \frac{2}{75}.
\end{align*}
\end{solution}

\begin{exercise}\label{exer:4-13-stick}
将长度为 1\,m 的木棒随机地截成两段，则两段长度的相关系数为（\quad）。
\end{exercise}
\begin{solution}
应选 \(-1\)。

随机截成两段后，若把其中一段长度记为 \(X\)，另一段长度自然就是 \(Y=1-X=-X+1\)。

这是一条斜率为负的直线。根据相关系数的性质：若 \(P\{Y=aX+b\}=1\)，则当 \(a>0\) 时 \(\rho_{XY}=1\)，当 \(a<0\) 时 \(\rho_{XY}=-1\)。

这里 \(a=-1<0\)，所以 \(\rho_{XY}=-1\)。

直观上也很好理解：一段越长，另一段就必然越短，而且关系完全线性，所以是完全负相关。
\end{solution}

\begin{exercise}\label{exer:4-coin-rho}
掷硬币 2 次，正面出现的次数记为 \(X\)，反面出现的次数记为 \(Y\)，则 \(X\) 与 \(Y\) 的相关系数 \(\rho_{XY} =\,\underline{\hspace{2cm}}\)。
\end{exercise}
\begin{solution}
应填 \(-1\)。

设两次掷硬币中正面出现次数为 \(X\)，反面出现次数为 \(Y\)。由于总共只掷了 2 次，所以无论结果怎样，总有 \(X+Y=2\)，因此 \(Y=2-X=-X+2\)。

这说明 \(Y\) 与 \(X\) 之间存在确定的线性关系，而且斜率是 \(-1\)。根据相关系数的性质，确定线性关系 \(Y=aX+b\) 中，若 \(a<0\)，则 \(\rho_{XY}=-1\)。

所以这里 \(\rho_{XY}=-1\)。

本质原因是：正面次数一旦增加，反面次数就必须等量减少，两者完全负相关。
\end{solution}


\begin{exercise}\label{exer:4-14-linear}
如果存在常数 \(a, b\ (a \neq 0)\)，使 \(P(Y = aX + b) = 1\)，且 \(0 < DX < +\infty\)，则 \(\rho_{XY} =\,\underline{\hspace{2cm}}\)。
\end{exercise}

\begin{solution}
由题意，\(Y=aX+b\) 几乎处处成立，因此可以直接由相关系数定义来算。

先求协方差：
\[
\mathrm{Cov}(X,Y)=\mathrm{Cov}(X,aX+b)=a\,\mathrm{Cov}(X,X)=aDX.
\]

再求 \(Y\) 的方差：
\[
DY=D(aX+b)=a^2DX.
\]

于是
\[
\rho_{XY}
=\frac{\mathrm{Cov}(X,Y)}{\sqrt{DX}\sqrt{DY}}
=\frac{aDX}{\sqrt{DX}\sqrt{a^2DX}}
=\frac{a}{|a|}.
\]

所以
\[
\rho_{XY} = \begin{cases} 1, & a > 0, \\ -1, & a < 0. \end{cases}
\]

也就是说，随机变量之间只要存在确定的线性关系，相关系数就会取到 \(\pm1\)。
\end{solution}



\begin{exercise}\label{exer:4-s13-event}
设 \(A, B\) 为两个随机事件，且 \(P(A) = \dfrac{1}{4}\)，\(P(B|A) = \dfrac{1}{3}\)，\(P(A|B) = \dfrac{1}{2}\)。令
\[
X = \begin{cases} 1, & A \text{ 发生}, \\ 0, & A \text{ 不发生}, \end{cases}\quad
Y = \begin{cases} 1, & B \text{ 发生}, \\ 0, & B \text{ 不发生}. \end{cases}
\]
求：(1) \((X,Y)\) 的概率分布；(2) \(\rho_{XY}\)；(3) \(Z = X^2 + Y^2\) 的概率分布。
\end{exercise}
\begin{solution}
(1) 先计算基本概率：
\[
P(AB) = P(A)\,P(B|A) = \frac{1}{4} \times \frac{1}{3} = \frac{1}{12},\quad
P(B) = \frac{P(AB)}{P(A|B)} = \frac{1/12}{1/2} = \frac{1}{6}.
\]
于是
\begin{align*}
P\{X=0,Y=0\} &= P(\overline{A}\,\overline{B}) = 1 - P(A) - P(B) + P(AB) = 1 - \frac{1}{4} - \frac{1}{6} + \frac{1}{12} = \frac{2}{3}, \\
P\{X=0,Y=1\} &= P(\overline{A}\,B) = P(B) - P(AB) = \frac{1}{6} - \frac{1}{12} = \frac{1}{12}, \\
P\{X=1,Y=0\} &= P(A\,\overline{B}) = P(A) - P(AB) = \frac{1}{4} - \frac{1}{12} = \frac{1}{6}, \\
P\{X=1,Y=1\} &= P(AB) = \frac{1}{12}.
\end{align*}
(2) 边缘分布：\(X \sim (0:\frac{3}{4},\ 1:\frac{1}{4})\)，\(Y \sim (0:\frac{5}{6},\ 1:\frac{1}{6})\)。
\[
DX = \frac{3}{16},\quad DY = \frac{5}{36},\quad E(XY) = \frac{1}{12},
\]
\[
\mathrm{Cov}(X,Y) = \frac{1}{12} - \frac{1}{4} \times \frac{1}{6} = \frac{1}{24},\quad
\rho_{XY} = \frac{1/24}{\sqrt{3/16}\sqrt{5/36}} = \frac{1}{\sqrt{15}} = \frac{\sqrt{15}}{15}.
\]
(3) \(Z\) 的可能取值为 \(0, 1, 2\)：
\[
P\{Z=0\} = \frac{2}{3},\quad P\{Z=1\} = \frac{1}{12}+\frac{1}{6} = \frac{1}{4},\quad P\{Z=2\} = \frac{1}{12}.
\]
\end{solution}

\begin{exercise}\label{exer:4-7-joint}
设随机变量 \(X\) 与 \(Y\) 的概率分布分别为
\[
\begin{array}{c|ccc}
X & 0 & 1 & 2 \\
\hline
P & \frac{1}{3} & \frac{1}{3} & \frac{1}{3}
\end{array},\quad
\begin{array}{c|ccc}
Y & -1 & 0 & 1 \\
\hline
P & \frac{1}{3} & \frac{1}{3} & \frac{1}{3}
\end{array},
\]
且 \(P\{X^2 = Y^2\} = 1\)。\par
(1) 求 \((X,Y)\) 的概率分布；\par
(2) 求 \(Z = XY\) 的概率分布；\par
(3) 求 $\rho_{XY}$。
\end{exercise}

\begin{solution}
(1) 由 \(P\{X^2 = Y^2\} = 1\) 知
\[
P\{X=0,Y=-1\} = P\{X=0,Y=1\} = P\{X=1,Y=0\} = 0,
\]
于是
\[
P\{X=0,Y=0\} = \frac{1}{3},\quad
P\{X=1,Y=1\} = \frac{1}{3},\quad
P\{X=1,Y=-1\} = \frac{1}{3}.
\]
\((X,Y)\) 的联合分布为
\[
\begin{array}{c|ccc|c}
X \backslash Y & -1 & 0 & 1 & P_i \\
\hline
0 & 0 & \frac{1}{3} & 0 & \frac{1}{3} \\[4pt]
1 & \frac{1}{3} & 0 & \frac{1}{3} & \frac{2}{3} \\
\hline
P_j & \frac{1}{3} & \frac{1}{3} & \frac{1}{3} & 1
\end{array}
\]
(2) \(Z = XY\) 的可能取值为 \(-1, 0, 1\)：
\begin{align*}
    P\{Z=-1\} &= P\{X=1,Y=-1\} = \frac{1}{3},\\
    P\{Z=0\} &= P\{X=0,Y=0\} = \frac{1}{3},\\
    P\{Z=1\} &= P\{X=1,Y=1\} = \frac{1}{3}.
\end{align*}
(3) \(EX = \dfrac{2}{3}\)，\(DX = \dfrac{2}{9}\)；\(EY = 0\)，\(DY = \dfrac{2}{3}\)；\(E(XY) = 0\)。
故 \(\mathrm{Cov}(X,Y) = 0\)，\(\rho_{XY} = 0\)。此例说明 \(X\) 与 \(Y\) 不相关，但由联合分布可知 \(X\) 与 \(Y\) 并不独立。
\end{solution}

\begin{exercise}\label{exer:4-2-pm1}
设离散型随机变量 \(X_1, X_2\) 都只取 \(1\) 和 \(-1\)，已知 \(P(X_1 = -1) = P(X_2 = -1) = \dfrac{1}{2}\)，且 \(P(X_2 = -1 \mid X_1 = -1) = \dfrac{1}{3}\)。\par
(1) 求 \((X_1, X_2)\) 的联合概率函数；\par
(2) 求 \(P(X_1 + X_2 = 0)\)；\par
(3) 分别求 \(\mathrm{Cov}(X_1, X_2)\) 和 \(\rho_{X_1X_2}\)。
\end{exercise}
\begin{solution}
(1) \(P(X_2=-1,X_1=-1) = P(X_2=-1 \mid X_1=-1)\,P(X_1=-1) = \dfrac{1}{3} \times \dfrac{1}{2} = \dfrac{1}{6}\)，\\
\(P(X_2=1,X_1=-1) = P(X_1=-1) - P(X_2=-1,X_1=-1) = \dfrac{1}{3}\)。\\
由对称性：\(P(X_2=1,X_1=1) = \dfrac{1}{6}\)，\(P(X_2=-1,X_1=1) = \dfrac{1}{3}\)。
\[
\begin{array}{c|cc}
X_1 \backslash X_2 & -1 & 1 \\
\hline
-1 & \frac{1}{6} & \frac{1}{3} \\[4pt]
1 & \frac{1}{3} & \frac{1}{6}
\end{array}
\]
(2) \(P(X_1+X_2=0) = \dfrac{1}{3} + \dfrac{1}{3} = \dfrac{2}{3}\)。\\
(3) \(EX_1 = 0\)，\(EX_2 = 0\)，\(DX_1 = DX_2 = 1\)。\(E(X_1 X_2) = 1 \times \left(\dfrac{1}{6}+\dfrac{1}{6}\right) + (-1) \times \left(\dfrac{1}{3}+\dfrac{1}{3}\right) = -\dfrac{1}{3}\)。\\
\(\mathrm{Cov}(X_1,X_2) = -\dfrac{1}{3}\)，\(\rho_{X_1X_2} = -\dfrac{1}{3}\)。
\end{solution}


\begin{exercise}\label{exer:4-4-sym}
设随机变量 \(X\) 的概率密度为 \(f(x)\)，\(DX = 2\)，而随机变量 \(Y\) 的概率密度为 \(f(-y)\)，它们的相关系数 \(\rho_{XY} = \dfrac{1}{4}\)，记 \(Z = X + Y\)，求 \(EZ\)，\(DZ\)。
\end{exercise}
\begin{solution}
这道题的关键是先读懂密度 \(f(-y)\) 的含义：它表示 \(Y\) 的分布是把 \(X\) 关于原点镜像后的结果，因此 \(Y\) 与 \(-X\) 同分布。

先求 \(EZ\)。由换元 \(t=-y\)，有
\[
\begin{aligned}
EY
&=\int_{-\infty}^{+\infty} y\,f(-y)\,\mathrm{d}y\\
&=\int_{+\infty}^{-\infty} (-t)\,f(t)\,(-\,\mathrm{d}t)\\
&=-\int_{-\infty}^{+\infty} t\,f(t)\,\mathrm{d}t\\
&=-EX.
\end{aligned}
\]
所以
\[
EZ=EX+EY=EX-EX=0.
\]

再求 \(DY\)。同理，
\[
\begin{aligned}
EY^2
&=\int_{-\infty}^{+\infty} y^2 f(-y)\,\mathrm{d}y\\
&=\int_{-\infty}^{+\infty} t^2 f(t)\,\mathrm{d}t\\
&=EX^2.
\end{aligned}
\]
又因为 \(EY=-EX\)，故
\[
DY=EY^2-(EY)^2=EX^2-(EX)^2=DX=2.
\]

题目给的是相关系数，不是协方差，所以还要先换回来：
\[
\mathrm{Cov}(X,Y)=\rho_{XY}\sqrt{DX}\sqrt{DY}=\frac14\times \sqrt{2}\times \sqrt{2}=\frac12.
\]

最后利用 \(D(X+Y)=DX+DY+2\,\mathrm{Cov}(X,Y)\)，得到 \(DZ=2+2+2\times \frac12=5\)。因此 \(EZ=0,\ DZ=5\)。
\end{solution}

\begin{exercise}\label{exer:4-5-conv}
随机变量 \(X, Y\) 相互独立，均服从参数为 \(\lambda\) 的指数分布。\par
(1) 求 \(Z = X + Y\) 的概率密度函数；\par
(2) 求 \(EZ\)，\(DZ\)；\par
(3) 求 \(\rho_{XZ}\)。
\end{exercise}
\begin{solution}
(1) \(F_Z(z) = P(X+Y \leqslant z) = \displaystyle\int_0^z \mathrm{d}x \int_0^{z-x} \lambda^2 e^{-\lambda x} e^{-\lambda y}\,\mathrm{d}y = 1 - (1+\lambda z)e^{-\lambda z}\)，故
\[
f_Z(z) = \begin{cases} \lambda^2 z\,e^{-\lambda z}, & z > 0, \\[4pt] 0, & z \leqslant 0. \end{cases}
\]
(2) \(EZ = EX + EY = \dfrac{2}{\lambda}\)，\(DZ = DX + DY = \dfrac{2}{\lambda^2}\)。\\
(3) \(\mathrm{Cov}(X,Z) = \mathrm{Cov}(X, X+Y) = DX + \mathrm{Cov}(X,Y) = DX = \dfrac{1}{\lambda^2}\)，故
\[
\rho_{XZ} = \frac{\mathrm{Cov}(X,Z)}{\sqrt{DX}\sqrt{DZ}} = \frac{1/\lambda^2}{\sqrt{1/\lambda^2}\sqrt{2/\lambda^2}} = \frac{1}{\sqrt{2}}.
\]
\end{solution}


\subsection{不相关性的判定}

\begin{exercise}\label{exer:4-s4-x3}
设随机变量 \(X\) 在 \([-1,1]\) 上服从均匀分布，\(Y = X^3\)，则 \(X\) 与 \(Y\)（\quad）。\par
A. 不相关且相互独立\quad B. 不相关且不相互独立\quad
C. 相关且相互独立\quad D. 相关且不相互独立
\end{exercise}

\begin{solution}
应选 D。\(X\) 与 \(Y\) 之间存在确定性函数关系 \(Y = X^3\)，故不可能独立，排除 A、C。计算协方差：\(EX = 0\)，
\[
E(XY) = E(X^4) = \int_{-1}^{1} x^4 \cdot \frac{1}{2}\,\mathrm{d}x = \frac{1}{5} \neq 0,
\]
故 \(\mathrm{Cov}(X,Y) = E(XY) - EX \cdot EY = \dfrac{1}{5} \neq 0\)，\(X\) 与 \(Y\) 相关。

\begin{note}
本题常见误区：有人误将 \(E(XY)\) 写成 \(\int x^3 \cdot (1/2)\,\mathrm{d}x = 0\)（奇函数），但这计算的是 \(E(X^3)\) 而非 \(E(XY) = E(X^4)\)。注意 \(XY = X \cdot X^3 = X^4\)，是偶函数。
\end{note}
\end{solution}

\begin{exercise}\label{exer:4-s15-normal3}
设 \(X \sim N(0,\sigma^2)\ (\sigma > 0)\)，\(Y = X^3\)。\par
(1) 求 \(EY\)；\par
(2) 判断 \(X\) 与 \(Y\) 是否相关？是否独立？说明理由。
\end{exercise}

\begin{solution}
(1) \(X\) 的密度 \(f(x)\) 为偶函数，\(x^3\) 为奇函数，故
\[
EY = E(X^3) = \int_{-\infty}^{+\infty} x^3 f(x)\,\mathrm{d}x = 0.
\]

(2) \(E(XY) = E(X^4)\)。对 \(X \sim N(0,\sigma^2)\)，利用 Gamma 函数计算：
\[
E(X^4) = \int_{-\infty}^{+\infty} x^4 \cdot \frac{1}{\sqrt{2\pi}\,\sigma}\,e^{-\frac{x^2}{2\sigma^2}}\,\mathrm{d}x = \frac{2}{\sqrt{2\pi}\,\sigma}\int_0^{+\infty} x^4\,e^{-\frac{x^2}{2\sigma^2}}\,\mathrm{d}x = 3\sigma^4 \neq 0.
\]
故 \(\mathrm{Cov}(X,Y) = E(XY) - EX \cdot EY = 3\sigma^4 \neq 0\)，\(X\) 与 \(Y\) 相关，从而不独立。
\end{solution}

\begin{exercise}\label{exer:4-11-even}
设随机变量 \(X\) 的密度函数 \(f(x)\) 是偶函数，令 \(Y = X^2\)，假定 \(EX\)、\(EY\) 都存在，求证 \(\mathrm{Cov}(X,Y) = 0\)。
\end{exercise}

\begin{solution}
由于 \(f(x)\) 为偶函数，\(x^3\) 为奇函数，故
\[
\int_{-\infty}^{+\infty} x^3 f(x)\,\mathrm{d}x = 0.
\]
同时 \(xf(x)\) 也是奇函数，所以 \(EX=0\)。于是
\[
\mathrm{Cov}(X,Y) = E(XY) - EX \cdot EY = E(X^3) - EX \cdot E(X^2) = 0 - 0\cdot E(X^2) = 0.
\]
\end{solution}

\begin{exercise}\label{exer:4-12-sincos}
设 \(\theta\) 服从 \([-\pi,\pi]\) 上的均匀分布，又 \(X = \sin\theta\)，\(Y = \cos\theta\)，则 \(\rho_{XY} =\,\underline{\hspace{2cm}}\)。
\end{exercise}

\begin{solution}
应填 0。
\[
EX = \int_{-\pi}^{\pi} \frac{\sin\theta}{2\pi}\,\mathrm{d}\theta = 0,\quad
EY = \int_{-\pi}^{\pi} \frac{\cos\theta}{2\pi}\,\mathrm{d}\theta = 0,
\]
故
\[
\rho_{XY} = \frac{EXY - EX \cdot EY}{\sqrt{DX}\sqrt{DY}} = \frac{\frac{1}{2\pi}\int_{-\pi}^{\pi}\sin\theta\cos\theta\,\mathrm{d}\theta}{\sqrt{DX}\sqrt{DY}} = 0.
\]

\begin{note}
此例中 \(X^2 + Y^2 = \sin^2\theta + \cos^2\theta = 1\)，说明 \(X\) 与 \(Y\) 之间存在非线性关系（不独立），但 \(\rho_{XY} = 0\)（不相关）。这是"不相关不等于独立"的典型例子。
\end{note}
\end{solution}

\begin{exercise}\label{exer:4-18-iid}
已知随机变量 \(X, Y\) 独立同分布，记 \(U = X - Y\)，\(V = X + Y\)，则 \(U, V\)（\quad）。\par
A. 不独立\quad B. 独立\quad C. 相关系数不为 0\quad D. 相关系数为 0
\end{exercise}

\begin{solution}
应选 D。\(EU = EX - EY = 0\)，\(EV = EX + EY = 2EX\)。

\(E(UV) = E(X^2 - Y^2) = EX^2 - EY^2 = 0\)（因为 \(X, Y\) 同分布），故
\[
\rho_{UV} = \frac{E(UV) - EU \cdot EV}{\sqrt{DU}\sqrt{DV}} = 0.
\]
由于未给出具体分布，无法判断是否独立，故只能选 D。

\begin{note}
若 \(X, Y\) 同为正态分布，则 \(U, V\) 服从联合正态分布，不相关即独立，此时应选 B。但本题未说明分布类型，故只能选 D。
\end{note}
\end{solution}


\begin{exercise}\label{exer:4-19-interval}
在区间 \([0.5, 1]\) 中任取一个值 \(x\)，再在区间 \([-x, x]\) 中任取一个值 \(y\)，构成二维随机变量 \((X,Y)\)。\par
(1) 求 \((X,Y)\) 的联合密度函数 \(f(x,y)\)，关于 \(Y\) 的边缘密度函数 \(f_Y(y)\)，并判断 \(X, Y\) 是否独立；\par
(2) 求 \(\mathrm{Cov}(X,Y)\)，并据此判断 \(X, Y\) 是否相关。
\end{exercise}

\begin{solution}
(1) \((X,Y)\) 的取值区域为 \(0.5 < x < 1\)，\(-x < y < x\)。

\(X \sim U(0.5, 1)\)，故 \(f_X(x) = \dfrac{1}{0.5} = 2\ (0.5 < x < 1)\)。给定 \(X = x\)，\(Y \sim U(-x, x)\)，故 \(f(y|x) = \dfrac{1}{2x}\ (-x < y < x)\)。

联合密度函数：
\[
f(x,y) = f_X(x)\,f(y|x) = 2 \times \frac{1}{2x} = \frac{1}{x},\quad 0.5 < x < 1,\ -x < y < x.
\]

关于 \(Y\) 的边缘密度函数。对固定的 \(y\)，\(x\) 的范围由 \(\max(0.5, |y|) < x < 1\) 确定：
\[
f_Y(y) = \begin{cases}
\displaystyle\int_y^1 \frac{1}{x}\,\mathrm{d}x = -\ln y, & 0.5 \leqslant y < 1, \\[8pt]
\displaystyle\int_{0.5}^1 \frac{1}{x}\,\mathrm{d}x = \ln 2, & -0.5 \leqslant y < 0.5, \\[8pt]
\displaystyle\int_{-y}^1 \frac{1}{x}\,\mathrm{d}x = -\ln(-y), & -1 \leqslant y < -0.5.
\end{cases}
\]

\(f_X(x) = \displaystyle\int_{-x}^{x} \frac{1}{x}\,\mathrm{d}y = 2\)。由于 \(f(x,y) = \dfrac{1}{x} \neq 2\,f_Y(y) = f_X(x)\,f_Y(y)\)，故 \(X\) 与 \(Y\) 不独立。

(2) \(EX = \displaystyle\int_{0.5}^{1} 2x\,\mathrm{d}x = \dfrac{3}{4}\)。由对称性（\(f_Y(y)\) 关于 \(y = 0\) 对称），\(EY = 0\)。

\[
\mathrm{Cov}(X,Y) = EXY - EX \cdot EY = EXY = \int_{0.5}^{1}\int_{-x}^{x} xy \cdot \frac{1}{x}\,\mathrm{d}y\,\mathrm{d}x = \int_{0.5}^{1} x\!\left[\int_{-x}^{x} y\,\mathrm{d}y\right]\mathrm{d}x = 0.
\]

所以 \(X, Y\) 不相关（但由 (1) 知二者不独立）。
\end{solution}

\subsection{正态分布的数字特征}

\begin{exercise}\label{exer:4-s2-z1z2}
设随机变量 \((X,Y)\) 服从二维正态分布，\(X\) 与 \(Y\) 的期望均为 \(\mu\)，方差均为 \(\sigma^2\)，相关系数 \(\rho_{XY} = 0\)。记 \(Z_1 = 2X+Y\)，\(Z_2 = 2X-Y\)，则 \(\rho_{Z_1Z_2} =\,\underline{\hspace{2cm}}\)。
\end{exercise}
\begin{solution}
由 \(\rho_{XY} = 0\) 知 \(X\) 与 \(Y\) 相互独立。
\[
DZ_1 = 4DX + DY = 5\sigma^2,\quad DZ_2 = 4DX + DY = 5\sigma^2,
\]
\[
\mathrm{Cov}(Z_1,Z_2) = \mathrm{Cov}(2X+Y,\, 2X-Y) = 4DX - DY = 3\sigma^2,
\]
\[
\rho_{Z_1Z_2} = \frac{3\sigma^2}{\sqrt{5\sigma^2}\sqrt{5\sigma^2}} = \frac{3}{5}.
\]
\end{solution}



\begin{exercise}\label{exer:4-15-absnormal}
设两个随机变量 \(X, Y\) 相互独立，且都服从均值为 0、方差为 0.5 的正态分布，求 \(|X-Y|\) 的数学期望与方差。
\end{exercise}

\begin{solution}
设 \(Z=X-Y\)。由于 \(X,Y\) 相互独立，且都服从均值为 0、方差为 0.5 的正态分布，所以线性组合 \(Z\) 仍服从正态分布，并且 \(EZ=EX-EY=0,\ DZ=DX+DY=0.5+0.5=1\)，因此 \(Z\sim N(0,1)\)。

题目要求的是 \(|X-Y|=|Z|\) 的数学期望与方差，先求 \(E|Z|\)：
\[
E(|X-Y|)=E|Z|=\int_{-\infty}^{+\infty}|z|\,\frac{1}{\sqrt{2\pi}}\,e^{-\frac{z^2}{2}}\,\mathrm{d}z.
\]
被积函数关于 0 对称，于是
\[
E|Z|=\frac{2}{\sqrt{2\pi}}\int_0^{+\infty} z\,e^{-\frac{z^2}{2}}\,\mathrm{d}z
=\sqrt{\frac{2}{\pi}}.
\]

再求方差。注意
\[
|Z|^2=Z^2,
\]
所以
\[
E(|X-Y|^2)=E(|Z|^2)=E(Z^2)=DZ+(EZ)^2=1+0=1.
\]
因此
\[
D(|X-Y|)=D(|Z|)=E(|Z|^2)-[E|Z|]^2=1-\frac{2}{\pi}.
\]

所以
\[
E(|X-Y|)=\sqrt{\frac{2}{\pi}},\qquad D(|X-Y|)=1-\frac{2}{\pi}.
\]
\end{solution}

\begin{exercise}\label{exer:4-16-multi}
已知随机变量 \(X \sim N(0,1)\)，\(Y \sim N(1,4)\)，且 \(X\) 与 \(Y\) 的相关系数为 \(\dfrac{1}{3}\)，设 \(Z = 2X + 3Y\)，求 \(\rho_{XZ}\)。
\end{exercise}
\begin{solution}
由 \(\rho_{XY} = \dfrac{\mathrm{Cov}(X,Y)}{1 \times 2} = \dfrac{1}{3}\)，得 \(\mathrm{Cov}(X,Y) = \dfrac{2}{3}\)。
\[
\mathrm{Cov}(X,Z) = \mathrm{Cov}(X, 2X+3Y) = 2DX + 3\,\mathrm{Cov}(X,Y) = 2 + 3 \times \frac{2}{3} = 4,
\]
\[
DZ = 4DX + 9DY + 12\,\mathrm{Cov}(X,Y) = 4 + 36 + 12 \times \frac{2}{3} = 48,
\]
故
\[
\rho_{XZ} = \frac{\mathrm{Cov}(X,Z)}{\sqrt{DX}\sqrt{DZ}} = \frac{4}{\sqrt{48}} = \frac{4}{4\sqrt{3}} = \frac{1}{\sqrt{3}}.
\]
\end{solution}


\begin{exercise}\label{exer:4-rho-xz}
已知随机变量 \(X \sim N(1,3^2)\)，\(Y \sim N(0,4^2)\)，且 \(\rho_{XY} = \dfrac{1}{2}\)。设 \(Z = \dfrac{X}{3} + \dfrac{Y}{2}\)，求：(1) \(EZ\) 和 \(DZ\)；(2) \(\rho_{XZ}\)。
\end{exercise}
\begin{solution}
(1) \(\mathrm{Cov}(X,Y) = \dfrac{1}{2} \times 3 \times 4 = 6\)。
\[
EZ = \frac{1}{3} \times 1 + \frac{1}{2} \times 0 = \frac{1}{3},\quad
DZ = \frac{1}{9} \times 9 + \frac{1}{4} \times 16 + 2 \times \frac{1}{3} \times \frac{1}{2} \times 6 = 1 + 4 + 2 = 7.
\]
(2) \(\mathrm{Cov}(X,Z) = \dfrac{1}{3}\,DX + \dfrac{1}{2}\,\mathrm{Cov}(X,Y) = 3 + 3 = 6\)，故
\[
\rho_{XZ} = \frac{\mathrm{Cov}(X,Z)}{\sqrt{DX}\sqrt{DZ}} = \frac{6}{\sqrt{9 \times 7}} = \frac{2\sqrt{7}}{7}.
\]
\end{solution}



\begin{exercise}\label{exer:4-17-bivariate}
设二维随机变量 \((X,Y) \sim N(1, 3^2;\, 0, 4^2,\, 0)\)，则 \(E(XY) =\,\underline{\hspace{2cm}}\)。
\end{exercise}

\begin{solution}
应填 \(0\)。

题目给出
\[
(X,Y)\sim N(1,3^2;\,0,4^2,\,0),
\]
这里最后一个参数 0 表示相关系数
\[
\rho_{XY}=0.
\]
因此
\[
\mathrm{Cov}(X,Y)=\rho_{XY}\sqrt{DX}\sqrt{DY}=0.
\]

如果只是一般随机变量，\(\rho_{XY}=0\) 只能说明不相关，未必独立；但题目这里是\textbf{二维正态分布}，此时“不相关”等价于“独立”。所以可以推出
\[
X\ \text{与}\ Y\ \text{相互独立}.
\]

于是
\[
E(XY) = EX \cdot EY = 1 \times 0 = 0.
\]

因此答案是 \(0\)。
\end{solution}




\chapter{大数定律和中心极限定理}

\begin{note}
\textbf{【高频题型模板】求 $DX$ 的三条路径：}

考试中求方差的题目，根据已知条件选择最快路径：
\begin{enumerate}
    \item \textbf{公式法（最常用）：} $DX=EX^2-(EX)^2$。先分别算 $EX$ 和 $EX^2$，再相减。
    \item \textbf{定义法：} $DX=E(X-EX)^2$。当 $EX=0$ 时特别简单：$DX=EX^2$。
    \item \textbf{性质法：} 利用 $D(aX+b)=a^2DX$、独立时 $D(X\pm Y)=DX+DY$。
\end{enumerate}

\textbf{【高频题型模板】求 $E[g(X,Y)]$：}

不需要先求 $Z=g(X,Y)$ 的分布！直接用期望公式：
\[
E[g(X,Y)]=\iint g(x,y)\,f(x,y)\,\mathrm{d}x\,\mathrm{d}y\quad\text{（连续型）}
\]
\[
E[g(X,Y)]=\sum_i\sum_j g(x_i,y_j)\,p_{ij}\quad\text{（离散型）}
\]
\textbf{常见错误：} 先费力求出 $Z$ 的分布再算 $EZ$——浪费时间且容易出错。
\end{note}

\begin{note}
\textbf{自学导读：本章在整本书中的地位。}

前四章你学的是"给定一个分布，算概率、算期望、算方差"。本章回答一个更深的问题：\textbf{当样本量很大时，会发生什么？}

两个核心结论：
\begin{enumerate}
    \item \textbf{大数定律}：大量独立随机变量的\textbf{平均值}会稳定在期望附近——这解释了为什么"多次实验取平均"是可靠的。
    \item \textbf{中心极限定理}：大量独立随机变量的\textbf{和}（不管每个变量本身服从什么分布）近似服从正态分布——这解释了为什么正态分布在自然界中无处不在。
\end{enumerate}

\textbf{为什么这对后面的统计学至关重要？} 从第 6 章开始，你将面对"从样本推断总体"的问题。大数定律保证了样本均值是总体均值的好估计；中心极限定理保证了即使总体不是正态的，样本均值也近似正态——这是后面构造置信区间和假设检验的理论基础。

\textbf{前置知识检查：} 本章只需要第 4 章的期望 $EX$、方差 $DX$ 及其性质（线性性、独立时方差可加）。如果你对 $E(\bar{X})=\mu$、$D(\bar{X})=\sigma^2/n$ 还不熟，建议先回顾第 4 章再继续。
\end{note}

\begin{note}
\textbf{自学导读：本章在整本书中的位置。}

前四章你学会了如何\textbf{精确}描述随机变量——给出分布、算期望、算方差。但现实中我们经常面对这样的问题：
\begin{itemize}
    \item 抛硬币 1000 次，正面比例是否"接近" $1/2$？（大数定律）
    \item 100 个零件的总重量大约服从什么分布？能否近似计算概率？（中心极限定理）
    \item 只知道随机变量非负且知道平均值，能否先控制它取大值的概率？（马尔可夫不等式）
    \item 只知道均值和方差，不知道具体分布，能否给概率一个粗略的上界？（切比雪夫不等式）
\end{itemize}

\textbf{核心思想只有一句话：} 大量独立随机变量的平均效果是\textbf{可预测的}——平均值趋向期望（大数定律），总和近似正态（中心极限定理）。

\textbf{为什么这一章是后续统计学的基石？} 第 6--9 章的全部内容（抽样分布、估计、检验）都建立在"样本均值近似正态"这一事实之上。没有中心极限定理，统计推断就失去了理论根基。所以本章虽然短，但地位极其关键。

\textbf{学习建议：}
\begin{enumerate}
    \item 先理解"依概率收敛"的直觉：$n$ 越大，偏离的概率越小；
    \item 大数定律：记住结论"样本均值 $\xrightarrow{P}$ 总体均值"；
    \item 中心极限定理：记住结论"标准化后的和 $\xrightarrow{d}$ 标准正态"，并学会用它近似计算概率。
\end{enumerate}
\end{note}

\[
\begin{cases}
\text{马尔可夫不等式：} \displaystyle P\{X \geqslant a\} \leqslant \frac{EX}{a}\ (X\ge 0) \\[6pt]
\text{切比雪夫不等式：} \displaystyle P\left\{|X - EX| \geqslant \varepsilon\right\} \leq \frac{DX}{\varepsilon^2} \\[6pt]
\text{大数定理}
\begin{cases}
\text{切比雪夫大数定律} \\
\text{独立同分布的大数定理} \\
\text{辛钦大数定理}
\end{cases} \\[6pt]
\text{中心极限定理}
\begin{cases}
\text{列维--林德伯格中心极限定理} \\
\text{拉普拉斯中心极限定理}
\end{cases}
\end{cases}
\]

\begin{introduction}
    \item 马尔可夫不等式~\ref{thm:markov}
    \item 依概率收敛的定义与性质~\ref{def:convergence-prob}
    \item 切比雪夫不等式的应用~\ref{thm:chebyshev}
    \item 切比雪夫大数定律~\ref{thm:chebyshev-lln}
    \item 伯努利大数定律~\ref{thm:bernoulli-lln}
    \item 辛钦大数定律~\ref{thm:khinchin-lln}
    \item 列维--林德伯格中心极限定理~\ref{thm:levy-lindberg}
    \item 棣莫弗--拉普拉斯中心极限定理~\ref{thm:demoivre-laplace}
\end{introduction}

\begin{note}
\textbf{使用提醒：} 自学本章时，先分清题目究竟在问哪一类问题。若题目关心“平均数、频率最终靠近谁”，主角是大数定律；若题目关心“和或均值的概率怎样近似计算”，主角是中心极限定理；若题目只给出非负性和期望，马尔可夫不等式可以先给一个最粗的右尾上界；若进一步给出均值和方差，又暂时拿不到精确分布，则切比雪夫不等式常常是更合适的保底工具。先分工，再套公式，通常比先背定理名字更稳。
\end{note}

\begin{note}
\textbf{不要把本章学成“定理名单”。} 更自然的判断顺序是先问自己两个问题：
\begin{enumerate}
    \item 题目到底在问“平均值最终靠近谁”，还是在问“和/均值的概率怎样近似计算”？
    \item 题目已知的信息够不够强：只是给了均值方差，还是已经说明独立同分布，甚至给出了具体分布？
\end{enumerate}

若目标是“说明平均值会稳定”，就往大数定律方向想；若目标是“近似分布、近似概率”，就往中心极限定理方向想。至于用哪一个版本，不必死背名字，而要看\textbf{手里有什么条件，想得到什么结论}：切比雪夫不等式是信息最少时的保底工具，大数定律负责说明“平均起来会稳”，中心极限定理负责进一步说明“稳了以后怎样波动”。
\end{note}

\begin{note}
\textbf{一张最实用的判断图：}
\begin{enumerate}
    \item 只给 \(X\ge 0\) 和 \(EX\)，问右尾概率上界，优先想马尔可夫不等式；
    \item 只给 \(EX\)、\(DX\)，问概率上界或下界，优先想切比雪夫不等式；
    \item 题目问“频率”“均值”最后靠近谁，优先想大数定律；
    \item 题目问“和”“均值”的近似分布或近似概率，且有独立同分布和矩条件，优先想中心极限定理。
\end{enumerate}

真正要记住的不是定理名，而是：\textbf{先看对象，再看条件，最后才套公式。}
\end{note}

\section{依概率收敛}

\begin{note}
\textbf{自学导读：} "依概率收敛"是一个新的极限概念，和高数中的数列极限不同。数列极限 $a_n\to a$ 是确定性的；而这里 $X_n$ 是随机变量，每次实验的值都不同，所以不能说"$X_n$ 的值趋近于 $a$"。正确的说法是：随着 $n$ 增大，$X_n$ 偏离 $a$ 超过任意小量 $\varepsilon$ 的\textbf{概率}趋于 0。

直觉：掷骰子 $n$ 次取平均，$n=10$ 时平均值可能偏离 3.5 很远；$n=10000$ 时平均值几乎肯定在 3.5 附近。这就是依概率收敛。
\end{note}

大数定律和中心极限定理研究的是大量随机变量之和在极限状态下的行为规律，而"依概率收敛"是描述这一极限行为的基本工具。

\begin{definition}[依概率收敛] \label{def:convergence-prob}
设 \(X\) 为随机变量，\(\{X_n\}\ (n=1,2,\dots)\) 为随机变量序列。若对任意 \(\varepsilon > 0\)，有
\[
\lim_{n\to\infty} P\{|X_n - X| \geqslant \varepsilon\} = 0
\quad \text{或等价地} \quad
\lim_{n\to\infty} P\{|X_n - X| < \varepsilon\} = 1,
\]
则称序列 \(\{X_n\}\) \textbf{依概率收敛}于 \(X\)，记为
\[
X_n \stackrel{P}{\longrightarrow} X\ (n\to\infty).
\]
特别地，若 \(X = a\) 为常数，则称 \(\{X_n\}\) 依概率收敛于常数 \(a\)。
\end{definition}

\begin{property}[概率收敛运算性质]\label{prop:conv-prob-operation}
依概率收敛具有以下运算性质。设 \(X_n \stackrel{P}{\longrightarrow} a\)，\(Y_n \stackrel{P}{\longrightarrow} b\)，\(g(x,y)\) 为连续函数，则
\[
g(X_n, Y_n) \stackrel{P}{\longrightarrow} g(a, b).
\]
更一般地，对 \(m\) 元连续函数 \(g(x_1, x_2, \dots, x_m)\)，若 \(X_{n}^{(i)} \stackrel{P}{\longrightarrow} a_i\ (i=1,2,\dots,m)\)，则
\[
g\!\left(X_{n}^{(1)}, X_{n}^{(2)}, \dots, X_{n}^{(m)}\right) \stackrel{P}{\longrightarrow} g(a_1, a_2, \dots, a_m).
\]
\end{property}

\begin{note}
\textbf{自学追问：} 这个定义到底在说什么？固定任意一个误差带 \(\varepsilon>0\) 后，只要 \(n\) 足够大，\(X_n\) 偏离 \(X\) 超过 \(\varepsilon\) 的概率就会变得任意小。也就是说，它强调的是“较大偏差越来越不常见”。

\textbf{易错点：} 依概率收敛是概率意义下的收敛，并不要求每一次试验中 \(X_n\) 都逐项贴近 \(X\)。它允许偶尔仍出现较大偏差，只要求这类偏差发生的概率趋于 0。

依概率收敛是重要的弱收敛方式之一。通常有“几乎必然收敛 \(\Rightarrow\) 依概率收敛 \(\Rightarrow\) 依分布收敛”，但反向一般不成立。本章中的大数定律，描述的正是随机变量序列的算术平均值依概率收敛到其期望这一事实。
\end{note}

\begin{note}
\textbf{为什么大数定律通常写成“依概率收敛”，而不是“每次实验都越来越接近”？} 因为随机试验天生允许偶然波动。即使掷一万次骰子，样本平均数大多数时候会很接近 \(3.5\)，但仍可能偶尔偏得比较远；概率论研究的是这种“大偏差会不会越来越少见”，而不是要求每一条样本路径都像数列那样逐项贴近。

换句话说，本章的极限定理都在做同一件事：不是消灭随机性，而是说明\textbf{样本量变大后，随机性的影响会以可控的方式缩小}。理解了这一点，再看后面“大数定律讲平均值稳定，CLT 讲剩余波动的形状”，就会顺很多。
\end{note}


\section{马尔可夫不等式}

\begin{note}
\textbf{直觉理解：} 马尔可夫不等式说的是一句很朴素的话：\textbf{如果一个非负随机变量的平均值不大，那么它取到很大数值的概率也不可能太大。}

例如，若 \(X\ge 0\) 且 \(EX=2\)，那么 \(X\) 取到 \(10\) 以上的概率至多只能是 \(0.2\)。否则，若 \(X\ge10\) 的情况出现得太频繁，平均值就不可能只有 2。

\textbf{什么时候用它？} 当题目只告诉你随机变量\textbf{非负}且只知道 \(EX\)，还没有方差信息时，马尔可夫不等式往往是最先能用的统一工具。
\end{note}

\begin{theorem}[马尔可夫不等式] \label{thm:markov}
设随机变量 \(X\ge 0\)，且数学期望 \(EX\) 存在，则对任意 \(a>0\)，有
\[
P\{X\ge a\}\le \frac{EX}{a}.
\]
\end{theorem}
\begin{proof}
仅对连续型给出证明（离散型类似）。设 \(X\) 的概率密度为 \(f(x)\)，则
\begin{align*}
EX &= \int_0^{+\infty}x\,f(x)\,\mathrm{d}x \\
&\ge \int_a^{+\infty}x\,f(x)\,\mathrm{d}x \\
&\ge a\int_a^{+\infty}f(x)\,\mathrm{d}x \\
&= a\,P\{X\ge a\}.
\end{align*}
两边同除以 \(a\) 即得
\[
P\{X\ge a\}\le \frac{EX}{a}.
\]
\end{proof}

\begin{note}
\textbf{为什么这里一定要求 \(X\ge 0\)？} 因为马尔可夫不等式比较的是“平均大小”和“右尾事件 \(X\ge a\) 的概率”。若 \(X\) 允许取负值，大量负值会把平均值拉小，但这并不能阻止它偶尔取很大的正值，所以单凭 \(EX\) 就无法有效控制 \(P\{X\ge a\}\)。

\textbf{和切比雪夫不等式的关系：} 切比雪夫不等式其实可以看成把马尔可夫不等式用在非负随机变量 \(Y=(X-EX)^2\) 上。因为事件 \(\{|X-EX|\ge \varepsilon\}\) 等价于
\[
\{(X-EX)^2\ge \varepsilon^2\},
\]
再对 \(Y\) 使用马尔可夫不等式，就得到
\[
P\{|X-EX|\ge \varepsilon\}\le \frac{E[(X-EX)^2]}{\varepsilon^2}=\frac{DX}{\varepsilon^2}.
\]
所以从逻辑上说，\textbf{马尔可夫不等式更基础，切比雪夫不等式更有针对性。}
\end{note}

\begin{exercise}\label{exer:5-markov-basic}
设随机变量 \(X\ge 0\)，且 \(EX=3\)。由马尔可夫不等式可得
\[
P(X\ge 12)\le \underline{\hspace{2cm}}.
\]
\end{exercise}
\begin{solution}
由马尔可夫不等式 \(P(X\ge a)\le \dfrac{EX}{a}\)，令 \(a=12\)，代入 \(EX=3\)，得
\[
P(X\ge 12)\le \frac{3}{12}=\frac14.
\]

所以应填 \(\dfrac14\)。
\end{solution}


\section{切比雪夫不等式}

\begin{note}
\textbf{直觉理解：} 切比雪夫不等式说的是一句大白话——\textbf{方差小的随机变量不太可能偏离均值太远}。具体地：偏离均值超过 $k$ 个标准差的概率不超过 $1/k^2$。

比如 $DX=1$，那么 $P(|X-EX|\geqslant 3) \leqslant 1/9 \approx 11\%$。这个估计对\textbf{任何}分布都成立（不需要知道具体是什么分布），所以它是一个"万能但粗糙"的工具。

\textbf{什么时候用它？} 当题目只告诉你 $EX$ 和 $DX$，没给具体分布时，切比雪夫不等式是你唯一的武器。
\end{note}

\begin{theorem}[切比雪夫不等式] \label{thm:chebyshev}
设随机变量 \(X\) 的数学期望 \(EX\) 和方差 \(DX\) 都存在，则对任意 \(\varepsilon > 0\)，有
\[
P\{|X - EX| \geqslant \varepsilon\} \leqslant \frac{DX}{\varepsilon^2},
\]
或等价地，
\[
P\{|X - EX| < \varepsilon\} \geqslant 1 - \frac{DX}{\varepsilon^2}.
\]
\end{theorem}
\begin{proof}
仅对连续型给出证明（离散型类似）。设 \(X\) 的概率密度为 \(f(x)\)，则
\begin{align*}
DX &= E[(X-EX)^2] = \int_{-\infty}^{+\infty}(x-EX)^2\,f(x)\,\mathrm{d}x \\
&\geqslant \int_{|x-EX|\geqslant\varepsilon}(x-EX)^2\,f(x)\,\mathrm{d}x
\geqslant \varepsilon^2 \int_{|x-EX|\geqslant\varepsilon}f(x)\,\mathrm{d}x \\
&= \varepsilon^2\,P\{|X-EX| \geqslant \varepsilon\}.
\end{align*}
两边同除以 \(\varepsilon^2\) 即得 \(P\{|X-EX| \geqslant \varepsilon\} \leqslant DX/\varepsilon^2\)。
\end{proof}

\begin{note}
\textbf{使用提醒：} 当题目只给出期望和方差，或者虽然给出分布但不值得去做精确积分时，切比雪夫不等式可以提供一个统一的“保底估计”。它反映的是：方差越小，随机变量偏离均值太远的可能性就越小。

\textbf{易错点：} 求 \(P\{|X-EX|\ge \varepsilon\}\) 时用原式给上界；求 \(P\{|X-EX|<\varepsilon\}\) 时要改用等价形式给下界。它给出的通常是比较粗的界，不要把它误当成精确概率。
\end{note}

\begin{exercise}\label{exer:4-s10-cheby}
设 \(X\) 为随机变量，\(P(|X - EX| < \varepsilon) \geqslant 0.9\)，\(DX = 0.009\)，用切比雪夫不等式估计 \(\varepsilon\) 的取值范围。
\end{exercise}
\begin{solution}
题目已经给出了 \(P(|X-EX|<\varepsilon)\ge 0.9\)，这正对应切比雪夫不等式的等价形式 \(P(|X-EX|<\varepsilon)\ge 1-\dfrac{DX}{\varepsilon^2}\)。把 \(DX=0.009\) 代入，可得
\[
P(|X-EX|<\varepsilon)\ge 1-\frac{0.009}{\varepsilon^2}.
\]
要使这个下界至少达到题目要求的 \(0.9\)，应有
\[
1-\frac{0.009}{\varepsilon^2}\ge 0.9 \implies \frac{0.009}{\varepsilon^2}\le 0.1 \implies \varepsilon^2\ge \frac{0.009}{0.1}=0.09 \implies \varepsilon\ge 0.3.
\]
因此 \(\varepsilon\) 的取值范围为 \(\varepsilon\ge 0.3\)。
\end{solution}


\section{大数定律}

大数定律是概率论中最重要的极限定理之一，它从理论上严格证明了"频率稳定于概率"这一经验事实。切比雪夫不等式（定理~\ref{thm:chebyshev}）是大数定律的理论基础——正是利用这一不等式对概率进行粗略估计，才能证明各种形式的大数定律。

在应用切比雪夫不等式时，需要特别注意不等号的方向：原式为
\[
P\{|X - EX| \geq \varepsilon\} \leq \frac{DX}{\varepsilon^2},
\]
其等价形式为
\[
P\{|X - EX| < \varepsilon\} \geq 1 - \frac{DX}{\varepsilon^2}.
\]
公式中的随机变量 \(X\) 可以替换为任何满足方差存在条件的随机变量。

\begin{exercise}\label{exer:5-chebyshev-uniform}
设随机变量 \(X \sim U(-1,2)\)，由切比雪夫不等式得 \(P(|X - 0.5| < 1) >\,\underline{\hspace{2cm}}\)。
\end{exercise}
\begin{solution}
应填 \(\dfrac14\)。

使用切比雪夫不等式前，先把题目识别成
\[
P(|X-EX|<\varepsilon)
\]
的形式。

由于 \(X\sim U(-1,2)\)，所以
\[
EX=\frac{-1+2}{2}=\frac12,\qquad DX=\frac{(b-a)^2}{12}=\frac{(2-(-1))^2}{12}=\frac34.
\]

题目中的事件正好是 \(|X-0.5|<1\)，即 \(\varepsilon=1\)。由切比雪夫不等式的等价形式
\[
P(|X-EX|<\varepsilon)\ge 1-\frac{DX}{\varepsilon^2},
\]
得到
\[
P(|X - 0.5| < 1) \geq 1 - \frac{DX}{\varepsilon^2} = 1 - \frac{3}{4} = \frac{1}{4}.
\]

所以应填 \(\dfrac14\)。
\end{solution}


\begin{exercise}\label{exer:5-chebyshev-basic}
设随机变量 \(X\) 的方差 \(DX = 2\)，则由切比雪夫不等式，有 \(P(|X - EX| \geq 2) \leq\,\underline{\hspace{2cm}}\)。
\end{exercise}
\begin{solution}
本题就是切比雪夫不等式的直接代入。

切比雪夫不等式为
\[
P(|X-EX|\ge \varepsilon)\le \frac{DX}{\varepsilon^2}.
\]
题目给出 \(DX=2\)，且这里
\[
\varepsilon=2.
\]
因此
\[
P(|X - EX| \geq 2) \leq \frac{DX}{\varepsilon^2} = \frac{2}{4} = \frac{1}{2}.
\]

所以应填 \(\dfrac12\)。
\end{solution}

\begin{exercise}\label{exer:5-chebyshev-sum}
设 \(EX = -2\)，\(EY = 2\)，\(DX = 1\)，\(DY = 4\)，\(\rho(X,Y) = 0.5\)。求 \(D(X+Y) =\,\underline{\hspace{2cm}}\)，并根据切比雪夫不等式估计 \(P\{|X+Y| \geq 6\}\,\underline{\hspace{2cm}}\)。
\end{exercise}
\begin{solution}
先由相关系数求协方差。因为
\[
\rho(X,Y)=\frac{\mathrm{Cov}(X,Y)}{\sqrt{DX}\sqrt{DY}},
\]
所以
\[
\mathrm{Cov}(X,Y)=\rho(X,Y)\sqrt{DX}\sqrt{DY}
=0.5\times \sqrt{1}\times \sqrt{4}=1.
\]

于是
\[
D(X+Y) = DX + DY + 2\,\mathrm{Cov}(X,Y) = 1 + 4 + 2 = 7.
\]

再看第二问。由期望的线性性质，
\[
E(X+Y)=EX+EY=-2+2=0.
\]
因此事件 \(\{|X+Y|\ge 6\}\) 正好就是 \(\{|(X+Y)-E(X+Y)|\ge 6\}\)，这已经是切比雪夫不等式的标准形式，其中 \(\varepsilon=6\)。所以
\[
P\{|X+Y| \geq 6\} \leq \frac{D(X+Y)}{36} = \frac{7}{36}.
\]

故应填
\[
D(X+Y)=7,\qquad P\{|X+Y|\ge 6\}\le \frac{7}{36}.
\]
\end{solution}



\begin{exercise}\label{exer:5-chebyshev-pdf}
设随机变量 \(X\) 的概率密度为
\[
f(x) = \begin{cases} \dfrac{1}{2}\,x^2\,e^{-x}, & x > 0, \\[4pt] 0, & x \leq 0, \end{cases}
\]
利用切比雪夫不等式估计 \(P(1 < X < 5) \geq\,\underline{\hspace{2cm}}\)。
\end{exercise}
\begin{solution}
利用公式 \(\displaystyle\int_0^{+\infty} e^{-x}\,x^k\,\mathrm{d}x = k!\)（即 Gamma 函数）计算期望和方差：
\[
EX = \int_0^{+\infty} \frac{1}{2}\,x^3\,e^{-x}\,\mathrm{d}x = \frac{1}{2}\,\Gamma(4) = 3,\quad
EX^2 = \int_0^{+\infty} \frac{1}{2}\,x^4\,e^{-x}\,\mathrm{d}x = \frac{1}{2}\,\Gamma(5) = 12,
\]
\[
DX = EX^2 - (EX)^2 = 12 - 9 = 3.
\]
注意到 \(P(1 < X < 5) = P(|X - 3| < 2)\)，由切比雪夫不等式：
\[
P(|X - 3| < 2) \geq 1 - \frac{DX}{2^2} = 1 - \frac{3}{4} = \frac{1}{4}.
\]

\begin{note}
本题没有直接给出切比雪夫不等式的标准形式，需要先计算 \(EX\) 和 \(DX\) 来识别 \(|X - EX| < \varepsilon\) 的结构。此外，题目中不等号方向为"\(\geq\)"（求概率下界），因此应使用等价形式 \(P\{|X-EX|<\varepsilon\} \geq 1 - DX/\varepsilon^2\)。
\end{note}
\end{solution}



\begin{exercise}\label{exer:5-chebyshev-exact}
设随机变量 \(X\) 的概率密度为
\[
f(x) = \begin{cases} 2x, & 0 < x < 1, \\ 0, & \text{其他}, \end{cases}
\]
则 \(P(|X - EX| \geq 2\sqrt{DX})\) 等于（\quad）。\par
A. \(\dfrac{9-8\sqrt{2}}{9}\)\quad
B. \(\dfrac{6+4\sqrt{2}}{9}\)\quad
C. \(\dfrac{6-8\sqrt{2}}{9}\)\quad
D. \(\dfrac{6-4\sqrt{2}}{9}\)
\end{exercise}
\begin{solution}
应选 D。
\[
EX = \int_0^1 2x^2\,\mathrm{d}x = \frac{2}{3},\quad
EX^2 = \int_0^1 2x^3\,\mathrm{d}x = \frac{1}{2},\quad
DX = \frac{1}{2} - \frac{4}{9} = \frac{1}{18}.
\]
\(2\sqrt{DX} = 2 \times \dfrac{\sqrt{2}}{6} = \dfrac{\sqrt{2}}{3}\)，不等式 \(|X - 2/3| \geq \sqrt{2}/3\) 的解为
\(X \geq \dfrac{2+\sqrt{2}}{3}\) 或 \(X \leq \dfrac{2-\sqrt{2}}{3}\)。
由于 \(\dfrac{2+\sqrt{2}}{3} > 1\)（超出取值范围），第一部分概率为 0。只需计算：
\[
P\!\left(X \leq \frac{2-\sqrt{2}}{3}\right) = \int_0^{\frac{2-\sqrt{2}}{3}} 2x\,\mathrm{d}x = \left(\frac{2-\sqrt{2}}{3}\right)^{\!2} = \frac{6-4\sqrt{2}}{9}.
\]
\begin{note}
作为对比，切比雪夫不等式仅给出粗略估计 \(P(|X-EX| \geq 2\sqrt{DX}) \leq 1/4 = 0.25\)，而精确值为 \(\dfrac{6-4\sqrt{2}}{9} \approx 0.139\)。
\end{note}
\end{solution}



\subsection{切比雪夫大数定律}

\begin{note}
\textbf{为什么切比雪夫大数定律要假设“相互独立 + 方差一致有界”？} 因为证明的核心只有一件事：让样本平均
\[
\overline{X}_n=\frac1n\sum_{i=1}^n X_i
\]
的方差趋于 0。

若 \(X_1,\dots,X_n\) 相互独立，就有
\[
D\overline{X}_n=\frac1{n^2}\sum_{i=1}^n DX_i.
\]
这里“独立”的作用是消掉协方差项；“方差一致有界”的作用是保证上式最多像 \(C/n\) 一样衰减到 0。也就是说，这两个条件都不是为了形式整齐，而是为了保证\textbf{没有哪几个样本的波动会长期主宰平均值}。若后面的样本方差越来越大，例如与 \(i^2\) 同量级，那么平均值就未必还能稳定下来。
\end{note}

\begin{theorem}[切比雪夫大数定律] \label{thm:chebyshev-lln}
设 \(X_1, X_2, \dots, X_n, \dots\) 是相互独立的随机变量序列，若方差 \(DX_i\ (i \geq 1)\) 存在且一致有界，即存在常数 \(C > 0\)，使得
\[
DX_i \leq C \quad (\text{对一切 } i \geq 1),
\]
则对任意 \(\varepsilon > 0\)，有
\[
\lim_{n\to\infty} P\!\left\{\left|\frac{1}{n}\sum_{i=1}^{n} X_i - \frac{1}{n}\sum_{i=1}^{n} EX_i\right| < \varepsilon\right\} = 1.
\]
\end{theorem}
\begin{proof}
令 \(\overline{X}_n = \dfrac{1}{n}\sum_{i=1}^{n} X_i\)。由期望的线性性质，\(E\overline{X}_n = \dfrac{1}{n}\sum EX_i\)。由独立性，
\[
D\overline{X}_n = \frac{1}{n^2}\sum_{i=1}^{n} DX_i \leq \frac{nC}{n^2} = \frac{C}{n} \xrightarrow{n\to\infty} 0.
\]
对 \(\overline{X}_n\) 应用切比雪夫不等式：
\[
P\{|\overline{X}_n - E\overline{X}_n| \geqslant \varepsilon\} \leqslant \frac{D\overline{X}_n}{\varepsilon^2} \leqslant \frac{C}{n\varepsilon^2} \xrightarrow{n\to\infty} 0.
\]
\end{proof}

\begin{note}
切比雪夫大数定律的意义在于：当独立随机变量的方差一致有界时，\(n\) 个随机变量的算术平均值依概率收敛于其期望的算术平均值。

\textbf{使用提醒：} 定理正文写成
\[
\frac1n\sum_{i=1}^n EX_i
\]
而不是直接写 \(\mu\)，是因为它允许 \(X_1,X_2,\dots\) 的期望彼此不同。只有在每个 \(EX_i\) 都等于同一个常数 \(\mu\) 时，结论才简化为
\[
\frac{1}{n}\sum_{i=1}^{n} X_i \stackrel{P}{\longrightarrow} \mu.
\]
\end{note}

\begin{exercise}\label{exer:5-lln-exponential}
设 \(X_1, X_2, \dots, X_n, \dots\) 是相互独立的随机变量序列，\(X_n\) 服从参数为 \(n\ (n \geq 1)\) 的指数分布，则下列随机变量序列中\textbf{不}满足切比雪夫大数定律条件的是（\quad）。\par
A. \(X_1, \dfrac{1}{2}X_2, \dots, \dfrac{1}{n}X_n, \dots\)\quad
B. \(X_1, X_2, \dots, X_n, \dots\)\par
C. \(X_1, 2X_2, \dots, nX_n, \dots\)\quad
D. \(X_1, 2^2X_2, \dots, n^2X_n, \dots\)
\end{exercise}

\begin{solution}
应选 D。

切比雪夫大数定律要求：\(\{X_n\}\) 相互独立，方差存在且一致有界（\(DX_n \leq C\)）。

由题设，\(DX_n = \dfrac{1}{n^2}\)（指数分布 \(E(\lambda)\) 的方差为 \(\frac{1}{\lambda^2}\)，此处 \(\lambda = n\)）。

逐一检验：
\begin{itemize}
    \item A: \(D\!\left(\dfrac{1}{n}X_n\right) = \dfrac{1}{n^2}\,DX_n = \dfrac{1}{n^4} \leq 1\) \checkmark
    \item B: \(DX_n = \dfrac{1}{n^2} \leq 1\) \checkmark
    \item C: \(D(nX_n) = n^2\,DX_n = 1 \leq 1\) \checkmark
    \item D: \(D(n^2 X_n) = n^4\,DX_n = n^2\)，当 \(n \to \infty\) 时无上界 \(\times\)
\end{itemize}
故选 D。
\end{solution}

\subsection{伯努利大数定律}

\begin{theorem}[伯努利大数定律] \label{thm:bernoulli-lln}
设 \(\mu_n\) 是 \(n\) 重伯努利试验中事件 \(A\) 发生的次数，\(p\) 是事件 \(A\) 在每次试验中发生的概率（\(0 < p < 1\)），则对任意 \(\varepsilon > 0\)，有
\[
\lim_{n\to\infty} P\!\left\{\left|\frac{\mu_n}{n} - p\right| < \varepsilon\right\} = 1,
\]
即频率 \(\frac{\mu_n}{n}\) 依概率收敛于概率 \(p\)。
\end{theorem}

\begin{note}
\textbf{为什么这里强调每次试验都有同一个概率 \(p\)？} 因为伯努利大数定律要说明的是“频率最终稳定到某个固定常数”。若第 \(i\) 次试验成功概率是 \(p_i\)，而不是同一个 \(p\)，那么频率更自然地靠近的是平均概率
\[
\frac1n\sum_{i=1}^n p_i,
\]
而不是某个统一的 \(p\)。因此这里的“同分布”不是装饰条件，而是为了让“极限到底是谁”这件事先被说清楚。
\end{note}
\begin{proof}
令 \(X_i\) 表示第 \(i\) 次试验中 \(A\) 发生的次数，则 \(X_i \sim B(1,p)\)，\(EX_i = p\)，\(DX_i = p(1-p) \leq \dfrac{1}{4}\)。且 \(\mu_n = \sum_{i=1}^{n} X_i\)，\(\dfrac{\mu_n}{n} = \overline{X}_n\)。

由于 \(\{X_i\}\) 相互独立，方差一致有界（\(\leq 1/4\)），由切比雪夫大数定律直接可得
\[
\overline{X}_n = \frac{\mu_n}{n} \xrightarrow{P} p.
\]
\end{proof}

\begin{note}
伯努利大数定律是切比雪夫大数定律的特例（取 \(X_i \sim B(1,p)\)），它为"频率稳定于概率"这一经验事实提供了严格的数学证明，也是用频率定义概率的理论依据。

\textbf{自学追问：} 为什么“频率”会和样本均值扯上关系？因为把第 \(i\) 次试验中事件 \(A\) 是否发生记成 0-1 变量 \(X_i\) 后，发生则 \(X_i=1\)，不发生则 \(X_i=0\)。于是
\[
\frac{\mu_n}{n}=\frac{1}{n}\sum_{i=1}^{n}X_i=\bar X_n,
\]
也就是说，频率本质上就是这一组 0-1 样本的算术平均值。

在数理统计中，经验分布函数
\[
F_n(x) = \frac{\text{样本值中} \leq x \text{的个数}}{n}
\]
正是事件 \(\{X \leq x\}\) 在 \(n\) 次试验中出现的频率。由伯努利大数定律，当 \(n\) 充分大时，\(F_n(x)\) 可作为总体分布函数 \(F(x)\) 的近似，且 \(n\) 越大近似越好。
\end{note}

\begin{exercise}\label{exer:5-empirical}
设 \((2, 1, 5, 2, 1, 3, 1)\) 是来自总体 \(X\) 的简单随机样本值，求总体 \(X\) 的经验分布函数 \(F_7(x)\)。
\end{exercise}

\begin{solution}
将样本值按从小到大排列：\(1, 1, 1, 2, 2, 3, 5\)。

\[
F_7(x) = \begin{cases}
0, & x < 1, \\
\dfrac{3}{7}, & 1 \leq x < 2, \\[6pt]
\dfrac{5}{7}, & 2 \leq x < 3, \\[6pt]
\dfrac{6}{7}, & 3 \leq x < 5, \\[6pt]
1, & x \geq 5.
\end{cases}
\]
\end{solution}

\subsection{辛钦大数定律}

\begin{note}
\textbf{切比雪夫型大数定律和辛钦大数定律不要机械比较“谁更强”。} 两者处理的是两种不同的信息结构：

\textbf{切比雪夫型结论}允许 \(X_1,X_2,\dots\) 的分布彼此不同，所以极限写成
\[
\frac1n\sum_{i=1}^n EX_i;
\]
代价是要额外控制方差一致有界。

\textbf{辛钦型结论}把条件收紧成“独立同分布”，于是极限自然就是同一个 \(\mu\)；回报是它只要求期望存在，不再要求方差存在。

所以它们不是简单的“一个包含另一个”，而更像是一种交换：\textbf{用“同分布”换取“更弱的矩条件”}。
\end{note}

\begin{theorem}[辛钦大数定律] \label{thm:khinchin-lln}
设 \(X_1, X_2, \dots, X_n, \dots\) 是独立同分布的随机变量序列，若数学期望 \(EX_i = \mu\ (i = 1, 2, \dots)\) 存在，则对任意 \(\varepsilon > 0\)，有
\[
\lim_{n\to\infty} P\!\left\{\left|\frac{1}{n}\sum_{i=1}^{n} X_i - \mu\right| < \varepsilon\right\} = 1,
\]
即 \(\dfrac{1}{n}\displaystyle\sum_{i=1}^{n} X_i \stackrel{P}{\longrightarrow} \mu\)。
\end{theorem}

\begin{proof}[证明思路]
在方差 \(DX_k = \sigma^2\ (k=1,2,\dots)\) 存在的条件下给出证明。因为
\[
E\!\left(\frac{1}{n}\sum_{k=1}^{n} X_k\right) = \frac{1}{n}\sum_{k=1}^{n} EX_k = \mu,
\]
又由独立性得
\[
D\!\left(\frac{1}{n}\sum_{k=1}^{n} X_k\right) = \frac{1}{n^2}\sum_{k=1}^{n} DX_k = \frac{\sigma^2}{n}.
\]
由切比雪夫不等式：
\[
P\!\left\{\left|\frac{1}{n}\sum_{k=1}^{n} X_k - \mu\right| \geq \varepsilon\right\}
\leq \frac{\sigma^2/n}{\varepsilon^2} \xrightarrow{n\to\infty} 0.
\]
\end{proof}

\begin{note}
辛钦大数定律不要求方差存在（上述证明是方差存在时的特例），仅要求期望存在，适用范围更广。

\textbf{为什么“只要求期望存在”已经是关键门槛？} 因为样本均值想稳定到的对象就是总体均值 \(\mu\)。若连 \(EX_1\) 都不存在，那么“平均值最后应该靠近谁”本身就没有清楚的中心，自然也无从谈起收敛到某个固定均值。典型地，柯西分布没有数学期望，它的样本均值就不会像正态情形那样稳定。

\textbf{使用提醒：} 若题目问的不是 \(\frac1n\sum X_i\)，而是
\[
\frac1n\sum h(X_i)
\]
这类“样本函数平均值”的极限，只要 \(X_i\) 独立同分布且 \(E|h(X_1)|<\infty\)，就可以令 \(Y_i=h(X_i)\)，再把辛钦大数定律直接用在 \(Y_i\) 上。

进一步地，若 \(X_1, X_2, \dots\) 独立同分布且 \(E|X_1|^k<\infty\)，则把 \(Y_i=X_i^k\) 视为新的独立同分布随机变量，再应用辛钦大数定律可得
\[
\frac{1}{n}\sum_{i=1}^{n} X_i^k \stackrel{P}{\longrightarrow} E(X_1^k).
\]

\textbf{易错点：} 不要因为证明思路里用了方差，就误以为辛钦大数定律要求方差存在。定理本身只要求期望存在；方差存在只是一个更容易说明的特殊情形。

在数理统计中，判断估计量是否具有\textbf{一致性}（相合性）时，常常需要利用依概率收敛和大数定律。
\end{note}

\begin{exercise}\label{exer:5-khinchin-x2}
设随机变量序列 \(X_1, X_2, \dots, X_n, \dots\) 独立同分布，\(X_1\) 的概率密度为
\[
f(x) = \begin{cases} 1 - |x|, & |x| < 1, \\ 0, & \text{其他}, \end{cases}
\]
则当 \(n \to \infty\) 时，\(\dfrac{1}{n}\displaystyle\sum_{i=1}^{n} X_i^2\) 依概率收敛于（\quad）。\par
A. \(\dfrac{1}{8}\)\quad B. \(\dfrac{1}{6}\)\quad C. \(\dfrac{1}{3}\)\quad D. \(\dfrac{1}{2}\)
\end{exercise}
\begin{solution}
应选 B。

本题虽然表面上不是样本均值 \(\dfrac1n\sum X_i\)，但它是 \(\dfrac1n\sum_{i=1}^n X_i^2\)，即“样本平方的平均值”。因此可以令 \(Y_i=X_i^2\)。
由于 \(X_1,X_2,\dots\) 独立同分布，\(Y_1,Y_2,\dots\) 也独立同分布；又因为 \(X_i\) 的取值落在 \([-1,1]\) 上，故 \(E(Y_i)=E(X_i^2)\) 一定存在。于是由辛钦大数定律，
\[
\frac{1}{n}\sum_{i=1}^{n} X_i^2 \stackrel{P}{\longrightarrow} E(X_1^2).
\]

下面只需求 \(E(X_1^2)\)。由密度关于 0 对称，可写成
\[
E(X_1^2) = \int_{-1}^{1} x^2(1-|x|)\,\mathrm{d}x = 2\int_0^1 x^2(1-x)\,\mathrm{d}x
= 2\!\left[\frac{x^3}{3} - \frac{x^4}{4}\right]_0^1 = 2\!\left(\frac{1}{3}-\frac{1}{4}\right) = \frac{1}{6}.
\]

因此极限为 \(\dfrac16\)，应选 B。
\end{solution}

\begin{exercise}\label{exer:5-lln-dice}
将一枚骰子重复掷 \(n\) 次，则当 \(n \to \infty\) 时，\(n\) 次掷出点数的算术平均值 \(\overline{X}_n\) 依概率收敛于\,\underline{\hspace{2cm}}。
\end{exercise}
\begin{solution}
应填 \(\dfrac{7}{2}\)。

设第 \(i\) 次掷出的点数为 \(X_i\)。由于每次掷骰子互不影响，且每次都等可能取 \(1,2,3,4,5,6\)，所以
\[
X_1, X_2, \dots, X_n
\]
是独立同分布随机变量序列。

题目问的是算术平均值
\[
\overline{X}_n=\frac1n\sum_{i=1}^n X_i,
\]
这正是辛钦大数定律的标准应用场景。先求单次掷骰子的期望：
\[
P(X_i=k)=\frac16,\qquad k=1,2,\dots,6,
\]
因此
\[
EX_i = \frac{1 + 2 + 3 + 4 + 5 + 6}{6} = \frac{7}{2}.
\]

由辛钦大数定律，
\[
\overline{X}_n = \frac{1}{n}\sum_{i=1}^{n} X_i \stackrel{P}{\longrightarrow} EX_i = \frac{7}{2}.
\]

所以应填 \(\dfrac{7}{2}\)。
\end{solution}



\section{中心极限定理}

中心极限定理研究的是：在什么条件下，大量独立随机变量之和的分布趋近于正态分布。它是正态分布在概率论中占据核心地位的根本原因，也是大样本统计推断的理论基础。

从自学角度看，本章两类极限定理的分工可以这样记忆：\textbf{大数定律回答“平均值最终靠近谁”，中心极限定理回答“围绕这个中心的波动近似长什么样”}。前者关注收敛目标，后者关注大样本下的分布形状与概率近似。

\begin{note}
\textbf{CLT 比大数定律多回答了一层问题。} 大数定律告诉我们
\[
\overline{X}_n-\mu \stackrel{P}{\longrightarrow} 0,
\]
但它没有告诉我们：误差通常有多大？怎么近似算概率？中心极限定理补的正是这一块。它说明当 \(n\) 很大时，\(\overline{X}_n-\mu\) 虽然越来越小，但其典型波动量级大约是 \(1/\sqrt{n}\)；等价地，\(\sum_{i=1}^n X_i-n\mu\) 的典型波动量级大约是 \(\sqrt{n}\)。

这也解释了为什么 CLT 的标准化分母不是 \(n\)，而是 \(\sqrt{n}\sigma\)：我们不是在看“总量本身多大”，而是在看\textbf{围绕中心的随机起伏到底有多大}。
\end{note}

\subsection{列维--林德伯格中心极限定理}

\begin{theorem}[列维--林德伯格中心极限定理] \label{thm:levy-lindberg}
设 \(X_1, X_2, \dots, X_n, \dots\) 是独立同分布的随机变量序列，且数学期望 \(EX_i = \mu\) 和方差 \(DX_i = \sigma^2 > 0\) 存在 \((i = 1, 2, \dots)\)，则对任意实数 \(x\)，有
\[
\lim_{n\to\infty} P\!\left\{\frac{\sum_{i=1}^{n} X_i - n\mu}{\sqrt{n}\,\sigma} \leq x\right\} = \Phi(x).
\]
\end{theorem}

\begin{note}
应用列维--林德伯格中心极限定理时，需满足三个条件：\textbf{独立、同分布、期望和方差存在}，缺一不可。

\textbf{自学追问：} 为什么总要先做标准化？因为不同题目中的和 \(\sum X_i\) 其中心位置和波动大小都不同，只有先减去 \(n\mu\)、再除以 \(\sqrt{n}\sigma\)，才能统一化成标准正态分布函数 \(\Phi\) 可以处理的形式。

当 \(n\) 足够大时，\(\displaystyle\sum_{i=1}^{n} X_i\) 近似服从正态分布 \(N(n\mu,\, n\sigma^2)\)，因此
\[
P\!\left\{a < \sum_{i=1}^{n} X_i < b\right\} \approx \Phi\!\left(\frac{b - n\mu}{\sqrt{n}\,\sigma}\right) - \Phi\!\left(\frac{a - n\mu}{\sqrt{n}\,\sigma}\right).
\]

\textbf{使用提醒：} 若题目问的是样本均值 \(\overline{X}_n\)，也完全可以使用中心极限定理，因为
\[
\frac{\overline{X}_n-\mu}{\sigma/\sqrt{n}}
=\frac{\sum_{i=1}^n X_i-n\mu}{\sqrt{n}\,\sigma}.
\]
也就是说，“和”的标准化与“均值”的标准化本质上是一回事。

\textbf{易错点：} 中心极限定理近似的是”和或均值的分布”，不是单个 \(X_i\) 的分布；另外，若样本量不大，或者题目不满足独立同分布、方差存在等条件，就不能机械套用。
\end{note}

\begin{note}
\textbf{为什么 CLT 的条件正好是“独立同分布 + 均值方差存在”？}
\begin{itemize}
    \item \textbf{独立}：避免很多变量同涨同跌，使总和里出现大规模同步偏移；
    \item \textbf{同分布}：保证每一项贡献的尺度相近，研究的是“很多相似的小扰动叠加”；
    \item \textbf{方差存在}：保证总波动能被 \(\sqrt{n}\) 这个量级刻画；若尾部太重、方差发散，正态近似通常会失效；
    \item \textbf{\(\sigma^2>0\)}：排除退化情形。若方差为 0，则每个 \(X_i\) 都几乎处处等于常数，根本不存在需要近似的“波动形状”。
\end{itemize}

所以 CLT 不是“只要样本多就一定正态”，而是“很多独立、尺度相近、单个波动不过分极端的小贡献叠加后，才会出现正态形状”。把这个图像想清楚，比背条件列表更不容易混淆。
\end{note}

\begin{note}
\textbf{把 CLT 题翻译成标准形式时，先抓三样东西：中心、尺度、类型。}
\begin{itemize}
    \item \textbf{中心}：通常是 \(n\mu\) 或 \(np\)；
    \item \textbf{尺度}：通常是 \(\sqrt{n}\sigma\) 或 \(\sqrt{np(1-p)}\)；
    \item \textbf{类型}：若原变量是计数型离散变量，最后常要顺手检查是否需要连续性修正。
\end{itemize}

这一步不是机械技巧，而是在回答：正态曲线是在“模仿”哪一个离散分布的整体形状。
\end{note}

\begin{note}
\textbf{过渡示例（CLT 的机械操作）：} 设 $X_1,\dots,X_{100}$ iid，$EX_i=2$，$DX_i=4$。求 $P\{180<\sum X_i<220\}$。

\textbf{步骤：}
\begin{enumerate}
    \item 确认条件：iid，期望方差存在 $\checkmark$
    \item 写出和的参数：$E(\sum X_i)=100\times 2=200$，$D(\sum X_i)=100\times 4=400$，$\sigma_{\text{和}}=20$
    \item 标准化：$P\left\{\frac{180-200}{20}<\frac{\sum X_i-200}{20}<\frac{220-200}{20}\right\}=P\{-1<Z<1\}$
    \item 查表：$\Phi(1)-\Phi(-1)=2\Phi(1)-1=2\times 0.8413-1=0.6826$
\end{enumerate}
\textbf{口诀：} 求参数 $\to$ 标准化 $\to$ 查 $\Phi$ 表。每道 CLT 题都是这三步。
\end{note}

\begin{exercise}\label{exer:5-clt-condition}
设随机变量 \(X_1, X_2, \dots, X_n\) 相互独立，\(S_n = X_1 + X_2 + \cdots + X_n\)。根据列维--林德伯格定理，当 \(n\) 充分大时，\(S_n\) 近似服从正态分布，只要（\quad）。\par
A. 有相同的期望和方差\quad B. 服从同一离散型分布\quad
C. 服从同一指数分布\quad D. 服从同一连续型分布
\end{exercise}

\begin{solution}
应选 C。

这类选择题最稳妥的做法不是凭感觉，而是按中心极限定理的三个条件逐项检查：
\[
\text{独立}+\text{同分布}+\text{期望、方差存在}.
\]
题干已经说明 \(X_1,\dots,X_n\) 相互独立，所以只需继续检查“是否同分布”以及“期望、方差是否一定存在”。

\begin{itemize}
    \item A：有相同的期望和方差 \(\Rightarrow\) 未必同分布，不满足条件。
    \item B：服从同一离散型分布 \(\Rightarrow\) 方差未必存在（反例：若 \(P(X=k)=c/k^3\ (k=1,2,\dots)\)，则 \(EX^2=\sum k^2\cdot c/k^3=\sum c/k\) 发散，方差不存在）。
    \item C：指数分布 \(E(\lambda)\) 的期望和方差分别为 \(1/\lambda\) 和 \(1/\lambda^2\)，均存在，满足全部三个条件 \checkmark
    \item D：服从同一连续型分布 \(\Rightarrow\) 方差未必存在（如柯西分布）。
\end{itemize}

因此，只有 C 能保证中心极限定理的全部前提成立。
\end{solution}

\begin{exercise}\label{exer:5-clt-exp100}
设随机变量 \(X_1, X_2, \dots, X_{100}\) 相互独立，均服从参数为 \(\dfrac{1}{2}\) 的指数分布，利用中心极限定理计算
\[
P\!\left(\sum_{i=1}^{100} X_i < 240\right) =\,\underline{\hspace{2cm}}.
\]
\end{exercise}

\begin{solution}
先识别题型：题目要求的是 100 个独立同分布随机变量之和的概率，而 \(n=100\) 已经较大，适合用中心极限定理。

\(X_i \sim E\!\left(\dfrac{1}{2}\right)\)，因此
\[
\mu=EX_i=2,\qquad \sigma^2=DX_i=4.
\]

设
\[
S=\sum_{i=1}^{100} X_i.
\]
则
\[
ES=100\times 2=200,\qquad DS=100\times 4=400.
\]

由列维--林德伯格中心极限定理，
\[
\frac{S-200}{20}\approx N(0,1),
\]
即 \(S\) 近似服从正态分布 \(N(200,400)\)，于是
\[
P(S < 240) = P\!\left(\frac{S - 200}{20} < \frac{240 - 200}{20}\right) \approx \Phi(2).
\]

因此应填 \(\Phi(2)\)。
\end{solution}

\begin{exercise}\label{exer:5-clt-exp36}
设某种元件的寿命（小时）服从指数分布，其概率密度函数为
\[
f(x) = \begin{cases} \dfrac{1}{100}\,e^{-x/100}, & x > 0, \\[4pt] 0, & x \leq 0, \end{cases}
\]
各元件的寿命相互独立。随机取 36 个元件，求寿命总和大于 4200 小时的概率近似值。
\end{exercise}

\begin{solution}
由密度函数知 \(\lambda = \dfrac{1}{100}\)。设第 \(i\) 只元件的寿命为 \(X_i\)，则 \(X_i \sim E\!\left(\dfrac{1}{100}\right)\)。

\(EX_i = 100\)，\(DX_i = 10000\)。令 \(Y = \displaystyle\sum_{i=1}^{36} X_i\)，则
\[
EY = 3600,\quad DY = 36 \times 10000 = 360000.
\]

由列维--林德伯格定理，\(Y\) 近似服从 \(N(3600, 360000)\)，故
\[
P(Y > 4200) = 1 - P(Y \leq 4200) \approx 1 - \Phi\!\left(\frac{4200 - 3600}{600}\right) = 1 - \Phi(1) \approx 0.1587.
\]
\end{solution}

\begin{exercise}\label{exer:5-clt-box}
生产线生产的产品成箱包装，每箱质量是随机的。假设平均每箱重 50 千克，标准差为 5 千克。若用载重为 5 吨的汽车承运，试用中心极限定理说明每辆车最多可以装多少箱，才能保证不超载的概率大于 0.977。\ (\(\Phi(2) = 0.977\))
\end{exercise}

\begin{solution}
设 \(X_i\ (i=1,2,\dots,n)\) 为第 \(i\) 箱的质量（千克），\(X_i\) 独立同分布，\(EX_i = 50\)，\(\sqrt{DX_i} = 5\)。

\(n\) 箱总质量 \(T_n = \displaystyle\sum_{i=1}^{n} X_i\)，\(ET_n = 50n\)，\(\sqrt{DT_n} = 5\sqrt{n}\)。

由列维--林德伯格定理，\(T_n\) 近似服从 \(N(50n,\, 25n)\)：
\[
P(T_n \leq 5000) = P\!\left(\frac{T_n - 50n}{5\sqrt{n}} \leq \frac{5000 - 50n}{5\sqrt{n}}\right) \approx \Phi\!\left(\frac{1000 - 10n}{\sqrt{n}}\right) > 0.977 = \Phi(2),
\]
即 \(\dfrac{1000 - 10n}{\sqrt{n}} > 2\)。

令 \(t = \sqrt{n}\)，则 \(1000 - 10t^2 > 2t\)，即 \(10t^2 + 2t - 1000 < 0\)。解得 \(t < \dfrac{-2 + \sqrt{4+40000}}{20} = \dfrac{-2+2\sqrt{10001}}{20} \approx 9.950\)，故 \(n < 99.0\)。

验证边界值：\(n=98\) 时 \(\dfrac{1000-980}{\sqrt{98}} = \dfrac{20}{9.90} \approx 2.02 > 2\) ；\(n=99\) 时 \(\dfrac{1000-990}{\sqrt{99}} = \dfrac{10}{9.95} \approx 1.005 < 2\) 。因此最多装 \(\mathbf{98}\) 箱。
\end{solution}

\begin{exercise}\label{exer:5-clt-exp-norm}
设 \(X_1, X_2, \dots, X_n, \dots\) 为独立同分布随机变量序列，且均服从参数为 \(\lambda\ (\lambda > 1)\) 的指数分布，记 \(\Phi(x)\) 为标准正态分布函数，则（\quad）。\par
A. \(\displaystyle\lim_{n\to\infty} P\!\left\{\frac{\sum_{i=1}^n X_i - n\lambda}{\lambda\sqrt{n}} \leq x\right\} = \Phi(x)\)~~~~~~
B. \(\displaystyle\lim_{n\to\infty} P\!\left\{\frac{\sum_{i=1}^n X_i - n\lambda}{\sqrt{\lambda n}} \leq x\right\} = \Phi(x)\)\par
C. \(\displaystyle\lim_{n\to\infty} P\!\left\{\frac{\lambda\sum_{i=1}^n X_i - n}{\sqrt{n}} \leq x\right\} = \Phi(x)\)~~~~~~
D. \(\displaystyle\lim_{n\to\infty} P\!\left\{\frac{\sum_{i=1}^n X_i - \lambda}{\sqrt{n\lambda}} \leq x\right\} = \Phi(x)\)
\end{exercise}
\begin{solution}
应选 C。

这里最容易混淆的是：指数分布 \(E(\lambda)\) 中的 \(\lambda\) 是\textbf{速率参数}，不是均值。故
\[
EX_i=\frac1\lambda,\qquad DX_i=\frac1{\lambda^2}.
\]

设
\[
S_n=\sum_{i=1}^{n} X_i.
\]
则
\[
ES_n=\frac{n}{\lambda},\qquad DS_n=\frac{n}{\lambda^2}.
\]

由列维--林德伯格中心极限定理，正确的标准化应为
\[
\frac{S_n - ES_n}{\sqrt{DS_n}} = \frac{S_n - n/\lambda}{\sqrt{n}/\lambda}
= \frac{\lambda S_n - n}{\sqrt{n}} \to N(0,1),
\]
这恰好对应选项 C。

其余选项的问题分别是：
\begin{itemize}
    \item A、B 把中心位置错写成了 \(n\lambda\)，而正确的均值应是 \(n/\lambda\)；
    \item B、D 的标准化分母也与真实方差不符。
\end{itemize}

所以应选 C。
\end{solution}

\begin{exercise}\label{exer:5-clt-poisson}
设 \(X_1, X_2, \dots, X_n\) 相互独立，且均服从泊松分布 \(P(2)\)，则由列维--林德伯格中心极限定理
\[
\lim_{n\to\infty} P\!\left(\frac{\sum_{i=1}^n X_i - 2n}{\sqrt{n}} < 2\right) = \text{（\quad）}.
\]
A. \(\Phi(\sqrt{2})\)\quad B. \(\Phi(2)\)\quad C. \(\Phi(1)\)\quad D. \(1\)
\end{exercise}
\begin{solution}
应选 A。\(X_i \sim P(2)\)，故 \(EX_i = 2\)，\(DX_i = 2\)。由中心极限定理：
\[
\frac{\sum X_i - n \cdot 2}{\sqrt{n \cdot 2}} \to N(0,1).
\]
题目中式子变形：
\[
P\!\left(\frac{\sum X_i - 2n}{\sqrt{n}} < 2\right)
= P\!\left(\frac{\sum X_i - 2n}{\sqrt{2n}} < \frac{2}{\sqrt{2}}\right)
= P(Z < \sqrt{2}) = \Phi(\sqrt{2}).
\]
\end{solution}



\begin{exercise}\label{exer:5-clt-moment}
假设 \(X_1, X_2, \dots, X_n\) 是来自总体 \(X\) 的简单随机样本，已知 \(E(X^k) = \alpha_k\ (k = 1,2,3,4)\)。证明：当 \(n\) 充分大时，随机变量
\[
Z_n = \frac{1}{n}\sum_{i=1}^{n} X_i^2
\]
近似服从正态分布，并指出其分布参数。
\end{exercise}
\begin{solution}
由题设，\(X_1, X_2, \dots, X_n\) 独立同分布，故 \(X_1^2, X_2^2, \dots, X_n^2\) 也独立同分布，且
\[
E(X_i^2) = \alpha_2,\quad D(X_i^2) = E(X_i^4) - [E(X_i^2)]^2 = \alpha_4 - \alpha_2^2.
\]
由列维--林德伯格中心极限定理，
\[
\frac{\sum_{i=1}^{n} X_i^2 - n\alpha_2}{\sqrt{n(\alpha_4 - \alpha_2^2)}} \to N(0,1),
\]
等价于
\[
\frac{Z_n - \alpha_2}{\sqrt{(\alpha_4 - \alpha_2^2)/n}} \to N(0,1),
\]
因此当 \(n\) 充分大时，\(Z_n\) 近似服从正态分布，参数为
\[
\mu = \alpha_2,\quad \sigma^2 = \frac{\alpha_4 - \alpha_2^2}{n}.
\]
\end{solution}


\subsection{棣莫弗--拉普拉斯中心极限定理}

\begin{note}
\textbf{二项分布什么时候更应该想“正态近似”，什么时候更应该想“泊松近似”？} 关键不只是 \(n\) 大不大，而是分布形状是否已经足够接近对称。

当 \(p\) 很小而 \(np\) 只是中等大小时，“成功”是稀有事件，分布大量集中在 0 附近，右尾较长，偏斜明显，此时泊松近似往往更自然；当 \(np\) 和 \(n(1-p)\) 都不小，成功和失败两边都积累了足够多的小随机扰动，二项分布在中心附近会变得更对称，这时正态近似才更稳。

所以不要把“\(n\) 大”机械翻译成“必用 CLT”。真正该问的是：\textbf{这个计数变量现在更像“稀有事件计数”，还是更像“许多小波动叠加后的对称总量”？}
\end{note}

\begin{theorem}[棣莫弗--拉普拉斯中心极限定理] \label{thm:demoivre-laplace}
设随机变量 \(Y_n \sim B(n, p)\ (0 < p < 1,\ n \geq 1)\)，则对任意实数 \(x\)，有
\[
\lim_{n\to\infty} P\!\left\{\frac{Y_n - np}{\sqrt{np(1-p)}} \leq x\right\} = \Phi(x).
\]
\end{theorem}

\begin{note}
棣莫弗--拉普拉斯定理是列维--林德伯格定理在二项分布上的特例。若将 \(Y_n\) 看作 \(n\) 个独立同分布的 \(B(1,p)\) 随机变量之和，即 \(Y_n = \sum_{i=1}^{n} X_i\)，其中 \(X_i \sim \begin{pmatrix} 1 & 0 \\ p & 1-p \end{pmatrix}\)，则 \(\mu = p\)，\(\sigma^2 = p(1-p)\)，代入列维--林德伯格定理即得。

\textbf{使用提醒：} 只要题目本质上是在问“\(n\) 次独立成败试验中的成功总次数”，就应先把它识别成二项分布 \(B(n,p)\)，再把中心 \(np\) 和标准差 \(\sqrt{np(1-p)}\) 写出来。

\textbf{易错点：} 标准化分母必须是 \(\sqrt{np(1-p)}\)，不能只写成 \(\sqrt{np}\) 或 \(\sqrt{n}\)。只有在 \(p=\dfrac12\) 等特殊情形下，它才会进一步化简。若题目要求更高精度，做区间概率时还要注意连续性修正。

\textbf{二项分布概率计算的三种方法}：设 \(X \sim B(n, p)\)。
\begin{enumerate}
    \item \(n\) 较小（\(n \leq 10\)）：直接用公式 \(P\{X=k\} = \mathrm{C}_n^k\,p^k(1-p)^{n-k}\)。
    \item \(n\) 较大且 \(p\) 较小（\(n > 10,\ p < 0.1,\ np\) 适中）：泊松近似 \(\displaystyle P\{X=k\} \approx \dfrac{\lambda^k}{k!}\,e^{-\lambda}\)，其中 \(\lambda = np\)。
    \item \(n\) 较大且 \(p\) 不太大（\(np \geq 10\)）：正态近似（中心极限定理）
    \[
    P\{a < X < b\} \approx \Phi\!\left(\frac{b-np}{\sqrt{np(1-p)}}\right) - \Phi\!\left(\frac{a-np}{\sqrt{np(1-p)}}\right).
    \]
\end{enumerate}
\end{note}

\begin{note}
\textbf{为什么二项分布常常要做连续性修正？} 因为二项分布只在整数点上取值，而正态分布是连续曲线。直接拿连续曲线去替代离散柱状图时，边界点的概率质量会被“挤”到两侧。

所以当题目写成 \(P\{a \leq X \leq b\}\) 这类整数区间时，常常会先把它想成
\[
P\{a-0.5 < X < b+0.5\},
\]
再去做正态近似。样本量越大，这个修正越不显眼；但在 \(np\) 不算特别大时，它往往能明显改善答案。
\end{note}

\begin{exercise}\label{exer:5-dml-identify}
从废品率为 0.03 的大量产品中随机抽取 1000 个，根据中心极限定理，废品数 \(X\) 近似服从的分布为\,\underline{\hspace{2cm}}（要求写出分布的参数）。
\end{exercise}
\begin{solution}
废品数 \(X\) 本质上是 1000 次独立伯努利试验中“废品”出现的总次数，因此
\[
X\sim B(1000,0.03).
\]

当 \(n\) 很大时，二项分布可以用棣莫弗--拉普拉斯中心极限定理作正态近似，其近似参数取二项分布的均值和方差：
\[
EX=np=1000\times 0.03=30,\qquad DX=np(1-p)=1000\times 0.03\times 0.97=29.1.
\]

因此 \(X\approx N(30,\ 29.1)\)。

所以应填“近似服从 \(N(30,29.1)\)”。
\end{solution}

\begin{exercise}\label{exer:5-clt-tap}
某宿舍楼按 400 名学生的规模建造，每人每天有约 10\% 的时间占用一个水龙头（各人独立）。为能以 95\% 的概率保证用水需要，至少应安装多少个水龙头？\ (\(\Phi(1.645) = 0.95\))
\end{exercise}
\begin{solution}
先把题目建模成二项分布。对每个学生来说，在某一时刻“是否正在用水”只有两种结果：
\begin{itemize}
    \item 用水，概率为 \(0.1\)；
    \item 不用水，概率为 \(0.9\)。
\end{itemize}
又因为 400 名学生彼此独立，所以同一时刻的用水人数 \(X\) 服从 \(X\sim B(400,0.1)\)。
因此 \(EX=np=40,\ DX=np(1-p)=400\times 0.1\times 0.9=36\)。

题目要保证“安装的水龙头数至少覆盖同时用水人数”，即要求 \(P(X\le N)\ge 0.95\)。
当 \(n=400\) 较大时，可用中心极限定理把 \(X\) 近似为正态分布 \(N(40,36)\)。于是
\[
P(X \leq N) \approx \Phi\!\left(\frac{N - 40}{6}\right) \geq 0.95 = \Phi(1.645),
\]
从而
\[
\frac{N-40}{6}\ge 1.645 \implies N\ge 40+6\times 1.645=49.87.
\]
由于水龙头数量必须取整数，故至少应安装 \(N=50\) 个水龙头。
\end{solution}

\begin{exercise}\label{exer:5-clt-quality}
某生产线产品合格率为 0.8。要使一批产品的合格率落在 76\% 与 84\% 之间的概率不小于 90\%，问这批产品至少要生产多少件？\ (\(\Phi(1.64) = 0.95\))
\end{exercise}
\begin{solution}
设生产 \(n\) 件，合格数 \(X \sim B(n, 0.8)\)。要求
\[
P\!\left(0.76 \leq \frac{X}{n} \leq 0.84\right) \geq 0.90.
\]
标准化：\(E(X/n) = 0.8\)，\(\mathrm{Var}(X/n) = 0.16/n\)。
\[
P\!\left(\left|\frac{X}{n} - 0.8\right| \leq 0.04\right)
= P\!\left(|Z| \leq \frac{0.04}{\sqrt{0.16/n}}\right)
= P(|Z| \leq 0.1\sqrt{n}) = 2\Phi(0.1\sqrt{n}) - 1 \geq 0.90.
\]
\(\Phi(0.1\sqrt{n}) \geq 0.95\)，\(0.1\sqrt{n} \geq 1.64\)，\(n \geq 268.96\)，取 \(n = 269\)。
\end{solution}


\begin{exercise}\label{exer:5-dml-coin}
设 \(X_n\) 表示将一枚硬币随意投掷 \(n\) 次出现正面的次数，则（\quad）。\par
A. \(\displaystyle\lim_{n\to\infty} P\!\left\{\frac{X_n - n}{\sqrt{n}} \leq x\right\} = \Phi(x)\)~~~~
B. \(\displaystyle\lim_{n\to\infty} P\!\left\{\frac{2X_n - n}{\sqrt{n}} \leq x\right\} = \Phi(x)\)\par
C. \(\displaystyle\lim_{n\to\infty} P\!\left\{\frac{X_n - 2n}{\sqrt{n}} \leq x\right\} = \Phi(x)\)~~~~
D. \(\displaystyle\lim_{n\to\infty} P\!\left\{\frac{2X_n - 2n}{\sqrt{n}} \leq x\right\} = \Phi(x)\)
\end{exercise}

\begin{solution}
应选 B。

题目中的 \(X_n\) 表示“\(n\) 次投掷中出现正面的次数”，因此
\[
X_n\sim B\!\left(n,\frac12\right).
\]
这时
\[
np=\frac n2,\qquad np(1-p)=n\cdot \frac12\cdot \frac12=\frac n4.
\]

由棣莫弗--拉普拉斯定理，
\[
\frac{X_n - n/2}{\sqrt{n/4}} = \frac{X_n - n/2}{\sqrt{n}/2} = \frac{2X_n - n}{\sqrt{n}} \to N(0,1).
\]

因此
\[
\lim_{n\to\infty} P\!\left\{\frac{2X_n-n}{\sqrt n}\le x\right\}=\Phi(x),
\]
正好对应选项 B。

A、C、D 都把中心位置或标准化因子写错了，所以不能选。
\end{solution}


\begin{exercise}\label{exer:5-dml-burn}
设某种电气元件不能承受超负荷试验的概率为 0.05。现对 100 个这样的元件进行超负荷试验，以 \(X\) 表示不能承受试验而烧毁的元件数。根据中心极限定理，求 \(P\{5 \leq X \leq 10\}\) 的近似值。\ (\(\Phi(2.29) = 0.989\))
\end{exercise}
\begin{solution}
\(X \sim B(100, 0.05)\)，\(EX = np = 5\)，\(DX = np(1-p) = 4.75\)。由棣莫弗--拉普拉斯中心极限定理：
\[
\begin{aligned}
P\{5 \leq X \leq 10\}
&= P\!\left\{\frac{5 - 5}{\sqrt{4.75}} \leq \frac{X - 5}{\sqrt{4.75}} \leq \frac{10 - 5}{\sqrt{4.75}}\right\} \\
&= P\!\left\{0 \leq \frac{X - 5}{\sqrt{4.75}} \leq 2.29\right\}
\approx \Phi(2.29) - \Phi(0) = 0.989 - 0.5 = 0.489.
\end{aligned}
\]
\end{solution}

\begin{note}
上式未使用连续性修正。若使用连续性修正 \(P\{5 \leq X \leq 10\} \approx P\{4.5 < X < 10.5\}\)，则结果约为 0.585，与不修正的结果 0.489 有较大差异。考研题目中是否使用连续性修正取决于出题者意图，一般以题目给出的参考数据为准。
\end{note}


\begin{exercise}\label{exer:5-clt-bernoulli-sum}
设 \(X_1, X_2, \dots, X_{100}\) 为来自总体 \(X\) 的简单随机样本，其中 \(P\{X=0\} = P\{X=1\} = \dfrac{1}{2}\)。利用中心极限定理求 \(P\!\left\{\displaystyle\sum_{i=1}^{100} X_i \leq 55\right\}\) 的近似值（\quad）。\par
A. \(1 - \Phi(1)\)\quad B. \(\Phi(1)\)\quad C. \(1 - \Phi(0.2)\)\quad D. \(\Phi(0.2)\)
\end{exercise}
\begin{solution}
应选 B。

因为 \(P\{X=0\}=P\{X=1\}=\dfrac12\)，所以每个样本变量都服从伯努利分布
\[
X_i\sim B\!\left(1,\frac12\right),
\]
并且 \(X_1,\dots,X_{100}\) 相互独立同分布。

设
\[
S=\sum_{i=1}^{100} X_i.
\]
则
\[
ES=100\times \frac12=50,\qquad DS=100\times \frac14=25.
\]

由于样本量 \(n=100\) 较大，可用中心极限定理把 \(S\) 近似为正态分布 \(N(50,25)\)。于是
\[
P(S \leq 55) = P\!\left(\frac{S - 50}{5} \leq \frac{55 - 50}{5}\right) \approx \Phi(1).
\]

因此应选 B。
\end{solution}

\begin{exercise}\label{exer:5-climb}
在一次集体登山活动中，每个人意外受伤的概率为 1\%，各人是否受伤相互独立。\par
(1) 为保证没有人受伤的概率大于 0.90，应如何控制参加人数？\par
(2) 若有 100 人参加，求受伤人数不超过 2 人的概率的近似值（用中心极限定理，结果可用 \(\Phi\) 表示）。
\end{exercise}
\begin{solution}
(1) 设人数为 \(n\)，受伤人数 \(X \sim B(n, 0.01)\)。
要求 \(P(X = 0) = 0.99^n > 0.9\)。取对数：
\[
n\ln 0.99 > \ln 0.9 \implies n < \frac{\ln 0.9}{\ln 0.99} \approx 10.48,
\]
故人数控制在 \(\mathbf{10}\) 人及以内。\\
(2) \(n = 100\)，\(X \sim B(100, 0.01)\)，\(EX = 1\)，\(DX = 0.99\)。
\[
P(X \leq 2) = P\!\left(\frac{X-1}{\sqrt{0.99}} \leq \frac{1}{\sqrt{0.99}}\right) \approx \Phi\!\left(\frac{1}{\sqrt{0.99}}\right) = \Phi(\sqrt{100/99}).
\]
\begin{note}
本题 \(np = 1 < 5\)，正态近似精度较差。更精确的方法是泊松近似：\(X \approx \text{Poisson}(1)\)，\(P(X \leq 2) = e^{-1}(1+1+1/2) = 5/(2e) \approx 0.920\)。正态近似仅作为方法展示。
\end{note}
\end{solution}



\begin{exercise}\label{exer:5-normal-route}
从甲地前往乙地有两条路线：第一条穿过市区，用时 \(X_1 \sim N(50, 100)\)（分钟）；第二条走环城公路，用时 \(X_2 \sim N(60, 16)\)（分钟）。若只有 65 分钟可用，应走哪条路线？
\end{exercise}
\begin{solution}
\[
P(X_1 \leq 65) = \Phi\!\left(\frac{65-50}{10}\right) = \Phi(1.5) = 0.9332,
\]
\[
P(X_2 \leq 65) = \Phi\!\left(\frac{65-60}{4}\right) = \Phi(1.25) = 0.8944.
\]
\(0.9332 > 0.8944\)，故走第一条路线较好。
\begin{note}
本题本质是比较两个正态分布在同一阈值下的累积概率，不涉及中心极限定理。但作为正态分布在实际决策中的应用，与CLT应用题有很好的互补性。
\end{note}
\end{solution}


\begin{exercise}\label{exer:5-clt-blood}
某地区某种疾病的患病率为 0.1，现对 10000 人进行血检普查（各人独立）。\par
(1) 用中心极限定理估计患病人数在 \([970, 1030]\) 内的概率（用 \(\Phi\) 表示）；\par
(2) 有两套化验方案：方案一逐一化验（10000 次）；方案二将每 \(k\) 人分为一组，先混合化验，阴性则只需 1 次，阳性则再逐一化验另半血样（共 \(k+1\) 次）。问 \(k\) 满足什么条件时方案二更优？\(k = 4\) 时方案二的平均化验次数是多少？
\end{exercise}
\begin{solution}
(1) 患病人数 \(X \sim B(10000, 0.1)\)，\(EX = 1000\)，\(DX = 900\)，\(\sigma = 30\)。
\[
P(970 \leq X \leq 1030) \approx \Phi\!\left(\frac{30}{30}\right) - \Phi\!\left(\frac{-30}{30}\right) = \Phi(1) - \Phi(-1) = 2\Phi(1) - 1.
\]
\begin{note}
严格使用连续性修正时为 \(P(969.5 < X < 1030.5)\)，因 \(n\) 很大（10000），修正可忽略。
\end{note}
(2) 每组 \(k\) 人全为阴性的概率为 \(0.9^k\)，每组化验次数 \(Y_i\) 满足：
\[
E(Y_i) = 1 \cdot 0.9^k + (k+1)(1-0.9^k) = k + 1 - k \cdot 0.9^k.
\]
总期望 \(E(T_2) = \dfrac{10000}{k}(k+1-k \cdot 0.9^k) = 10000\!\left(1 + \dfrac{1}{k} - 0.9^k\right)\)。\\
方案二更优的条件为 \(E(T_2) < 10000\)，即 \(k \cdot 0.9^k > 1\)。
当 \(k = 4\) 时：
\[
E(Y_i) = 5 - 4 \times 0.9^4 = 5 - 4 \times 0.6561 = 2.3756,
\]
\[
E(T_2) = \frac{10000}{4} \times 2.3756 = 5939.
\]
\end{solution}

\begin{note}
\textbf{本章最容易混淆的几个结论，可以按“问题--条件--结论”一起看：}
\begin{enumerate}
    \item \textbf{马尔可夫不等式}：问题是“非负随机变量取大值的概率能不能估个上界”；条件是 \(X\ge0\) 且 \(EX\) 存在；结论是 \(P\{X\ge a\}\le EX/a\)。
    \item \textbf{切比雪夫不等式}：问题是“偏离均值的概率能不能估个界”；条件是 \(EX\)、\(DX\) 存在；结论是给出上界或下界，但通常不精确。
    \item \textbf{切比雪夫大数定律}：问题是“平均值会不会稳定”；条件是相互独立且方差一致有界；结论是样本平均依概率收敛到平均期望。
    \item \textbf{伯努利大数定律}：问题是“频率会不会稳定到概率”；条件是重复进行同一伯努利试验；结论是 \(\mu_n/n \stackrel{P}{\longrightarrow} p\)。
    \item \textbf{辛钦大数定律}：问题仍是“样本均值会不会稳定”；条件是独立同分布且期望存在；结论是 \(\bar X_n \stackrel{P}{\longrightarrow} \mu\)。
    \item \textbf{列维--林德伯格中心极限定理}：问题是“和或均值的分布能否近似计算”；条件是独立同分布且均值方差存在；结论是标准化后的和或均值近似服从标准正态分布。
    \item \textbf{棣莫弗--拉普拉斯定理}：问题是“二项分布的成功次数如何近似”；条件是 \(Y_n \sim B(n,p)\) 且 \(n\) 较大；结论是 \(Y_n\) 近似服从 \(N(np,\ np(1-p))\)。
\end{enumerate}

如果两条结论看起来都能用，优先选那个\textbf{直接回答题目目标}的：问极限先想大数定律，问近似概率先想中心极限定理。
\end{note}

\begin{note}
\textbf{自学检查点——大数定律与中心极限定理：}
\begin{enumerate}
    \item 马尔可夫不等式的条件是什么？（答：\(X\ge 0\) 且 \(EX\) 存在）
    \item 切比雪夫不等式和马尔可夫不等式有什么关系？（答：切比雪夫可看成把马尔可夫用于 \((X-EX)^2\)）
    \item 大数定律和中心极限定理分别回答什么问题？（答：大数定律回答"均值收敛到哪里"；CLT 回答"和/均值的分布近似什么形状"）
    \item 辛钦大数定律的条件是什么？（答：iid + 期望存在）
    \item CLT 的标准化公式：$\sum_{i=1}^n X_i$ 近似服从什么分布？（答：$N(n\mu, n\sigma^2)$）
    \item 二项分布 $B(n,p)$ 用 CLT 近似时，标准化分母是什么？（答：$\sqrt{np(1-p)}$）
    \item "$n$ 足够大"在实际中一般取多少？（答：通常 $n\geqslant 30$；对二项分布要求 $np\geqslant 5$ 且 $n(1-p)\geqslant 5$）
\end{enumerate}
\end{note}


\chapter{样本与抽样分布}

\begin{note}
\textbf{自学导读：从概率论到数理统计的转折点。}

\textbf{思维方式的根本转变：} 前 5 章的逻辑是"已知分布 $\to$ 求概率/期望"（从总体到样本）。从本章开始，逻辑\textbf{反过来}："已知样本数据 $\to$ 推断总体分布的参数"（从样本到总体）。这就是\textbf{统计推断}。

\textbf{为什么需要抽样分布？} 你从总体中抽了 $n$ 个样本 $X_1,\dots,X_n$，算出样本均值 $\bar{X}$。但 $\bar{X}$ 本身也是随机变量——每次抽样得到的 $\bar{X}$ 都不同。要判断 $\bar{X}$ 离真实均值 $\mu$ 有多远，你必须知道 $\bar{X}$ 服从什么分布。这就是"抽样分布"要解决的问题。

\textbf{本章三大分布的直觉：}
\begin{itemize}
    \item $\chi^2$ 分布：$n$ 个标准正态变量的平方和——衡量"总偏差有多大"；
    \item $t$ 分布：标准正态除以 $\chi^2/n$ 的根号——当方差未知时代替正态分布；
    \item $F$ 分布：两个 $\chi^2$ 的比——比较两组数据的波动大小。
\end{itemize}

\textbf{前置知识：} 本章需要第 3 章的独立性、第 4 章的期望/方差性质、第 5 章的中心极限定理。如果你对"独立同分布"和"$E(\bar{X})=\mu$, $D(\bar{X})=\sigma^2/n$"还不清楚，请先回顾。
\end{note}

{\small
\[
\begin{cases}
\text{基本概念}
\begin{cases}
\text{总体与个体} \\
\text{样本}
\begin{cases}
\text{代表性：与总体具有相同概率分布} \\
\text{独立性：个体之间相互独立}
\end{cases}
\end{cases} \\[6pt]
\text{统计量}
\begin{cases}
\text{概念：不含其它未知参数的样本函数} \\[6pt]
\text{常见统计量：}
\begin{cases}
\text{样本均值：}\bar{X}=\displaystyle\frac{1}{n}\sum_{i=1}^n X_i \\[6pt]
\text{样本方差：}S^2=\displaystyle\frac{1}{n-1}\sum_{i=1}^n (X_i-\bar{X})^2=\frac{1}{n-1}\left(\sum_{i=1}^n X_i^2-n\bar{X}^2\right) \\[6pt]
\text{样本}k\text{阶原点矩：}A_k=\displaystyle\frac{1}{n}\sum_{i=1}^n X_i^k \\[6pt]
\text{样本}k\text{阶中心距：}B_k=\displaystyle\frac{1}{n}\sum_{i=1}^n (X_i-\bar{X})^k,k=2,3\cdots
\end{cases} \\[6pt]
\text{常用统计量的数字特征}
\begin{cases}
E\bar{X}=EX,D\bar{X}=\displaystyle\frac{DX}{n} \\[6pt]
E(S^2)=DX
\end{cases}
\end{cases} \\[6pt]
\text{抽样分布}
\begin{cases}
\chi^2\text{分布}
\begin{cases}
\text{形式：}n\text{个独立标准正态分布的平方和，自由度为}n \\[6pt]
\text{性质：(1)独立可加性；(2)若}X\sim\chi^2(n)\text{,则}EX=n,DX=2n
\end{cases} \\[6pt]
t\text{分布}
\begin{cases}
\text{形式：}Z=\displaystyle\frac{X}{\sqrt{Y/n}}\text{,其中}X,Y\text{独立，} \\[6pt]
X\sim N(0,1),Y\sim\chi^2(n) \\[6pt]
\text{性质：若}Z\sim t(n)\text{,则}EZ=0,DZ=\displaystyle\frac{n}{n-2}
\end{cases} \\[6pt]
F\text{分布}
\begin{cases}
\text{形式：}F=\displaystyle\frac{X/n_1}{Y/n_2}\sim F(n_1,n_2)\text{,其中}X,Y\text{独立，} \\[6pt]
X\sim\chi^2(n_1),Y\sim\chi^2(n_2) \\[6pt]
\text{性质：若}X\sim t(n)\text{,则}X^2\sim F(1,n)
\end{cases} \\[6pt]
\text{单正态总体抽样}
\begin{cases}
\displaystyle\frac{\bar{X}-\mu}{\sigma/\sqrt{n}}\sim N(0,1),\\
\frac{(n-1)S^2}{\sigma^2}\sim\chi^2(n-1),\quad
\bar{X}\text{与}S^2\text{独立}\\
\displaystyle\sum_{i=1}^n\left(\frac{X_i-\mu}{\sigma}\right)^2\sim\chi^2(n),\quad
\frac{\bar{X}-\mu}{S/\sqrt{n}}\sim t(n-1)
\end{cases} \\[6pt]
\text{双正态总体抽样}\\
X\sim N(\mu_1,\sigma_1^2),Y\sim N(\mu_2,\sigma_2^2)
\end{cases}
\end{cases}
\]
}

\begin{introduction}
    \item 总体、个体与样本的概念~\ref{def:population}
    \item 统计量与常用统计量~\ref{def:statistic}
    \item 三大抽样分布~\ref{def:chi2}, \ref{def:t-dist}, \ref{def:f-dist}
    \item 正态总体下的抽样分布定理~\ref{thm:normal-single}, \ref{thm:normal-two}
    \item 分布的分位数~\ref{def:quantile}
\end{introduction}

\begin{note}
\textbf{使用提醒：} 本章主线可以概括为“先有样本，再造统计量，再研究统计量自己的分布”。后面区间估计、假设检验里的 \(z\)、\(t\)、\(\chi^2\)、\(F\) 统计量，之所以能放心使用，前提都是在本章先把它们的抽样分布认清。做题若还停留在“原总体服从什么分布”，却没有追问“我手里的统计量服从什么分布”，往往就会在第 6 章之后持续卡住。
\end{note}

\section{基本概念}

\subsection{总体与样本}

\begin{definition}[总体与个体] \label{def:population}
研究对象的全体称为\textbf{总体}（population），组成总体的每一个元素称为\textbf{个体}（individual）。在对总体进行统计研究时，通常将总体与某个随机变量 \(X\) 等同起来，称“总体 \(X\)”，总体的分布即指随机变量 \(X\) 的分布。
\end{definition}

\begin{note}
直观地说，总体并不是“已经抽到的那几组数据”，而是\textbf{所有可能观测结果的来源}。以后提到“总体均值”“总体方差”时，说的是这个来源本身的参数，而不是样本算出来的数值。最容易混淆的一点是：\textbf{总体参数属于总体，统计量属于样本}，两者角色不同。
\end{note}

\begin{definition}[简单随机样本] \label{def:sample}
设 \(X\) 是具有分布函数 \(F\) 的随机变量，若 \(X_1, X_2, \dots, X_n\) 是与 \(X\) 具有同一分布函数 \(F\) 且相互独立的随机变量，则称 \(X_1, X_2, \dots, X_n\) 为从总体 \(X\)（或分布函数 \(F\)）得到的\textbf{容量为 \(n\) 的简单随机样本}，简称\textbf{样本}。它们的观察值 \(x_1, x_2, \dots, x_n\) 称为\textbf{样本值}。
\end{definition}

\begin{note}
简单随机样本有两个缺一不可的条件：一是\textbf{同分布}，二是\textbf{相互独立}。少了前者，就不是来自同一个总体；少了后者，后续很多结论（例如样本联合分布可分解、样本均值方差公式、三大抽样分布）都不能直接使用。另一个常见混淆是：“样本”指随机变量 \(X_1,\dots,X_n\)，“样本值”指一次观测得到的具体数字 \(x_1,\dots,x_n\)。
\end{note}

\begin{property}[样本的分布] \label{prop:sample-dist}
设总体 \(X\) 的分布函数为 \(F(x)\)，则样本 \(X_1, X_2, \dots, X_n\) 的联合分布函数为
\[
F(x_1, x_2, \dots, x_n) = \prod_{i=1}^{n} F(x_i).
\]
对于离散型总体，联合分布为 \(P\{X_1 = x_1, \dots, X_n = x_n\} = \prod_{i=1}^{n} P\{X_i = x_i\}\)；对于连续型总体，联合密度为 \(f(x_1, \dots, x_n) = \prod_{i=1}^{n} f(x_i)\)。
\end{property}


\subsection{统计量}

\begin{definition}[统计量] \label{def:statistic}
设 \(X_1, X_2, \dots, X_n\) 为来自总体 \(X\) 的一个样本，\(g(X_1, X_2, \dots, X_n)\) 是 \(X_1, X_2, \dots, X_n\) 的函数。若 \(g\) 中\textbf{不含任何未知参数}，则称 \(g(X_1, X_2, \dots, X_n)\) 为一个\textbf{统计量}。也就是说，统计量只能由样本和已知常数构成，不能把待估参数偷偷写进去。
\end{definition}

\begin{note}
统计量是样本的函数，也是随机变量。判断一个表达式是否为统计量时，最可靠的判据只有一个：\textbf{看它是否含有未知参数}。已知参数可以出现，未知参数不可以出现；至于表达式长短、是否有求和、是否有最大值最小值，都不是本质标准。
\end{note}

\begin{note}
\textbf{使用提醒：} \(\overline{X}\)、\(S^2\)、\(\max X_i\)、\(\min X_i\) 都可以是统计量；关键不在形式，而在是否含未知参数。比如 \(\sum (X_i-\mu)^2\) 在 \(\mu\) 已知时是统计量，但 \(\sum (X_i-\sigma)^2\) 若 \(\sigma\) 未知，就不是统计量。
\end{note}

\begin{note}
\textbf{易错点：} 统计量和估计量不是同义词。估计量一定是统计量，因为它必须是不含未知参数的样本函数；但统计量未必专门拿来估参数，很多统计量只是后续构造抽样分布、枢轴量或检验统计量的中间工具。把这两个概念分开，后面学点估计与假设检验时会清楚很多。
\end{note}

\begin{exercise}[统计量的判定] \label{ex:statistic_01}
设总体 \(X \sim N(\mu,\sigma^2)\)，其中 \(\mu\) 已知，\(\sigma^2\) 未知，\(X_1,X_2,\dots,X_n\) 是简单随机样本，则下列样本函数中\textbf{不是}统计量的是（\quad）。\par
A. \(\dfrac{1}{n}\displaystyle\sum_{i=1}^{n}X_i\)\qquad
B. \(\max\limits_{1\leq i\leq n}\{X_i\}\)\par
C. \(\displaystyle\sum_{i=1}^{n}\!\left(\dfrac{X_i-\mu}{\sigma}\right)^{\!2}\)\qquad
D. \(\dfrac{1}{n}\displaystyle\sum_{i=1}^{n}(X_i-\mu)^2\)
\end{exercise}
\begin{solution}
应选 C。

判断这类题时，不必先算值，只要逐项检查“是否仅由样本和已知常数构成”即可。

A 中 \(\dfrac{1}{n}\sum_{i=1}^{n}X_i=\overline{X}\) 只由样本构成，是统计量。

B 中 \(\max\limits_{1\le i\le n}\{X_i\}\) 虽然形式特殊，但本质上仍只是样本的函数，也是统计量。

C 中
\[
\sum_{i=1}^{n}\left(\frac{X_i-\mu}{\sigma}\right)^2
\]
含有未知参数 \(\sigma\)。题目虽然说明 \(\mu\) 已知，但 \(\sigma^2\) 未知，所以这里的 \(\sigma\) 不能出现在统计量里，因此它\textbf{不是}统计量。

D 中 \(\dfrac{1}{n}\sum_{i=1}^{n}(X_i-\mu)^2\) 只含样本和已知参数 \(\mu\)，仍是统计量。

所以只有 C 不满足统计量定义。
\end{solution}


\begin{exercise}\label{exer:6-statistic-judge}
设总体 \(X \sim N(\mu, \sigma^2)\)，其中 \(\mu\) 未知，\(\sigma^2\) 已知，\(X_1, X_2, X_3\) 是简单随机样本，则下列表达式中不是统计量的是（\quad）。\par
A. \(X_1 + X_2 + X_3\)\qquad B. \(\min(X_1, X_2, X_3)\)\qquad C. \(\displaystyle\sum_{i=1}^{3}\frac{X_i^2}{\sigma^2}\)\qquad D. \(\overline{X} + 2\mu\)
\end{exercise}

\begin{solution}
应选 D。

这类题的判别标准只有一句：\textbf{统计量必须是不含未知参数的样本函数}。题目中 \(\sigma^2\) 已知、\(\mu\) 未知，所以凡是出现 \(\mu\) 的表达式都不能直接算作统计量，而出现已知的 \(\sigma^2\) 不构成障碍。

\textbf{逐项判断：}
A 中 \(X_1+X_2+X_3\) 只由样本构成，是统计量。

B 中 \(\min(X_1,X_2,X_3)\) 也是样本的函数，仍是统计量。

C 中 \(\sum\limits_{i=1}^{3}\dfrac{X_i^2}{\sigma^2}\) 虽然出现了参数 \(\sigma^2\)，但它是\textbf{已知参数}，所以不影响统计量判定，仍是统计量。

D 中 \(\overline{X}+2\mu\) 含有未知参数 \(\mu\)，因此不是统计量。

因此只有 D 不满足统计量定义，故应选 D。
\end{solution}

\begin{exercise}\label{exer:6-statistic-sigma}
设 \((X_1, X_2, X_3)\) 为取自总体 \(N(0, \sigma^2)\) 的样本，\(\sigma^2\) 为未知参数。\(T_1 = X_1+X_2+X_3\)，\(T_2 = \dfrac{X_3 - X_1}{2}\)，\(T_3 = \dfrac{1}{\sigma}\sum_{i=1}^{3}X_i\)，\(T_4 = \max\{X_1, X_2, X_3\}\)，其中不是统计量的是（\quad）。\par
A. \(T_1\)\qquad B. \(T_2\)\qquad C. \(T_3\)\qquad D. \(T_4\)
\end{exercise}
\begin{solution}
应选 C。

这道题与上一题完全同型，关键仍是检查\textbf{是否含未知参数}。题目中总体为 \(N(0,\sigma^2)\)，未知的是 \(\sigma\)，因此只要表达式里还显式保留着 \(\sigma\)，就不是统计量。

\textbf{逐项判断：}
\[
T_1=X_1+X_2+X_3,\qquad T_2=\frac{X_3-X_1}{2},\qquad T_4=\max\{X_1,X_2,X_3\}
\]
都只由样本值构成，所以都是统计量。

\(T_3=\dfrac{1}{\sigma}\sum_{i=1}^{3}X_i\) 中出现了未知参数 \(\sigma\)，故不是统计量。

因此答案选 C。
\end{solution}

\begin{exercise}\label{exer:6-stat-mu}
设 \((X_1, X_2, X_3)\) 为取自总体的样本，\(\sigma^2\) 已知，\(\mu\) 未知，\(T_1=X_1+X_2+X_3\)，\(T_2=(X_3-\mu)/2\)，\(T_3=\sum_{i=1}^{3}X_i\)，\(T_4=\max\{X_1,X_2,X_3\}\)，其中不是统计量的是（\quad）。\par
A. \(T_1\)\quad B. \(T_2\)\quad C. \(T_3\)\quad D. \(T_4\)
\end{exercise}
\begin{solution}
应选 B。

本题里 \(\sigma^2\) 已知、\(\mu\) 未知，所以判别关键就是“有没有出现未知参数 \(\mu\) ”。不要因为出现减法或最大值就误以为不再是统计量，真正决定性的仍然只是参数是否未知。

\textbf{逐项判断：}
A 中 \(T_1=X_1+X_2+X_3\) 只含样本，是统计量。

B 中 \(T_2=\dfrac{X_3-\mu}{2}\) 含有未知参数 \(\mu\)，所以不是统计量。

C 中 \(T_3=\sum_{i=1}^{3}X_i\) 只含样本，是统计量。

D 中 \(T_4=\max\{X_1,X_2,X_3\}\) 虽然用了最大值符号，但仍只是样本函数，也是统计量。

因此只有 B 不是统计量，应选 B。
\end{solution}


\begin{definition}[常用统计量] \label{def:common-stat}
设 \(X_1, X_2, \dots, X_n\) 是来自总体 \(X\) 的样本，\(\overline{X}\) 和 \(S^2\) 分别为样本均值与样本方差。
\begin{enumerate}
    \item \textbf{样本均值}：\(\displaystyle\overline{X} = \frac{1}{n}\sum_{i=1}^{n}X_i\)；
    \item \textbf{样本方差}：\(\displaystyle S^2 = \frac{1}{n-1}\sum_{i=1}^{n}(X_i - \overline{X})^2 = \frac{1}{n-1}\!\left(\sum_{i=1}^{n}X_i^2 - n\overline{X}^2\right)\)；
    \item \textbf{样本标准差}：\(S = \sqrt{S^2}\)；
    \item \textbf{样本 \(k\) 阶原点矩}：\(\displaystyle A_k = \frac{1}{n}\sum_{i=1}^{n}X_i^k\ (k=1,2,\dots)\)；
    \item \textbf{样本 \(k\) 阶中心矩}：\(\displaystyle B_k = \frac{1}{n}\sum_{i=1}^{n}(X_i - \overline{X})^k\ (k=2,3,\dots)\)。
\end{enumerate}
\end{definition}

\begin{note}
这几个统计量里最容易混淆的是样本方差 \(S^2\) 的分母。这里采用
\[
S^2=\frac{1}{n-1}\sum_{i=1}^{n}(X_i-\overline{X})^2,
\]
目的是让它成为总体方差的无偏估计。若把分母误写成 \(n\)，虽然形式更像”平均”，但后续很多无偏性和抽样分布结论都要跟着改变，不能混用。
\end{note}

\begin{note}
\textbf{直觉理解”为什么除以 $n-1$”：} $n$ 个数据 $X_1,\dots,X_n$ 本来有 $n$ 个”自由度”（可以独立变化的方向）。但一旦你用 $\bar{X}$ 代替了真实均值 $\mu$，就多加了一个约束 $\sum(X_i-\bar{X})=0$，相当于”用掉了一个自由度”。剩下只有 $n-1$ 个独立的偏差，所以除以 $n-1$ 才是正确的”平均”。

这也解释了为什么 $\frac{(n-1)S^2}{\sigma^2}\sim\chi^2(n-1)$ 而不是 $\chi^2(n)$——自由度正好反映了被”用掉”的那一个。
\end{note}

\begin{property}[统计量的数字特征] \label{prop:stat-moment}
设总体 \(X\) 的期望为 \(EX = \mu\)，方差为 \(DX = \sigma^2\)，\(X_1, X_2, \dots, X_n\) 为容量为 \(n\) 的样本，则
\[
EX_i = \mu,\quad DX_i = \sigma^2,\quad E\overline{X} = \mu,\quad D\overline{X} = \frac{\sigma^2}{n},\quad E(S^2) = \sigma^2.
\]
\end{property}
\begin{proof}
\(E\overline{X} = \mu\) 和 \(D\overline{X} = \sigma^2/n\) 由期望和方差的性质直接可得。关键在于证明 \(E(S^2) = \sigma^2\)。令 \(Y_i = X_i - \mu\)，则
\[
\sum_{i=1}^{n}(X_i - \overline{X})^2 = \sum_{i=1}^{n}(Y_i - \overline{Y})^2 = \sum_{i=1}^{n} Y_i^2 - n\overline{Y}^2.
\]
取期望：
\[
E\!\left[\sum(Y_i - \overline{Y})^2\right] = n\sigma^2 - n\cdot\frac{\sigma^2}{n} = (n-1)\sigma^2.
\]
因此
\[
E(S^2) = E\!\left[\frac{1}{n-1}\sum_{i=1}^{n}(X_i-\overline{X})^2\right] = \frac{(n-1)\sigma^2}{n-1} = \sigma^2.
\]
\end{proof}

\begin{note}
由大数定律，样本均值和样本方差依概率收敛于总体均值和总体方差：
\(\overline{X} \stackrel{P}{\longrightarrow} \mu\)，\(S^2 \stackrel{P}{\longrightarrow} \sigma^2\)，\(A_k \stackrel{P}{\longrightarrow} EX^k\)。
\end{note}

\begin{exercise}\label{exer:6-yi-var}
设 \(X_1, X_2, \dots, X_n\ (n > 2)\) 为独立同分布的随机变量，均服从 \(N(0,1)\)，记 \(\overline{X} = \frac{1}{n}\sum_{i=1}^{n}X_i\)，\(Y_i = X_i - \overline{X}\ (i=1,2,\dots,n)\)。则 \(DY_i =\,\underline{\hspace{2cm}}\)。
\end{exercise}

\begin{solution}
应填 \(\dfrac{n-1}{n}\)。

先把 \(Y_i\) 写成样本的线性组合：
\[
Y_i=X_i-\overline{X}
=X_i-\frac{1}{n}\sum_{k=1}^{n}X_k
=\left(1-\frac{1}{n}\right)X_i-\frac{1}{n}\sum_{k\ne i}X_k.
\]

由于 \(X_1,\dots,X_n\) 相互独立，且都服从 \(N(0,1)\)，所以每个 \(X_k\) 的方差都是 1，不同项之间协方差为 0。于是
\[
DY_i
=\left(1-\frac{1}{n}\right)^2 DX_i
\frac{1}{n^2}\sum_{k\ne i}DX_k.
\]
代入 \(DX_k=1\) 以及“除 \(i\) 外共有 \(n-1\) 项”，得
\[
DY_i=\left(1-\frac{1}{n}\right)^2+\frac{n-1}{n^2}
=1-\frac{2}{n}+\frac{1}{n^2}+\frac{n-1}{n^2}
=1-\frac{1}{n}
=\frac{n-1}{n}.
\]

所以 \(DY_i=\dfrac{n-1}{n}\)。
\end{solution}

\begin{exercise}\label{exer:6-sk-sq}
已知总体 \(X\) 的期望 \(EX = 0\)，方差 \(DX = \sigma^2\)，从总体 \(X\) 中抽取容量为 \(n\) 的简单随机样本，其样本均值和样本方差分别为 \(\overline{X}, S^2\)。记 \(S_k^2 = k\overline{X}^2 + kS^2\ (k=1,2,3,4)\)，则（\quad）。\par
A. \(E(S_1^2) = \sigma^2\)\quad B. \(E(S_2^2) = \sigma^2\)\quad C. \(E(S_3^2) = \sigma^2\)\quad D. \(E(S_4^2) = \sigma^2\)
\end{exercise}

\begin{solution}
四个选项均不能精确成立。

\(E(\overline{X}^2) = D\overline{X} + (E\overline{X})^2 = \dfrac{\sigma^2}{n} + 0 = \dfrac{\sigma^2}{n}\)，\(E(S^2) = \sigma^2\)。
\[
E(S_k^2) = k\,E(\overline{X}^2) + k\,E(S^2) = k \cdot \frac{\sigma^2}{n} + k\sigma^2 = k\sigma^2\!\left(1+\frac{1}{n}\right) = k\sigma^2 \cdot \frac{n+1}{n}.
\]
令 \(E(S_k^2) = \sigma^2\)，得 \(k(n+1) = n\)，即 \(k = \dfrac{n}{n+1}\)。对于任意正整数 \(n\)：
\begin{itemize}
    \item \(k=1\)：需 \(n\to\infty\)，无有限解；
    \item \(k=2\)：需 \(n=-2\)，不可能；
    \item \(k=3\)：需 \(n=-3/2\)，不可能；
    \item \(k=4\)：需 \(n=-4/3\)，不可能。
\end{itemize}
故对任意正整数 \(n\)，四个选项均不精确成立。若本题出自某套考研真题，建议核实原题条件。
\begin{note}
此题需注意 \(E(\overline{X}^2) = D\overline{X} = \sigma^2/n\)（因为 \(E\overline{X}=0\)），而非 \((EX)^2=0\)。
\end{note}
\end{solution}


\section{三大抽样分布}

\begin{note}
\textbf{自学导读：为什么要"发明"$\chi^2$、$t$、$F$ 这三个分布？}

你可能会问：正态分布不够用吗？答案是：当总体方差 $\sigma^2$ \textbf{已知}时，$\frac{\bar{X}-\mu}{\sigma/\sqrt{n}}\sim N(0,1)$，正态确实够用。但实际中 $\sigma^2$ 几乎总是未知的，你只能用样本方差 $S^2$ 代替。一旦代替，分母变成了随机的 $S/\sqrt{n}$，整个统计量就不再服从正态分布了——它服从 $t$ 分布。

类似地：
\begin{itemize}
    \item 想检验方差是否等于某个值？需要 $(n-1)S^2/\sigma^2$ 的分布——这就是 $\chi^2$ 分布；
    \item 想比较两组数据的方差谁大？需要两个 $\chi^2$ 的比——这就是 $F$ 分布。
\end{itemize}

所以这三个分布不是凭空造出来的，而是"正态总体 + 方差未知"这个现实约束的必然产物。后面第 7--9 章的所有公式都建立在它们之上。
\end{note}

\subsection{\texorpdfstring{\(\chi^2\)}{χ²} 分布}

\begin{definition}[\texorpdfstring{\(\chi^2\)}{χ²} 分布] \label{def:chi2}
设 \(X_1, X_2, \dots, X_n\) 相互独立且都服从标准正态分布 \(N(0,1)\)，令
\[
\chi^2 = X_1^2 + X_2^2 + \cdots + X_n^2,
\]
则称 \(\chi^2\) 服从自由度为 \(n\) 的 \(\chi^2\) 分布，记为 \(\chi^2 \sim \chi^2(n)\)。
\end{definition}

\begin{note}
\textbf{自学追问：} 为什么是“平方和”？因为标准正态变量既有正负，又以 0 为中心，平方以后就只保留“偏离 0 的大小”，而把多个独立偏离量加起来，就得到一个专门衡量总体波动大小的分布。这里的自由度 \(n\) 可以先直观理解为“有多少个独立的标准正态平方项参与了相加”。
\end{note}

\begin{property}\label{prop:chi2-property} 
\begin{enumerate}
    \item \(E[\chi^2(n)] = n\)，\(D[\chi^2(n)] = 2n\)。
    \item \textbf{可加性}：若 \(X \sim \chi^2(m)\)，\(Y \sim \chi^2(n)\)，且 \(X, Y\) 相互独立，则 \(X + Y \sim \chi^2(m+n)\)。
\end{enumerate}
\end{property}

\begin{exercise}\label{exer:6-chi2-ez}
设 \((X_1, X_2, \dots, X_{10})\) 与 \((Y_1, Y_2, \dots, Y_{10})\) 分别是来自正态总体 \(N(0,1)\) 的两组独立样本，\(Z = \sum_{i=1}^{10}(X_i-\overline{X})^2 + \sum_{i=1}^{10}(Y_i-\overline{Y})^2\)，则 \(EZ =\,\underline{\hspace{2cm}}\)。
\end{exercise}
\begin{solution}
应填 \(18\)。

先分别看两部分。因为 \(X_1,\dots,X_{10}\) 来自正态总体 \(N(0,1)\)，所以由单正态总体抽样分布，
\[
\sum_{i=1}^{10}(X_i-\overline{X})^2=\frac{(10-1)S_X^2}{1}\sim \chi^2(9).
\]
同理，
\[
\sum_{i=1}^{10}(Y_i-\overline{Y})^2\sim \chi^2(9).
\]

又由于 \((X_1,\dots,X_{10})\) 与 \((Y_1,\dots,Y_{10})\) 是两组独立样本，所以这两个卡方变量也相互独立。利用 \(\chi^2\) 分布的可加性可得
\[
Z=\sum_{i=1}^{10}(X_i-\overline{X})^2+\sum_{i=1}^{10}(Y_i-\overline{Y})^2\sim\chi^2(18).
\]

这里每组样本都因为减去一个样本均值而损失 1 个自由度，所以单组对应 \(\chi^2(9)\)；两组独立样本的自由度相加，就得到 \(\chi^2(18)\)。

最后利用 \(\chi^2(\nu)\) 的期望等于自由度 \(\nu\)，便有
\[
EZ=18.
\]

所以答案是 \(18\)。
\end{solution}


\begin{exercise}\label{exer:6-chi2-basic}
设 \(X_1, X_2, X_3, X_4\) 是来自正态总体 \(N(0, 2^2)\) 的样本，令 \(Y = (X_1+X_2)^2 + (X_3-X_4)^2\)。当 \(C =\,\underline{\hspace{2cm}}\) 时，\(CY \sim \chi^2(2)\)。
\end{exercise}

\begin{solution}
应填 \(\dfrac{1}{8}\)。

\textbf{思路是把 \(Y\) 改写成两个独立标准正态变量的平方和。}

由于 \(X_1,X_2,X_3,X_4\) 相互独立，且每个变量都服从 \(N(0,2^2)\)，所以
\[
X_1+X_2 \sim N(0,8),\qquad X_3-X_4 \sim N(0,8).
\]
又因为前一个线性组合只含 \(X_1,X_2\)，后一个只含 \(X_3,X_4\)，它们由两组独立样本构成，因此彼此独立。于是标准化后有
\[
\frac{X_1+X_2}{\sqrt{8}} \sim N(0,1),\qquad \frac{X_3-X_4}{\sqrt{8}} \sim N(0,1).
\]
\[
\left(\frac{X_1+X_2}{\sqrt{8}}\right)^{\!2} + \left(\frac{X_3-X_4}{\sqrt{8}}\right)^{\!2} \sim \chi^2(2) \implies C = \frac{1}{8}.
\]
\[
\text{即}\quad \frac{Y}{8}\sim \chi^2(2),
\]
故 \(C=\dfrac{1}{8}\)。
\end{solution}

\begin{exercise}\label{exer:6-chi2-ab}
设 \(X_1, X_2, X_3, X_4\) 是来自正态总体 \(N(0, 2^2)\) 的简单随机样本，记
\[
X = a(X_1 - 2X_2)^2 + b(3X_3 - 4X_4)^2,
\]
若统计量 \(X\) 服从自由度为 2 的 \(\chi^2\) 分布，则 \(ab =\,\underline{\hspace{2cm}}\)。
\end{exercise}

\begin{solution}
应填 \(\dfrac{1}{2000}\)。

\textbf{仍然先把两个平方项各自标准化成标准正态平方。}

\(X_1 - 2X_2\) 与 \(3X_3 - 4X_4\) 都是正态变量的线性组合，且
\[
X_1 - 2X_2 \sim N(0,\, 1^2\cdot 4 + (-2)^2\cdot 4) = N(0, 20),
\]
\[
3X_3 - 4X_4 \sim N(0,\, 3^2\cdot 4 + (-4)^2\cdot 4) = N(0, 100).
\]
又因为这两个线性组合分别只含 \((X_1,X_2)\) 与 \((X_3,X_4)\)，而两组样本彼此独立，所以它们相互独立。

标准化后得到两个独立的标准正态变量：
\[
\frac{X_1-2X_2}{\sqrt{20}}\sim N(0,1),\qquad \frac{3X_3-4X_4}{10}\sim N(0,1).
\]
因此由 \(\chi^2\) 分布的定义，
\[
\frac{(X_1-2X_2)^2}{20} + \frac{(3X_3-4X_4)^2}{100} \sim \chi^2(2),
\]
与题设
\[
X = a(X_1 - 2X_2)^2 + b(3X_3 - 4X_4)^2
\]
对照可知 \(a=\frac{1}{20},\ b=\frac{1}{100}\)，所以 \(ab=\frac{1}{2000}\)。
\end{solution}


\subsection{\texorpdfstring{\(t\)}{t} 分布}

\begin{definition}[\texorpdfstring{\(t\)}{t} 分布] \label{def:t-dist}
设 \(X \sim N(0,1)\)，\(Y \sim \chi^2(n)\)，且 \(X, Y\) 相互独立，令
\[
t = \frac{X}{\sqrt{Y/n}},
\]
则称 \(t\) 服从自由度为 \(n\) 的 \(t\) 分布，记为 \(t \sim t(n)\)。
\end{definition}

\begin{note}
\textbf{使用提醒：} \(t\) 分布本质上是“标准化后的均值”再除以“样本标准差”的结果，因此它总在“总体正态、方差未知”的场景里出现。自由度来自分母里的 \(\chi^2(n)\)：分母若由样本方差构造出来，常常会因为估计了一个均值而损失 1 个自由度。
\end{note}

\begin{property}\label{prop:t-dist-property}
\begin{enumerate}
    \item \(t\) 分布的概率密度关于 \(x = 0\) 对称，\(Et = 0\ (n \geq 2)\)，\(Dt = \dfrac{n}{n-2}\ (n > 2)\)。
    \item 当 \(n\) 足够大时，\(t\) 分布近似于标准正态分布；\(n\) 越大，曲线越"瘦"越"高"。
    \item 由对称性：\(t_{1-\alpha}(n) = -t_\alpha(n)\)。
    \item 若 \(t \sim t(n)\)，则 \(t^2 \sim F(1,n)\)。
\end{enumerate}
\end{property}

\begin{note}
\textbf{直觉理解 $t$ 分布为什么比正态分布"胖"：} 分子 $X\sim N(0,1)$ 是确定的标准正态；分母 $\sqrt{Y/n}$ 本身也是随机的。当分母偶然偏小时，整个比值会被"放大"，产生极端值。所以 $t$ 分布的尾巴比正态分布更厚。

当自由度 $n\to\infty$ 时，大数定律保证 $Y/n\to 1$（因为 $EY=n$），分母趋于常数 1，$t$ 分布就退化为标准正态。实际中 $n\geqslant 30$ 时两者已经非常接近，这就是为什么"大样本时 $t$ 检验和 $Z$ 检验结果几乎一样"。
\end{note}

\begin{exercise}[\texorpdfstring{\(t\)}{t} 分布的构造] \label{ex:sampling_04}
设 \(X_1,X_2,X_3,X_4\) 为来自 \(N(0,1/2)\) 的简单随机样本，则统计量
\[
Y=\frac{X_1-X_2}{\sqrt{X_3^2+X_4^2}}
\]
服从（\quad）。\par
A. \(N(0,1)\)\qquad B. \(\chi^2(2)\)\qquad C. \(t(2)\)\qquad D. \(F(2,2)\)
\end{exercise}
\begin{solution}
应选 C。\textbf{分子}：\(X_1-X_2\sim N(0,1/2+1/2)=N(0,1)\)，故 \(X_1-X_2\sim N(0,1)\)。\\
\textbf{分母}：\(\sqrt{2}\,X_i\sim N(0,1)\)，从而
\[
(\sqrt{2}\,X_3)^2+(\sqrt{2}\,X_4)^2 = 2(X_3^2+X_4^2)\sim\chi^2(2).
\]
由 \(t\) 分布的定义，分子（标准正态）除以分母（\(\chi^2(2)/2\) 开方）：
\[
Y = \frac{X_1-X_2}{\sqrt{X_3^2+X_4^2}} = \frac{(X_1-X_2)/1}{\sqrt{2(X_3^2+X_4^2)/2}}\sim t(2).
\]
\end{solution}


\begin{exercise}\label{exer:6-t-new}
设 \(X_1, X_2, \dots, X_n\) 为来自正态分布 \(N(\mu, \sigma^2)\) 的简单随机样本，\(\overline{X}\) 和 \(S^2\) 分别为样本均值和样本方差。又设 \(X_{n+1} \sim N(\mu, \sigma^2)\) 且与 \(X_1, \dots, X_n\) 相互独立，则统计量
\[
\frac{X_{n+1} - \overline{X}}{S}\sqrt{\frac{n}{n+1}}
\]
服从\,\underline{\hspace{2cm}}分布（注明自由度）。
\end{exercise}

\begin{solution}
应填 \(t(n-1)\)。

\(\overline{X} \sim N(\mu, \sigma^2/n)\)，\(X_{n+1} \sim N(\mu, \sigma^2)\)，且相互独立，故
\[
X_{n+1} - \overline{X} \sim N\!\left(0,\, \left(1+\frac{1}{n}\right)\sigma^2\right),\quad \frac{X_{n+1}-\overline{X}}{\sigma\sqrt{1+1/n}} \sim N(0,1).
\]

又 \(\dfrac{(n-1)S^2}{\sigma^2} \sim \chi^2(n-1)\)，且与 \(\overline{X}\) 独立，故
\[
\frac{X_{n+1}-\overline{X}}{S}\sqrt{\frac{n}{n+1}} = \frac{(X_{n+1}-\overline{X})\big/\!\left(\sigma\sqrt{1+\frac{1}{n}}\right)}{\sqrt{(n-1)S^2/[\sigma^2(n-1)]}} \sim t(n-1).
\]
\end{solution}

\begin{exercise}\label{exer:6-t-y-sq}
设随机变量 \(X \sim t(n)\)，\(Y \sim F(1,n)\)，给定 \(\alpha\ (0 < \alpha < 0.5)\)，常数 \(c\) 满足 \(P\{X > c\} = \alpha\)，则 \(P\{Y > c^2\} =\,\underline{\hspace{2cm}}\)。\par
A. \(\alpha\)\quad B. \(1-\alpha\)\quad C. \(2\alpha\)\quad D. \(1-2\alpha\)
\end{exercise}
\begin{solution}
应选 C。

由 \(X \sim t(n)\) 知
\[
X^2 \sim F(1,n).
\]
因此 \(Y\) 与 \(X^2\) 同分布，题目所求概率可改写为
\[
P\{Y > c^2\} = P\{X^2 > c^2\}.
\]
\[
\text{而 } X^2>c^2 \Longleftrightarrow X>c \text{ 或 } X<-c.
\]
又因为 \(t\) 分布关于 0 对称，且题设给出 \(P\{X>c\}=\alpha\)，所以
\[
P\{X<-c\}=\alpha.
\]
故
\[
P\{Y > c^2\} = P\{X > c\} + P\{X < -c\} = \alpha + \alpha = 2\alpha.
\]
\begin{note}
此题只能断言 \(Y\) 与 \(X^2\) 同分布（即 \(P(Y>c^2)=P(X^2>c^2)\)），不能说 \(Y=X^2\)——服从相同分布的随机变量不一定相同。
\end{note}
\end{solution}

\subsection{\texorpdfstring{\(F\)}{F} 分布}

\begin{definition}[\texorpdfstring{\(F\)}{F} 分布] \label{def:f-dist}
设 \(X \sim \chi^2(n_1)\)，\(Y \sim \chi^2(n_2)\)，且 \(X, Y\) 相互独立，令
\[
F = \frac{X/n_1}{Y/n_2},
\]
则称 \(F\) 服从自由度为 \((n_1, n_2)\) 的 \(F\) 分布，记为 \(F \sim F(n_1, n_2)\)，其中 \(n_1\) 为第一自由度，\(n_2\) 为第二自由度。
\end{definition}

\begin{note}
\textbf{自学追问：} \(F\) 分布为什么总像“比值”？因为它比较的是两个独立波动量的相对大小：分子是一份标准化后的方差型量，分母也是一份标准化后的方差型量。于是 \(F\) 分布天然适合回答“两个方差是否一样大”“这个统计量相对另一个放大了多少”这一类问题。
\end{note}

\begin{property}\label{prop:f-dist-property}
\begin{enumerate}
    \item 若 \(F \sim F(n_1, n_2)\)，则 \(\dfrac{1}{F} \sim F(n_2, n_1)\)。
    \item \(F_{1-\alpha}(n_1, n_2) = \dfrac{1}{F_\alpha(n_2, n_1)}\)，常用于求 \(F\) 分布表中未列出的分位数。
    \item 若 \(t \sim t(n)\)，则 \(t^2 \sim F(1, n)\)。
\end{enumerate}
\end{property}

\begin{exercise}\label{exer:6-f-const}
设 \(X_1, X_2, X_3, X_4\) 相互独立，均服从 \(N(0,1)\)。\(Z = c \times \dfrac{X_1^2/1}{X_2^2+X_3^2+X_4^2}\) 服从 \(F\) 分布，则常数 \(c =\,\underline{\hspace{2cm}}\)。\par
A. \(3\)\quad B. \(\dfrac{1}{3}\)\quad C. \(\dfrac{2}{3}\)\quad D. \(1\)
\end{exercise}
\begin{solution}
应选 A。

因为 \(X_1,X_2,X_3,X_4\) 相互独立，且都服从 \(N(0,1)\)，所以
\[
X_1^2 \sim \chi^2(1),
\qquad
X_2^2+X_3^2+X_4^2 \sim \chi^2(3),
\]
并且这两个卡方变量相互独立。

\(F\) 分布的标准形式是
\[
\frac{\chi^2(n_1)/n_1}{\chi^2(n_2)/n_2}\sim F(n_1,n_2).
\]
因此本题要写成
\[
\frac{X_1^2/1}{(X_2^2+X_3^2+X_4^2)/3} \sim F(1,3) \implies c = 3.
\]

所以常数应取
\[
c=3.
\]
\end{solution}

\begin{exercise}\label{exer:6-t2-f}
设随机变量 \(T \sim t(15)\)，则 \(T^2\) 服从的分布为\,\underline{\hspace{2cm}}。
\end{exercise}
\begin{solution}
应填 \(F(1,15)\)。

先回忆标准结论：若
\[
T = \frac{U}{\sqrt{V/n}},
\]
其中 \(U \sim N(0,1)\)，\(V \sim \chi^2(n)\) 且相互独立，则 \(T \sim t(n)\)。

于是平方后
\[
T^2 = \frac{U^2/1}{V/n} \sim F(1,n),
\]
因为 \(U^2 \sim \chi^2(1)\)，所以这正好是 \(F\) 分布的标准形式。

本题给定 \(T \sim t(15)\)，即 \(n=15\)，故
\[
T^2 \sim F(1,15).
\]
\end{solution}

\begin{exercise}\label{exer:6-inv-t2}
设随机变量 \(X \sim t(n)\ (n > 1)\)，则 \(Y = \dfrac{1}{X^2}\) 服从的分布为（\quad）。\par
A. \(\chi^2(n)\)\quad B. \(\chi^2(n-1)\)\quad C. \(F(n,1)\)\quad D. \(F(1,n)\)
\end{exercise}
\begin{solution}
应选 C。

先把 \(X\sim t(n)\) 写成标准形式：
\[
X = \frac{U}{\sqrt{V/n}},
\]
其中 \(U \sim N(0,1)\)，\(V \sim \chi^2(n)\) 且相互独立。

于是
\[
X^2 = \frac{U^2/1}{V/n} \sim F(1,n).
\]
题目要求的是 \(Y=\dfrac{1}{X^2}\)，也就是把一个 \(F(1,n)\) 变量取倒数。根据 \(F\) 分布的性质“若 \(F\sim F(n_1,n_2)\)，则 \(1/F\sim F(n_2,n_1)\)”，可得
\[
Y = \frac{1}{X^2} = \frac{V/n}{U^2/1} \sim F(n, 1).
\]
\[
\text{故应选 } C.
\]
\end{solution}

\begin{exercise}[\texorpdfstring{\(F\)}{F} 分布系数的确定] \label{ex:sampling_03}
设总体 \(X\sim N(\mu_1,4),\ Y\sim N(\mu_2,5)\)，\(X\) 与 \(Y\) 独立，\(X_1,\dots,X_8\) 和 \(Y_1,\dots,Y_{10}\) 分别来自 \(X\) 和 \(Y\)，\(S_X^2,S_Y^2\) 为样本方差，则（\quad）。\par
A. \(\dfrac{2S_X^2}{5S_Y^2}\sim F(7,9)\)\qquad
B. \(\dfrac{5S_X^2}{2S_Y^2}\sim F(7,9)\)\par
C. \(\dfrac{4S_X^2}{5S_Y^2}\sim F(7,9)\)\qquad
D. \(\dfrac{5S_X^2}{4S_Y^2}\sim F(7,9)\)
\end{exercise}

\begin{solution}
应选 D。由正态总体抽样分布：
\[
\frac{7S_X^2}{\sigma_1^2}=\frac{7S_X^2}{4}\sim\chi^2(7),\qquad \frac{9S_Y^2}{\sigma_2^2}=\frac{9S_Y^2}{5}\sim\chi^2(9).
\]
构造 \(F\) 统计量时，系数需使分子分母分别除以各自自由度：
\[
\frac{7S_X^2/4}{9S_Y^2/5}\cdot\frac{1}{7/9}=\frac{9\cdot 7S_X^2}{7\cdot 9S_Y^2}\cdot\frac{5}{4}=\frac{5S_X^2}{4S_Y^2}\sim F(7,9).
\]
\textbf{一般技巧}：若 \(\dfrac{(m-1)S_1^2}{\sigma_1^2}\sim\chi^2(m{-}1)\)，\(\dfrac{(n-1)S_2^2}{\sigma_2^2}\sim\chi^2(n{-}1)\)，则
\[
\frac{S_1^2/S_2^2}{\sigma_1^2/\sigma_2^2}\sim F(m{-}1,n{-}1).
\]
此处 \(\sigma_1^2/\sigma_2^2=4/5\)，故 \(\dfrac{S_1^2/S_2^2}{4/5}=\dfrac{5S_X^2}{4S_Y^2}\sim F(7,9)\)。
\end{solution}



\section{正态总体下的抽样分布}

\begin{note}
\textbf{自学导读：本节是第 7--9 章的"公式源头"。}

下面的三条定理是整个数理统计的核心工具。后面点估计、区间估计、假设检验中出现的所有公式，都可以追溯到这里。建议你：
\begin{enumerate}
    \item 先记住结论的\textbf{形式}：$\bar{X}$ 标准化后是什么分布？$S^2$ 乘以什么系数后是什么分布？
    \item 再理解\textbf{为什么}：$\bar{X}$ 是 $n$ 个正态变量的线性组合，所以仍是正态；$(n-1)S^2/\sigma^2$ 是 $n-1$ 个标准正态平方和，所以是 $\chi^2(n-1)$；用 $S$ 代替 $\sigma$ 后，分母变成随机的，所以从 $N(0,1)$ 变成了 $t(n-1)$。
    \item 特别注意 $\bar{X}$ 与 $S^2$ \textbf{相互独立}这一结论——它不直观，但后面构造 $t$ 统计量和 $F$ 统计量时必须用到。
\end{enumerate}
\end{note}

\subsection{单个正态总体}

\begin{theorem}[单正态总体抽样分布] \label{thm:normal-single}
设总体 \(X \sim N(\mu, \sigma^2)\)，\(X_1, X_2, \dots, X_n\) 为来自 \(X\) 的样本，\(\overline{X}, S^2\) 分别为样本均值和样本方差，则
\begin{enumerate}
    \item \(\displaystyle\frac{\overline{X} - \mu}{\sigma/\sqrt{n}} \sim N(0,1)\)，即 \(\overline{X} \sim N\!\left(\mu, \frac{\sigma^2}{n}\right)\)；
    \item \(\displaystyle\frac{(n-1)S^2}{\sigma^2} = \sum_{i=1}^{n}\!\left(\frac{X_i - \overline{X}}{\sigma}\right)^{\!2} \sim \chi^2(n-1)\)，且 \(\overline{X}\) 与 \(S^2\) 相互独立；
    \item \(\displaystyle\frac{\overline{X} - \mu}{S/\sqrt{n}} \sim t(n-1)\)（当 \(\sigma\) 未知时，用 \(S\) 替代 \(\sigma\)）。
\end{enumerate}
\end{theorem}

\begin{note}
这个定理是后面所有 \(t\) 检验、\(\chi^2\) 检验、区间估计公式的出发点，因此要特别清楚它\textbf{什么时候能直接用}：前提必须是总体本身服从正态分布。若总体不正态，则“\(\overline{X}\) 与 \(S^2\) 独立”“\((n-1)S^2/\sigma^2\sim\chi^2(n-1)\)”这两条一般不能直接照搬。

定理中结论(2)的证明涉及矩阵知识，不要求掌握，但其结论必须牢记。\textbf{注意自由度}：\(\sum(X_i-\overline{X})^2/\sigma^2\) 的自由度为 \(n-1\)（而非 \(n\)），因为 \(\overline{X}\) 已经消耗了一个自由度。进一步有 \(\displaystyle\frac{n(\overline{X}-\mu)^2}{S^2} \sim F(1,n-1)\)。常见误用是把“总体均值未知”和“总体均值恰好等于 0”混成一回事，从而把自由度写错。
\textbf{自学追问：} 为什么这个定理几乎”统治”第 6 章？因为区间估计、\(t\) 检验、\(\chi^2\) 检验、\(F\) 检验的统计量都由它直接导出。
\end{note}

\begin{note}
\textbf{$\bar X$ 与 $S^2$ 独立的直觉解释：}

这个结论违反直觉——$S^2$ 的公式里明明含有 $\bar X$，怎么可能和 $\bar X$ 独立？

\textbf{几何视角：} 把样本 $(X_1,\dots,X_n)$ 看作 $\mathbb{R}^n$ 中的一个点。$\bar X$ 是该点在”全1方向” $\mathbf{1}=(1,\dots,1)$ 上的投影长度（除以 $\sqrt{n}$）；而 $S^2$ 度量的是该点到这条直线的距离。在正态分布下，球对称性保证了”沿某方向的投影”与”到该方向的距离”相互独立。

\textbf{为什么仅正态成立？} 因为只有正态分布的联合密度具有球对称性（等密度面是球面）。换成其他分布（如均匀分布），等密度面不再是球面，投影和距离就不再独立。

\textbf{考试意义：} 凡是需要构造 $t=\dfrac{\bar X-\mu}{S/\sqrt{n}}$ 的地方，都隐含使用了”分子的 $\bar X$ 与分母的 $S$ 独立”这一前提。若题目总体不是正态，这个 $t$ 统计量就不服从 $t$ 分布。
\end{note}

\begin{exercise}\label{exer:6-s2-chi2}
设 \(X_1, X_2, \dots, X_n\) 是总体 \(X \sim N(\mu, \sigma^2)\) 的简单随机样本，\(\overline{X}\) 为样本均值，则下列结论正确的是（\quad）。\par
A. \(\displaystyle\frac{1}{\sigma^2}\sum_{i=1}^{n}(X_i-\overline{X})^2 \sim \chi^2(n-1)\)\qquad B. \(\displaystyle\frac{1}{n}\sum_{i=1}^{n}(X_i-\overline{X})^2 \sim \chi^2(n-1)\)\par
C. \(\displaystyle\frac{1}{n-1}\sum_{i=1}^{n}(X_i-\overline{X})^2 \sim \chi^2(n-1)\)\qquad D. \(\displaystyle\frac{1}{\sigma^2}\sum_{i=1}^{n}(X_i-\overline{X})^2 \sim \chi^2(n)\)
\end{exercise}
\begin{solution}
应选 A。

这道题的核心判据只有一句：
\[
\frac{(n-1)S^2}{\sigma^2}
=\frac{1}{\sigma^2}\sum_{i=1}^{n}(X_i-\overline{X})^2
\sim \chi^2(n-1).
\]

于是 A 与定理完全一致，正确。

B 错在漏掉了 \(\sigma^2\)。\(\dfrac{1}{n}\sum(X_i-\overline{X})^2\) 本身只是一个样本方差型统计量，不会直接服从 \(\chi^2(n-1)\)。

C 也错在没有除以 \(\sigma^2\)。\(\dfrac{1}{n-1}\sum(X_i-\overline{X})^2=S^2\)，它不是卡方分布，而是 \(\chi^2(n-1)\) 再乘上 \(\dfrac{\sigma^2}{n-1}\) 得到的。

D 错在自由度写成了 \(n\)。只要出现 \(\overline{X}\)，就意味着损失了一个自由度，因此应为 \(n-1\)。

所以正确选项是 A。
\end{solution}

\begin{exercise}[样本方差的期望] \label{ex:statistic_05}
设总体 \(X\) 的概率密度为 \(f(x)=\dfrac{1}{2}e^{-|x|}\)（\(-\infty<x<+\infty\)），\(X_1,X_2,\dots,X_n\) 为简单随机样本，样本方差为 \(S^2\)，则 \(E(S^2)=\)\underline{\qquad}。
\end{exercise}
\begin{solution}
由 \(E(S^2)=DX\)，先计算总体矩：
\[
EX=\int_{-\infty}^{+\infty}x\cdot\frac{1}{2}e^{-|x|}\,\mathrm{d}x=0 \quad(\text{奇函数}).
\]
\[
EX^2=\int_{-\infty}^{+\infty}x^2\cdot\frac{1}{2}e^{-|x|}\,\mathrm{d}x
=2\int_0^{+\infty}x^2\cdot\frac{1}{2}e^{-x}\,\mathrm{d}x
=\int_0^{+\infty}x^2\,e^{-x}\,\mathrm{d}x=\Gamma(3)=2.
\]
\[
DX=EX^2-(EX)^2=2,\qquad E(S^2)=DX=2.
\]
\begin{note}
这里利用了拉普拉斯分布 \(f(x)=\frac{1}{2}e^{-|x|}\) 的矩公式，其本质是 \(\text{Exp}(1)\) 分布的偶函数延拓。
\end{note}
\end{solution}


\begin{exercise}\label{exer:6-normal-judge}
设 \(X_1, X_2, \dots, X_n\) 是来自总体 \(X \sim N(\mu, \sigma^2)\) 的简单随机样本，\(\overline{X}\) 与 \(S^2\) 分别为其样本均值和样本方差，则下列结论正确的是（\quad）。\par
A. \(2X_2 - X_1 \sim N(\mu, \sigma^2)\)\quad B. \(\dfrac{n(\overline{X}-\mu)^2}{S^2} \sim F(1,n-1)\)\par
C. \(\dfrac{S^2}{\sigma^2} \sim \chi^2(n-1)\)\quad D. \(\dfrac{\overline{X}-\mu}{S}\sqrt{n-1} \sim t(n-1)\)
\end{exercise}
\begin{solution}
应选 B。

这类题的关键，是先回忆正态总体抽样分布的三个标准结论：
\[
\overline{X}\sim N\!\left(\mu,\frac{\sigma^2}{n}\right),\qquad
\frac{(n-1)S^2}{\sigma^2}\sim\chi^2(n-1),\qquad
\frac{\sqrt{n}(\overline{X}-\mu)}{S}\sim t(n-1).
\]
然后逐项判断。

对 A，\(2X_2-X_1\) 仍服从正态分布，但
\[
E(2X_2-X_1)=2\mu-\mu=\mu,
\]
\[
D(2X_2-X_1)=4DX_2+DX_1=4\sigma^2+\sigma^2=5\sigma^2.
\]
所以它应服从 \(N(\mu,5\sigma^2)\)，不是 \(N(\mu,\sigma^2)\)，故 A 错。

对 B，由
\[
\frac{\sqrt{n}(\overline{X}-\mu)}{S}\sim t(n-1),
\]
两边平方可得
\[
\frac{n(\overline{X}-\mu)^2}{S^2}\sim F(1,n-1),
\]
所以 B 对。

对 C，正确的卡方结论应为
\[
\frac{(n-1)S^2}{\sigma^2}\sim\chi^2(n-1),
\]
题中少了系数 \(n-1\)，故 C 错。

对 D，\(t\) 统计量的正确形式应为
\[
\frac{\sqrt{n}(\overline{X}-\mu)}{S}\sim t(n-1),
\]
前面的系数应是 \(\sqrt{n}\)，不是 \(\sqrt{n-1}\)，故 D 错。

综上，正确选项是 B。
\end{solution}



\begin{exercise}\label{exer:6-esum-sq}
设总体 \(X \sim N(1, 4)\)，\(X_1, X_2, \dots, X_n\) 为来自总体 \(X\) 的样本，\(\overline{X}\) 为样本均值。则
\[
E\!\left(\sum_{i=1}^{n}(X_i-\overline{X})^2\right)=\underline{\hspace{2cm}}.
\]
A. \(4n\)\quad B. \(4(n-1)\)\quad C. \(2n\)\quad D. \(2(n-1)\)
\end{exercise}

\begin{solution}
应选 B。

题目要求的是
\[
E\!\left(\sum_{i=1}^{n}(X_i-\overline{X})^2\right),
\]
而不是 \(E(S^2)\)。最方便的做法是直接套正态总体下的抽样分布结论：
\[
\frac{1}{\sigma^2}\sum_{i=1}^{n}(X_i-\overline{X})^2 \sim \chi^2(n-1).
\]

因为题中总体 \(X\sim N(1,4)\)，所以 \(\sigma^2=4\)。两边取期望，利用 \(\chi^2(n-1)\) 的期望是 \(n-1\)，得
\[
E\!\left[\frac{1}{4}\sum_{i=1}^{n}(X_i-\overline{X})^2\right]=n-1.
\]
于是
\[
E\!\left(\sum_{i=1}^{n}(X_i-\overline{X})^2\right)=4(n-1).
\]

因此应选 B。
\end{solution}

\begin{exercise}\label{exer:6-y-t-dist}
设总体 \(X \sim N(0, 3^2)\)，\(X_1, X_2, \dots, X_8\) 是来自总体 \(X\) 的简单随机样本，则统计量
\[
Y = \frac{X_1 + X_2 + X_3 + X_4}{\sqrt{X_5^2 + X_6^2 + X_7^2 + X_8^2}}
\]
服从（\quad）。\par
A. \(\chi^2(4)\)\quad B. \(t(4)\)\quad C. \(F(4,4)\)\quad D. \(N(0, 4^2)\)
\end{exercise}

\begin{solution}
应选 B。

要判断 \(Y\) 的分布，关键是把它整理成
\[
\frac{U}{\sqrt{V/\nu}}
\]
的形式，其中 \(U\sim N(0,1)\)、\(V\sim\chi^2(\nu)\)，且 \(U,V\) 相互独立。

\[
X_1+X_2+X_3+X_4\sim N(0,4\times 9)=N(0,36),
\]
\[
U = \frac{X_1+X_2+X_3+X_4}{\sqrt{4 \times 9}} = \frac{X_1+X_2+X_3+X_4}{6} \sim N(0,1),
\]
\[
V = \frac{X_5^2+X_6^2+X_7^2+X_8^2}{9} \sim \chi^2(4),
\]
这是因为 \(\dfrac{X_i}{3}\sim N(0,1)\ (i=5,6,7,8)\)，四个标准正态平方和服从自由度为 4 的卡方分布。

又由于 \(U\) 只由 \(X_1,\dots,X_4\) 构成，\(V\) 只由 \(X_5,\dots,X_8\) 构成，而原样本彼此独立，所以 \(U,V\) 相互独立。故
\[
Y = \frac{(X_1+X_2+X_3+X_4)/6}{\sqrt{(X_5^2+\cdots+X_8^2)/4}} = \frac{U}{\sqrt{V/4}} \sim t(4).
\]

因此 \(Y\sim t(4)\)，应选 B。
\end{solution}


\begin{exercise}\label{exer:6-cov-t}
设 \(X_1, X_2\) 是取自正态总体 \(N(\mu, \sigma^2)\) 的简单随机样本，\(Y_1=X_1+X_2\)，\(Y_2=X_1-X_2\)。则 \(\mathrm{Cov}(Y_1, Y_2) =\,\underline{\hspace{2cm}}\)；已知 \((Y_1, Y_2)\) 服从二维正态分布，若 \(c\) 为非零常数，则当 \(c =\,\underline{\hspace{2cm}}\) 时，\(\dfrac{c(Y_1-2\mu)}{|Y_2|}\) 服从自由度为\,\underline{\hspace{2cm}}\,的\,\underline{\hspace{2cm}}\,分布。
\end{exercise}
\begin{solution}
第一空填 \(0\)：
\[
\mathrm{Cov}(Y_1, Y_2) = \mathrm{Cov}(X_1+X_2,\, X_1-X_2) = DX_1 - DX_2 = \sigma^2 - \sigma^2 = 0.
\]
第二空填 \(1\)，第三空填 \(1\)，第四空填 \(t\)。由 \((Y_1,Y_2)\) 服从二维正态分布且协方差为 0，知 \(Y_1, Y_2\) 相互独立。\(Y_1 \sim N(2\mu, 2\sigma^2)\)，故 \(\dfrac{Y_1-2\mu}{\sqrt{2}\,\sigma} \sim N(0,1)\)；\(Y_2 \sim N(0, 2\sigma^2)\)，故 \(\left(\dfrac{Y_2}{\sqrt{2}\,\sigma}\right)^{\!2} \sim \chi^2(1)\)。
构造 \(t\) 分布：
\[
t = \frac{(Y_1-2\mu)/(\sqrt{2}\,\sigma)}{\sqrt{(Y_2/(\sqrt{2}\,\sigma))^2}} = \frac{Y_1-2\mu}{|Y_2|} \sim t(1),
\]
故 \(c = 1\)。
\end{solution}


\begin{exercise}\label{exer:6-f-xbar}
设 \((X_1, X_2, \dots, X_{10})\) 是取自正态总体 \(N(0, 4)\) 的一个样本，\(\overline{X}\) 为样本均值，\(S^2\) 为样本方差。当 \(c =\,\underline{\hspace{2cm}}\) 时，\(c\,\dfrac{\overline{X}^2}{S^2}\) 服从 \(F\) 分布。
\end{exercise}
\begin{solution}
应填 \(10\)。

因为总体 \(X\sim N(0,4)\)，样本量 \(n=10\)，所以
\[
\overline{X}\sim N\!\left(0,\frac{4}{10}\right).
\]
标准化得
\[
\frac{\sqrt{10}\,\overline{X}}{2}\sim N(0,1),
\]
因此
\[
\frac{10\overline{X}^2}{4}\sim\chi^2(1).
\]

另一方面，在正态总体下，
\[
\frac{(n-1)S^2}{\sigma^2}=\frac{9S^2}{4}\sim\chi^2(9),
\]
并且 \(\overline{X}\) 与 \(S^2\) 相互独立。

故
\[
F = \frac{\chi^2(1)/1}{\chi^2(9)/9} = \frac{10\overline{X}^2/4}{S^2/4} = 10\,\frac{\overline{X}^2}{S^2} \sim F(1, 9).
\]

所以所求常数为
\[
c=10.
\]
\end{solution}

\begin{exercise}\label{exer:6-t-prob}
设 \(X_1, X_2, \dots, X_8\) 是来自正态总体 \(N(0, 4)\) 的简单随机样本。常数 \(c =\,\underline{\hspace{2cm}}\) 时，统计量 \(Y = \dfrac{c(X_1-X_2)}{\sqrt{\sum_{i=3}^{8}X_i^2}} \sim\,\underline{\hspace{2cm}}\)；\(P\!\left\{|\overline{X}| \leq \dfrac{\sqrt{2}}{2}\right\} =\,\underline{\hspace{2cm}}\)。
\end{exercise}
\begin{solution}
第一空填 \(\sqrt{3}\)，第二空填 \(t(6)\)，第三空填 \(2\Phi(1)-1\)。

先看统计量 \(Y\)。因为 \(X_1,X_2\) 独立，且都服从 \(N(0,4)\)，所以
\[
X_1-X_2\sim N(0,4+4)=N(0,8),
\]
从而
\[
\frac{X_1-X_2}{\sqrt{8}}\sim N(0,1).
\]

再看分母。对 \(i=3,\dots,8\)，有 \(\dfrac{X_i}{2}\sim N(0,1)\)，故
\[
\sum_{i=3}^{8}X_i^2
=4\sum_{i=3}^{8}\left(\frac{X_i}{2}\right)^2
\sim 4\chi^2(6).
\]
又因为分子只含 \(X_1,X_2\)，分母只含 \(X_3,\dots,X_8\)，所以它们相互独立。

由 \(t\) 分布定义：
\[
t = \frac{(X_1-X_2)/\sqrt{8}}{\sqrt{\sum_{i=3}^{8}X_i^2/(4 \times 6)}} = \sqrt{\frac{24}{8}}\,\frac{X_1-X_2}{\sqrt{\sum_{i=3}^{8}X_i^2}} = \sqrt{3}\,\frac{X_1-X_2}{\sqrt{\sum_{i=3}^{8}X_i^2}} \sim t(6).
\]
故
\[
c=\sqrt{3},\qquad Y\sim t(6).
\]

再算概率。因为
\[
\overline{X}\sim N\!\left(0,\frac{4}{8}\right)=N\!\left(0,\frac12\right),
\]
所以
\[
\sqrt{2}\,\overline{X}\sim N(0,1).
\]
\[
P\!\left\{|\overline{X}| \leq \frac{\sqrt{2}}{2}\right\} = P\!\left\{-1 \leq \sqrt{2}\,\overline{X} \leq 1\right\} = 2\Phi(1)-1.
\]
\end{solution}

\begin{exercise}\label{exer:6-xbar-s2-prob}
设总体 \(X \sim N(1, \sigma^2)\)，\(X_1, X_2, \dots, X_6\) 是取自 \(X\) 的简单随机样本，\(\overline{X}, S\) 分别为样本均值与样本标准差。求 \(P\{\overline{X} > 1\}\) 和 \(P\{\overline{X} < 1,\ S^2 < 1.8472\sigma^2\}\)。\ (\(\chi^2_{0.05}(5) = 9.236\))
\end{exercise}
\begin{solution}
\[
\overline{X}\sim N\!\left(1,\frac{\sigma^2}{6}\right).
\]
由于正态分布关于均值 1 对称，所以 \(P\{\overline{X}>1\}=0.5\)。

又因为总体正态时 \(\overline{X}\) 与 \(S^2\) 相互独立，所以
\[
P\{\overline{X} < 1,\ S^2 < 1.8472\sigma^2\} = P\{\overline{X} < 1\}\,P\{S^2 < 1.8472\sigma^2\}.
\]
其中 \(P\{\overline{X}<1\}=0.5\)。

再由
\[
\frac{5S^2}{\sigma^2}\sim\chi^2(5),
\]
以及 \(\chi^2_{0.05}(5)=9.236\) 表示 \(P\{\chi^2(5)>9.236\}=0.05\)，
可得
\[
P\{S^2 < 1.8472\sigma^2\} = P\!\left\{\frac{5S^2}{\sigma^2} < 9.236\right\} = 1 - 0.05 = 0.95.
\]

因此
\[
P\{\overline{X} < 1,\ S^2 < 1.8472\sigma^2\}=0.5\times 0.95=0.475.
\]

故所求两个概率分别为 \(0.5\) 和 \(0.475\)。
\end{solution}


\subsection{两个正态总体}

\begin{theorem}[双正态总体抽样分布] \label{thm:normal-two}
设总体 \(X \sim N(\mu_1, \sigma_1^2)\)，\(Y \sim N(\mu_2, \sigma_2^2)\)，\(X_1, \dots, X_{n_1}\) 和 \(Y_1, \dots, Y_{n_2}\) 分别为来自 \(X\) 和 \(Y\) 的样本，两组样本相互独立。记 \(\overline{X}, S_1^2\) 和 \(\overline{Y}, S_2^2\) 分别为两组的样本均值与样本方差，则
\begin{enumerate}
    \item \(\displaystyle\frac{S_1^2/\sigma_1^2}{S_2^2/\sigma_2^2} \sim F(n_1-1, n_2-1)\)；
    \item 当 \(\sigma_1^2 = \sigma_2^2 = \sigma^2\) 时，
    \[
    \frac{(\overline{X}-\overline{Y})-(\mu_1-\mu_2)}{S_w\sqrt{\frac{1}{n_1}+\frac{1}{n_2}}} \sim t(n_1+n_2-2),
    \]
    其中 \(S_w^2 = \dfrac{(n_1-1)S_1^2+(n_2-1)S_2^2}{n_1+n_2-2}\) 称为\textbf{合并样本方差}。
\end{enumerate}
\end{theorem}

\begin{note}
这一定理能直接使用的前提有两层：一是两个总体都必须正态，二是两组样本必须彼此独立。第 (1) 条主要用于方差比问题，第 (2) 条只在“\(\sigma_1^2=\sigma_2^2\) 相等但未知”时成立。最常见误用是：方差不相等时仍强行套合并方差 \(S_w^2\) 的 \(t\) 统计量；或者两组样本并非独立却仍按双样本 \(t/F\) 分布处理。
\end{note}

\begin{exercise}[样本均值差的分布] \label{ex:sampling_02}
设总体 \(X\) 与 \(Y\) 相互独立且都服从 \(N(\mu,\sigma^2)\)，\(\overline{X},\overline{Y}\) 分别是来自 \(X,Y\)（容量均为 \(n\)）的样本均值，则当 \(n\) 固定时，\(P\{|\overline{X}-\overline{Y}|>\sigma\}\) 随 \(\sigma\) 的增大而（\quad）。\par
A. 单调增大\qquad B. 单调减小\qquad C. 保持不变\qquad D. 增减不定
\end{exercise}
\begin{solution}
应选 C。由独立性，\(\overline{X}\sim N\!\left(\mu,\frac{\sigma^2}{n}\right)\)，\(\overline{Y}\sim N\!\left(\mu,\frac{\sigma^2}{n}\right)\)，故
\[
\overline{X}-\overline{Y}\sim N\!\left(0,\frac{2\sigma^2}{n}\right).
\]
标准化后：
\[
P\{|\overline{X}-\overline{Y}|>\sigma\}
= P\!\left\{\left|\frac{\overline{X}-\overline{Y}}{\sigma\sqrt{2/n}}\right|>\sqrt{\frac{n}{2}}\right\}
= 2\!\left[1-\Phi\!\left(\sqrt{\frac{n}{2}}\right)\right],
\]
该值仅与 \(n\) 有关，与 \(\sigma\) 无关，故选 C。
\end{solution}


\begin{exercise}\label{exer:6-y-f-dist}
设 \(X_1, X_2, \dots, X_{15}\) 是来自正态总体 \(N(0, 9)\) 的简单随机样本，则统计量
\[
Y = \frac{2X_1^2 + X_2^2 + \cdots + X_{10}^2}{X_{11}^2 + X_{12}^2 + \cdots + X_{15}^2}
\]
服从（\quad）。\par
A. \(\chi^2(10)\)\quad B. \(t(15)\)\quad C. \(F(10,5)\)\quad D. \(F(5,10)\)
\end{exercise}

\begin{solution}
应选 C。
\[
\chi_1^2 = \frac{2X_1^2+X_2^2+\cdots+X_{10}^2}{9}
\]
注意到 \(2X_1^2/9 = 2(X_1/3)^2\)，而 \(X_1/3 \sim N(0,1)\)，故 \((X_1/3)^2 \sim \chi^2(1)\)，从而 \(2(X_1/3)^2\) 服从自由度为 2 的 \(\chi^2\) 分布的尺度变换。更直接的方法是：
\[
Y = \frac{(2X_1^2+X_2^2+\cdots+X_{10}^2)/9}{(X_{11}^2+\cdots+X_{15}^2)/9} = \frac{U/10}{V/5},
\]
其中
\[
U = \dfrac{2X_1^2+\cdots+X_{10}^2}{9},\qquad
V = \dfrac{X_{11}^2+\cdots+X_{15}^2}{9}.
\]
分子等价于 10 个独立 \(N(0,1)\) 变量的平方和，所以 \(U \sim \chi^2(10)\)；同理 \(V \sim \chi^2(5)\)。

故 \(Y = \dfrac{U/10}{V/5} \sim F(10, 5)\)。

\begin{note}
\(X_1/3 \sim N(0,1)\)，故 \((X_1/3)^2 \sim \chi^2(1)\)。由伽马分布的可加性，\(2(X_1/3)^2 \sim \chi^2(2)\)（因为 \(\chi^2(1)\) 即 \(\text{Ga}(1/2,1/2)\)，两倍服从 \(\text{Ga}(1,1/2)=\chi^2(2)\)）。因此 \(2X_1^2/9\) 贡献 2 个自由度，分子共 10 个自由度。
\end{note}
\end{solution}


\section{分布的分位数}

\begin{note}
\textbf{自学导读：分位数是查表的"接口"。}

后面区间估计和假设检验中，你会反复看到 $z_{\alpha/2}$、$t_{\alpha/2}(n-1)$、$\chi^2_{\alpha/2}(n-1)$ 这类符号。它们的含义统一且简单：\textbf{上 $\alpha$ 分位数 $=$ 右尾面积恰好为 $\alpha$ 的那个临界点}。

直觉：想象密度曲线下的面积为 1。从右边切掉面积 $\alpha$ 的那个"刀口"位置，就是上 $\alpha$ 分位数。

\textbf{易错点：} 不同教材可能用"上侧"或"下侧"定义。本书统一用\textbf{上侧}：$P(X>a_\alpha)=\alpha$。做题时先确认教材/试卷用哪种定义，再查表。
\end{note}

\begin{definition}[上 \texorpdfstring{\(\alpha\)}{α} 分位数] \label{def:quantile}
设 \(X\) 为随机变量，\(0 < \alpha < 1\)。若实数 \(a_\alpha\) 满足 \(P\{X > a_\alpha\} = \alpha\)（即上侧尾部面积为 \(\alpha\)），则称 \(a_\alpha\) 为 \(X\) 的\textbf{上 \(\alpha\) 分位数}。
\end{definition}

三大抽样分布的上 \(\alpha\) 分位数分别为：
\begin{enumerate}
    \item \(\chi^2\) 分布：\(P\{\chi^2 > \chi_\alpha^2(n)\} = \alpha\)；
    \item \(t\) 分布：\(P\{t > t_\alpha(n)\} = \alpha\)，由对称性 \(t_{1-\alpha}(n) = -t_\alpha(n)\)；
    \item \(F\) 分布：\(P\{F > F_\alpha(m,n)\} = \alpha\)，且 \(F_{1-\alpha}(m,n) = \dfrac{1}{F_\alpha(n,m)}\)。
\end{enumerate}

\begin{note}
不同教材对分位数的定义可能不同：有的用上侧（\(P\{X > a\} = \alpha\)），有的用下侧（\(P\{X \leq a\} = \alpha\)）。使用时需注意对应关系：
\begin{itemize}
    \item 上侧 \(\alpha\) 分位数 \(a_\alpha\) 对应下侧 \(1-\alpha\) 分位数；
    \item 两种定义下，\(\chi^2\) 分布满足 \(\chi^2_{\alpha,\text{上侧}}(n) = \chi^2_{1-\alpha,\text{下侧}}(n)\)。
\end{itemize}
\end{note}

\begin{note}
\textbf{使用提醒：} 查分位数表前，建议固定做三步检查：先认分布族（标准正态、\(t\)、\(\chi^2\) 还是 \(F\)），再认自由度，最后认清教材用的是上侧定义还是下侧定义。尤其 \(F\) 分布除了尾部方向外，还要额外检查第一、第二自由度的顺序，次序一换，分位数就会跟着变。
\end{note}

\begin{exercise}\label{exer:6-quantile-s2}
设总体 \(X \sim N(\mu, 16)\)，\(X_1, X_2, \dots, X_{10}\) 是取自 \(X\) 的样本，\(S^2\) 为样本方差。求常数 \(a =\,\underline{\hspace{2cm}}\)，使 \(P(S^2 \geq a) = 0.1\) 成立。\par
A. \(\dfrac{16}{9}\chi_{0.1}^2(9)\)\quad B. \(\dfrac{16}{9}\chi_{0.05}^2(10)\)\quad C. \(\chi_{0.1}^2(10)\)\quad D. \(\chi_{0.05}^2(9)\)
\end{exercise}

\begin{solution}
应选 A。

本题先识别样本方差的抽样分布。由于总体 \(X \sim N(\mu,16)\)，样本容量 \(n=10\)，所以
\[
\frac{(n-1)S^2}{\sigma^2}=\frac{9S^2}{16}\sim\chi^2(9).
\]
题目要求
\[
P(S^2 \geq a)=0.1,
\]
等价于
\[
P\!\left(\frac{9}{16}S^2 \geq \frac{9}{16}a\right)=0.1.
\]
而上 \(0.1\) 分位数 \(\chi_{0.1}^2(9)\) 的定义正是右尾概率等于 \(0.1\) 的点，因此必须有
\[
\frac{9}{16}a = \chi_{0.1}^2(9)\implies a = \frac{16}{9}\chi_{0.1}^2(9).
\]
故应选 A。
\end{solution}

\begin{exercise}\label{exer:6-quantile-range}
设随机变量 \(X \sim \chi^2(2)\)，用 \(\chi_\alpha^2(2)\) 表示自由度为 2 的 \(\chi^2\) 分布的上 \(\alpha\) 分位数。已知 \(P(x < X < y) = 0.95\)，\(P(X > y) = 0.02\)，请用分位数表示 \(x\) 和 \(y\)。
\end{exercise}

\begin{solution}
先看右端点 \(y\)。题目已给 \(P(X>y)=0.02\)。
而 \(\chi^2_\alpha(2)\) 的定义正是“使右尾概率等于 \(\alpha\) 的点”，所以 \(y=\chi^2_{0.02}(2)\)。

再求左端点 \(x\)。由 \(P(x<X<y)=0.95\) 以及右尾概率 \(P(X>y)=0.02\)，可知剩下的左侧概率为
\[
P(X\le x)=1-0.95-0.02=0.03.
\]
这一步最容易算错成 \(0.05\) 或 \(0.02\)，本质上是在把总概率 1 分成三段：左边、中间、右边。

于是 \(P(X>x)=1-0.03=0.97\)，按上 \(\alpha\) 分位数的定义，满足 \(P(X>x)=0.97\) 的点就是 \(x=\chi^2_{0.97}(2)\)。

所以
\[
x=\chi^2_{0.97}(2),\qquad y=\chi^2_{0.02}(2).
\]
\end{solution}


\section{综合练习}

\begin{exercise}\label{exer:6-multi-normal}
设 \(X \sim N(0,1)\)，\(X_1, X_2, \dots, X_n\) 是取自 \(X\) 的样本，\(\overline{X}\) 为样本均值，\(S^2\) 为样本方差，则（\quad）。\par
A. \(\overline{X} \sim N(0,1)\)\qquad B. \(n\overline{X} \sim N(0,1)\)\par
C. \(\displaystyle\sum_{i=1}^{n}X_i^2 \sim \chi^2(n)\)\qquad D. \(\dfrac{\overline{X}}{S} \sim t(n-1)\)
\end{exercise}

\begin{solution}
应选 C。

\textbf{这类综合判断题的做法是逐项对照单正态总体抽样分布的三个标准结论。}

A: 由 \(\overline{X}\sim N(0,1/n)\)，可知它不是 \(N(0,1)\)，故 A 错。

B: 因为 \(n\overline{X}=\sum_{i=1}^n X_i\)，所以
\[
n\overline{X}\sim N(0,n),
\]
也不是 \(N(0,1)\)，故 B 错。

C: 由于 \(X_1,\dots,X_n\) 相互独立且都服从 \(N(0,1)\)，所以
\[
\sum_{i=1}^n X_i^2\sim \chi^2(n),
\]
这正是 \(\chi^2\) 分布的定义，故 C 对。

D: \(t\) 分布的标准形式应为
\[
\frac{\overline{X}-\mu}{S/\sqrt{n}}\sim t(n-1).
\]
题中只写 \(\overline{X}/S\)，既少了中心化的 \(\mu\)，也少了 \(\sqrt{n}\) 的缩放，因此 D 错。
\end{solution}

\begin{exercise}\label{exer:6-multi-harbin}
设 \(X_1, X_2, \dots, X_n\) 是来自总体 \(X \sim N(0,1)\) 的简单随机样本，则下列统计量的分布不正确的是（\quad）。\par
A. \(\displaystyle\sum_{i=1}^{n}X_i^2 \sim \chi^2(n)\)\qquad B. \(\displaystyle\frac{\sqrt{n-1}\,X_n}{\sqrt{\sum_{i=1}^{n-1}X_i^2}} \sim t(n-1)\)\par
C. \(\displaystyle\frac{1}{n}\sum_{i=1}^{n}X_i \sim N(0,1)\)\qquad D. \(\displaystyle\frac{(n-1)\sum_{i=1}^{2}X_i^2}{2\sum_{i=3}^{n}X_i^2} \sim F(2,n-2)\)
\end{exercise}

\begin{solution}
应选 C。

\textbf{仍按“定义或标准形式”逐项核对。}

A: 因为 \(X_1,\dots,X_n\) 独立同分布于 \(N(0,1)\)，所以
\[
\sum_{i=1}^{n}X_i^2 \sim \chi^2(n),
\]
故 A 正确。

B: \(X_n \sim N(0,1)\)，且
\[
\sum_{i=1}^{n-1}X_i^2 \sim \chi^2(n-1).
\]
又因 \(X_n\) 与 \(X_1,\dots,X_{n-1}\) 独立，所以
\[
\frac{X_n}{\sqrt{\sum_{i=1}^{n-1}X_i^2/(n-1)}} \sim t(n-1),
\]
即 B 正确。

C: \(\dfrac{1}{n}\sum X_i=\overline{X}\sim N(0,1/n)\)，不是 \(N(0,1)\)，所以 C 错。

D: 前两项平方和与后 \(n-2\) 项平方和分别服从
\[
\sum_{i=1}^{2}X_i^2\sim \chi^2(2),\qquad \sum_{i=3}^{n}X_i^2\sim \chi^2(n-2),
\]
且二者独立，因此
\[
\frac{\sum_{i=1}^{2}X_i^2/2}{\sum_{i=3}^{n}X_i^2/(n-2)} \sim F(2,n-2),
\]
故 D 正确。

因此分布不正确的是 C。
\end{solution}


\begin{exercise}\label{exer:6-multi-dist}
设 \(X_1, X_2, \dots, X_n\) 是总体 \(N(0,1)\) 的简单随机样本，则 \(a =\,\underline{\hspace{2cm}}\) 时 \(a(X_1-X_2)^2 \sim\,\underline{\hspace{2cm}}\)；\(\dfrac{2(X_1-X_2)}{\sqrt{\sum_{i=3}^{n}X_i^2}} \sim\,\underline{\hspace{2cm}}\)。
\end{exercise}
\begin{solution}
第一空填 \(\dfrac{1}{2}\)，第二空填 \(\chi^2(1)\)，第三空填 \(t(8)\)。

先处理前两空。因为 \(X_1,X_2\) 相互独立且都服从 \(N(0,1)\)，所以
\[
X_1-X_2\sim N(0,1+1)=N(0,2).
\]
把它标准化：
\[
\frac{X_1-X_2}{\sqrt2}\sim N(0,1).
\]
对标准正态平方，就得到自由度为 1 的卡方分布：
\[
\left(\frac{X_1-X_2}{\sqrt2}\right)^2=\frac{(X_1-X_2)^2}{2}\sim\chi^2(1).
\]
因此
\[
a=\frac12,\qquad a(X_1-X_2)^2\sim\chi^2(1).
\]

再看第三空。分子已经知道
\[
\frac{X_1-X_2}{\sqrt2}\sim N(0,1).
\]
分母中的
\[
\sum_{i=3}^{n}X_i^2
\]
是 \(n-2\) 个独立标准正态平方和，所以
\[
\sum_{i=3}^{n}X_i^2\sim\chi^2(n-2).
\]
而且它与分子由不同样本项构成，因此相互独立。于是按 \(t\) 分布定义，
\[
\frac{(X_1-X_2)/\sqrt2}{\sqrt{\sum_{i=3}^{n}X_i^2/(n-2)}}\sim t(n-2).
\]
把左式整理成题目给定形式：
\[
\frac{(X_1-X_2)/\sqrt2}{\sqrt{\sum_{i=3}^{n}X_i^2/(n-2)}}
=\sqrt{\frac{n-2}{2}}\cdot\frac{X_1-X_2}{\sqrt{\sum_{i=3}^{n}X_i^2}}.
\]
题目系数是 2，所以令
\[
\sqrt{\frac{n-2}{2}}=2 \implies n-2=8,\qquad n=10.
\]
因此该统计量服从 \(t(n-2)=t(8)\)。

综上，答案依次为 \(\dfrac12\)、\(\chi^2(1)\)、\(t(8)\)。
\end{solution}


\begin{exercise}\label{exer:6-ET-DT}
设总体 \(X \sim N(\mu, \sigma^2)\)，\(X_1, X_2, \dots, X_n\) 是取自 \(X\) 的简单随机样本，\(\overline{X}, S^2\) 分别为样本均值和样本方差。记 \(T = \overline{X}^2 - \dfrac{1}{n}S^2\)。\par
(1) 证明 \(ET = \mu^2\)；\par
(2) 当 \(\mu = 0\)，\(\sigma = 1\) 时，求 \(DT\)。
\end{exercise}
\begin{solution}
(1) \(E(\overline{X}) = \mu\)，\(D(\overline{X}) = \sigma^2/n\)，\(E(S^2) = \sigma^2\)。
\[
E(\overline{X}^2) = D(\overline{X}) + [E(\overline{X})]^2 = \frac{\sigma^2}{n} + \mu^2.
\]
\[
ET = E\!\left(\overline{X}^2 - \frac{1}{n}S^2\right) = \left(\frac{\sigma^2}{n} + \mu^2\right) - \frac{\sigma^2}{n} = \mu^2. 
\]
(2) 当 \(\mu=0\)，\(\sigma=1\) 时，\(X \sim N(0,1)\)。\(\sqrt{n}\,\overline{X} \sim N(0,1)\)，故 \(n\overline{X}^2 \sim \chi^2(1)\)。\((n-1)S^2 \sim \chi^2(n-1)\)。由 \(\overline{X}\) 与 \(S^2\) 独立：
\[
D(\overline{X}^2) = D\!\left(\frac{\chi^2(1)}{n}\right) = \frac{2}{n^2},\quad D\!\left(\frac{S^2}{n}\right) = \frac{1}{n^2}\cdot\frac{2(n-1)}{(n-1)^2} = \frac{2}{n^2(n-1)}.
\]
\[
DT = D(\overline{X}^2) + D\!\left(\frac{S^2}{n}\right) = \frac{2}{n^2} + \frac{2}{n^2(n-1)} = \frac{2n}{n^2(n-1)} = \frac{2}{n(n-1)}.
\]
\end{solution}

\begin{note}
\textbf{自学检查点——抽样分布：}
\begin{enumerate}
    \item 样本方差 $S^2$ 的分母为什么是 $n-1$ 而不是 $n$？（答：为了保证 $E(S^2)=\sigma^2$，即无偏性。直觉上，$\bar{X}$ 已经用掉了一个"自由度"来估计 $\mu$）
    \item $\frac{\bar{X}-\mu}{\sigma/\sqrt{n}}$ 服从什么分布？$\frac{\bar{X}-\mu}{S/\sqrt{n}}$ 呢？（答：前者 $N(0,1)$，后者 $t(n-1)$）
    \item $\chi^2(n)$ 的期望和方差？（答：$n$ 和 $2n$）
    \item $t_\alpha(n)$ 和 $t_{1-\alpha}(n)$ 的关系？（答：$t_{1-\alpha}(n)=-t_\alpha(n)$，由对称性）
    \item $F_\alpha(m,n)$ 和 $F_{1-\alpha}(n,m)$ 的关系？（答：$F_{1-\alpha}(m,n)=1/F_\alpha(n,m)$）
\end{enumerate}
\end{note}



\chapter{点估计}

\begin{note}
\textbf{自学导读：统计推断的第一步——用样本"猜"参数。}

从本章开始，你正式进入统计推断。核心问题：总体分布的形式已知（比如正态），但参数 $\mu$、$\sigma^2$ 未知，你手里只有一组样本数据 $x_1,\dots,x_n$。怎么用这些数据"猜"出参数的值？

\textbf{两种猜法：}
\begin{itemize}
    \item \textbf{矩估计（MME）}：用样本矩代替总体矩，解方程。思路朴素——"样本均值 $\approx$ 总体均值"，直接令 $\bar{X}=\mu$ 就得到 $\hat\mu$。
    \item \textbf{最大似然估计（MLE）}：问"什么参数值最可能产生我观察到的这组数据？"然后选使似然函数最大的参数值。
\end{itemize}

\textbf{自学提醒：} MLE 的计算步骤是固定的（写似然函数 $\to$ 取对数 $\to$ 求导令为零 $\to$ 解方程），不要被不同分布的具体形式吓住——步骤永远一样，只是代入的密度函数不同。

\textbf{前置知识：} 第 2 章的常见分布（正态、指数、均匀等）的密度函数，第 4 章的期望/方差公式，第 6 章的样本均值和样本方差概念。
\end{note}

{\small
\[
\text{点估计}
\begin{cases}
\text{具体内容: 总体分布未知, 寻求适当的统计量估计分布中的未知参数} \\
\text{估计方法}
\begin{cases}
\text{矩估计}
\begin{cases}
\text{思想: 样本矩替换总体矩} \\
\text{掌握: 一阶矩、二阶矩}
\end{cases} \\
\text{最大似然估计}
\begin{cases}
\text{思想: 样本似然函数取最大值时对应的未知参数} \\
\text{掌握: 最大似然估计步骤}
\end{cases}
\end{cases} \\
\text{估计量的评选标准}
\begin{cases}
\text{无偏性: } E(\hat{\theta}) = \theta \\
\text{有效性: 若 } \hat{\theta}_1,\hat{\theta}_2 \text{ 为 } \theta \text{ 的无偏估计量,} \\
\qquad D(\hat{\theta}_1) \leq D(\hat{\theta}_2), \text{ 则 } \hat{\theta}_1 \text{ 比 } \hat{\theta}_2 \text{ 有效} \\
\text{一致性(相合性): } \lim\limits_{n \to \infty} P\left\{|\hat{\theta} - \theta| < \varepsilon\right\} = 1, \\
\qquad \text{也可记为 } \hat{\theta} \xrightarrow{P} \theta \\
\text{均方误差: } \hat{\theta} \text{ 为 } \theta \text{ 的估计量,} \\
\qquad \text{称 } M(\hat{\theta}) = E(\hat{\theta} - \theta)^2 \text{ 为 } \hat{\theta} \text{ 关于 } \theta \text{ 的均方误差}, \\
\quad \text{若 } \hat{\theta}_1,\hat{\theta}_2 \text{ 为 } \theta \text{ 的无偏估计量, 且 } M(\hat{\theta}_1) \leq M(\hat{\theta}_2), \\
\quad \text{则 } \hat{\theta}_1 \text{ 在均方误差下比 } \hat{\theta}_2 \text{ 有效}
\end{cases}
\end{cases}
\]
}

\begin{introduction}
    \item 点估计的基本概念~\ref{def:point-est}
    \item 矩估计法~\ref{subsec:mme}
    \item 最大似然估计法~\ref{subsec:mle}
    \item 估计量的评价标准~\ref{subsec:criteria}
\end{introduction}

\begin{note}
\textbf{使用提醒：} 点估计题通常先回答三个问题：要估哪个参数、准备用什么统计量、最后按什么标准评价它。矩估计更像“用样本矩替换总体矩再反解参数”，最大似然估计更像“让当前样本出现得尽可能合理”；而无偏性、有效性、一致性、均方误差，则是在得到估计量之后再判断它“好不好”的不同标尺。
\end{note}

\section{点估计的基本概念}

\begin{definition}[点估计] \label{def:point-est}
设总体 \(X\) 的分布函数形式已知，但含有一个或多个未知参数 \(\theta\)。利用总体 \(X\) 的一个样本 \[X_1, X_2, \dots, X_n\]
构造一个适当的统计量 \[\hat{\theta}(X_1, X_2, \dots, X_n)\] 
来估计未知参数 \(\theta\)，则称其为 \(\theta\) 的\textbf{估计量}。将样本观测值 \(x_1, x_2, \dots, x_n\)代入估计量，得到的数值 \[\hat{\theta}(x_1, x_2, \dots, x_n)\] 称为 \(\theta\) 的\textbf{估计值}。
\end{definition}

\begin{note}
估计量是随机变量（样本的函数），估计值是具体的数值。常用的点估计方法有\textbf{矩估计法}和\textbf{最大似然估计法}。
\end{note}

\begin{note}
\textbf{自学追问：} 为什么要把“估计量”和“估计值”分开？因为前者是随机变量，后者是具体数值；不分开就没法谈无偏性、方差和一致性。
\end{note}

\section{矩估计法（MME）} \label{subsec:mme}

\begin{definition}[矩] \label{def:moments}
设 \(X, Y\) 为随机变量，\(k, l\) 为正整数。
\begin{enumerate}
    \item 称 \(E(X^k)\) 为 \(X\) 的\textbf{\(k\) 阶原点矩}。
    \item 称 \(E[(X - EX)^k]\) 为 \(X\) 的\textbf{\(k\) 阶中心矩}。
    \item 称 \(E(X^k Y^l)\) 为 \(X\) 与 \(Y\) 的\textbf{\(k+l\) 阶混合原点矩}。
    \item 称 \(E[(X - EX)^k (Y - EY)^l]\) 为 \(X\) 与 \(Y\) 的\textbf{\(k+l\) 阶混合中心矩}。
\end{enumerate}
\end{definition}

\begin{note}
从矩的定义可以看出：数学期望 \(EX\) 是 \(X\) 的一阶原点矩，方差 \(DX\) 是 \(X\) 的二阶中心矩，协方差 \(\mathrm{Cov}(X,Y)\) 是 \(X\) 与 \(Y\) 的二阶混合中心矩。
\end{note}

矩估计法的核心思想是\textbf{用样本矩替换总体矩}：由大数定律，样本矩依概率收敛于总体矩，因此可以用样本矩来估计总体矩，从而解出未知参数。

\begin{property}[矩估计的步骤] \label{prop:mme-steps}
设总体 \(X\) 含有 \(k\) 个未知参数 \(\theta_1, \theta_2, \dots, \theta_k\)，\(X_1, X_2, \dots, X_n\) 为样本。
\begin{enumerate}
    \item 计算总体的前 \(k\) 阶矩 \(\mu_j = E(X^j)\ (j=1,2,\dots,k)\)，它们是 \(\theta_1, \dots, \theta_k\) 的函数。
    \item 计算对应的样本矩 \(A_j = \dfrac{1}{n}\displaystyle\sum_{i=1}^{n}X_i^j\ (j=1,2,\dots,k)\)。
    \item 令 \(\mu_j = A_j\ (j=1,2,\dots,k)\)，建立 \(k\) 个方程。
    \item 解方程组，得到 \(\hat{\theta}_1, \hat{\theta}_2, \dots, \hat{\theta}_k\) 作为矩估计量。
\end{enumerate}
\end{property}

\begin{note}
求矩估计时，以低阶为原则——取能解出参数的最小 \(j\) 值。例如，若一个参数可用一阶矩解出，就不必使用二阶矩。对于两个参数，通常用 \(EX\) 和 \(EX^2\)（或等价地 \(EX\) 和 \(DX\)）建立两个方程。
\end{note}

\begin{exercise}[矩估计量的方差] \label{ex:mom_03}
设总体 \(X\) 的概率密度为
\[
f(x)=
\begin{cases}
\dfrac{6x}{\theta^3}(\theta-x), & 0<x<\theta,\\
0, & \text{其他},
\end{cases}
\]
\(X_1,X_2,\dots,X_n\) 是简单随机样本。求：(1) \(\theta\) 的矩估计量；(2) 该矩估计量的方差。
\end{exercise}
\begin{solution}
(1) 计算总体期望：
\[
EX=\int_0^\theta\frac{6x^2}{\theta^3}(\theta-x)\,\mathrm{d}x=\frac{6}{\theta^3}\!\left[\theta\cdot\frac{\theta^3}{3}-\frac{\theta^4}{4}\right]=\frac{6}{\theta^3}\cdot\frac{\theta^4}{12}=\frac{\theta}{2}.
\]
令 \(\theta/2=\overline{X}\)，得 \(\hat{\theta}=2\overline{X}\)。\\
(2) 先求 \(DX\)：
\[
EX^2=\int_0^\theta\frac{6x^3}{\theta^3}(\theta-x)\,\mathrm{d}x=\frac{6}{\theta^3}\!\left[\theta\cdot\frac{\theta^4}{4}-\frac{\theta^5}{5}\right]=\frac{6}{\theta^3}\cdot\frac{\theta^5}{20}=\frac{3\theta^2}{10}.
\]
\[
DX=EX^2-(EX)^2=\frac{3\theta^2}{10}-\frac{\theta^2}{4}=\frac{6\theta^2-5\theta^2}{20}=\frac{\theta^2}{20}.
\]
因此
\[
D\hat{\theta}=D(2\overline{X})=4D\overline{X}=\frac{4}{n}DX=\frac{4}{n}\cdot\frac{\theta^2}{20}=\frac{\theta^2}{5n}.
\]
\end{solution}



\section{最大似然估计法（MLE）} \label{subsec:mle}

最大似然估计的基本思想：在参数可能的取值范围内，选取使\textbf{样本似然函数}最大的参数值作为估计值。对离散型分布，可以理解为“让当前观测结果出现的概率尽可能大”；对连续型分布，则应理解为“让观测值对应的密度值尽可能大”，不能机械地把它说成概率本身。

\begin{property}[最大似然估计的步骤] \label{prop:mle-steps}
设总体 \(X\) 的分布（离散型为 \(p(x;\theta)\)，连续型为 \(f(x;\theta)\)）含有未知参数 \(\theta\)，样本观测值为 \(x_1, x_2, \dots, x_n\)。
\begin{enumerate}
    \item 写出\textbf{似然函数}：
    \[
    L(\theta) = \prod_{i=1}^{n}p(x_i;\theta)\quad\text{（离散型）}\qquad\text{或}\qquad L(\theta) = \prod_{i=1}^{n}f(x_i;\theta)\quad\text{（连续型）}.
    \]
    \item 取对数：\(\ln L(\theta)\)（将乘积化为求和，便于求导）。
    \item 对 \(\theta\) 求导并令导数为零，解出 \(\hat{\theta}\)：
    \begin{itemize}
        \item \textbf{有驻点}：令 \(\dfrac{\mathrm{d}\ln L}{\mathrm{d}\theta} = 0\)，解出 \(\hat{\theta}\)；
        \item \textbf{单调型（无驻点）}：\(\hat{\theta}\) 在参数范围的\textbf{边界}取得（通常是 \(\max\{X_i\}\) 或 \(\min\{X_i\}\)）；
        \item \textbf{常数型}：\(\hat{\theta}\) 不唯一。
    \end{itemize}
\end{enumerate}
\end{property}

\begin{note}
\textbf{易错点：} 最大似然估计不一定都来自“令导数等于 0”。若似然函数在参数范围内单调，极大值往往直接落在边界点（如样本最大值或最小值）；若参数还影响支持区间，甚至应先写出参数约束，再谈求导。做 MLE 题时，\textbf{先看参数范围和似然是否单调，再看有没有驻点}，通常更不容易漏解。
\end{note}



\begin{property}[最大似然估计的不变性] \label{prop:mle-inv}
设 \(\hat{\theta}\) 是 \(\theta\) 的最大似然估计，\(u = u(\theta)\) 具有反函数，则 \(\hat{u} = u(\hat{\theta})\) 是 \(u(\theta)\) 的最大似然估计。例如，若 \(\hat{\theta}\) 是 \(\theta\) 的 MLE，则 \(e^{\hat{\theta}}\) 是 \(e^\theta\) 的 MLE。
\end{property}


\begin{exercise}[最大似然估计的不变性] \label{ex:mle_04}
设总体 \(X\) 的概率密度为
\[
f(x)=
\begin{cases}
\dfrac{2x}{\alpha^2}, & 0\leq x\leq\alpha,\\
0, & \text{其他},
\end{cases}
\]
其中 \(\alpha>1\) 未知。\(X_1,X_2,\dots,X_n\) 为简单随机样本。\par
(1) 求 \(p=P\{0<X<\sqrt{\alpha}\}\)；\par
(2) 求 \(p\) 的最大似然估计量 \(\hat{p}\)。
\end{exercise}
\begin{solution}
(1)
\[
p=P\{0<X<\sqrt{\alpha}\}=\int_0^{\sqrt{\alpha}}\frac{2x}{\alpha^2}\,\mathrm{d}x=\frac{1}{\alpha^2}\cdot(\sqrt{\alpha})^2=\frac{1}{\alpha}.
\]
(2) 当 \(0\leq x_i\leq\alpha\ (i=1,\dots,n)\) 时，
\[
L(\alpha)=\prod_{i=1}^{n}\frac{2x_i}{\alpha^2}=\frac{2^n\prod_{i=1}^{n}x_i}{\alpha^{2n}}.
\]
\(L(\alpha)\) 关于 \(\alpha\) 单调递减，约束为 \(\alpha\geq\max\{X_1,\dots,X_n\}\)，故
\[
\hat{\alpha}=\max\{X_1,X_2,\dots,X_n\}.
\]
由 MLE 的\textbf{不变性}，\(p=g(\alpha)=1/\alpha\) 的 MLE 为
\[
\hat{p}=\frac{1}{\hat{\alpha}}=\frac{1}{\max\{X_1,X_2,\dots,X_n\}}.
\]
\end{solution}

\begin{note}
\textbf{MLE 单调型的处理步骤详解：} 上题是典型的"似然函数单调"情形。遇到这类题时，按以下步骤操作：
\begin{enumerate}
    \item 写出似然函数 $L(\theta)$，注意\textbf{参数的约束范围}（通常来自密度函数的支撑集，如 $0\leqslant x\leqslant\theta$ 意味着 $\theta\geqslant\max\{x_i\}$）；
    \item 判断 $L(\theta)$（或 $\ln L(\theta)$）关于 $\theta$ 是单调递增还是递减；
    \item 若单调递减：$\theta$ 越小 $L$ 越大，但 $\theta$ 不能小于约束下界 $\to$ $\hat\theta = \max\{X_i\}$；
    \item 若单调递增：$\theta$ 越大 $L$ 越大，但 $\theta$ 不能超过约束上界 $\to$ $\hat\theta = \min\{X_i\}$。
\end{enumerate}
\textbf{口诀：} 递减取 max，递增取 min。
\end{note}

\begin{note}
\textbf{图解"支撑含参数"型 MLE（以 $U(0,\theta)$ 为例）：}

设 $X_1,\dots,X_n\overset{\text{iid}}{\sim}U(0,\theta)$，似然函数为
\[
L(\theta)=\frac{1}{\theta^n}\cdot\mathbf{1}\{\theta\geq X_{(n)}\}.
\]
\textbf{关键理解：} 当 $\theta < X_{(n)}$ 时，$L(\theta)=0$（不可能产生比 $\theta$ 还大的观测值）。当 $\theta\geq X_{(n)}$ 时，$L(\theta)=1/\theta^n$ 关于 $\theta$ 严格递减。

想象 $L(\theta)$ 的图像：$\theta$ 轴上，左半段（$\theta<X_{(n)}$）为零，右半段（$\theta\geq X_{(n)}$）是下降曲线。最大值恰在"断崖边缘" $\theta=X_{(n)}$ 处取得。

\textbf{做题模板：} 凡是密度中出现 $0\leq x\leq\theta$ 或 $\theta\leq x$ 这类条件的：
\begin{enumerate}
    \item 写出似然函数时，\textbf{必须带上示性函数}（约束条件）；
    \item 在约束允许的范围内判断单调性；
    \item 最大值在约束边界取得 $\Rightarrow$ $\hat\theta=X_{(n)}$ 或 $X_{(1)}$。
\end{enumerate}
\end{note}

\begin{exercise}[离散分布的最大似然估计] \label{ex:mle_01}
设总体 \(X\) 的概率分布为
\[
P\{X=1\}=\frac{1-\theta}{2},\quad P\{X=2\}=P\{X=3\}=\frac{1+\theta}{4}.
\]
利用来自总体 \(X\) 的样本值 \(1,3,2,2,1,3,1,2\)，可得 \(\theta\) 的最大似然估计值为（\quad）。\par
A. \(\dfrac{1}{4}\)\qquad B. \(\dfrac{3}{8}\)\qquad C. \(\dfrac{1}{2}\)\qquad D. \(\dfrac{5}{8}\)
\end{exercise}
\begin{solution}
应选 A。样本中 \(X=1\) 出现 3 次，\(X=2\) 出现 3 次，\(X=3\) 出现 2 次。似然函数：
\[
L(\theta)=\left(\frac{1-\theta}{2}\right)^{\!3}\left(\frac{1+\theta}{4}\right)^{\!5}.
\]
对数似然函数：
\[
\ln L(\theta)=-3\ln 2-5\ln 4+3\ln(1-\theta)+5\ln(1+\theta).
\]
令导数为零：
\[
\frac{\mathrm{d}}{\mathrm{d}\theta}\ln L(\theta)=-\frac{3}{1-\theta}+\frac{5}{1+\theta}=0
\implies -3(1+\theta)+5(1-\theta)=0
\implies \theta=\frac{1}{4}.
\]
\end{solution}

\begin{exercise}[矩估计与最大似然估计] \label{ex:est_02}
已知总体 \(X\) 的概率密度为
\[
f(x)=
\begin{cases}
(1+\theta)x^\theta, & 0<x<1,\\
0, & \text{其他},
\end{cases}
\]
其中 \(\theta>-1\) 未知。\(X_1,X_2,\dots,X_n\) 为简单随机样本，求 \(\theta\) 的矩估计量和最大似然估计量。
\end{exercise}
\begin{solution}
\textbf{矩估计。}
\[
\mu_1=EX=\int_0^1(1+\theta)x^{\theta+1}\,\mathrm{d}x=(1+\theta)\left[\frac{x^{\theta+2}}{\theta+2}\right]_0^1=\frac{1+\theta}{2+\theta}.
\]
令 \(\mu_1=\overline{X}\)，即 \(\dfrac{1+\theta}{2+\theta}=\overline{X}\)，解得
\[
1+\theta=(2+\theta)\overline{X}\implies \theta-\theta\overline{X}=2\overline{X}-1\implies \hat{\theta}_{\text{矩}}=\frac{2\overline{X}-1}{1-\overline{X}}.
\]
\textbf{最大似然估计。}似然函数
\[
L(\theta)=\prod_{i=1}^{n}(1+\theta)x_i^\theta=(1+\theta)^n\left(\prod_{i=1}^{n}x_i\right)^{\!\theta},\quad 0<x_i<1.
\]
\[
\ln L(\theta)=n\ln(1+\theta)+\theta\sum_{i=1}^{n}\ln x_i.
\]
\[
\frac{\mathrm{d}}{\mathrm{d}\theta}\ln L(\theta)=\frac{n}{1+\theta}+\sum_{i=1}^{n}\ln x_i=0
\implies \hat{\theta}_{\text{MLE}}=-\frac{n}{\displaystyle\sum_{i=1}^{n}\ln X_i}-1.
\]
\end{solution}


\begin{exercise}\label{exer:7-poisson-bacteria}
为检验某种自然水消毒设备的效果，从消毒后的水中随机抽取 50 升，化验每升水中大肠杆菌的个数（服从泊松分布），结果如下：
\begin{center}
\begin{tabular}{c|ccccccc}
大肠杆菌数/升 & 0 & 1 & 2 & 3 & 4 & 5 & 6 \\
\hline
升数 & 17 & 20 & 10 & 2 & 1 & 0 & 0
\end{tabular}
\end{center}
问：平均每升水中大肠杆菌的个数为多少时，出现上述情况的概率最大？
\end{exercise}
\begin{solution}
设每升水中大肠杆菌个数为 \(X_i\sim P(\lambda)\)，且各升之间可视为相互独立。题目所问“平均每升多少个时，出现当前观测结果的概率最大”，本质上就是要求参数 \(\lambda\) 的最大似然估计。

对于泊松总体，最大似然估计量是样本均值，因此只需根据频数表计算
\[
\overline{x} = \frac{0\times17+1\times20+2\times10+3\times2+4\times1}{50} = \frac{50}{50} = 1.
\]
所以
\[
\hat{\lambda}=1.
\]
也就是说，当平均每升大肠杆菌个数取 \(1\) 时，观察到题目所给结果的似然最大。
\end{solution}

\begin{exercise}\label{exer:7-normal-sigma2}
设 \((X_1, X_2, \dots, X_n)\) 为取自总体 \(N(0, \sigma^2)\) 的样本，则 \(\sigma^2\) 的最大似然估计是（\quad）。\par
A. \(S^2\)\qquad B. \(\dfrac{n-1}{n}S^2\)\qquad C. \(\overline{X}^2\)\qquad D. \(\dfrac{1}{n}\displaystyle\sum_{i=1}^{n}X_i^2\)
\end{exercise}
\begin{solution}
应选 D。关键在于本题中总体均值 \(\mu=0\) 是已知的，因此似然函数应直接围绕 \(\sum x_i^2\) 展开，而不是围绕样本方差 \(S^2\) 展开。

样本的似然函数为
\[
L(\sigma^2) = \frac{1}{(\sqrt{2\pi\sigma^2})^n}\exp\!\left(-\frac{1}{2\sigma^2}\sum_{i=1}^{n}x_i^2\right)~~~~~
\ln L = -\frac{n}{2}\ln(2\pi) - \frac{n}{2}\ln\sigma^2 - \frac{1}{2\sigma^2}\sum_{i=1}^{n}x_i^2.
\]
令 \(\dfrac{\partial\ln L}{\partial\sigma^2} = -\dfrac{n}{2\sigma^2} + \dfrac{1}{2(\sigma^2)^2}\displaystyle\sum_{i=1}^{n}x_i^2 = 0\)，得
\[
\hat{\sigma}^2 = \frac{1}{n}\sum_{i=1}^{n}X_i^2.
\]
因此最大似然估计量正是选项 D。
\end{solution}


\begin{proposition}\label{prop:mle-mme-equiv} 
    对于某些单参数指数族模型，如果其概率质量函数或密度函数可写成
    \[
    P(X=x; \theta)=h(x)c(\theta)\exp\{w(\theta)t(x)\},
    \]
    其中 \(t(x)\) 体现了样本对参数的“自然统计量”结构。若它还能化简为 \(t(x)=x\)，那么一阶矩方程与最大似然方程往往会写成同一个形式，因此矩估计与最大似然估计常常一致。几何分布与本文前面的二项分布就是这样的典型例子。
\end{proposition}
\begin{note}
这个结论\textbf{什么时候可以直接用}：当模型确实能化成单参数指数族，且自然统计量正好就是 \(t(x)=x\) 时，可以迅速判断“一阶矩方程”和“似然方程”会落到同一条等式上。不要把它误推广成“任何单参数分布都满足 MME=MLE”，更不能在多参数模型、\(t(x)\neq x\) 或边界型似然问题里机械套用。
\end{note}
\begin{proof}
\begin{enumerate}
    \item 写出对数似然函数：$$\ln L(\theta) = \sum_{i=1}^n \ln P(X_i; \theta) = \sum_{i=1}^n \ln h(X_i) + n\ln c(\theta) + w(\theta)\sum_{i=1}^n X_i$$
    \item 求导寻找最大似然估计（MLE）：将对数似然函数对 $\theta$ 求导并令其等于0：$$\frac{\mathrm{d}\ln L}{\mathrm{d}\theta} = n\frac{c'(\theta)}{c(\theta)} + w'(\theta)\sum_{i=1}^n X_i = 0$$等式两边同除以 $n \cdot w'(\theta)$，我们得到 MLE 的核心方程：$$-\frac{c'(\theta)}{c(\theta)w'(\theta)} = \frac{1}{n}\sum_{i=1}^n X_i = \overline{X}$$
    \item 再来看看理论期望（总体一阶矩）：在数理统计中有一个基础定理（正则条件下对数似然导数的期望为0），即 $E\left[ \frac{\mathrm{d}\ln P(X; \theta)}{\mathrm{d}\theta} \right] = 0$。代入单样本的导数 $\frac{c'(\theta)}{c(\theta)} + w'(\theta)X$，取期望得到：$$\frac{c'(\theta)}{c(\theta)} + w'(\theta)E(X) = 0 \implies E(X) = -\frac{c'(\theta)}{c(\theta)w'(\theta)}$$
    \item 将第 2 步中的 \(\overline{X}\) 与第 3 步得到的 \(E(X)\) 对应起来，可以发现两者都归结为同一个方程：$$E(X)=\overline{X}$$这正是一阶矩估计的基本方程。
\end{enumerate}
结论：在本题这类“单参数指数族且 \(t(x)=x\)”的模型中，最大似然方程会化成 \(\overline{X}\) 与总体一阶矩的对应关系，因此 MLE 与 MME 往往会得到同一个估计量。但这只是这类模型下的典型现象，不能推广到所有分布。
\end{proof}

\begin{exercise}\label{exer:7-geometric}
设总体 \(X\) 的分布律为 \(P(X=x) = p(1-p)^{x-1}\ (x=1,2,\dots)\)，其中 \(p\ (0<p<1)\) 未知。\(X_1, \dots, X_n\) 为样本，求 \(p\) 的矩估计量和最大似然估计量。
\end{exercise}
\begin{solution}
几何分布的数学期望
\begin{align*}
E(X) = \sum_{x=1}^\infty x \cdot p(1-p)^{x-1} = p \sum_{x=1}^\infty x(1-p)^{x-1} = p \cdot \frac{1}{[1-(1-p)]^2} = p \cdot \frac{1}{p^2} = \frac{1}{p}
\end{align*}
矩估计的核心思想是用样本矩代替总体矩，建立方程求解未知参数。对于单参数问题，我们用一阶原点矩（期望）来求解
即：
$$E(X) = \overline{X},~~~\text{其中}\overline{X} = \frac{1}{n}\sum_{i=1}^n X_i\implies \frac1p=\overline{X}\implies \hat{p}_{\text{矩}} = \frac1{\overline{X}}$$
\textbf{最大似然估计}是想找到一个参数值 $\hat{p}$，使得当前样本出现的概率最大（即似然函数取最大值），似然函数是样本的联合概率（离散型），即所有样本观测值概率的乘积：
\begin{align*}
L(p) &= \prod_{i=1}^n P(X=x_i) = \prod_{i=1}^n \left[ p(1-p)^{x_i-1} \right]=  p^n \cdot (1-p)^{\left( \sum_{i=1}^n x_i \right) - n}\\
\ln L(p) &= n\ln p + \left( \sum_{i=1}^n X_i - n \right) \ln(1-p)\\
\frac{\mathrm{d}}{\mathrm{d}p} \ln L(p) &= \frac{n}{p} - \frac{\sum_{i=1}^n X_i - n}{1-p} = 0\implies n(1-p) = p\left( \sum_{i=1}^n X_i - n \right),~n = p\sum_{i=1}^n X_i
\end{align*}
故 \(\hat{p} = \frac{1}{\overline{X}}\)
\end{solution}

\begin{exercise}\label{exer:7-discrete-3val}
设总体 \(X\) 的概率分布为 \(P\{X=0\} = \theta^2\)，\(P\{X=1\} = 2\theta(1-\theta)\)，\(P\{X=2\} = (1-\theta)^2\)，其中 \(\theta\ (0 < \theta < 1/2)\) 是未知参数。利用样本值 2, 0, 2, 2, 1, 0, 2, 2，求 \(\theta\) 的矩估计值和最大似然估计值。
\end{exercise}
\begin{solution}
\textbf{矩估计}：
\[
EX = 2\theta(1-\theta) + 2(1-\theta)^2 = 2-2\theta.
\]
\(\overline{x} = \frac{11}{8}\)，令 \(2-2\theta = \frac{11}{8}\)，解得 \(\hat{\theta}_{\text{矩}} = \frac{5}{16}\)。\\
\textbf{最大似然估计}：样本值中 \(X=0\) 出现2次，\(X=1\) 出现1次，\(X=2\) 出现5次。
\[
L(\theta) = \theta^4\cdot2\theta(1-\theta)\cdot(1-\theta)^{10} = 2\theta^5(1-\theta)^{11}.
\]
\[
\ln L = \ln 2 + 5\ln\theta + 11\ln(1-\theta),\quad \frac{\mathrm{d}\ln L}{\mathrm{d}\theta} = \frac{5}{\theta} - \frac{11}{1-\theta} = 0.
\]
解得 \(\hat{\theta}_{\text{MLE}} = \frac{5}{16}\)。
\end{solution}

\begin{exercise}\label{exer:7-discrete-123}
设总体 \(X\) 的分布列如下：
\begin{center}
\begin{tabular}{c|ccc}
\(X\) & 1 & 2 & 3 \\
\hline
\(P\) & \(\theta^2\) & \(2\theta(1-\theta)\) & \((1-\theta)^2\)
\end{tabular}
\end{center}
其中 \(\theta\) 为未知参数，样本观测值为 \(X_1=1, X_2=2, X_3=1\)。求 \(\theta\) 的矩估计值和最大似然估计值。
\end{exercise}
\begin{solution}
\textbf{矩估计}：
\[
EX = 1\cdot\theta^2 + 2\cdot2\theta(1-\theta) + 3(1-\theta)^2 = 3-2\theta.
\]
\(\overline{x} = \frac{4}{3}\)，令 \(3-2\theta = \frac{4}{3}\)，得 \(\hat{\theta}_{\text{矩}} = \frac56\)。\\
\textbf{最大似然估计}：
\[
L(\theta) = \theta^2 \cdot 2\theta(1-\theta) \cdot \theta^2 = 2\theta^5(1-\theta).
\]
\(\ln L = \ln 2 + 5\ln\theta + \ln(1-\theta)\)，令 \(\dfrac{5}{\theta} - \dfrac{1}{1-\theta} = 0\)，得 \(\hat{\theta}_{\text{MLE}} = \frac56\)。
\end{solution}

\begin{note}
    几何分布和二项分布都属于指数族分布。下面给出一些别的例子。
\end{note}

\begin{exercise}\label{exer:7-discrete-4val}
设总体 \(X\) 的概率分布为
\begin{center}
\begin{tabular}{c|cccc}
\(X\) & 0 & 1 & 2 & 3 \\
\hline
\(P\) & \(\theta^2\) & \(2\theta(1-\theta)\) & \(\theta^2\) & \(1-2\theta\)
\end{tabular}
\end{center}
其中 \(\theta\ (0 < \theta < 1/2)\) 是未知参数。从总体 \(X\) 中抽取容量为 8 的样本，其样本值为 3, 1, 3, 0, 3, 1, 2, 3。求 \(\theta\) 的矩估计值和最大似然估计值。
\end{exercise}
\begin{solution}
\textbf{矩估计}：总体仅含一个参数，用一阶矩。
\[
EX = 0\cdot\theta^2 + 1\cdot2\theta(1-\theta) + 2\cdot\theta^2 + 3(1-2\theta) = 3-4\theta.
\]
令 \(3-4\theta = \overline{X}\)，\(\overline{x} = \dfrac{3+1+3+0+3+1+2+3}{8} = 2\)，故$\hat{\theta}_{\text{矩}} = \dfrac{3-2}{4} = \dfrac{1}{4}$.\\
\textbf{最大似然估计}：样本值中 \(X=0\) 出现1次，\(X=1\) 出现2次，\(X=2\) 出现1次，\(X=3\) 出现4次。
\[
L(\theta) = \theta^2\cdot[2\theta(1-\theta)]^2\cdot\theta^2\cdot(1-2\theta)^4 = 4\theta^6(1-\theta)^2(1-2\theta)^4.
\]
\[
\ln L(\theta) = \ln 4 + 6\ln\theta + 2\ln(1-\theta) + 4\ln(1-2\theta).
\]
令 \(\dfrac{\mathrm{d}\ln L}{\mathrm{d}\theta} = \dfrac{6}{\theta} - \dfrac{2}{1-\theta} - \dfrac{8}{1-2\theta} = 0\)，化简得 \(12\theta^2 - 14\theta + 3 = 0\)。解得$\theta = \dfrac{14\pm\sqrt{196-144}}{24} = \dfrac{7\pm\sqrt{13}}{12}$因为 \(\dfrac{7+\sqrt{13}}{12} > \dfrac{1}{2}\)（不合题意），故
\[
\hat{\theta}_{\text{MLE}} = \frac{7-\sqrt{13}}{12}.
\]
\end{solution}


\begin{exercise}\label{exer:7-discrete-0123}
设总体 \(X\) 的分布列为
\begin{center}
\begin{tabular}{c|cccc}
\(X\) & 0 & 1 & 2 & 3 \\
\hline
\(P\) & \(\dfrac{\theta(1-\theta)}{2}\) & \(\theta^2\) & \(\dfrac{\theta(1-\theta)}{2}\) & \(1-\theta\)
\end{tabular}
\end{center}
其中 \(\theta\ (0 < \theta < 1)\) 未知。样本观测值为 0, 0, 1, 2, 2, 3, 3, 3, 0, 2。求 \(\theta\) 的矩估计值与最大似然估计值。
\end{exercise}
\begin{solution}
频数统计：\(X=0\) 出现3次，\(X=1\) 出现1次，\(X=2\) 出现3次，\(X=3\) 出现3次。\(\overline{x} = 1.6\)。
\textbf{矩估计}：
\[
EX = \theta^2 + \theta(1-\theta) + 3(1-\theta) = 3-2\theta.
\]
令 \(3-2\theta = 1.6\)，得 \(\hat{\theta}_{\text{矩}} = 0.7\)。\\
\textbf{最大似然估计}：
\[
L(\theta) = \left[\frac{\theta(1-\theta)}{2}\right]^{\!3}\theta^2\!\left[\frac{\theta(1-\theta)}{2}\right]^{\!3}(1-\theta)^3 = \frac{1}{64}\,\theta^8(1-\theta)^9.
\]
\(\ln L = 8\ln\theta + 9\ln(1-\theta) - \ln 64\)，令 \(\dfrac{8}{\theta} - \dfrac{9}{1-\theta} = 0\)，得 \(\hat{\theta}_{\text{MLE}} = \dfrac{8}{17}\)。
\end{solution}

\begin{remark}
以上是离散型的，我们再来看一下连续型的。
\end{remark}

\begin{exercise}\label{exer:7-poisson-tri}
设总体 \(X\) 服从泊松分布 \(P(\theta)\)，总体 \(Y\) 具有密度函数 \(f(y,\lambda) = \dfrac{6y}{\lambda^3}(\lambda-y)\ (0 < y < \lambda)\)。\(X_1, \dots, X_n\) 及 \(Y_1, \dots, Y_n\) 分别为来自 \(X\) 和 \(Y\) 的样本。\par
(1) 求 \(\theta\) 的最大似然估计 \(\hat{\theta}\)；~~(2) 求 \(\lambda\) 的矩估计 \(\hat{\lambda}\)；~~(3) 求 \(\hat{\lambda}\) 的方差。
\end{exercise}
\begin{solution}
(1) 即 \(P(X=x) = e^{-\theta}\dfrac{\theta^x }{x!},\ \theta > 0\)，泊松分布的似然函数：
\[
L(\theta) = \prod_{i=1}^{n}\frac{\theta^{x_i}e^{-\theta}}{x_i!} = e^{-n\theta}\,\theta^{\sum x_i}\prod_{i=1}^{n}\frac{1}{x_i!}.
\]
\[
\ln L = -n\theta + \sum x_i\ln\theta + C,\quad \frac{\partial\ln L}{\partial\theta} = -n + \frac{\sum x_i}{\theta} = 0\implies \hat{\theta} = \overline{X}
\]
(2) \[EY = \displaystyle\int_0^\lambda y\cdot\frac{6y}{\lambda^3}(\lambda-y)\,\mathrm{d}y = \frac{6}{\lambda^3}\!\left[\lambda\int_0^\lambda y^2\,\mathrm{d}y - \int_0^\lambda y^3\,\mathrm{d}y\right] = \frac{6}{\lambda^3}\!\left(\frac{\lambda^4}{3}-\frac{\lambda^4}{4}\right) = \frac{\lambda}{2}\]
令 \(\dfrac{\lambda}{2} = \overline{Y}\)，得 \(\hat{\lambda} = 2\overline{Y}\)。\\
(3) \[EY^2 = \displaystyle\int_0^\lambda \frac{6y^3}{\lambda^3}(\lambda-y)\,\mathrm{d}y = \frac{6}{\lambda^3}\!\left(\frac{\lambda^5}{4}-\frac{\lambda^5}{5}\right) = \frac{3\lambda^2}{10}\]
\(DY = EY^2 - (EY)^2 = \dfrac{3\lambda^2}{10} - \dfrac{\lambda^2}{4} = \dfrac{\lambda^2}{20}\)。
\(D(\hat{\lambda}) = D(2\overline{Y}) = \dfrac{4}{n}\,DY = \dfrac{\lambda^2}{5n}\)。
\end{solution}

\begin{note}
对于单调型似然函数（\(\ln L\) 关于 \(\theta\) 单调），\(\hat{\theta}\) 通常在参数范围的边界取得。常见模式：若 \(L(\theta)\) 关于 \(\theta\) 单调递减，则 \(\hat{\theta} = \max\{X_i\}\)；若单调递增，则 \(\hat{\theta} = \min\{X_i\}\)。
\end{note}

\begin{exercise}[单减取最大，单增取最小]\label{exer:7-unif-theta}
设总体 \(X\) 的概率密度为
\[
f(x;\theta) = \begin{cases} \dfrac{1}{1-\theta}, & \theta \leq x \leq 1, \\[4pt] 0, & \text{其他}, \end{cases}
\]
其中 \(\theta\) 为未知参数，\(X_1, X_2, \dots, X_n\) 为简单随机样本。\par
(1) 求 \(\theta\) 的矩估计量；\par
(2) 求 \(\theta\) 的最大似然估计量。
\end{exercise}
\begin{solution}
(1) 总体 \(X \sim U(\theta, 1)\)，\(EX = \dfrac{1+\theta}{2}\)。令 \(\dfrac{1+\theta}{2} = \overline{X}\)，得 \(\hat{\theta}_{\text{矩}} = 2\overline{X}-1\)。\\
(2) 似然函数为\footnote{密度为 0 的区间直接被过滤了，否则整个乘积为 0，似然函数取最小值，不可能对应最大似然估计。}
\[
L(\theta) = \frac{1}{(1-\theta)^n},\quad \text{约束条件：}\theta \leq \min\{x_1, \dots, x_n\},\ x_i \leq 1.
\]
\(L(\theta)\) 关于 \(\theta\) 单调递增（分母 \(1-\theta\) 递减），故当 \(\theta = \min\{X_1, \dots, X_n\}\) 时 \(L\) 最大：
\[
\hat{\theta}_{\text{MLE}} = \min\{X_1, X_2, \dots, X_n\}.
\]
\end{solution}




\begin{exercise}\label{exer:7-rayleigh}
设某种产品的使用寿命 \(T\) 的分布函数 \(F(t)\) 满足方程
\[
F'(t) + \frac{2t}{\theta^2}[F(t)-1] = 0,\quad t \geq 0,\quad F(0) = 0,
\]
其中 \(\theta > 1\) 为未知参数。产品性能 \(Q(\theta) = \dfrac{\theta^2}{2}\ln\theta - \dfrac{3}{4}\theta^2 + \theta\)。\par
(1) 求 \(P\{T > t\}\) 与 \(P\{T > s+t \mid T > s\}\)（\(s,t > 0\)）；\par
(2) 任取 \(n\) 个产品做寿命试验，测得寿命 \(t_1, \dots, t_n\)，求 \(\theta\) 的最大似然估计值 \(\hat{\theta}\)；\par
(3) 求 \(Q\) 的最大似然估计值 \(\hat{Q}\)。
\end{exercise}

\begin{solution}
\textbf{先求 \(F(t)\)}。微分方程 \(F'(t) + \dfrac{2t}{\theta^2}F(t) = \dfrac{2t}{\theta^2}\) 是一阶线性 ODE，积分因子为 \(e^{t^2/\theta^2}\)：
\[
F(t)\,e^{t^2/\theta^2} = \int e^{t^2/\theta^2}\,\mathrm{d}\!\left(\frac{t^2}{\theta^2}\right) + C = e^{t^2/\theta^2} + C.
\]
由 \(F(0)=0\) 得 \(C=-1\)，故
\[
F(t) = 1 - e^{-t^2/\theta^2}\quad(t \geq 0).
\]
(1) \(P\{T > t\} = e^{-t^2/\theta^2}\)。
\[
P\{T > s+t \mid T > s\} = \frac{P\{T > s+t\}}{P\{T > s\}} = \frac{e^{-(s+t)^2/\theta^2}}{e^{-s^2/\theta^2}} = e^{-(t^2+2st)/\theta^2}.
\]
(2) 密度 \(f(t;\theta) = F'(t) = \dfrac{2t}{\theta^2}\,e^{-t^2/\theta^2}\ (t > 0)\)。
\[
L(\theta) = \theta^{-2n}\cdot 2^n\!\left(\prod_{i=1}^{n}t_i\right)e^{-\sum t_i^2/\theta^2},\quad
\ln L = n\ln 2 - 2n\ln\theta + \sum\ln t_i - \frac{1}{\theta^2}\sum t_i^2.
\]
令 \(\dfrac{\mathrm{d}\ln L}{\mathrm{d}\theta} = -\dfrac{2n}{\theta} + \dfrac{2}{\theta^3}\sum t_i^2 = 0\)，得 \(\hat{\theta} = \displaystyle\left(\dfrac{1}{n}\sum_{i=1}^{n}t_i^2\right)^{\!1/2}\)。\\
(3) 由不变性原则，\(Q'(\theta) = \theta\ln\theta - \theta + 1>0\)，故 \(Q(\theta)\) 有反函数。
\[
\hat{Q} = Q(\hat{\theta}) = \frac{\hat{\theta}^2}{2}\ln\hat{\theta} - \frac{3}{4}\hat{\theta}^2 + \hat{\theta}.
\]
\end{solution}



\section{估计量的评价标准} \label{subsec:criteria}

\begin{definition}[无偏性] \label{def:unbiased}
若 \(E\hat{\theta} = \theta\)，则称 \(\hat{\theta}\) 为 \(\theta\) 的\textbf{无偏估计量}。也就是说，从重复抽样的长期平均效果看，这个估计量既不系统性偏大，也不系统性偏小。
\end{definition}

\begin{note}
无偏性评价的是“平均是否对准真值”，不是“每次估得是否都接近真值”。一个无偏估计量仍可能波动很大；一个有偏估计量也可能因为方差更小而在均方误差意义下更好。因此无偏只是评价标准之一，不是唯一标准。
\end{note}

\begin{exercise}[无偏估计的构造与验证] \label{ex:unbiased_05}
设 \(X_1,X_2,\dots,X_n\) 是来自 \(N(\mu,\sigma^2)\) 的简单随机样本，记
\[
\overline{X}=\frac{1}{n}\sum_{i=1}^{n}X_i,\quad S^2=\frac{1}{n-1}\sum_{i=1}^{n}(X_i-\overline{X})^2,\quad T=\overline{X}^2-\frac{1}{n}S^2.
\]
证明：\(T\) 是 \(\mu^2\) 的无偏估计量。
\end{exercise}
\begin{solution}
要证明 \(T\) 是 \(\mu^2\) 的无偏估计量，只需验证 \(ET=\mu^2\)。

先对 \(T\) 逐项取期望：
\[
ET=E\!\left(\overline{X}^2-\frac{1}{n}S^2\right)
=E(\overline{X}^2)-\frac{1}{n}E(S^2).
\]

其中
\[
E(\overline{X})=\mu,\qquad D(\overline{X})=\frac{\sigma^2}{n},\qquad E(S^2)=\sigma^2.
\]
因此
\[
E(\overline{X}^2)=D(\overline{X})+[E(\overline{X})]^2
=\frac{\sigma^2}{n}+\mu^2.
\]
代回去可得
\[
ET=\left(\frac{\sigma^2}{n}+\mu^2\right)-\frac{1}{n}\sigma^2=\mu^2.
\]

这就说明 \(T\) 的期望恰好等于待估参数 \(\mu^2\)，所以 \(T\) 是 \(\mu^2\) 的无偏估计量。

\begin{note}
此题的技巧是：对 \(\overline{X}^2\) 使用 \(E(\overline{X}^2)=D(\overline{X})+(E\overline{X})^2\) 展开，再利用 \(E(S^2)=\sigma^2\) 消去方差项。此结论在数学三中也常以"求 \(ET\)"的形式出现。
\end{note}
\end{solution}



\begin{proposition}[系数和为1]\label{prop:unbiased-coeff}
若估计量形式为：
\[
\hat \mu = a_1 X_1 + a_2 X_2 + \dots + a_n X_n
\]
则
\[
\hat\mu \text{ 是 } \mu \text{ 的无偏估计} \iff a_1+a_2+\dots+a_n=1
\]
\end{proposition}
\begin{proof}
\[
E(\hat{\mu}) = E\!\left(\sum_{i=1}^{n}a_i X_i\right) = \sum_{i=1}^{n}a_i\,EX_i = \sum_{i=1}^{n}a_i\,\mu = \left(\sum_{i=1}^{n}a_i\right)\mu.
\]
故 \(E(\hat{\mu})=\mu \iff \displaystyle\sum_{i=1}^{n}a_i = 1\)。
\end{proof}

\begin{exercise}\label{exer:7-unbiased-judge}
设 \(X_1, X_2, X_3\) 是总体的样本，则下列统计量不是总体均值 \(\mu\) 的无偏估计量的是（\quad）。\par
A. \(\dfrac{1}{3}(X_1+X_2+X_3)\)~~~ B. \(\dfrac{1}{3}X_1+\dfrac{1}{6}X_2+\dfrac{1}{2}X_3\)~~~
C. \(\dfrac{1}{6}X_1+\dfrac{1}{3}X_2+\dfrac{1}{2}X_3\)~~~ D. \(\dfrac{1}{6}X_1+\dfrac{1}{6}X_2+\dfrac{1}{3}X_3\)
\end{exercise}
\begin{solution}
应选 D。

对于形如
\[
\hat\mu=a_1X_1+a_2X_2+a_3X_3
\]
的估计量，要成为总体均值 \(\mu\) 的无偏估计，充要条件是系数和等于 1，即 \(a_1+a_2+a_3=1\)。

逐项检查：

A 的系数和是 \(\frac13+\frac13+\frac13=1\)，所以无偏。

B 的系数和是 \(\frac13+\frac16+\frac12=1\)，所以无偏。

C 的系数和是 \(\frac16+\frac13+\frac12=1\)，所以无偏。

D 的系数和是 \(\frac16+\frac16+\frac13=\frac23\ne1\)，因此它不是 \(\mu\) 的无偏估计量。

故应选 D。
\end{solution}

\begin{exercise}\label{exer:7-unbiased-c}
设 \(X_1, X_2, \dots, X_n\) 是来自分布 \(f(x;\theta) = \dfrac{3x^2}{\theta^3},0 \leq x \leq \theta\) 的简单随机样本。若统计量 \(g(X_1, \dots, X_n) = c\overline{X}\) 为 \(\theta\) 的无偏估计，则常数 \(c\)为？
\end{exercise}
\begin{solution}
应填 \(\dfrac{4}{3}\)。

要使 \(g(X_1,\dots,X_n)=c\overline{X}\) 成为 \(\theta\) 的无偏估计，必须满足
\[
E(c\overline{X})=\theta.
\]
因此先求总体期望 \(EX\)。

由密度函数
\[
f(x;\theta)=\frac{3x^2}{\theta^3},\qquad 0\le x\le \theta,
\]
可得
\[
EX=\int_0^\theta x\cdot \frac{3x^2}{\theta^3}\,\mathrm dx
=\frac{3}{\theta^3}\int_0^\theta x^3\,\mathrm dx
=\frac{3}{\theta^3}\cdot\frac{\theta^4}{4}
=\frac{3\theta}{4}.
\]

由于 \(E(\overline{X})=EX=\dfrac{3\theta}{4}\)，于是
\[
E(c\overline{X})=cE(\overline{X})=c\cdot\frac{3\theta}{4}.
\]
令它等于 \(\theta\)，得
\[
c\cdot\frac{3\theta}{4}=\theta\implies c=\frac{4}{3}.
\]

所以常数 \(c=\dfrac{4}{3}\)。
\end{solution}

\begin{exercise}\label{exer:7-exp-mle}
设总体 \(X\) 具有概率密度
\[
f(x,\theta) = \begin{cases} \dfrac{1}{\theta}\,e^{-x/\theta}, & x > 0\ (\theta > 0), \\ 0, & x \leq 0, \end{cases}
\]
其中 \(\theta\) 未知，\(X_1, \dots, X_n\) 为样本。求 \(\theta\) 的最大似然估计 \(\hat{\theta}\)，并证明 \(\hat{\theta}\) 是 \(\theta\) 的无偏估计。
\end{exercise}
\begin{solution}
先求最大似然估计。

由样本独立性，似然函数为
\[
L(\theta)=\prod_{i=1}^n \frac{1}{\theta}e^{-X_i/\theta}
=\theta^{-n}\exp\!\left(-\frac{1}{\theta}\sum_{i=1}^n X_i\right),\qquad \theta>0.
\]

取对数得
\[
\ln L(\theta)=-n\ln\theta-\frac{1}{\theta}\sum_{i=1}^n X_i.
\]

对 \(\theta\) 求导：
\[
\frac{\mathrm{d}\ln L}{\mathrm{d}\theta}
=-\frac{n}{\theta}+\frac{1}{\theta^2}\sum_{i=1}^n X_i.
\]
令导数为 0，得
\[
-\frac{n}{\theta}+\frac{1}{\theta^2}\sum_{i=1}^n X_i=0
\quad\Longrightarrow\quad
\hat{\theta}=\frac{1}{n}\sum_{i=1}^n X_i=\overline{X}.
\]

再验证它确实对应极大值。二阶导数为
\[
\frac{\mathrm{d}^2\ln L}{\mathrm{d}\theta^2}
=\frac{n}{\theta^2}-\frac{2}{\theta^3}\sum_{i=1}^n X_i.
\]
在 \(\theta=\overline{X}\) 处，
\[
\left.\frac{\mathrm{d}^2\ln L}{\mathrm{d}\theta^2}\right|_{\theta=\overline{X}}
=\frac{n}{\overline{X}^2}-\frac{2n\overline{X}}{\overline{X}^3}
=-\frac{n}{\overline{X}^2}<0,
\]
所以 \(\hat{\theta}=\overline{X}\) 的确是最大似然估计。

下面证明无偏性。指数分布满足
\[
EX=\theta,
\]
因此
\[
E(\hat{\theta})=E(\overline{X})=\frac{1}{n}\sum_{i=1}^n E(X_i)=\frac{1}{n}\cdot n\theta=\theta.
\]

故 \(\hat{\theta}=\overline{X}\) 是 \(\theta\) 的无偏估计。
\end{solution}


\begin{exercise}\label{exer:7-unbiased-diff}
设 \(X_1, X_2, \dots, X_n\) 是来自正态总体 \(X \sim N(\mu, \sigma^2)\) 的简单随机样本，为使
\[
D = k\sum_{i=1}^{n-1}(X_{i+1}-X_i)^2
\]
成为总体方差 \(\sigma^2\) 的无偏估计量，则 \(k =\,\underline{\hspace{2cm}}\)。
\end{exercise}
\begin{solution}
要使 \(D\) 成为 \(\sigma^2\) 的无偏估计，只需令 \(ED=\sigma^2\)。

先看每一项 \((X_{i+1}-X_i)^2\) 的期望。由于 \(X_i,X_{i+1}\) 相互独立，且都服从 \(N(\mu,\sigma^2)\)，所以
\[
E(X_{i+1}-X_i)=\mu-\mu=0,
\]
\[
D(X_{i+1}-X_i)=DX_{i+1}+DX_i=2\sigma^2.
\]
从而
\[
E\bigl[(X_{i+1}-X_i)^2\bigr]
=D(X_{i+1}-X_i)+\bigl(E(X_{i+1}-X_i)\bigr)^2
=2\sigma^2.
\]

于是
\[
ED
=k\sum_{i=1}^{n-1}E\bigl[(X_{i+1}-X_i)^2\bigr]
=k(n-1)\cdot 2\sigma^2
=2k(n-1)\sigma^2.
\]

令它等于 \(\sigma^2\)，得到
\[
2k(n-1)\sigma^2=\sigma^2 \quad\Longrightarrow\quad k=\frac{1}{2(n-1)}.
\]

故所求 \(k=\frac{1}{2(n-1)}\)。
\end{solution}




\begin{exercise}\label{exer:7-two-normal-mle}
设 \(X_1, \dots, X_m\) 是取自总体 \(N(1, \sigma^2)\) 的简单随机样本，\(Y_1, \dots, Y_n\) 是取自总体 \(N(2, \sigma^2)\) 的简单随机样本。\par
(1) 求联合密度函数；\par
(2) 求 \(\sigma^2\) 的最大似然估计 \(\hat{\sigma}^2\)；\par
(3) \(\hat{\sigma}^2\) 是否为 \(\sigma^2\) 的无偏估计？
\end{exercise}
\begin{solution}
(1) 设 \(T = \displaystyle\sum_{i=1}^{m}(X_i-1)^2 + \sum_{j=1}^{n}(Y_j-2)^2\)，则
\[
f = (2\pi\sigma^2)^{-(m+n)/2}\,\exp\!\left\{-\frac{T}{2\sigma^2}\right\}.
\]
(2) \(\ln f = -\dfrac{m+n}{2}\ln(2\pi\sigma^2) - \dfrac{T}{2\sigma^2}\)，令 \(\dfrac{\partial\ln f}{\partial\sigma^2} = -\dfrac{m+n}{2\sigma^2}+\dfrac{T}{2(\sigma^2)^2}=0\)，得
\[
\hat{\sigma}^2 = \frac{1}{m+n}\!\left[\sum_{i=1}^{m}(X_i-1)^2+\sum_{j=1}^{n}(Y_j-2)^2\right].
\]
(3) 是。\(E\!\left[\sum(X_i-1)^2\right] = m\sigma^2\)，\(E\!\left[\sum(Y_j-2)^2\right] = n\sigma^2\)，故
\[
E(\hat{\sigma}^2) = \frac{(m+n)\sigma^2}{m+n} = \sigma^2. 
\]
\end{solution}


\begin{exercise}\label{exer:7-pareto}
设总体 \(X\) 的分布函数为
\[
F(x,\alpha,\beta) = \begin{cases} 1-\left(\dfrac{\alpha}{x}\right)^{\!\beta}, & x > \alpha, \\ 0, & x \leq \alpha, \end{cases}
\]
其中 \(\alpha > 0\)，\(\beta > 1\)，\(X_1, \dots, X_n\) 为样本。\par
(1) 当 \(\alpha=1\) 时，求 \(\beta\) 的矩估计和最大似然估计；\par
(2) 当 \(\beta=2\) 时，求 \(\alpha\) 的最大似然估计并讨论它的无偏性。
\end{exercise}
\begin{solution}
密度 \(f(x) = \beta\,\alpha^\beta\,x^{-(\beta+1)}\ (x > \alpha)\)。\\
(1) \(\alpha=1\) 时，\(f(x) = \beta\,x^{-(\beta+1)}\ (x > 1)\)。
\[
EX = \int_1^\infty x\cdot\beta\,x^{-(\beta+1)}\,\mathrm{d}x = \frac{\beta}{\beta-1}.
\]
矩估计：令 \(\dfrac{\beta}{\beta-1} = \overline{X}\)，得 \(\hat{\beta}_{\text{矩}} = \dfrac{\overline{X}}{\overline{X}-1}\)。\\
似然函数 \[L(\beta) = \beta^n\!\left(\prod X_i\right)^{-(\beta+1)},~~~\ln L = n\ln\beta - (\beta+1)\sum\ln X_i\]
令 \(n/\beta - \sum\ln X_i = 0\)，得 \(\hat{\beta}_{\text{MLE}} = \dfrac{n}{\sum\ln X_i}\)。\\
(2) \(\beta=2\) 时，\(f(x) = 2\alpha^2 x^{-3}\ (x > \alpha)\)。约束 \(\alpha \leq X_{(1)}\)。\\
\(\ln L = n\ln 2 + 2n\ln\alpha - 3\sum\ln X_i\)，\(\dfrac{\partial\ln L}{\partial\alpha} = \dfrac{2n}{\alpha} > 0\)，故 \(L\) 关于 \(\alpha\) 递增，\(\hat{\alpha} = X_{(1)}\)。\\
对于 Pareto(\(\alpha\),2)，\(E(X_{(1)}) = \dfrac{2n\alpha}{2n-1} \neq \alpha\)，故 \(\hat{\alpha}\) 是有偏估计。
\end{solution}


\begin{exercise}\label{exer:7-mono-3x2}
设总体 \(X\) 的概率密度为 \(f(x,\theta) = \dfrac{3x^2}{\theta^3}\ (0 \leq x \leq \theta)\)，其中 \(\theta > 0\) 未知，\(X_1, \dots, X_n\) 为简单随机样本。\par
(1) 求 \(\theta\) 的矩估计量 \(\hat{\theta}_1\) 和最大似然估计量 \(\hat{\theta}_2\)；\par
(2) 判断 \(\hat{\theta}_2\) 是否是 \(\theta\) 的无偏估计量。
\end{exercise}
\begin{solution}
(1) \(EX = \displaystyle\int_0^\theta \frac{3x^3}{\theta^3}\,\mathrm{d}x = \frac{3\theta}{4}\)。令 \(\dfrac{3\theta}{4} = \overline{X}\)，得 \(\hat{\theta}_1 = \dfrac{4}{3}\overline{X}\)。似然函数 \(L(\theta) = \dfrac{3^n\prod x_i^2}{\theta^{3n}}\) 关于 \(\theta\) 单减，故
\[
\hat{\theta}_2 = \max\{X_1, X_2, \dots, X_n\}.
\]
(2) 由于样本同分布，$P(X_i \le x) = P(X \le x) = F_X(x)$，对于样本最大值，分布函数公式为：
\begin{align*}
F_{\max}(x) &= P\left(\max\{X_1,X_2,\dots,X_n\} \leq x\right) = P\left(X_1 \leq x,\ X_2 \leq x,\ \dots,\ X_n \leq x\right)\\
&= \prod_{i=1}^n P(X_i \leq x) = \prod_{i=1}^n F_X(x) = \left[F_X(x)\right]^n
\end{align*}
求总体 $X$ 的分布函数 $F_X(x)$
$$F_X(x) =
\begin{cases}
0, & x<0 \\
\displaystyle \int_{0}^{x} \frac{3t^2}{\theta^3} dt = \frac{x^3}{\theta^3}, & 0\leq x\leq\theta \\
1, & x>\theta
\end{cases}$$
代入最大值分布公式，当 $0\leq x\leq\theta$ 时：
$$F_{\hat{\theta}_2}(x) = \left(F_X(x)\right)^n = \left(\frac{x^3}{\theta^3}\right)^n = \left(\frac{x}{\theta}\right)^{3n}$$
当 $x<0$ 时 $F=0$，$x>\theta$ 时 $F=1$。密度是分布函数的导数：
$$f_{\hat{\theta}_2}(x) = F'_{\hat{\theta}_2}(x) =\begin{cases}\displaystyle \frac{3n x^{3n-1}}{\theta^{3n}}, & 0\leq x\leq\theta \\0, & \text{其他}\end{cases}$$
根据期望公式 $E(\hat{\theta}_2) = \int_{-\infty}^{+\infty} x f_{\hat{\theta}_2}(x) dx$：
$$\begin{aligned}
E(\hat{\theta}_2) &= \int_{0}^{\theta} x \cdot \frac{3n x^{3n-1}}{\theta^{3n}} dx = \frac{3n}{\theta^{3n}} \int_{0}^{\theta} x^{3n} dx = \frac{3n}{\theta^{3n}} \cdot \left. \frac{x^{3n+1}}{3n+1} \right|_{0}^{\theta} = \frac{3n}{\theta^{3n}} \cdot \frac{\theta^{3n+1}}{3n+1} = \frac{3n}{3n+1} \theta
\end{aligned}$$
故 \(\hat{\theta}_2\) 是\textbf{有偏}估计量。
\begin{note}
作为对比，\(\hat{\theta}_1 = \dfrac{4}{3}\overline{X}\) 满足 \(E(\hat{\theta}_1) = \dfrac{4}{3}\cdot\dfrac{3\theta}{4} = \theta\)，是无偏的。
\end{note}
\end{solution}



\begin{exercise}\label{exer:7-exp4-shift}
设总体 \(X\) 的密度函数为
\[
f_X(x,\theta) = \begin{cases} 4\,e^{-4(x-\theta)}, & x \geq \theta, \\ 0, & \text{其他}, \end{cases}
\]
\(X_1, \dots, X_n\) 为来自总体 \(X\) 的样本。\par
(1) 求 \(\theta\) 的最大似然估计 \(\hat{\theta}\)；\par
(2) \(\hat{\theta}\) 是否为无偏估计？如果不是，修正为无偏估计。
\end{exercise}
\begin{solution}
(1) 构造似然函数：
$$L(\theta) = \prod_{i=1}^n 4 e^{-4(X_i-\theta)} = 4^n e^{-4\sum_{i=1}^n X_i + 4n\theta}, \quad \theta \leq X_{(1)} = \min\{X_i\}$$
$L(\theta)$ 关于 $\theta$ 严格单调递增，因此在定义域右端点取最大值，故最大似然估计量：
$$\hat{\theta} = X_{(1)} = \min\{X_1,X_2,\dots,X_n\}$$
(2) 我们先做一个变量平移：令 $Y_i = X_i - \theta$，则当 $X_i \geq \theta$ 时，$Y_i \geq 0$，代入密度函数：
  $$  f_{Y_i}(y) = f_X(y+\theta,\theta) = 4 e^{-4y}, \quad y \geq 0  $$
  这正是参数为 $\lambda=4$ 的指数分布：$Y_i \sim \text{Exp}(4)$，且 $Y_1,Y_2,\dots,Y_n$ 独立同分布。
样本最小值 $X_{(1)} = \min\{X_i\} = \min\{Y_i+\theta\} = \theta + \min\{Y_i\}$，因此：
  $$  X_{(1)} - \theta = \min\{Y_1,Y_2,\dots,Y_n\}  $$
  即：$\boldsymbol{X_{(1)} - \theta}$ 是 $n$ 个独立 $\text{Exp}(4)$ 变量的最小值。单个 $\text{Exp}(4)$ 的分布函数：$F_Y(y) = 1 - e^{-4y},\ y\geq0$，$n$ 个独立变量的最小值的分布函数：
  $$  F_{\min Y}(y) = P(\min Y_i \leq y) = 1 - P(\min Y_i > y) = 1 - [1-F_Y(y)]^n = 1 - e^{-4ny}  $$
  这正是参数为 $\lambda=4n$ 的指数分布：$\min Y_i \sim \text{Exp}(4n)$。根据指数分布的期望公式：若 $Z \sim \text{Exp}(\lambda)$，则 $E(Z) = \frac{1}{\lambda}$。因此：
  $$  E\left(\min Y_i\right) = \frac{1}{4n}  $$
由 $X_{(1)} = \theta + \min Y_i$，两边取期望：
\begin{align*}
E(\hat{\theta}) = E(X_{(1)}) = E\left(\theta + \min Y_i\right) = \theta + E\left(\min Y_i\right) = \theta + \frac{1}{4n}\neq \theta
\end{align*}
因此 $\hat{\theta} = X_{(1)}$ 是有偏估计量，偏差为 $\frac{1}{4n}$，我们只需要把偏差减掉，构造新的估计量：
$$
\hat{\theta}_{\text{unbiased}} = X_{(1)} - \frac{1}{4n}
$$
验证无偏性：
$$
E\left(\hat{\theta}_{\text{unbiased}}\right) = E\left(X_{(1)} - \frac{1}{4n}\right) = \left(\theta + \frac{1}{4n}\right) - \frac{1}{4n} = \theta
$$
修正后的估计量是 $\theta$ 的无偏估计。
\end{solution}
\begin{note}
题目中的 $f_X(x,\theta)$ 是位置参数型指数分布（也叫平移指数分布）：标准指数分布 $\text{Exp}(\lambda)$ 是 $f(x) = \lambda e^{-\lambda x},\ x\geq0$；平移 $\theta$ 后，分布变为 $f(x) = \lambda e^{-\lambda(x-\theta)},\ x\geq\theta$，其中 $\theta$ 是位置参数（分布的左端点）。在这类“下端点未知、似然函数随 \(\theta\) 单调变化”的典型边界参数模型中，极大似然估计往往落在样本最小值处；本题正是这一类代表情形。
\end{note}

\begin{proposition}[样本方差的无偏性]\label{prop:sample-variance-unbiased}
    样本方差 $S^2$ 是总体方差 $\sigma^2$ 的无偏估计，这个结论对“任意”具有有限方差的总体分布都成立。
\end{proposition}
\begin{note}
这条结论的使用范围非常广：\textbf{不要求总体正态}，只要求总体方差存在即可，因此它比卡方分布结论更稳。不要误把它和“\((n-1)S^2/\sigma^2\sim\chi^2(n-1)\)”混为一谈；后者必须建立在总体正态的前提下。
\end{note}
\begin{proof}
首先给出样本方差的定义
\[
S^2=\frac{1}{n-1}\sum_{i=1}^n (X_i-\overline{X})^2.
\]
将 \(X_i-\overline{X}\) 改写为 \((X_i-\mu)-(\overline{X}-\mu)\)，则
\[
\sum_{i=1}^n (X_i-\overline{X})^2
=\sum_{i=1}^n (X_i-\mu)^2-n(\overline{X}-\mu)^2.
\]
于是取期望得
\[
E\!\left[\sum_{i=1}^n (X_i-\overline{X})^2\right]
=\sum_{i=1}^n E(X_i-\mu)^2-nE(\overline{X}-\mu)^2.
\]
由于 \(E(X_i-\mu)^2=\sigma^2\)，且
\[
E(\overline{X}-\mu)^2=D(\overline{X})=\frac{\sigma^2}{n},
\]
所以
\[
E\!\left[\sum_{i=1}^n (X_i-\overline{X})^2\right]
=n\sigma^2-n\cdot\frac{\sigma^2}{n}=(n-1)\sigma^2.
\]
因此
\[
E(S^2)=\frac{1}{n-1}E\!\left[\sum_{i=1}^n (X_i-\overline{X})^2\right]=\sigma^2.
\]
\end{proof}

\begin{exercise}\label{exer:7-sigma2-unbiased}
设 \(X_1, \dots, X_n\) 和 \(Y_1, \dots, Y_m\) 分别是来自总体 \(X \sim N(\mu_1, \sigma^2)\) 和 \(Y \sim N(\mu_2, \sigma^2)\) 的简单随机样本，样本方差分别为 \(S_X^2, S_Y^2\)。则 \(\sigma^2\) 的无偏估计是（\quad）。\par
A. \(S_X^2 + S_Y^2\)\quad B. \((m-1)S_X^2 + (n-1)S_Y^2\)\quad C. \(\dfrac{S_X^2+S_Y^2}{m+n-2}\)\quad D. \(\dfrac{(m-1)S_X^2+(n-1)S_Y^2}{m+n-2}\)
\end{exercise}
\begin{solution}
应选 D。

判断哪个是 \(\sigma^2\) 的无偏估计，最直接的方法就是分别求期望。

由于样本方差对总体方差都是无偏的，所以
\[
E(S_X^2)=\sigma^2,\qquad E(S_Y^2)=\sigma^2.
\]

于是各选项对应的期望分别为
\[
E(S_X^2+S_Y^2)=2\sigma^2,
\]
\[
E\!\bigl((m-1)S_X^2+(n-1)S_Y^2\bigr)=(m+n-2)\sigma^2,
\]
\[
E\!\left(\frac{S_X^2+S_Y^2}{m+n-2}\right)=\frac{2\sigma^2}{m+n-2},
\]
\[
E\!\left(\frac{(m-1)S_X^2+(n-1)S_Y^2}{m+n-2}\right)
=\frac{(m-1)\sigma^2+(n-1)\sigma^2}{m+n-2}.
\]

最后一个式子化简为
\[
E\!\left(\frac{(m-1)S_X^2+(n-1)S_Y^2}{m+n-2}\right) = \frac{(m-1)\sigma^2+(n-1)\sigma^2}{m+n-2} = \sigma^2.
\]

因此只有 D 的期望恰好等于 \(\sigma^2\)，故 D 是 \(\sigma^2\) 的无偏估计。

这其实就是“无偏估计量做加权平均仍可保持无偏”的一个简单例子。
\end{solution}


\begin{proposition}[泊松分布无偏估计的不唯一性]\label{prop:poisson-unbiased}
    设 $X_1, X_2, \dots, X_n$ 是来自泊松总体 $P(\lambda)$ 的简单随机样本，那么样本均值 $\overline{X}$ 和样本方差 $S^2$ 都是未知参数 $\lambda$ 的无偏估计。
\end{proposition}
\begin{note}
这条命题的重点不只是“会算”，更在于提醒我们：\textbf{无偏估计一般不唯一}。以后做选择题时，如果题目只问“哪个是无偏估计”，不要先入为主地认为答案只能有样本均值这一种；但若进一步比较优劣，还需要看方差或均方误差，不能只看无偏性。
\end{note}
\begin{proof}
泊松总体 \(P(\lambda)\) 满足
\[
E(X)=\lambda,\qquad D(X)=\lambda.
\]
因此对样本均值 \(\overline{X}\)，有
\[
E(\overline{X})=\frac{1}{n}\sum_{i=1}^n E(X_i)=\lambda,
\]
故 \(\overline{X}\) 是 \(\lambda\) 的无偏估计量。

另一方面，由前述“样本方差的无偏性”，可知
\[
E(S^2)=D(X)=\lambda.
\]
所以 \(S^2\) 也是 \(\lambda\) 的无偏估计量。

这说明：对同一个参数 \(\lambda\)，无偏估计量并不唯一。
\end{proof}

\begin{exercise}
    设 $X_1,X_2,\dots,X_n$ 是取自总体 $P(\lambda)$ 的样本，$\overline{X}$ 为样本均值，则 $\lambda^2$ 的一个无偏估计为（\quad）。\\
A. $\overline{X}^2 - \overline{X}$\qquad B. $\overline{X}^2 + \overline{X}$\qquad C. $\overline{X}^2 - \dfrac{1}{n}\overline{X}$\qquad D. $\overline{X}^2 + \dfrac{1}{n}\overline{X}$
\end{exercise}
\begin{solution}
应选 C。

要求的是 \(\lambda^2\) 的无偏估计，因此关键是先把 \(E(\overline{X}^2)\) 算出来。

因为总体服从泊松分布 \(P(\lambda)\)，所以
\[
EX=\lambda,\qquad DX=\lambda.
\]
从而样本均值满足
\[
E(\overline{X})=\lambda,\qquad D(\overline{X})=\frac{\lambda}{n}.
\]

利用恒等式 \(E(T^2)=DT+[E(T)]^2\)，取 \(T=\overline{X}\)，得
\[
E(\overline{X}^2)=D(\overline{X})+[E(\overline{X})]^2
=\frac{\lambda}{n}+\lambda^2.
\]

现在只需把多出来的 \(\dfrac{\lambda}{n}\) 消掉。又因为
\[
E\!\left(\frac{1}{n}\overline{X}\right)=\frac{1}{n}E(\overline{X})=\frac{\lambda}{n},
\]
故
\[
E\!\left(\overline{X}^2-\frac{1}{n}\overline{X}\right)
=\left(\frac{\lambda}{n}+\lambda^2\right)-\frac{\lambda}{n}
=\lambda^2.
\]

因此
\[
\overline{X}^2-\frac{1}{n}\overline{X}
\]
是 \(\lambda^2\) 的无偏估计，故应选 C。
\end{solution}

\begin{proposition}[均匀分布参数的无偏估计]\label{prop:uniform-unbiased}
    设 $X_1, X_2, \dots, X_n$ 是来自均匀总体 $U(0, \theta)$ 的简单随机样本，次序统计量最大值 \[X_{(n)} = \max(X_1, \dots, X_n)\]是 $\theta$ 的极大似然估计量，但它是有偏的。经过修正后的统计量 $T = \frac{n+1}{n}X_{(n)}$ 才是 $\theta$ 的无偏估计。
\end{proposition}
\begin{note}
这个结论适用于“支撑区间上端点未知”的均匀分布模型。直接把 \(X_{(n)}\) 当作无偏估计会系统性偏小，因为样本最大值通常达不到真实上界 \(\theta\)。考试里一旦看到 \(U(0,\theta)\) 或类似边界参数模型，就要警惕“MLE 往往在边界上，但常常有偏”。
\end{note}
\begin{proof}
    首先求出最大值 $X_{(n)}$ 的分布。因为样本相互独立，所以 $X_{(n)} \le x$ 等价于所有的 $X_i \le x$：$$F_{X_{(n)}}(x) = P(X_{(n)} \le x) = \prod_{i=1}^n P(X_i \le x) = \left(\frac{x}{\theta}\right)^n \quad (0 < x < \theta)$$对其求导，得到 $X_{(n)}$ 的概率密度函数：$$f_{X_{(n)}}(x) = \frac{n}{\theta^n}x^{n-1} \quad (0 < x < \theta)$$接下来求 $X_{(n)}$ 的数学期望：$$E(X_{(n)}) = \int_0^\theta x \cdot \frac{n}{\theta^n}x^{n-1} dx = \frac{n}{\theta^n} \int_0^\theta x^n dx = \frac{n}{\theta^n} \cdot \frac{\theta^{n+1}}{n+1} = \frac{n}{n+1}\theta$$由于 $E(X_{(n)}) \ne \theta$，所以极大似然估计量 $X_{(n)}$ 是有偏的（它总是倾向于低估总体的上限 $\theta$）。
    为了消除偏差，我们在等式两边同乘 $\frac{n+1}{n}$：$$E\left(\frac{n+1}{n}X_{(n)}\right) = \frac{n+1}{n} \cdot \frac{n}{n+1}\theta = \theta$$因此，修正后的估计量 $T = \frac{n+1}{n}X_{(n)}$ 是总体参数 $\theta$ 的无偏估计量。
\end{proof}


\begin{definition}[有效性] \label{def:efficiency}
若 \(\hat{\theta}_1, \hat{\theta}_2\) 都是 \(\theta\) 的无偏估计量，且 \(D\hat{\theta}_1 < D\hat{\theta}_2\)，则称 \(\hat{\theta}_1\) 比 \(\hat{\theta}_2\) \textbf{有效}。直观上说，在两者都“不偏”的前提下，谁波动更小，谁就更值得优先采用。
\end{definition}

\begin{note}
有效性比较有一个前提：\textbf{只能在无偏估计量之间比较}。如果某个估计量本身有偏，就不能直接拿“方差更小”来宣布它更有效，而应转去比较均方误差。很多同学会把“有效”误理解成“任何意义下都更好”，这是不准确的。
\end{note}





\begin{exercise}\label{exer:7-efficiency}
设总体 \(X \sim N(\mu, \sigma^2)\)，\(X_1, X_2, X_3\) 是来自总体 \(X\) 的简单随机样本，记
\[
\hat{\mu}_1 = \frac{1}{5}X_1 + \frac{3}{10}X_2 + \frac{1}{2}X_3,\quad
\hat{\mu}_2 = \frac{1}{3}X_1 + \frac{1}{4}X_2 + \frac{5}{12}X_3,\quad
\hat{\mu}_3 = \frac{1}{3}X_1 + \frac{1}{6}X_2 + \frac{1}{2}X_3.
\]
证明 \(\hat{\mu}_1, \hat{\mu}_2, \hat{\mu}_3\) 都是 \(\mu\) 的无偏估计，并指出哪一个最有效。
\end{exercise}

\begin{solution}
先验证无偏性。线性估计量是否对 \(\mu\) 无偏，只需看各系数之和是否为 1。

对 \(\hat{\mu}_1\)：
\[
E\hat{\mu}_1
=\left(\frac15+\frac{3}{10}+\frac12\right)\mu
=\mu.
\]

对 \(\hat{\mu}_2\)：
\[
E\hat{\mu}_2
=\left(\frac13+\frac14+\frac{5}{12}\right)\mu
=\mu.
\]

对 \(\hat{\mu}_3\)：
\[
E\hat{\mu}_3
=\left(\frac13+\frac16+\frac12\right)\mu
=\mu.
\]

所以 \(\hat{\mu}_1,\hat{\mu}_2,\hat{\mu}_3\) 都是 \(\mu\) 的无偏估计。

下面比较谁最有效。由于 \(X_1,X_2,X_3\) 相互独立，且方差都等于 \(\sigma^2\)，所以各线性组合的方差等于“系数平方和乘 \(\sigma^2\)”。

\[
D\hat{\mu}_1
=\left[\left(\frac15\right)^2+\left(\frac{3}{10}\right)^2+\left(\frac12\right)^2\right]\sigma^2
=\left(\frac{1}{25}+\frac{9}{100}+\frac{1}{4}\right)\sigma^2
=\frac{38}{100}\sigma^2,
\]
\[
D\hat{\mu}_2
=\left[\left(\frac13\right)^2+\left(\frac14\right)^2+\left(\frac{5}{12}\right)^2\right]\sigma^2
=\left(\frac{1}{9}+\frac{1}{16}+\frac{25}{144}\right)\sigma^2
=\frac{50}{144}\sigma^2,
\]
\[
D\hat{\mu}_3
=\left[\left(\frac13\right)^2+\left(\frac16\right)^2+\left(\frac12\right)^2\right]\sigma^2
=\left(\frac{1}{9}+\frac{1}{36}+\frac{1}{4}\right)\sigma^2
=\frac{14}{36}\sigma^2.
\]

把它们化成小数：
\[
\frac{38}{100}=0.380,\qquad
\frac{50}{144}\approx 0.347,\qquad
\frac{14}{36}\approx 0.389.
\]
所以
\[
\frac{50}{144}<\frac{38}{100}<\frac{14}{36}.
\]

因此 \(\hat{\mu}_2\) 的方差最小，所以 \(\hat{\mu}_2\) 最有效。
\end{solution}

\begin{exercise}\label{exer:7-eff-N01}
设总体 \(X \sim N(0,1)\)，\(X_1, X_2\) 是从总体抽取的一个样本，下列四个无偏估计量中最有效的是（\quad）。\par
A. \(\mu_1 = X_1\)\qquad B. \(\mu_2 = \dfrac{2}{3}X_1+\dfrac{1}{3}X_2\)\qquad C. \(\mu_3 = \dfrac{1}{4}X_1+\dfrac{3}{4}X_2\)\qquad D. \(\mu_4 = \dfrac{1}{2}X_1+\dfrac{1}{2}X_2\)
\end{exercise}
\begin{solution}
应选 D。

题目已经说明这四个估计量都是无偏的，因此比较“最有效”只需比较方差大小。

又因为总体 \(X\sim N(0,1)\)，所以 \(DX_1=DX_2=1\)，
且 \(X_1,X_2\) 相互独立。于是线性组合的方差等于各项系数平方乘方差后相加。

对 A：\(D\mu_1=DX_1=1\)。

对 B：
\[
D\mu_2=D\!\left(\frac23X_1+\frac13X_2\right)
=\left(\frac23\right)^2+\left(\frac13\right)^2
=\frac59.
\]

对 C：
\[
D\mu_3=\left(\frac14\right)^2+\left(\frac34\right)^2
=\frac{1}{16}+\frac{9}{16}
=\frac58.
\]

对 D：
\[
D\mu_4=\left(\frac12\right)^2+\left(\frac12\right)^2=\frac12.
\]

比较可得
\[
\frac12<\frac59<\frac58<1.
\]
所以 \(\mu_4\) 的方差最小，也就是最有效。

故应选 D。
\end{solution}


\begin{exercise}\label{exer:7-u0-compare}
设样本 \(X_1, X_2, X_3\) 来自总体 \(U(0,\theta)\)，则估计量 \(\hat{\theta}_1=2\overline{X}\)，\(\hat{\theta}_2=X_1+X_2\)，\(\hat{\theta}_3=\dfrac{2}{5}X_1+\dfrac{4}{5}(X_2+X_3)\) 中，哪一个较优？
\end{exercise}
\begin{solution}
先判断三个估计量是否都能作为 \(\theta\) 的无偏估计量来比较。对总体 \(U(0,\theta)\)，有
\[
EX=\frac{\theta}{2},\qquad DX=\frac{\theta^2}{12}.
\]
\begin{note}
对 \(U(0,\theta)\) 来说，\(EX=\theta/2\)，所以线性估计量 \(a_1X_1+a_2X_2+a_3X_3\) 无偏的条件不是“系数和为 1”，而是 \(a_1+a_2+a_3=2\)。这正是这道题最容易误判的地方。
\end{note}
因此
\[
E(\hat{\theta}_1)=E(2\overline{X})=2EX=\theta,
\]
\[
E(\hat{\theta}_2)=E(X_1+X_2)=2EX=\theta,
\]
\[
E(\hat{\theta}_3)=\left(\frac25+\frac45+\frac45\right)EX=2EX=\theta.
\]
所以三者都对 \(\theta\) 无偏，可以继续比较方差。
\[
D\hat{\theta}_1 = 4D(\overline{X}) = \frac{4\theta^2}{12\cdot3} = \frac{\theta^2}{9},
\quad
D\hat{\theta}_2 = 2\cdot\frac{\theta^2}{12} = \frac{\theta^2}{6},
\]
\[
D\hat{\theta}_3 = \left(\frac{4}{25}+\frac{16}{25}+\frac{16}{25}\right)\frac{\theta^2}{12} = \frac{36\theta^2}{300} = \frac{3\theta^2}{25}.
\]
\[
\frac{\theta^2}{9}<\frac{3\theta^2}{25}<\frac{\theta^2}{6}.
\]
故 \(\hat{\theta}_1=2\overline{X}\) 的方差最小，因此它在这三个无偏估计量中最有效，也就是较优的估计量。
\end{solution}



\begin{definition}[一致性（相合性）] \label{def:consistent}
若 \(\hat{\theta}\xrightarrow{P}\theta\)（即对任意 \(\varepsilon > 0\)，\(\lim_{n\to\infty}P\{|\hat{\theta}-\theta| < \varepsilon\} = 1\)），则称 \(\hat{\theta}\) 为 \(\theta\) 的\textbf{一致估计量}（或\textbf{相合估计量}）。这表示样本量越来越大时，估计量偏离真值的概率会越来越小。
\end{definition}

\begin{note}
验证一致性常用的两种方法：
\begin{enumerate}
    \item \textbf{切比雪夫不等式}：若 \(D\hat{\theta} \to 0\)（当 \(n\to\infty\)），则 \(P\{|\hat{\theta}-\theta| \geq \varepsilon\} \leq D\hat{\theta}/\varepsilon^2 \to 0\)。
    \item \textbf{辛钦大数定律}：若 \(\hat{\theta}\) 是样本值的函数且可表示为独立同分布随机变量的均值，则可直接得出依概率收敛。
\end{enumerate}
\textbf{不要混淆无偏与一致}：无偏看的是每个固定 \(n\) 下期望是否等于真值；一致看的是 \(n\to\infty\) 时估计量本身是否逼近真值。二者彼此都不能自动推出对方。
\end{note}


\begin{exercise}\label{exer:7-laplace}
设总体 \(X\) 的概率密度为 \(f(x) = \dfrac{1}{2\theta}\,e^{-|x|/\theta}\ (-\infty < x < \infty)\)，其中 \(\theta > 0\) 未知，\(X_1, \dots, X_n\) 为样本。\par
(1) 求 \(\theta\) 的最大似然估计量 \(\hat{\theta}\)；\par
(2) 证明 \(\hat{\theta}\) 是 \(\theta\) 的无偏估计量；\par
(3) 证明 \(\hat{\theta}\) 是 \(\theta\) 的一致估计量。
\end{exercise}
\begin{solution}
(1) 由样本独立性，似然函数为
\[
L(\theta)=\prod_{i=1}^{n}\frac{1}{2\theta}e^{-|X_i|/\theta}
=(2\theta)^{-n}e^{-\sum|X_i|/\theta}.
\]
取对数得
\[
\ln L(\theta)=-n\ln(2\theta)-\frac{1}{\theta}\sum_{i=1}^{n}|X_i|.
\]
对 \(\theta\) 求导并令其为 0：
\[
\frac{\mathrm d}{\mathrm d\theta}\ln L(\theta)
=-\frac{n}{\theta}+\frac{1}{\theta^2}\sum_{i=1}^{n}|X_i|=0,
\]
故
\[
\hat{\theta}=\frac{1}{n}\sum_{i=1}^{n}|X_i|.
\]

(2) 为证明无偏性，先计算单个样本的 \(E|X|\)：
\[
E|X|
=\frac{1}{2\theta}\int_{-\infty}^{+\infty}|x|e^{-|x|/\theta}\,\mathrm{d}x
=\frac{1}{\theta}\int_0^\infty xe^{-x/\theta}\,\mathrm{d}x.
\]
令 \(t=x/\theta\)，则 \(x=\theta t,\ \mathrm dx=\theta\,\mathrm dt\)，于是
\[
E|X|=\theta\int_0^\infty te^{-t}\,\mathrm{d}t=\theta\,\Gamma(2)=\theta.
\]
因此
\[
E(\hat{\theta})=E\!\left(\frac1n\sum_{i=1}^{n}|X_i|\right)=\frac1n\sum_{i=1}^{n}E|X_i|=\theta,
\]
所以 \(\hat{\theta}\) 是 \(\theta\) 的无偏估计量。

(3) 由于 \(|X_1|,|X_2|,\dots,|X_n|\) 独立同分布，且 \(E|X_i|=\theta\)，由辛钦大数定律：
\[
\hat{\theta} = \frac{1}{n}\sum|X_i| \xrightarrow{P} E|X| = \theta. 
\]
故 \(\hat{\theta}\) 是 \(\theta\) 的一致估计量。
\end{solution}


\begin{exercise}\label{exer:7-consistent}
设 \(X_1, X_2, \dots, X_n\) 是来自总体 \(X\) 的简单随机样本，\(EX = \mu\)，\(DX = \sigma^2\)。证明统计量
\[
Y = \frac{2}{n(n+1)}\sum_{i=1}^{n}iX_i
\]
是 \(\mu\) 的无偏相合估计量。
\end{exercise}
\begin{solution}
\textbf{无偏性}：
\[
EY = \frac{2}{n(n+1)}\sum_{i=1}^{n}i\,EX_i = \frac{2\mu}{n(n+1)}\cdot\frac{n(n+1)}{2} = \mu.
\]
\textbf{一致性}：由 \(\sum i^2 = \dfrac{n(n+1)(2n+1)}{6}\)，
\[
DY = \left[\frac{2}{n(n+1)}\right]^{\!2}\sigma^2\sum_{i=1}^{n}i^2 = \frac{4\sigma^2}{n^2(n+1)^2}\cdot\frac{n(n+1)(2n+1)}{6} = \frac{2(2n+1)}{3n(n+1)}\,\sigma^2.
\]
由切比雪夫不等式，对任意 \(\varepsilon > 0\)：
\[
P\{|Y-\mu| < \varepsilon\} \geq 1 - \frac{DY}{\varepsilon^2} = 1 - \frac{2(2n+1)\sigma^2}{3n(n+1)\varepsilon^2}.
\]
令 \(n\to\infty\)，右端 \(\to 1\)，故 \(Y\xrightarrow{P}\mu\)，\(Y\) 是 \(\mu\) 的相合估计量。
\end{solution}



\begin{definition}[均方误差] \label{def:mse}
设 \(\hat{\theta}\) 为参数 \(\theta\) 的估计量，称
\[
\mathrm{MSE}(\hat{\theta}) = E(\hat{\theta} - \theta)^2
\]
为 \(\hat{\theta}\) 关于 \(\theta\) 的\textbf{均方误差}（Mean Squared Error, MSE）。
\end{definition}

\begin{note}
均方误差把“偏差”和“波动”合在一起衡量，因此它适合比较\textbf{不一定无偏}的估计量。实际使用时，一个估计量即使略有偏差，只要方差降得足够多，也可能拥有更小的均方误差。这正是“有偏但更稳定”的估计量依然有价值的原因。
\end{note}

\begin{note}
\textbf{使用提醒：} 比较估计量时可以按这个顺序想：若题目强调“无偏”，先筛掉有偏候选量，再比较方差看有效性；若候选量本身可能有偏，就转去比较均方误差；若题目关心“大样本下会不会越来越贴近真值”，再看一致性。把这些评价标准放回各自的使用场景，才不会出现“明明在比 MSE，却只盯着无偏性”的误判。
\end{note}

\begin{property}[均方误差的分解] \label{prop:mse-decomp}
均方误差可分解为偏差的平方与方差之和：
\[
\mathrm{MSE}(\hat{\theta}) = [E(\hat{\theta}) - \theta]^2 + D(\hat{\theta}) = \mathrm{Bias}^2(\hat{\theta}) + D(\hat{\theta}).
\]
\begin{proof}
\[
E(\hat{\theta}-\theta)^2 = E\!\left[(\hat{\theta}-E\hat{\theta})+(E\hat{\theta}-\theta)\right]^2 = D(\hat{\theta}) + [E(\hat{\theta})-\theta]^2.
\]
交叉项 \(2E[(\hat{\theta}-E\hat{\theta})(E\hat{\theta}-\theta)] = 0\)。
\end{proof}
\end{property}

\begin{note}
均方误差综合考虑了估计量的\textbf{偏差}和\textbf{方差}：
\begin{itemize}
    \item 若 \(\hat{\theta}\) 是无偏估计（\(E\hat{\theta}=\theta\)），则 \(\mathrm{MSE}(\hat{\theta}) = D(\hat{\theta})\)，此时均方误差最小等价于方差最小（即有效性）。
    \item 若 \(\hat{\theta}_1, \hat{\theta}_2\) 为 \(\theta\) 的估计量（不一定无偏），且 \(\mathrm{MSE}(\hat{\theta}_1) \leq \mathrm{MSE}(\hat{\theta}_2)\)，则称 \(\hat{\theta}_1\) 在均方误差意义下比 \(\hat{\theta}_2\) 有效。
    \item 一个有偏但方差很小的估计量，其均方误差可能\textbf{优于}无偏但方差较大的估计量。
\end{itemize}
\end{note}


\begin{exercise}\label{exer:7-shifted-exp}
设总体 \(X\) 的概率密度为
\[
f(x) = \begin{cases} e^{\theta-x}, & x > \theta, \\ 0, & x \leq \theta, \end{cases}
\]
\(\theta\) 未知，\(X_1, \dots, X_n\) 为样本。\par
(1) 求 \(\theta\) 的矩估计量 \(\hat{\theta}_M\) 和最大似然估计量 \(\hat{\theta}_L\)；\par
(2) 比较 \(\hat{\theta}_M\) 与 \(\hat{\theta}_L\) 的优劣。
\end{exercise}
\begin{solution}
(1) \(EX = \displaystyle\int_\theta^\infty x\,e^{\theta-x}\,\mathrm{d}x = \theta + 1\)。令 \(\theta+1 = \overline{X}\)，得 \(\hat{\theta}_M = \overline{X}-1\)。
\(L(\theta) = e^{n\theta-\sum X_i}\)，关于 \(\theta\) 递增（\(\theta \leq X_{(1)}\)），故 \(\hat{\theta}_L = X_{(1)}\)。\\
(2) \(\hat{\theta}_M\) 无偏。\(DX = E(X^2)-(EX)^2 = \displaystyle\int_\theta^\infty(x-\theta-1)^2e^{\theta-x}\,\mathrm{d}x\)。令 \(t=x-\theta\)：\(\displaystyle\int_0^\infty(t-1)^2e^{-t}\,\mathrm{d}t = \Gamma(3)-2\Gamma(2)+\Gamma(1) = 2-2+1 = 1\)，故 \(D(\hat{\theta}_M) = D(\overline{X}) = 1/n\)，\(\mathrm{MSE}_M = 1/n\)。\\
\(X_{(1)}-\theta\) 为 \(n\) 个 \(\mathrm{Exp}(1)\) 的最小值，\(E(X_{(1)}-\theta) = 1/n\)，\(D(X_{(1)}-\theta) = 1/n^2\)。故 \(E(\hat{\theta}_L) = \theta + 1/n\)，\(D(\hat{\theta}_L) = 1/n^2\)，
\[
\mathrm{MSE}_L = \left(\frac{1}{n}\right)^{\!2} + \frac{1}{n^2} = \frac{2}{n^2}.
\]
当 \(n > 2\) 时，\(2/n^2 < 1/n\)，\(\hat{\theta}_L\) 的均方误差更小，更优。
\end{solution}

\begin{note}
\textbf{自学检查点——点估计：}
\begin{enumerate}
    \item 矩估计的核心思想是什么？（答：用样本矩代替总体矩，解方程）
    \item MLE 的四步是什么？（答：写似然函数 $\to$ 取对数 $\to$ 求导令为零 $\to$ 解方程）
    \item 似然函数单调递增时，MLE 怎么取？（答：取参数允许范围的上界，通常是 $X_{(n)}$ 或 $X_{(1)}$）
    \item 无偏性的定义？（答：$E(\hat\theta)=\theta$）
    \item $S^2$ 是 $\sigma^2$ 的无偏估计，那 $S$ 是 $\sigma$ 的无偏估计吗？（答：不是！$E(S)\neq\sigma$，因为期望和开根号不能交换）
\end{enumerate}
\end{note}


\begin{note}
\textbf{【高频题型模板】MLE 完整解题流程：}
\begin{enumerate}
    \item \textbf{写似然函数：} $L(\theta)=\displaystyle\prod_{i=1}^n f(x_i;\theta)$（连续型）或 $\prod P(X=x_i;\theta)$（离散型）。
    \item \textbf{取对数：} $\ln L(\theta)=\sum_{i=1}^n \ln f(x_i;\theta)$。
    \item \textbf{判断类型：}
    \begin{itemize}
        \item 若 $\ln L$ 对 $\theta$ 可导且有驻点 $\Rightarrow$ 令 $\dfrac{\mathrm{d}\ln L}{\mathrm{d}\theta}=0$ 解出 $\hat\theta$；
        \item 若 $\ln L$ 单调（支撑含参数）$\Rightarrow$ 在约束边界取极值（递减取 max，递增取 min）。
    \end{itemize}
    \item \textbf{验证：} 确认是极大值（二阶导 $<0$ 或单调性判断）。
\end{enumerate}
\textbf{常见追问：} 求出 $\hat\theta$ 后，题目可能继续问 $g(\theta)$ 的 MLE $\Rightarrow$ 直接用不变性：$\widehat{g(\theta)}=g(\hat\theta)$。
\end{note}


\chapter{区间估计}

\begin{note}
\textbf{自学导读：从"一个数"到"一个区间"。}

第 7 章的点估计给出一个数 $\hat\mu = \bar{x}$，但你无法知道这个数离真值 $\mu$ 有多近。区间估计的想法是：给出一个区间 $[\bar{x}-\Delta,\ \bar{x}+\Delta]$，并保证"真值 $\mu$ 落在这个区间内的概率 $\geqslant 1-\alpha$"。

\textbf{核心逻辑（三步走）：}
\begin{enumerate}
    \item 找一个包含未知参数 $\theta$ 和样本的统计量（叫\textbf{枢轴量}），其分布已知且不含未知参数；
    \item 根据置信水平 $1-\alpha$，从分布表中查出临界值；
    \item 解不等式，把 $\theta$ 夹在中间，得到置信区间。
\end{enumerate}

\textbf{自学提醒：} 不要死记公式表。记住"$\sigma$ 已知用 $Z$，$\sigma$ 未知用 $t$，估方差用 $\chi^2$"这三条规则，再按三步走推导，公式自然出来。

\textbf{与第 6 章的联系：} 枢轴量的分布正是第 6 章学过的抽样分布定理——$\frac{\bar{X}-\mu}{\sigma/\sqrt{n}}\sim N(0,1)$，$\frac{\bar{X}-\mu}{S/\sqrt{n}}\sim t(n-1)$，$\frac{(n-1)S^2}{\sigma^2}\sim\chi^2(n-1)$。本章只是把这些结论"反过来用"。
\end{note}


% 单个正态总体（完整保留所有内容，优化行距和排版）
{\small
\[
\begin{cases}
\text{单个正态 } X\sim N(\mu,\sigma^2) \\
\text{置信水平为 } 1-\alpha
\end{cases}
\Rightarrow
\begin{cases}
\text{总体均值 } \mu:
\begin{cases}
\sigma^2 \text{ 已知}: & \left( \overline{X} - \dfrac{\sigma}{\sqrt{n}} z_{\frac{\alpha}{2}}, \overline{X} + \dfrac{\sigma}{\sqrt{n}} z_{\frac{\alpha}{2}} \right)\\[4pt]
\sigma^2 \text{ 未知}: & \left( \overline{X} - \dfrac{S}{\sqrt{n}} t_{\frac{\alpha}{2}}(n-1), \overline{X} + \dfrac{S}{\sqrt{n}} t_{\frac{\alpha}{2}}(n-1) \right)
\end{cases}\\[8pt]
\text{总体方差 } \sigma^2:
\begin{cases}
\mu \text{ 已知}: & \left( \dfrac{\sum_{i=1}^n (X_i-\mu)^2}{\chi_{\frac{\alpha}{2}}^2(n)}, \dfrac{\sum_{i=1}^n (X_i-\mu)^2}{\chi_{1-\frac{\alpha}{2}}^2(n)} \right)\\[4pt]
\mu \text{ 未知}: & \left( \dfrac{(n-1)S^2}{\chi_{\frac{\alpha}{2}}^2(n-1)}, \dfrac{(n-1)S^2}{\chi_{1-\frac{\alpha}{2}}^2(n-1)} \right)
\end{cases}
\end{cases}
\]
}

\vspace{1em}
\[
\begin{aligned}
&\text{两个正态 } \\
&X_1 \sim N(\mu_1,\sigma_1^2);\\
& X_2 \sim N(\mu_2,\sigma_2^2) \\
&\text{置信水平为 }1-\alpha
\end{aligned}
\;
\left\{
\begin{aligned}
&\text{总体均值差 } (\mu_1-\mu_2): \\
&\qquad
\begin{cases}
\sigma_1^2,\sigma_2^2 \text{ 已知}: \\[2pt]
\quad (\bar X_1-\bar X_2) \pm z_{\alpha/2} \sqrt{\sigma_1^2/n_1+\sigma_2^2/n_2};\\[8pt]
\sigma_1^2=\sigma_2^2=\sigma^2,\;\sigma^2 \text{ 未知}: \\[2pt]
\quad(\bar X_1-\bar X_2) \pm t_{\alpha/2}(n_1+n_2-2)\,S_w\sqrt{1/n_1+1/n_2};
\end{cases}
\\[12pt]
&\text{总体方差比 } \bigl( \sigma_1^2/\sigma_2^2 \bigr): \\
&\qquad
\begin{cases}
\mu_1,\mu_2 \text{ 已知}: \\[2pt]
\quad\Bigl( \frac{\frac1{n_1}\sum_{i=1}^{n_1}(X_i-\mu_1)^2}{\frac1{n_2}\sum_{j=1}^{n_2}(X_j-\mu_2)^2}
\cdot\frac1{F_{\alpha/2}(n_1,n_2)},\;
\frac{\frac1{n_1}\sum_{i=1}^{n_1}(X_i-\mu_1)^2}{\frac1{n_2}\sum_{j=1}^{n_2}(X_j-\mu_2)^2}
\cdot\frac1{F_{1-\alpha/2}(n_1,n_2)} \Bigr) \\[8pt]
\mu_1,\mu_2 \text{ 未知}: \\[2pt]
\quad\Bigl( \frac{S_1^2}{S_2^2}\cdot\frac1{F_{\alpha/2}(n_1-1,n_2-1)},\;
\frac{S_1^2}{S_2^2}\cdot\frac1{F_{1-\alpha/2}(n_1-1,n_2-1)} \Bigr)
\end{cases}
\end{aligned}
\right.
\]


\begin{introduction}
  \item 区间估计的基本概念：双侧置信区间与单侧置信区间
  \item 构造置信区间的基本原理与枢轴量法
  \item 正态总体参数（均值、方差）的区间估计
\end{introduction}

\begin{note}
\textbf{使用提醒：} 区间估计题的选型可以按“估谁 + 总体条件 + 单侧/双侧”三步走。先认是均值、方差、均值差还是方差比；再看总体是否正态、方差已知还是未知；最后根据“至少/至多/是否等于”判断该构造单侧区间还是双侧区间。把这三步先理顺，再去选 \(z/t/\chi^2/F\) 会稳很多。
\end{note}

\section{区间估计的基本概念}

对于未知参数 $\theta$，除了求出它的点估计外，我们往往还希望给出一个合理的范围，并知道这个范围包含参数 $\theta$ 真值的可信程度。这样的范围通常以区间的形式给出，这种估计形式被称为\textbf{区间估计}。

\begin{definition}[置信区间与置信水平] \label{def:confidence_interval}
设 $\theta$ 是总体 $X$ 的分布函数中的一个未知参数，对于给定的显著性水平 $\alpha\ (0<\alpha<1)$，如果由样本 $X_1,X_2,\dots,X_n$ 确定的两个统计量 $\underline{\theta}=\underline{\theta}(X_1,X_2,\dots,X_n)$ 和 $\overline{\theta}=\overline{\theta}(X_1,X_2,\dots,X_n)$，满足 $\underline{\theta} < \overline{\theta}$，且对于任意的 $\theta \in \Theta$ 均有：
\begin{equation}
P\{\underline{\theta}<\theta<\overline{\theta}\} \ge 1-\alpha
\end{equation}
则称随机区间 $(\underline{\theta},\overline{\theta})$ 是 $\theta$ 的置信度（或置信水平）为 $1-\alpha$ 的\textbf{双侧置信区间}。其中，$\underline{\theta}$ 和 $\overline{\theta}$ 分别称为置信度为 $1-\alpha$ 的\textbf{置信下限}和\textbf{置信上限}。

特别地，如果满足 $P\{\theta < \underline{\theta}\} = P\{\theta > \overline{\theta}\} = \frac{\alpha}{2}$，则称这种置信区间为\textbf{等尾置信区间}。
\end{definition}

\begin{note}
置信区间最容易混淆的地方是“概率到底作用在谁身上”。参数 \(\theta\) 是固定但未知的，真正有随机性的是由样本算出来的区间端点 \(\underline{\theta},\overline{\theta}\)。因此“置信水平 \(1-\alpha\)”的含义是：\textbf{在重复抽样中，这种构造方法给出的区间有 \(1-\alpha\) 的比例能覆盖真值}，而不是“对当前这一个已经算出来的区间，\(\theta\) 落在其中的概率是 \(1-\alpha\)”。
\end{note}

\begin{note}
\textbf{自学追问：} 为什么每道区间题都先找枢轴量？因为要先构造一个“分布已知、又含参数”的量，再把不等式反解成参数区间。换句话说，置信区间的计算本质上是“先标准化，再反解”。
\end{note}

在实际应用中，有时我们只关心参数的下界或上界（例如零件的寿命越长越好，废品率越低越好），此时就需要使用单侧置信区间。

\begin{definition}[单侧置信区间] \label{def:one_sided_confidence_interval}
对于给定值 $\alpha\ (0<\alpha<1)$，若由样本确定的统计量 $\underline{\theta}$ 满足 $P\{\theta > \underline{\theta}\} \ge 1-\alpha$，则称 $(\underline{\theta},+\infty)$ 是 $\theta$ 的置信水平为 $1-\alpha$ 的\textbf{单侧置信区间}，$\underline{\theta}$ 称为\textbf{单侧置信下限}。

同理，若由样本确定的统计量 $\overline{\theta}$ 满足 $P\{\theta < \overline{\theta}\} \ge 1-\alpha$，则称 $(-\infty,\overline{\theta})$ 是 $\theta$ 的置信水平为 $1-\alpha$ 的\textbf{单侧置信区间}，$\overline{\theta}$ 称为\textbf{单侧置信上限}。
\end{definition}

\begin{note}
若总体 $X$ 为连续型随机变量，我们通常要求上述概率严格等于 $1-\alpha$；若 $X$ 为离散型随机变量，则使其尽可能地接近（且不小于） $1-\alpha$。单侧区间适用于“只关心下界”或“只关心上界”的场景，例如寿命至少多大、次品率至多多大。不要把单侧区间中的 \(\alpha\) 仍按双侧区间那样平均分到两端。
\end{note}

\begin{note}
\textbf{自学追问：} 什么时候该选下限，什么时候该选上限？若题目关心“参数至少有多大”“不能太小”，就构造单侧下限；若关心“参数至多有多大”“不能太大”，就构造单侧上限。方向由实际问题中的“至少/至多、下不低于/上不超过”这些关键词决定。
\end{note}

\section{构造置信区间的基本原理}

构造置信区间最常用的方法是\textbf{枢轴量法}。枢轴量是一个由样本和未知参数共同构成的函数，其最大的特点是\textbf{它的分布完全已知，不依赖于任何未知参数}。下面我们以正态分布均值的区间估计为例，说明枢轴量法的基本步骤。

\begin{note}
\textbf{易错点：} 枢轴量不是“把未知参数消掉后的纯样本函数”，而是“虽然还含参数，但它的分布已经与未知参数无关”的量。若你构造出来的统计量里，分布形式仍然要靠未知参数才能写清楚，那它就还不能直接拿来反解置信区间。
\end{note}

利用上述枢轴量法，我们推导出了现实中最常见的正态总体（包括单个和两个正态总体）在不同情况下的参数置信区间。为了便于查阅，我们将相关公式归纳为表~\ref{tab:confidence_intervals}。

\begin{table}[htbp]
\centering
\caption{常见的正态总体参数置信区间（置信度为 $1-\alpha$）}
\label{tab:confidence_intervals}
\resizebox{\textwidth}{!}{
\begin{tabular}{lllc}
\toprule
\textbf{总体条件} & \textbf{待估参数} & \textbf{枢轴量及分布} & \textbf{置信区间} \\
\midrule
\multicolumn{4}{c}{\textbf{单个正态总体 $X \sim N(\mu, \sigma^2)$}} \\
\midrule
$\sigma^2$ 已知 & $\mu$ & $\frac{\overline{X}-\mu}{\sigma/\sqrt{n}} \sim N(0,1)$ & $\left(\overline{X}-\frac{\sigma}{\sqrt{n}}z_{\alpha/2},\ \overline{X}+\frac{\sigma}{\sqrt{n}}z_{\alpha/2}\right)$ \\
$\sigma^2$ 未知 & $\mu$ & $\frac{\overline{X}-\mu}{S/\sqrt{n}} \sim t(n-1)$ & $\left(\overline{X}-\frac{S}{\sqrt{n}}t_{\alpha/2}(n-1),\ \overline{X}+\frac{S}{\sqrt{n}}t_{\alpha/2}(n-1)\right)$ \\
$\mu$ 未知 & $\sigma^2$ & $\frac{(n-1)S^2}{\sigma^2} \sim \chi^2(n-1)$ & $\left(\frac{(n-1)S^2}{\chi^2_{\alpha/2}(n-1)},\ \frac{(n-1)S^2}{\chi^2_{1-\alpha/2}(n-1)}\right)$ \\
$\mu$ 已知 & $\sigma^2$ & $\frac{\sum(X_i-\mu)^2}{\sigma^2} \sim \chi^2(n)$ & $\left(\frac{\sum(X_i-\mu)^2}{\chi^2_{\alpha/2}(n)},\ \frac{\sum(X_i-\mu)^2}{\chi^2_{1-\alpha/2}(n)}\right)$ \\
\midrule
\multicolumn{4}{c}{\textbf{两个正态总体 $X \sim N(\mu_1, \sigma_1^2), Y \sim N(\mu_2, \sigma_2^2)$}} \\
\midrule
$\sigma_1^2,\sigma_2^2$ 已知 & $\mu_1-\mu_2$ & 正态分布 $N(0,1)$ & $\left((\overline{X}-\overline{Y}) \pm z_{\alpha/2}\sqrt{\frac{\sigma_1^2}{n_1}+\frac{\sigma_2^2}{n_2}}\right)$ \\
$\sigma_1^2=\sigma_2^2$ 未知 & $\mu_1-\mu_2$ & $t$ 分布 $t(n_1+n_2-2)$ & $\left((\overline{X}-\overline{Y}) \pm t_{\alpha/2}(n_1+n_2-2)S_w\sqrt{\frac{1}{n_1}+\frac{1}{n_2}}\right)$ \\
$\mu_1,\mu_2$ 已知 & $\frac{\sigma_1^2}{\sigma_2^2}$ & $F$ 分布 $F(n_1,n_2)$ & $\left(\frac{\sum(X_i-\mu_1)^2/n_1}{\sum(Y_j-\mu_2)^2/n_2}\frac{1}{F_{\alpha/2}(n_1,n_2)},\ \frac{\sum(X_i-\mu_1)^2/n_1}{\sum(Y_j-\mu_2)^2/n_2} F_{\alpha/2}(n_2,n_1)\right)$ \\
$\mu_1,\mu_2$ 未知 & $\frac{\sigma_1^2}{\sigma_2^2}$ & $F$ 分布 $F(n_1-1,n_2-1)$ & $\left(\frac{S_1^2}{S_2^2}\frac{1}{F_{\alpha/2}(n_1-1,n_2-1)},\ \frac{S_1^2}{S_2^2}F_{\alpha/2}(n_2-1,n_1-1)\right)$ \\
\bottomrule
\end{tabular}
}
\end{table}

\begin{remark}
上述表格中公式的推导核心在于”寻找正确的枢轴量”，因此读者不应当采用死记硬背的方式来记忆公式，而是应当通过枢轴量和分位数的定义，自行推导得出。
\end{remark}

\begin{note}
\textbf{枢轴量选取的决策流程（做题三问）：}
\begin{enumerate}
    \item \textbf{估什么？} 确定待估参数是 $\mu$ 还是 $\sigma^2$（或两总体的差/比）。
    \item \textbf{知什么？} 看题目条件中哪些参数已知、哪些未知：
    \begin{itemize}
        \item 估 $\mu$ 且 $\sigma^2$ 已知 $\Longrightarrow$ 用 $Z=\dfrac{\bar X-\mu}{\sigma/\sqrt{n}}\sim N(0,1)$；
        \item 估 $\mu$ 且 $\sigma^2$ 未知 $\Longrightarrow$ 用 $T=\dfrac{\bar X-\mu}{S/\sqrt{n}}\sim t(n-1)$；
        \item 估 $\sigma^2$ 且 $\mu$ 未知 $\Longrightarrow$ 用 $\chi^2=\dfrac{(n-1)S^2}{\sigma^2}\sim\chi^2(n-1)$；
        \item 估 $\sigma^2$ 且 $\mu$ 已知 $\Longrightarrow$ 用 $\chi^2=\dfrac{\sum(X_i-\mu)^2}{\sigma^2}\sim\chi^2(n)$。
    \end{itemize}
    \item \textbf{反解。} 对枢轴量写出 $P(a\le W\le b)=1-\alpha$，再把不等式解成”参数 $\in$ (下界, 上界)”的形式。
\end{enumerate}
记住口诀：\textbf{估谁含谁，知啥用啥，分布定型后反解}。
\end{note}

\begin{exercise}[正态总体均值的区间估计] \label{ex:pivotal_quantity}
设总体 $X\sim N(\mu,\sigma^2)$，其中方差 $\sigma^2$ 已知，均值 $\mu$ 未知。现有一组来自 $X$ 的样本 $X_1,X_2,\dots,X_n$，求 $\mu$ 的置信水平为 $1-\alpha$ 的置信区间。
\end{exercise}

\begin{solution}
这道题是枢轴量法最标准的起点。思路分三步：先找枢轴量，再写出中心概率式，最后把不等式解成关于 \(\mu\) 的区间。

第一步，构造关于 \(\mu\) 的枢轴量。因为总体正态且 \(\sigma^2\) 已知，所以
\[
\overline{X}\sim N\!\left(\mu,\frac{\sigma^2}{n}\right).
\]
把它标准化可得
\[
W=\frac{\overline{X}-\mu}{\sigma/\sqrt{n}}\sim N(0,1).
\]
这里 \(W\) 的分布已经完全已知，不含未知参数，因此它就是所需的枢轴量。

第二步，根据置信水平 \(1-\alpha\) 写出等尾概率式：
\[
P\left\{-z_{\alpha/2}<\frac{\overline{X}-\mu}{\sigma/\sqrt{n}}<z_{\alpha/2}\right\}=1-\alpha.
\]

第三步，把中间的不等式对 \(\mu\) 进行变形。两边同乘 \(\sigma/\sqrt{n}\)，再移项，得到
\[
P\left\{\overline{X}-\frac{\sigma}{\sqrt{n}}z_{\alpha/2}<\mu<\overline{X}+\frac{\sigma}{\sqrt{n}}z_{\alpha/2}\right\}=1-\alpha.
\]

因此，\(\mu\) 的置信水平为 \(1-\alpha\) 的置信区间为
\[
\left(\overline{X}-\frac{\sigma}{\sqrt{n}}z_{\alpha/2},\ \overline{X}+\frac{\sigma}{\sqrt{n}}z_{\alpha/2}\right).
\]

\begin{note}
本题之所以用 \(z_{\alpha/2}\)，关键在于\(\sigma\) 已知且采用双侧等尾区间。以后若 \(\sigma\) 未知，就不能继续使用这个正态枢轴量，而要改用 \(t\) 分布。
\end{note}
\end{solution}

除了上述 $\sigma^2$ 已知时 $\mu$ 的区间估计外，我们还可以利用类似的方法推导其他常见参数的置信区间。

\begin{exercise}[正态总体均值的区间估计（方差未知）] \label{ex:ci_mu_unknown_var}
设总体 $X\sim N(\mu,\sigma^2)$，其中均值 $\mu$ 和方差 $\sigma^2$ 均未知。现有一组来自 $X$ 的独立样本 $X_1,X_2,\dots,X_n$，样本均值为 $\overline{X}$，样本方差为 $S^2$。求 $\mu$ 的置信水平为 $1-\alpha$ 的置信区间。
\end{exercise}

\begin{solution}
与上一题相比，唯一的关键变化是：\(\sigma^2\) 现在未知，所以不能再直接用
\[
\frac{\overline{X}-\mu}{\sigma/\sqrt{n}}.
\]
这时要改用 \(t\) 统计量。

因为总体正态，所以有经典结论
\[
W=\frac{\overline{X}-\mu}{S/\sqrt{n}}\sim t(n-1).
\]
这里之所以能用 \(t\) 分布，靠的是“总体正态”以及“\(\overline{X}\) 与 \(S^2\) 独立”这两个前提。

在置信水平 \(1-\alpha\) 下，取双侧等尾区间：
\[
P\left\{-t_{\alpha/2}(n-1)<\frac{\overline{X}-\mu}{S/\sqrt{n}}<t_{\alpha/2}(n-1)\right\}=1-\alpha.
\]

把不等式解出 \(\mu\)，得到
\[
P\left\{\overline{X}-\frac{S}{\sqrt{n}}t_{\alpha/2}(n-1)<\mu<\overline{X}+\frac{S}{\sqrt{n}}t_{\alpha/2}(n-1)\right\}=1-\alpha.
\]

因此 \(\mu\) 的置信区间为
\[
\left(\overline{X}-\frac{S}{\sqrt{n}}t_{\alpha/2}(n-1),\ \overline{X}+\frac{S}{\sqrt{n}}t_{\alpha/2}(n-1)\right).
\]

\begin{note}
这道题最容易错的地方有两个：一是把 \(z_{\alpha/2}\) 误写成 \(t_{\alpha/2}(n-1)\) 之外的正态分位点；二是把自由度写成 \(n\)。只要样本方差 \(S^2\) 是围绕 \(\overline{X}\) 构造的，\(t\) 分布自由度就是 \(n-1\)。
\end{note}
\end{solution}

\begin{exercise}[正态总体方差的区间估计（均值未知）] \label{ex:ci_var_unknown_mu}
设总体 $X\sim N(\mu,\sigma^2)$，其中均值 $\mu$ 和方差 $\sigma^2$ 均未知。现有一组来自 $X$ 的独立样本 $X_1,X_2,\dots,X_n$，样本方差为 $S^2$。求方差 $\sigma^2$ 的置信水平为 $1-\alpha$ 的置信区间。
\end{exercise}

\begin{solution}
这一题要估计的是方差 \(\sigma^2\)，所以应优先寻找与样本方差有关的枢轴量。

在总体正态且总体均值未知的条件下，有
\[
W=\frac{(n-1)S^2}{\sigma^2}\sim\chi^2(n-1).
\]
这个统计量的分布不再含未知参数，因此可以直接用来构造区间。

由于 \(\chi^2\) 分布不是对称分布，双侧区间不能像正态和 \(t\) 分布那样写成“正负同一个分位点”的形式，而应分别取左右尾各 \(\alpha/2\)：
\[
P\left\{\chi^2_{1-\alpha/2}(n-1)<\frac{(n-1)S^2}{\sigma^2}<\chi^2_{\alpha/2}(n-1)\right\}=1-\alpha.
\]

现在解这个不等式。注意中间量与 \(\sigma^2\) 在分母上，变形时要格外小心方向：
\[
P\left\{\frac{(n-1)S^2}{\chi^2_{\alpha/2}(n-1)}<\sigma^2<\frac{(n-1)S^2}{\chi^2_{1-\alpha/2}(n-1)}\right\}=1-\alpha.
\]

因此，\(\sigma^2\) 的置信区间为
\[
\left(\frac{(n-1)S^2}{\chi^2_{\alpha/2}(n-1)},\ \frac{(n-1)S^2}{\chi^2_{1-\alpha/2}(n-1)}\right).
\]

\begin{note}
本题最容易错在上下界次序。因为 \(\chi^2_{\alpha/2}(n-1)\) 比 \(\chi^2_{1-\alpha/2}(n-1)\) 大，所以它反而出现在分母较大的那一侧，最终形成较小的下界和较大的上界。若机械照搬“左分位写左边、右分位写右边”，就很容易写反。
\end{note}
\end{solution}

对于两个正态总体 $X \sim N(\mu_1, \sigma_1^2)$ 和 $Y \sim N(\mu_2, \sigma_2^2)$，我们同样可以利用类似的方法构造枢轴量，来估计它们的均值差或方差比。

\begin{exercise}[两正态总体均值差的区间估计（方差相等且未知）] \label{ex:ci_diff_mu_unknown_var}
设有两个独立的正态总体 $X \sim N(\mu_1, \sigma^2)$ 和 $Y \sim N(\mu_2, \sigma^2)$，方差相等但具体数值未知。分别抽取容量为 $n_1$ 和 $n_2$ 的独立样本，求均值差 $\mu_1-\mu_2$ 的置信水平为 $1-\alpha$ 的置信区间。
\end{exercise}

\begin{solution}
我们要估计的是均值差 \(\mu_1-\mu_2\)，最自然的点估计量是
\[
\overline{X}-\overline{Y}.
\]

题目特别强调“两总体方差相等但未知”，这时不能分别用 \(S_1,S_2\) 生搬硬套，而要先构造\textbf{合并样本方差}
\[
S_w^2=\frac{(n_1-1)S_1^2+(n_2-1)S_2^2}{n_1+n_2-2}.
\]

在两个总体正态、两组样本独立且 \(\sigma_1^2=\sigma_2^2\) 的前提下，有
\[
W=\frac{(\overline{X}-\overline{Y})-(\mu_1-\mu_2)}{S_w\sqrt{\frac{1}{n_1}+\frac{1}{n_2}}}\sim t(n_1+n_2-2).
\]

于是按双侧等尾原则写出
\[
P\left\{-t_{\alpha/2}(n_1+n_2-2)<\frac{(\overline{X}-\overline{Y})-(\mu_1-\mu_2)}{S_w\sqrt{\frac{1}{n_1}+\frac{1}{n_2}}}<t_{\alpha/2}(n_1+n_2-2)\right\}=1-\alpha.
\]

解出 \(\mu_1-\mu_2\)，便得到区间
\[
\left((\overline{X}-\overline{Y})-t_{\alpha/2}(n_1+n_2-2)S_w\sqrt{\frac{1}{n_1}+\frac{1}{n_2}},\ 
(\overline{X}-\overline{Y})+t_{\alpha/2}(n_1+n_2-2)S_w\sqrt{\frac{1}{n_1}+\frac{1}{n_2}}\right).
\]

\begin{note}
这里最关键的适用条件不是“方差未知”，而是“\textbf{两总体方差相等且未知}”。若方差不相等，就不能直接用这一套合并方差 \(t\) 公式。
\end{note}
\end{solution}

\begin{exercise}[两正态总体方差比的区间估计（均值未知）] \label{ex:ci_ratio_var_unknown_mu}
设有两个独立的正态总体 $X \sim N(\mu_1, \sigma_1^2)$ 和 $Y \sim N(\mu_2, \sigma_2^2)$，其均值均未知。求方差比 $\sigma_1^2/\sigma_2^2$ 的置信水平为 $1-\alpha$ 的置信区间。
\end{exercise}

\begin{solution}
要估计方差比 \(\sigma_1^2/\sigma_2^2\)，最自然的样本对应物是 \(S_1^2/S_2^2\)。在两个总体正态且均值未知的条件下，有
\[
\frac{S_1^2/\sigma_1^2}{S_2^2/\sigma_2^2}
=\frac{S_1^2}{S_2^2}\cdot\frac{\sigma_2^2}{\sigma_1^2}
\sim F(n_1-1,n_2-1).
\]
这就是所需的枢轴量。

在置信水平 \(1-\alpha\) 下，写成双侧形式：
\[
P\left\{F_{1-\alpha/2}(n_1-1,n_2-1)<\frac{S_1^2}{S_2^2}\cdot\frac{\sigma_2^2}{\sigma_1^2}<F_{\alpha/2}(n_1-1,n_2-1)\right\}=1-\alpha.
\]

接下来把它整理成关于 \(\sigma_1^2/\sigma_2^2\) 的不等式。变形后得到
\[
P\left\{\frac{S_1^2}{S_2^2}\cdot\frac{1}{F_{\alpha/2}(n_1-1,n_2-1)}
<\frac{\sigma_1^2}{\sigma_2^2}
<\frac{S_1^2}{S_2^2}\cdot\frac{1}{F_{1-\alpha/2}(n_1-1,n_2-1)}\right\}=1-\alpha.
\]

再利用 \(F\) 分布分位点的倒数关系
\[
F_{1-\alpha/2}(n_1-1,n_2-1)=\frac{1}{F_{\alpha/2}(n_2-1,n_1-1)},
\]
可把上界写得更常用一些。于是最终区间为
\[
\left(\frac{S_1^2}{S_2^2}\cdot\frac{1}{F_{\alpha/2}(n_1-1,n_2-1)},\ 
\frac{S_1^2}{S_2^2}\cdot F_{\alpha/2}(n_2-1,n_1-1)\right).
\]

\begin{note}
这题常错在两个地方：一是把自由度误写成 \(n_1,n_2\)；二是忘记 \(F\) 分布左右分位点并不对称，需要借助倒数关系来整理上界。
\end{note}
\end{solution}

\begin{exercise}[正态总体方差的区间估计（均值已知）] \label{ex:ci_var_known_mu}
设总体 $X\sim N(\mu,\sigma^2)$，其中均值 $\mu$ 已知，方差 $\sigma^2$ 未知。现有一组来自 $X$ 的独立样本 $X_1,X_2,\dots,X_n$。求方差 $\sigma^2$ 的置信水平为 $1-\alpha$ 的置信区间。
\end{exercise}

\begin{solution}
由于均值 $\mu$ 已知，我们在构造关于方差 $\sigma^2$ 的枢轴量时，不应再使用样本均值 $\overline{X}$ 带来的自由度损失，而是直接使用真实均值 $\mu$。

首先，对每个样本进行标准化处理。因为 $X_i \sim N(\mu, \sigma^2)$，所以：
\begin{equation*}
\frac{X_i-\mu}{\sigma} \sim N(0,1), \quad i=1,2,\dots,n
\end{equation*}

根据卡方分布的定义，有限个相互独立的标准正态随机变量的平方和服从 $\chi^2$ 分布。将上述 $n$ 个相互独立的随机变量平方后求和，得到我们所需的枢轴量 $W$：
\begin{equation*}
W = \sum_{i=1}^n \left(\frac{X_i-\mu}{\sigma}\right)^2 = \frac{\sum_{i=1}^n (X_i-\mu)^2}{\sigma^2} \sim \chi^2(n)
\end{equation*}
此时 $W$ 服从自由度为 $n$ 的卡方分布（注意：因为使用的是已知的总体均值 $\mu$ 而不是样本均值 $\overline{X}$，所以自由度没有减少，仍为 $n$），且 $W$ 的分布不依赖于未知参数 $\sigma^2$，是一个理想的枢轴量。

利用 $\chi^2(n)$ 分布的上 $\alpha/2$ 分位点 $\chi^2_{\alpha/2}(n)$ 和上 $1-\alpha/2$ 分位点 $\chi^2_{1-\alpha/2}(n)$，在等尾概率的情况下有：
\begin{equation*}
P\left\{\chi^2_{1-\alpha/2}(n) < \frac{\sum_{i=1}^n (X_i-\mu)^2}{\sigma^2} < \chi^2_{\alpha/2}(n)\right\}=1-\alpha
\end{equation*}

接下来从不等式中解出 $\sigma^2$。对各项取倒数（此时不等号反向）：
\begin{equation*}
P\left\{\frac{1}{\chi^2_{\alpha/2}(n)} < \frac{\sigma^2}{\sum_{i=1}^n (X_i-\mu)^2} < \frac{1}{\chi^2_{1-\alpha/2}(n)}\right\}=1-\alpha
\end{equation*}

同时乘以 $\sum_{i=1}^n (X_i-\mu)^2$，得到：
\begin{equation*}
P\left\{\frac{\sum_{i=1}^n (X_i-\mu)^2}{\chi^2_{\alpha/2}(n)} < \sigma^2 < \frac{\sum_{i=1}^n (X_i-\mu)^2}{\chi^2_{1-\alpha/2}(n)}\right\}=1-\alpha
\end{equation*}

由此得到方差 $\sigma^2$ 的置信水平为 $1-\alpha$ 的置信区间为：
\begin{equation*}
\left(\frac{\sum_{i=1}^n(X_i-\mu)^2}{\chi^2_{\alpha/2}(n)},\ \frac{\sum_{i=1}^n(X_i-\mu)^2}{\chi^2_{1-\alpha/2}(n)}\right)
\end{equation*}
\begin{note}
与“\(\mu\) 未知”时相比，这里自由度是 \(n\) 而不是 \(n-1\)。原因不在于题型换了，而在于我们没有再用样本去估计总体均值，因此没有损失那一个自由度。
\end{note}
\end{solution}


\begin{exercise}[两正态总体均值差的区间估计（方差已知）] \label{ex:ci_diff_mu_known_var}
设有两个独立的正态总体 $X \sim N(\mu_1, \sigma_1^2)$ 和 $Y \sim N(\mu_2, \sigma_2^2)$，其中方差 $\sigma_1^2, \sigma_2^2$ 均已知，但均值 $\mu_1, \mu_2$ 未知。分别抽取容量为 $n_1$ 和 $n_2$ 的独立样本。求均值差 $\mu_1-\mu_2$ 的置信水平为 $1-\alpha$ 的置信区间。
\end{exercise}

\begin{solution}
我们已知样本均值的分布分别为：
\begin{equation*}
\overline{X} \sim N\left(\mu_1, \frac{\sigma_1^2}{n_1}\right), \quad \overline{Y} \sim N\left(\mu_2, \frac{\sigma_2^2}{n_2}\right)
\end{equation*}

要估计均值差 $\mu_1-\mu_2$，最自然的无偏估计量是两样本均值之差 $\overline{X}-\overline{Y}$。由于两样本相互独立，根据正态分布的可加性，$\overline{X}-\overline{Y}$ 也服从正态分布，其数学期望和方差分别为：
\begin{align*}
E(\overline{X}-\overline{Y}) &= E(\overline{X}) - E(\overline{Y}) = \mu_1 - \mu_2 \\
D(\overline{X}-\overline{Y}) &= D(\overline{X}) + D(\overline{Y}) = \frac{\sigma_1^2}{n_1} + \frac{\sigma_2^2}{n_2}
\end{align*}
因此有：
\begin{equation*}
\overline{X}-\overline{Y} \sim N\left(\mu_1-\mu_2,\ \frac{\sigma_1^2}{n_1} + \frac{\sigma_2^2}{n_2}\right)
\end{equation*}

对其进行标准化处理，便可构造出完全不含未知参数的枢轴量 $W$：
\begin{equation*}
W = \frac{(\overline{X}-\overline{Y}) - (\mu_1-\mu_2)}{\sqrt{\frac{\sigma_1^2}{n_1} + \frac{\sigma_2^2}{n_2}}} \sim N(0,1)
\end{equation*}

利用标准正态分布的对称性，以及上 $\alpha/2$ 分位点 $z_{\alpha/2}$，可以建立概率等式：
\begin{equation*}
P\left\{-z_{\alpha/2} < \frac{(\overline{X}-\overline{Y}) - (\mu_1-\mu_2)}{\sqrt{\frac{\sigma_1^2}{n_1} + \frac{\sigma_2^2}{n_2}}} < z_{\alpha/2}\right\} = 1-\alpha
\end{equation*}

解此不等式，分离出待估参数 $\mu_1-\mu_2$：
\begin{equation*}
-z_{\alpha/2}\sqrt{\frac{\sigma_1^2}{n_1} + \frac{\sigma_2^2}{n_2}} < (\overline{X}-\overline{Y}) - (\mu_1-\mu_2) < z_{\alpha/2}\sqrt{\frac{\sigma_1^2}{n_1} + \frac{\sigma_2^2}{n_2}}
\end{equation*}
同时乘以 $-1$ 并加上 $\overline{X}-\overline{Y}$，由于对称性，正负边界的形式依然保持：
\begin{equation*}
P\left\{(\overline{X}-\overline{Y}) - z_{\alpha/2}\sqrt{\frac{\sigma_1^2}{n_1} + \frac{\sigma_2^2}{n_2}} < \mu_1-\mu_2 < (\overline{X}-\overline{Y}) + z_{\alpha/2}\sqrt{\frac{\sigma_1^2}{n_1} + \frac{\sigma_2^2}{n_2}}\right\} = 1-\alpha
\end{equation*}

故 $\mu_1-\mu_2$ 的置信水平为 $1-\alpha$ 的置信区间为：
\begin{equation*}
\left((\overline{X}-\overline{Y}) - z_{\alpha/2}\sqrt{\frac{\sigma_1^2}{n_1}+\frac{\sigma_2^2}{n_2}},\ (\overline{X}-\overline{Y}) + z_{\alpha/2}\sqrt{\frac{\sigma_1^2}{n_1}+\frac{\sigma_2^2}{n_2}}\right)
\end{equation*}
\end{solution}


\begin{exercise}[两正态总体方差比的区间估计（均值已知）] \label{ex:ci_ratio_var_known_mu}
设有两个独立的正态总体 $X \sim N(\mu_1, \sigma_1^2)$ 和 $Y \sim N(\mu_2, \sigma_2^2)$，其中均值 $\mu_1, \mu_2$ 均已知，但方差 $\sigma_1^2, \sigma_2^2$ 未知。分别抽取容量为 $n_1$ 和 $n_2$ 的独立样本。求方差比 $\sigma_1^2/\sigma_2^2$ 的置信水平为 $1-\alpha$ 的置信区间。
\end{exercise}

\begin{solution}
由于总体的真实均值 $\mu_1, \mu_2$ 均已知，我们可以充分利用它们构造 $\chi^2$ 统计量，而无需借助于消耗自由度的样本方差 $S^2$。

根据前面的推导，对于总体 $X$ 和总体 $Y$，我们分别有：
\begin{equation*}
U = \frac{\sum_{i=1}^{n_1}(X_i-\mu_1)^2}{\sigma_1^2} \sim \chi^2(n_1) \quad \text{和} \quad V = \frac{\sum_{j=1}^{n_2}(Y_j-\mu_2)^2}{\sigma_2^2} \sim \chi^2(n_2)
\end{equation*}

根据 $F$ 分布的定义，由两个独立的 $\chi^2$ 分布除以它们各自的自由度后求比值，即可得到服从 $F$ 分布的随机变量。因此，我们构造如下枢轴量：
\begin{equation*}
W = \frac{U/n_1}{V/n_2} = \frac{\frac{1}{n_1\sigma_1^2}\sum_{i=1}^{n_1}(X_i-\mu_1)^2}{\frac{1}{n_2\sigma_2^2}\sum_{j=1}^{n_2}(Y_j-\mu_2)^2} \sim F(n_1, n_2)
\end{equation*}
为了方便书写和解不等式，我们将 $W$ 重写为含有未知参数 $\sigma_1^2/\sigma_2^2$ 和纯样本数值的形式：
\begin{equation*}
W = \frac{\frac{1}{n_1}\sum_{i=1}^{n_1}(X_i-\mu_1)^2}{\frac{1}{n_2}\sum_{j=1}^{n_2}(Y_j-\mu_2)^2} \cdot \frac{\sigma_2^2}{\sigma_1^2} \sim F(n_1, n_2)
\end{equation*}
此时 $W$ 的分布完全确定，不含任何未知参数，是一个合适的枢轴量。

利用 $F$ 分布的分位点，我们有：
\begin{equation*}
P\left\{F_{1-\alpha/2}(n_1, n_2) < \frac{\frac{1}{n_1}\sum_{i=1}^{n_1}(X_i-\mu_1)^2}{\frac{1}{n_2}\sum_{j=1}^{n_2}(Y_j-\mu_2)^2} \cdot \frac{\sigma_2^2}{\sigma_1^2} < F_{\alpha/2}(n_1, n_2)\right\}=1-\alpha
\end{equation*}

对上述不等式的每一项取倒数（中间项的 $\sigma_2^2/\sigma_1^2$ 变为 $\sigma_1^2/\sigma_2^2$，且不等号方向反向）：
\begin{equation*}
P\left\{\frac{1}{F_{\alpha/2}(n_1, n_2)} < \frac{\frac{1}{n_2}\sum_{j=1}^{n_2}(Y_j-\mu_2)^2}{\frac{1}{n_1}\sum_{i=1}^{n_1}(X_i-\mu_1)^2} \cdot \frac{\sigma_1^2}{\sigma_2^2} < \frac{1}{F_{1-\alpha/2}(n_1, n_2)}\right\}=1-\alpha
\end{equation*}

同时乘以 $\frac{\frac{1}{n_1}\sum_{i=1}^{n_1}(X_i-\mu_1)^2}{\frac{1}{n_2}\sum_{j=1}^{n_2}(Y_j-\mu_2)^2}$ 使得中间只留下待估的方差比：
\begin{equation*}
P\left\{\frac{\frac{1}{n_1}\sum(X_i-\mu_1)^2}{\frac{1}{n_2}\sum(Y_j-\mu_2)^2} \cdot \frac{1}{F_{\alpha/2}(n_1, n_2)} < \frac{\sigma_1^2}{\sigma_2^2} < \frac{\frac{1}{n_1}\sum(X_i-\mu_1)^2}{\frac{1}{n_2}\sum(Y_j-\mu_2)^2} \cdot \frac{1}{F_{1-\alpha/2}(n_1, n_2)}\right\}=1-\alpha
\end{equation*}

最后，利用 $F$ 分布的倒数性质 $F_{1-\alpha/2}(n_1, n_2) = \frac{1}{F_{\alpha/2}(n_2, n_1)}$，可以将上限的右侧转化为更直观的形式。最终得到 $\sigma_1^2/\sigma_2^2$ 的置信区间为：
\begin{equation*}
\left(\frac{\frac{1}{n_1}\sum_{i=1}^{n_1}(X_i-\mu_1)^2}{\frac{1}{n_2}\sum_{j=1}^{n_2}(Y_j-\mu_2)^2}\frac{1}{F_{\alpha/2}(n_1,n_2)},\ \frac{\frac{1}{n_1}\sum_{i=1}^{n_1}(X_i-\mu_1)^2}{\frac{1}{n_2}\sum_{j=1}^{n_2}(Y_j-\mu_2)^2} F_{\alpha/2}(n_2,n_1)\right)
\end{equation*}
\end{solution}


\section{典型例题}

\begin{exercise} \label{ex:ci_length}
设总体 $X \sim N(\mu,\sigma^2)$（$\sigma^2$ 已知），则在给定样本容量 $n$ 及置信度 $1-\alpha$ 的情况下，未知参数 $\mu$ 的置信区间长度随着样本均值 $\overline{X}$ 的增加而（\qquad）。
\begin{enumerate}
  \item[(A)] 增加
  \item[(B)] 减少
  \item[(C)] 不变
  \item[(D)] 不能确定增或减
\end{enumerate}
\end{exercise}

\begin{solution}
应选 (C)。

在给定条件下，由表~\ref{tab:confidence_intervals} 可知，未知参数 $\mu$ 的置信区间为：
\begin{equation*}
\left( \overline{X} - \frac{\sigma}{\sqrt{n}} z_{\alpha/2},\ \overline{X} + \frac{\sigma}{\sqrt{n}} z_{\alpha/2} \right)
\end{equation*}
该置信区间的长度为
\[
L = 2\frac{\sigma}{\sqrt{n}} z_{\alpha/2}.
\]
显然，区间长度 \(L\) 只与总体标准差 \(\sigma\)、样本容量 \(n\) 以及对应置信水平的分位点 \(z_{\alpha/2}\) 有关，而与样本均值 \(\overline{X}\) 毫无关系。

也就是说，\(\overline{X}\) 的变化只会让整个区间左右平移，不会改变区间的宽度。因此，随着 \(\overline{X}\) 的变化，置信区间的长度保持不变。
\end{solution}

\begin{exercise} \label{ex:ci_large_sample}
设总体 $X$ 的方差为 $1$，根据来自 $X$ 的容量为 $100$ 的简单随机样本，测得样本均值为 $5$。已知标准正态分布函数值 $\Phi(1.96)=0.975$。求 $X$ 的数学期望 $\mu$ 的置信度近似等于 $0.95$ 的置信区间。
\end{exercise}

\begin{solution}
题目中并未说明总体 $X$ 是否为正态分布，但是由于样本容量 $n=100 \ge 30$，属于大样本。根据中心极限定理，样本均值 $\overline{X}$ 近似服从正态分布，即：
\begin{equation*}
Z = \frac{\overline{X}-\mu}{\sigma/\sqrt{n}} = \frac{\overline{X}-\mu}{1/\sqrt{100}} \stackrel{\text{近似}}{\sim} N(0,1)
\end{equation*}
又因为 $\Phi(1.96)=0.975$，可知 $z_{0.025} = 1.96$，故 $P\{-1.96 < Z < 1.96\} \approx 0.95$。

由此可得 $\mu$ 的近似置信区间为：
\begin{equation*}
\left( \overline{X} - \frac{1}{10} \times 1.96,\ \overline{X} + \frac{1}{10} \times 1.96 \right)
\end{equation*}
代入已知数据 $\overline{X}=5$，计算出置信区间的下限为 $4.804$，上限为 $5.196$。

所以 $\mu$ 的置信区间近似为 $(4.804,\ 5.196)$。
\end{solution}


\begin{exercise} \label{prob:ci_1}
对于正态总体 $X \sim N(\mu, \sigma^2)$，$\sigma^2$ \textbf{已知}，记 $L$ 是 $\mu$ 的置信度为 $0.95$ 的置信区间的长度，那么当样本容量 $n$ \underline{\qquad\qquad} 时，$L \leq \sigma$。\\
（参考数据：$z_{0.05}=1.645$，$z_{0.025}=1.96$，$t_{0.025}(15)=1.753$，$t_{0.025}(16)=2.12$）
\end{exercise}
\begin{solution}
因为 $\sigma^2$ 已知，使用正态分布作为枢轴量。置信区间长度为：
\begin{equation*}
L = 2 \cdot \frac{\sigma}{\sqrt{n}} z_{0.025}
\end{equation*}
代入 $z_{0.025}=1.96$，并令 $L \leq \sigma$：
\begin{equation*}
2 \cdot 1.96 \cdot \frac{\sigma}{\sqrt{n}} \leq \sigma \implies \frac{3.92}{\sqrt{n}} \leq 1 \implies \sqrt{n} \geq 3.92 \implies n \geq 15.3664
\end{equation*}
由于样本容量 $n$ 必须为正整数，故 $n \ge 16$。
\end{solution}

\begin{exercise} \label{prob:ci_2}
生产一个零件所需时间 $X \sim N(\mu, \sigma^2)$。现观察了 $25$ 个零件，测得样本均值 $\overline{X}=5.5$，样本标准差 $S=1.73$。求 $\mu$ 的置信度为 $0.95$ 的置信区间 \underline{\qquad\qquad}。\\
（参考数据：$t_{0.05}(24)=1.7109,\ t_{0.025}(24)=2.0639,\ \Phi(1.96)=0.975$）
\end{exercise}
\begin{solution}
由于总体方差 $\sigma^2$ 未知，需要使用 $t$ 分布。自由度 $df = n-1 = 24$，查表得 $t_{0.025}(24)=2.0639$。
置信区间公式为：
\begin{equation*}
\left( \overline{X}-\frac{S}{\sqrt{n}}t_{0.025}(24),\ \overline{X}+\frac{S}{\sqrt{n}}t_{0.025}(24) \right)
\end{equation*}
计算边际误差：
\begin{equation*}
\frac{1.73}{5} \times 2.0639 \approx 0.7141
\end{equation*}
下限为 $5.5 - 0.7141 = 4.7859$，上限为 $5.5 + 0.7141 = 6.2141$。故置信区间为 $(4.7859,\ 6.2141)$。
\end{solution}

\begin{exercise} \label{prob:ci_3}
设总体 $X \sim N(\mu, \sigma^2)$，有一个容量为 $10$ 的样本，测得样本方差 $S^2=2$。求 $\sigma^2$ 的置信水平为 $0.95$ 的置信区间 \underline{\qquad\qquad}。\\
（参考数据：$\chi^2_{0.975}(9)=2.70$，$\chi^2_{0.025}(9)=19.023$）
\end{exercise}
\begin{solution}
由于 $\mu$ 未知，求 $\sigma^2$ 的置信区间使用 $\chi^2$ 分布，自由度 $df = 10-1 = 9$。
方差的置信区间公式为：
\begin{equation*}
\left( \frac{(n-1)S^2}{\chi^2_{0.025}(9)},\ \frac{(n-1)S^2}{\chi^2_{0.975}(9)} \right)
\end{equation*}
计算分子：$(n-1)S^2 = 9 \times 2 = 18$。代入分母得到：
\begin{equation*}
\frac{18}{19.023} \approx 0.9462,\quad \frac{18}{2.70} \approx 6.6667
\end{equation*}
故置信区间为 $(0.9462,\ 6.6667)$。
\end{solution}

\begin{exercise} \label{prob:ci_4}
已知 $X\sim N(\mu,\sigma^2)$，样本容量 $n=9$，样本均值 $\overline{X}=6$，求 $\mu$ 在置信水平 $0.95$ 下的单侧置信下限 \underline{\qquad\qquad}。\\
（参考数据：$\Phi(1.96)=0.975,\ \Phi(1.645)=0.95$）
\end{exercise}
\begin{solution}
题目没有给出 \(\sigma\) 的具体数值，因此结果只能保留含 \(\sigma\) 的形式。又因为参考数据给的是标准正态分位点，所以这里对应的是“\(\sigma\) 已知时 \(\mu\) 的单侧置信区间”。

单侧置信下限的要求是
\[
P\{\mu \ge \underline{\mu}\}=0.95.
\]
利用标准正态分布的枢轴量，有：
\begin{equation*}
P\left( \frac{\overline{X}-\mu}{\sigma/\sqrt{n}} \le z_{0.05} \right) = 0.95
\end{equation*}
解得 $\mu \ge \overline{X} - z_{0.05}\frac{\sigma}{\sqrt{n}}$。代入数据 $\overline{X}=6, n=9, z_{0.05}=1.645$：
\begin{equation*}
\underline{\mu} = 6 - 1.645 \frac{\sigma}{3} \approx 6 - 0.5483\sigma
\end{equation*}
\end{solution}

\begin{exercise} \label{prob:ci_5}
对 $\alpha\in(0,1)$，由样本 $X_1,\dots,X_n$ 得到统计量 $\underline{\theta}$，若对一切参数 $\theta$ 均有 \underline{\qquad\qquad}，则称 $(\underline{\theta},+\infty)$ 为 $\theta$ 的置信水平为 $1-\alpha$ 的单侧置信区间，$\underline{\theta}$ 称为单侧置信下限。
\end{exercise}
\begin{solution}
应填
\[
P\{\theta>\underline{\theta}\}\ge 1-\alpha.
\]

理由直接来自单侧置信下限的定义。区间 \((\underline{\theta},+\infty)\) 要成为参数 \(\theta\) 的置信水平为 \(1-\alpha\) 的单侧置信区间，意思就是：由样本算出的下限 \(\underline{\theta}\) 不应把真值“压到区间外面去”，并且这种覆盖成功的概率至少要达到 \(1-\alpha\)。

换成事件语言，就是
\[
P\{\theta>\underline{\theta}\}\ge 1-\alpha.
\]
如果教材专门讨论连续型情形，也常把它写成等号
\[
P\{\theta>\underline{\theta}\}=1-\alpha,
\]
但从定义本身看，填空的标准形式应为“\(\ge 1-\alpha\)”。
\end{solution}



\begin{exercise} \label{prob:ci_6}
设总体 $X\sim N(\mu,\sigma^2)$，有一个容量 $n=9$ 的样本，已知 $\sum_{i=1}^9 X_i=198$，$\sum_{i=1}^9 X_i^2=4372$。求 $\mu$ 与 $\sigma$ 的置信水平为 $0.95$ 的双侧置信区间（结果保留四位小数）。\\
（参考数据：$t_{0.025}(8)=2.306$，$\chi^2_{0.975}(8)=2.180$，$\chi^2_{0.025}(8)=17.535$）
\end{exercise}
\begin{solution}
首先计算样本均值和样本方差：
\begin{equation*}
\overline{X} = \frac{1}{n}\sum_{i=1}^9 X_i = \frac{198}{9} = 22
\end{equation*}
\begin{equation*}
S^2 = \frac{1}{n-1}\left(\sum_{i=1}^9 X_i^2 - n\overline{X}^2\right) = \frac{1}{8}\left(4372 - 9 \times 22^2\right) = \frac{16}{8} = 2
\end{equation*}
于是样本标准差 $S = \sqrt{2} \approx 1.4142$。

求 $\mu$ 的置信区间（方差未知，用 $t$ 分布）：
自由度 $df=8$，分位点 $t_{0.025}(8)=2.306$。边际误差为 $\frac{\sqrt{2}}{3} \times 2.306 \approx 1.0871$。
下限：$22 - 1.0871 = 20.9129$，上限：$22 + 1.0871 = 23.0871$。
即 $\mu$ 的置信区间为 $(20.9129,\ 23.0871)$。

求 $\sigma$ 的置信区间：
先求出 $\sigma^2$ 的置信区间：
\begin{equation*}
\left( \frac{(n-1)S^2}{\chi^2_{0.025}(8)},\ \frac{(n-1)S^2}{\chi^2_{0.975}(8)} \right) = \left( \frac{16}{17.535},\ \frac{16}{2.180} \right) \approx (0.9125,\ 7.3395)
\end{equation*}
对上下限开平方，得到标准差 $\sigma$ 的置信区间：
\begin{equation*}
(\sqrt{0.9125},\ \sqrt{7.3395}) \approx (0.9552,\ 2.7091)
\end{equation*}
综上，$\mu$ 的置信区间为 $(20.9129,\ 23.0871)$，$\sigma$ 的置信区间为 $(0.9552,\ 2.7091)$。
\end{solution}

\begin{note}
\textbf{自学检查点——区间估计：}
\begin{enumerate}
    \item 置信水平 $1-\alpha=0.95$ 的正确含义是什么？（答：重复抽样 100 次，大约 95 次构造出的区间能覆盖真值）
    \item 估计正态总体均值 $\mu$，$\sigma$ 已知用什么分布？未知呢？（答：已知用 $Z$，未知用 $t(n-1)$）
    \item 估计正态总体方差 $\sigma^2$ 用什么分布？（答：$\chi^2(n-1)$）
    \item 置信区间越宽，置信水平越高还是越低？（答：越高。宽区间更容易覆盖真值）
    \item 增大样本量 $n$ 对置信区间有什么影响？（答：区间变窄，精度提高）
\end{enumerate}
\end{note}

\begin{note}
\textbf{【高频题型模板】区间估计大题的标准解法：}
\begin{enumerate}
    \item \textbf{判断题型：} 估 $\mu$ 还是 $\sigma^2$？$\sigma$ 已知还是未知？单总体还是双总体？
    \item \textbf{选枢轴量：} 根据判断结果，从表中选对应的枢轴量（$Z$/$t$/$\chi^2$/$F$）。
    \item \textbf{写概率式：} $P\{a\leq W\leq b\}=1-\alpha$，查表得分位数 $a,b$。
    \item \textbf{反解参数：} 把不等式变形为"参数 $\in$ (下界, 上界)"。
    \item \textbf{代入数据：} 算出 $\bar x$、$s^2$ 等样本统计量，代入得数值区间。
\end{enumerate}
\textbf{注意：} 对 $\chi^2$ 和 $F$ 分布（非对称），双侧区间的上下分位数不对称：用 $\chi^2_{\alpha/2}$ 和 $\chi^2_{1-\alpha/2}$，不能像正态那样只查一个值。
\end{note}


\chapter{假设检验}

\begin{note}
\textbf{自学导读：区间估计的"反面"——用数据做决策。}

区间估计回答"参数大概在哪个范围"；假设检验回答"参数是否等于某个特定值"。两者用的数学工具完全相同（同样的枢轴量、同样的分布），只是问法不同。

\textbf{核心逻辑：}
\begin{enumerate}
    \item 先假设 $H_0$（原假设）为真——比如"$\mu=\mu_0$"；
    \item 在 $H_0$ 为真的前提下，算出检验统计量应该落在什么范围（接受域）；
    \item 如果实际算出的统计量落在接受域\textbf{之外}（拒绝域），说明"如果 $H_0$ 为真，出现这么极端的数据概率 $<\alpha$"——小概率事件发生了，所以拒绝 $H_0$。
\end{enumerate}

\textbf{自学提醒：}
\begin{itemize}
    \item 假设检验和区间估计是\textbf{一体两面}：$\mu_0$ 落在置信区间内 $\iff$ 不拒绝 $H_0:\mu=\mu_0$。
    \item 选统计量的规则和第 8 章完全一样：$\sigma$ 已知用 $Z$，未知用 $t$，检验方差用 $\chi^2$，比较两组方差用 $F$。
    \item "不拒绝 $H_0$"$\neq$"$H_0$ 为真"。只是说数据不够极端，无法推翻 $H_0$。
\end{itemize}
\end{note}

{\small
\[
\begin{cases}
\text{假设检验}
\begin{cases}
\text{基本思想: 随机试验中小概率事件一般不会发生} \\
\text{基本步骤: 提出假设} \Rightarrow \text{构建统计量} \Rightarrow \text{写出拒绝域} \Rightarrow \text{检验} \\
\text{主要检验: 双边检验或单边检验} \\
\text{显著性检验: 仅考虑控制第一类错误的假设检验，限定第一类错误发生的最大概率为} \alpha
\end{cases} \\
\text{单个正态总体}
\begin{cases}
\text{总体均值}
\begin{cases}
\text{总体方差已知, } Z\text{检验，也称} U\text{检验} \\
\text{总体方差未知, } t\text{检验}
\end{cases} \\
\text{总体方差}
\begin{cases}
\text{总体均值已知, } \chi^2(n)\text{检验} \\
\text{总体均值未知, } \chi^2(n-1)\text{检验}
\end{cases}
\end{cases} \\
\text{两个正态总体} \\
(X \sim N(\mu_1,\sigma_1^2), Y \sim N(\mu_2,\sigma_2^2))
\begin{cases}
\text{总体均值差} (\mu_1-\mu_2)
\begin{cases}
\sigma_1^2=\sigma_2^2=\sigma^2\text{已知, } Z\text{检验} \\
\sigma_1^2=\sigma_2^2=\sigma^2\text{未知, } t(n_1+n_2-2)\text{检验}
\end{cases} \\
\text{总体方差比} \left( \frac{\sigma_1^2}{\sigma_2^2} \right)
\begin{cases}
\mu_1,\mu_2\text{已知, } F(n_1,n_2)\text{检验} \\
\mu_1,\mu_2\text{未知, } F(n_1-1,n_2-1)\text{检验}
\end{cases}
\end{cases}
\end{cases}
\]
}

\begin{introduction}
    \item 原假设与备择假设~\ref{def:hypothesis-test}
    \item 两类错误与显著性水平~\ref{def:two-errors}
    \item 假设检验的基本步骤~\ref{sec:ht-steps}
    \item 正态总体参数的假设检验~\ref{sec:ht-normal}
    \item 置信区间与假设检验的关系~\ref{sec:ht-ci}
\end{introduction}

\begin{note}
\textbf{使用提醒：} 假设检验选方法时，不要先背拒绝域，先按四个问题定位：检验哪个参数、单样本还是双样本、总体方差已知还是未知、备择假设是左尾右尾还是双尾。前两问决定该用哪类统计量，第三问常常决定是 \(z\) 还是 \(t\)、是 \(\chi^2\) 还是 \(F\)，最后一问才决定拒绝域落在哪个尾部。
\end{note}

\section{基本概念}

\begin{definition}[假设检验] \label{def:hypothesis-test}
关于总体分布中的未知参数或分布类型的每一种论断称为\textbf{统计假设}。根据样本信息去推断该假设是否成立的统计推断问题称为\textbf{假设检验}。其中需要检验的假设称为\textbf{原假设}（或零假设），记为 \(H_0\)；与 \(H_0\) 对立的假设称为\textbf{备择假设}，记为 \(H_1\)。
\end{definition}

\begin{note}
假设检验的本质不是“证明 \(H_0\) 真”，而是“根据样本判断是否有足够证据拒绝 \(H_0\) ”。因此最后的结论通常只能说“拒绝 \(H_0\)”或“不能拒绝 \(H_0\)”，而不宜表述成“接受 \(H_0\) 一定正确”。
\end{note}

\begin{note}
\textbf{自学追问：} 为什么总要先设 \(H_0\)？因为检验需要在 \(H_0\) 下给出可计算的分布，从而控制第一类错误。结论也要只写“拒绝”或“不能拒绝”，不要把“不能拒绝”误写成“证明了 \(H_0\) 真”。
\end{note}

\begin{definition}[两类错误] \label{def:two-errors}
\begin{enumerate}
    \item \textbf{第一类错误（弃真）}：\(H_0\) 为真时，错误地拒绝了 \(H_0\)。
    \[
    \alpha = P\{\text{拒绝}\, H_0 \mid H_0\text{ 为真}\}.
    \]
    \item \textbf{第二类错误（取伪）}：\(H_0\) 不真时，错误地未拒绝 \(H_0\)。
    \[
    \beta = P\{\text{未拒绝}\, H_0 \mid H_1\text{ 为真}\}.
    \]
\end{enumerate}
\end{definition}

\begin{note}
两类错误不能同时都压到很小而又不付出代价。样本量固定时，如果把拒绝域设得更苛刻以减小第一类错误 \(\alpha\)，往往会增大第二类错误 \(\beta\)。所以假设检验本质上是在不同错误风险之间做权衡。
\end{note}

\begin{definition}[显著性水平与显著性检验] \label{def:significance-level}
在假设检验中，只控制犯第一类错误的概率不超过 \(\alpha\)（\(\alpha\) 通常取 0.1, 0.05, 0.01），而不考虑第二类错误的概率，这种检验称为\textbf{显著性检验}，\(\alpha\) 称为\textbf{显著性水平}。
\end{definition}

\begin{note}
\(\alpha\) 不是“原假设为真的概率”，也不是“犯任何错误的总概率”。它专指：\textbf{在 \(H_0\) 真的前提下，却把 \(H_0\) 错误拒绝掉的最大容许概率}。做题时只要看到“显著性水平”，就要立刻联想到第一类错误。
\end{note}

\begin{note}
\textbf{使用提醒：} 把 \(\alpha\) 取得更小，意味着检验更谨慎、更不容易拒绝 \(H_0\)；这通常会降低第一类错误，却常常提高第二类错误。它不意味着”结论更真”，只意味着”拒绝 \(H_0\) 的门槛更高”。
\end{note}

\begin{note}
\textbf{$\alpha$ 与 $\beta$ 的跷跷板关系（直觉图景）：}

想象两条钟形曲线：一条以 $\mu_0$（$H_0$ 真时）为中心，另一条以 $\mu_1$（$H_1$ 真时）为中心，两条曲线有重叠区域。拒绝域是一条竖线（临界值）右边的区域。
\begin{itemize}
    \item 把竖线\textbf{右移}（$\alpha\downarrow$，更难拒绝）$\Longrightarrow$ 左边曲线尾巴变小（$\alpha$ 减小），但右边曲线被截掉的面积变大（$\beta$ 增大）。
    \item 把竖线\textbf{左移}（$\alpha\uparrow$，更容易拒绝）$\Longrightarrow$ $\alpha$ 增大，$\beta$ 减小。
\end{itemize}
\textbf{唯一能同时减小 $\alpha$ 和 $\beta$ 的方法：增大样本量 $n$}——这会让两条曲线都变窄（方差 $\sigma^2/n$ 减小），重叠区域缩小。

\textbf{考试常考判断题：}
\begin{itemize}
    \item “$\alpha$ 越小检验越好” $\times$（$\beta$ 可能很大）
    \item “增大 $n$ 可同时减小 $\alpha$ 和 $\beta$” \checkmark
    \item “$\alpha+\beta=1$” $\times$（两者没有这种简单关系）
\end{itemize}
\end{note}

\begin{note}
\textbf{小概率原理}是假设检验的理论基础：概率很小的事件在一次试验中认为不会发生。若在 \(H_0\) 成立的条件下，样本观测值导致某个概率 \(\leq \alpha\) 的事件发生，则拒绝 \(H_0\)。
\end{note}

\begin{definition}[双边检验与单边检验]\label{def:ht-types}
\begin{itemize}
    \item \textbf{双边检验}：\(H_0\colon \mu = \mu_0\)，\(H_1\colon \mu \neq \mu_0\)。
    \item \textbf{右边检验}：\(H_0\colon \mu \leq \mu_0\)，\(H_1\colon \mu > \mu_0\)。
    \item \textbf{左边检验}：\(H_0\colon \mu \geq \mu_0\)，\(H_1\colon \mu < \mu_0\)。
\end{itemize}
\end{definition}

\begin{note}
原假设 \(H_0\) 与备择假设 \(H_1\) 的地位不对等：\(H_0\) 是受保护的（控制弃真概率），因此选择 \(H_0\) 需要谨慎，一般遵循以下原则：
\begin{enumerate}
    \item 将后果较严重的错误设为第一类错误。例如检验药品真假，\(H_0\colon\) 药品为假，\(H_1\colon\) 药品为真（误判假药为真的后果更严重）。
    \item 若无严重后果，常选 \(H_0\) 为"维持现状"。例如检验新技术是否提高收益，\(H_0\colon\) 新技术未提高收益。
    \item 用\textbf{反证法}思想：若要验证 \(x\) 显著高于 \(y\)，则设 \(H_0\colon x \leq y\)，拒绝 \(H_0\) 后便有充分理由认为 \(x > y\)。
\end{enumerate}
\textbf{检验类型由备择假设决定}：若 \(H_1\) 写”\(\neq\)”就是双边检验，写”\(>\)”就是右边检验，写”\(<\)”就是左边检验。不要只盯着 \(H_0\) 看方向。
\end{note}

\begin{note}
\textbf{考试中如何快速确定 $H_0$ 和 $H_1$（三步法）：}
\begin{enumerate}
    \item \textbf{找”想证明的结论”。} 题目中”是否显著提高/降低/不同”所指向的方向，就是 $H_1$。因为假设检验的逻辑是反证法——把想证的放在 $H_1$，拒绝 $H_0$ 后才能下结论。
    \item \textbf{$H_0$ 取等号。} $H_0$ 永远包含”$=$”（$=$、$\le$、$\ge$），因为只有在等号成立时才能精确算出统计量的分布。
    \item \textbf{定方向。}
    \begin{itemize}
        \item 题目说”是否有显著差异/变化” $\Longrightarrow$ 双边：$H_0\colon\theta=\theta_0$，$H_1\colon\theta\neq\theta_0$；
        \item 题目说”是否显著提高/大于” $\Longrightarrow$ 右边：$H_0\colon\theta\le\theta_0$，$H_1\colon\theta>\theta_0$；
        \item 题目说”是否显著降低/小于” $\Longrightarrow$ 左边：$H_0\colon\theta\ge\theta_0$，$H_1\colon\theta<\theta_0$。
    \end{itemize}
\end{enumerate}
\textbf{常见陷阱：} 题目说”检验均值是否为 $\mu_0$”，很多同学误以为 $H_0\colon\mu\neq\mu_0$。记住：$H_0$ 永远是”没变化”的那一方，$H_1$ 才是”有变化”。
\end{note}

\section{假设检验的步骤} \label{sec:ht-steps}

\begin{property}[假设检验的四步流程]\label{prop:ht-steps}
\begin{enumerate}
    \item 根据问题提出 \(H_0\) 和 \(H_1\)，确定检验类型（双边/单边）。
    \item 给定显著性水平 \(\alpha\) 和样本容量 \(n\)。
    \item 选择\textbf{检验统计量}，确定\textbf{拒绝域}。
    \item 由样本计算统计量的观测值，判断是否落入拒绝域，作出结论。
\end{enumerate}
\end{property}

\begin{note}
\textbf{拒绝域的形式与备择假设 \(H_1\) 一致}，便于记忆：
\begin{itemize}
    \item \(H_1\colon \theta > \theta_0\)（右边检验）\(\Longrightarrow\) 拒绝域在\textbf{右尾}；
    \item \(H_1\colon \theta < \theta_0\)（左边检验）\(\Longrightarrow\) 拒绝域在\textbf{左尾}；
    \item \(H_1\colon \theta \neq \theta_0\)（双边检验）\(\Longrightarrow\) 拒绝域在\textbf{双尾}。
\end{itemize}
不能拒绝 \(H_0\) 并不意味着 \(H_0\) 一定为真，只是没有足够证据拒绝它。
\end{note}

\begin{note}
\textbf{假设检验与区间估计的统一视角：} 两者用的数学工具完全相同——同样的枢轴量、同样的分布表。区别仅在于"问法"：
\begin{itemize}
    \item 区间估计问"参数在哪个范围？" $\Longrightarrow$ 反解不等式得区间；
    \item 假设检验问"参数是否等于某值？" $\Longrightarrow$ 把该值代入枢轴量，看观测值是否落入小概率区域。
\end{itemize}
\textbf{快速判断法：} 若 $\mu_0$ 落在 $\mu$ 的 $1-\alpha$ 置信区间内，则在显著性水平 $\alpha$ 下不拒绝 $H_0\colon\mu=\mu_0$；反之拒绝。这是两章内容的桥梁。
\end{note}


\section{正态总体参数的假设检验} \label{sec:ht-normal}

\begin{exercise}[检验统计量的表达式] \label{ex:ht_stat}
设 \((X_1,\dots,X_n)\) 是来自 \(N(\mu,\sigma^2)\) 的样本（\(\mu,\sigma^2\) 未知），\(\overline{X}=\dfrac{1}{n}\displaystyle\sum_{i=1}^{n}X_i\)，\(Q^2=\displaystyle\sum_{i=1}^{n}(X_i-\overline{X})^2\)。检验 \(H_0\colon\mu=0\) 时，应选取的检验统计量为\underline{\qquad}（用 \(\overline{X},Q^2\) 表示）。
\end{exercise}
\begin{solution}
\textbf{先判题型。} 总体是正态总体，但 \(\sigma^2\) 未知，而题目检验的是均值 \(\mu\)，因此这里不能用 \(z\) 检验，而应使用单样本 \(t\) 检验。

在原假设 \(H_0\colon \mu=\mu_0\) 下，标准统计量应写成
\[
T=\frac{\overline{X}-\mu_0}{S/\sqrt{n}}.
\]
本题中 \(\mu_0=0\)，又由
\[
S^2=\frac{Q^2}{n-1},\qquad S=\frac{Q}{\sqrt{n-1}},
\]
代入即可把统计量改写成题目要求的 \(\overline{X},Q^2\) 形式：
\[
T=\frac{\overline{X}-\mu_0}{S/\sqrt{n}}\bigg|_{\mu_0=0}
=\frac{\overline{X}}{Q/\sqrt{n(n-1)}}=\frac{\overline{X}\sqrt{n(n-1)}}{Q}\sim t(n-1).
\]
\[
\text{故应选取的检验统计量为}\quad T=\frac{\overline{X}\sqrt{n(n-1)}}{Q}.
\]
\end{solution}


以下汇总了单正态总体与双正态总体在不同条件下的常用检验统计量及其拒绝域。
\begin{table}[htbp]
\centering
\caption{正态总体均值与方差的假设检验（显著性水平为 $\alpha$）} \label{tab:hypothesis_tests}
\resizebox{\textwidth}{!}{
\begin{tabular}{lllll}
\toprule
\textbf{参数} & \textbf{条件} & \textbf{原假设 $H_0$} & \textbf{检验统计量} & \textbf{拒绝域（对应右单边、左单边、双边）} \\
\midrule
\multirow{3}{*}{$\mu$} & \multirow{3}{*}{$\sigma^2$ 已知} & $\mu \leq \mu_0$ ($H_1: \mu > \mu_0$) & \multirow{3}{*}{$Z = \frac{\overline{X}-\mu_0}{\sigma/\sqrt{n}} \sim N(0,1)$} & $z \geq z_\alpha$ \\
& & $\mu \geq \mu_0$ ($H_1: \mu < \mu_0$) & & $z \leq -z_\alpha$ \\
& & $\mu = \mu_0$ ($H_1: \mu \neq \mu_0$) & & $|z| \geq z_{\frac{\alpha}{2}}$ \\
\midrule
\multirow{3}{*}{$\mu$} & \multirow{3}{*}{$\sigma^2$ 未知} & $\mu \leq \mu_0$ ($H_1: \mu > \mu_0$) & \multirow{3}{*}{$t = \frac{\overline{X}-\mu_0}{S/\sqrt{n}} \sim t(n-1)$} & $t \geq t_\alpha(n-1)$ \\
& & $\mu \geq \mu_0$ ($H_1: \mu < \mu_0$) & & $t \leq -t_\alpha(n-1)$ \\
& & $\mu = \mu_0$ ($H_1: \mu \neq \mu_0$) & & $|t| \geq t_{\frac{\alpha}{2}}(n-1)$ \\
\midrule
\multirow{3}{*}{$\mu_1-\mu_2$} & \multirow{3}{*}{$\sigma_1^2,\sigma_2^2$ 已知} & $\mu_1-\mu_2 \leq \delta$ & \multirow{3}{*}{$Z = \frac{(\overline{X}-\overline{Y})-\delta}{\sqrt{\frac{\sigma_1^2}{n_1}+\frac{\sigma_2^2}{n_2}}} \sim N(0,1)$} & $z \geq z_\alpha$ \\
& & $\mu_1-\mu_2 \geq \delta$ & & $z \leq -z_\alpha$ \\
& & $\mu_1-\mu_2 = \delta$ & & $|z| \geq z_{\frac{\alpha}{2}}$ \\
\midrule
\multirow{3}{*}{$\mu_1-\mu_2$} & \multirow{3}{*}{$\sigma_1^2=\sigma_2^2$ 未知} & $\mu_1-\mu_2 \leq \delta$ & \multirow{3}{*}{$t = \frac{(\overline{X}-\overline{Y})-\delta}{S_w\sqrt{\frac{1}{n_1}+\frac{1}{n_2}}} \sim t(n_1+n_2-2)$} & $t \geq t_\alpha(n_1+n_2-2)$ \\
& & $\mu_1-\mu_2 \geq \delta$ & & $t \leq -t_\alpha(n_1+n_2-2)$ \\
& & $\mu_1-\mu_2 = \delta$ & & $|t| \geq t_{\frac{\alpha}{2}}(n_1+n_2-2)$ \\
\midrule
\multirow{3}{*}{$\sigma^2$} & \multirow{3}{*}{$\mu$ 未知} & $\sigma^2 \leq \sigma_0^2$ & \multirow{3}{*}{$\chi^2 = \frac{(n-1)S^2}{\sigma_0^2} \sim \chi^2(n-1)$} & $\chi^2 \geq \chi^2_\alpha(n-1)$ \\
& & $\sigma^2 \geq \sigma_0^2$ & & $\chi^2 \leq \chi^2_{1-\alpha}(n-1)$ \\
& & $\sigma^2 = \sigma_0^2$ & & $\chi^2 \geq \chi^2_{\frac{\alpha}{2}}(n-1)$ 或 $\chi^2 \leq \chi^2_{1-\frac{\alpha}{2}}(n-1)$ \\
\midrule
\multirow{3}{*}{$\sigma_1^2 / \sigma_2^2$} & \multirow{3}{*}{$\mu_1,\mu_2$ 未知} & $\sigma_1^2 \leq \sigma_2^2$ & \multirow{3}{*}{$F = \frac{S_1^2}{S_2^2} \sim F(n_1-1,n_2-1)$} & $F \geq F_\alpha(n_1-1,n_2-1)$ \\
& & $\sigma_1^2 \geq \sigma_2^2$ & & $F \leq F_{1-\alpha}(n_1-1,n_2-1)$ \\
& & $\sigma_1^2 = \sigma_2^2$ & & $F \geq F_{\frac{\alpha}{2}}(n_1-1,n_2-1)$ 或 $F \leq F_{1-\frac{\alpha}{2}}(n_1-1,n_2-1)$ \\
\bottomrule
\end{tabular}
}
\end{table}
\begin{remark} \label{rem:rejection_region}
\textbf{记忆捷径}：仔细观察上表可知，拒绝域的不等号方向总是与备择假设$H_1$的形式（不等号方向）保持一致。
\end{remark}

\begin{note}
\textbf{检验统计量选取的决策流程（与区间估计完全对称）：}
\begin{enumerate}
    \item \textbf{检什么？} 检验 $\mu$ 还是 $\sigma^2$（或两总体的差/比）。
    \item \textbf{知什么？} 与区间估计相同的判断：
    \begin{itemize}
        \item 检 $\mu$ 且 $\sigma^2$ 已知 $\Longrightarrow$ $Z$ 检验；
        \item 检 $\mu$ 且 $\sigma^2$ 未知 $\Longrightarrow$ $t$ 检验；
        \item 检 $\sigma^2$（$\mu$ 未知）$\Longrightarrow$ $\chi^2$ 检验；
        \item 检两总体方差比 $\Longrightarrow$ $F$ 检验。
    \end{itemize}
    \item \textbf{代 $H_0$ 的值。} 把 $H_0$ 中的参数值代入枢轴量，得到检验统计量。
    \item \textbf{看 $H_1$ 定拒绝域方向。}
\end{enumerate}
口诀：\textbf{区间估计反着来——代入 $\theta_0$ 后看落不落在尾巴上}。
\end{note}

\subsection{单个正态总体均值的检验}

\begin{property}[Z 检验（\(\sigma^2\) 已知）]\label{prop:z-test}
检验统计量 \(Z = \dfrac{\overline{X}-\mu_0}{\sigma/\sqrt{n}} \sim N(0,1)\)。

\begin{center}
\begin{tabular}{ccc}
\toprule
\(H_1\) & 拒绝域 \\
\midrule
\(\mu > \mu_0\) & \(z \geq z_\alpha\) \\[4pt]
\(\mu < \mu_0\) & \(z \leq -z_\alpha\) \\[4pt]
\(\mu \neq \mu_0\) & \(|z| \geq z_{\alpha/2}\) \\
\bottomrule
\end{tabular}
\end{center}
\end{property}

\begin{exercise}[辛烷等级的左边 \texorpdfstring{\(Z\)}{Z} 检验] \label{ex:ht_z_left}
一种燃料的辛烷等级服从正态分布，平均值为 98，标准差为 0.8。从一批新燃料中随机抽取 25 桶，样本均值为 97.7。假定新燃料辛烷等级标准差仍为 0.8，问新燃料辛烷等级平均值是否比原燃料偏低（\(\alpha=0.05\)）？
（\(z_{0.05}=1.64\)）
\end{exercise}
\begin{solution}
\(H_0\colon\mu\geq98\)，\(H_1\colon\mu<98\)（左边检验，\(\alpha=0.05\)）。
\(\sigma=0.8\) 已知，使用 \(Z\) 检验：
\[
Z=\frac{\overline{X}-98}{0.8/\sqrt{25}}\sim N(0,1),\quad\text{拒绝域}\ z\leq -z_{0.05}=-1.64.
\]
\[
z=\frac{97.7-98}{0.8/5}=\frac{-0.3}{0.16}=-1.875<-1.64.
\]
落入拒绝域，拒绝 \(H_0\)，认为新燃料辛烷等级显著偏低。
\end{solution}



\begin{exercise}\label{exer:9-z-test}
某工厂生产灯泡的寿命 \(X \sim N(\mu, 200^2)\)。过去灯泡平均寿命为 1500 小时，采用新工艺后抽取 25 只，测得平均寿命为 1675 小时。问在显著性水平 \(\alpha = 0.05\) 下，新工艺是否显著提高了灯泡寿命？（\(z_{0.05}=1.65\)）
\end{exercise}
\begin{solution}
\textbf{第一步：提出假设。}问题问"是否显著提高"，故用右边检验：
\[
H_0\colon \mu \leq 1500,\quad H_1\colon \mu > 1500.
\]
\textbf{第二步：选择检验统计量。}总体方差 \(\sigma^2=200^2\) 已知，使用 \(Z\) 检验：
\[
Z = \frac{\overline{X}-\mu_0}{\sigma/\sqrt{n}} = \frac{\overline{X}-1500}{200/\sqrt{25}} \sim N(0,1).
\]
\textbf{第三步：确定拒绝域。}右边检验，\(\alpha=0.05\)，拒绝域为 \(z \geq z_{0.05}=1.65\)。\\
\textbf{第四步：计算并判断。}将样本数据代入：
\[
z = \frac{1675-1500}{200/5} = \frac{175}{40} = 4.375.
\]
\(4.375 > 1.65\)，\(z\) 落入拒绝域，拒绝 \(H_0\)，即在 5\% 显著性水平下认为新工艺显著提高了灯泡寿命。
\end{solution}


\begin{property}[t 检验（\(\sigma^2\) 未知）]\label{prop:t-test}
检验统计量 \(t = \dfrac{\overline{X}-\mu_0}{S/\sqrt{n}} \sim t(n-1)\)。
\begin{center}
\begin{tabular}{ccc}
\toprule
\(H_1\) & 拒绝域 \\
\midrule
\(\mu > \mu_0\) & \(t \geq t_\alpha(n-1)\) \\[4pt]
\(\mu < \mu_0\) & \(t \leq -t_\alpha(n-1)\) \\[4pt]
\(\mu \neq \mu_0\) & \(|t| \geq t_{\alpha/2}(n-1)\) \\
\bottomrule
\end{tabular}
\end{center}
\end{property}


\begin{exercise}\label{exer:9-t-test-two}
某厂生产灯泡寿命 \(X \sim N(\mu, \sigma^2)\)（\(\mu, \sigma^2\) 均未知），抽取 25 只测得 \(\overline{x}=1700\)，\(s=490\)。在 \(\alpha=0.01\) 下，能否认为平均寿命为 2000 小时？
（\(\;t_{0.005}(24)=2.7969\)）
\end{exercise}
\begin{solution}
\textbf{提出假设。}问题问"能否认为平均寿命为 2000"，属于双边检验：
\[
H_0\colon \mu = 2000,\quad H_1\colon \mu \neq 2000.
\]
\textbf{选择检验统计量。}\(\sigma^2\) 未知，使用 \(t\) 检验：
\[
t = \frac{\overline{X}-\mu_0}{S/\sqrt{n}} \sim t(n-1) = t(24).
\]
\textbf{确定拒绝域。}双边检验，\(\alpha=0.01\)，拒绝域为 \(|t| \geq t_{\alpha/2}(24) = t_{0.005}(24)=2.7969\)。\\
\textbf{计算并判断。}已知 \(\overline{x}=1700,\ s=490,\ n=25\)：
\[
t = \frac{1700-2000}{490/\sqrt{25}} = \frac{-300}{98} = -3.061,\quad |t| = 3.061.
\]
\(|t| = 3.061 > 2.7969\)，落入拒绝域，拒绝 \(H_0\)，即在 1\% 显著性水平下不能认为平均寿命为 2000 小时。
\end{solution}


\begin{exercise}\label{exer:9-t-test-right}
某高校某专业 25 名学生高数成绩服从正态分布，样本均值 \(\overline{x}=69.5\)，标准差 \(s=15\)。在 \(\alpha=0.05\) 下，检验该专业平均成绩是否\textbf{显著高于} 70 分。（\(t_{0.05}(24)=1.711\)）
\end{exercise}
\begin{solution}
\textbf{提出假设。}问题问"是否\textbf{显著高于} 70 分"，故用右边检验：
\[
H_0\colon \mu \leq 70,\quad H_1\colon \mu > 70.
\]
\textbf{选择检验统计量。}\(\sigma^2\) 未知，使用 \(t\) 检验：
\[
t = \frac{\overline{X}-\mu_0}{S/\sqrt{n}} \sim t(24).
\]
\textbf{确定拒绝域。}右边检验，\(\alpha=0.05\)，拒绝域为 \(t \geq t_{0.05}(24)=1.711\)。\\
\textbf{计算并判断。}已知 \(\overline{x}=69.5,\ s=15,\ n=25\)：
\[
t = \frac{69.5-70}{15/\sqrt{25}} = \frac{-0.5}{3} = -0.167.
\]
\(-0.167 < 1.711\)，未落入拒绝域（实际上统计量为负，远小于正的临界值），故不能拒绝 \(H_0\)，即在 5\% 显著性水平下没有充分证据认为平均成绩显著高于 70 分。
\begin{note}
从直觉上看，样本均值 \(\overline{x}=69.5\) 反而\textbf{低于} 70 分，自然不可能得出"显著高于"的结论。这提醒我们：在使用单边检验时，应先看样本数据的方向是否与 \(H_1\) 一致；若方向相反，通常就已经说明样本没有支持 \(H_1\) 的趋势。
\end{note}
\end{solution}

\begin{exercise}[考试成绩均值的 \texorpdfstring{\(t\)}{t} 检验] \label{ex:ht_t_two}
设某次考试成绩服从正态分布，随机抽取 36 位考生，算得 \(\overline{x}=66.5\) 分，\(s^2=15^2\)。问在 \(\alpha=0.05\) 下，是否可以认为全体考生平均成绩为 70 分？（\(t_{0.025}(35)=2.0301\)）
\end{exercise}
\begin{solution}
\(H_0\colon\mu=70\)，\(H_1\colon\mu\neq70\)（双边 \(t\) 检验，\(\alpha=0.05\)）。
\[
t=\frac{\overline{X}-70}{S/\sqrt{n}}\sim t(35),\quad\text{拒绝域}\ |t|\geq t_{0.025}(35)=2.0301.
\]
\[
t=\frac{66.5-70}{15/\sqrt{36}}=\frac{-3.5}{2.5}=-1.4,\quad |t|=1.4<2.0301.
\]
未落入拒绝域，故不能拒绝 \(H_0\)；在该显著性水平下，没有充分证据认为平均成绩偏离 70 分。
\end{solution}

\begin{exercise}[电池寿命方差的双边 \texorpdfstring{\(\chi^2\)}{χ²} 检验] \label{ex:ht_chi2_battery}
某厂电池寿命 \(X\sim N(\mu,\sigma^2)\)，正常生产时 \(\sigma^2=5000\)。现因生产条件波动较大，抽取 26 只电池测得样本方差 \(s^2=7200\)。能否判断寿命波动性有显著变化（\(\alpha=0.02\)）？
（\(\chi^2_{0.01}(25)=44.314\)，\(\chi^2_{0.99}(25)=11.524\)）
\end{exercise}
\begin{solution}
\(H_0\colon\sigma^2=5000\)，\(H_1\colon\sigma^2\neq5000\)（双边 \(\chi^2\) 检验，\(\alpha=0.02\)）。
\[
\chi^2=\frac{(n-1)s^2}{\sigma_0^2}=\frac{25\times7200}{5000}=36.
\]
\textbf{不拒绝域}为
\[
[\,\chi^2_{0.99}(25),\ \chi^2_{0.01}(25)\,]=[11.524,\ 44.314].
\]
\(11.524<36<44.314\)，未落入拒绝域，故不能拒绝 \(H_0\)，即暂时没有充分证据判断波动性有显著变化。
\end{solution}


\begin{exercise}[机床零件标准差的 \texorpdfstring{\(\chi^2\)}{χ²} 检验] \label{ex:ht_chi2_two}
某机床加工零件尺寸 \(X\sim N(\mu,\sigma^2)\)，抽取 10 个零件测得长度为
\[
5.5,\ 6.0,\ 6.2,\ 5.8,\ 5.6,\ 6.0,\ 6.4,\ 6.4,\ 6.0,\ 6.1.
\]
能否认为标准差为 0.4（\(\alpha=0.05\)）？
（\(\chi^2_{0.025}(9)=19.023\)，\(\chi^2_{0.975}(9)=2.700\)）
\end{exercise}
\begin{solution}
\(H_0\colon\sigma^2=0.16\)，\(H_1\colon\sigma^2\neq0.16\)（双边 \(\chi^2\) 检验）。\\
\(\overline{x}=60.0/10=6.0\)，离差平方和：
\[
\sum_{i=1}^{10}(x_i-6.0)^2=0.25+0+0.04+0.04+0.16+0+0.16+0.16+0+0.01=0.82.
\]
\[
\chi^2=\frac{(n-1)s^2}{\sigma_0^2}=\frac{0.82}{0.16}=5.125.
\]
拒绝域：\(\chi^2\leq 2.700\) 或 \(\chi^2\geq 19.023\)。
\(2.700<5.125<19.023\)，未落入拒绝域，故不能拒绝 \(H_0\)；在该显著性水平下，没有充分证据认为标准差偏离 \(0.4\)。
\end{solution}


\subsection{两个正态总体均值差的检验}

\begin{property}[两个正态总体均值差的检验]\label{prop:two-sample-mean-test}
设 \(X_1,\dots,X_{n_1}\) 来自 \(N(\mu_1,\sigma_1^2)\)，\(Y_1,\dots,Y_{n_2}\) 来自 \(N(\mu_2,\sigma_2^2)\)，两组独立。\\
\textbf{情形1}：\(\sigma_1^2, \sigma_2^2\) 已知。检验统计量
\[
Z = \frac{(\overline{X}-\overline{Y})-\delta}{\sqrt{\sigma_1^2/n_1+\sigma_2^2/n_2}} \sim N(0,1).
\]
\textbf{情形2}：\(\sigma_1^2=\sigma_2^2=\sigma^2\) 未知。检验统计量
\[
t = \frac{(\overline{X}-\overline{Y})-\delta}{S_w\sqrt{1/n_1+1/n_2}} \sim t(n_1+n_2-2),
\]
其中 \(S_w^2 = \dfrac{(n_1-1)S_1^2+(n_2-1)S_2^2}{n_1+n_2-2}\) 为\textbf{合并样本方差}。\\
两种情形下，\(H_1\colon \mu_1-\mu_2 > \delta\) 对应右尾拒绝域，\(H_1\colon \mu_1-\mu_2 < \delta\) 对应左尾，\(H_1\colon \mu_1-\mu_2 \neq \delta\) 对应双尾。
\end{property}

\subsection{单个正态总体方差的检验}

\begin{property}[$\chi^2$ 检验]\label{prop:chi2-test}
\textbf{\(\mu\) 未知}时，检验统计量
\[
\chi^2 = \frac{(n-1)S^2}{\sigma_0^2} \sim \chi^2(n-1).
\]
\textbf{\(\mu\) 已知}时，检验统计量
\[
\chi^2 = \frac{\sum_{i=1}^{n}(X_i-\mu)^2}{\sigma_0^2} \sim \chi^2(n).
\]
\begin{center}
\begin{tabular}{cc}
\toprule
\(H_1\) & 拒绝域（\(\mu\) 未知） \\
\midrule
\(\sigma^2 > \sigma_0^2\) & \(\chi^2 \geq \chi^2_\alpha(n-1)\) \\[4pt]
\(\sigma^2 < \sigma_0^2\) & \(\chi^2 \leq \chi^2_{1-\alpha}(n-1)\) \\[4pt]
\(\sigma^2 \neq \sigma_0^2\) & \(\chi^2 \geq \chi^2_{\alpha/2}(n-1)\) 或 \(\chi^2 \leq \chi^2_{1-\alpha/2}(n-1)\) \\
\bottomrule
\end{tabular}
\end{center}
\end{property}

\begin{exercise}\label{exer:9-chi2-two}
某厂铜线折断力 \(X \sim N(\mu, \sigma^2)\)，抽查 10 根，\(\overline{x}=575.2\)，\(s^2=68.16\)。在 \(\alpha=0.05\) 下检验 \(H_0\colon \sigma^2=64\)，\(H_1\colon \sigma^2\neq64\)。
（\(\chi^2_{0.025}(9)=19.0\)，\(\chi^2_{0.975}(9)=2.7\)）
\end{exercise}
\begin{solution}
\textbf{提出假设。}检验方差是否为 64，使用双边检验：
\[
H_0\colon \sigma^2 = 64,\quad H_1\colon \sigma^2 \neq 64.
\]
\textbf{选择检验统计量。}\(\mu\) 未知，使用 \(\chi^2\) 检验：
\[
\chi^2 = \frac{(n-1)S^2}{\sigma_0^2} = \frac{(10-1)S^2}{64} \sim \chi^2(9).
\]
\textbf{确定拒绝域。}双边检验，\(\alpha=0.05\)，拒绝域为：
\[
\chi^2 \leq \chi^2_{0.975}(9) = 2.7 \quad\text{或}\quad \chi^2 \geq \chi^2_{0.025}(9) = 19.0.
\]
不拒绝域为 \((2.7,\ 19.0)\)。\\
\textbf{计算并判断。}
\[
\chi^2 = \frac{9 \times 68.16}{64} = \frac{613.44}{64} = 9.585.
\]
\(2.7 < 9.585 < 19.0\)，未落入拒绝域，故不能拒绝 \(H_0\)，即在 5\% 显著性水平下没有充分证据认为这批铜线折断力的方差偏离 64。
\end{solution}


\begin{exercise}\label{exer:9-chi2-left}
测定溶液水分，10 个测定值 \(s=0.037\%\)，总体服从正态分布。在 \(\alpha=0.05\) 下检验 \(H_0\colon \sigma \geq 0.04\%\)，\(H_1\colon \sigma < 0.04\%\)。
（\(\chi^2_{0.95}(9)=3.325\)）
\end{exercise}
\begin{solution}
\textbf{提出假设。}题目给出 \(H_0\colon \sigma \geq 0.04\%\)，需等价转化为关于 \(\sigma^2\) 的假设：
\[
H_0\colon \sigma^2 \geq (0.04\%)^2 = 1.6\times10^{-6},\quad H_1\colon \sigma^2 < (0.04\%)^2.
\]
这是\textbf{左边检验}。\\
\textbf{选择检验统计量。}\(\mu\) 未知，使用 \(\chi^2\) 检验：
\[
\chi^2 = \frac{(n-1)S^2}{\sigma_0^2} = \frac{9\,S^2}{(0.04\%)^2} \sim \chi^2(9).
\]
\textbf{确定拒绝域。}左边检验，\(\alpha=0.05\)，拒绝域为 \(\chi^2 \leq \chi^2_{1-\alpha}(9) = \chi^2_{0.95}(9) = 3.325\)。\\
\textbf{计算并判断。}已知 \(s=0.037\%\)：
\[
\chi^2 = \frac{9 \times (0.037\%)^2}{(0.04\%)^2} = 9 \times \left(\frac{0.037}{0.04}\right)^{\!2} = 9 \times 0.8556 = 7.701.
\]
\(7.701 > 3.325\)，未落入拒绝域，故不能拒绝 \(H_0\)，即在 5\% 显著性水平下没有充分证据认为 \(\sigma < 0.04\%\)。
\end{solution}



\subsection{两个正态总体方差比的检验}

\begin{property}[F 检验]\label{prop:f-test}
\(\mu_1, \mu_2\) 未知时，检验统计量
\[
F = \frac{S_1^2}{S_2^2} \sim F(n_1-1, n_2-1).
\]
\begin{center}
\begin{tabular}{cc}
\toprule
\(H_1\) & 拒绝域 \\
\midrule
\(\sigma_1^2 > \sigma_2^2\) & \(F \geq F_\alpha(n_1-1, n_2-1)\) \\[4pt]
\(\sigma_1^2 < \sigma_2^2\) & \(F \leq F_{1-\alpha}(n_1-1, n_2-1)\) \\[4pt]
\(\sigma_1^2 \neq \sigma_2^2\) & \(F \geq F_{\alpha/2}\) 或 \(F \leq F_{1-\alpha/2}\) \\
\bottomrule
\end{tabular}
\end{center}
其中 \(F_{1-\alpha}(n_1-1,n_2-1) = \dfrac{1}{F_\alpha(n_2-1,n_1-1)}\)。
\end{property}

\begin{exercise}\label{exer:9-f-test}
两个小麦品种从播种到抽穗天数如下：
\begin{center}
\begin{tabular}{ccccccccccc|c}
\toprule
品种1 & 101 & 100 & 99 & 99 & 98 & 100 & 98 & 99 & 99 & 99 & 行和 992 \\
\midrule
品种2 & 100 & 98 & 100 & 99 & 98 & 99 & 98 & 98 & 99 & 100 & 行和 989 \\
\bottomrule
\end{tabular}
\end{center}
在 \(\alpha=0.05\) 下：(1) 检验方差是否有显著差异；(2) 检验均值是否有显著差异。
（\(F_{0.025}(9,9)=4.03\)，\(F_{0.975}(9,9)=0.248\)，\(t_{0.025}(18)=2.101\)）
\end{exercise}
\begin{solution}
\textbf{准备工作}：\(n_1=n_2=10\)，\(\overline{x}_1=992/10=99.2\)，\(\overline{x}_2=989/10=98.9\)。
计算两组的离差平方和：
\[
\sum_{i=1}^{10}(x_{1i}-99.2)^2 = (101{-}99.2)^2\!\times\!1 + (100{-}99.2)^2\!\times\!2 + (99{-}99.2)^2\!\times\!5 + (98{-}99.2)^2\!\times\!2 = 7.6,
\]
\[
\sum_{j=1}^{10}(x_{2j}-98.9)^2 = (100{-}98.9)^2\!\times\!3 + (99{-}98.9)^2\!\times\!3 + (98{-}98.9)^2\!\times\!4 = 6.9.
\]
\[
S_1^2 = \frac{7.6}{9} \approx 0.844,\quad S_2^2 = \frac{6.9}{9} \approx 0.767.
\]
\textbf{(1) 检验方差是否有显著差异（F 检验）。}提出假设：\(H_0\colon \sigma_1^2=\sigma_2^2\)，\(H_1\colon \sigma_1^2\neq\sigma_2^2\)（双边检验）。
检验统计量 \(F = S_1^2/S_2^2 \sim F(9,9)\)（在 \(H_0\) 成立时）。拒绝域：\(F \geq F_{0.025}(9,9)=4.03\) 或 \(F \leq F_{0.975}(9,9)=0.248\)。
\[
F = \frac{0.844}{0.767} \approx 1.100.
\]
\(0.248 < 1.100 < 4.03\)，未落入拒绝域，故不能拒绝 \(H_0\)，即没有充分证据认为两品种方差存在显著差异。因此后续可以用等方差的 pooled \(t\) 检验。\\
\textbf{(2) 检验均值是否有显著差异（t 检验）。}提出假设：\(H_0\colon \mu_1=\mu_2\)，\(H_1\colon \mu_1\neq\mu_2\)（双边检验）。由(1)知方差相等，使用 pooled \(t\) 检验。先计算合并样本方差：
\[
S_w^2 = \frac{(n_1-1)S_1^2+(n_2-1)S_2^2}{n_1+n_2-2} = \frac{9\times0.844+9\times0.767}{18} = \frac{7.6+6.9}{18} = \frac{14.5}{18} \approx 0.806.
\]
检验统计量：
\[
t = \frac{\overline{X}_1-\overline{X}_2}{S_w\sqrt{1/n_1+1/n_2}} = \frac{99.2-98.9}{\sqrt{0.806\times(0.1+0.1)}} = \frac{0.3}{\sqrt{0.1612}} = \frac{0.3}{0.401} \approx 0.747.
\]
拒绝域：\(|t| \geq t_{0.025}(18)=2.101\)。\(|0.747| < 2.101\)，未落入拒绝域，故不能拒绝 \(H_0\)，即没有充分证据认为两品种均值存在显著差异。
\begin{note}
在实际应用中，检验两总体均值差时，通常\textbf{先}做 F 检验判断方差是否相等，再根据结果选择合适的 \(t\) 检验方法。
\end{note}
\end{solution}

\begin{exercise}[绿地面积的 \texorpdfstring{\(t\)}{t} 检验与 \texorpdfstring{\(\chi^2\)}{χ²} 检验] \label{ex:ht_combo}
为改建绿地，5 位学生独立测量面积（单位：km\(^2\)）：
\[
1.23,\ 1.22,\ 1.20,\ 1.26,\ 1.23.
\]
设测量误差服从正态分布（\(\alpha=0.05\)）\\
(1) 以前面积 \(\mu_0=1.23\) km\(^2\)，是否有必要修改？\\
(2) 若要求标准差不超过 \(\sigma=0.015\)，能否认为标准差显著偏大？
（\(t_{0.025}(4)=2.7764\)，\(\chi^2_{0.05}(4)=9.488\)）
\end{exercise}
\begin{solution}
\(\overline{x}=1.228\)，离差平方和 \(\sum(x_i-\overline{x})^2=0.00188\)，\(s^2=0.000470\)，\(s\approx0.0217\)。\\
(1) \(H_0\colon\mu=1.23\)，\(H_1\colon\mu\neq1.23\)（双边 \(t\) 检验）。
\[
t=\frac{1.228-1.23}{0.0217/\sqrt{5}}\approx-0.206,\quad |t|<2.7764.
\]
故不能拒绝 \(H_0\)，在该显著性水平下没有必要修改。\\
(2) \(H_0\colon\sigma\leq0.015\)，\(H_1\colon\sigma>0.015\)（右边 \(\chi^2\) 检验）。
\[
\chi^2=\frac{4\times0.000470}{0.015^2}=\frac{0.00188}{0.000225}\approx8.356<9.488.
\]
故不能拒绝 \(H_0\)，不能认为标准差显著偏大。
\end{solution}


\begin{exercise}[\texorpdfstring{\(F\)}{F} 检验与合并 \texorpdfstring{\(t\)}{t} 检验] \label{ex:ht_f_pooled}
在施肥与不施肥的两块苗床上各抽取 6 株苗木，苗高数据（cm）如下：
\begin{center}
\begin{tabular}{ccccccc|c}
\toprule
施肥 & 77.3 & 79.1 & 81.0 & 79.1 & 82.1 & 77.3 & 行和 475.9 \\
\midrule
不施肥 & 75.5 & 76.2 & 78.1 & 72.4 & 77.4 & 76.7 & 行和 456.3 \\
\bottomrule
\end{tabular}
\end{center}
设苗高服从正态分布（\(\alpha=0.05\)）。\\
(1) 检验两组方差是否相同；(2) 检验施肥是否显著高于不施肥。
（\(F_{0.025}(5,5)=7.15\)，\(F_{0.975}(5,5)=0.14\)，\(t_{0.05}(10)=1.812\)）
\end{exercise}
\begin{solution}
\(n_1=n_2=6\)，\(\overline{x}_1=475.9/6\approx79.317\)，\(\overline{x}_2=456.3/6=76.05\)。
离差平方和：
\[
\sum(x_{1i}-\overline{x}_1)^2\approx18.808,\quad S_1^2=\frac{18.808}{5}=3.762;
\]
\[
\sum(x_{2j}-\overline{x}_2)^2=20.095,\quad S_2^2=\frac{20.095}{5}=4.019.
\]
\textbf{(1)} \(H_0\colon\sigma_1^2=\sigma_2^2\)，\(H_1\colon\sigma_1^2\neq\sigma_2^2\)。
\[
F=\frac{S_1^2}{S_2^2}=\frac{3.762}{4.019}\approx0.936.
\]
拒绝域：\(F\geq7.15\) 或 \(F\leq0.14\)。\(0.14<0.936<7.15\)，未落入拒绝域，故不能拒绝 \(H_0\)，即方差无显著差异。\\
\textbf{(2)} 方差齐性成立，用 pooled \(t\) 检验。\(H_0\colon\mu_1\leq\mu_2\)，\(H_1\colon\mu_1>\mu_2\)。
\[
S_w^2=\frac{18.808+20.095}{10}=3.890,\quad S_w\approx1.972.
\]
\[
t=\frac{79.317-76.05}{1.972\sqrt{1/6+1/6}}=\frac{3.267}{1.972/\sqrt{3}}\approx\frac{3.267}{1.139}\approx2.869.
\]
拒绝域 \(t\geq t_{0.05}(10)=1.812\)。\(2.869>1.812\)，拒绝 \(H_0\)，施肥苗高显著高于不施肥。
\end{solution}


\section{置信区间与假设检验的关系} \label{sec:ht-ci}

\begin{property}[置信区间与假设检验的等价性]\label{prop:ci-test-equiv}
设 \((\theta_L, \theta_U)\) 为参数 \(\theta\) 的置信水平为 \(1-\alpha\) 的置信区间，则对于双边检验
\[
H_0\colon \theta=\theta_0,\quad H_1\colon \theta\neq\theta_0,
\]
“不能拒绝 \(H_0\)”等价于 “\(\theta_0\) 落在置信区间 \((\theta_L, \theta_U)\) 内”。对于单边检验也有类似的对应关系：若做右边检验 \(H_0\colon \theta\leq\theta_0\) 并构造下单侧置信区间 \((\underline{\theta},+\infty)\)，则拒绝域 \(z\geq z_\alpha\) 等价于 \(\theta_0<\underline{\theta}\)；若做左边检验 \(H_0\colon \theta\geq\theta_0\) 并构造上单侧置信区间 \((-\infty,\overline{\theta})\)，则拒绝域对应于 \(\theta_0>\overline{\theta}\)。
\end{property}

\begin{note}
\textbf{易错点：} “置信区间和假设检验等价”有严格前提：参数相同、显著性水平匹配、而且尾部形式也要匹配。双侧 \(1-\alpha\) 置信区间对应双边显著性水平为 \(\alpha\) 的检验；右边检验要配下单侧区间，左边检验要配上单侧区间，不能拿双侧区间去机械替代单边检验。
\end{note}

\begin{exercise}\label{exer:9-ci-test}
零件长度 \(X \sim N(\mu, \sigma^2)\)（\(\mu, \sigma^2\) 均未知），\(n=16\)，\(\overline{x}=10\)，\(s^2=0.16\)。\\
(1) 求 \(\mu\) 的置信水平为 0.95 的置信区间；\\
(2) 在 \(\alpha=0.05\) 下检验 \(H_0\colon \mu=9.7\)，\(H_1\colon \mu\neq9.7\)。（\(t_{0.025}(15)=2.132\)）
\end{exercise}
\begin{solution}
\textbf{(1) 求置信区间。}\(\mu\) 未知，\(\sigma^2\) 未知，使用 \(t\) 分布。
已知 \(n=16\)，\(\overline{x}=10\)，\(s=\sqrt{0.16}=0.4\)，\(t_{0.025}(15)=2.132\)。置信区间为：
\[
\left(\overline{x} - \frac{s}{\sqrt{n}}\,t_{\alpha/2}(n{-}1),\quad \overline{x} + \frac{s}{\sqrt{n}}\,t_{\alpha/2}(n{-}1)\right).
\]
代入计算：
\[
\frac{s}{\sqrt{n}}\,t_{\alpha/2}(15) = \frac{0.4}{4}\times2.132 = 0.2132.
\]
\[
\text{置信区间} = (10-0.2132,\ 10+0.2132) = (9.7868,\ 10.2132).
\]
\textbf{(2) 假设检验。}提出假设：\(H_0\colon \mu=9.7\)，\(H_1\colon \mu\neq9.7\)（双边检验，\(\alpha=0.05\)）。
检验统计量 \(t = \dfrac{\overline{X}-9.7}{S/\sqrt{n}} \sim t(15)\)，拒绝域 \(|t| \geq t_{0.025}(15)=2.132\)。
\[
t = \frac{10-9.7}{0.4/4} = \frac{0.3}{0.1} = 3.0 > 2.132.
\]
落入拒绝域，拒绝 \(H_0\)，认为 \(\mu\neq9.7\)。
\textbf{等价性验证}：由命题~\ref{prop:ci-test-equiv}，双边检验中“不能拒绝 \(H_0\)”等价于 \(\mu_0\) 落在置信区间内。此处 \(\mu_0=9.7\)，而置信区间为 \((9.7868,\ 10.2132)\)，\(9.7 < 9.7868\)，\(\mu_0\) 在区间外，故应拒绝 \(H_0\)，与上述检验结论一致。
\end{solution}

\begin{exercise}[置信区间与假设检验的综合] \label{ex:ci_ht_01}
根据环保条例，排放工业废水中某有害物质含量不得超过 0.5。设该含量 \(X\sim N(\mu,\sigma^2)\)，现取 5 份水样，测得含量分别为
\[
0.530,\ 0.542,\ 0.510,\ 0.495,\ 0.515.
\]
(1) 求 \(\mu\) 的置信度为 95\% 的置信区间；\\
(2) 能否据此说明有害物质含量超过了规定标准（\(\alpha=0.05\)）？
（\(t_{0.05}(4)=2.132\)，\(t_{0.025}(4)=2.776\)）
\end{exercise}
\begin{solution}
\textbf{数据预处理。} \(n=5\)
\[
\overline{x}=\frac{0.530+0.542+0.510+0.495+0.515}{5}=0.5184.
\]
\[
\sum_{i=1}^{5}(x_i-0.5184)^2=0.0116^2+0.0236^2+0.0084^2+0.0234^2+0.0034^2=0.001321.
\]
\[
s^2=\frac{0.001321}{4}=0.000330,\quad s\approx 0.0182.
\]
(1) \(\sigma^2\) 未知，\(\mu\) 的 95\% 置信区间为
\begin{align*}
&\left(\overline{x}-\frac{s}{\sqrt{n}}\,t_{0.025}(4),\;\overline{x}+\frac{s}{\sqrt{n}}\,t_{0.025}(4)\right)\\
=&\left(0.5184-\frac{0.0182}{\sqrt{5}}\times2.776,\;0.5184+\frac{0.0182}{\sqrt{5}}\times2.776\right)=(0.4958,\;0.5410).
\end{align*}
(2) 右边检验：\(H_0\colon\mu\leq 0.5\)，\(H_1\colon\mu>0.5\)。
\[
t=\frac{0.5184-0.5}{0.0182/\sqrt{5}}=\frac{0.0184}{0.00813}\approx 2.264.
\]
拒绝域 \(t\geq t_{0.05}(4)=2.132\)。\(2.264>2.132\)，落入拒绝域，拒绝 \(H_0\)，即认为有害物质含量超过了规定标准。
\begin{note}
也可以利用(1)的结论：\(\mu_0=0.5\) 不在置信区间 \((0.4958,0.5410)\) 内，但此区间为\textbf{双侧}置信区间，对应\textbf{双边}检验。而本题(2)是\textbf{单边}检验，二者不完全等价，需单独计算。
\end{note}
\end{solution}


\section{两类错误}

\begin{note}
在固定样本容量 \(n\) 的条件下，\(\alpha\) 与 \(\beta\) 存在此消彼长的关系：\(\alpha\) 小则 \(\beta\) 大，\(\beta\) 小则 \(\alpha\) 大。\textbf{它们不满足} \(\beta=1-\alpha\)。
实际中，总是在控制 \(\alpha\)（使弃真概率较小）的条件下尽量使 \(\beta\) 小，因为人们通常认为错误地拒绝 \(H_0\) 比错误地未能拒绝 \(H_0\) 的后果更严重。
\end{note}

\begin{exercise}[离散总体第一类错误的计算] \label{ex:error_discrete}
设总体 \(X\) 的分布律为 \(P\{X=1\}=\theta\)，\(P\{X=0\}=1-\theta\)（\(0<\theta<1\)），\(X_1,X_2,X_3\) 为样本。对于检验
\[
H_0\colon\theta=0.1,\quad H_1\colon\theta>0.1,
\]
拒绝域为 \(\{X_1=1,\ X_2=0,\ X_3=1\}\)。求犯第一类错误的概率。
\end{exercise}
\begin{solution}
第一类错误指的是：\textbf{\(H_0\) 为真时却拒绝了 \(H_0\)} 的概率。

本题原假设为 \(H_0:\theta=0.1\)，而拒绝域给成 \(\{X_1=1,\ X_2=0,\ X_3=1\}\)。
因此犯第一类错误的概率，就是在 \(\theta=0.1\) 成立时，样本恰好落入这个拒绝域的概率。

当 \(\theta=0.1\) 时，
\[
P(X_i=1)=0.1,\qquad P(X_i=0)=0.9.
\]
又因为 \(X_1,X_2,X_3\) 是样本，相互独立，所以
\[
\alpha=P\{X_1=1,X_2=0,X_3=1\mid\theta=0.1\}=0.1\times0.9\times0.1=0.009.
\]

故犯第一类错误的概率为
\[
\alpha=0.009.
\]
\end{solution}


\begin{exercise}[正态总体两类错误的计算] \label{ex:error_normal}
设 \(X_1,X_2,\dots,X_{20}\) 为总体 \(N(\mu,1)\) 的样本，考虑假设检验
\[
H_0\colon\mu=2\quad\text{vs}\quad H_1\colon\mu=3,
\]
拒绝域为 \(W=\{\overline{X}\geq 2.6\}\)。求犯第一类错误的概率 \(\alpha\) 和犯第二类错误的概率 \(\beta\)。
\end{exercise}
\begin{solution}
\(\overline{X}\sim N(\mu,\,1/n)\)，此处 \(n=20\)，\(\overline{X}\sim N(\mu,\,1/20)\)。\\
\textbf{第一类错误}：\(\alpha=P\{\overline{X}\geq2.6\mid\mu=2\}\)。
标准化：令 \(Z=\dfrac{\overline{X}-2}{1/\sqrt{20}}\sim N(0,1)\)，
\[
\alpha=P\!\left\{Z\geq\frac{2.6-2}{1/\sqrt{20}}\right\}=P\{Z\geq 0.6\sqrt{20}\}=P\{Z\geq 1.2\sqrt{5}\}=1-\Phi(1.2\sqrt{5}).
\]
\textbf{第二类错误}：\(\beta=P\{\overline{X}<2.6\mid\mu=3\}\)。
令 \(Z'=\dfrac{\overline{X}-3}{1/\sqrt{20}}\sim N(0,1)\)，
\[
\beta=P\!\left\{Z'<\frac{2.6-3}{1/\sqrt{20}}\right\}=P\{Z'<-0.4\sqrt{20}\}=P\{Z'<-0.8\sqrt{5}\}=\Phi(-0.8\sqrt{5}).
\]
\end{solution}


\begin{exercise}\label{exer:9-error-calc}
设连续型随机变量 \(U\) 的密度分别为：
\[
H_0\colon f(x)=\begin{cases}1/2, & 0\leq x\leq2,\\0, & \text{其他},\end{cases}
\quad
H_1\colon f(x)=\begin{cases}x/2, & 0\leq x\leq2,\\0, & \text{其他}.\end{cases}
\]
检验规则：当 \(U>3/2\) 时拒绝 \(H_0\)。求犯第一类错误的概率 \(\alpha\) 和犯第二类错误的概率 \(\beta\)。\par
A. \(3/4,\;7/16\)\qquad B. \(7/16,\;3/4\)\qquad C. \(9/16,\;1/4\)\qquad D. \(1/4,\;9/16\)
\end{exercise}
\begin{solution}
应选 D。根据两类错误的定义：
\begin{itemize}
    \item \textbf{第一类错误}\(\alpha = P\{\text{拒绝}\,H_0 \mid H_0\text{ 为真}\} = P\!\left\{U>\tfrac{3}{2}\,\middle|\,H_0\right\}\)。
    在 \(H_0\) 下，\(U\) 服从 \([0,2]\) 上的均匀分布，密度 \(f(x)=1/2\)：
    \[
    \alpha = \int_{3/2}^{2}\frac{1}{2}\,\mathrm{d}x = \frac{1}{2}\times\left(2-\frac{3}{2}\right) = \frac{1}{2}\times\frac{1}{2} = \frac{1}{4}.
    \]
    \item \textbf{第二类错误}\(\beta = P\{\text{未拒绝}\,H_0 \mid H_1\text{ 为真}\} = P\!\left\{U\leq\tfrac{3}{2}\,\middle|\,H_1\right\}\)。
    在 \(H_1\) 下，\(U\) 的密度 \(f(x)=x/2\)（\(0\leq x\leq2\)）：
    \[
    \beta = \int_0^{3/2}\frac{x}{2}\,\mathrm{d}x = \frac{1}{2}\times\frac{x^2}{2}\Bigg|_0^{3/2} = \frac{1}{4}\times\frac{9}{4} = \frac{9}{16}.
    \]
\end{itemize}
故 \(\alpha=1/4\)，\(\beta=9/16\)，选 D。
\end{solution}


\begin{exercise}\label{exer:9-error-concept1}
在假设检验中，\(H_0\) 不成立时，检验统计量值未落入拒绝域，从而没有拒绝 \(H_0\)，这种错误称为第\textbf{二}类错误。
\end{exercise}
\begin{solution}
第二类错误（取伪）：\(H_0\) 不成立（\(H_1\) 为真）时，检验统计量值未落入拒绝域，从而错误地未拒绝 \(H_0\)。
\begin{note}
要区分两类错误的关键：
\begin{itemize}
    \item 第一类错误：\(H_0\) 本来为真，却被错误拒绝；
    \item 第二类错误：\(H_0\) 本来不真，却没有被拒绝。
\end{itemize}
\end{note}
\end{solution}


\begin{exercise}\label{exer:9-error-concept2}
显著性水平 \(\alpha\) 的含义是（\quad）。\par
A. \(H_0\) 成立，经检验被拒绝的概率\qquad B. \(H_0\) 成立，经检验不能拒绝的概率\par
C. \(H_0\) 不成立，经检验被拒绝的概率\qquad D. \(H_0\) 不成立，经检验不能拒绝的概率
\end{exercise}
\begin{solution}
应选 A。显著性水平 \(\alpha\) 的定义是：在 \(H_0\) 成立（为真）的条件下，检验结果拒绝 \(H_0\) 的概率，即
\[
\alpha = P\{\text{拒绝}\,H_0 \mid H_0\text{ 为真}\}.
\]
逐项分析：
\begin{itemize}
    \item A: \(H_0\) 成立，经检验被拒绝 \(\longrightarrow\) 恰好是 \(\alpha\) 的定义 \(\checkmark\)。
    \item B: \(H_0\) 成立，不能拒绝 \(\longrightarrow\) 这是 \(1-\alpha\)，不是 \(\alpha\) \(\times\)。
    \item C: \(H_0\) 不成立，被拒绝 \(\longrightarrow\) 这是\textbf{检验功效} \(1-\beta\) \(\times\)。
    \item D: \(H_0\) 不成立，不能拒绝 \(\longrightarrow\) 这是第二类错误概率 \(\beta\) \(\times\)。
\end{itemize}
\end{solution}

\begin{note}
\textbf{自学检查点——假设检验：}
\begin{enumerate}
    \item 假设检验的三步是什么？（答：建立假设 $\to$ 选统计量和拒绝域 $\to$ 代入数据做判断）
    \item 第一类错误和第二类错误分别是什么？（答：第一类 = 弃真（$H_0$ 真却拒绝）；第二类 = 取伪（$H_0$ 假却不拒绝））
    \item 显著性水平 $\alpha$ 控制的是哪类错误？（答：第一类）
    \item "不拒绝 $H_0$"等于"$H_0$ 为真"吗？（答：不等于！只是证据不足以拒绝）
    \item 检验均值 $\mu$（$\sigma$ 未知）用什么统计量？（答：$t=\frac{\bar{X}-\mu_0}{S/\sqrt{n}}\sim t(n-1)$）
    \item 双边检验的拒绝域形式？（答：$|T|>t_{\alpha/2}(n-1)$）
\end{enumerate}
\end{note}

\begin{note}
\textbf{【高频题型模板】假设检验大题的标准四步：}
\begin{enumerate}
    \item \textbf{建立假设：} 找"想证明的结论"放 $H_1$，$H_0$ 取等号方向。写清 $H_0$ 和 $H_1$。
    \item \textbf{选统计量：} 与区间估计完全相同的判断逻辑（检 $\mu$/$\sigma^2$，$\sigma$ 知/不知）。把 $H_0$ 中的参数值代入枢轴量，得到检验统计量。
    \item \textbf{定拒绝域：} 看 $H_1$ 的方向——"$>$"右尾、"$<$"左尾、"$\neq$"双尾。查表得临界值。
    \item \textbf{算值下结论：} 代入样本数据算出统计量观测值，判断是否落入拒绝域。
\end{enumerate}
\textbf{答题格式（阅卷得分点）：}
\begin{itemize}
    \item 明确写出 $H_0$ 和 $H_1$（1分）
    \item 写出检验统计量及其分布（2分）
    \item 写出拒绝域并标明临界值（2分）
    \item 代入计算并给出"拒绝/不拒绝 $H_0$"的结论（1分）
\end{itemize}
\end{note}



\chapter{经典概率模型与随机过程初探}

本章不再额外引入新的主理论，而是把前面已经学过的古典概型、条件概率、随机变量分布、二维区域概率和数学期望重新组织成几个经典模型。它们常出现在期末综合题、提高题和考研概率统计的入门题中，表面上是“新故事”，实质上仍然是在调用前文已经学过的方法。

\begin{note}
\textbf{使用提醒：} 学这一章时，重点不是记住几个背景故事，而是看清每个模型的“第一步”。
\begin{itemize}
    \item 抽签公平与生日问题：先想古典概型、对立事件和不放回或独立抽样；
    \item 赌徒破产：先写“第一步之后会发生什么”，再补边界条件；
    \item 会面问题：先把题意翻译成二维独立均匀分布下的区域面积；
    \item 优惠券收集：先把总过程拆成若干阶段，再调用几何分布和期望线性性。
\end{itemize}
本章主线只要求做到“识别模型、选对方法、能写出第一步”，不把高等概率论中的随机过程理论作为考试要求。
\end{note}

\section{古典概型与对立事件模型}

这一节处理两类常见综合题：一类是“不放回抽取但只关心某个位置”的题，另一类是“至少出现一次重复”的题。前者回扣第一章的全概率公式和第二章的超几何分布思想，后者的关键则是学会先求对立事件。

\begin{note}
\textbf{回扣前文：} 本节会反复调用第 1 章的古典概型、乘法公式与全概率公式。若题目问“前面若干次抽取中有多少次命中”，还应联想到第 2 章的超几何分布。
\end{note}

\subsection{抽签模型}

抽签公平看似只是常识，但它能训练一种很重要的识别能力：当题目里有“顺序在变、组成在变、但中间结果未知”的结构时，往往要把“未知结果”用全概率公式平均掉。

\begin{definition}[抽签模型]\label{def:drawing-lots-model}
设有 \(N\) 个球，其中 \(M\) 个为红球，其余 \(N-M\) 个为白球，并设 \(1\le n\le N\)。\(n\) 个人依次不放回地各随机抽取一球。记
\[
A_k=\{\text{第 }k\text{ 个人抽中红球}\},\qquad k=1,2,\dots,n.
\]
\end{definition}

\begin{theorem}[抽签的公平性]\label{thm:drawing-lots-fairness}
无论抽取顺序如何，每个人抽中红球的概率都相同，即
\[
P(A_k)=\frac{M}{N},\qquad k=1,2,\dots,n.
\]
\end{theorem}

\begin{proof}
对第一个人，显然有
\[
P(A_1)=\frac{M}{N}.
\]

下面看第二个人。由于第二个人是否抽中红球，取决于第一个人抽走的是红球还是白球，所以先按事件组 \(\{A_1,\overline{A}_1\}\) 分解：
\[
P(A_2)=P(A_2|A_1)P(A_1)+P(A_2|\overline{A}_1)P(\overline{A}_1).
\]
其中
\[
P(A_2|A_1)=\frac{M-1}{N-1},\qquad P(A_2|\overline{A}_1)=\frac{M}{N-1}.
\]
代入得
\begin{align*}
P(A_2)
&=\frac{M-1}{N-1}\cdot\frac{M}{N}+\frac{M}{N-1}\cdot\frac{N-M}{N} \\
&=\frac{M(M-1)+M(N-M)}{N(N-1)} \\
&=\frac{M(N-1)}{N(N-1)}=\frac{M}{N}.
\end{align*}

对任意第 \(k\) 个人，也可以直接用“随机排列”的对称性来理解：把 \(N\) 个球随机排成一列后，第 \(k\) 个位置是红球的概率，和第一个位置、第二个位置是红球的概率并没有区别，因此仍为 \(\frac{M}{N}\)。
\end{proof}

\begin{note}
\textbf{使用提醒：} 若题目只问“第 \(k\) 次抽中”的概率，可直接用抽签公平性；但若题目进一步问“前 \(m\) 次中恰有 \(r\) 次命中”，那就不再只是公平性问题，而要转化为第 2 章中的超几何分布模型。
\end{note}

\begin{exercise}[抽签公平与不放回抽样]\label{ex:classic_lots}
袋中有 5 个红球和 7 个白球，甲、乙、丙依次不放回地各取 1 球。求：
\begin{enumerate}
    \item 乙取到红球的概率；
    \item 已知甲未取到红球，乙取到红球的概率；
    \item 甲、乙两人中恰有一人取到红球的概率。
\end{enumerate}
\end{exercise}

\begin{solution}
先记
\[
A=\{\text{甲取到红球}\},\qquad B=\{\text{乙取到红球}\}.
\]

(1) 这一问只关心“第二个位置是否为红球”，可以直接用抽签公平性：
\[
P(B)=\frac{5}{12}.
\]

(2) 这里附加了条件“甲未取到红球”，说明第一个人已经取走了一个白球。此时袋中剩下 11 个球，其中 5 个红球、6 个白球，所以
\[
P(B|\overline A)=\frac{5}{11}.
\]

(3) “甲、乙中恰有一人取到红球”可以分成两个互斥情形：
\[
A\overline B\qquad\text{和}\qquad \overline A B.
\]
因此
\[
P(\text{恰有一人取到红球})=P(A\overline B)+P(\overline A B).
\]
分别计算：
\[
P(A\overline B)=\frac{5}{12}\cdot\frac{7}{11},\qquad
P(\overline A B)=\frac{7}{12}\cdot\frac{5}{11}.
\]
于是
\[
P(A\overline B)+P(\overline A B)
=\frac{5}{12}\cdot\frac{7}{11}+\frac{7}{12}\cdot\frac{5}{11}
=\frac{35}{66}.
\]

所以三问答案依次为
\[
\frac{5}{12},\qquad \frac{5}{11},\qquad \frac{35}{66}.
\]
\end{solution}

\subsection{生日碰撞模型}

这类题的共同特点是：题目问“至少有一次重复”，而直接求“至少一次”很难，因此标准入口往往是先求“完全没有重复”的概率，再取对立事件。

\begin{definition}[生日碰撞模型]\label{def:birthday-model}
设有 \(n\) 个人，忽略闰年，假定每个人的生日独立且等可能地落在 365 天中的任一天。我们关心事件
\[
C_n=\{\text{至少有两个人生日相同}\}.
\]
\end{definition}

\begin{theorem}[生日问题的基本公式]\label{thm:birthday-model}
当 \(1\le n\le 365\) 时，
\[
P(\overline{C_n})=\frac{365\cdot364\cdots(365-n+1)}{365^n},
\]
从而
\[
P(C_n)=1-\frac{365\cdot364\cdots(365-n+1)}{365^n}.
\]
当 \(n>365\) 时，由抽屉原理知 \(P(C_n)=1\)。
\end{theorem}

\begin{proof}
先求对立事件 \(\overline{C_n}\)，即“\(n\) 个人的生日互不相同”。

第 1 个人的生日可任意取，因此概率为 1；第 2 个人要避开前 1 个人的生日，概率为 \(\frac{364}{365}\)；第 3 个人要避开前 2 个人的生日，概率为 \(\frac{363}{365}\)；依此类推，第 \(n\) 个人要避开前 \(n-1\) 个人的生日，概率为 \(\frac{365-n+1}{365}\)。

由乘法公式，
\[
P(\overline{C_n})
=1\cdot\frac{364}{365}\cdot\frac{363}{365}\cdots\frac{365-n+1}{365}
=\frac{365\cdot364\cdots(365-n+1)}{365^n}.
\]
因此
\[
P(C_n)=1-P(\overline{C_n})
=1-\frac{365\cdot364\cdots(365-n+1)}{365^n}.
\]
\end{proof}

\begin{note}
\textbf{易错点：} 这类题最常见的误判是直接去算“至少一对相同”。真正容易写出的第一步，是先算“所有人都不同”，因为这个事件天然能按先后顺序用乘法公式展开。

\textbf{考研延伸：} 当 \(n\) 相对 365 不太大时，可用
\[
\prod_{k=0}^{n-1}\left(1-\frac{k}{365}\right)
\approx \exp\!\left(-\frac{n(n-1)}{2\cdot365}\right)
\]
近似，从而得到
\[
P(C_n)\approx 1-\exp\!\left(-\frac{n(n-1)}{730}\right).
\]
这个近似只作视野补充，不作为本章主线要求。
\end{note}

\begin{exercise}[生日碰撞概率]\label{ex:birthday_collision}
某班有 30 名同学，假设每名同学的生日独立且等可能地落在 365 天中的任一天。求班里至少有两名同学生日相同的概率表达式。
\end{exercise}

\begin{solution}
这道题的关键不是直接求“至少有两名同学生日相同”，而是先求它的对立事件。

设
\[
C=\{\text{30 名同学中至少有两人同生日}\}.
\]
则其对立事件为
\[
\overline C=\{\text{30 名同学的生日两两不同}\}.
\]

对事件 \(\overline C\)，第 1 名同学的生日不受限制；第 2 名同学要避开前 1 人的生日；第 3 名同学要避开前 2 人的生日；一直到第 30 名同学要避开前 29 人的生日。由乘法公式，
\[
P(\overline C)
=\frac{365}{365}\cdot\frac{364}{365}\cdot\frac{363}{365}\cdots\frac{336}{365}
=\frac{365\cdot364\cdots336}{365^{30}}.
\]

因此所求概率为
\[
P(C)=1-P(\overline C)
=1-\frac{365\cdot364\cdots336}{365^{30}}.
\]

所以班里至少有两名同学生日相同的概率表达式是
\[
1-\frac{365\cdot364\cdots336}{365^{30}}.
\]
\end{solution}

\section{条件分解与递推模型}

有些题目无法直接枚举所有可能路径，但每走一步之后，问题又会回到“同一类问题，只是状态变了”。这时最自然的入口就是“第一步分析法”：先按第一步的结果分解，再把未知量写成递推关系。

\begin{note}
\textbf{回扣前文：} 这一节本质上仍然是在用第 1 章的全概率公式。差别只在于：分解之后，未知量没有被一次算完，而是转化成了新的未知量，于是形成递推方程。
\end{note}

\subsection{赌徒破产问题}

赌徒破产问题是“第一步分析法”的标准模型。它的价值不只在于得到一个公式，更在于训练我们写出递推式和边界条件。

\begin{definition}[赌徒破产模型]\label{def:gambler-ruin-model}
一名赌徒有初始本金 \(i\) 元，目标是达到 \(N\) 元就停止；若本金降到 0 元，则视为破产并停止，其中 \(0\le i\le N\)。每次赌博，赌徒以概率 \(p\) 赢 1 元，以概率 \(q=1-p\) 输 1 元。记
\[
P_i=\{\text{从本金 }i\text{ 出发，最终先达到 }N\text{ 而不是先破产的概率}\}.
\]
\end{definition}

\begin{theorem}[公平赌局中的达标概率]\label{thm:gambler-fair}
若赌局公平，即 \(p=q=\frac12\)，则
\[
P_i=\frac{i}{N},\qquad i=0,1,\dots,N.
\]
\end{theorem}

\begin{proof}
对于 \(i=1,2,\dots,N-1\)，第一步只可能有两种结果：
\begin{itemize}
    \item 以概率 \(\frac12\) 赢 1 元，此后问题变成“从 \(i+1\) 元出发是否先达到 \(N\)”；
    \item 以概率 \(\frac12\) 输 1 元，此后问题变成“从 \(i-1\) 元出发是否先达到 \(N\)”。
\end{itemize}
因此由全概率公式，
\[
P_i=\frac12P_{i+1}+\frac12P_{i-1},\qquad i=1,2,\dots,N-1.
\]
整理得 \(P_{i+1}-P_i=P_i-P_{i-1}\)，
说明相邻两项之差相等，所以数列 \(\{P_i\}\) 是等差数列。

再利用边界条件 \(P_0=0,\ P_N=1\)，便知等差数列从 0 线性增长到 1，因此 \(P_i=\frac{i}{N}\)。
\end{proof}

\begin{proposition}[偏置赌局的扩展公式]
若 \(p\neq q\)，则
\[
P_i=\frac{1-(q/p)^i}{1-(q/p)^N},\qquad i=0,1,\dots,N.
\]
\end{proposition}

\begin{proof}
同样先做第一步分析：
\[
P_i=pP_{i+1}+qP_{i-1},\qquad i=1,2,\dots,N-1,
\]
且边界条件仍为
\[
P_0=0,\qquad P_N=1.
\]

把上式写成
\[
p(P_{i+1}-P_i)=q(P_i-P_{i-1}),
\]
即
\[
P_{i+1}-P_i=\frac{q}{p}(P_i-P_{i-1}).
\]
这说明相邻差值构成等比数列。由此累加可得
\[
P_i=\frac{1-(q/p)^i}{1-q/p}\cdot P_1.
\]
再代入边界条件 \(P_N=1\)，得
\[
1=\frac{1-(q/p)^N}{1-q/p}\cdot P_1\implies P_1=\frac{1-q/p}{1-(q/p)^N},
\]
代回即可得到
\[
P_i=\frac{1-(q/p)^i}{1-(q/p)^N}.
\]
\end{proof}

\begin{note}
\textbf{易错点：} 做这类题时，顺序一定是“先写边界条件，再写递推式”。边界条件写错，后面即使递推整理正确，也会把整个答案带偏。这里最基本的边界就是
\[
P_0=0,\qquad P_N=1.
\]
\end{note}

\begin{exercise}[公平赌局中的达标概率]\label{ex:gambler_fair}
某赌徒现有 3 元，与庄家进行公平赌博，每次赢 1 元或输 1 元。若资金达到 10 元就停止，若输光也停止。求该赌徒最终达到 10 元的概率。
\end{exercise}

\begin{solution}
这是一道标准的公平赌徒破产题。

先识别参数：\(i=3,\ N=10,\ p=q=\frac12\)。由于赌局公平，可直接调用定理~\ref{thm:gambler-fair} 中 \(P_i=\frac{i}{N}\)，于是 \(P_3=\frac{3}{10}\)。

若想从模型角度理解，这个结果表示：在公平赌局下，最终先达到目标的概率，恰好等于“当前本金占目标本金的比例”。

因此该赌徒最终达到 10 元的概率为 \(\frac{3}{10}\)。
\end{solution}

\begin{exercise}[偏置赌局中的达标概率]\label{ex:gambler_bias}
某赌徒现有 2 元，每次以概率 0.4 赢 1 元，以概率 0.6 输 1 元。若资金达到 5 元就停止，若输光也停止。求该赌徒最终达到 5 元的概率。
\end{exercise}

\begin{solution}
先判断题型。由于每一步赢的概率和输的概率不相等，所以这不是公平赌局，而是偏置赌局，应使用扩展公式。

题中参数为
\[
i=2,\qquad N=5,\qquad p=0.4=\frac25,\qquad q=0.6=\frac35.
\]
因此
\[
\frac{q}{p}=\frac{3/5}{2/5}=\frac32.
\]

由偏置赌局的公式
\[
P_i=\frac{1-(q/p)^i}{1-(q/p)^N},
\]
可得
\[
P_2=\frac{1-(3/2)^2}{1-(3/2)^5}.
\]

下面把它化简。先算分子与分母：
\[
1-\left(\frac32\right)^2=1-\frac94=-\frac54,
\]
\[
1-\left(\frac32\right)^5=1-\frac{243}{32}=-\frac{211}{32}.
\]
于是
\[
P_2=\frac{-5/4}{-211/32}
=\frac{5}{4}\cdot\frac{32}{211}
=\frac{40}{211}.
\]

所以该赌徒最终达到 5 元的概率为
\[
\frac{40}{211}\approx0.1896.
\]

这个结果也符合直觉：由于每一步输钱的概率更大，所以最终达标的概率应明显小于 \(\frac25\) 或 \(\frac12\)。
\end{solution}

\section{二维连续模型与几何概率}

本节的代表题看起来像“相约见面”的应用题，但本质仍是第 3 章的二维连续型随机变量。关键不在积分技巧，而在于能否把题目正确翻译成平面区域。

\begin{note}
\textbf{回扣前文：} 若两个随机时刻相互独立且都在某个区间上均匀分布，就应立刻想到第 3 章中的二维独立均匀分布。做题第一步通常是设
\[
X\sim U(0,T),\qquad Y\sim U(0,T),\qquad X\perp Y,
\]
再把条件改写成 \((X,Y)\) 平面中的区域。
\end{note}

\subsection{会面问题}

会面问题是几何概率中最常见、也最适合考试的模型之一，因为它把“文字条件”直接转化成“平面区域”的面积计算。

\begin{definition}[会面模型]\label{def:meeting-model}
甲、乙两人约定在时长为 \(T\) 的时间段内独立且均匀地到达某地，分别记到达时刻为 \(X,Y\)。若两人都会等待 \(\tau\) 个时间单位后离开，其中 \(0\le \tau\le T\)，则称事件
\[
M=\{\text{两人能够见面}\}
\]
为会面事件。
\end{definition}

\begin{theorem}[会面概率公式]\label{thm:meeting-model}
在定义~\ref{def:meeting-model} 的条件下，
\[
P(M)=1-\left(1-\frac{\tau}{T}\right)^2.
\]
\end{theorem}

\begin{proof}
两人能够见面，当且仅当两人的到达时刻之差不超过等待时间，即 \(|X-Y|\le \tau\)。由于 \(X,Y\) 独立且都服从 \(U(0,T)\)，所以点 \((X,Y)\) 在正方形 \([0,T]\times[0,T]\) 上均匀分布，总面积为 \(T^2\)。条件 \(|X-Y|\le \tau\) 表示正方形中夹在两条直线 \(y=x+\tau\) 与 \(y=x-\tau\) 之间的带状区域。其补集由两个全等直角三角形组成，每个三角形的直角边长都是 \(T-\tau\)，所以补集总面积为
\[
2\cdot\frac12(T-\tau)^2=(T-\tau)^2.
\]

因此
\[
P(M)=1-\frac{(T-\tau)^2}{T^2}
=1-\left(1-\frac{\tau}{T}\right)^2.
\]
\end{proof}

\begin{note}
\textbf{易错点：} 这类题不是一维均匀分布题，而是\textbf{二维独立均匀分布}题。若只写一个随机变量，再试图直接在数轴上做概率，往往会漏掉另一人的到达时刻。

\textbf{选读：} 布丰投针和会面问题的方法本质相同，都是把“随机位置 + 几何条件”转化为二维区域概率。若继续使用第 3 章二维连续型随机变量的工具，可得短针情形下
\[
P(\text{针与平行线相交})=\frac{2l}{\pi d}.
\]
但这类模型在本章只作视野补充，不作为考试主线。
\end{note}

\begin{exercise}[会面问题]\label{ex:meeting_problem}
甲、乙二人在某 30 分钟时间段内独立且均匀地到达约定地点。若两人各自最多等待 5 分钟，求他们能够见面的概率。
\end{exercise}

\begin{solution}
先识别模型。这是一道标准的会面问题，应把两人的到达时刻建模成二维独立均匀分布。

设
\[
X=\text{甲的到达时刻},\qquad Y=\text{乙的到达时刻}.
\]
则
\[
X\sim U(0,30),\qquad Y\sim U(0,30),\qquad X\perp Y.
\]

两人能够见面，当且仅当其中一人到达后，另一人在 5 分钟内到达，也就是 \(|X-Y|\le 5\)。由定理~\ref{thm:meeting-model}，取 \(T=30,\ \tau=5\)，便有
\[
P(\text{见面})=1-\left(1-\frac{5}{30}\right)^2
=1-\left(\frac56\right)^2
=1-\frac{25}{36}
=\frac{11}{36}.
\]

如果从图形上理解，就是在边长为 30 的正方形中，去掉两块边长为 25 的直角三角形，其余带状区域所占的面积比例正好是 \(\frac{11}{36}\)。

所以他们能够见面的概率为
\[
\frac{11}{36}.
\]
\end{solution}

\section{分布分解与数学期望模型}

有些题不是问一个事件会不会发生，而是问“平均要等多久”。这时常见做法不是直接对总过程建模，而是先把总时间拆成若干阶段，再逐段计算期望。

\begin{note}
\textbf{回扣前文：} 本节直接调用第 2 章的几何分布和第 4 章的数学期望线性性。做题时若看到“集齐所有类型”“直到第一次出现新种类”这类表述，应优先考虑按阶段拆分等待时间。
\end{note}

\subsection{优惠券收集模型}

优惠券收集模型常出现在卡片、盲盒、贴纸等题目中。它训练的不是复杂计算，而是把“总等待时间”拆成“每多收集一种所需时间”的能力。

\begin{definition}[优惠券收集模型]\label{def:coupon-collector}
设一共有 \(n\) 种不同的优惠券。每次购买商品时，独立且等概率地随机得到其中一种。记 \(T\) 为从一张都没有开始，到集齐全部 \(n\) 种优惠券所需的购买次数。
\end{definition}

\begin{theorem}[集齐全部优惠券所需次数的期望]\label{thm:coupon-collector}
在定义~\ref{def:coupon-collector} 的条件下，
\[
E(T)=n\left(1+\frac12+\frac13+\cdots+\frac1n\right).
\]
\end{theorem}

\begin{proof}
把总过程拆成 \(n\) 个阶段。设
\[
T_k=\text{从已经收集到 }k-1\text{ 种，到第一次收集到第 }k\text{ 种所需的次数},
\]
其中 \(k=1,2,\dots,n\)。则 \(T=T_1+T_2+\cdots+T_n\)。

先看每个阶段的分布。
\begin{itemize}
    \item 当 \(k=1\) 时，第一张必然带来第 1 种，因此 \(T_1=1\)；
    \item 当已经有 \(k-1\) 种时，还缺 \(n-k+1\) 种，所以下一次买到“新种类”的概率为 \(p_k=\dfrac{n-k+1}{n}\)。
    因而 \(T_k\) 服从参数为 \(p_k\) 的几何分布。
\end{itemize}

由几何分布的期望公式，
\[
E(T_k)=\frac{1}{p_k}=\frac{n}{n-k+1}.
\]
因此
\begin{align*}
E(T)
&=E(T_1)+E(T_2)+\cdots+E(T_n) \\
&=\frac{n}{n}+\frac{n}{n-1}+\cdots+\frac{n}{1} \\
&=n\left(1+\frac12+\frac13+\cdots+\frac1n\right).
\end{align*}

这里用到的关键是期望的线性性。即使不去额外证明各阶段相互独立，上式仍然成立。
\end{proof}

\begin{note}
\textbf{易错点：} 总等待时间之所以能拆成若干阶段相加，不是因为“每个阶段都很像”，而是因为每个阶段都有明确的终点：新增一种此前没有出现过的类型。写出阶段变量后，期望才能顺着线性性加起来。
\end{note}

\begin{exercise}[已有部分收藏时的优惠券收集]\label{ex:coupon_partial}
某商店随机附赠 5 种卡片之一，每次购买时得到各种卡片的概率相同，且各次购买相互独立。若小王目前已经拥有其中 2 种不同卡片，求他平均还需购买多少次才能集齐全部 5 种。
\end{exercise}

\begin{solution}
这道题不需要从头重新集齐，而是从“已经有 2 种”开始继续收集。因此最自然的做法，是只把后面还需要经历的几个阶段拆出来。

设
\[
T=\text{从已有 2 种开始，到集齐 5 种还需的购买次数}.
\]
再设
\[
T_3=\text{从 2 种到 3 种所需次数},\quad
T_4=\text{从 3 种到 4 种所需次数},\quad
T_5=\text{从 4 种到 5 种所需次数}.
\]
则 \(T=T_3+T_4+T_5\)。

现在逐段看成功概率。

当已经有 2 种时，还缺 3 种，所以下一次买到“新种类”的概率为 \(p_3=\frac35\)。
因此
\[
E(T_3)=\frac{1}{p_3}=\frac53.
\]

当已经有 3 种时，还缺 2 种，所以
\[
p_4=\frac25,\qquad E(T_4)=\frac{1}{p_4}=\frac52.
\]

当已经有 4 种时，还缺 1 种，所以
\[
p_5=\frac15,\qquad E(T_5)=\frac{1}{p_5}=5.
\]

由期望的线性性，
\[
E(T)=E(T_3)+E(T_4)+E(T_5)=\frac53+\frac52+5.
\]
通分化简得
\[
E(T)=\frac{10}{6}+\frac{15}{6}+\frac{30}{6}=\frac{55}{6}.
\]

所以小王平均还需购买
\[
\frac{55}{6}\approx 9.17
\]
次才能集齐全部 5 种卡片。
\end{solution}

\section{选读：从递推到随机过程}

前面几个模型虽然背景不同，但都有一个共同点：\textbf{当前状态会影响下一步会发生什么}。只不过在本章主线里，我们只关心一个概率、一个面积或一个期望；若进一步研究“整个演化过程怎样随时间变化”，问题就会自然走向随机过程。

\begin{note}
\textbf{选读 1：波利亚罐子。} 若每次抽到某种颜色后，不仅放回原球，还额外加入同色球，那么“当前颜色比例”就会反过来影响未来抽样概率。它可以看作带有正反馈的抽样模型，和本章前面的条件分解、递推思想一脉相承，只是研究对象从“一个数值”变成了“整条演化路径”。

\textbf{选读 2：埃伦费斯特模型。} 若把系统状态记为“某容器中当前有多少个粒子”，那么下一步如何变化只由当前状态决定。这类模型说明：当题目开始追踪“状态怎样一步步演化”时，前面学过的第一步分析法会进一步发展成马尔可夫链等随机过程工具。

\textbf{使用边界：} 本书在这一章不要求掌握鞅、平稳分布等高阶概念。知道这些模型说明”递推与条件分解还可以继续向随机过程发展”，就足够了。
\end{note}

%% ============================================================
%% 附录：自学工具箱
%% ============================================================

\chapter*{附录：自学工具箱}
\addcontentsline{toc}{chapter}{附录：自学工具箱}

\section*{A. 易混淆概念对比表}
\addcontentsline{toc}{section}{A. 易混淆概念对比表}

\begin{table}[htbp]
\centering
\caption{易混淆概念辨析}
\label{tab:confusion_pairs}
\begin{tabular}{p{3.5cm}p{5.5cm}p{5.5cm}}
\toprule
\textbf{概念对} & \textbf{概念 A} & \textbf{概念 B} \\
\midrule
分布函数 vs 密度函数 & $F(x)=P\{X\leq x\}$，单调不减，值域$[0,1]$ & $f(x)=F'(x)$，可以大于1，面积=1 \\
\midrule
概率 vs 概率密度 & 连续型 $P\{X=a\}=0$（单点概率为零） & $f(a)$ 可以很大，但它不是概率 \\
\midrule
估计量 vs 估计值 & 随机变量（样本的函数），如 $\hat\theta=\bar X$ & 具体数值，如 $\hat\theta=2.5$ \\
\midrule
无偏性 vs 有效性 & $E(\hat\theta)=\theta$（打靶中心对准靶心） & $D(\hat\theta)$ 最小（散布最小） \\
\midrule
置信度 vs 概率 & “95\%置信区间”指方法的长期覆盖率 & 不是”参数有95\%概率在区间内” \\
\midrule
拒绝$H_0$ vs 证明$H_1$真 & 只是”有足够证据反对$H_0$” & 不等于”$H_1$一定正确” \\
\midrule
不拒绝$H_0$ vs 接受$H_0$ & “证据不足，无法拒绝” & 不等于”$H_0$一定为真” \\
\midrule
$\alpha$ vs $p$值 & $\alpha$是预设的门槛（如0.05） & $p$值是数据算出的”极端程度” \\
\midrule
$\sigma^2$ vs $S^2$ & 总体方差（常数，通常未知） & 样本方差（随机变量，除以$n-1$） \\
\midrule
独立 vs 不相关 & 联合分布可拆：$f(x,y)=f_X f_Y$ & 仅线性无关：$\mathrm{Cov}=0$ \\
\midrule
$n$ vs $n-1$（自由度） & $\sum(X_i-\mu)^2/\sigma^2\sim\chi^2(n)$ & $\sum(X_i-\bar X)^2/\sigma^2\sim\chi^2(n-1)$ \\
\bottomrule
\end{tabular}
\end{table}

\section*{B. 典型错误示范}
\addcontentsline{toc}{section}{B. 典型错误示范}

\begin{note}
\textbf{错误 1：求方差时忘记减去均值的平方}

\textbf{错误解法：} 设 $X\sim U(0,1)$，求 $DX$。”$DX=EX^2=\int_0^1 x^2\,\mathrm{d}x=1/3$。”

\textbf{错在哪：} $DX=EX^2-(EX)^2$，不是 $EX^2$。

\textbf{正确：} $EX=1/2$，$EX^2=1/3$，$DX=1/3-1/4=1/12$。
\end{note}

\begin{note}
\textbf{错误 2：独立性判断只看边缘}

\textbf{错误解法：} “已知 $f_X(x)$ 和 $f_Y(y)$，所以 $f(x,y)=f_X(x)f_Y(y)$，独立。”

\textbf{错在哪：} 只有当联合密度\textbf{在整个支撑域上}都等于边缘之积时才独立。若支撑域不是矩形（如 $0<x<y<1$），即使边缘密度都存在，也不独立。

\textbf{正确：} 必须验证 $f(x,y)=f_X(x)\cdot f_Y(y)$ 对\textbf{所有} $(x,y)$ 成立。
\end{note}

\begin{note}
\textbf{错误 3：MLE 对不可导似然函数求导}

\textbf{错误解法：} $X\sim U(0,\theta)$，”$\ln L=-n\ln\theta$，令 $\frac{\mathrm{d}}{\mathrm{d}\theta}\ln L=-n/\theta=0$，无解，故 MLE 不存在。”

\textbf{错在哪：} 忽略了约束 $\theta\geq X_{(n)}$。在约束范围内 $\ln L$ 单调递减，最大值在左端点取得。

\textbf{正确：} $\hat\theta=X_{(n)}=\max\{X_1,\dots,X_n\}$。
\end{note}

\begin{note}
\textbf{错误 4：区间估计用错分位数}

\textbf{错误解法：} $\sigma$ 未知，求 $\mu$ 的95\%置信区间。”查 $z_{0.025}=1.96$，区间为 $\bar x\pm 1.96\cdot S/\sqrt{n}$。”

\textbf{错在哪：} $\sigma$ 未知时应用 $t$ 分布，不是 $z$ 分布。

\textbf{正确：} 查 $t_{0.025}(n-1)$，区间为 $\bar x\pm t_{0.025}(n-1)\cdot S/\sqrt{n}$。
\end{note}

\begin{note}
\textbf{错误 5：假设检验中 $H_0$ 和 $H_1$ 写反}

\textbf{错误解法：} “检验均值是否大于5”，写 $H_0\colon\mu>5$，$H_1\colon\mu\leq 5$。

\textbf{错在哪：} $H_0$ 必须包含等号。”想证明的”放 $H_1$。

\textbf{正确：} $H_0\colon\mu\leq 5$，$H_1\colon\mu>5$（右边检验）。
\end{note}

\begin{note}
\textbf{错误 6：条件密度的 $y$ 范围写成常数}

\textbf{错误解法：} 支撑域为 $0<x<y<1$，求 $f_{Y|X}(y|x)$。”$y\in(0,1)$。”

\textbf{错在哪：} 固定 $x$ 后，$y$ 的范围是 $(x,1)$，不是 $(0,1)$。

\textbf{正确：} $f_{Y|X}(y|x)$ 定义在 $x<y<1$ 上。
\end{note}

\section*{C. 计算技巧集锦}
\addcontentsline{toc}{section}{C. 计算技巧集锦}

\begin{note}
\textbf{技巧 1：凑 Gamma 函数}

$\Gamma(n)=(n-1)!$，$\Gamma(1/2)=\sqrt\pi$。凡是遇到 $\int_0^{+\infty}x^{n-1}e^{-\lambda x}\,\mathrm{d}x$ 的积分，直接写成 $\dfrac{\Gamma(n)}{\lambda^n}$。

\textbf{应用：} 求指数分布的 $k$ 阶矩：$EX^k=\int_0^\infty x^k\lambda e^{-\lambda x}\,\mathrm{d}x=\dfrac{k!}{\lambda^k}$。
\end{note}

\begin{note}
\textbf{技巧 2：正态分布的对称性}

设 $X\sim N(0,1)$，则 $E(X^{2k+1})=0$（奇数阶矩为零），因为被积函数是奇函数。

\textbf{推论：} $\mathrm{Cov}(X,X^2)=EX^3-(EX)(EX^2)=0$。这就是”$X$ 与 $X^2$ 不相关但不独立”的经典反例。
\end{note}

\begin{note}
\textbf{技巧 3：交换积分次序}

求 $E[g(X,Y)]$ 时，若按原次序积分困难，画出积分区域后交换 $\mathrm{d}x\,\mathrm{d}y$ 的次序。

\textbf{典型场景：} $E(\min(X,Y))$、$E(|X-Y|)$ 等——先画出 $x>y$ 和 $x<y$ 的区域，分区积分。
\end{note}

\begin{note}
\textbf{技巧 4：利用已知分布”凑常数”}

若 $f(x)=cx^2 e^{-x}$（$x>0$），要求 $c$。认出 $x^2 e^{-x}$ 是 $\Gamma(3,1)$ 密度的核心部分，而 $\Gamma(3)=2!=2$，故 $c=1/2$。

\textbf{口诀：} 看到 $x^{n-1}e^{-\lambda x}$ 凑 Gamma；看到 $e^{-x^2}$ 凑正态；看到 $x^{a-1}(1-x)^{b-1}$ 凑 Beta。
\end{note}

\begin{note}
\textbf{技巧 5：方差的”拆分-独立”法}

$D(\sum a_i X_i)=\sum a_i^2 DX_i + 2\sum_{i<j}a_i a_j\mathrm{Cov}(X_i,X_j)$。

若独立，交叉项全消：$D(\sum a_i X_i)=\sum a_i^2 DX_i$。

\textbf{快速应用：} $D(\bar X)=D\left(\frac{1}{n}\sum X_i\right)=\frac{1}{n^2}\cdot n\sigma^2=\frac{\sigma^2}{n}$。
\end{note}

\begin{note}
\textbf{技巧 6：离散型求期望的”补集法”}

$P\{X\geq 1\}=1-P\{X=0\}$。对几何分布、泊松分布等，求”至少一次”的概率时用补集更快。

求期望时：若 $X$ 取非负整数值，$EX=\sum_{k=1}^\infty P\{X\geq k\}$（尾概率求和公式）。
\end{note}

\section*{D. 考前冲刺清单}
\addcontentsline{toc}{section}{D. 考前冲刺清单}

\begin{note}
以下每条都应能在 10 秒内回答。若卡住，立刻翻回对应章节复习。

\textbf{第 1 章（概率）}
\begin{enumerate}
    \item 全概率公式的结构？（分母展开为各原因的加权和）
    \item 贝叶斯公式 = 全概率公式的哪一项除以总和？
    \item $P(A\cup B)=P(A)+P(B)-P(AB)$
    \item 独立的定义？$P(AB)=P(A)P(B)$
    \item 条件概率 $P(B|A)=P(AB)/P(A)$
\end{enumerate}

\textbf{第 2 章（一维随机变量）}
\begin{enumerate}
    \item 泊松分布的 $EX$ 和 $DX$？（都等于 $\lambda$）
    \item 指数分布的无记忆性公式？$P\{X>s+t|X>s\}=P\{X>t\}$
    \item 标准化公式？$P\{X\leq a\}=\Phi\left(\frac{a-\mu}{\sigma}\right)$
    \item CDF 法求 $Y=g(X)$ 的密度：写CDF→积分→求导
    \item $F(x)$ 的三个性质？（单调不减、右连续、$F(-\infty)=0,F(+\infty)=1$）
\end{enumerate}

\textbf{第 3 章（二维随机变量）}
\begin{enumerate}
    \item 边缘密度怎么求？（对另一个变量积分）
    \item 独立性判据？$f(x,y)=f_X(x)f_Y(y)$ 恒成立
    \item 条件密度 = 联合密度 / 边缘密度
    \item $Z=X+Y$ 的卷积公式？$f_Z(z)=\int f_X(x)f_Y(z-x)\,\mathrm{d}x$
    \item 支撑域不是矩形 $\Rightarrow$ 几乎一定不独立
\end{enumerate}

\textbf{第 4 章（数字特征）}
\begin{enumerate}
    \item $DX=EX^2-(EX)^2$
    \item $D(aX+b)=a^2 DX$
    \item $\mathrm{Cov}(X,Y)=EXY-EX\cdot EY$
    \item 独立 $\Rightarrow$ 不相关；反之一般不成立（除非二维正态）
    \item $\rho_{XY}=\mathrm{Cov}(X,Y)/(\sqrt{DX}\sqrt{DY})$，$|\rho|\leq 1$
\end{enumerate}

\textbf{第 5 章（大数定律与CLT）}
\begin{enumerate}
    \item 切比雪夫不等式？$P\{|X-EX|\geq\varepsilon\}\leq DX/\varepsilon^2$
    \item 辛钦大数定律的条件？（iid + 期望存在）
    \item CLT 的结论？$\frac{\sum X_i - n\mu}{\sqrt{n}\sigma}\xrightarrow{d}N(0,1)$
    \item CLT 的使用场景？（$n$ 大时近似计算概率）
    \item 二项分布的正态近似？$B(n,p)\approx N(np,np(1-p))$
\end{enumerate}

\textbf{第 6 章（抽样分布）}
\begin{enumerate}
    \item $\bar X\sim N(\mu,\sigma^2/n)$
    \item $(n-1)S^2/\sigma^2\sim\chi^2(n-1)$
    \item $\bar X$ 与 $S^2$ 独立（仅正态总体）
    \item $t$ 分布 = 标准正态 / $\sqrt{\chi^2/n}$
    \item $F$ 分布 = 两个独立 $\chi^2$ 之比（各除以自由度）
\end{enumerate}

\textbf{第 7 章（点估计）}
\begin{enumerate}
    \item 矩估计核心？令 $EX=\bar X$（样本矩代替总体矩）
    \item MLE 步骤？写似然→取对数→求导（或判断单调）→解
    \item MLE 不变性？$\widehat{g(\theta)}=g(\hat\theta)$
    \item 无偏性定义？$E(\hat\theta)=\theta$
    \item $S^2$ 无偏估计 $\sigma^2$，但 $S$ 不是 $\sigma$ 的无偏估计
\end{enumerate}

\textbf{第 8 章（区间估计）}
\begin{enumerate}
    \item 估 $\mu$（$\sigma$ 已知）用什么？（$Z$ 分位数）
    \item 估 $\mu$（$\sigma$ 未知）用什么？（$t$ 分位数）
    \item 估 $\sigma^2$ 用什么？（$\chi^2$ 分位数）
    \item 置信区间越宽 $\Leftrightarrow$ 置信水平越高
    \item 增大 $n$ $\Rightarrow$ 区间变窄
\end{enumerate}

\textbf{第 9 章（假设检验）}
\begin{enumerate}
    \item $H_0$ 永远包含等号
    \item 拒绝域方向与 $H_1$ 一致
    \item $\alpha$ = 弃真概率（第一类错误）
    \item $\alpha\downarrow$ 则 $\beta\uparrow$（样本量固定时）
    \item 不拒绝 $H_0$ $\neq$ 证明 $H_0$ 为真
\end{enumerate}
\end{note}

\section*{E. 章节知识依赖图}
\addcontentsline{toc}{section}{E. 章节知识依赖图}

\begin{note}
\textbf{概念如何串联——从底层到顶层：}

\begin{verbatim}
Ch1 概率公理 ──→ Ch2 随机变量 ──→ Ch3 多维随机变量
    │                  │                    │
    │                  ↓                    ↓
    │              Ch4 数字特征 ←──── 协方差/相关系数
    │                  │
    ↓                  ↓
Ch5 大数定律/CLT ←── “均值收敛””和趋正态”
         │
         ↓
Ch6 抽样分布定理 ←── 正态总体 + X̄与S²独立
    │         │
    ↓         ↓
Ch7 点估计    Ch8 区间估计 ←→ Ch9 假设检验
(矩法/MLE)   (枢轴量反解)     (枢轴量代入H₀)
\end{verbatim}

\textbf{关键依赖链：}
\begin{itemize}
    \item Ch6 的三大定理（$\bar X$ 正态、$S^2$ 卡方、$t$ 统计量）是 Ch7--9 所有公式的\textbf{唯一来源}。
    \item Ch8 和 Ch9 用的数学工具完全相同（同一个枢轴量），只是”问法”不同：区间估计反解参数范围，假设检验代入 $\theta_0$ 看是否落入尾部。
    \item Ch4 的协方差/相关系数为 Ch3 的独立性判断提供了\textbf{必要条件}（不相关），但不是充分条件。
    \item Ch5 的 CLT 解释了为什么 Ch6 中 $\bar X$ 即使总体不正态，$n$ 大时也近似正态。
\end{itemize}

\textbf{学习顺序建议：} 严格按 Ch1→Ch2→Ch3→Ch4→Ch5→Ch6→Ch7→Ch8→Ch9 顺序学习。跳章会导致后续公式”不知道从哪来”。如果时间紧张，Ch5 可以快速过（只记结论），但 Ch6 绝不能跳——它是统计部分的地基。
\end{note}


\end{document}
